\documentclass[10pt,a4paper]{article}
\usepackage[a4paper,left=18mm,right=18mm,top=20mm,bottom=19mm,headheight=14pt]{geometry}
\usepackage[utf8]{inputenc}
\usepackage[T1]{fontenc}
\usepackage{lmodern}
\usepackage{xcolor}
\definecolor{ink}{HTML}{16394D}
\definecolor{teal}{HTML}{087D82}
\definecolor{muted}{HTML}{536D77}
\usepackage[most]{tcolorbox}
\usepackage{hyperref}
\hypersetup{colorlinks=true,linkcolor=teal,urlcolor=teal}
\usepackage{fancyhdr,needspace,enumitem,microtype}
\usepackage{titlesec}
\setlength{\parindent}{0pt}
\setlength{\parskip}{3.7pt}
\titleformat{\section}{\Large\bfseries\sffamily\color{ink}}{\thesection}{0.5em}{}
\titleformat{\subsection}{\normalsize\bfseries\sffamily\color{teal}}{}{0em}{}
\setlist[enumerate]{topsep=2pt,itemsep=1.5pt}
\pagestyle{fancy}
\fancyhf{}
\lhead{\footnotesize Kompyuter tizimlari va dasturiy muhit}
\rhead{\footnotesize ATTESTATSIYA 2026}
\cfoot{\footnotesize\thepage}
\renewcommand{\headrulewidth}{0.3pt}
\emergencystretch=3em
\sloppy
\begin{document}
\begin{titlepage}
\vspace*{27mm}
{\color{teal}\rule{\textwidth}{2pt}}\par
\vspace{21mm}
{\large\sffamily INFORMATIKA ATTESTATSIYASI \quad 2026}\par
\vspace{13mm}
{\fontsize{24}{30}\selectfont\bfseries\sffamily\color{ink}KOMPYUTER TIZIMLARI\\VA DASTURIY MUHIT}\par
\vspace{8mm}
{\large Izohli testlar kitobi}\par
\vspace{18mm}
{\Huge\bfseries\color{teal}609}\par
{\large\bfseries Test va uning tagidagi izohli javob}\par
\vspace{12mm}
67 ta published \qquad 542 ta review\par
606 ta Y1 \qquad 2 ta Y2 \qquad 1 ta Y3\par
\vfill
{\color{muted}2026 spetsifikatsiyasi 2-bo'lim: texnik va dasturiy ta'minot (4--5-savollar).\\
Microsoft Office keyingi mustaqil kitobda.}\par
\vspace{10mm}
{\color{teal}\rule{\textwidth}{2pt}}\par
\end{titlepage}
\section*{Kitobdan foydalanish}
\addcontentsline{toc}{section}{Kitobdan foydalanish}
Har bir testdan so'ng alohida rangli blokda to'g'ri javob va izoh beriladi.
Savollar bazaning \textit{published} yoki \textit{review} holatiga ko'ra belgilangan.
Review - tekshiruvdagi savol, mutaxassis tomonidan yakuniy tasdiqlangan degani emas.
\par\textbf{Manba:} attestatsiya100\% Supabase: \texttt{questions}, \texttt{question\_options},
\texttt{question\_keys}, \texttt{lessons}.
\par 609 ta amaldagi savol (67 published, 542 review). 200 ta arxivlangan savol kiritilmagan.
605 ta Y1 savolida 4 variant; bitta Y1 savolida 5 variant manbadagi holicha saqlanadi.
\par\textbf{Spetsifikatsiya:} katta 2-bo'lim 4--10-savollarga tegishli.
Bu kitob qurilmalar, operatsion tizimlar, fayllar va dasturlar (4--5) qismini qamraydi.
Word, Excel va PowerPoint (6--10) keyingi alohida kitobga ajratiladi.
\tableofcontents
\clearpage
\section{Kompyuter qurilmalari va ularning vazifalari}
\subsection{Nashr qilingan testlar}

% DB_ID 3b291efa-c5ea-5c09-9492-a10687870c44
\needspace{5\baselineskip}
{\bfseries\sffamily\color{ink} 1.} {scriptsizecolor{muted} Y1 | published | 3b291efa}
\par
Moslashtiring: 1) Nettop; 2) Tablet; 3) Laptop; 4) Monoblok. a) shaxsiy kompyuter; b) monitor va tizim bloki birlashtirilgan ixcham kompyuter; c) bloknot tipidagi shaxsiy kompyuter; d) planshetli shaxsiy kompyuter; e) ixcham qutidagi, internet ilovalari uchun mo'ljallangan kompyuter.
\begin{enumerate}[label=\Alph*),leftmargin=9mm]
\item 1-e, 2-d, 3-a, 4-b
\item 1-e, 2-d, 3-c, 4-b
\item 1-e, 2-a, 3-c, 4-b
\item 1-a, 2-d, 3-c, 4-b
\end{enumerate}
\begin{tcolorbox}[enhanced,breakable,colback=teal!4,colframe=teal!60!black,boxrule=0.45pt,arc=1.2mm,left=2.7mm,right=2.7mm,top=1.8mm,bottom=1.8mm,before skip=2mm,after skip=2mm]
{\bfseries\sffamily\color{teal}To'g'ri javob:} \textbf{B) 1-e, 2-d, 3-c, 4-b}
\par\smallskip
{\bfseries Izoh:} Nettop - e, tablet - d, laptop - c, monoblok - b.
\end{tcolorbox}
\vspace{2mm}\hrule height 0.2pt\vspace{3mm}

% DB_ID 41098c41-e0ff-43e0-b393-2b22bbb37006
\needspace{5\baselineskip}
{\bfseries\sffamily\color{ink} 2.} {scriptsizecolor{muted} Y1 | published | 41098c41}
\par
"Texnik ta'minot" atamasiga eng aniq ta'rifni toping.
\begin{enumerate}[label=\Alph*),leftmargin=9mm]
\item kompyuter bajaradigan dasturlar, operatsion tizimlar va algoritmik vositalar majmui
\item kompyuter tizimining jismoniy qurilmalari majmui
\item tizimni yaratadigan, boshqaradigan va undan foydalanadigan insonlar hamda ularning bilimlari
\item qayd etilgan belgi, qiymat yoki kuzatuv bo'lib, kontekstsiz mazmuni yetarli bo'lmasligi mumkin
\end{enumerate}
\begin{tcolorbox}[enhanced,breakable,colback=teal!4,colframe=teal!60!black,boxrule=0.45pt,arc=1.2mm,left=2.7mm,right=2.7mm,top=1.8mm,bottom=1.8mm,before skip=2mm,after skip=2mm]
{\bfseries\sffamily\color{teal}To'g'ri javob:} \textbf{B) kompyuter tizimining jismoniy qurilmalari majmui}
\par\smallskip
{\bfseries Izoh:} Texnik ta'minot - kompyuter tizimining jismoniy qurilmalari majmui.
\end{tcolorbox}
\vspace{2mm}\hrule height 0.2pt\vspace{3mm}

% DB_ID 45a17526-4c3a-5284-a7b9-4e140296006d
\needspace{5\baselineskip}
{\bfseries\sffamily\color{ink} 3.} {scriptsizecolor{muted} Y1 | published | 45a17526}
\par
Shaxsiy kompyuterning odatiy asosiy tarkibiy qismlarini tanlang: 1) monitor; 2) veb-kamera; 3) sichqoncha; 4) tizim bloki; 5) klaviatura; 6) skaner; 7) printer.
\begin{enumerate}[label=\Alph*),leftmargin=9mm]
\item 1, 4, 5
\item 1, 3, 5, 7
\item 2, 3, 6, 7
\item 1, 3, 4, 5
\item 1, 3, 6
\end{enumerate}
\begin{tcolorbox}[enhanced,breakable,colback=teal!4,colframe=teal!60!black,boxrule=0.45pt,arc=1.2mm,left=2.7mm,right=2.7mm,top=1.8mm,bottom=1.8mm,before skip=2mm,after skip=2mm]
{\bfseries\sffamily\color{teal}To'g'ri javob:} \textbf{D) 1, 3, 4, 5}
\par\smallskip
{\bfseries Izoh:} Odatiy asosiy to'plam monitor, tizim bloki, klaviatura va sichqonchadan iborat: 1, 3, 4, 5.
\end{tcolorbox}
\vspace{2mm}\hrule height 0.2pt\vspace{3mm}

% DB_ID 5271d862-5e87-444d-888b-beb24115b817
\needspace{5\baselineskip}
{\bfseries\sffamily\color{ink} 4.} {scriptsizecolor{muted} Y1 | published | 5271d862}
\par
Kompyuterda bir vaqtning o'zida bir nechta dastur ishlayapti. Elektr ta'minoti uzilganda odatda tarkibini yo'qotadigan tezkor xotira nimaga xizmat qiladi?
\begin{enumerate}[label=\Alph*),leftmargin=9mm]
\item Ishlayotgan dastur va ma'lumotni vaqtincha saqlaydi
\item Operatsion tizim yuklovchisini tok o'chganda ham saqlaydi
\item Foydalanuvchi fayllarini uzoq muddat diskda saqlaydi
\item Kompyuterning tarmoq manzilini avtomatik ravishda taqsimlaydi
\end{enumerate}
\begin{tcolorbox}[enhanced,breakable,colback=teal!4,colframe=teal!60!black,boxrule=0.45pt,arc=1.2mm,left=2.7mm,right=2.7mm,top=1.8mm,bottom=1.8mm,before skip=2mm,after skip=2mm]
{\bfseries\sffamily\color{teal}To'g'ri javob:} \textbf{A) Ishlayotgan dastur va ma'lumotni vaqtincha saqlaydi}
\par\smallskip
{\bfseries Izoh:} RAM ishlayotgan dasturlar va foydalanilayotgan ma'lumotlarni vaqtincha saqlaydi. Odatdagi RAM energiyaga bog'liq; uzoq muddatli saqlash uchun SSD/HDD kabi qurilmalar ishlatiladi.
\end{tcolorbox}
\vspace{2mm}\hrule height 0.2pt\vspace{3mm}

% DB_ID 6a9607dc-841f-4edb-b002-c89f06015e3a
\needspace{5\baselineskip}
{\bfseries\sffamily\color{ink} 5.} {scriptsizecolor{muted} Y1 | published | 6a9607dc}
\par
Quyidagilardan qaysi biri kiritish qurilmasi?
\begin{enumerate}[label=\Alph*),leftmargin=9mm]
\item Skaner
\item Monitor
\item Printer
\item Karnay
\end{enumerate}
\begin{tcolorbox}[enhanced,breakable,colback=teal!4,colframe=teal!60!black,boxrule=0.45pt,arc=1.2mm,left=2.7mm,right=2.7mm,top=1.8mm,bottom=1.8mm,before skip=2mm,after skip=2mm]
{\bfseries\sffamily\color{teal}To'g'ri javob:} \textbf{A) Skaner}
\par\smallskip
{\bfseries Izoh:} Skaner tashqi ma'lumotni kompyuterga kiritadi.
\end{tcolorbox}
\vspace{2mm}\hrule height 0.2pt\vspace{3mm}

% DB_ID 9428a109-834f-4c2f-905c-079c9b72f798
\needspace{5\baselineskip}
{\bfseries\sffamily\color{ink} 6.} {scriptsizecolor{muted} Y1 | published | 9428a109}
\par
Qurilmalarni tavsiflari bilan moslang: 1) planshet; 2) laptop; 3) smartfon; 4) desktop; 5) monoblok. A) ko'p funksiyali mobil telefon; B) portativ ko'ngilochar qurilma; C) monitor bilan bitta korpusda yig'ilgan kompyuter; D) ixcham portativ kompyuter; E) stolusti shaxsiy kompyuter.
\begin{enumerate}[label=\Alph*),leftmargin=9mm]
\item 1-A; 2-D; 3-B; 4-C; 5-E
\item 1-B; 2-C; 3-A; 4-E; 5-D
\item 1-B; 2-D; 3-A; 4-E; 5-C
\item 1-D; 2-B; 3-A; 4-E; 5-C
\end{enumerate}
\begin{tcolorbox}[enhanced,breakable,colback=teal!4,colframe=teal!60!black,boxrule=0.45pt,arc=1.2mm,left=2.7mm,right=2.7mm,top=1.8mm,bottom=1.8mm,before skip=2mm,after skip=2mm]
{\bfseries\sffamily\color{teal}To'g'ri javob:} \textbf{C) 1-B; 2-D; 3-A; 4-E; 5-C}
\par\smallskip
{\bfseries Izoh:} Planshet-B, laptop-D, smartfon-A, desktop-E, monoblok-C; moslik 1-B, 2-D, 3-A, 4-E, 5-C.
\end{tcolorbox}
\vspace{2mm}\hrule height 0.2pt\vspace{3mm}

% DB_ID d3f134de-7762-4246-9e8a-fab31cdc3dd2
\needspace{5\baselineskip}
{\bfseries\sffamily\color{ink} 7.} {scriptsizecolor{muted} Y1 | published | d3f134de}
\par
Klaviatura tugmalarini vazifalari bilan moslang. \newline 1. Caps Lock \newline 2. Ctrl+C \newline 3. Alt+F4 \newline 4. Enter \newline a. Yangi qatorga o'tish \newline b. Tanlangan matnni nusxalash \newline c. Doimiy katta harf rejimini yoqish/o'chirish \newline d. Dastur yoki oynani yopish
\begin{enumerate}[label=\Alph*),leftmargin=9mm]
\item 1-c, 2-b, 3-d, 4-a
\item 1-b, 2-a, 3-c, 4-d
\item 1-d, 2-c, 3-b, 4-a
\item 1-c, 2-d, 3-a, 4-b
\end{enumerate}
\begin{tcolorbox}[enhanced,breakable,colback=teal!4,colframe=teal!60!black,boxrule=0.45pt,arc=1.2mm,left=2.7mm,right=2.7mm,top=1.8mm,bottom=1.8mm,before skip=2mm,after skip=2mm]
{\bfseries\sffamily\color{teal}To'g'ri javob:} \textbf{A) 1-c, 2-b, 3-d, 4-a}
\par\smallskip
{\bfseries Izoh:} Caps Lock - katta harf rejimi; Ctrl+C - nusxalash; Alt+F4 - oynani yopish; Enter - yangi qator. [Yutuqli\_test, test kodi 12; Telegram xabarlar 6446 va 6556.]
\end{tcolorbox}
\vspace{2mm}\hrule height 0.2pt\vspace{3mm}

% DB_ID d7f0b639-ffbb-487d-bfc2-55c1e1604584
\needspace{5\baselineskip}
{\bfseries\sffamily\color{ink} 8.} {scriptsizecolor{muted} Y1 | published | d7f0b639}
\par
SSD va HDD qurilmalarining umumiy vazifasi qaysi?
\begin{enumerate}[label=\Alph*),leftmargin=9mm]
\item Uzoq muddatli ma'lumot saqlash
\item Matnni ekranga chiqarish
\item Ovoz yozish
\item Buyruqlarni bajarish
\end{enumerate}
\begin{tcolorbox}[enhanced,breakable,colback=teal!4,colframe=teal!60!black,boxrule=0.45pt,arc=1.2mm,left=2.7mm,right=2.7mm,top=1.8mm,bottom=1.8mm,before skip=2mm,after skip=2mm]
{\bfseries\sffamily\color{teal}To'g'ri javob:} \textbf{A) Uzoq muddatli ma'lumot saqlash}
\par\smallskip
{\bfseries Izoh:} SSD va HDD ikkalasi ham doimiy saqlash qurilmalaridir.
\end{tcolorbox}
\vspace{2mm}\hrule height 0.2pt\vspace{3mm}

% DB_ID e67b1532-e60e-4ade-8a21-74ede21819d6
\needspace{5\baselineskip}
{\bfseries\sffamily\color{ink} 9.} {scriptsizecolor{muted} Y1 | published | e67b1532}
\par
CPU, RAM, videokarta va boshqa modullar elektr hamda ma'lumot uzatish yo'llari orqali bir tizimga ulanadi. Asosiy plataning bosh vazifasi qaysi?
\begin{enumerate}[label=\Alph*),leftmargin=9mm]
\item Asosiy komponentlarni o'zaro bog'lash
\item Doimiy fayl xotirasini kengaytirish
\item Monitor ekranida grafikani hisoblash
\item Tarmoq paketlarini serverlarga marshrutlash
\end{enumerate}
\begin{tcolorbox}[enhanced,breakable,colback=teal!4,colframe=teal!60!black,boxrule=0.45pt,arc=1.2mm,left=2.7mm,right=2.7mm,top=1.8mm,bottom=1.8mm,before skip=2mm,after skip=2mm]
{\bfseries\sffamily\color{teal}To'g'ri javob:} \textbf{A) Asosiy komponentlarni o'zaro bog'lash}
\par\smallskip
{\bfseries Izoh:} Asosiy plata protsessor, tezkor xotira va boshqa komponentlar o'rtasidagi fizik hamda elektr aloqani tashkil qiladi. Grafikani GPU, paketlarni yo'naltirishni router bajaradi.
\end{tcolorbox}
\vspace{2mm}\hrule height 0.2pt\vspace{3mm}

% DB_ID f624fcf0-eb3f-4b1b-807b-e362afe010bf
\needspace{5\baselineskip}
{\bfseries\sffamily\color{ink} 10.} {scriptsizecolor{muted} Y1 | published | f624fcf0}
\par
Kompyuterda dastur buyruqlari bajarilib, sonlar ustida arifmetik va mantiqiy amallar amalga oshiriladi. Ushbu ishni bajaradigan asosiy komponentning vazifasi qaysi?
\begin{enumerate}[label=\Alph*),leftmargin=9mm]
\item Buyruqlarni bajarish va ma'lumotni qayta ishlash
\item Fayllarni elektr o'chganida ham uzoq muddat saqlash
\item Ekranga uzatiladigan tasvir signalini ranglarga aylantirish
\item Tarmoq orqali uzatiladigan paketlarni manziliga qarab yo'naltirish
\end{enumerate}
\begin{tcolorbox}[enhanced,breakable,colback=teal!4,colframe=teal!60!black,boxrule=0.45pt,arc=1.2mm,left=2.7mm,right=2.7mm,top=1.8mm,bottom=1.8mm,before skip=2mm,after skip=2mm]
{\bfseries\sffamily\color{teal}To'g'ri javob:} \textbf{A) Buyruqlarni bajarish va ma'lumotni qayta ishlash}
\par\smallskip
{\bfseries Izoh:} CPU dastur buyruqlarini bajaradi va ma'lumotni qayta ishlaydi. Doimiy saqlash SSD/HDD vazifasi; grafik tasvir bilan GPU va tarmoq marshrutlash bilan router shug'ullanadi.
\end{tcolorbox}
\vspace{2mm}\hrule height 0.2pt\vspace{3mm}

% DB_ID 0b34062b-f6d4-47e3-965b-b16208c58b81
\needspace{5\baselineskip}
{\bfseries\sffamily\color{ink} 11.} {scriptsizecolor{muted} Y1 | published | 0b34062b}
\par
Quyidagi ta'rif qaysi atamaga tegishli? Kompyuter tizimining jismoniy qurilmalari majmui.
\begin{enumerate}[label=\Alph*),leftmargin=9mm]
\item Dasturiy ta'minot
\item Insoniy ta'minot
\item Ma'lumot
\item Texnik ta'minot
\end{enumerate}
\begin{tcolorbox}[enhanced,breakable,colback=teal!4,colframe=teal!60!black,boxrule=0.45pt,arc=1.2mm,left=2.7mm,right=2.7mm,top=1.8mm,bottom=1.8mm,before skip=2mm,after skip=2mm]
{\bfseries\sffamily\color{teal}To'g'ri javob:} \textbf{D) Texnik ta'minot}
\par\smallskip
{\bfseries Izoh:} Ta'rif Texnik ta'minot tushunchasiga tegishli.
\end{tcolorbox}
\vspace{2mm}\hrule height 0.2pt\vspace{3mm}

% DB_ID 3413264a-e535-5cbd-9691-859b0dd172a1
\needspace{5\baselineskip}
{\bfseries\sffamily\color{ink} 12.} {scriptsizecolor{muted} Y1 | published | 3413264a}
\par
Hisoblash vositalarini darslikdagi rivojlanish davrlari bilan moslang: 1) Abak; 2) ENIAC; 3) Charlz Bebbijning bug' orqali boshqariladigan mashinasi; 4) Shikkard mashinasi. a) elektron hisoblash mashinalari; b) sun'iy intellektga asoslangan mashinalar; c) mexanik mashinalargacha; d) elektromexanik mashinalar; e) mexanik mashinalar.
\begin{enumerate}[label=\Alph*),leftmargin=9mm]
\item 1-c, 2-d, 3-a, 4-e
\item 1-b, 2-a, 3-d, 4-e
\item 1-c, 2-b, 3-a, 4-e
\item 1-c, 2-a, 3-d, 4-e
\end{enumerate}
\begin{tcolorbox}[enhanced,breakable,colback=teal!4,colframe=teal!60!black,boxrule=0.45pt,arc=1.2mm,left=2.7mm,right=2.7mm,top=1.8mm,bottom=1.8mm,before skip=2mm,after skip=2mm]
{\bfseries\sffamily\color{teal}To'g'ri javob:} \textbf{D) 1-c, 2-a, 3-d, 4-e}
\par\smallskip
{\bfseries Izoh:} 10-sinf Informatika va AT darsligidagi tasnif bo'yicha abak - mexanik mashinalargacha, Shikkard mashinasi - mexanik, Bebbijning 1822-yilgi bug' mashinasi - elektromexanik, ENIAC esa elektron davrga kiradi.
\end{tcolorbox}
\vspace{2mm}\hrule height 0.2pt\vspace{3mm}

% DB_ID dfe9f1f8-e616-4644-8670-9e3d75f0cb87
\needspace{5\baselineskip}
{\bfseries\sffamily\color{ink} 13.} {scriptsizecolor{muted} Y1 | published | dfe9f1f8}
\par
Sensorli planshet ekranida foydalanuvchi barmog'i bilan tugma bosadi, ayni paytda ekranda javob tasviri paydo bo'ladi. Ushbu ekran qaysi qurilma turiga kiradi?
\begin{enumerate}[label=\Alph*),leftmargin=9mm]
\item Ham kiritish, ham chiqarish
\item Faqat ma'lumotni kirituvchi qurilma
\item Faqat axborotni chiqaruvchi qurilma
\item Faqat ma'lumotni saqlovchi qurilma
\end{enumerate}
\begin{tcolorbox}[enhanced,breakable,colback=teal!4,colframe=teal!60!black,boxrule=0.45pt,arc=1.2mm,left=2.7mm,right=2.7mm,top=1.8mm,bottom=1.8mm,before skip=2mm,after skip=2mm]
{\bfseries\sffamily\color{teal}To'g'ri javob:} \textbf{A) Ham kiritish, ham chiqarish}
\par\smallskip
{\bfseries Izoh:} Sensorli ekran teginishni kiritish signali sifatida qabul qiladi va tasvirni chiqaradi. Shu sababli bir qurilmada kiritish va chiqarish funksiyalari mavjud.
\end{tcolorbox}
\vspace{2mm}\hrule height 0.2pt\vspace{3mm}

% DB_ID 00baa524-f77b-44b7-b950-6021830f5664
\needspace{5\baselineskip}
{\bfseries\sffamily\color{ink} 14.} {scriptsizecolor{muted} Y1 | published | 00baa524}
\par
Maktab arxivida 500 GB faylni elektr o'chirilgandan keyin ham saqlab qolish kerak. Qaysi qurilma eng mos?
\begin{enumerate}[label=\Alph*),leftmargin=9mm]
\item SSD yoki HDD
\item RAM
\item CPU registri
\item Monitor
\end{enumerate}
\begin{tcolorbox}[enhanced,breakable,colback=teal!4,colframe=teal!60!black,boxrule=0.45pt,arc=1.2mm,left=2.7mm,right=2.7mm,top=1.8mm,bottom=1.8mm,before skip=2mm,after skip=2mm]
{\bfseries\sffamily\color{teal}To'g'ri javob:} \textbf{A) SSD yoki HDD}
\par\smallskip
{\bfseries Izoh:} Uzoq muddatli saqlash uchun SSD/HDD kabi doimiy saqlash qurilmasi kerak.
\end{tcolorbox}
\vspace{2mm}\hrule height 0.2pt\vspace{3mm}

% DB_ID 07a2a5ac-624e-4cd8-8e31-bdd0267bb1ee
\needspace{5\baselineskip}
{\bfseries\sffamily\color{ink} 15.} {scriptsizecolor{muted} Y1 | published | 07a2a5ac}
\par
Qurilmalarni vazifalari bilan moslang: 1) klaviatura 2) printer 3) protsessor 4) asosiy plata. a) xotirada ma'lumot saqlaydi b) hisoblaydi va boshqaradi c) buyruqlarni kiritadi d) qog'ozga chop etadi e) qurilmalar ulanadigan slotlarga ega.
\begin{enumerate}[label=\Alph*),leftmargin=9mm]
\item 1-a,2-c,3-e,4-d
\item 1-d,2-e,3-b,4-c
\item 1-b,2-d,3-c,4-e
\item 1-c,2-d,3-b,4-e
\end{enumerate}
\begin{tcolorbox}[enhanced,breakable,colback=teal!4,colframe=teal!60!black,boxrule=0.45pt,arc=1.2mm,left=2.7mm,right=2.7mm,top=1.8mm,bottom=1.8mm,before skip=2mm,after skip=2mm]
{\bfseries\sffamily\color{teal}To'g'ri javob:} \textbf{D) 1-c,2-d,3-b,4-e}
\par\smallskip
{\bfseries Izoh:} Klaviatura c, printer d, CPU b, asosiy plata e.
\end{tcolorbox}
\vspace{2mm}\hrule height 0.2pt\vspace{3mm}

% DB_ID 0875bd23-8c26-4118-a949-a7557c4ee5d1
\needspace{5\baselineskip}
{\bfseries\sffamily\color{ink} 16.} {scriptsizecolor{muted} Y1 | published | 0875bd23}
\par
Onlayn darsda o'qituvchining ovozini kompyuterga kiritish uchun qaysi qurilma kerak?
\begin{enumerate}[label=\Alph*),leftmargin=9mm]
\item Mikrofon
\item Monitor
\item Printer
\item Videokarta
\end{enumerate}
\begin{tcolorbox}[enhanced,breakable,colback=teal!4,colframe=teal!60!black,boxrule=0.45pt,arc=1.2mm,left=2.7mm,right=2.7mm,top=1.8mm,bottom=1.8mm,before skip=2mm,after skip=2mm]
{\bfseries\sffamily\color{teal}To'g'ri javob:} \textbf{A) Mikrofon}
\par\smallskip
{\bfseries Izoh:} Mikrofon tovushni kiritish qurilmasi sifatida ishlaydi.
\end{tcolorbox}
\vspace{2mm}\hrule height 0.2pt\vspace{3mm}

% DB_ID 099a6a05-0c0f-49b7-bc40-466a8315e5a0
\needspace{5\baselineskip}
{\bfseries\sffamily\color{ink} 17.} {scriptsizecolor{muted} Y1 | published | 099a6a05}
\par
Grafik dasturda murakkab 3D tasvirlarni tez qayta ishlash talabi oshdi. Qaysi komponentni kuchaytirish eng bevosita yordam beradi?
\begin{enumerate}[label=\Alph*),leftmargin=9mm]
\item Grafik protsessor/videokarta
\item Klaviatura
\item Printer
\item Optik sichqoncha
\end{enumerate}
\begin{tcolorbox}[enhanced,breakable,colback=teal!4,colframe=teal!60!black,boxrule=0.45pt,arc=1.2mm,left=2.7mm,right=2.7mm,top=1.8mm,bottom=1.8mm,before skip=2mm,after skip=2mm]
{\bfseries\sffamily\color{teal}To'g'ri javob:} \textbf{A) Grafik protsessor/videokarta}
\par\smallskip
{\bfseries Izoh:} 3D grafik hisoblashlarni grafik protsessor tezlashtiradi.
\end{tcolorbox}
\vspace{2mm}\hrule height 0.2pt\vspace{3mm}

% DB_ID 1aaac57a-edf5-4a71-8386-5003dfb842a8
\needspace{5\baselineskip}
{\bfseries\sffamily\color{ink} 18.} {scriptsizecolor{muted} Y1 | published | 1aaac57a}
\par
Qog'ozdagi suratni kompyuterga raqamli ko'rinishda kiritish kerak. Qaysi qurilma tanlanadi?
\begin{enumerate}[label=\Alph*),leftmargin=9mm]
\item Skaner
\item Printer
\item Proyektor
\item Karnay
\end{enumerate}
\begin{tcolorbox}[enhanced,breakable,colback=teal!4,colframe=teal!60!black,boxrule=0.45pt,arc=1.2mm,left=2.7mm,right=2.7mm,top=1.8mm,bottom=1.8mm,before skip=2mm,after skip=2mm]
{\bfseries\sffamily\color{teal}To'g'ri javob:} \textbf{A) Skaner}
\par\smallskip
{\bfseries Izoh:} Skaner qog'ozdagi tasvirni raqamlashtirib kompyuterga kiritadi.
\end{tcolorbox}
\vspace{2mm}\hrule height 0.2pt\vspace{3mm}

% DB_ID 22db14f8-96df-4d72-ba61-2a3c9fe80b48
\needspace{5\baselineskip}
{\bfseries\sffamily\color{ink} 19.} {scriptsizecolor{muted} Y1 | published | 22db14f8}
\par
Tugmalarni vazifalar bilan moslang: 1) Delete 2) Back Space 3) Shift 4) Enter. a) bosh harfni yozish yoki son tugmasidagi belgini kiritish; b) faqat bosh harflarni yozish; c) kursordan o'ngdagi belgini o'chirish; d) amalni tasdiqlash; e) kursordan chapdagi belgini o'chirish.
\begin{enumerate}[label=\Alph*),leftmargin=9mm]
\item 1-e,2-d,3-a,4-b
\item 1-c,2-b,3-e,4-d
\item 1-c,2-e,3-a,4-d
\item 1-e,2-c,3-b,4-a
\end{enumerate}
\begin{tcolorbox}[enhanced,breakable,colback=teal!4,colframe=teal!60!black,boxrule=0.45pt,arc=1.2mm,left=2.7mm,right=2.7mm,top=1.8mm,bottom=1.8mm,before skip=2mm,after skip=2mm]
{\bfseries\sffamily\color{teal}To'g'ri javob:} \textbf{C) 1-c,2-e,3-a,4-d}
\par\smallskip
{\bfseries Izoh:} Delete o'ngdagi, Backspace chapdagi belgini o'chiradi; Shift registrni o'zgartiradi, Enter tasdiqlaydi. [2025 kuzgi demo YECHIMLARI, sahifa 2.]
\end{tcolorbox}
\vspace{2mm}\hrule height 0.2pt\vspace{3mm}

% DB_ID 3b4f4490-8292-424a-b785-416ed6479368
\needspace{5\baselineskip}
{\bfseries\sffamily\color{ink} 20.} {scriptsizecolor{muted} Y1 | published | 3b4f4490}
\par
Kompyuterning quyidagi qurilmalaridan qaysi biri markaziy protsessor (CPU) tarkibiga \textbf{kirmaydi}?
\begin{enumerate}[label=\Alph*),leftmargin=9mm]
\item Boshqarish qurilmasi
\item Operativ xotira
\item Registrlar
\item Arifmetik-mantiqiy qurilma
\end{enumerate}
\begin{tcolorbox}[enhanced,breakable,colback=teal!4,colframe=teal!60!black,boxrule=0.45pt,arc=1.2mm,left=2.7mm,right=2.7mm,top=1.8mm,bottom=1.8mm,before skip=2mm,after skip=2mm]
{\bfseries\sffamily\color{teal}To'g'ri javob:} \textbf{B) Operativ xotira}
\par\smallskip
{\bfseries Izoh:} Operativ xotira protsessordan tashqaridagi asosiy xotira hisoblanadi; boshqaruv qurilmasi, ALU va registrlar protsessor tarkibiga kiradi.
\end{tcolorbox}
\vspace{2mm}\hrule height 0.2pt\vspace{3mm}

% DB_ID 42a91902-8dfa-4efa-a4f3-4908fc47e833
\needspace{5\baselineskip}
{\bfseries\sffamily\color{ink} 21.} {scriptsizecolor{muted} Y1 | published | 42a91902}
\par
Maktab serveri, monitor, klaviatura va tarmoq qurilmalari ro'yxatga olindi. Bu vaziyatdagi asosiy tushunchani aniqlang.
\begin{enumerate}[label=\Alph*),leftmargin=9mm]
\item Dasturiy ta'minot
\item Texnik ta'minot
\item Insoniy ta'minot
\item Ma'lumot
\end{enumerate}
\begin{tcolorbox}[enhanced,breakable,colback=teal!4,colframe=teal!60!black,boxrule=0.45pt,arc=1.2mm,left=2.7mm,right=2.7mm,top=1.8mm,bottom=1.8mm,before skip=2mm,after skip=2mm]
{\bfseries\sffamily\color{teal}To'g'ri javob:} \textbf{B) Texnik ta'minot}
\par\smallskip
{\bfseries Izoh:} Vaziyat Texnik ta'minotga misol: kompyuter tizimining jismoniy qurilmalari majmui.
\end{tcolorbox}
\vspace{2mm}\hrule height 0.2pt\vspace{3mm}

% DB_ID 6de75396-6746-4d02-92d2-9cf57a77eee9
\needspace{5\baselineskip}
{\bfseries\sffamily\color{ink} 22.} {scriptsizecolor{muted} Y2 | published | 6de75396}
\par
Kompyuter qurilmalarini asosiy vazifalari bilan juftlang.
\textbf{Chap ustun:}
\begin{enumerate}[label=\arabic*),leftmargin=9mm]
\item CPU
\item RAM
\item SSD
\item Skaner
\end{enumerate}
\textbf{O'ng ustun:}
\begin{enumerate}[label=\Alph*),leftmargin=9mm]
\item Ma'lumotni uzoq muddat saqlash
\item Buyruqlarni bajarish va ma'lumotni qayta ishlash
\item Qog'ozdagi tasvirni kompyuterga kiritish
\item Ishlayotgan dastur va ma'lumotni vaqtincha saqlash
\end{enumerate}
\begin{tcolorbox}[enhanced,breakable,colback=teal!4,colframe=teal!60!black,boxrule=0.45pt,arc=1.2mm,left=2.7mm,right=2.7mm,top=1.8mm,bottom=1.8mm,before skip=2mm,after skip=2mm]
{\bfseries\sffamily\color{teal}To'g'ri javob:} \textbf{1-B, 2-D, 3-A, 4-C}
\par\smallskip
{\bfseries Izoh:} CPU - qayta ishlash; RAM - vaqtinchalik xotira; SSD - doimiy saqlash; skaner - kiritish qurilmasi.
\end{tcolorbox}
\vspace{2mm}\hrule height 0.2pt\vspace{3mm}

% DB_ID ce7d8782-dae5-4b90-bf9e-1aa73fdf56df
\needspace{5\baselineskip}
{\bfseries\sffamily\color{ink} 23.} {scriptsizecolor{muted} Y1 | published | ce7d8782}
\par
Video tahrirlash paytida ishlayotgan loyiha va dastur ma'lumotlarini tezkor saqlash uchun qaysi komponent eng bevosita ishlatiladi?
\begin{enumerate}[label=\Alph*),leftmargin=9mm]
\item RAM
\item Printer
\item Mikrofon
\item Skaner
\end{enumerate}
\begin{tcolorbox}[enhanced,breakable,colback=teal!4,colframe=teal!60!black,boxrule=0.45pt,arc=1.2mm,left=2.7mm,right=2.7mm,top=1.8mm,bottom=1.8mm,before skip=2mm,after skip=2mm]
{\bfseries\sffamily\color{teal}To'g'ri javob:} \textbf{A) RAM}
\par\smallskip
{\bfseries Izoh:} Ishlayotgan dastur va uning joriy ma'lumotlari RAMda vaqtincha saqlanadi.
\end{tcolorbox}
\vspace{2mm}\hrule height 0.2pt\vspace{3mm}

% DB_ID f6a02ff8-9b5d-4188-9eca-6d2fe628b0c5
\needspace{5\baselineskip}
{\bfseries\sffamily\color{ink} 24.} {scriptsizecolor{muted} Y1 | published | f6a02ff8}
\par
Bir qurilma foydalanuvchiga tasvir ko'rsatadi va barmoq tegishini ham qabul qiladi. Bu qurilma qanday tasniflanadi?
\begin{enumerate}[label=\Alph*),leftmargin=9mm]
\item Kiritish va chiqarish qurilmasi
\item Faqat saqlash qurilmasi
\item Faqat qayta ishlash qurilmasi
\item Faqat tarmoq qurilmasi
\end{enumerate}
\begin{tcolorbox}[enhanced,breakable,colback=teal!4,colframe=teal!60!black,boxrule=0.45pt,arc=1.2mm,left=2.7mm,right=2.7mm,top=1.8mm,bottom=1.8mm,before skip=2mm,after skip=2mm]
{\bfseries\sffamily\color{teal}To'g'ri javob:} \textbf{A) Kiritish va chiqarish qurilmasi}
\par\smallskip
{\bfseries Izoh:} Sensorli ekran bir vaqtning o'zida chiqarish va kiritish vazifasini bajaradi.
\end{tcolorbox}
\vspace{2mm}\hrule height 0.2pt\vspace{3mm}
\subsection{Tekshiruvdagi testlar}

% DB_ID 042cee6f-6793-23f3-b773-327de7192e4d
\needspace{5\baselineskip}
{\bfseries\sffamily\color{ink} 25.} {scriptsizecolor{muted} Y1 | review | 042cee6f}
\par
Chiqarish va ikki vazifali qurilmalar doirasida "aktuator" atamasining eng aniq ta'rifini toping.
\begin{enumerate}[label=\Alph*),leftmargin=9mm]
\item boshqaruv signalini fizik harakat yoki ta'sirga aylantiruvchi qurilma
\item ekranda yoki audio orqali vaqtinchalik taqdim etilgan natija
\item qog'oz yoki boshqa fizik muhitdagi natija
\item toner va elektrofotografik jarayon bilan tez chop etuvchi printer
\end{enumerate}
\begin{tcolorbox}[enhanced,breakable,colback=teal!4,colframe=teal!60!black,boxrule=0.45pt,arc=1.2mm,left=2.7mm,right=2.7mm,top=1.8mm,bottom=1.8mm,before skip=2mm,after skip=2mm]
{\bfseries\sffamily\color{teal}To'g'ri javob:} \textbf{A) boshqaruv signalini fizik harakat yoki ta'sirga aylantiruvchi qurilma}
\par\smallskip
{\bfseries Izoh:} boshqaruv signalini fizik harakat yoki ta'sirga aylantiruvchi qurilma. aktuator - boshqaruv signalini fizik harakat yoki ta'sirga aylantiruvchi qurilma. Manba: MAR, PDF 8-10; Qo'shimcha sahifa raqami (qaysi PDFga tegishli ekani aniqlanmagan; darslik sahifasi deb tasdiqlanmagan): 46
\end{tcolorbox}
\vspace{2mm}\hrule height 0.2pt\vspace{3mm}

% DB_ID 063a6168-4be9-2914-1556-df77247bd810
\needspace{5\baselineskip}
{\bfseries\sffamily\color{ink} 26.} {scriptsizecolor{muted} Y1 | review | 063a6168}
\par
Kiritish qurilmalari doirasida "OMR" atamasining eng aniq ta'rifini toping.
\begin{enumerate}[label=\Alph*),leftmargin=9mm]
\item tasvirdagi bosma matnni tahrirlanadigan belgilar sifatida tanish
\item magnit siyoh bilan bosilgan maxsus belgilarni o'qish
\item chiziqli koddagi identifikatorni avtomatik kiritish
\item oldindan belgilangan joylardagi belgilarni aniqlash
\end{enumerate}
\begin{tcolorbox}[enhanced,breakable,colback=teal!4,colframe=teal!60!black,boxrule=0.45pt,arc=1.2mm,left=2.7mm,right=2.7mm,top=1.8mm,bottom=1.8mm,before skip=2mm,after skip=2mm]
{\bfseries\sffamily\color{teal}To'g'ri javob:} \textbf{D) oldindan belgilangan joylardagi belgilarni aniqlash}
\par\smallskip
{\bfseries Izoh:} oldindan belgilangan joylardagi belgilarni aniqlash. OMR - oldindan belgilangan joylardagi belgilarni aniqlash. Manba: CAM10/11, PDF 37-45; Qo'shimcha sahifa raqami (qaysi PDFga tegishli ekani aniqlanmagan; darslik sahifasi deb tasdiqlanmagan): 37
\end{tcolorbox}
\vspace{2mm}\hrule height 0.2pt\vspace{3mm}

% DB_ID 0a21bdfc-4d4e-e568-8a6e-249b25b6a73e
\needspace{5\baselineskip}
{\bfseries\sffamily\color{ink} 27.} {scriptsizecolor{muted} Y1 | review | 0a21bdfc}
\par
Attestatsiyadagi nozik farqlar va integratsion masalalar doirasida "integratsion savol" atamasining eng aniq ta'rifini toping.
\begin{enumerate}[label=\Alph*),leftmargin=9mm]
\item obyektni muhim belgilari bo'yicha guruhga ajratish
\item bir nechta mavzuni bitta vaziyatda bog'laydigan topshiriq
\item variantlarni baholash uchun qo'llanadigan belgi
\item yechim bajarishi shart bo'lgan chegara yoki talab
\end{enumerate}
\begin{tcolorbox}[enhanced,breakable,colback=teal!4,colframe=teal!60!black,boxrule=0.45pt,arc=1.2mm,left=2.7mm,right=2.7mm,top=1.8mm,bottom=1.8mm,before skip=2mm,after skip=2mm]
{\bfseries\sffamily\color{teal}To'g'ri javob:} \textbf{B) bir nechta mavzuni bitta vaziyatda bog'laydigan topshiriq}
\par\smallskip
{\bfseries Izoh:} bir nechta mavzuni bitta vaziyatda bog'laydigan topshiriq. integratsion savol - bir nechta mavzuni bitta vaziyatda bog'laydigan topshiriq. Manba: ICT5, PDF 25-49; Qo'shimcha sahifa raqami (qaysi PDFga tegishli ekani aniqlanmagan; darslik sahifasi deb tasdiqlanmagan): 118
\end{tcolorbox}
\vspace{2mm}\hrule height 0.2pt\vspace{3mm}

% DB_ID 0f442919-4e84-f6e6-ce25-8bacc2bc8df4
\needspace{5\baselineskip}
{\bfseries\sffamily\color{ink} 28.} {scriptsizecolor{muted} Y1 | review | 0f442919}
\par
Chiqarish va ikki vazifali qurilmalar doirasida "haptik chiqish" atamasining eng aniq ta'rifini toping.
\begin{enumerate}[label=\Alph*),leftmargin=9mm]
\item ekranda yoki audio orqali vaqtinchalik taqdim etilgan natija
\item tebranish yoki kuch orqali beriladigan sezgi chiqishi
\item qog'oz yoki boshqa fizik muhitdagi natija
\item toner va elektrofotografik jarayon bilan tez chop etuvchi printer
\end{enumerate}
\begin{tcolorbox}[enhanced,breakable,colback=teal!4,colframe=teal!60!black,boxrule=0.45pt,arc=1.2mm,left=2.7mm,right=2.7mm,top=1.8mm,bottom=1.8mm,before skip=2mm,after skip=2mm]
{\bfseries\sffamily\color{teal}To'g'ri javob:} \textbf{B) tebranish yoki kuch orqali beriladigan sezgi chiqishi}
\par\smallskip
{\bfseries Izoh:} tebranish yoki kuch orqali beriladigan sezgi chiqishi. haptik chiqish - tebranish yoki kuch orqali beriladigan sezgi chiqishi. Manba: CAM9, PDF 16-32; Qo'shimcha sahifa raqami (qaysi PDFga tegishli ekani aniqlanmagan; darslik sahifasi deb tasdiqlanmagan): 47
\end{tcolorbox}
\vspace{2mm}\hrule height 0.2pt\vspace{3mm}

% DB_ID 110b3a44-e1ab-9ac6-66f8-32a09cd381e2
\needspace{5\baselineskip}
{\bfseries\sffamily\color{ink} 29.} {scriptsizecolor{muted} Y1 | review | 110b3a44}
\par
Kompyuter tizimi va uning ishlash modeli doirasida "umumiy maqsadli tizim" atamasining eng aniq ta'rifini toping.
\begin{enumerate}[label=\Alph*),leftmargin=9mm]
\item bir maqsad uchun hamkorlik qiluvchi texnik vosita, dastur, ma'lumot va foydalanuvchi majmuasi
\item turli dasturlar yordamida ko'p xil vazifani bajaradigan kompyuter tizimi
\item kompyuter tizimining fizik qismlari
\item kompyuterga vazifalarni bajarishni ko'rsatuvchi dasturlar
\end{enumerate}
\begin{tcolorbox}[enhanced,breakable,colback=teal!4,colframe=teal!60!black,boxrule=0.45pt,arc=1.2mm,left=2.7mm,right=2.7mm,top=1.8mm,bottom=1.8mm,before skip=2mm,after skip=2mm]
{\bfseries\sffamily\color{teal}To'g'ri javob:} \textbf{B) turli dasturlar yordamida ko'p xil vazifani bajaradigan kompyuter tizimi}
\par\smallskip
{\bfseries Izoh:} turli dasturlar yordamida ko'p xil vazifani bajaradigan kompyuter tizimi. umumiy maqsadli tizim - turli dasturlar yordamida ko'p xil vazifani bajaradigan kompyuter tizimi. Manba: MAR, PDF 4-10; Qo'shimcha sahifa raqami (qaysi PDFga tegishli ekani aniqlanmagan; darslik sahifasi deb tasdiqlanmagan): 8
\end{tcolorbox}
\vspace{2mm}\hrule height 0.2pt\vspace{3mm}

% DB_ID 16108f16-8233-45ba-40d1-19d35b3b2829
\needspace{5\baselineskip}
{\bfseries\sffamily\color{ink} 30.} {scriptsizecolor{muted} Y1 | review | 16108f16}
\par
Ikkilamchi saqlash qurilmalari doirasida "HDD" atamasining eng aniq ta'rifini toping.
\begin{enumerate}[label=\Alph*),leftmargin=9mm]
\item magnit plastinkali, katta sig'imli ikkilamchi xotira
\item flesh hujayralarga asoslangan harakatsiz ikkilamchi xotira
\item lazer orqali o'qiladigan va yoziladigan disk
\item ketma-ket murojaatli, arxiv uchun qulay magnit vosita
\end{enumerate}
\begin{tcolorbox}[enhanced,breakable,colback=teal!4,colframe=teal!60!black,boxrule=0.45pt,arc=1.2mm,left=2.7mm,right=2.7mm,top=1.8mm,bottom=1.8mm,before skip=2mm,after skip=2mm]
{\bfseries\sffamily\color{teal}To'g'ri javob:} \textbf{A) magnit plastinkali, katta sig'imli ikkilamchi xotira}
\par\smallskip
{\bfseries Izoh:} magnit plastinkali, katta sig'imli ikkilamchi xotira. HDD - magnit plastinkali, katta sig'imli ikkilamchi xotira. Manba: CAM9, PDF 33-47; Qo'shimcha sahifa raqami (qaysi PDFga tegishli ekani aniqlanmagan; darslik sahifasi deb tasdiqlanmagan): 30
\end{tcolorbox}
\vspace{2mm}\hrule height 0.2pt\vspace{3mm}

% DB_ID 16e3da95-94e5-e0f2-37d7-a8a635e5afb9
\needspace{5\baselineskip}
{\bfseries\sffamily\color{ink} 31.} {scriptsizecolor{muted} Y1 | review | 16e3da95}
\par
Chiqarish va ikki vazifali qurilmalar doirasida "lazer printer" atamasining eng aniq ta'rifini toping.
\begin{enumerate}[label=\Alph*),leftmargin=9mm]
\item ekranda yoki audio orqali vaqtinchalik taqdim etilgan natija
\item qog'oz yoki boshqa fizik muhitdagi natija
\item toner va elektrofotografik jarayon bilan tez chop etuvchi printer
\item mayda siyoh tomchilarini purkab tasvir hosil qiluvchi printer
\end{enumerate}
\begin{tcolorbox}[enhanced,breakable,colback=teal!4,colframe=teal!60!black,boxrule=0.45pt,arc=1.2mm,left=2.7mm,right=2.7mm,top=1.8mm,bottom=1.8mm,before skip=2mm,after skip=2mm]
{\bfseries\sffamily\color{teal}To'g'ri javob:} \textbf{C) toner va elektrofotografik jarayon bilan tez chop etuvchi printer}
\par\smallskip
{\bfseries Izoh:} toner va elektrofotografik jarayon bilan tez chop etuvchi printer. lazer printer - toner va elektrofotografik jarayon bilan tez chop etuvchi printer. Manba: ICT5, PDF 29-36; Qo'shimcha sahifa raqami (qaysi PDFga tegishli ekani aniqlanmagan; darslik sahifasi deb tasdiqlanmagan): 44
\end{tcolorbox}
\vspace{2mm}\hrule height 0.2pt\vspace{3mm}

% DB_ID 18ace082-e12b-17c6-81ff-04ec5d11e106
\needspace{5\baselineskip}
{\bfseries\sffamily\color{ink} 32.} {scriptsizecolor{muted} Y1 | review | 18ace082}
\par
Unumdorlik, nosozliklarni aniqlash va xavfsiz xizmat doirasida "profilaktika" atamasining eng aniq ta'rifini toping.
\begin{enumerate}[label=\Alph*),leftmargin=9mm]
\item tizim umumiy tezligini cheklayotgan eng sust resurs
\item nosozlik yuz berish ehtimolini oldindan kamaytiruvchi xizmat
\item RAM va ikkilamchi xotira o'rtasida sahifalarni almashish
\item qizish sabab protsessor tezligining avtomatik pasayishi
\end{enumerate}
\begin{tcolorbox}[enhanced,breakable,colback=teal!4,colframe=teal!60!black,boxrule=0.45pt,arc=1.2mm,left=2.7mm,right=2.7mm,top=1.8mm,bottom=1.8mm,before skip=2mm,after skip=2mm]
{\bfseries\sffamily\color{teal}To'g'ri javob:} \textbf{B) nosozlik yuz berish ehtimolini oldindan kamaytiruvchi xizmat}
\par\smallskip
{\bfseries Izoh:} nosozlik yuz berish ehtimolini oldindan kamaytiruvchi xizmat. profilaktika - nosozlik yuz berish ehtimolini oldindan kamaytiruvchi xizmat. Manba: CAM10/11, PDF 24-61; Qo'shimcha sahifa raqami (qaysi PDFga tegishli ekani aniqlanmagan; darslik sahifasi deb tasdiqlanmagan): 113
\end{tcolorbox}
\vspace{2mm}\hrule height 0.2pt\vspace{3mm}

% DB_ID 320ce098-4b7a-7045-947b-a2cecd37e430
\needspace{5\baselineskip}
{\bfseries\sffamily\color{ink} 33.} {scriptsizecolor{muted} Y1 | review | 320ce098}
\par
Unumdorlik, nosozliklarni aniqlash va xavfsiz xizmat doirasida "log" atamasining eng aniq ta'rifini toping.
\begin{enumerate}[label=\Alph*),leftmargin=9mm]
\item tizim hodisalari va xatolari qaydi
\item tizim umumiy tezligini cheklayotgan eng sust resurs
\item RAM va ikkilamchi xotira o'rtasida sahifalarni almashish
\item qizish sabab protsessor tezligining avtomatik pasayishi
\end{enumerate}
\begin{tcolorbox}[enhanced,breakable,colback=teal!4,colframe=teal!60!black,boxrule=0.45pt,arc=1.2mm,left=2.7mm,right=2.7mm,top=1.8mm,bottom=1.8mm,before skip=2mm,after skip=2mm]
{\bfseries\sffamily\color{teal}To'g'ri javob:} \textbf{A) tizim hodisalari va xatolari qaydi}
\par\smallskip
{\bfseries Izoh:} tizim hodisalari va xatolari qaydi. log - tizim hodisalari va xatolari qaydi. Manba: CAM9, PDF 12-47; Qo'shimcha sahifa raqami (qaysi PDFga tegishli ekani aniqlanmagan; darslik sahifasi deb tasdiqlanmagan): 112
\end{tcolorbox}
\vspace{2mm}\hrule height 0.2pt\vspace{3mm}

% DB_ID 3502c271-2feb-0b68-f577-5c7e878976f8
\needspace{5\baselineskip}
{\bfseries\sffamily\color{ink} 34.} {scriptsizecolor{muted} Y1 | review | 3502c271}
\par
Birlamchi xotira: registr, kesh, RAM va ROM doirasida "RAM" atamasining eng aniq ta'rifini toping.
\begin{enumerate}[label=\Alph*),leftmargin=9mm]
\item elektrsiz ham ma'lumotni saqlovchi o'zgarmas xotira
\item tez-tez kerak bo'ladigan ma'lumotni CPUga yaqin saqlovchi xotira
\item faol dastur va ma'lumotlar uchun o'zgaruvchan tezkor xotira
\item RAM yetishmaganda ikkilamchi xotiradan ajratilgan sekinroq maydon
\end{enumerate}
\begin{tcolorbox}[enhanced,breakable,colback=teal!4,colframe=teal!60!black,boxrule=0.45pt,arc=1.2mm,left=2.7mm,right=2.7mm,top=1.8mm,bottom=1.8mm,before skip=2mm,after skip=2mm]
{\bfseries\sffamily\color{teal}To'g'ri javob:} \textbf{C) faol dastur va ma'lumotlar uchun o'zgaruvchan tezkor xotira}
\par\smallskip
{\bfseries Izoh:} faol dastur va ma'lumotlar uchun o'zgaruvchan tezkor xotira. RAM - faol dastur va ma'lumotlar uchun o'zgaruvchan tezkor xotira. Manba: CAM9, PDF 11-14; Qo'shimcha sahifa raqami (qaysi PDFga tegishli ekani aniqlanmagan; darslik sahifasi deb tasdiqlanmagan): 23
\end{tcolorbox}
\vspace{2mm}\hrule height 0.2pt\vspace{3mm}

% DB_ID 36167f49-5af7-fa61-8950-5c0a7956aa06
\needspace{5\baselineskip}
{\bfseries\sffamily\color{ink} 35.} {scriptsizecolor{muted} Y1 | review | 36167f49}
\par
Birlamchi xotira: registr, kesh, RAM va ROM doirasida "firmware" atamasining eng aniq ta'rifini toping.
\begin{enumerate}[label=\Alph*),leftmargin=9mm]
\item faol dastur va ma'lumotlar uchun o'zgaruvchan tezkor xotira
\item elektrsiz ham ma'lumotni saqlovchi o'zgarmas xotira
\item qurilma ichiga o'rnatilgan past darajali boshqaruv dasturi
\item tez-tez kerak bo'ladigan ma'lumotni CPUga yaqin saqlovchi xotira
\end{enumerate}
\begin{tcolorbox}[enhanced,breakable,colback=teal!4,colframe=teal!60!black,boxrule=0.45pt,arc=1.2mm,left=2.7mm,right=2.7mm,top=1.8mm,bottom=1.8mm,before skip=2mm,after skip=2mm]
{\bfseries\sffamily\color{teal}To'g'ri javob:} \textbf{C) qurilma ichiga o'rnatilgan past darajali boshqaruv dasturi}
\par\smallskip
{\bfseries Izoh:} qurilma ichiga o'rnatilgan past darajali boshqaruv dasturi. firmware - qurilma ichiga o'rnatilgan past darajali boshqaruv dasturi. Manba: CAM9, PDF 11-14; Qo'shimcha sahifa raqami (qaysi PDFga tegishli ekani aniqlanmagan; darslik sahifasi deb tasdiqlanmagan): 27
\end{tcolorbox}
\vspace{2mm}\hrule height 0.2pt\vspace{3mm}

% DB_ID 3644c102-34f9-ab15-d4fc-16178fb01bd8
\needspace{5\baselineskip}
{\bfseries\sffamily\color{ink} 36.} {scriptsizecolor{muted} Y1 | review | 3644c102}
\par
Tizimli blok, bosh plata va ulanishlar doirasida "slot" atamasining eng aniq ta'rifini toping.
\begin{enumerate}[label=\Alph*),leftmargin=9mm]
\item ichki komponent o'rnatiladigan ulagich
\item ichki komponentlarni bog'lovchi asosiy bosma plata
\item ma'lumot va boshqaruv signallarini tashuvchi aloqa yo'llari
\item tashqi qurilma ulanadigan nuqta
\end{enumerate}
\begin{tcolorbox}[enhanced,breakable,colback=teal!4,colframe=teal!60!black,boxrule=0.45pt,arc=1.2mm,left=2.7mm,right=2.7mm,top=1.8mm,bottom=1.8mm,before skip=2mm,after skip=2mm]
{\bfseries\sffamily\color{teal}To'g'ri javob:} \textbf{A) ichki komponent o'rnatiladigan ulagich}
\par\smallskip
{\bfseries Izoh:} ichki komponent o'rnatiladigan ulagich. slot - ichki komponent o'rnatiladigan ulagich. Manba: ICT10, PDF 12-18; Qo'shimcha sahifa raqami (qaysi PDFga tegishli ekani aniqlanmagan; darslik sahifasi deb tasdiqlanmagan): 12
\end{tcolorbox}
\vspace{2mm}\hrule height 0.2pt\vspace{3mm}

% DB_ID 39bb77a5-3db1-51a7-7002-207eb7e21c2c
\needspace{5\baselineskip}
{\bfseries\sffamily\color{ink} 37.} {scriptsizecolor{muted} Y1 | review | 39bb77a5}
\par
Tizimli blok, bosh plata va ulanishlar doirasida "shina" atamasining eng aniq ta'rifini toping.
\begin{enumerate}[label=\Alph*),leftmargin=9mm]
\item ichki komponentlarni bog'lovchi asosiy bosma plata
\item ichki komponent o'rnatiladigan ulagich
\item tashqi qurilma ulanadigan nuqta
\item ma'lumot va boshqaruv signallarini tashuvchi aloqa yo'llari
\end{enumerate}
\begin{tcolorbox}[enhanced,breakable,colback=teal!4,colframe=teal!60!black,boxrule=0.45pt,arc=1.2mm,left=2.7mm,right=2.7mm,top=1.8mm,bottom=1.8mm,before skip=2mm,after skip=2mm]
{\bfseries\sffamily\color{teal}To'g'ri javob:} \textbf{D) ma'lumot va boshqaruv signallarini tashuvchi aloqa yo'llari}
\par\smallskip
{\bfseries Izoh:} ma'lumot va boshqaruv signallarini tashuvchi aloqa yo'llari. shina - ma'lumot va boshqaruv signallarini tashuvchi aloqa yo'llari. Manba: CAM10/11, PDF 32-36; Qo'shimcha sahifa raqami (qaysi PDFga tegishli ekani aniqlanmagan; darslik sahifasi deb tasdiqlanmagan): 11
\end{tcolorbox}
\vspace{2mm}\hrule height 0.2pt\vspace{3mm}

% DB_ID 3c87f919-73f6-8417-4153-cbc7b54e14aa
\needspace{5\baselineskip}
{\bfseries\sffamily\color{ink} 38.} {scriptsizecolor{muted} Y1 | review | 3c87f919}
\par
Kompyuter tizimi va uning ishlash modeli doirasida "dasturiy ta'minot" atamasining eng aniq ta'rifini toping.
\begin{enumerate}[label=\Alph*),leftmargin=9mm]
\item bir maqsad uchun hamkorlik qiluvchi texnik vosita, dastur, ma'lumot va foydalanuvchi majmuasi
\item kompyuter tizimining fizik qismlari
\item kompyuterga vazifalarni bajarishni ko'rsatuvchi dasturlar
\item kiritish, qayta ishlash, chiqarish va saqlash modeli
\end{enumerate}
\begin{tcolorbox}[enhanced,breakable,colback=teal!4,colframe=teal!60!black,boxrule=0.45pt,arc=1.2mm,left=2.7mm,right=2.7mm,top=1.8mm,bottom=1.8mm,before skip=2mm,after skip=2mm]
{\bfseries\sffamily\color{teal}To'g'ri javob:} \textbf{C) kompyuterga vazifalarni bajarishni ko'rsatuvchi dasturlar}
\par\smallskip
{\bfseries Izoh:} kompyuterga vazifalarni bajarishni ko'rsatuvchi dasturlar. dasturiy ta'minot - kompyuterga vazifalarni bajarishni ko'rsatuvchi dasturlar. Manba: ICT5, PDF 25-32; Qo'shimcha sahifa raqami (qaysi PDFga tegishli ekani aniqlanmagan; darslik sahifasi deb tasdiqlanmagan): 5
\end{tcolorbox}
\vspace{2mm}\hrule height 0.2pt\vspace{3mm}

% DB_ID 3edaa125-8f08-d462-0d5f-4dbbb5446b20
\needspace{5\baselineskip}
{\bfseries\sffamily\color{ink} 39.} {scriptsizecolor{muted} Y1 | review | 3edaa125}
\par
Kiritish qurilmalari doirasida "OCR" atamasining eng aniq ta'rifini toping.
\begin{enumerate}[label=\Alph*),leftmargin=9mm]
\item oldindan belgilangan joylardagi belgilarni aniqlash
\item magnit siyoh bilan bosilgan maxsus belgilarni o'qish
\item tasvirdagi bosma matnni tahrirlanadigan belgilar sifatida tanish
\item chiziqli koddagi identifikatorni avtomatik kiritish
\end{enumerate}
\begin{tcolorbox}[enhanced,breakable,colback=teal!4,colframe=teal!60!black,boxrule=0.45pt,arc=1.2mm,left=2.7mm,right=2.7mm,top=1.8mm,bottom=1.8mm,before skip=2mm,after skip=2mm]
{\bfseries\sffamily\color{teal}To'g'ri javob:} \textbf{C) tasvirdagi bosma matnni tahrirlanadigan belgilar sifatida tanish}
\par\smallskip
{\bfseries Izoh:} tasvirdagi bosma matnni tahrirlanadigan belgilar sifatida tanish. OCR - tasvirdagi bosma matnni tahrirlanadigan belgilar sifatida tanish. Manba: CAM9, PDF 16-32; Qo'shimcha sahifa raqami (qaysi PDFga tegishli ekani aniqlanmagan; darslik sahifasi deb tasdiqlanmagan): 36
\end{tcolorbox}
\vspace{2mm}\hrule height 0.2pt\vspace{3mm}

% DB_ID 418305e3-6ee3-978b-9d78-354865cc9cce
\needspace{5\baselineskip}
{\bfseries\sffamily\color{ink} 40.} {scriptsizecolor{muted} Y1 | review | 418305e3}
\par
Markaziy protsessor va bajarish sikli doirasida "CU" atamasining eng aniq ta'rifini toping.
\begin{enumerate}[label=\Alph*),leftmargin=9mm]
\item ko'rsatmalarni bajarib, ma'lumotlarni qayta ishlovchi markaziy qurilma
\item arifmetik va mantiqiy amallarni bajaruvchi qism
\item ko'rsatmalar oqimi va boshqa qismlar ishini boshqaruvchi qism
\item CPU ichidagi juda tez va kichik xotira
\end{enumerate}
\begin{tcolorbox}[enhanced,breakable,colback=teal!4,colframe=teal!60!black,boxrule=0.45pt,arc=1.2mm,left=2.7mm,right=2.7mm,top=1.8mm,bottom=1.8mm,before skip=2mm,after skip=2mm]
{\bfseries\sffamily\color{teal}To'g'ri javob:} \textbf{C) ko'rsatmalar oqimi va boshqa qismlar ishini boshqaruvchi qism}
\par\smallskip
{\bfseries Izoh:} ko'rsatmalar oqimi va boshqa qismlar ishini boshqaruvchi qism. CU - ko'rsatmalar oqimi va boshqa qismlar ishini boshqaruvchi qism. Manba: ICT10, PDF 17-20; Qo'shimcha sahifa raqami (qaysi PDFga tegishli ekani aniqlanmagan; darslik sahifasi deb tasdiqlanmagan): 19
\end{tcolorbox}
\vspace{2mm}\hrule height 0.2pt\vspace{3mm}

% DB_ID 43ac4680-7139-43e4-c595-f92b9b21b5eb
\needspace{5\baselineskip}
{\bfseries\sffamily\color{ink} 41.} {scriptsizecolor{muted} Y1 | review | 43ac4680}
\par
Ikkilamchi saqlash qurilmalari doirasida "magnit lenta" atamasining eng aniq ta'rifini toping.
\begin{enumerate}[label=\Alph*),leftmargin=9mm]
\item magnit plastinkali, katta sig'imli ikkilamchi xotira
\item flesh hujayralarga asoslangan harakatsiz ikkilamchi xotira
\item lazer orqali o'qiladigan va yoziladigan disk
\item ketma-ket murojaatli, arxiv uchun qulay magnit vosita
\end{enumerate}
\begin{tcolorbox}[enhanced,breakable,colback=teal!4,colframe=teal!60!black,boxrule=0.45pt,arc=1.2mm,left=2.7mm,right=2.7mm,top=1.8mm,bottom=1.8mm,before skip=2mm,after skip=2mm]
{\bfseries\sffamily\color{teal}To'g'ri javob:} \textbf{D) ketma-ket murojaatli, arxiv uchun qulay magnit vosita}
\par\smallskip
{\bfseries Izoh:} ketma-ket murojaatli, arxiv uchun qulay magnit vosita. magnit lenta - ketma-ket murojaatli, arxiv uchun qulay magnit vosita. Manba: ICT7, PDF 23-38; Qo'shimcha sahifa raqami (qaysi PDFga tegishli ekani aniqlanmagan; darslik sahifasi deb tasdiqlanmagan): 33
\end{tcolorbox}
\vspace{2mm}\hrule height 0.2pt\vspace{3mm}

% DB_ID 4a6ca2a6-4405-d2ae-0d7e-bd40301891f5
\needspace{5\baselineskip}
{\bfseries\sffamily\color{ink} 42.} {scriptsizecolor{muted} Y1 | review | 4a6ca2a6}
\par
Birlamchi xotira: registr, kesh, RAM va ROM doirasida "UEFI" atamasining eng aniq ta'rifini toping.
\begin{enumerate}[label=\Alph*),leftmargin=9mm]
\item faol dastur va ma'lumotlar uchun o'zgaruvchan tezkor xotira
\item elektrsiz ham ma'lumotni saqlovchi o'zgarmas xotira
\item tez-tez kerak bo'ladigan ma'lumotni CPUga yaqin saqlovchi xotira
\item qurilmani boshlang'ich sozlab, OT yuklanishini boshlovchi firmware muhiti
\end{enumerate}
\begin{tcolorbox}[enhanced,breakable,colback=teal!4,colframe=teal!60!black,boxrule=0.45pt,arc=1.2mm,left=2.7mm,right=2.7mm,top=1.8mm,bottom=1.8mm,before skip=2mm,after skip=2mm]
{\bfseries\sffamily\color{teal}To'g'ri javob:} \textbf{D) qurilmani boshlang'ich sozlab, OT yuklanishini boshlovchi firmware muhiti}
\par\smallskip
{\bfseries Izoh:} qurilmani boshlang'ich sozlab, OT yuklanishini boshlovchi firmware muhiti. UEFI - qurilmani boshlang'ich sozlab, OT yuklanishini boshlovchi firmware muhiti. Manba: CAM10/11, PDF 33-36; Qo'shimcha sahifa raqami (qaysi PDFga tegishli ekani aniqlanmagan; darslik sahifasi deb tasdiqlanmagan): 29
\end{tcolorbox}
\vspace{2mm}\hrule height 0.2pt\vspace{3mm}

% DB_ID 51714bd9-4ebe-3392-999f-f7a1b479d7e1
\needspace{5\baselineskip}
{\bfseries\sffamily\color{ink} 43.} {scriptsizecolor{muted} Y1 | review | 51714bd9}
\par
Unumdorlik, nosozliklarni aniqlash va xavfsiz xizmat doirasida "thermal throttling" atamasining eng aniq ta'rifini toping.
\begin{enumerate}[label=\Alph*),leftmargin=9mm]
\item tizim umumiy tezligini cheklayotgan eng sust resurs
\item RAM va ikkilamchi xotira o'rtasida sahifalarni almashish
\item qizish sabab protsessor tezligining avtomatik pasayishi
\item belgi va dalillar asosida sababni tizimli aniqlash
\end{enumerate}
\begin{tcolorbox}[enhanced,breakable,colback=teal!4,colframe=teal!60!black,boxrule=0.45pt,arc=1.2mm,left=2.7mm,right=2.7mm,top=1.8mm,bottom=1.8mm,before skip=2mm,after skip=2mm]
{\bfseries\sffamily\color{teal}To'g'ri javob:} \textbf{C) qizish sabab protsessor tezligining avtomatik pasayishi}
\par\smallskip
{\bfseries Izoh:} qizish sabab protsessor tezligining avtomatik pasayishi. thermal throttling - qizish sabab protsessor tezligining avtomatik pasayishi. Manba: ICT10, PDF 12-29; Qo'shimcha sahifa raqami (qaysi PDFga tegishli ekani aniqlanmagan; darslik sahifasi deb tasdiqlanmagan): 110
\end{tcolorbox}
\vspace{2mm}\hrule height 0.2pt\vspace{3mm}

% DB_ID 57c9c04f-ed4f-226b-8bf3-882c1a7d1280
\needspace{5\baselineskip}
{\bfseries\sffamily\color{ink} 44.} {scriptsizecolor{muted} Y1 | review | 57c9c04f}
\par
Unumdorlik, nosozliklarni aniqlash va xavfsiz xizmat doirasida "bottleneck" atamasining eng aniq ta'rifini toping.
\begin{enumerate}[label=\Alph*),leftmargin=9mm]
\item tizim umumiy tezligini cheklayotgan eng sust resurs
\item RAM va ikkilamchi xotira o'rtasida sahifalarni almashish
\item qizish sabab protsessor tezligining avtomatik pasayishi
\item belgi va dalillar asosida sababni tizimli aniqlash
\end{enumerate}
\begin{tcolorbox}[enhanced,breakable,colback=teal!4,colframe=teal!60!black,boxrule=0.45pt,arc=1.2mm,left=2.7mm,right=2.7mm,top=1.8mm,bottom=1.8mm,before skip=2mm,after skip=2mm]
{\bfseries\sffamily\color{teal}To'g'ri javob:} \textbf{A) tizim umumiy tezligini cheklayotgan eng sust resurs}
\par\smallskip
{\bfseries Izoh:} tizim umumiy tezligini cheklayotgan eng sust resurs. bottleneck - tizim umumiy tezligini cheklayotgan eng sust resurs. Manba: CAM9, PDF 12-47; Qo'shimcha sahifa raqami (qaysi PDFga tegishli ekani aniqlanmagan; darslik sahifasi deb tasdiqlanmagan): 108
\end{tcolorbox}
\vspace{2mm}\hrule height 0.2pt\vspace{3mm}

% DB_ID 5b4f75fe-97e8-8879-2ba2-2196b45250d0
\needspace{5\baselineskip}
{\bfseries\sffamily\color{ink} 45.} {scriptsizecolor{muted} Y1 | review | 5b4f75fe}
\par
Tizimli blok, bosh plata va ulanishlar doirasida "bosh plata" atamasining eng aniq ta'rifini toping.
\begin{enumerate}[label=\Alph*),leftmargin=9mm]
\item ma'lumot va boshqaruv signallarini tashuvchi aloqa yo'llari
\item ichki komponent o'rnatiladigan ulagich
\item ichki komponentlarni bog'lovchi asosiy bosma plata
\item tashqi qurilma ulanadigan nuqta
\end{enumerate}
\begin{tcolorbox}[enhanced,breakable,colback=teal!4,colframe=teal!60!black,boxrule=0.45pt,arc=1.2mm,left=2.7mm,right=2.7mm,top=1.8mm,bottom=1.8mm,before skip=2mm,after skip=2mm]
{\bfseries\sffamily\color{teal}To'g'ri javob:} \textbf{C) ichki komponentlarni bog'lovchi asosiy bosma plata}
\par\smallskip
{\bfseries Izoh:} ichki komponentlarni bog'lovchi asosiy bosma plata. bosh plata - ichki komponentlarni bog'lovchi asosiy bosma plata. Manba: CAM9, PDF 9-14; Qo'shimcha sahifa raqami (qaysi PDFga tegishli ekani aniqlanmagan; darslik sahifasi deb tasdiqlanmagan): 10
\end{tcolorbox}
\vspace{2mm}\hrule height 0.2pt\vspace{3mm}

% DB_ID 6feb5ab9-1705-77a6-91d0-5f87976d1a29
\needspace{5\baselineskip}
{\bfseries\sffamily\color{ink} 46.} {scriptsizecolor{muted} Y1 | review | 6feb5ab9}
\par
Chiqarish va ikki vazifali qurilmalar doirasida "siyohpurkagich" atamasining eng aniq ta'rifini toping.
\begin{enumerate}[label=\Alph*),leftmargin=9mm]
\item ekranda yoki audio orqali vaqtinchalik taqdim etilgan natija
\item qog'oz yoki boshqa fizik muhitdagi natija
\item toner va elektrofotografik jarayon bilan tez chop etuvchi printer
\item mayda siyoh tomchilarini purkab tasvir hosil qiluvchi printer
\end{enumerate}
\begin{tcolorbox}[enhanced,breakable,colback=teal!4,colframe=teal!60!black,boxrule=0.45pt,arc=1.2mm,left=2.7mm,right=2.7mm,top=1.8mm,bottom=1.8mm,before skip=2mm,after skip=2mm]
{\bfseries\sffamily\color{teal}To'g'ri javob:} \textbf{D) mayda siyoh tomchilarini purkab tasvir hosil qiluvchi printer}
\par\smallskip
{\bfseries Izoh:} mayda siyoh tomchilarini purkab tasvir hosil qiluvchi printer. siyohpurkagich - mayda siyoh tomchilarini purkab tasvir hosil qiluvchi printer. Manba: ICT10, PDF 9-13; Qo'shimcha sahifa raqami (qaysi PDFga tegishli ekani aniqlanmagan; darslik sahifasi deb tasdiqlanmagan): 45
\end{tcolorbox}
\vspace{2mm}\hrule height 0.2pt\vspace{3mm}

% DB_ID 73d7c4a0-61a3-5443-0b7c-19f1c11c1c3a
\needspace{5\baselineskip}
{\bfseries\sffamily\color{ink} 47.} {scriptsizecolor{muted} Y1 | review | 73d7c4a0}
\par
Kiritish qurilmalari doirasida "MICR" atamasining eng aniq ta'rifini toping.
\begin{enumerate}[label=\Alph*),leftmargin=9mm]
\item magnit siyoh bilan bosilgan maxsus belgilarni o'qish
\item tasvirdagi bosma matnni tahrirlanadigan belgilar sifatida tanish
\item oldindan belgilangan joylardagi belgilarni aniqlash
\item chiziqli koddagi identifikatorni avtomatik kiritish
\end{enumerate}
\begin{tcolorbox}[enhanced,breakable,colback=teal!4,colframe=teal!60!black,boxrule=0.45pt,arc=1.2mm,left=2.7mm,right=2.7mm,top=1.8mm,bottom=1.8mm,before skip=2mm,after skip=2mm]
{\bfseries\sffamily\color{teal}To'g'ri javob:} \textbf{A) magnit siyoh bilan bosilgan maxsus belgilarni o'qish}
\par\smallskip
{\bfseries Izoh:} magnit siyoh bilan bosilgan maxsus belgilarni o'qish. MICR - magnit siyoh bilan bosilgan maxsus belgilarni o'qish. Manba: CAM5, PDF 60-70; Qo'shimcha sahifa raqami (qaysi PDFga tegishli ekani aniqlanmagan; darslik sahifasi deb tasdiqlanmagan): 38
\end{tcolorbox}
\vspace{2mm}\hrule height 0.2pt\vspace{3mm}

% DB_ID 79598b21-6c6b-78f3-71da-4becd68ac243
\needspace{5\baselineskip}
{\bfseries\sffamily\color{ink} 48.} {scriptsizecolor{muted} Y1 | review | 79598b21}
\par
Kompyuter tizimi va uning ishlash modeli doirasida "texnik ta'minot" atamasining eng aniq ta'rifini toping.
\begin{enumerate}[label=\Alph*),leftmargin=9mm]
\item bir maqsad uchun hamkorlik qiluvchi texnik vosita, dastur, ma'lumot va foydalanuvchi majmuasi
\item kompyuter tizimining fizik qismlari
\item kompyuterga vazifalarni bajarishni ko'rsatuvchi dasturlar
\item kiritish, qayta ishlash, chiqarish va saqlash modeli
\end{enumerate}
\begin{tcolorbox}[enhanced,breakable,colback=teal!4,colframe=teal!60!black,boxrule=0.45pt,arc=1.2mm,left=2.7mm,right=2.7mm,top=1.8mm,bottom=1.8mm,before skip=2mm,after skip=2mm]
{\bfseries\sffamily\color{teal}To'g'ri javob:} \textbf{B) kompyuter tizimining fizik qismlari}
\par\smallskip
{\bfseries Izoh:} kompyuter tizimining fizik qismlari. texnik ta'minot - kompyuter tizimining fizik qismlari. Manba: CAM10/11, PDF 31-35; Qo'shimcha sahifa raqami (qaysi PDFga tegishli ekani aniqlanmagan; darslik sahifasi deb tasdiqlanmagan): 4
\end{tcolorbox}
\vspace{2mm}\hrule height 0.2pt\vspace{3mm}

% DB_ID 7df3bd4e-90e8-5038-e76d-6b65d8e09a52
\needspace{5\baselineskip}
{\bfseries\sffamily\color{ink} 49.} {scriptsizecolor{muted} Y1 | review | 7df3bd4e}
\par
Markaziy protsessor va bajarish sikli doirasida "ALU" atamasining eng aniq ta'rifini toping.
\begin{enumerate}[label=\Alph*),leftmargin=9mm]
\item ko'rsatmalarni bajarib, ma'lumotlarni qayta ishlovchi markaziy qurilma
\item arifmetik va mantiqiy amallarni bajaruvchi qism
\item ko'rsatmalar oqimi va boshqa qismlar ishini boshqaruvchi qism
\item CPU ichidagi juda tez va kichik xotira
\end{enumerate}
\begin{tcolorbox}[enhanced,breakable,colback=teal!4,colframe=teal!60!black,boxrule=0.45pt,arc=1.2mm,left=2.7mm,right=2.7mm,top=1.8mm,bottom=1.8mm,before skip=2mm,after skip=2mm]
{\bfseries\sffamily\color{teal}To'g'ri javob:} \textbf{B) arifmetik va mantiqiy amallarni bajaruvchi qism}
\par\smallskip
{\bfseries Izoh:} arifmetik va mantiqiy amallarni bajaruvchi qism. ALU - arifmetik va mantiqiy amallarni bajaruvchi qism. Manba: CAM10/11, PDF 32-34; Qo'shimcha sahifa raqami (qaysi PDFga tegishli ekani aniqlanmagan; darslik sahifasi deb tasdiqlanmagan): 18
\end{tcolorbox}
\vspace{2mm}\hrule height 0.2pt\vspace{3mm}

% DB_ID 84abce40-792b-c137-765f-fd00d66118ab
\needspace{5\baselineskip}
{\bfseries\sffamily\color{ink} 50.} {scriptsizecolor{muted} Y1 | review | 84abce40}
\par
Markaziy protsessor va bajarish sikli doirasida "takt chastotasi" atamasining eng aniq ta'rifini toping.
\begin{enumerate}[label=\Alph*),leftmargin=9mm]
\item sekunddagi takt impulslar soni
\item ko'rsatmalarni bajarib, ma'lumotlarni qayta ishlovchi markaziy qurilma
\item arifmetik va mantiqiy amallarni bajaruvchi qism
\item ko'rsatmalar oqimi va boshqa qismlar ishini boshqaruvchi qism
\end{enumerate}
\begin{tcolorbox}[enhanced,breakable,colback=teal!4,colframe=teal!60!black,boxrule=0.45pt,arc=1.2mm,left=2.7mm,right=2.7mm,top=1.8mm,bottom=1.8mm,before skip=2mm,after skip=2mm]
{\bfseries\sffamily\color{teal}To'g'ri javob:} \textbf{A) sekunddagi takt impulslar soni}
\par\smallskip
{\bfseries Izoh:} sekunddagi takt impulslar soni. takt chastotasi - sekunddagi takt impulslar soni. Manba: CAM9, PDF 10-14; Qo'shimcha sahifa raqami (qaysi PDFga tegishli ekani aniqlanmagan; darslik sahifasi deb tasdiqlanmagan): 21
\end{tcolorbox}
\vspace{2mm}\hrule height 0.2pt\vspace{3mm}

% DB_ID 88d1b97f-d0bc-3dcd-8e3a-babd912fc484
\needspace{5\baselineskip}
{\bfseries\sffamily\color{ink} 51.} {scriptsizecolor{muted} Y1 | review | 88d1b97f}
\par
Chiqarish va ikki vazifali qurilmalar doirasida "qattiq nusxa" atamasining eng aniq ta'rifini toping.
\begin{enumerate}[label=\Alph*),leftmargin=9mm]
\item ekranda yoki audio orqali vaqtinchalik taqdim etilgan natija
\item qog'oz yoki boshqa fizik muhitdagi natija
\item toner va elektrofotografik jarayon bilan tez chop etuvchi printer
\item mayda siyoh tomchilarini purkab tasvir hosil qiluvchi printer
\end{enumerate}
\begin{tcolorbox}[enhanced,breakable,colback=teal!4,colframe=teal!60!black,boxrule=0.45pt,arc=1.2mm,left=2.7mm,right=2.7mm,top=1.8mm,bottom=1.8mm,before skip=2mm,after skip=2mm]
{\bfseries\sffamily\color{teal}To'g'ri javob:} \textbf{B) qog'oz yoki boshqa fizik muhitdagi natija}
\par\smallskip
{\bfseries Izoh:} qog'oz yoki boshqa fizik muhitdagi natija. qattiq nusxa - qog'oz yoki boshqa fizik muhitdagi natija. Manba: CAM10/11, PDF 37-45; Qo'shimcha sahifa raqami (qaysi PDFga tegishli ekani aniqlanmagan; darslik sahifasi deb tasdiqlanmagan): 43
\end{tcolorbox}
\vspace{2mm}\hrule height 0.2pt\vspace{3mm}

% DB_ID 8b6e0fd5-6531-cbe8-3793-863e53c18953
\needspace{5\baselineskip}
{\bfseries\sffamily\color{ink} 52.} {scriptsizecolor{muted} Y1 | review | 8b6e0fd5}
\par
Tizimli blok, bosh plata va ulanishlar doirasida "quvvat bloki" atamasining eng aniq ta'rifini toping.
\begin{enumerate}[label=\Alph*),leftmargin=9mm]
\item ichki komponentlarni bog'lovchi asosiy bosma plata
\item ma'lumot va boshqaruv signallarini tashuvchi aloqa yo'llari
\item tarmoq elektrini komponentlarga mos doimiy tokka aylantiruvchi qism
\item ichki komponent o'rnatiladigan ulagich
\end{enumerate}
\begin{tcolorbox}[enhanced,breakable,colback=teal!4,colframe=teal!60!black,boxrule=0.45pt,arc=1.2mm,left=2.7mm,right=2.7mm,top=1.8mm,bottom=1.8mm,before skip=2mm,after skip=2mm]
{\bfseries\sffamily\color{teal}To'g'ri javob:} \textbf{C) tarmoq elektrini komponentlarga mos doimiy tokka aylantiruvchi qism}
\par\smallskip
{\bfseries Izoh:} tarmoq elektrini komponentlarga mos doimiy tokka aylantiruvchi qism. quvvat bloki - tarmoq elektrini komponentlarga mos doimiy tokka aylantiruvchi qism. Manba: CAM9, PDF 9-14; Qo'shimcha sahifa raqami (qaysi PDFga tegishli ekani aniqlanmagan; darslik sahifasi deb tasdiqlanmagan): 14
\end{tcolorbox}
\vspace{2mm}\hrule height 0.2pt\vspace{3mm}

% DB_ID 8c3d0f4c-cc6c-4748-9025-6399cd3dff14
\needspace{5\baselineskip}
{\bfseries\sffamily\color{ink} 53.} {scriptsizecolor{muted} Y1 | review | 8c3d0f4c}
\par
Birlamchi xotira: registr, kesh, RAM va ROM doirasida "kesh" atamasining eng aniq ta'rifini toping.
\begin{enumerate}[label=\Alph*),leftmargin=9mm]
\item tez-tez kerak bo'ladigan ma'lumotni CPUga yaqin saqlovchi xotira
\item faol dastur va ma'lumotlar uchun o'zgaruvchan tezkor xotira
\item elektrsiz ham ma'lumotni saqlovchi o'zgarmas xotira
\item RAM yetishmaganda ikkilamchi xotiradan ajratilgan sekinroq maydon
\end{enumerate}
\begin{tcolorbox}[enhanced,breakable,colback=teal!4,colframe=teal!60!black,boxrule=0.45pt,arc=1.2mm,left=2.7mm,right=2.7mm,top=1.8mm,bottom=1.8mm,before skip=2mm,after skip=2mm]
{\bfseries\sffamily\color{teal}To'g'ri javob:} \textbf{A) tez-tez kerak bo'ladigan ma'lumotni CPUga yaqin saqlovchi xotira}
\par\smallskip
{\bfseries Izoh:} tez-tez kerak bo'ladigan ma'lumotni CPUga yaqin saqlovchi xotira. kesh - tez-tez kerak bo'ladigan ma'lumotni CPUga yaqin saqlovchi xotira. Manba: ICT10, PDF 15-20; Qo'shimcha sahifa raqami (qaysi PDFga tegishli ekani aniqlanmagan; darslik sahifasi deb tasdiqlanmagan): 25
\end{tcolorbox}
\vspace{2mm}\hrule height 0.2pt\vspace{3mm}

% DB_ID 8dc05c1e-0f97-5b7e-9c6e-569c256c986a
\needspace{5\baselineskip}
{\bfseries\sffamily\color{ink} 54.} {scriptsizecolor{muted} Y1 | review | 8dc05c1e}
\par
Attestatsiyadagi nozik farqlar va integratsion masalalar doirasida "distraktor" atamasining eng aniq ta'rifini toping.
\begin{enumerate}[label=\Alph*),leftmargin=9mm]
\item obyektni muhim belgilari bo'yicha guruhga ajratish
\item variantlarni baholash uchun qo'llanadigan belgi
\item ishonarli, ammo noto'g'ri javob varianti
\item yechim bajarishi shart bo'lgan chegara yoki talab
\end{enumerate}
\begin{tcolorbox}[enhanced,breakable,colback=teal!4,colframe=teal!60!black,boxrule=0.45pt,arc=1.2mm,left=2.7mm,right=2.7mm,top=1.8mm,bottom=1.8mm,before skip=2mm,after skip=2mm]
{\bfseries\sffamily\color{teal}To'g'ri javob:} \textbf{C) ishonarli, ammo noto'g'ri javob varianti}
\par\smallskip
{\bfseries Izoh:} ishonarli, ammo noto'g'ri javob varianti. distraktor - ishonarli, ammo noto'g'ri javob varianti. Manba: MAR, PDF 4-20; Qo'shimcha sahifa raqami (qaysi PDFga tegishli ekani aniqlanmagan; darslik sahifasi deb tasdiqlanmagan): 119
\end{tcolorbox}
\vspace{2mm}\hrule height 0.2pt\vspace{3mm}

% DB_ID 8e47d714-286f-764d-2794-420a0427ab2c
\needspace{5\baselineskip}
{\bfseries\sffamily\color{ink} 55.} {scriptsizecolor{muted} Y1 | review | 8e47d714}
\par
Tizimli blok, bosh plata va ulanishlar doirasida "kengaytirish kartasi" atamasining eng aniq ta'rifini toping.
\begin{enumerate}[label=\Alph*),leftmargin=9mm]
\item ichki komponentlarni bog'lovchi asosiy bosma plata
\item ma'lumot va boshqaruv signallarini tashuvchi aloqa yo'llari
\item ichki komponent o'rnatiladigan ulagich
\item kompyuterga qo'shimcha funksiya berish uchun slotga o'rnatiladigan plata
\end{enumerate}
\begin{tcolorbox}[enhanced,breakable,colback=teal!4,colframe=teal!60!black,boxrule=0.45pt,arc=1.2mm,left=2.7mm,right=2.7mm,top=1.8mm,bottom=1.8mm,before skip=2mm,after skip=2mm]
{\bfseries\sffamily\color{teal}To'g'ri javob:} \textbf{D) kompyuterga qo'shimcha funksiya berish uchun slotga o'rnatiladigan plata}
\par\smallskip
{\bfseries Izoh:} kompyuterga qo'shimcha funksiya berish uchun slotga o'rnatiladigan plata. kengaytirish kartasi - kompyuterga qo'shimcha funksiya berish uchun slotga o'rnatiladigan plata. Manba: CAM10/11, PDF 32-36; Qo'shimcha sahifa raqami (qaysi PDFga tegishli ekani aniqlanmagan; darslik sahifasi deb tasdiqlanmagan): 15
\end{tcolorbox}
\vspace{2mm}\hrule height 0.2pt\vspace{3mm}

% DB_ID 8ed51b7b-f3fa-678c-c40e-58264c2853b4
\needspace{5\baselineskip}
{\bfseries\sffamily\color{ink} 56.} {scriptsizecolor{muted} Y1 | review | 8ed51b7b}
\par
Chiqarish va ikki vazifali qurilmalar doirasida "yumshoq nusxa" atamasining eng aniq ta'rifini toping.
\begin{enumerate}[label=\Alph*),leftmargin=9mm]
\item ekranda yoki audio orqali vaqtinchalik taqdim etilgan natija
\item qog'oz yoki boshqa fizik muhitdagi natija
\item toner va elektrofotografik jarayon bilan tez chop etuvchi printer
\item mayda siyoh tomchilarini purkab tasvir hosil qiluvchi printer
\end{enumerate}
\begin{tcolorbox}[enhanced,breakable,colback=teal!4,colframe=teal!60!black,boxrule=0.45pt,arc=1.2mm,left=2.7mm,right=2.7mm,top=1.8mm,bottom=1.8mm,before skip=2mm,after skip=2mm]
{\bfseries\sffamily\color{teal}To'g'ri javob:} \textbf{A) ekranda yoki audio orqali vaqtinchalik taqdim etilgan natija}
\par\smallskip
{\bfseries Izoh:} ekranda yoki audio orqali vaqtinchalik taqdim etilgan natija. yumshoq nusxa - ekranda yoki audio orqali vaqtinchalik taqdim etilgan natija. Manba: CAM9, PDF 16-32; Qo'shimcha sahifa raqami (qaysi PDFga tegishli ekani aniqlanmagan; darslik sahifasi deb tasdiqlanmagan): 42
\end{tcolorbox}
\vspace{2mm}\hrule height 0.2pt\vspace{3mm}

% DB_ID 9c1dbd2d-b59a-c6f2-21df-55f98b9fdb79
\needspace{5\baselineskip}
{\bfseries\sffamily\color{ink} 57.} {scriptsizecolor{muted} Y1 | review | 9c1dbd2d}
\par
Ikkilamchi saqlash qurilmalari doirasida "USB flesh" atamasining eng aniq ta'rifini toping.
\begin{enumerate}[label=\Alph*),leftmargin=9mm]
\item portativ solid-state saqlash qurilmasi
\item magnit plastinkali, katta sig'imli ikkilamchi xotira
\item flesh hujayralarga asoslangan harakatsiz ikkilamchi xotira
\item lazer orqali o'qiladigan va yoziladigan disk
\end{enumerate}
\begin{tcolorbox}[enhanced,breakable,colback=teal!4,colframe=teal!60!black,boxrule=0.45pt,arc=1.2mm,left=2.7mm,right=2.7mm,top=1.8mm,bottom=1.8mm,before skip=2mm,after skip=2mm]
{\bfseries\sffamily\color{teal}To'g'ri javob:} \textbf{A) portativ solid-state saqlash qurilmasi}
\par\smallskip
{\bfseries Izoh:} portativ solid-state saqlash qurilmasi. USB flesh - portativ solid-state saqlash qurilmasi. Manba: MAR, PDF 7-9; Qo'shimcha sahifa raqami (qaysi PDFga tegishli ekani aniqlanmagan; darslik sahifasi deb tasdiqlanmagan): 34
\end{tcolorbox}
\vspace{2mm}\hrule height 0.2pt\vspace{3mm}

% DB_ID 9d3c65dd-9d83-9a83-4caf-90c4064f58da
\needspace{5\baselineskip}
{\bfseries\sffamily\color{ink} 58.} {scriptsizecolor{muted} Y1 | review | 9d3c65dd}
\par
Markaziy protsessor va bajarish sikli doirasida "CPU" atamasining eng aniq ta'rifini toping.
\begin{enumerate}[label=\Alph*),leftmargin=9mm]
\item ko'rsatmalarni bajarib, ma'lumotlarni qayta ishlovchi markaziy qurilma
\item arifmetik va mantiqiy amallarni bajaruvchi qism
\item ko'rsatmalar oqimi va boshqa qismlar ishini boshqaruvchi qism
\item CPU ichidagi juda tez va kichik xotira
\end{enumerate}
\begin{tcolorbox}[enhanced,breakable,colback=teal!4,colframe=teal!60!black,boxrule=0.45pt,arc=1.2mm,left=2.7mm,right=2.7mm,top=1.8mm,bottom=1.8mm,before skip=2mm,after skip=2mm]
{\bfseries\sffamily\color{teal}To'g'ri javob:} \textbf{A) ko'rsatmalarni bajarib, ma'lumotlarni qayta ishlovchi markaziy qurilma}
\par\smallskip
{\bfseries Izoh:} ko'rsatmalarni bajarib, ma'lumotlarni qayta ishlovchi markaziy qurilma. CPU - ko'rsatmalarni bajarib, ma'lumotlarni qayta ishlovchi markaziy qurilma. Manba: CAM9, PDF 10-14; Qo'shimcha sahifa raqami (qaysi PDFga tegishli ekani aniqlanmagan; darslik sahifasi deb tasdiqlanmagan): 17
\end{tcolorbox}
\vspace{2mm}\hrule height 0.2pt\vspace{3mm}

% DB_ID a2330198-124c-0cfa-26a3-d0aa33ceeb73
\needspace{5\baselineskip}
{\bfseries\sffamily\color{ink} 59.} {scriptsizecolor{muted} Y1 | review | a2330198}
\par
Markaziy protsessor va bajarish sikli doirasida "kesh" atamasining eng aniq ta'rifini toping.
\begin{enumerate}[label=\Alph*),leftmargin=9mm]
\item ko'rsatmalarni bajarib, ma'lumotlarni qayta ishlovchi markaziy qurilma
\item CPU va RAM o'rtasidagi tezkor oraliq xotira
\item arifmetik va mantiqiy amallarni bajaruvchi qism
\item ko'rsatmalar oqimi va boshqa qismlar ishini boshqaruvchi qism
\end{enumerate}
\begin{tcolorbox}[enhanced,breakable,colback=teal!4,colframe=teal!60!black,boxrule=0.45pt,arc=1.2mm,left=2.7mm,right=2.7mm,top=1.8mm,bottom=1.8mm,before skip=2mm,after skip=2mm]
{\bfseries\sffamily\color{teal}To'g'ri javob:} \textbf{B) CPU va RAM o'rtasidagi tezkor oraliq xotira}
\par\smallskip
{\bfseries Izoh:} CPU va RAM o'rtasidagi tezkor oraliq xotira. kesh - CPU va RAM o'rtasidagi tezkor oraliq xotira. Manba: CAM10/11, PDF 32-34; Qo'shimcha sahifa raqami (qaysi PDFga tegishli ekani aniqlanmagan; darslik sahifasi deb tasdiqlanmagan): 22
\end{tcolorbox}
\vspace{2mm}\hrule height 0.2pt\vspace{3mm}

% DB_ID a97a7de1-a010-d7df-4acc-dba0bdc2f228
\needspace{5\baselineskip}
{\bfseries\sffamily\color{ink} 60.} {scriptsizecolor{muted} Y1 | review | a97a7de1}
\par
Attestatsiyadagi nozik farqlar va integratsion masalalar doirasida "asoslash" atamasining eng aniq ta'rifini toping.
\begin{enumerate}[label=\Alph*),leftmargin=9mm]
\item obyektni muhim belgilari bo'yicha guruhga ajratish
\item variantlarni baholash uchun qo'llanadigan belgi
\item yechim bajarishi shart bo'lgan chegara yoki talab
\item tanlovni vaziyatdagi dalil bilan bog'lash
\end{enumerate}
\begin{tcolorbox}[enhanced,breakable,colback=teal!4,colframe=teal!60!black,boxrule=0.45pt,arc=1.2mm,left=2.7mm,right=2.7mm,top=1.8mm,bottom=1.8mm,before skip=2mm,after skip=2mm]
{\bfseries\sffamily\color{teal}To'g'ri javob:} \textbf{D) tanlovni vaziyatdagi dalil bilan bog'lash}
\par\smallskip
{\bfseries Izoh:} tanlovni vaziyatdagi dalil bilan bog'lash. asoslash - tanlovni vaziyatdagi dalil bilan bog'lash. Manba: SPEC, PDF 3; Qo'shimcha sahifa raqami (qaysi PDFga tegishli ekani aniqlanmagan; darslik sahifasi deb tasdiqlanmagan): 120
\end{tcolorbox}
\vspace{2mm}\hrule height 0.2pt\vspace{3mm}

% DB_ID aa8229b6-7d5f-aeff-f5c5-a50ffb522007
\needspace{5\baselineskip}
{\bfseries\sffamily\color{ink} 61.} {scriptsizecolor{muted} Y1 | review | aa8229b6}
\par
Attestatsiyadagi nozik farqlar va integratsion masalalar doirasida "tasnif" atamasining eng aniq ta'rifini toping.
\begin{enumerate}[label=\Alph*),leftmargin=9mm]
\item variantlarni baholash uchun qo'llanadigan belgi
\item yechim bajarishi shart bo'lgan chegara yoki talab
\item obyektni muhim belgilari bo'yicha guruhga ajratish
\item bir nechta mavzuni bitta vaziyatda bog'laydigan topshiriq
\end{enumerate}
\begin{tcolorbox}[enhanced,breakable,colback=teal!4,colframe=teal!60!black,boxrule=0.45pt,arc=1.2mm,left=2.7mm,right=2.7mm,top=1.8mm,bottom=1.8mm,before skip=2mm,after skip=2mm]
{\bfseries\sffamily\color{teal}To'g'ri javob:} \textbf{C) obyektni muhim belgilari bo'yicha guruhga ajratish}
\par\smallskip
{\bfseries Izoh:} obyektni muhim belgilari bo'yicha guruhga ajratish. tasnif - obyektni muhim belgilari bo'yicha guruhga ajratish. Manba: SPEC, PDF 3; Qo'shimcha sahifa raqami (qaysi PDFga tegishli ekani aniqlanmagan; darslik sahifasi deb tasdiqlanmagan): 115
\end{tcolorbox}
\vspace{2mm}\hrule height 0.2pt\vspace{3mm}

% DB_ID bb1d2004-c474-cc9d-959d-c0288053fcc2
\needspace{5\baselineskip}
{\bfseries\sffamily\color{ink} 62.} {scriptsizecolor{muted} Y1 | review | bb1d2004}
\par
Birlamchi xotira: registr, kesh, RAM va ROM doirasida "virtual xotira" atamasining eng aniq ta'rifini toping.
\begin{enumerate}[label=\Alph*),leftmargin=9mm]
\item faol dastur va ma'lumotlar uchun o'zgaruvchan tezkor xotira
\item RAM yetishmaganda ikkilamchi xotiradan ajratilgan sekinroq maydon
\item elektrsiz ham ma'lumotni saqlovchi o'zgarmas xotira
\item tez-tez kerak bo'ladigan ma'lumotni CPUga yaqin saqlovchi xotira
\end{enumerate}
\begin{tcolorbox}[enhanced,breakable,colback=teal!4,colframe=teal!60!black,boxrule=0.45pt,arc=1.2mm,left=2.7mm,right=2.7mm,top=1.8mm,bottom=1.8mm,before skip=2mm,after skip=2mm]
{\bfseries\sffamily\color{teal}To'g'ri javob:} \textbf{B) RAM yetishmaganda ikkilamchi xotiradan ajratilgan sekinroq maydon}
\par\smallskip
{\bfseries Izoh:} RAM yetishmaganda ikkilamchi xotiradan ajratilgan sekinroq maydon. virtual xotira - RAM yetishmaganda ikkilamchi xotiradan ajratilgan sekinroq maydon. Manba: MAR, PDF 7-8; Qo'shimcha sahifa raqami (qaysi PDFga tegishli ekani aniqlanmagan; darslik sahifasi deb tasdiqlanmagan): 26
\end{tcolorbox}
\vspace{2mm}\hrule height 0.2pt\vspace{3mm}

% DB_ID bbdcd1a7-15f2-947a-4459-afc509b40f59
\needspace{5\baselineskip}
{\bfseries\sffamily\color{ink} 63.} {scriptsizecolor{muted} Y1 | review | bbdcd1a7}
\par
Kiritish qurilmalari doirasida "shtrix-kod o'quvchi" atamasining eng aniq ta'rifini toping.
\begin{enumerate}[label=\Alph*),leftmargin=9mm]
\item tasvirdagi bosma matnni tahrirlanadigan belgilar sifatida tanish
\item chiziqli koddagi identifikatorni avtomatik kiritish
\item oldindan belgilangan joylardagi belgilarni aniqlash
\item magnit siyoh bilan bosilgan maxsus belgilarni o'qish
\end{enumerate}
\begin{tcolorbox}[enhanced,breakable,colback=teal!4,colframe=teal!60!black,boxrule=0.45pt,arc=1.2mm,left=2.7mm,right=2.7mm,top=1.8mm,bottom=1.8mm,before skip=2mm,after skip=2mm]
{\bfseries\sffamily\color{teal}To'g'ri javob:} \textbf{B) chiziqli koddagi identifikatorni avtomatik kiritish}
\par\smallskip
{\bfseries Izoh:} chiziqli koddagi identifikatorni avtomatik kiritish. shtrix-kod o'quvchi - chiziqli koddagi identifikatorni avtomatik kiritish. Manba: ICT5, PDF 28-34; Qo'shimcha sahifa raqami (qaysi PDFga tegishli ekani aniqlanmagan; darslik sahifasi deb tasdiqlanmagan): 39
\end{tcolorbox}
\vspace{2mm}\hrule height 0.2pt\vspace{3mm}

% DB_ID bf005cda-18cc-b191-19cf-908f7ed6a5f6
\needspace{5\baselineskip}
{\bfseries\sffamily\color{ink} 64.} {scriptsizecolor{muted} Y1 | review | bf005cda}
\par
Birlamchi xotira: registr, kesh, RAM va ROM doirasida "ROM" atamasining eng aniq ta'rifini toping.
\begin{enumerate}[label=\Alph*),leftmargin=9mm]
\item faol dastur va ma'lumotlar uchun o'zgaruvchan tezkor xotira
\item tez-tez kerak bo'ladigan ma'lumotni CPUga yaqin saqlovchi xotira
\item RAM yetishmaganda ikkilamchi xotiradan ajratilgan sekinroq maydon
\item elektrsiz ham ma'lumotni saqlovchi o'zgarmas xotira
\end{enumerate}
\begin{tcolorbox}[enhanced,breakable,colback=teal!4,colframe=teal!60!black,boxrule=0.45pt,arc=1.2mm,left=2.7mm,right=2.7mm,top=1.8mm,bottom=1.8mm,before skip=2mm,after skip=2mm]
{\bfseries\sffamily\color{teal}To'g'ri javob:} \textbf{D) elektrsiz ham ma'lumotni saqlovchi o'zgarmas xotira}
\par\smallskip
{\bfseries Izoh:} elektrsiz ham ma'lumotni saqlovchi o'zgarmas xotira. ROM - elektrsiz ham ma'lumotni saqlovchi o'zgarmas xotira. Manba: CAM10/11, PDF 33-36; Qo'shimcha sahifa raqami (qaysi PDFga tegishli ekani aniqlanmagan; darslik sahifasi deb tasdiqlanmagan): 24
\end{tcolorbox}
\vspace{2mm}\hrule height 0.2pt\vspace{3mm}

% DB_ID c37909d3-bdb6-a59a-1770-f191dac0cd69
\needspace{5\baselineskip}
{\bfseries\sffamily\color{ink} 65.} {scriptsizecolor{muted} Y1 | review | c37909d3}
\par
Ikkilamchi saqlash qurilmalari doirasida "SSD" atamasining eng aniq ta'rifini toping.
\begin{enumerate}[label=\Alph*),leftmargin=9mm]
\item magnit plastinkali, katta sig'imli ikkilamchi xotira
\item flesh hujayralarga asoslangan harakatsiz ikkilamchi xotira
\item lazer orqali o'qiladigan va yoziladigan disk
\item ketma-ket murojaatli, arxiv uchun qulay magnit vosita
\end{enumerate}
\begin{tcolorbox}[enhanced,breakable,colback=teal!4,colframe=teal!60!black,boxrule=0.45pt,arc=1.2mm,left=2.7mm,right=2.7mm,top=1.8mm,bottom=1.8mm,before skip=2mm,after skip=2mm]
{\bfseries\sffamily\color{teal}To'g'ri javob:} \textbf{B) flesh hujayralarga asoslangan harakatsiz ikkilamchi xotira}
\par\smallskip
{\bfseries Izoh:} flesh hujayralarga asoslangan harakatsiz ikkilamchi xotira. SSD - flesh hujayralarga asoslangan harakatsiz ikkilamchi xotira. Manba: CAM10/11, PDF 34-37; Qo'shimcha sahifa raqami (qaysi PDFga tegishli ekani aniqlanmagan; darslik sahifasi deb tasdiqlanmagan): 31
\end{tcolorbox}
\vspace{2mm}\hrule height 0.2pt\vspace{3mm}

% DB_ID c8a44969-6f19-6d65-e3d7-864e04117ad1
\needspace{5\baselineskip}
{\bfseries\sffamily\color{ink} 66.} {scriptsizecolor{muted} Y1 | review | c8a44969}
\par
Kompyuter tizimi va uning ishlash modeli doirasida "IPOS" atamasining eng aniq ta'rifini toping.
\begin{enumerate}[label=\Alph*),leftmargin=9mm]
\item bir maqsad uchun hamkorlik qiluvchi texnik vosita, dastur, ma'lumot va foydalanuvchi majmuasi
\item kompyuter tizimining fizik qismlari
\item kompyuterga vazifalarni bajarishni ko'rsatuvchi dasturlar
\item kiritish, qayta ishlash, chiqarish va saqlash modeli
\end{enumerate}
\begin{tcolorbox}[enhanced,breakable,colback=teal!4,colframe=teal!60!black,boxrule=0.45pt,arc=1.2mm,left=2.7mm,right=2.7mm,top=1.8mm,bottom=1.8mm,before skip=2mm,after skip=2mm]
{\bfseries\sffamily\color{teal}To'g'ri javob:} \textbf{D) kiritish, qayta ishlash, chiqarish va saqlash modeli}
\par\smallskip
{\bfseries Izoh:} kiritish, qayta ishlash, chiqarish va saqlash modeli. IPOS - kiritish, qayta ishlash, chiqarish va saqlash modeli. Manba: ICT10, PDF 6-12; Qo'shimcha sahifa raqami (qaysi PDFga tegishli ekani aniqlanmagan; darslik sahifasi deb tasdiqlanmagan): 6
\end{tcolorbox}
\vspace{2mm}\hrule height 0.2pt\vspace{3mm}

% DB_ID cbe7b3cd-b2ad-2da9-6f05-34c3b367a183
\needspace{5\baselineskip}
{\bfseries\sffamily\color{ink} 67.} {scriptsizecolor{muted} Y1 | review | cbe7b3cd}
\par
Unumdorlik, nosozliklarni aniqlash va xavfsiz xizmat doirasida "paging" atamasining eng aniq ta'rifini toping.
\begin{enumerate}[label=\Alph*),leftmargin=9mm]
\item tizim umumiy tezligini cheklayotgan eng sust resurs
\item RAM va ikkilamchi xotira o'rtasida sahifalarni almashish
\item qizish sabab protsessor tezligining avtomatik pasayishi
\item belgi va dalillar asosida sababni tizimli aniqlash
\end{enumerate}
\begin{tcolorbox}[enhanced,breakable,colback=teal!4,colframe=teal!60!black,boxrule=0.45pt,arc=1.2mm,left=2.7mm,right=2.7mm,top=1.8mm,bottom=1.8mm,before skip=2mm,after skip=2mm]
{\bfseries\sffamily\color{teal}To'g'ri javob:} \textbf{B) RAM va ikkilamchi xotira o'rtasida sahifalarni almashish}
\par\smallskip
{\bfseries Izoh:} RAM va ikkilamchi xotira o'rtasida sahifalarni almashish. paging - RAM va ikkilamchi xotira o'rtasida sahifalarni almashish. Manba: CAM10/11, PDF 24-61; Qo'shimcha sahifa raqami (qaysi PDFga tegishli ekani aniqlanmagan; darslik sahifasi deb tasdiqlanmagan): 109
\end{tcolorbox}
\vspace{2mm}\hrule height 0.2pt\vspace{3mm}

% DB_ID d1d3053f-702e-f0e6-5163-9870f41325b6
\needspace{5\baselineskip}
{\bfseries\sffamily\color{ink} 68.} {scriptsizecolor{muted} Y1 | review | d1d3053f}
\par
Ikkilamchi saqlash qurilmalari doirasida "bulutli saqlash" atamasining eng aniq ta'rifini toping.
\begin{enumerate}[label=\Alph*),leftmargin=9mm]
\item magnit plastinkali, katta sig'imli ikkilamchi xotira
\item tarmoq orqali masofadagi server xotirasidan foydalanish xizmati
\item flesh hujayralarga asoslangan harakatsiz ikkilamchi xotira
\item lazer orqali o'qiladigan va yoziladigan disk
\end{enumerate}
\begin{tcolorbox}[enhanced,breakable,colback=teal!4,colframe=teal!60!black,boxrule=0.45pt,arc=1.2mm,left=2.7mm,right=2.7mm,top=1.8mm,bottom=1.8mm,before skip=2mm,after skip=2mm]
{\bfseries\sffamily\color{teal}To'g'ri javob:} \textbf{B) tarmoq orqali masofadagi server xotirasidan foydalanish xizmati}
\par\smallskip
{\bfseries Izoh:} tarmoq orqali masofadagi server xotirasidan foydalanish xizmati. bulutli saqlash - tarmoq orqali masofadagi server xotirasidan foydalanish xizmati. Manba: CAM9, PDF 33-47; Qo'shimcha sahifa raqami (qaysi PDFga tegishli ekani aniqlanmagan; darslik sahifasi deb tasdiqlanmagan): 35
\end{tcolorbox}
\vspace{2mm}\hrule height 0.2pt\vspace{3mm}

% DB_ID da1b375f-fe6b-65cb-8cfb-08f406fb1546
\needspace{5\baselineskip}
{\bfseries\sffamily\color{ink} 69.} {scriptsizecolor{muted} Y1 | review | da1b375f}
\par
Attestatsiyadagi nozik farqlar va integratsion masalalar doirasida "mezon" atamasining eng aniq ta'rifini toping.
\begin{enumerate}[label=\Alph*),leftmargin=9mm]
\item obyektni muhim belgilari bo'yicha guruhga ajratish
\item yechim bajarishi shart bo'lgan chegara yoki talab
\item bir nechta mavzuni bitta vaziyatda bog'laydigan topshiriq
\item variantlarni baholash uchun qo'llanadigan belgi
\end{enumerate}
\begin{tcolorbox}[enhanced,breakable,colback=teal!4,colframe=teal!60!black,boxrule=0.45pt,arc=1.2mm,left=2.7mm,right=2.7mm,top=1.8mm,bottom=1.8mm,before skip=2mm,after skip=2mm]
{\bfseries\sffamily\color{teal}To'g'ri javob:} \textbf{D) variantlarni baholash uchun qo'llanadigan belgi}
\par\smallskip
{\bfseries Izoh:} variantlarni baholash uchun qo'llanadigan belgi. mezon - variantlarni baholash uchun qo'llanadigan belgi. Manba: CAM9, PDF 8-47; Qo'shimcha sahifa raqami (qaysi PDFga tegishli ekani aniqlanmagan; darslik sahifasi deb tasdiqlanmagan): 116
\end{tcolorbox}
\vspace{2mm}\hrule height 0.2pt\vspace{3mm}

% DB_ID dcadcb93-bdc4-8792-265f-7d2483cc6b79
\needspace{5\baselineskip}
{\bfseries\sffamily\color{ink} 70.} {scriptsizecolor{muted} Y1 | review | dcadcb93}
\par
Markaziy protsessor va bajarish sikli doirasida "registr" atamasining eng aniq ta'rifini toping.
\begin{enumerate}[label=\Alph*),leftmargin=9mm]
\item ko'rsatmalarni bajarib, ma'lumotlarni qayta ishlovchi markaziy qurilma
\item arifmetik va mantiqiy amallarni bajaruvchi qism
\item ko'rsatmalar oqimi va boshqa qismlar ishini boshqaruvchi qism
\item CPU ichidagi juda tez va kichik xotira
\end{enumerate}
\begin{tcolorbox}[enhanced,breakable,colback=teal!4,colframe=teal!60!black,boxrule=0.45pt,arc=1.2mm,left=2.7mm,right=2.7mm,top=1.8mm,bottom=1.8mm,before skip=2mm,after skip=2mm]
{\bfseries\sffamily\color{teal}To'g'ri javob:} \textbf{D) CPU ichidagi juda tez va kichik xotira}
\par\smallskip
{\bfseries Izoh:} CPU ichidagi juda tez va kichik xotira. registr - CPU ichidagi juda tez va kichik xotira. Manba: MAR, PDF 6-7; Qo'shimcha sahifa raqami (qaysi PDFga tegishli ekani aniqlanmagan; darslik sahifasi deb tasdiqlanmagan): 20
\end{tcolorbox}
\vspace{2mm}\hrule height 0.2pt\vspace{3mm}

% DB_ID e5f3abf8-5d6f-92fa-96e3-ddd140243e77
\needspace{5\baselineskip}
{\bfseries\sffamily\color{ink} 71.} {scriptsizecolor{muted} Y1 | review | e5f3abf8}
\par
Unumdorlik, nosozliklarni aniqlash va xavfsiz xizmat doirasida "diagnostika" atamasining eng aniq ta'rifini toping.
\begin{enumerate}[label=\Alph*),leftmargin=9mm]
\item tizim umumiy tezligini cheklayotgan eng sust resurs
\item RAM va ikkilamchi xotira o'rtasida sahifalarni almashish
\item qizish sabab protsessor tezligining avtomatik pasayishi
\item belgi va dalillar asosida sababni tizimli aniqlash
\end{enumerate}
\begin{tcolorbox}[enhanced,breakable,colback=teal!4,colframe=teal!60!black,boxrule=0.45pt,arc=1.2mm,left=2.7mm,right=2.7mm,top=1.8mm,bottom=1.8mm,before skip=2mm,after skip=2mm]
{\bfseries\sffamily\color{teal}To'g'ri javob:} \textbf{D) belgi va dalillar asosida sababni tizimli aniqlash}
\par\smallskip
{\bfseries Izoh:} belgi va dalillar asosida sababni tizimli aniqlash. diagnostika - belgi va dalillar asosida sababni tizimli aniqlash. Manba: MAR, PDF 10-20; Qo'shimcha sahifa raqami (qaysi PDFga tegishli ekani aniqlanmagan; darslik sahifasi deb tasdiqlanmagan): 111
\end{tcolorbox}
\vspace{2mm}\hrule height 0.2pt\vspace{3mm}

% DB_ID e65e8a6f-57b5-8c51-97a7-d0c1b40716f2
\needspace{5\baselineskip}
{\bfseries\sffamily\color{ink} 72.} {scriptsizecolor{muted} Y1 | review | e65e8a6f}
\par
Kompyuter tizimi va uning ishlash modeli doirasida "o'rnatilgan tizim" atamasining eng aniq ta'rifini toping.
\begin{enumerate}[label=\Alph*),leftmargin=9mm]
\item kattaroq qurilma ichida muayyan vazifani bajaruvchi kompyuter tizimi
\item bir maqsad uchun hamkorlik qiluvchi texnik vosita, dastur, ma'lumot va foydalanuvchi majmuasi
\item kompyuter tizimining fizik qismlari
\item kompyuterga vazifalarni bajarishni ko'rsatuvchi dasturlar
\end{enumerate}
\begin{tcolorbox}[enhanced,breakable,colback=teal!4,colframe=teal!60!black,boxrule=0.45pt,arc=1.2mm,left=2.7mm,right=2.7mm,top=1.8mm,bottom=1.8mm,before skip=2mm,after skip=2mm]
{\bfseries\sffamily\color{teal}To'g'ri javob:} \textbf{A) kattaroq qurilma ichida muayyan vazifani bajaruvchi kompyuter tizimi}
\par\smallskip
{\bfseries Izoh:} kattaroq qurilma ichida muayyan vazifani bajaruvchi kompyuter tizimi. o'rnatilgan tizim - kattaroq qurilma ichida muayyan vazifani bajaruvchi kompyuter tizimi. Manba: ICT9, PDF 14; Qo'shimcha sahifa raqami (qaysi PDFga tegishli ekani aniqlanmagan; darslik sahifasi deb tasdiqlanmagan): 7
\end{tcolorbox}
\vspace{2mm}\hrule height 0.2pt\vspace{3mm}

% DB_ID eae1ef66-8216-7d15-4b06-e108d68aec40
\needspace{5\baselineskip}
{\bfseries\sffamily\color{ink} 73.} {scriptsizecolor{muted} Y1 | review | eae1ef66}
\par
Kiritish qurilmalari doirasida "biometrik kiritish" atamasining eng aniq ta'rifini toping.
\begin{enumerate}[label=\Alph*),leftmargin=9mm]
\item tasvirdagi bosma matnni tahrirlanadigan belgilar sifatida tanish
\item oldindan belgilangan joylardagi belgilarni aniqlash
\item magnit siyoh bilan bosilgan maxsus belgilarni o'qish
\item insonning noyob jismoniy yoki xulqiy belgisini kiritish
\end{enumerate}
\begin{tcolorbox}[enhanced,breakable,colback=teal!4,colframe=teal!60!black,boxrule=0.45pt,arc=1.2mm,left=2.7mm,right=2.7mm,top=1.8mm,bottom=1.8mm,before skip=2mm,after skip=2mm]
{\bfseries\sffamily\color{teal}To'g'ri javob:} \textbf{D) insonning noyob jismoniy yoki xulqiy belgisini kiritish}
\par\smallskip
{\bfseries Izoh:} insonning noyob jismoniy yoki xulqiy belgisini kiritish. biometrik kiritish - insonning noyob jismoniy yoki xulqiy belgisini kiritish. Manba: MAR, PDF 8-9; Qo'shimcha sahifa raqami (qaysi PDFga tegishli ekani aniqlanmagan; darslik sahifasi deb tasdiqlanmagan): 41
\end{tcolorbox}
\vspace{2mm}\hrule height 0.2pt\vspace{3mm}

% DB_ID ec58fe7b-7e1b-3301-a4e6-a3b2996daf18
\needspace{5\baselineskip}
{\bfseries\sffamily\color{ink} 74.} {scriptsizecolor{muted} Y1 | review | ec58fe7b}
\par
Attestatsiyadagi nozik farqlar va integratsion masalalar doirasida "cheklov" atamasining eng aniq ta'rifini toping.
\begin{enumerate}[label=\Alph*),leftmargin=9mm]
\item yechim bajarishi shart bo'lgan chegara yoki talab
\item obyektni muhim belgilari bo'yicha guruhga ajratish
\item variantlarni baholash uchun qo'llanadigan belgi
\item bir nechta mavzuni bitta vaziyatda bog'laydigan topshiriq
\end{enumerate}
\begin{tcolorbox}[enhanced,breakable,colback=teal!4,colframe=teal!60!black,boxrule=0.45pt,arc=1.2mm,left=2.7mm,right=2.7mm,top=1.8mm,bottom=1.8mm,before skip=2mm,after skip=2mm]
{\bfseries\sffamily\color{teal}To'g'ri javob:} \textbf{A) yechim bajarishi shart bo'lgan chegara yoki talab}
\par\smallskip
{\bfseries Izoh:} yechim bajarishi shart bo'lgan chegara yoki talab. cheklov - yechim bajarishi shart bo'lgan chegara yoki talab. Manba: CAM10/11, PDF 23-61; Qo'shimcha sahifa raqami (qaysi PDFga tegishli ekani aniqlanmagan; darslik sahifasi deb tasdiqlanmagan): 117
\end{tcolorbox}
\vspace{2mm}\hrule height 0.2pt\vspace{3mm}

% DB_ID efd18602-ed14-118b-d17a-c6fd34dfd897
\needspace{5\baselineskip}
{\bfseries\sffamily\color{ink} 75.} {scriptsizecolor{muted} Y1 | review | efd18602}
\par
Kompyuter tizimi va uning ishlash modeli doirasida "kompyuter tizimi" atamasining eng aniq ta'rifini toping.
\begin{enumerate}[label=\Alph*),leftmargin=9mm]
\item bir maqsad uchun hamkorlik qiluvchi texnik vosita, dastur, ma'lumot va foydalanuvchi majmuasi
\item kompyuter tizimining fizik qismlari
\item kompyuterga vazifalarni bajarishni ko'rsatuvchi dasturlar
\item kiritish, qayta ishlash, chiqarish va saqlash modeli
\end{enumerate}
\begin{tcolorbox}[enhanced,breakable,colback=teal!4,colframe=teal!60!black,boxrule=0.45pt,arc=1.2mm,left=2.7mm,right=2.7mm,top=1.8mm,bottom=1.8mm,before skip=2mm,after skip=2mm]
{\bfseries\sffamily\color{teal}To'g'ri javob:} \textbf{A) bir maqsad uchun hamkorlik qiluvchi texnik vosita, dastur, ma'lumot va foydalanuvchi majmuasi}
\par\smallskip
{\bfseries Izoh:} bir maqsad uchun hamkorlik qiluvchi texnik vosita, dastur, ma'lumot va foydalanuvchi majmuasi. Manba: CAM9, PDF 8-15; Qo'shimcha sahifa raqami (qaysi PDFga tegishli ekani aniqlanmagan; darslik sahifasi deb tasdiqlanmagan): 3
\end{tcolorbox}
\vspace{2mm}\hrule height 0.2pt\vspace{3mm}

% DB_ID f2463ae5-d178-cd51-1206-1ce43186a195
\needspace{5\baselineskip}
{\bfseries\sffamily\color{ink} 76.} {scriptsizecolor{muted} Y1 | review | f2463ae5}
\par
Tizimli blok, bosh plata va ulanishlar doirasida "port" atamasining eng aniq ta'rifini toping.
\begin{enumerate}[label=\Alph*),leftmargin=9mm]
\item ichki komponentlarni bog'lovchi asosiy bosma plata
\item tashqi qurilma ulanadigan nuqta
\item ma'lumot va boshqaruv signallarini tashuvchi aloqa yo'llari
\item ichki komponent o'rnatiladigan ulagich
\end{enumerate}
\begin{tcolorbox}[enhanced,breakable,colback=teal!4,colframe=teal!60!black,boxrule=0.45pt,arc=1.2mm,left=2.7mm,right=2.7mm,top=1.8mm,bottom=1.8mm,before skip=2mm,after skip=2mm]
{\bfseries\sffamily\color{teal}To'g'ri javob:} \textbf{B) tashqi qurilma ulanadigan nuqta}
\par\smallskip
{\bfseries Izoh:} tashqi qurilma ulanadigan nuqta. port - tashqi qurilma ulanadigan nuqta. Manba: MAR, PDF 5-8; Qo'shimcha sahifa raqami (qaysi PDFga tegishli ekani aniqlanmagan; darslik sahifasi deb tasdiqlanmagan): 13
\end{tcolorbox}
\vspace{2mm}\hrule height 0.2pt\vspace{3mm}

% DB_ID f2c7220f-0564-52c8-cfb2-d0417e7af8af
\needspace{5\baselineskip}
{\bfseries\sffamily\color{ink} 77.} {scriptsizecolor{muted} Y1 | review | f2c7220f}
\par
Kiritish qurilmalari doirasida "sensor" atamasining eng aniq ta'rifini toping.
\begin{enumerate}[label=\Alph*),leftmargin=9mm]
\item tasvirdagi bosma matnni tahrirlanadigan belgilar sifatida tanish
\item oldindan belgilangan joylardagi belgilarni aniqlash
\item fizik kattalikni o'lchab signalga aylantiruvchi qurilma
\item magnit siyoh bilan bosilgan maxsus belgilarni o'qish
\end{enumerate}
\begin{tcolorbox}[enhanced,breakable,colback=teal!4,colframe=teal!60!black,boxrule=0.45pt,arc=1.2mm,left=2.7mm,right=2.7mm,top=1.8mm,bottom=1.8mm,before skip=2mm,after skip=2mm]
{\bfseries\sffamily\color{teal}To'g'ri javob:} \textbf{C) fizik kattalikni o'lchab signalga aylantiruvchi qurilma}
\par\smallskip
{\bfseries Izoh:} fizik kattalikni o'lchab signalga aylantiruvchi qurilma. sensor - fizik kattalikni o'lchab signalga aylantiruvchi qurilma. Manba: ICT10, PDF 8-12; Qo'shimcha sahifa raqami (qaysi PDFga tegishli ekani aniqlanmagan; darslik sahifasi deb tasdiqlanmagan): 40
\end{tcolorbox}
\vspace{2mm}\hrule height 0.2pt\vspace{3mm}

% DB_ID f7d08fde-aee5-0c52-f407-ccbf3e671178
\needspace{5\baselineskip}
{\bfseries\sffamily\color{ink} 78.} {scriptsizecolor{muted} Y1 | review | f7d08fde}
\par
Ikkilamchi saqlash qurilmalari doirasida "optik disk" atamasining eng aniq ta'rifini toping.
\begin{enumerate}[label=\Alph*),leftmargin=9mm]
\item magnit plastinkali, katta sig'imli ikkilamchi xotira
\item flesh hujayralarga asoslangan harakatsiz ikkilamchi xotira
\item lazer orqali o'qiladigan va yoziladigan disk
\item ketma-ket murojaatli, arxiv uchun qulay magnit vosita
\end{enumerate}
\begin{tcolorbox}[enhanced,breakable,colback=teal!4,colframe=teal!60!black,boxrule=0.45pt,arc=1.2mm,left=2.7mm,right=2.7mm,top=1.8mm,bottom=1.8mm,before skip=2mm,after skip=2mm]
{\bfseries\sffamily\color{teal}To'g'ri javob:} \textbf{C) lazer orqali o'qiladigan va yoziladigan disk}
\par\smallskip
{\bfseries Izoh:} lazer orqali o'qiladigan va yoziladigan disk. optik disk - lazer orqali o'qiladigan va yoziladigan disk. Manba: ICT10, PDF 14-18; Qo'shimcha sahifa raqami (qaysi PDFga tegishli ekani aniqlanmagan; darslik sahifasi deb tasdiqlanmagan): 32
\end{tcolorbox}
\vspace{2mm}\hrule height 0.2pt\vspace{3mm}

% DB_ID 04626ba5-3e94-05ce-1294-d185e8656c99
\needspace{5\baselineskip}
{\bfseries\sffamily\color{ink} 79.} {scriptsizecolor{muted} Y1 | review | 04626ba5}
\par
Markaziy protsessor va bajarish sikli mavzusida "CPU" qatnashgan atama-ta'rif juftliklaridan to'g'risini toping.
\begin{enumerate}[label=\Alph*),leftmargin=9mm]
\item CPU - arifmetik va mantiqiy amallarni bajaruvchi qism
\item CPU - ko'rsatmalar oqimi va boshqa qismlar ishini boshqaruvchi qism
\item CPU - ko'rsatmalarni bajarib, ma'lumotlarni qayta ishlovchi markaziy qurilma
\item CPU - CPU ichidagi juda tez va kichik xotira
\end{enumerate}
\begin{tcolorbox}[enhanced,breakable,colback=teal!4,colframe=teal!60!black,boxrule=0.45pt,arc=1.2mm,left=2.7mm,right=2.7mm,top=1.8mm,bottom=1.8mm,before skip=2mm,after skip=2mm]
{\bfseries\sffamily\color{teal}To'g'ri javob:} \textbf{C) CPU - ko'rsatmalarni bajarib, ma'lumotlarni qayta ishlovchi markaziy qurilma}
\par\smallskip
{\bfseries Izoh:} CPU - ko'rsatmalarni bajarib, ma'lumotlarni qayta ishlovchi markaziy qurilma. To'g'ri juftlik: CPU - ko'rsatmalarni bajarib, ma'lumotlarni qayta ishlovchi markaziy qurilma. Manba: CAM9, PDF 10-14; Qo'shimcha sahifa raqami (qaysi PDFga tegishli ekani aniqlanmagan; darslik sahifasi deb tasdiqlanmagan): 17
\end{tcolorbox}
\vspace{2mm}\hrule height 0.2pt\vspace{3mm}

% DB_ID 046af31c-b925-4105-924b-623f0590bc68
\needspace{5\baselineskip}
{\bfseries\sffamily\color{ink} 80.} {scriptsizecolor{muted} Y1 | review | 046af31c}
\par
Birlamchi xotira: registr, kesh, RAM va ROM mavzusidagi vazifada faol dastur va ma'lumotlar uchun o'zgaruvchan tezkor xotira talab qilindi. Qaysi tushuncha mos?
\begin{enumerate}[label=\Alph*),leftmargin=9mm]
\item ROM
\item kesh
\item virtual xotira
\item RAM
\end{enumerate}
\begin{tcolorbox}[enhanced,breakable,colback=teal!4,colframe=teal!60!black,boxrule=0.45pt,arc=1.2mm,left=2.7mm,right=2.7mm,top=1.8mm,bottom=1.8mm,before skip=2mm,after skip=2mm]
{\bfseries\sffamily\color{teal}To'g'ri javob:} \textbf{D) RAM}
\par\smallskip
{\bfseries Izoh:} RAM. Berilgan vazifa RAM tushunchasiga mos: faol dastur va ma'lumotlar uchun o'zgaruvchan tezkor xotira. Manba: CAM9, PDF 11-14; Qo'shimcha sahifa raqami (qaysi PDFga tegishli ekani aniqlanmagan; darslik sahifasi deb tasdiqlanmagan): 23
\end{tcolorbox}
\vspace{2mm}\hrule height 0.2pt\vspace{3mm}

% DB_ID 06542b03-1de5-4a37-71f7-d59f3554ffe8
\needspace{5\baselineskip}
{\bfseries\sffamily\color{ink} 81.} {scriptsizecolor{muted} Y1 | review | 06542b03}
\par
Tizimli blok, bosh plata va ulanishlar mavzusida "bosh plata" qatnashgan atama-ta'rif juftliklaridan to'g'risini toping.
\begin{enumerate}[label=\Alph*),leftmargin=9mm]
\item bosh plata - ichki komponentlarni bog'lovchi asosiy bosma plata
\item bosh plata - ma'lumot va boshqaruv signallarini tashuvchi aloqa yo'llari
\item bosh plata - ichki komponent o'rnatiladigan ulagich
\item bosh plata - tashqi qurilma ulanadigan nuqta
\end{enumerate}
\begin{tcolorbox}[enhanced,breakable,colback=teal!4,colframe=teal!60!black,boxrule=0.45pt,arc=1.2mm,left=2.7mm,right=2.7mm,top=1.8mm,bottom=1.8mm,before skip=2mm,after skip=2mm]
{\bfseries\sffamily\color{teal}To'g'ri javob:} \textbf{A) bosh plata - ichki komponentlarni bog'lovchi asosiy bosma plata}
\par\smallskip
{\bfseries Izoh:} bosh plata - ichki komponentlarni bog'lovchi asosiy bosma plata. To'g'ri juftlik: bosh plata - ichki komponentlarni bog'lovchi asosiy bosma plata. Manba: CAM9, PDF 9-14; Qo'shimcha sahifa raqami (qaysi PDFga tegishli ekani aniqlanmagan; darslik sahifasi deb tasdiqlanmagan): 10
\end{tcolorbox}
\vspace{2mm}\hrule height 0.2pt\vspace{3mm}

% DB_ID 09bea107-a473-17ff-e6ff-6783e3a2ed46
\needspace{5\baselineskip}
{\bfseries\sffamily\color{ink} 82.} {scriptsizecolor{muted} Y1 | review | 09bea107}
\par
Kompyuter tizimi va uning ishlash modeli mavzusidagi vazifada turli dasturlar yordamida ko'p xil vazifani bajaradigan kompyuter tizimi talab qilindi. Qaysi tushuncha mos?
\begin{enumerate}[label=\Alph*),leftmargin=9mm]
\item kompyuter tizimi
\item texnik ta'minot
\item umumiy maqsadli tizim
\item dasturiy ta'minot
\end{enumerate}
\begin{tcolorbox}[enhanced,breakable,colback=teal!4,colframe=teal!60!black,boxrule=0.45pt,arc=1.2mm,left=2.7mm,right=2.7mm,top=1.8mm,bottom=1.8mm,before skip=2mm,after skip=2mm]
{\bfseries\sffamily\color{teal}To'g'ri javob:} \textbf{C) umumiy maqsadli tizim}
\par\smallskip
{\bfseries Izoh:} umumiy maqsadli tizim. Berilgan vazifa umumiy maqsadli tizim tushunchasiga mos: turli dasturlar yordamida ko'p xil vazifani bajaradigan kompyuter tizimi. Manba: MAR, PDF 4-10; Qo'shimcha sahifa raqami (qaysi PDFga tegishli ekani aniqlanmagan; darslik sahifasi deb tasdiqlanmagan): 9
\end{tcolorbox}
\vspace{2mm}\hrule height 0.2pt\vspace{3mm}

% DB_ID 0a82400c-877b-6fc8-065b-b88356723850
\needspace{5\baselineskip}
{\bfseries\sffamily\color{ink} 83.} {scriptsizecolor{muted} Y1 | review | 0a82400c}
\par
Kiritish qurilmalari mavzusida "MICR" qatnashgan atama-ta'rif juftliklaridan to'g'risini toping.
\begin{enumerate}[label=\Alph*),leftmargin=9mm]
\item MICR - tasvirdagi bosma matnni tahrirlanadigan belgilar sifatida tanish
\item MICR - oldindan belgilangan joylardagi belgilarni aniqlash
\item MICR - magnit siyoh bilan bosilgan maxsus belgilarni o'qish
\item MICR - chiziqli koddagi identifikatorni avtomatik kiritish
\end{enumerate}
\begin{tcolorbox}[enhanced,breakable,colback=teal!4,colframe=teal!60!black,boxrule=0.45pt,arc=1.2mm,left=2.7mm,right=2.7mm,top=1.8mm,bottom=1.8mm,before skip=2mm,after skip=2mm]
{\bfseries\sffamily\color{teal}To'g'ri javob:} \textbf{C) MICR - magnit siyoh bilan bosilgan maxsus belgilarni o'qish}
\par\smallskip
{\bfseries Izoh:} MICR - magnit siyoh bilan bosilgan maxsus belgilarni o'qish. To'g'ri juftlik: MICR - magnit siyoh bilan bosilgan maxsus belgilarni o'qish. Manba: CAM5, PDF 60-70; Qo'shimcha sahifa raqami (qaysi PDFga tegishli ekani aniqlanmagan; darslik sahifasi deb tasdiqlanmagan): 38
\end{tcolorbox}
\vspace{2mm}\hrule height 0.2pt\vspace{3mm}

% DB_ID 0e942168-f4c7-0e90-2ef6-efe2e178d583
\needspace{5\baselineskip}
{\bfseries\sffamily\color{ink} 84.} {scriptsizecolor{muted} Y1 | review | 0e942168}
\par
Birlamchi xotira: registr, kesh, RAM va ROM mavzusida "kesh" qatnashgan atama-ta'rif juftliklaridan to'g'risini toping.
\begin{enumerate}[label=\Alph*),leftmargin=9mm]
\item kesh - faol dastur va ma'lumotlar uchun o'zgaruvchan tezkor xotira
\item kesh - elektrsiz ham ma'lumotni saqlovchi o'zgarmas xotira
\item kesh - tez-tez kerak bo'ladigan ma'lumotni CPUga yaqin saqlovchi xotira
\item kesh - RAM yetishmaganda ikkilamchi xotiradan ajratilgan sekinroq maydon
\end{enumerate}
\begin{tcolorbox}[enhanced,breakable,colback=teal!4,colframe=teal!60!black,boxrule=0.45pt,arc=1.2mm,left=2.7mm,right=2.7mm,top=1.8mm,bottom=1.8mm,before skip=2mm,after skip=2mm]
{\bfseries\sffamily\color{teal}To'g'ri javob:} \textbf{C) kesh - tez-tez kerak bo'ladigan ma'lumotni CPUga yaqin saqlovchi xotira}
\par\smallskip
{\bfseries Izoh:} kesh - tez-tez kerak bo'ladigan ma'lumotni CPUga yaqin saqlovchi xotira. To'g'ri juftlik: kesh - teztez kerak bo'ladigan ma'lumotni CPUga yaqin saqlovchi xotira. Manba: ICT10, PDF 15-20; Qo'shimcha sahifa raqami (qaysi PDFga tegishli ekani aniqlanmagan; darslik sahifasi deb tasdiqlanmagan): 25
\end{tcolorbox}
\vspace{2mm}\hrule height 0.2pt\vspace{3mm}

% DB_ID 0ee821dd-84eb-ea95-ed2c-12824e27f371
\needspace{5\baselineskip}
{\bfseries\sffamily\color{ink} 85.} {scriptsizecolor{muted} Y1 | review | 0ee821dd}
\par
Attestatsiyadagi nozik farqlar va integratsion masalalar mavzusidagi vazifada bir nechta mavzuni bitta vaziyatda bog'laydigan topshiriq talab qilindi. Qaysi tushuncha mos?
\begin{enumerate}[label=\Alph*),leftmargin=9mm]
\item tasnif
\item mezon
\item integratsion savol
\item cheklov
\end{enumerate}
\begin{tcolorbox}[enhanced,breakable,colback=teal!4,colframe=teal!60!black,boxrule=0.45pt,arc=1.2mm,left=2.7mm,right=2.7mm,top=1.8mm,bottom=1.8mm,before skip=2mm,after skip=2mm]
{\bfseries\sffamily\color{teal}To'g'ri javob:} \textbf{C) integratsion savol}
\par\smallskip
{\bfseries Izoh:} integratsion savol. Berilgan vazifa integratsion savol tushunchasiga mos: bir nechta mavzuni bitta vaziyatda bog'laydigan topshiriq. Manba: ICT5, PDF 25-49; Qo'shimcha sahifa raqami (qaysi PDFga tegishli ekani aniqlanmagan; darslik sahifasi deb tasdiqlanmagan): 118
\end{tcolorbox}
\vspace{2mm}\hrule height 0.2pt\vspace{3mm}

% DB_ID 10ca75f0-644b-2588-c6d2-5abb98691c24
\needspace{5\baselineskip}
{\bfseries\sffamily\color{ink} 86.} {scriptsizecolor{muted} Y1 | review | 10ca75f0}
\par
Kiritish qurilmalari mavzusidagi vazifada oldindan belgilangan joylardagi belgilarni aniqlash talab qilindi. Qaysi tushuncha mos?
\begin{enumerate}[label=\Alph*),leftmargin=9mm]
\item OMR
\item OCR
\item MICR
\item shtrix-kod o'quvchi
\end{enumerate}
\begin{tcolorbox}[enhanced,breakable,colback=teal!4,colframe=teal!60!black,boxrule=0.45pt,arc=1.2mm,left=2.7mm,right=2.7mm,top=1.8mm,bottom=1.8mm,before skip=2mm,after skip=2mm]
{\bfseries\sffamily\color{teal}To'g'ri javob:} \textbf{A) OMR}
\par\smallskip
{\bfseries Izoh:} OMR. Berilgan vazifa OMR tushunchasiga mos: oldindan belgilangan joylardagi belgilarni aniqlash. Manba: CAM10/11, PDF 37-45; Qo'shimcha sahifa raqami (qaysi PDFga tegishli ekani aniqlanmagan; darslik sahifasi deb tasdiqlanmagan): 37
\end{tcolorbox}
\vspace{2mm}\hrule height 0.2pt\vspace{3mm}

% DB_ID 1363527b-22c0-e843-e308-4afb09d10ef6
\needspace{5\baselineskip}
{\bfseries\sffamily\color{ink} 87.} {scriptsizecolor{muted} Y1 | review | 1363527b}
\par
Unumdorlik, nosozliklarni aniqlash va xavfsiz xizmat mavzusidagi vazifada belgi va dalillar asosida sababni tizimli aniqlash talab qilindi. Qaysi tushuncha mos?
\begin{enumerate}[label=\Alph*),leftmargin=9mm]
\item diagnostika
\item bottleneck
\item paging
\item thermal throttling
\end{enumerate}
\begin{tcolorbox}[enhanced,breakable,colback=teal!4,colframe=teal!60!black,boxrule=0.45pt,arc=1.2mm,left=2.7mm,right=2.7mm,top=1.8mm,bottom=1.8mm,before skip=2mm,after skip=2mm]
{\bfseries\sffamily\color{teal}To'g'ri javob:} \textbf{A) diagnostika}
\par\smallskip
{\bfseries Izoh:} diagnostika. Berilgan vazifa diagnostika tushunchasiga mos: belgi va dalillar asosida sababni tizimli aniqlash. Manba: MAR, PDF 10-20; Qo'shimcha sahifa raqami (qaysi PDFga tegishli ekani aniqlanmagan; darslik sahifasi deb tasdiqlanmagan): 111
\end{tcolorbox}
\vspace{2mm}\hrule height 0.2pt\vspace{3mm}

% DB_ID 137cf0a5-7a07-e9a6-2037-71d15f5bbda2
\needspace{5\baselineskip}
{\bfseries\sffamily\color{ink} 88.} {scriptsizecolor{muted} Y1 | review | 137cf0a5}
\par
Birlamchi xotira: registr, kesh, RAM va ROM mavzusidagi vazifada qurilma ichiga o'rnatilgan past darajali boshqaruv dasturi talab qilindi. Qaysi tushuncha mos?
\begin{enumerate}[label=\Alph*),leftmargin=9mm]
\item RAM
\item ROM
\item kesh
\item firmware
\end{enumerate}
\begin{tcolorbox}[enhanced,breakable,colback=teal!4,colframe=teal!60!black,boxrule=0.45pt,arc=1.2mm,left=2.7mm,right=2.7mm,top=1.8mm,bottom=1.8mm,before skip=2mm,after skip=2mm]
{\bfseries\sffamily\color{teal}To'g'ri javob:} \textbf{D) firmware}
\par\smallskip
{\bfseries Izoh:} firmware. Berilgan vazifa firmware tushunchasiga mos: qurilma ichiga o'rnatilgan past darajali boshqaruv dasturi. Manba: CAM9, PDF 11-14; Qo'shimcha sahifa raqami (qaysi PDFga tegishli ekani aniqlanmagan; darslik sahifasi deb tasdiqlanmagan): 28
\end{tcolorbox}
\vspace{2mm}\hrule height 0.2pt\vspace{3mm}

% DB_ID 1675719d-2cc2-3c86-8cb1-e211a20dfa54
\needspace{5\baselineskip}
{\bfseries\sffamily\color{ink} 89.} {scriptsizecolor{muted} Y1 | review | 1675719d}
\par
Markaziy protsessor va bajarish sikli mavzusida "CU" qatnashgan atama-ta'rif juftliklaridan to'g'risini toping.
\begin{enumerate}[label=\Alph*),leftmargin=9mm]
\item CU - ko'rsatmalar oqimi va boshqa qismlar ishini boshqaruvchi qism
\item CU - ko'rsatmalarni bajarib, ma'lumotlarni qayta ishlovchi markaziy qurilma
\item CU - arifmetik va mantiqiy amallarni bajaruvchi qism
\item CU - CPU ichidagi juda tez va kichik xotira
\end{enumerate}
\begin{tcolorbox}[enhanced,breakable,colback=teal!4,colframe=teal!60!black,boxrule=0.45pt,arc=1.2mm,left=2.7mm,right=2.7mm,top=1.8mm,bottom=1.8mm,before skip=2mm,after skip=2mm]
{\bfseries\sffamily\color{teal}To'g'ri javob:} \textbf{A) CU - ko'rsatmalar oqimi va boshqa qismlar ishini boshqaruvchi qism}
\par\smallskip
{\bfseries Izoh:} CU - ko'rsatmalar oqimi va boshqa qismlar ishini boshqaruvchi qism. To'g'ri juftlik: CU - ko'rsatmalar oqimi va boshqa qismlar ishini boshqaruvchi qism. Manba: ICT10, PDF 17-20; Qo'shimcha sahifa raqami (qaysi PDFga tegishli ekani aniqlanmagan; darslik sahifasi deb tasdiqlanmagan): 19
\end{tcolorbox}
\vspace{2mm}\hrule height 0.2pt\vspace{3mm}

% DB_ID 1b4637fd-f0cb-3f68-bea5-4383deee54fe
\needspace{5\baselineskip}
{\bfseries\sffamily\color{ink} 90.} {scriptsizecolor{muted} Y1 | review | 1b4637fd}
\par
Chiqarish va ikki vazifali qurilmalar mavzusida "yumshoq nusxa" qatnashgan atama-ta'rif juftliklaridan to'g'risini toping.
\begin{enumerate}[label=\Alph*),leftmargin=9mm]
\item yumshoq nusxa - qog'oz yoki boshqa fizik muhitdagi natija
\item yumshoq nusxa - toner va elektrofotografik jarayon bilan tez chop etuvchi printer
\item yumshoq nusxa - ekranda yoki audio orqali vaqtinchalik taqdim etilgan natija
\item yumshoq nusxa - mayda siyoh tomchilarini purkab tasvir hosil qiluvchi printer
\end{enumerate}
\begin{tcolorbox}[enhanced,breakable,colback=teal!4,colframe=teal!60!black,boxrule=0.45pt,arc=1.2mm,left=2.7mm,right=2.7mm,top=1.8mm,bottom=1.8mm,before skip=2mm,after skip=2mm]
{\bfseries\sffamily\color{teal}To'g'ri javob:} \textbf{C) yumshoq nusxa - ekranda yoki audio orqali vaqtinchalik taqdim etilgan natija}
\par\smallskip
{\bfseries Izoh:} yumshoq nusxa - ekranda yoki audio orqali vaqtinchalik taqdim etilgan natija. To'g'ri juftlik: yumshoq nusxa - ekranda yoki audio orqali vaqtinchalik taqdim etilgan natija. Manba: CAM9, PDF 16-32; Qo'shimcha sahifa raqami (qaysi PDFga tegishli ekani aniqlanmagan; darslik sahifasi deb tasdiqlanmagan): 42
\end{tcolorbox}
\vspace{2mm}\hrule height 0.2pt\vspace{3mm}

% DB_ID 1d76f77a-4406-83ba-fbac-805d659ae2b3
\needspace{5\baselineskip}
{\bfseries\sffamily\color{ink} 91.} {scriptsizecolor{muted} Y1 | review | 1d76f77a}
\par
Attestatsiyadagi nozik farqlar va integratsion masalalar mavzusida "cheklov" qatnashgan atama-ta'rif juftliklaridan to'g'risini toping.
\begin{enumerate}[label=\Alph*),leftmargin=9mm]
\item cheklov - obyektni muhim belgilari bo'yicha guruhga ajratish
\item cheklov - variantlarni baholash uchun qo'llanadigan belgi
\item cheklov - yechim bajarishi shart bo'lgan chegara yoki talab
\item cheklov - bir nechta mavzuni bitta vaziyatda bog'laydigan topshiriq
\end{enumerate}
\begin{tcolorbox}[enhanced,breakable,colback=teal!4,colframe=teal!60!black,boxrule=0.45pt,arc=1.2mm,left=2.7mm,right=2.7mm,top=1.8mm,bottom=1.8mm,before skip=2mm,after skip=2mm]
{\bfseries\sffamily\color{teal}To'g'ri javob:} \textbf{C) cheklov - yechim bajarishi shart bo'lgan chegara yoki talab}
\par\smallskip
{\bfseries Izoh:} cheklov - yechim bajarishi shart bo'lgan chegara yoki talab. To'g'ri juftlik: cheklov - yechim bajarishi shart bo'lgan chegara yoki talab. Manba: CAM10/11, PDF 23-61; Qo'shimcha sahifa raqami (qaysi PDFga tegishli ekani aniqlanmagan; darslik sahifasi deb tasdiqlanmagan): 117
\end{tcolorbox}
\vspace{2mm}\hrule height 0.2pt\vspace{3mm}

% DB_ID 2079932f-ed26-165f-a699-a09f4ae40c42
\needspace{5\baselineskip}
{\bfseries\sffamily\color{ink} 92.} {scriptsizecolor{muted} Y1 | review | 2079932f}
\par
Markaziy protsessor va bajarish sikli mavzusidagi vazifada ko'rsatmalarni bajarib, ma'lumotlarni qayta ishlovchi markaziy qurilma talab qilindi. Qaysi tushuncha mos?
\begin{enumerate}[label=\Alph*),leftmargin=9mm]
\item ALU
\item CPU
\item CU
\item registr
\end{enumerate}
\begin{tcolorbox}[enhanced,breakable,colback=teal!4,colframe=teal!60!black,boxrule=0.45pt,arc=1.2mm,left=2.7mm,right=2.7mm,top=1.8mm,bottom=1.8mm,before skip=2mm,after skip=2mm]
{\bfseries\sffamily\color{teal}To'g'ri javob:} \textbf{B) CPU}
\par\smallskip
{\bfseries Izoh:} CPU. Berilgan vazifa CPU tushunchasiga mos: ko'rsatmalarni bajarib, ma'lumotlarni qayta ishlovchi markaziy qurilma. Manba: CAM9, PDF 10-14; Qo'shimcha sahifa raqami (qaysi PDFga tegishli ekani aniqlanmagan; darslik sahifasi deb tasdiqlanmagan): 17
\end{tcolorbox}
\vspace{2mm}\hrule height 0.2pt\vspace{3mm}

% DB_ID 208372b4-fbb4-d3be-ca66-594534e5ed4c
\needspace{5\baselineskip}
{\bfseries\sffamily\color{ink} 93.} {scriptsizecolor{muted} Y1 | review | 208372b4}
\par
Attestatsiyadagi nozik farqlar va integratsion masalalar mavzusidagi vazifada yechim bajarishi shart bo'lgan chegara yoki talab talab qilindi. Qaysi tushuncha mos?
\begin{enumerate}[label=\Alph*),leftmargin=9mm]
\item tasnif
\item cheklov
\item mezon
\item integratsion savol
\end{enumerate}
\begin{tcolorbox}[enhanced,breakable,colback=teal!4,colframe=teal!60!black,boxrule=0.45pt,arc=1.2mm,left=2.7mm,right=2.7mm,top=1.8mm,bottom=1.8mm,before skip=2mm,after skip=2mm]
{\bfseries\sffamily\color{teal}To'g'ri javob:} \textbf{B) cheklov}
\par\smallskip
{\bfseries Izoh:} cheklov. Berilgan vazifa cheklov tushunchasiga mos: yechim bajarishi shart bo'lgan chegara yoki talab. Manba: CAM10/11, PDF 23-61; Qo'shimcha sahifa raqami (qaysi PDFga tegishli ekani aniqlanmagan; darslik sahifasi deb tasdiqlanmagan): 117
\end{tcolorbox}
\vspace{2mm}\hrule height 0.2pt\vspace{3mm}

% DB_ID 21ed7927-b9c8-c693-5ae4-74cec652c815
\needspace{5\baselineskip}
{\bfseries\sffamily\color{ink} 94.} {scriptsizecolor{muted} Y1 | review | 21ed7927}
\par
Ikkilamchi saqlash qurilmalari mavzusidagi vazifada lazer orqali o'qiladigan va yoziladigan disk talab qilindi. Qaysi tushuncha mos?
\begin{enumerate}[label=\Alph*),leftmargin=9mm]
\item HDD
\item SSD
\item magnit lenta
\item optik disk
\end{enumerate}
\begin{tcolorbox}[enhanced,breakable,colback=teal!4,colframe=teal!60!black,boxrule=0.45pt,arc=1.2mm,left=2.7mm,right=2.7mm,top=1.8mm,bottom=1.8mm,before skip=2mm,after skip=2mm]
{\bfseries\sffamily\color{teal}To'g'ri javob:} \textbf{D) optik disk}
\par\smallskip
{\bfseries Izoh:} optik disk. Berilgan vazifa optik disk tushunchasiga mos: lazer orqali o'qiladigan va yoziladigan disk. Manba: ICT10, PDF 14-18; Qo'shimcha sahifa raqami (qaysi PDFga tegishli ekani aniqlanmagan; darslik sahifasi deb tasdiqlanmagan): 32
\end{tcolorbox}
\vspace{2mm}\hrule height 0.2pt\vspace{3mm}

% DB_ID 250039ee-33ab-932a-07f7-f5acf8157fc3
\needspace{5\baselineskip}
{\bfseries\sffamily\color{ink} 95.} {scriptsizecolor{muted} Y1 | review | 250039ee}
\par
Unumdorlik, nosozliklarni aniqlash va xavfsiz xizmat mavzusida "log" qatnashgan atama-ta'rif juftliklaridan to'g'risini toping.
\begin{enumerate}[label=\Alph*),leftmargin=9mm]
\item log - tizim umumiy tezligini cheklayotgan eng sust resurs
\item log - RAM va ikkilamchi xotira o'rtasida sahifalarni almashish
\item log - tizim hodisalari va xatolari qaydi
\item log - qizish sabab protsessor tezligining avtomatik pasayishi
\end{enumerate}
\begin{tcolorbox}[enhanced,breakable,colback=teal!4,colframe=teal!60!black,boxrule=0.45pt,arc=1.2mm,left=2.7mm,right=2.7mm,top=1.8mm,bottom=1.8mm,before skip=2mm,after skip=2mm]
{\bfseries\sffamily\color{teal}To'g'ri javob:} \textbf{C) log - tizim hodisalari va xatolari qaydi}
\par\smallskip
{\bfseries Izoh:} log - tizim hodisalari va xatolari qaydi. To'g'ri juftlik: log - tizim hodisalari va xatolari qaydi. Manba: CAM9, PDF 12-47; Qo'shimcha sahifa raqami (qaysi PDFga tegishli ekani aniqlanmagan; darslik sahifasi deb tasdiqlanmagan): 113
\end{tcolorbox}
\vspace{2mm}\hrule height 0.2pt\vspace{3mm}

% DB_ID 278d43d6-681f-8149-ffcb-d7f06aee40b3
\needspace{5\baselineskip}
{\bfseries\sffamily\color{ink} 96.} {scriptsizecolor{muted} Y1 | review | 278d43d6}
\par
Chiqarish va ikki vazifali qurilmalar mavzusidagi vazifada mayda siyoh tomchilarini purkab tasvir hosil qiluvchi printer talab qilindi. Qaysi tushuncha mos?
\begin{enumerate}[label=\Alph*),leftmargin=9mm]
\item siyohpurkagich
\item yumshoq nusxa
\item qattiq nusxa
\item lazer printer
\end{enumerate}
\begin{tcolorbox}[enhanced,breakable,colback=teal!4,colframe=teal!60!black,boxrule=0.45pt,arc=1.2mm,left=2.7mm,right=2.7mm,top=1.8mm,bottom=1.8mm,before skip=2mm,after skip=2mm]
{\bfseries\sffamily\color{teal}To'g'ri javob:} \textbf{A) siyohpurkagich}
\par\smallskip
{\bfseries Izoh:} siyohpurkagich. Berilgan vazifa siyohpurkagich tushunchasiga mos: mayda siyoh tomchilarini purkab tasvir hosil qiluvchi printer. Manba: ICT10, PDF 9-13; Qo'shimcha sahifa raqami (qaysi PDFga tegishli ekani aniqlanmagan; darslik sahifasi deb tasdiqlanmagan): 45
\end{tcolorbox}
\vspace{2mm}\hrule height 0.2pt\vspace{3mm}

% DB_ID 28d66f92-5014-8d9b-07c4-2376c268f156
\needspace{5\baselineskip}
{\bfseries\sffamily\color{ink} 97.} {scriptsizecolor{muted} Y1 | review | 28d66f92}
\par
Attestatsiyadagi nozik farqlar va integratsion masalalar mavzusidagi vazifada variantlarni baholash uchun qo'llanadigan belgi talab qilindi. Qaysi tushuncha mos?
\begin{enumerate}[label=\Alph*),leftmargin=9mm]
\item mezon
\item tasnif
\item cheklov
\item integratsion savol
\end{enumerate}
\begin{tcolorbox}[enhanced,breakable,colback=teal!4,colframe=teal!60!black,boxrule=0.45pt,arc=1.2mm,left=2.7mm,right=2.7mm,top=1.8mm,bottom=1.8mm,before skip=2mm,after skip=2mm]
{\bfseries\sffamily\color{teal}To'g'ri javob:} \textbf{A) mezon}
\par\smallskip
{\bfseries Izoh:} mezon. Berilgan vazifa mezon tushunchasiga mos: variantlarni baholash uchun qo'llanadigan belgi. Manba: CAM9, PDF 8-47; Qo'shimcha sahifa raqami (qaysi PDFga tegishli ekani aniqlanmagan; darslik sahifasi deb tasdiqlanmagan): 116
\end{tcolorbox}
\vspace{2mm}\hrule height 0.2pt\vspace{3mm}

% DB_ID 2a922111-8774-4c80-c493-456efb8b9d59
\needspace{5\baselineskip}
{\bfseries\sffamily\color{ink} 98.} {scriptsizecolor{muted} Y1 | review | 2a922111}
\par
Attestatsiyadagi nozik farqlar va integratsion masalalar mavzusidagi vazifada tanlovni vaziyatdagi dalil bilan bog'lash talab qilindi. Qaysi tushuncha mos?
\begin{enumerate}[label=\Alph*),leftmargin=9mm]
\item asoslash
\item tasnif
\item mezon
\item cheklov
\end{enumerate}
\begin{tcolorbox}[enhanced,breakable,colback=teal!4,colframe=teal!60!black,boxrule=0.45pt,arc=1.2mm,left=2.7mm,right=2.7mm,top=1.8mm,bottom=1.8mm,before skip=2mm,after skip=2mm]
{\bfseries\sffamily\color{teal}To'g'ri javob:} \textbf{A) asoslash}
\par\smallskip
{\bfseries Izoh:} asoslash. Berilgan vazifa asoslash tushunchasiga mos: tanlovni vaziyatdagi dalil bilan bog'lash. Manba: SPEC, PDF 3; Qo'shimcha sahifa raqami (qaysi PDFga tegishli ekani aniqlanmagan; darslik sahifasi deb tasdiqlanmagan): 120
\end{tcolorbox}
\vspace{2mm}\hrule height 0.2pt\vspace{3mm}

% DB_ID 2c6b9023-91ee-8279-1632-bf5edf38f006
\needspace{5\baselineskip}
{\bfseries\sffamily\color{ink} 99.} {scriptsizecolor{muted} Y1 | review | 2c6b9023}
\par
Chiqarish va ikki vazifali qurilmalar mavzusidagi vazifada boshqaruv signalini fizik harakat yoki ta'sirga aylantiruvchi qurilma talab qilindi. Qaysi tushuncha mos?
\begin{enumerate}[label=\Alph*),leftmargin=9mm]
\item yumshoq nusxa
\item aktuator
\item qattiq nusxa
\item lazer printer
\end{enumerate}
\begin{tcolorbox}[enhanced,breakable,colback=teal!4,colframe=teal!60!black,boxrule=0.45pt,arc=1.2mm,left=2.7mm,right=2.7mm,top=1.8mm,bottom=1.8mm,before skip=2mm,after skip=2mm]
{\bfseries\sffamily\color{teal}To'g'ri javob:} \textbf{B) aktuator}
\par\smallskip
{\bfseries Izoh:} aktuator. Berilgan vazifa aktuator tushunchasiga mos: boshqaruv signalini fizik harakat yoki ta'sirga aylantiruvchi qurilma. Manba: MAR, PDF 8-10; Qo'shimcha sahifa raqami (qaysi PDFga tegishli ekani aniqlanmagan; darslik sahifasi deb tasdiqlanmagan): 46
\end{tcolorbox}
\vspace{2mm}\hrule height 0.2pt\vspace{3mm}

% DB_ID 2ce7bc40-8d4f-368b-2db0-d967ca8ed814
\needspace{5\baselineskip}
{\bfseries\sffamily\color{ink} 100.} {scriptsizecolor{muted} Y1 | review | 2ce7bc40}
\par
Kompyuter tizimi va uning ishlash modeli mavzusida "dasturiy ta'minot" qatnashgan atama-ta'rif juftliklaridan to'g'risini toping.
\begin{enumerate}[label=\Alph*),leftmargin=9mm]
\item dasturiy ta'minot - kompyuterga vazifalarni bajarishni ko'rsatuvchi dasturlar
\item dasturiy ta'minot - bir maqsad uchun hamkorlik qiluvchi texnik vosita, dastur, ma'lumot va foydalanuvchi majmuasi
\item dasturiy ta'minot - kompyuter tizimining fizik qismlari
\item dasturiy ta'minot - kiritish, qayta ishlash, chiqarish va saqlash modeli
\end{enumerate}
\begin{tcolorbox}[enhanced,breakable,colback=teal!4,colframe=teal!60!black,boxrule=0.45pt,arc=1.2mm,left=2.7mm,right=2.7mm,top=1.8mm,bottom=1.8mm,before skip=2mm,after skip=2mm]
{\bfseries\sffamily\color{teal}To'g'ri javob:} \textbf{A) dasturiy ta'minot - kompyuterga vazifalarni bajarishni ko'rsatuvchi dasturlar}
\par\smallskip
{\bfseries Izoh:} dasturiy ta'minot - kompyuterga vazifalarni bajarishni ko'rsatuvchi dasturlar. To'g'ri juftlik: dasturiy ta'minot - kompyuterga vazifalarni bajarishni ko'rsatuvchi dasturlar. Manba: ICT5, PDF 25-32; Qo'shimcha sahifa raqami (qaysi PDFga tegishli ekani aniqlanmagan; darslik sahifasi deb tasdiqlanmagan): 5
\end{tcolorbox}
\vspace{2mm}\hrule height 0.2pt\vspace{3mm}

% DB_ID 31bd5e6b-57c3-ef90-31d9-1593803a41ac
\needspace{5\baselineskip}
{\bfseries\sffamily\color{ink} 101.} {scriptsizecolor{muted} Y1 | review | 31bd5e6b}
\par
Chiqarish va ikki vazifali qurilmalar mavzusidagi vazifada tebranish yoki kuch orqali beriladigan sezgi chiqishi talab qilindi. Qaysi tushuncha mos?
\begin{enumerate}[label=\Alph*),leftmargin=9mm]
\item yumshoq nusxa
\item qattiq nusxa
\item haptik chiqish
\item lazer printer
\end{enumerate}
\begin{tcolorbox}[enhanced,breakable,colback=teal!4,colframe=teal!60!black,boxrule=0.45pt,arc=1.2mm,left=2.7mm,right=2.7mm,top=1.8mm,bottom=1.8mm,before skip=2mm,after skip=2mm]
{\bfseries\sffamily\color{teal}To'g'ri javob:} \textbf{C) haptik chiqish}
\par\smallskip
{\bfseries Izoh:} haptik chiqish. Berilgan vazifa haptik chiqish tushunchasiga mos: tebranish yoki kuch orqali beriladigan sezgi chiqishi. Manba: CAM9, PDF 16-32; Qo'shimcha sahifa raqami (qaysi PDFga tegishli ekani aniqlanmagan; darslik sahifasi deb tasdiqlanmagan): 47
\end{tcolorbox}
\vspace{2mm}\hrule height 0.2pt\vspace{3mm}

% DB_ID 4071cc4b-4a01-1a27-c47a-97c1a983956c
\needspace{5\baselineskip}
{\bfseries\sffamily\color{ink} 102.} {scriptsizecolor{muted} Y1 | review | 4071cc4b}
\par
Ikkilamchi saqlash qurilmalari mavzusida "USB flesh" qatnashgan atama-ta'rif juftliklaridan to'g'risini toping.
\begin{enumerate}[label=\Alph*),leftmargin=9mm]
\item USB flesh - magnit plastinkali, katta sig'imli ikkilamchi xotira
\item USB flesh - flesh hujayralarga asoslangan harakatsiz ikkilamchi xotira
\item USB flesh - portativ solid-state saqlash qurilmasi
\item USB flesh - lazer orqali o'qiladigan va yoziladigan disk
\end{enumerate}
\begin{tcolorbox}[enhanced,breakable,colback=teal!4,colframe=teal!60!black,boxrule=0.45pt,arc=1.2mm,left=2.7mm,right=2.7mm,top=1.8mm,bottom=1.8mm,before skip=2mm,after skip=2mm]
{\bfseries\sffamily\color{teal}To'g'ri javob:} \textbf{C) USB flesh - portativ solid-state saqlash qurilmasi}
\par\smallskip
{\bfseries Izoh:} USB flesh - portativ solidstate saqlash qurilmasi. To'g'ri juftlik: USB flesh - portativ solid-state saqlash qurilmasi. Manba: MAR, PDF 7-9; Qo'shimcha sahifa raqami (qaysi PDFga tegishli ekani aniqlanmagan; darslik sahifasi deb tasdiqlanmagan): 34
\end{tcolorbox}
\vspace{2mm}\hrule height 0.2pt\vspace{3mm}

% DB_ID 5298681e-547e-85f9-9528-5edb7e84f997
\needspace{5\baselineskip}
{\bfseries\sffamily\color{ink} 103.} {scriptsizecolor{muted} Y1 | review | 5298681e}
\par
Birlamchi xotira: registr, kesh, RAM va ROM mavzusidagi vazifada tez-tez kerak bo'ladigan ma'lumotni CPUga yaqin saqlovchi xotira talab qilindi. Qaysi tushuncha mos?
\begin{enumerate}[label=\Alph*),leftmargin=9mm]
\item RAM
\item kesh
\item ROM
\item virtual xotira
\end{enumerate}
\begin{tcolorbox}[enhanced,breakable,colback=teal!4,colframe=teal!60!black,boxrule=0.45pt,arc=1.2mm,left=2.7mm,right=2.7mm,top=1.8mm,bottom=1.8mm,before skip=2mm,after skip=2mm]
{\bfseries\sffamily\color{teal}To'g'ri javob:} \textbf{B) kesh}
\par\smallskip
{\bfseries Izoh:} kesh. Berilgan vazifa kesh tushunchasiga mos: tez-tez kerak bo'ladigan ma'lumotni CPUga yaqin saqlovchi xotira. Manba: ICT10, PDF 15-20; Qo'shimcha sahifa raqami (qaysi PDFga tegishli ekani aniqlanmagan; darslik sahifasi deb tasdiqlanmagan): 25
\end{tcolorbox}
\vspace{2mm}\hrule height 0.2pt\vspace{3mm}

% DB_ID 52ff18c9-7567-fee3-3beb-0cfdc12ab7df
\needspace{5\baselineskip}
{\bfseries\sffamily\color{ink} 104.} {scriptsizecolor{muted} Y1 | review | 52ff18c9}
\par
Unumdorlik, nosozliklarni aniqlash va xavfsiz xizmat mavzusidagi vazifada RAM va ikkilamchi xotira o'rtasida sahifalarni almashish talab qilindi. Qaysi tushuncha mos?
\begin{enumerate}[label=\Alph*),leftmargin=9mm]
\item bottleneck
\item thermal throttling
\item paging
\item diagnostika
\end{enumerate}
\begin{tcolorbox}[enhanced,breakable,colback=teal!4,colframe=teal!60!black,boxrule=0.45pt,arc=1.2mm,left=2.7mm,right=2.7mm,top=1.8mm,bottom=1.8mm,before skip=2mm,after skip=2mm]
{\bfseries\sffamily\color{teal}To'g'ri javob:} \textbf{C) paging}
\par\smallskip
{\bfseries Izoh:} paging. Berilgan vazifa paging tushunchasiga mos: RAM va ikkilamchi xotira o'rtasida sahifalarni almashish. Manba: CAM10/11, PDF 24-61; Qo'shimcha sahifa raqami (qaysi PDFga tegishli ekani aniqlanmagan; darslik sahifasi deb tasdiqlanmagan): 109
\end{tcolorbox}
\vspace{2mm}\hrule height 0.2pt\vspace{3mm}

% DB_ID 53c79e95-4f35-7d00-f634-05d38ab2d6e7
\needspace{5\baselineskip}
{\bfseries\sffamily\color{ink} 105.} {scriptsizecolor{muted} Y1 | review | 53c79e95}
\par
Birlamchi xotira: registr, kesh, RAM va ROM mavzusida "firmware" qatnashgan atama-ta'rif juftliklaridan to'g'risini toping.
\begin{enumerate}[label=\Alph*),leftmargin=9mm]
\item firmware - qurilma ichiga o'rnatilgan past darajali boshqaruv dasturi
\item firmware - faol dastur va ma'lumotlar uchun o'zgaruvchan tezkor xotira
\item firmware - elektrsiz ham ma'lumotni saqlovchi o'zgarmas xotira
\item firmware - tez-tez kerak bo'ladigan ma'lumotni CPUga yaqin saqlovchi xotira
\end{enumerate}
\begin{tcolorbox}[enhanced,breakable,colback=teal!4,colframe=teal!60!black,boxrule=0.45pt,arc=1.2mm,left=2.7mm,right=2.7mm,top=1.8mm,bottom=1.8mm,before skip=2mm,after skip=2mm]
{\bfseries\sffamily\color{teal}To'g'ri javob:} \textbf{A) firmware - qurilma ichiga o'rnatilgan past darajali boshqaruv dasturi}
\par\smallskip
{\bfseries Izoh:} firmware - qurilma ichiga o'rnatilgan past darajali boshqaruv dasturi. To'g'ri juftlik: firmware - qurilma ichiga o'rnatilgan past darajali boshqaruv dasturi. Manba: CAM9, PDF 11-14; Qo'shimcha sahifa raqami (qaysi PDFga tegishli ekani aniqlanmagan; darslik sahifasi deb tasdiqlanmagan): 28
\end{tcolorbox}
\vspace{2mm}\hrule height 0.2pt\vspace{3mm}

% DB_ID 5c3f9b05-41f8-cc6f-8488-3b23c175a140
\needspace{5\baselineskip}
{\bfseries\sffamily\color{ink} 106.} {scriptsizecolor{muted} Y1 | review | 5c3f9b05}
\par
Kompyuter tizimi va uning ishlash modeli mavzusidagi vazifada kattaroq qurilma ichida muayyan vazifani bajaruvchi kompyuter tizimi talab qilindi. Qaysi tushuncha mos?
\begin{enumerate}[label=\Alph*),leftmargin=9mm]
\item kompyuter tizimi
\item o'rnatilgan tizim
\item texnik ta'minot
\item dasturiy ta'minot
\end{enumerate}
\begin{tcolorbox}[enhanced,breakable,colback=teal!4,colframe=teal!60!black,boxrule=0.45pt,arc=1.2mm,left=2.7mm,right=2.7mm,top=1.8mm,bottom=1.8mm,before skip=2mm,after skip=2mm]
{\bfseries\sffamily\color{teal}To'g'ri javob:} \textbf{B) o'rnatilgan tizim}
\par\smallskip
{\bfseries Izoh:} o'rnatilgan tizim. Berilgan vazifa o'rnatilgan tizim tushunchasiga mos: kattaroq qurilma ichida muayyan vazifani bajaruvchi kompyuter tizimi. Manba: ICT9, PDF 14; Qo'shimcha sahifa raqami (qaysi PDFga tegishli ekani aniqlanmagan; darslik sahifasi deb tasdiqlanmagan): 7
\end{tcolorbox}
\vspace{2mm}\hrule height 0.2pt\vspace{3mm}

% DB_ID 5d437ed0-6cc9-1abf-a5a2-ceb4aa84d83a
\needspace{5\baselineskip}
{\bfseries\sffamily\color{ink} 107.} {scriptsizecolor{muted} Y1 | review | 5d437ed0}
\par
Unumdorlik, nosozliklarni aniqlash va xavfsiz xizmat mavzusidagi vazifada qizish sabab protsessor tezligining avtomatik pasayishi talab qilindi. Qaysi tushuncha mos?
\begin{enumerate}[label=\Alph*),leftmargin=9mm]
\item bottleneck
\item paging
\item diagnostika
\item thermal throttling
\end{enumerate}
\begin{tcolorbox}[enhanced,breakable,colback=teal!4,colframe=teal!60!black,boxrule=0.45pt,arc=1.2mm,left=2.7mm,right=2.7mm,top=1.8mm,bottom=1.8mm,before skip=2mm,after skip=2mm]
{\bfseries\sffamily\color{teal}To'g'ri javob:} \textbf{D) thermal throttling}
\par\smallskip
{\bfseries Izoh:} thermal throttling. Berilgan vazifa thermal throttling tushunchasiga mos: qizish sabab protsessor tezligining avtomatik pasayishi. Manba: ICT10, PDF 12-29; Qo'shimcha sahifa raqami (qaysi PDFga tegishli ekani aniqlanmagan; darslik sahifasi deb tasdiqlanmagan): 110
\end{tcolorbox}
\vspace{2mm}\hrule height 0.2pt\vspace{3mm}

% DB_ID 63fac448-9816-502a-cc21-4a94d4a74f8e
\needspace{5\baselineskip}
{\bfseries\sffamily\color{ink} 108.} {scriptsizecolor{muted} Y1 | review | 63fac448}
\par
Markaziy protsessor va bajarish sikli mavzusida "takt chastotasi" qatnashgan atama-ta'rif juftliklaridan to'g'risini toping.
\begin{enumerate}[label=\Alph*),leftmargin=9mm]
\item takt chastotasi - ko'rsatmalarni bajarib, ma'lumotlarni qayta ishlovchi markaziy qurilma
\item takt chastotasi - arifmetik va mantiqiy amallarni bajaruvchi qism
\item takt chastotasi - sekunddagi takt impulslar soni
\item takt chastotasi - ko'rsatmalar oqimi va boshqa qismlar ishini boshqaruvchi qism
\end{enumerate}
\begin{tcolorbox}[enhanced,breakable,colback=teal!4,colframe=teal!60!black,boxrule=0.45pt,arc=1.2mm,left=2.7mm,right=2.7mm,top=1.8mm,bottom=1.8mm,before skip=2mm,after skip=2mm]
{\bfseries\sffamily\color{teal}To'g'ri javob:} \textbf{C) takt chastotasi - sekunddagi takt impulslar soni}
\par\smallskip
{\bfseries Izoh:} takt chastotasi - sekunddagi takt impulslar soni. To'g'ri juftlik: takt chastotasi - sekunddagi takt impulslar soni. Manba: CAM9, PDF 10-14; Qo'shimcha sahifa raqami (qaysi PDFga tegishli ekani aniqlanmagan; darslik sahifasi deb tasdiqlanmagan): 21
\end{tcolorbox}
\vspace{2mm}\hrule height 0.2pt\vspace{3mm}

% DB_ID 689b2163-ca3a-6dfe-16d1-6936fda894e7
\needspace{5\baselineskip}
{\bfseries\sffamily\color{ink} 109.} {scriptsizecolor{muted} Y1 | review | 689b2163}
\par
Kiritish qurilmalari mavzusida "sensor" qatnashgan atama-ta'rif juftliklaridan to'g'risini toping.
\begin{enumerate}[label=\Alph*),leftmargin=9mm]
\item sensor - fizik kattalikni o'lchab signalga aylantiruvchi qurilma
\item sensor - tasvirdagi bosma matnni tahrirlanadigan belgilar sifatida tanish
\item sensor - oldindan belgilangan joylardagi belgilarni aniqlash
\item sensor - magnit siyoh bilan bosilgan maxsus belgilarni o'qish
\end{enumerate}
\begin{tcolorbox}[enhanced,breakable,colback=teal!4,colframe=teal!60!black,boxrule=0.45pt,arc=1.2mm,left=2.7mm,right=2.7mm,top=1.8mm,bottom=1.8mm,before skip=2mm,after skip=2mm]
{\bfseries\sffamily\color{teal}To'g'ri javob:} \textbf{A) sensor - fizik kattalikni o'lchab signalga aylantiruvchi qurilma}
\par\smallskip
{\bfseries Izoh:} sensor - fizik kattalikni o'lchab signalga aylantiruvchi qurilma. To'g'ri juftlik: sensor - fizik kattalikni o'lchab signalga aylantiruvchi qurilma. Manba: ICT10, PDF 8-12; Qo'shimcha sahifa raqami (qaysi PDFga tegishli ekani aniqlanmagan; darslik sahifasi deb tasdiqlanmagan): 40
\end{tcolorbox}
\vspace{2mm}\hrule height 0.2pt\vspace{3mm}

% DB_ID 6dd47ec3-6cc1-7ea3-d86a-c98690bfc4e9
\needspace{5\baselineskip}
{\bfseries\sffamily\color{ink} 110.} {scriptsizecolor{muted} Y1 | review | 6dd47ec3}
\par
Ikkilamchi saqlash qurilmalari mavzusidagi vazifada tarmoq orqali masofadagi server xotirasidan foydalanish xizmati talab qilindi. Qaysi tushuncha mos?
\begin{enumerate}[label=\Alph*),leftmargin=9mm]
\item HDD
\item SSD
\item bulutli saqlash
\item optik disk
\end{enumerate}
\begin{tcolorbox}[enhanced,breakable,colback=teal!4,colframe=teal!60!black,boxrule=0.45pt,arc=1.2mm,left=2.7mm,right=2.7mm,top=1.8mm,bottom=1.8mm,before skip=2mm,after skip=2mm]
{\bfseries\sffamily\color{teal}To'g'ri javob:} \textbf{C) bulutli saqlash}
\par\smallskip
{\bfseries Izoh:} bulutli saqlash. Berilgan vazifa bulutli saqlash tushunchasiga mos: tarmoq orqali masofadagi server xotirasidan foydalanish xizmati. Manba: CAM9, PDF 33-47; Qo'shimcha sahifa raqami (qaysi PDFga tegishli ekani aniqlanmagan; darslik sahifasi deb tasdiqlanmagan): 35
\end{tcolorbox}
\vspace{2mm}\hrule height 0.2pt\vspace{3mm}

% DB_ID 6fa5ccaa-c261-a093-c605-b9c2eed03cb0
\needspace{5\baselineskip}
{\bfseries\sffamily\color{ink} 111.} {scriptsizecolor{muted} Y1 | review | 6fa5ccaa}
\par
Tizimli blok, bosh plata va ulanishlar mavzusidagi vazifada tarmoq elektrini komponentlarga mos doimiy tokka aylantiruvchi qism talab qilindi. Qaysi tushuncha mos?
\begin{enumerate}[label=\Alph*),leftmargin=9mm]
\item bosh plata
\item shina
\item slot
\item quvvat bloki
\end{enumerate}
\begin{tcolorbox}[enhanced,breakable,colback=teal!4,colframe=teal!60!black,boxrule=0.45pt,arc=1.2mm,left=2.7mm,right=2.7mm,top=1.8mm,bottom=1.8mm,before skip=2mm,after skip=2mm]
{\bfseries\sffamily\color{teal}To'g'ri javob:} \textbf{D) quvvat bloki}
\par\smallskip
{\bfseries Izoh:} quvvat bloki. Berilgan vazifa quvvat bloki tushunchasiga mos: tarmoq elektrini komponentlarga mos doimiy tokka aylantiruvchi qism. Manba: CAM9, PDF 9-14; Qo'shimcha sahifa raqami (qaysi PDFga tegishli ekani aniqlanmagan; darslik sahifasi deb tasdiqlanmagan): 14
\end{tcolorbox}
\vspace{2mm}\hrule height 0.2pt\vspace{3mm}

% DB_ID 72122cd2-433a-6cdd-5551-71d945ac1089
\needspace{5\baselineskip}
{\bfseries\sffamily\color{ink} 112.} {scriptsizecolor{muted} Y1 | review | 72122cd2}
\par
Chiqarish va ikki vazifali qurilmalar mavzusidagi vazifada qog'oz yoki boshqa fizik muhitdagi natija talab qilindi. Qaysi tushuncha mos?
\begin{enumerate}[label=\Alph*),leftmargin=9mm]
\item yumshoq nusxa
\item lazer printer
\item qattiq nusxa
\item siyohpurkagich
\end{enumerate}
\begin{tcolorbox}[enhanced,breakable,colback=teal!4,colframe=teal!60!black,boxrule=0.45pt,arc=1.2mm,left=2.7mm,right=2.7mm,top=1.8mm,bottom=1.8mm,before skip=2mm,after skip=2mm]
{\bfseries\sffamily\color{teal}To'g'ri javob:} \textbf{C) qattiq nusxa}
\par\smallskip
{\bfseries Izoh:} qattiq nusxa. Berilgan vazifa qattiq nusxa tushunchasiga mos: qog'oz yoki boshqa fizik muhitdagi natija. Manba: CAM10/11, PDF 37-45; Qo'shimcha sahifa raqami (qaysi PDFga tegishli ekani aniqlanmagan; darslik sahifasi deb tasdiqlanmagan): 43
\end{tcolorbox}
\vspace{2mm}\hrule height 0.2pt\vspace{3mm}

% DB_ID 73c4af35-2124-b664-9905-79dcdee678f4
\needspace{5\baselineskip}
{\bfseries\sffamily\color{ink} 113.} {scriptsizecolor{muted} Y1 | review | 73c4af35}
\par
Kompyuter tizimi va uning ishlash modeli mavzusida "o'rnatilgan tizim" qatnashgan atama-ta'rif juftliklaridan to'g'risini toping.
\begin{enumerate}[label=\Alph*),leftmargin=9mm]
\item o'rnatilgan tizim - bir maqsad uchun hamkorlik qiluvchi texnik vosita, dastur, ma'lumot va foydalanuvchi majmuasi
\item o'rnatilgan tizim - kompyuter tizimining fizik qismlari
\item o'rnatilgan tizim - kattaroq qurilma ichida muayyan vazifani bajaruvchi kompyuter tizimi
\item o'rnatilgan tizim - kompyuterga vazifalarni bajarishni ko'rsatuvchi dasturlar
\end{enumerate}
\begin{tcolorbox}[enhanced,breakable,colback=teal!4,colframe=teal!60!black,boxrule=0.45pt,arc=1.2mm,left=2.7mm,right=2.7mm,top=1.8mm,bottom=1.8mm,before skip=2mm,after skip=2mm]
{\bfseries\sffamily\color{teal}To'g'ri javob:} \textbf{C) o'rnatilgan tizim - kattaroq qurilma ichida muayyan vazifani bajaruvchi kompyuter tizimi}
\par\smallskip
{\bfseries Izoh:} o'rnatilgan tizim - kattaroq qurilma ichida muayyan vazifani bajaruvchi kompyuter tizimi. To'g'ri juftlik: o'rnatilgan tizim - kattaroq qurilma ichida muayyan vazifani bajaruvchi kompyuter tizimi. Manba: ICT9, PDF 14; Qo'shimcha sahifa raqami (qaysi PDFga tegishli ekani aniqlanmagan; darslik sahifasi deb tasdiqlanmagan): 8
\end{tcolorbox}
\vspace{2mm}\hrule height 0.2pt\vspace{3mm}

% DB_ID 74314c96-8100-ace5-2bd0-bacac8119f8b
\needspace{5\baselineskip}
{\bfseries\sffamily\color{ink} 114.} {scriptsizecolor{muted} Y1 | review | 74314c96}
\par
Attestatsiyadagi nozik farqlar va integratsion masalalar mavzusidagi vazifada ishonarli, ammo noto'g'ri javob varianti talab qilindi. Qaysi tushuncha mos?
\begin{enumerate}[label=\Alph*),leftmargin=9mm]
\item tasnif
\item mezon
\item cheklov
\item distraktor
\end{enumerate}
\begin{tcolorbox}[enhanced,breakable,colback=teal!4,colframe=teal!60!black,boxrule=0.45pt,arc=1.2mm,left=2.7mm,right=2.7mm,top=1.8mm,bottom=1.8mm,before skip=2mm,after skip=2mm]
{\bfseries\sffamily\color{teal}To'g'ri javob:} \textbf{D) distraktor}
\par\smallskip
{\bfseries Izoh:} distraktor. Berilgan vazifa distraktor tushunchasiga mos: ishonarli, ammo noto'g'ri javob varianti. Manba: MAR, PDF 4-20; Qo'shimcha sahifa raqami (qaysi PDFga tegishli ekani aniqlanmagan; darslik sahifasi deb tasdiqlanmagan): 119
\end{tcolorbox}
\vspace{2mm}\hrule height 0.2pt\vspace{3mm}

% DB_ID 788af234-eaa0-b26f-231c-cb7026c313dc
\needspace{5\baselineskip}
{\bfseries\sffamily\color{ink} 115.} {scriptsizecolor{muted} Y1 | review | 788af234}
\par
Tizimli blok, bosh plata va ulanishlar mavzusidagi vazifada kompyuterga qo'shimcha funksiya berish uchun slotga o'rnatiladigan plata talab qilindi. Qaysi tushuncha mos?
\begin{enumerate}[label=\Alph*),leftmargin=9mm]
\item kengaytirish kartasi
\item bosh plata
\item shina
\item slot
\end{enumerate}
\begin{tcolorbox}[enhanced,breakable,colback=teal!4,colframe=teal!60!black,boxrule=0.45pt,arc=1.2mm,left=2.7mm,right=2.7mm,top=1.8mm,bottom=1.8mm,before skip=2mm,after skip=2mm]
{\bfseries\sffamily\color{teal}To'g'ri javob:} \textbf{A) kengaytirish kartasi}
\par\smallskip
{\bfseries Izoh:} kengaytirish kartasi. Berilgan vazifa kengaytirish kartasi tushunchasiga mos: kompyuterga qo'shimcha funksiya berish uchun slotga o'rnatiladigan plata. Manba: CAM10/11, PDF 32-36; Qo'shimcha sahifa raqami (qaysi PDFga tegishli ekani aniqlanmagan; darslik sahifasi deb tasdiqlanmagan): 15
\end{tcolorbox}
\vspace{2mm}\hrule height 0.2pt\vspace{3mm}

% DB_ID 78c14b43-b0c5-87f4-37dc-5dd4b0bbc303
\needspace{5\baselineskip}
{\bfseries\sffamily\color{ink} 116.} {scriptsizecolor{muted} Y1 | review | 78c14b43}
\par
Unumdorlik, nosozliklarni aniqlash va xavfsiz xizmat mavzusidagi vazifada nosozlik yuz berish ehtimolini oldindan kamaytiruvchi xizmat talab qilindi. Qaysi tushuncha mos?
\begin{enumerate}[label=\Alph*),leftmargin=9mm]
\item bottleneck
\item paging
\item profilaktika
\item thermal throttling
\end{enumerate}
\begin{tcolorbox}[enhanced,breakable,colback=teal!4,colframe=teal!60!black,boxrule=0.45pt,arc=1.2mm,left=2.7mm,right=2.7mm,top=1.8mm,bottom=1.8mm,before skip=2mm,after skip=2mm]
{\bfseries\sffamily\color{teal}To'g'ri javob:} \textbf{C) profilaktika}
\par\smallskip
{\bfseries Izoh:} profilaktika. Berilgan vazifa profilaktika tushunchasiga mos: nosozlik yuz berish ehtimolini oldindan kamaytiruvchi xizmat. Manba: CAM10/11, PDF 24-61; Qo'shimcha sahifa raqami (qaysi PDFga tegishli ekani aniqlanmagan; darslik sahifasi deb tasdiqlanmagan): 113
\end{tcolorbox}
\vspace{2mm}\hrule height 0.2pt\vspace{3mm}

% DB_ID 797f6859-1ae1-d1b6-60ae-e55177ac21bc
\needspace{5\baselineskip}
{\bfseries\sffamily\color{ink} 117.} {scriptsizecolor{muted} Y1 | review | 797f6859}
\par
Kompyuter tizimi va uning ishlash modeli mavzusidagi vazifada kompyuterga vazifalarni bajarishni ko'rsatuvchi dasturlar talab qilindi. Qaysi tushuncha mos?
\begin{enumerate}[label=\Alph*),leftmargin=9mm]
\item kompyuter tizimi
\item texnik ta'minot
\item IPOS
\item dasturiy ta'minot
\end{enumerate}
\begin{tcolorbox}[enhanced,breakable,colback=teal!4,colframe=teal!60!black,boxrule=0.45pt,arc=1.2mm,left=2.7mm,right=2.7mm,top=1.8mm,bottom=1.8mm,before skip=2mm,after skip=2mm]
{\bfseries\sffamily\color{teal}To'g'ri javob:} \textbf{D) dasturiy ta'minot}
\par\smallskip
{\bfseries Izoh:} dasturiy ta'minot. Berilgan vazifa dasturiy ta'minot tushunchasiga mos: kompyuterga vazifalarni bajarishni ko'rsatuvchi dasturlar. Manba: ICT5, PDF 25-32; Qo'shimcha sahifa raqami (qaysi PDFga tegishli ekani aniqlanmagan; darslik sahifasi deb tasdiqlanmagan): 5
\end{tcolorbox}
\vspace{2mm}\hrule height 0.2pt\vspace{3mm}

% DB_ID 7b5e6a3f-0b06-827c-b838-686aa6975aa5
\needspace{5\baselineskip}
{\bfseries\sffamily\color{ink} 118.} {scriptsizecolor{muted} Y1 | review | 7b5e6a3f}
\par
Birlamchi xotira: registr, kesh, RAM va ROM mavzusida "RAM" qatnashgan atama-ta'rif juftliklaridan to'g'risini toping.
\begin{enumerate}[label=\Alph*),leftmargin=9mm]
\item RAM - faol dastur va ma'lumotlar uchun o'zgaruvchan tezkor xotira
\item RAM - elektrsiz ham ma'lumotni saqlovchi o'zgarmas xotira
\item RAM - tez-tez kerak bo'ladigan ma'lumotni CPUga yaqin saqlovchi xotira
\item RAM - RAM yetishmaganda ikkilamchi xotiradan ajratilgan sekinroq maydon
\end{enumerate}
\begin{tcolorbox}[enhanced,breakable,colback=teal!4,colframe=teal!60!black,boxrule=0.45pt,arc=1.2mm,left=2.7mm,right=2.7mm,top=1.8mm,bottom=1.8mm,before skip=2mm,after skip=2mm]
{\bfseries\sffamily\color{teal}To'g'ri javob:} \textbf{A) RAM - faol dastur va ma'lumotlar uchun o'zgaruvchan tezkor xotira}
\par\smallskip
{\bfseries Izoh:} RAM - faol dastur va ma'lumotlar uchun o'zgaruvchan tezkor xotira. To'g'ri juftlik: RAM - faol dastur va ma'lumotlar uchun o'zgaruvchan tezkor xotira. Manba: CAM9, PDF 11-14; Qo'shimcha sahifa raqami (qaysi PDFga tegishli ekani aniqlanmagan; darslik sahifasi deb tasdiqlanmagan): 23
\end{tcolorbox}
\vspace{2mm}\hrule height 0.2pt\vspace{3mm}

% DB_ID 81304ce1-a1b8-67e6-1a35-50277be933f4
\needspace{5\baselineskip}
{\bfseries\sffamily\color{ink} 119.} {scriptsizecolor{muted} Y1 | review | 81304ce1}
\par
Attestatsiyadagi nozik farqlar va integratsion masalalar mavzusida "tasnif" qatnashgan atama-ta'rif juftliklaridan to'g'risini toping.
\begin{enumerate}[label=\Alph*),leftmargin=9mm]
\item tasnif - obyektni muhim belgilari bo'yicha guruhga ajratish
\item tasnif - variantlarni baholash uchun qo'llanadigan belgi
\item tasnif - yechim bajarishi shart bo'lgan chegara yoki talab
\item tasnif - bir nechta mavzuni bitta vaziyatda bog'laydigan topshiriq
\end{enumerate}
\begin{tcolorbox}[enhanced,breakable,colback=teal!4,colframe=teal!60!black,boxrule=0.45pt,arc=1.2mm,left=2.7mm,right=2.7mm,top=1.8mm,bottom=1.8mm,before skip=2mm,after skip=2mm]
{\bfseries\sffamily\color{teal}To'g'ri javob:} \textbf{A) tasnif - obyektni muhim belgilari bo'yicha guruhga ajratish}
\par\smallskip
{\bfseries Izoh:} tasnif - obyektni muhim belgilari bo'yicha guruhga ajratish. To'g'ri juftlik: tasnif - obyektni muhim belgilari bo'yicha guruhga ajratish. Manba: SPEC, PDF 3; Qo'shimcha sahifa raqami (qaysi PDFga tegishli ekani aniqlanmagan; darslik sahifasi deb tasdiqlanmagan): 115
\end{tcolorbox}
\vspace{2mm}\hrule height 0.2pt\vspace{3mm}

% DB_ID 85bd5e62-0a96-c838-c4ac-db6920d1c63d
\needspace{5\baselineskip}
{\bfseries\sffamily\color{ink} 120.} {scriptsizecolor{muted} Y1 | review | 85bd5e62}
\par
Attestatsiyadagi nozik farqlar va integratsion masalalar mavzusida "distraktor" qatnashgan atama-ta'rif juftliklaridan to'g'risini toping.
\begin{enumerate}[label=\Alph*),leftmargin=9mm]
\item distraktor - ishonarli, ammo noto'g'ri javob varianti
\item distraktor - obyektni muhim belgilari bo'yicha guruhga ajratish
\item distraktor - variantlarni baholash uchun qo'llanadigan belgi
\item distraktor - yechim bajarishi shart bo'lgan chegara yoki talab
\end{enumerate}
\begin{tcolorbox}[enhanced,breakable,colback=teal!4,colframe=teal!60!black,boxrule=0.45pt,arc=1.2mm,left=2.7mm,right=2.7mm,top=1.8mm,bottom=1.8mm,before skip=2mm,after skip=2mm]
{\bfseries\sffamily\color{teal}To'g'ri javob:} \textbf{A) distraktor - ishonarli, ammo noto'g'ri javob varianti}
\par\smallskip
{\bfseries Izoh:} distraktor - ishonarli, ammo noto'g'ri javob varianti. To'g'ri juftlik: distraktor - ishonarli, ammo noto'g'ri javob varianti. Manba: MAR, PDF 4-20; Qo'shimcha sahifa raqami (qaysi PDFga tegishli ekani aniqlanmagan; darslik sahifasi deb tasdiqlanmagan): 119
\end{tcolorbox}
\vspace{2mm}\hrule height 0.2pt\vspace{3mm}

% DB_ID 8b67f913-c362-3a6d-c10e-74c06fa25a3a
\needspace{5\baselineskip}
{\bfseries\sffamily\color{ink} 121.} {scriptsizecolor{muted} Y1 | review | 8b67f913}
\par
Ikkilamchi saqlash qurilmalari mavzusidagi vazifada portativ solid-state saqlash qurilmasi talab qilindi. Qaysi tushuncha mos?
\begin{enumerate}[label=\Alph*),leftmargin=9mm]
\item HDD
\item USB flesh
\item SSD
\item optik disk
\end{enumerate}
\begin{tcolorbox}[enhanced,breakable,colback=teal!4,colframe=teal!60!black,boxrule=0.45pt,arc=1.2mm,left=2.7mm,right=2.7mm,top=1.8mm,bottom=1.8mm,before skip=2mm,after skip=2mm]
{\bfseries\sffamily\color{teal}To'g'ri javob:} \textbf{B) USB flesh}
\par\smallskip
{\bfseries Izoh:} USB flesh. Berilgan vazifa USB flesh tushunchasiga mos: portativ solid-state saqlash qurilmasi. Manba: MAR, PDF 7-9; Qo'shimcha sahifa raqami (qaysi PDFga tegishli ekani aniqlanmagan; darslik sahifasi deb tasdiqlanmagan): 34
\end{tcolorbox}
\vspace{2mm}\hrule height 0.2pt\vspace{3mm}

% DB_ID 8c16dd1b-b39b-65ef-e587-1cb2d14aaf2b
\needspace{5\baselineskip}
{\bfseries\sffamily\color{ink} 122.} {scriptsizecolor{muted} Y1 | review | 8c16dd1b}
\par
Kompyuter tizimi va uning ishlash modeli mavzusida "kompyuter tizimi" qatnashgan atama-ta'rif juftliklaridan to'g'risini toping.
\begin{enumerate}[label=\Alph*),leftmargin=9mm]
\item kompyuter tizimi - kompyuter tizimining fizik qismlari
\item kompyuter tizimi - kompyuterga vazifalarni bajarishni ko'rsatuvchi dasturlar
\item kompyuter tizimi - bir maqsad uchun hamkorlik qiluvchi texnik vosita, dastur, ma'lumot va foydalanuvchi majmuasi
\item kompyuter tizimi - kiritish, qayta ishlash, chiqarish va saqlash modeli
\end{enumerate}
\begin{tcolorbox}[enhanced,breakable,colback=teal!4,colframe=teal!60!black,boxrule=0.45pt,arc=1.2mm,left=2.7mm,right=2.7mm,top=1.8mm,bottom=1.8mm,before skip=2mm,after skip=2mm]
{\bfseries\sffamily\color{teal}To'g'ri javob:} \textbf{C) kompyuter tizimi - bir maqsad uchun hamkorlik qiluvchi texnik vosita, dastur, ma'lumot va foydalanuvchi majmuasi}
\par\smallskip
{\bfseries Izoh:} kompyuter tizimi - bir maqsad uchun hamkorlik qiluvchi texnik vosita, dastur, ma'lumot va foydalanuvchi majmuasi. To'g'ri juftlik: kompyuter tizimi - bir maqsad uchun hamkorlik qiluvchi texnik vosita, dastur, ma'lumot va foydalanuvchi majmuasi. Manba: CAM9, PDF 8-15; Qo'shimcha sahifa raqami (qaysi PDFga tegishli ekani aniqlanmagan; darslik sahifasi deb tasdiqlanmagan): 3
\end{tcolorbox}
\vspace{2mm}\hrule height 0.2pt\vspace{3mm}

% DB_ID 90b497b2-e071-60ee-9a3a-d75cbb734818
\needspace{5\baselineskip}
{\bfseries\sffamily\color{ink} 123.} {scriptsizecolor{muted} Y1 | review | 90b497b2}
\par
Markaziy protsessor va bajarish sikli mavzusidagi vazifada CPU ichidagi juda tez va kichik xotira talab qilindi. Qaysi tushuncha mos?
\begin{enumerate}[label=\Alph*),leftmargin=9mm]
\item registr
\item CPU
\item ALU
\item CU
\end{enumerate}
\begin{tcolorbox}[enhanced,breakable,colback=teal!4,colframe=teal!60!black,boxrule=0.45pt,arc=1.2mm,left=2.7mm,right=2.7mm,top=1.8mm,bottom=1.8mm,before skip=2mm,after skip=2mm]
{\bfseries\sffamily\color{teal}To'g'ri javob:} \textbf{A) registr}
\par\smallskip
{\bfseries Izoh:} registr. Berilgan vazifa registr tushunchasiga mos: CPU ichidagi juda tez va kichik xotira. Manba: MAR, PDF 6-7; Qo'shimcha sahifa raqami (qaysi PDFga tegishli ekani aniqlanmagan; darslik sahifasi deb tasdiqlanmagan): 20
\end{tcolorbox}
\vspace{2mm}\hrule height 0.2pt\vspace{3mm}

% DB_ID 91246e52-e610-4060-20f1-a56e4b3980a5
\needspace{5\baselineskip}
{\bfseries\sffamily\color{ink} 124.} {scriptsizecolor{muted} Y1 | review | 91246e52}
\par
Birlamchi xotira: registr, kesh, RAM va ROM mavzusidagi vazifada RAM yetishmaganda ikkilamchi xotiradan ajratilgan sekinroq maydon talab qilindi. Qaysi tushuncha mos?
\begin{enumerate}[label=\Alph*),leftmargin=9mm]
\item RAM
\item ROM
\item virtual xotira
\item kesh
\end{enumerate}
\begin{tcolorbox}[enhanced,breakable,colback=teal!4,colframe=teal!60!black,boxrule=0.45pt,arc=1.2mm,left=2.7mm,right=2.7mm,top=1.8mm,bottom=1.8mm,before skip=2mm,after skip=2mm]
{\bfseries\sffamily\color{teal}To'g'ri javob:} \textbf{C) virtual xotira}
\par\smallskip
{\bfseries Izoh:} virtual xotira. Berilgan vazifa virtual xotira tushunchasiga mos: RAM yetishmaganda ikkilamchi xotiradan ajratilgan sekinroq maydon. Manba: MAR, PDF 7-8; Qo'shimcha sahifa raqami (qaysi PDFga tegishli ekani aniqlanmagan; darslik sahifasi deb tasdiqlanmagan): 26
\end{tcolorbox}
\vspace{2mm}\hrule height 0.2pt\vspace{3mm}

% DB_ID 937eb42a-019f-b019-f678-edeef7599b86
\needspace{5\baselineskip}
{\bfseries\sffamily\color{ink} 125.} {scriptsizecolor{muted} Y1 | review | 937eb42a}
\par
Ikkilamchi saqlash qurilmalari mavzusida "optik disk" qatnashgan atama-ta'rif juftliklaridan to'g'risini toping.
\begin{enumerate}[label=\Alph*),leftmargin=9mm]
\item optik disk - lazer orqali o'qiladigan va yoziladigan disk
\item optik disk - magnit plastinkali, katta sig'imli ikkilamchi xotira
\item optik disk - flesh hujayralarga asoslangan harakatsiz ikkilamchi xotira
\item optik disk - ketma-ket murojaatli, arxiv uchun qulay magnit vosita
\end{enumerate}
\begin{tcolorbox}[enhanced,breakable,colback=teal!4,colframe=teal!60!black,boxrule=0.45pt,arc=1.2mm,left=2.7mm,right=2.7mm,top=1.8mm,bottom=1.8mm,before skip=2mm,after skip=2mm]
{\bfseries\sffamily\color{teal}To'g'ri javob:} \textbf{A) optik disk - lazer orqali o'qiladigan va yoziladigan disk}
\par\smallskip
{\bfseries Izoh:} optik disk - lazer orqali o'qiladigan va yoziladigan disk. To'g'ri juftlik: optik disk - lazer orqali o'qiladigan va yoziladigan disk. Manba: ICT10, PDF 14-18; Qo'shimcha sahifa raqami (qaysi PDFga tegishli ekani aniqlanmagan; darslik sahifasi deb tasdiqlanmagan): 32
\end{tcolorbox}
\vspace{2mm}\hrule height 0.2pt\vspace{3mm}

% DB_ID 95409d6a-b7a1-35d3-72c2-a005801bda9c
\needspace{5\baselineskip}
{\bfseries\sffamily\color{ink} 126.} {scriptsizecolor{muted} Y1 | review | 95409d6a}
\par
Unumdorlik, nosozliklarni aniqlash va xavfsiz xizmat mavzusida "thermal throttling" qatnashgan atama-ta'rif juftliklaridan to'g'risini toping.
\begin{enumerate}[label=\Alph*),leftmargin=9mm]
\item thermal throttling - qizish sabab protsessor tezligining avtomatik pasayishi
\item thermal throttling - tizim umumiy tezligini cheklayotgan eng sust resurs
\item thermal throttling - RAM va ikkilamchi xotira o'rtasida sahifalarni almashish
\item thermal throttling - belgi va dalillar asosida sababni tizimli aniqlash
\end{enumerate}
\begin{tcolorbox}[enhanced,breakable,colback=teal!4,colframe=teal!60!black,boxrule=0.45pt,arc=1.2mm,left=2.7mm,right=2.7mm,top=1.8mm,bottom=1.8mm,before skip=2mm,after skip=2mm]
{\bfseries\sffamily\color{teal}To'g'ri javob:} \textbf{A) thermal throttling - qizish sabab protsessor tezligining avtomatik pasayishi}
\par\smallskip
{\bfseries Izoh:} thermal throttling - qizish sabab protsessor tezligining avtomatik pasayishi. To'g'ri juftlik: thermal throttling - qizish sabab protsessor tezligining avtomatik pasayishi. Manba: ICT10, PDF 12-29; Qo'shimcha sahifa raqami (qaysi PDFga tegishli ekani aniqlanmagan; darslik sahifasi deb tasdiqlanmagan): 110
\end{tcolorbox}
\vspace{2mm}\hrule height 0.2pt\vspace{3mm}

% DB_ID 95fda619-ad1f-c1d4-df23-91bb071b71a0
\needspace{5\baselineskip}
{\bfseries\sffamily\color{ink} 127.} {scriptsizecolor{muted} Y1 | review | 95fda619}
\par
Kiritish qurilmalari mavzusidagi vazifada fizik kattalikni o'lchab signalga aylantiruvchi qurilma talab qilindi. Qaysi tushuncha mos?
\begin{enumerate}[label=\Alph*),leftmargin=9mm]
\item OCR
\item OMR
\item MICR
\item sensor
\end{enumerate}
\begin{tcolorbox}[enhanced,breakable,colback=teal!4,colframe=teal!60!black,boxrule=0.45pt,arc=1.2mm,left=2.7mm,right=2.7mm,top=1.8mm,bottom=1.8mm,before skip=2mm,after skip=2mm]
{\bfseries\sffamily\color{teal}To'g'ri javob:} \textbf{D) sensor}
\par\smallskip
{\bfseries Izoh:} sensor. Berilgan vazifa sensor tushunchasiga mos: fizik kattalikni o'lchab signalga aylantiruvchi qurilma. Manba: ICT10, PDF 8-12; Qo'shimcha sahifa raqami (qaysi PDFga tegishli ekani aniqlanmagan; darslik sahifasi deb tasdiqlanmagan): 40
\end{tcolorbox}
\vspace{2mm}\hrule height 0.2pt\vspace{3mm}

% DB_ID 97f4c5a4-ef38-8b5f-afdc-725e7fa8c7bf
\needspace{5\baselineskip}
{\bfseries\sffamily\color{ink} 128.} {scriptsizecolor{muted} Y1 | review | 97f4c5a4}
\par
Kiritish qurilmalari mavzusida "OCR" qatnashgan atama-ta'rif juftliklaridan to'g'risini toping.
\begin{enumerate}[label=\Alph*),leftmargin=9mm]
\item OCR - tasvirdagi bosma matnni tahrirlanadigan belgilar sifatida tanish
\item OCR - oldindan belgilangan joylardagi belgilarni aniqlash
\item OCR - magnit siyoh bilan bosilgan maxsus belgilarni o'qish
\item OCR - chiziqli koddagi identifikatorni avtomatik kiritish
\end{enumerate}
\begin{tcolorbox}[enhanced,breakable,colback=teal!4,colframe=teal!60!black,boxrule=0.45pt,arc=1.2mm,left=2.7mm,right=2.7mm,top=1.8mm,bottom=1.8mm,before skip=2mm,after skip=2mm]
{\bfseries\sffamily\color{teal}To'g'ri javob:} \textbf{A) OCR - tasvirdagi bosma matnni tahrirlanadigan belgilar sifatida tanish}
\par\smallskip
{\bfseries Izoh:} OCR - tasvirdagi bosma matnni tahrirlanadigan belgilar sifatida tanish. To'g'ri juftlik: OCR - tasvirdagi bosma matnni tahrirlanadigan belgilar sifatida tanish. Manba: CAM9, PDF 16-32; Qo'shimcha sahifa raqami (qaysi PDFga tegishli ekani aniqlanmagan; darslik sahifasi deb tasdiqlanmagan): 36
\end{tcolorbox}
\vspace{2mm}\hrule height 0.2pt\vspace{3mm}

% DB_ID 99d05884-06c8-2d26-5db6-cb7d785df66d
\needspace{5\baselineskip}
{\bfseries\sffamily\color{ink} 129.} {scriptsizecolor{muted} Y1 | review | 99d05884}
\par
Tizimli blok, bosh plata va ulanishlar mavzusidagi vazifada ma'lumot va boshqaruv signallarini tashuvchi aloqa yo'llari talab qilindi. Qaysi tushuncha mos?
\begin{enumerate}[label=\Alph*),leftmargin=9mm]
\item shina
\item bosh plata
\item slot
\item port
\end{enumerate}
\begin{tcolorbox}[enhanced,breakable,colback=teal!4,colframe=teal!60!black,boxrule=0.45pt,arc=1.2mm,left=2.7mm,right=2.7mm,top=1.8mm,bottom=1.8mm,before skip=2mm,after skip=2mm]
{\bfseries\sffamily\color{teal}To'g'ri javob:} \textbf{A) shina}
\par\smallskip
{\bfseries Izoh:} shina. Berilgan vazifa shina tushunchasiga mos: ma'lumot va boshqaruv signallarini tashuvchi aloqa yo'llari. Manba: CAM10/11, PDF 32-36; Qo'shimcha sahifa raqami (qaysi PDFga tegishli ekani aniqlanmagan; darslik sahifasi deb tasdiqlanmagan): 11
\end{tcolorbox}
\vspace{2mm}\hrule height 0.2pt\vspace{3mm}

% DB_ID 9a0ec2dd-05db-0d13-d831-c2c83a6020f4
\needspace{5\baselineskip}
{\bfseries\sffamily\color{ink} 130.} {scriptsizecolor{muted} Y1 | review | 9a0ec2dd}
\par
Tizimli blok, bosh plata va ulanishlar mavzusidagi vazifada ichki komponentlarni bog'lovchi asosiy bosma plata talab qilindi. Qaysi tushuncha mos?
\begin{enumerate}[label=\Alph*),leftmargin=9mm]
\item shina
\item slot
\item port
\item bosh plata
\end{enumerate}
\begin{tcolorbox}[enhanced,breakable,colback=teal!4,colframe=teal!60!black,boxrule=0.45pt,arc=1.2mm,left=2.7mm,right=2.7mm,top=1.8mm,bottom=1.8mm,before skip=2mm,after skip=2mm]
{\bfseries\sffamily\color{teal}To'g'ri javob:} \textbf{D) bosh plata}
\par\smallskip
{\bfseries Izoh:} bosh plata. Berilgan vazifa bosh plata tushunchasiga mos: ichki komponentlarni bog'lovchi asosiy bosma plata. Manba: CAM9, PDF 9-14; Qo'shimcha sahifa raqami (qaysi PDFga tegishli ekani aniqlanmagan; darslik sahifasi deb tasdiqlanmagan): 10
\end{tcolorbox}
\vspace{2mm}\hrule height 0.2pt\vspace{3mm}

% DB_ID 9a9fd817-2bc6-788b-1147-2c6b278a191e
\needspace{5\baselineskip}
{\bfseries\sffamily\color{ink} 131.} {scriptsizecolor{muted} Y1 | review | 9a9fd817}
\par
Tizimli blok, bosh plata va ulanishlar mavzusidagi vazifada tashqi qurilma ulanadigan nuqta talab qilindi. Qaysi tushuncha mos?
\begin{enumerate}[label=\Alph*),leftmargin=9mm]
\item bosh plata
\item shina
\item port
\item slot
\end{enumerate}
\begin{tcolorbox}[enhanced,breakable,colback=teal!4,colframe=teal!60!black,boxrule=0.45pt,arc=1.2mm,left=2.7mm,right=2.7mm,top=1.8mm,bottom=1.8mm,before skip=2mm,after skip=2mm]
{\bfseries\sffamily\color{teal}To'g'ri javob:} \textbf{C) port}
\par\smallskip
{\bfseries Izoh:} port. Berilgan vazifa port tushunchasiga mos: tashqi qurilma ulanadigan nuqta. Manba: MAR, PDF 5-8; Qo'shimcha sahifa raqami (qaysi PDFga tegishli ekani aniqlanmagan; darslik sahifasi deb tasdiqlanmagan): 13
\end{tcolorbox}
\vspace{2mm}\hrule height 0.2pt\vspace{3mm}

% DB_ID 9b0d7230-23c3-8937-c3ad-99cbfc7e9805
\needspace{5\baselineskip}
{\bfseries\sffamily\color{ink} 132.} {scriptsizecolor{muted} Y1 | review | 9b0d7230}
\par
Unumdorlik, nosozliklarni aniqlash va xavfsiz xizmat mavzusidagi vazifada tizim umumiy tezligini cheklayotgan eng sust resurs talab qilindi. Qaysi tushuncha mos?
\begin{enumerate}[label=\Alph*),leftmargin=9mm]
\item paging
\item bottleneck
\item thermal throttling
\item diagnostika
\end{enumerate}
\begin{tcolorbox}[enhanced,breakable,colback=teal!4,colframe=teal!60!black,boxrule=0.45pt,arc=1.2mm,left=2.7mm,right=2.7mm,top=1.8mm,bottom=1.8mm,before skip=2mm,after skip=2mm]
{\bfseries\sffamily\color{teal}To'g'ri javob:} \textbf{B) bottleneck}
\par\smallskip
{\bfseries Izoh:} bottleneck. Berilgan vazifa bottleneck tushunchasiga mos: tizim umumiy tezligini cheklayotgan eng sust resurs. Manba: CAM9, PDF 12-47; Qo'shimcha sahifa raqami (qaysi PDFga tegishli ekani aniqlanmagan; darslik sahifasi deb tasdiqlanmagan): 108
\end{tcolorbox}
\vspace{2mm}\hrule height 0.2pt\vspace{3mm}

% DB_ID 9d1948d4-f820-ed86-be7e-3d91d1998b56
\needspace{5\baselineskip}
{\bfseries\sffamily\color{ink} 133.} {scriptsizecolor{muted} Y1 | review | 9d1948d4}
\par
Chiqarish va ikki vazifali qurilmalar mavzusidagi vazifada ekranda yoki audio orqali vaqtinchalik taqdim etilgan natija talab qilindi. Qaysi tushuncha mos?
\begin{enumerate}[label=\Alph*),leftmargin=9mm]
\item qattiq nusxa
\item yumshoq nusxa
\item lazer printer
\item siyohpurkagich
\end{enumerate}
\begin{tcolorbox}[enhanced,breakable,colback=teal!4,colframe=teal!60!black,boxrule=0.45pt,arc=1.2mm,left=2.7mm,right=2.7mm,top=1.8mm,bottom=1.8mm,before skip=2mm,after skip=2mm]
{\bfseries\sffamily\color{teal}To'g'ri javob:} \textbf{B) yumshoq nusxa}
\par\smallskip
{\bfseries Izoh:} yumshoq nusxa. Berilgan vazifa yumshoq nusxa tushunchasiga mos: ekranda yoki audio orqali vaqtinchalik taqdim etilgan natija. Manba: CAM9, PDF 16-32; Qo'shimcha sahifa raqami (qaysi PDFga tegishli ekani aniqlanmagan; darslik sahifasi deb tasdiqlanmagan): 42
\end{tcolorbox}
\vspace{2mm}\hrule height 0.2pt\vspace{3mm}

% DB_ID 9d34aa90-f68b-0a60-cf29-86b4a7e54658
\needspace{5\baselineskip}
{\bfseries\sffamily\color{ink} 134.} {scriptsizecolor{muted} Y1 | review | 9d34aa90}
\par
Tizimli blok, bosh plata va ulanishlar mavzusidagi vazifada ichki komponent o'rnatiladigan ulagich talab qilindi. Qaysi tushuncha mos?
\begin{enumerate}[label=\Alph*),leftmargin=9mm]
\item bosh plata
\item slot
\item shina
\item port
\end{enumerate}
\begin{tcolorbox}[enhanced,breakable,colback=teal!4,colframe=teal!60!black,boxrule=0.45pt,arc=1.2mm,left=2.7mm,right=2.7mm,top=1.8mm,bottom=1.8mm,before skip=2mm,after skip=2mm]
{\bfseries\sffamily\color{teal}To'g'ri javob:} \textbf{B) slot}
\par\smallskip
{\bfseries Izoh:} slot. Berilgan vazifa slot tushunchasiga mos: ichki komponent o'rnatiladigan ulagich. Manba: ICT10, PDF 12-18; Qo'shimcha sahifa raqami (qaysi PDFga tegishli ekani aniqlanmagan; darslik sahifasi deb tasdiqlanmagan): 12
\end{tcolorbox}
\vspace{2mm}\hrule height 0.2pt\vspace{3mm}

% DB_ID a36ff0cc-eb96-cfbd-4dc0-b85cc9eeb761
\needspace{5\baselineskip}
{\bfseries\sffamily\color{ink} 135.} {scriptsizecolor{muted} Y1 | review | a36ff0cc}
\par
Unumdorlik, nosozliklarni aniqlash va xavfsiz xizmat mavzusidagi vazifada tizim hodisalari va xatolari qaydi talab qilindi. Qaysi tushuncha mos?
\begin{enumerate}[label=\Alph*),leftmargin=9mm]
\item bottleneck
\item log
\item paging
\item thermal throttling
\end{enumerate}
\begin{tcolorbox}[enhanced,breakable,colback=teal!4,colframe=teal!60!black,boxrule=0.45pt,arc=1.2mm,left=2.7mm,right=2.7mm,top=1.8mm,bottom=1.8mm,before skip=2mm,after skip=2mm]
{\bfseries\sffamily\color{teal}To'g'ri javob:} \textbf{B) log}
\par\smallskip
{\bfseries Izoh:} log. Berilgan vazifa log tushunchasiga mos: tizim hodisalari va xatolari qaydi. Manba: CAM9, PDF 12-47; Qo'shimcha sahifa raqami (qaysi PDFga tegishli ekani aniqlanmagan; darslik sahifasi deb tasdiqlanmagan): 112
\end{tcolorbox}
\vspace{2mm}\hrule height 0.2pt\vspace{3mm}

% DB_ID a89b9690-7049-392e-42e8-f17b29d271c7
\needspace{5\baselineskip}
{\bfseries\sffamily\color{ink} 136.} {scriptsizecolor{muted} Y1 | review | a89b9690}
\par
Birlamchi xotira: registr, kesh, RAM va ROM mavzusidagi vazifada qurilmani boshlang'ich sozlab, OT yuklanishini boshlovchi firmware muhiti talab qilindi. Qaysi tushuncha mos?
\begin{enumerate}[label=\Alph*),leftmargin=9mm]
\item UEFI
\item RAM
\item ROM
\item kesh
\end{enumerate}
\begin{tcolorbox}[enhanced,breakable,colback=teal!4,colframe=teal!60!black,boxrule=0.45pt,arc=1.2mm,left=2.7mm,right=2.7mm,top=1.8mm,bottom=1.8mm,before skip=2mm,after skip=2mm]
{\bfseries\sffamily\color{teal}To'g'ri javob:} \textbf{A) UEFI}
\par\smallskip
{\bfseries Izoh:} UEFI. Berilgan vazifa UEFI tushunchasiga mos: qurilmani boshlang'ich sozlab, OT yuklanishini boshlovchi firmware muhiti. Manba: CAM10/11, PDF 33-36; Qo'shimcha sahifa raqami (qaysi PDFga tegishli ekani aniqlanmagan; darslik sahifasi deb tasdiqlanmagan): 29
\end{tcolorbox}
\vspace{2mm}\hrule height 0.2pt\vspace{3mm}

% DB_ID ab661d40-83d4-da4f-1a36-b25518aa5543
\needspace{5\baselineskip}
{\bfseries\sffamily\color{ink} 137.} {scriptsizecolor{muted} Y1 | review | ab661d40}
\par
Tizimli blok, bosh plata va ulanishlar mavzusida "quvvat bloki" qatnashgan atama-ta'rif juftliklaridan to'g'risini toping.
\begin{enumerate}[label=\Alph*),leftmargin=9mm]
\item quvvat bloki - tarmoq elektrini komponentlarga mos doimiy tokka aylantiruvchi qism
\item quvvat bloki - ichki komponentlarni bog'lovchi asosiy bosma plata
\item quvvat bloki - ma'lumot va boshqaruv signallarini tashuvchi aloqa yo'llari
\item quvvat bloki - ichki komponent o'rnatiladigan ulagich
\end{enumerate}
\begin{tcolorbox}[enhanced,breakable,colback=teal!4,colframe=teal!60!black,boxrule=0.45pt,arc=1.2mm,left=2.7mm,right=2.7mm,top=1.8mm,bottom=1.8mm,before skip=2mm,after skip=2mm]
{\bfseries\sffamily\color{teal}To'g'ri javob:} \textbf{A) quvvat bloki - tarmoq elektrini komponentlarga mos doimiy tokka aylantiruvchi qism}
\par\smallskip
{\bfseries Izoh:} quvvat bloki - tarmoq elektrini komponentlarga mos doimiy tokka aylantiruvchi qism. To'g'ri juftlik: quvvat bloki - tarmoq elektrini komponentlarga mos doimiy tokka aylantiruvchi qism. Manba: CAM9, PDF 9-14; Qo'shimcha sahifa raqami (qaysi PDFga tegishli ekani aniqlanmagan; darslik sahifasi deb tasdiqlanmagan): 14
\end{tcolorbox}
\vspace{2mm}\hrule height 0.2pt\vspace{3mm}

% DB_ID b1b7d86e-a91d-9a0c-63f2-ccca1ad012f5
\needspace{5\baselineskip}
{\bfseries\sffamily\color{ink} 138.} {scriptsizecolor{muted} Y1 | review | b1b7d86e}
\par
Attestatsiyadagi nozik farqlar va integratsion masalalar mavzusidagi vazifada obyektni muhim belgilari bo'yicha guruhga ajratish talab qilindi. Qaysi tushuncha mos?
\begin{enumerate}[label=\Alph*),leftmargin=9mm]
\item mezon
\item cheklov
\item integratsion savol
\item tasnif
\end{enumerate}
\begin{tcolorbox}[enhanced,breakable,colback=teal!4,colframe=teal!60!black,boxrule=0.45pt,arc=1.2mm,left=2.7mm,right=2.7mm,top=1.8mm,bottom=1.8mm,before skip=2mm,after skip=2mm]
{\bfseries\sffamily\color{teal}To'g'ri javob:} \textbf{D) tasnif}
\par\smallskip
{\bfseries Izoh:} tasnif. Berilgan vazifa tasnif tushunchasiga mos: obyektni muhim belgilari bo'yicha guruhga ajratish. Manba: SPEC, PDF 3; Qo'shimcha sahifa raqami (qaysi PDFga tegishli ekani aniqlanmagan; darslik sahifasi deb tasdiqlanmagan): 115
\end{tcolorbox}
\vspace{2mm}\hrule height 0.2pt\vspace{3mm}

% DB_ID ba5e9f38-b9c7-67a8-cb0d-4d661ac34a96
\needspace{5\baselineskip}
{\bfseries\sffamily\color{ink} 139.} {scriptsizecolor{muted} Y1 | review | ba5e9f38}
\par
Kiritish qurilmalari mavzusidagi vazifada insonning noyob jismoniy yoki xulqiy belgisini kiritish talab qilindi. Qaysi tushuncha mos?
\begin{enumerate}[label=\Alph*),leftmargin=9mm]
\item biometrik kiritish
\item OCR
\item OMR
\item MICR
\end{enumerate}
\begin{tcolorbox}[enhanced,breakable,colback=teal!4,colframe=teal!60!black,boxrule=0.45pt,arc=1.2mm,left=2.7mm,right=2.7mm,top=1.8mm,bottom=1.8mm,before skip=2mm,after skip=2mm]
{\bfseries\sffamily\color{teal}To'g'ri javob:} \textbf{A) biometrik kiritish}
\par\smallskip
{\bfseries Izoh:} biometrik kiritish. Berilgan vazifa biometrik kiritish tushunchasiga mos: insonning noyob jismoniy yoki xulqiy belgisini kiritish. Manba: MAR, PDF 8-9; Qo'shimcha sahifa raqami (qaysi PDFga tegishli ekani aniqlanmagan; darslik sahifasi deb tasdiqlanmagan): 41
\end{tcolorbox}
\vspace{2mm}\hrule height 0.2pt\vspace{3mm}

% DB_ID bb1e9000-8122-b53e-0788-880a1b58af46
\needspace{5\baselineskip}
{\bfseries\sffamily\color{ink} 140.} {scriptsizecolor{muted} Y1 | review | bb1e9000}
\par
Unumdorlik, nosozliklarni aniqlash va xavfsiz xizmat mavzusida "bottleneck" qatnashgan atama-ta'rif juftliklaridan to'g'risini toping.
\begin{enumerate}[label=\Alph*),leftmargin=9mm]
\item bottleneck - RAM va ikkilamchi xotira o'rtasida sahifalarni almashish
\item bottleneck - qizish sabab protsessor tezligining avtomatik pasayishi
\item bottleneck - tizim umumiy tezligini cheklayotgan eng sust resurs
\item bottleneck - belgi va dalillar asosida sababni tizimli aniqlash
\end{enumerate}
\begin{tcolorbox}[enhanced,breakable,colback=teal!4,colframe=teal!60!black,boxrule=0.45pt,arc=1.2mm,left=2.7mm,right=2.7mm,top=1.8mm,bottom=1.8mm,before skip=2mm,after skip=2mm]
{\bfseries\sffamily\color{teal}To'g'ri javob:} \textbf{C) bottleneck - tizim umumiy tezligini cheklayotgan eng sust resurs}
\par\smallskip
{\bfseries Izoh:} bottleneck - tizim umumiy tezligini cheklayotgan eng sust resurs. To'g'ri juftlik: bottleneck - tizim umumiy tezligini cheklayotgan eng sust resurs. Manba: CAM9, PDF 12-47; Qo'shimcha sahifa raqami (qaysi PDFga tegishli ekani aniqlanmagan; darslik sahifasi deb tasdiqlanmagan): 108
\end{tcolorbox}
\vspace{2mm}\hrule height 0.2pt\vspace{3mm}

% DB_ID bb6b5873-bb7a-4a6f-1f12-23ad95ff0ee3
\needspace{5\baselineskip}
{\bfseries\sffamily\color{ink} 141.} {scriptsizecolor{muted} Y1 | review | bb6b5873}
\par
Kompyuter tizimi va uning ishlash modeli mavzusidagi vazifada kiritish, qayta ishlash, chiqarish va saqlash modeli talab qilindi. Qaysi tushuncha mos?
\begin{enumerate}[label=\Alph*),leftmargin=9mm]
\item IPOS
\item kompyuter tizimi
\item texnik ta'minot
\item dasturiy ta'minot
\end{enumerate}
\begin{tcolorbox}[enhanced,breakable,colback=teal!4,colframe=teal!60!black,boxrule=0.45pt,arc=1.2mm,left=2.7mm,right=2.7mm,top=1.8mm,bottom=1.8mm,before skip=2mm,after skip=2mm]
{\bfseries\sffamily\color{teal}To'g'ri javob:} \textbf{A) IPOS}
\par\smallskip
{\bfseries Izoh:} IPOS. Berilgan vazifa IPOS tushunchasiga mos: kiritish, qayta ishlash, chiqarish va saqlash modeli. Manba: ICT10, PDF 6-12; Qo'shimcha sahifa raqami (qaysi PDFga tegishli ekani aniqlanmagan; darslik sahifasi deb tasdiqlanmagan): 6
\end{tcolorbox}
\vspace{2mm}\hrule height 0.2pt\vspace{3mm}

% DB_ID c3d1b32b-d825-9ca2-1dbf-ba4a69e45c82
\needspace{5\baselineskip}
{\bfseries\sffamily\color{ink} 142.} {scriptsizecolor{muted} Y1 | review | c3d1b32b}
\par
Kompyuter tizimi va uning ishlash modeli mavzusidagi vazifada bir maqsad uchun hamkorlik qiluvchi texnik vosita, dastur, ma'lumot va foydalanuvchi majmuasi talab qilindi. Qaysi tushuncha mos?
\begin{enumerate}[label=\Alph*),leftmargin=9mm]
\item texnik ta'minot
\item kompyuter tizimi
\item dasturiy ta'minot
\item IPOS
\end{enumerate}
\begin{tcolorbox}[enhanced,breakable,colback=teal!4,colframe=teal!60!black,boxrule=0.45pt,arc=1.2mm,left=2.7mm,right=2.7mm,top=1.8mm,bottom=1.8mm,before skip=2mm,after skip=2mm]
{\bfseries\sffamily\color{teal}To'g'ri javob:} \textbf{B) kompyuter tizimi}
\par\smallskip
{\bfseries Izoh:} kompyuter tizimi. Berilgan vazifa kompyuter tizimi tushunchasiga mos: bir maqsad uchun hamkorlik qiluvchi texnik vosita, dastur, ma'lumot va foydalanuvchi majmuasi. Manba: CAM9, PDF 8-15; Qo'shimcha sahifa raqami (qaysi PDFga tegishli ekani aniqlanmagan; darslik sahifasi deb tasdiqlanmagan): 3
\end{tcolorbox}
\vspace{2mm}\hrule height 0.2pt\vspace{3mm}

% DB_ID c90ea408-a8c0-4f34-5970-b872d00dda1b
\needspace{5\baselineskip}
{\bfseries\sffamily\color{ink} 143.} {scriptsizecolor{muted} Y1 | review | c90ea408}
\par
Ikkilamchi saqlash qurilmalari mavzusida "HDD" qatnashgan atama-ta'rif juftliklaridan to'g'risini toping.
\begin{enumerate}[label=\Alph*),leftmargin=9mm]
\item HDD - flesh hujayralarga asoslangan harakatsiz ikkilamchi xotira
\item HDD - lazer orqali o'qiladigan va yoziladigan disk
\item HDD - magnit plastinkali, katta sig'imli ikkilamchi xotira
\item HDD - ketma-ket murojaatli, arxiv uchun qulay magnit vosita
\end{enumerate}
\begin{tcolorbox}[enhanced,breakable,colback=teal!4,colframe=teal!60!black,boxrule=0.45pt,arc=1.2mm,left=2.7mm,right=2.7mm,top=1.8mm,bottom=1.8mm,before skip=2mm,after skip=2mm]
{\bfseries\sffamily\color{teal}To'g'ri javob:} \textbf{C) HDD - magnit plastinkali, katta sig'imli ikkilamchi xotira}
\par\smallskip
{\bfseries Izoh:} HDD - magnit plastinkali, katta sig'imli ikkilamchi xotira. To'g'ri juftlik: HDD - magnit plastinkali, katta sig'imli ikkilamchi xotira. Manba: CAM9, PDF 33-47; Qo'shimcha sahifa raqami (qaysi PDFga tegishli ekani aniqlanmagan; darslik sahifasi deb tasdiqlanmagan): 30
\end{tcolorbox}
\vspace{2mm}\hrule height 0.2pt\vspace{3mm}

% DB_ID cfe2c7bd-235a-688a-f8d1-b46e140a5150
\needspace{5\baselineskip}
{\bfseries\sffamily\color{ink} 144.} {scriptsizecolor{muted} Y1 | review | cfe2c7bd}
\par
Markaziy protsessor va bajarish sikli mavzusidagi vazifada sekunddagi takt impulslar soni talab qilindi. Qaysi tushuncha mos?
\begin{enumerate}[label=\Alph*),leftmargin=9mm]
\item CPU
\item takt chastotasi
\item ALU
\item CU
\end{enumerate}
\begin{tcolorbox}[enhanced,breakable,colback=teal!4,colframe=teal!60!black,boxrule=0.45pt,arc=1.2mm,left=2.7mm,right=2.7mm,top=1.8mm,bottom=1.8mm,before skip=2mm,after skip=2mm]
{\bfseries\sffamily\color{teal}To'g'ri javob:} \textbf{B) takt chastotasi}
\par\smallskip
{\bfseries Izoh:} takt chastotasi. Berilgan vazifa takt chastotasi tushunchasiga mos: sekunddagi takt impulslar soni. Manba: CAM9, PDF 10-14; Qo'shimcha sahifa raqami (qaysi PDFga tegishli ekani aniqlanmagan; darslik sahifasi deb tasdiqlanmagan): 21
\end{tcolorbox}
\vspace{2mm}\hrule height 0.2pt\vspace{3mm}

% DB_ID d66e6935-d098-fc9e-37ed-b203feeb3190
\needspace{5\baselineskip}
{\bfseries\sffamily\color{ink} 145.} {scriptsizecolor{muted} Y1 | review | d66e6935}
\par
Tizimli blok, bosh plata va ulanishlar mavzusida "slot" qatnashgan atama-ta'rif juftliklaridan to'g'risini toping.
\begin{enumerate}[label=\Alph*),leftmargin=9mm]
\item slot - ichki komponentlarni bog'lovchi asosiy bosma plata
\item slot - ma'lumot va boshqaruv signallarini tashuvchi aloqa yo'llari
\item slot - ichki komponent o'rnatiladigan ulagich
\item slot - tashqi qurilma ulanadigan nuqta
\end{enumerate}
\begin{tcolorbox}[enhanced,breakable,colback=teal!4,colframe=teal!60!black,boxrule=0.45pt,arc=1.2mm,left=2.7mm,right=2.7mm,top=1.8mm,bottom=1.8mm,before skip=2mm,after skip=2mm]
{\bfseries\sffamily\color{teal}To'g'ri javob:} \textbf{C) slot - ichki komponent o'rnatiladigan ulagich}
\par\smallskip
{\bfseries Izoh:} slot - ichki komponent o'rnatiladigan ulagich. To'g'ri juftlik: slot - ichki komponent o'rnatiladigan ulagich. Manba: ICT10, PDF 12-18; Qo'shimcha sahifa raqami (qaysi PDFga tegishli ekani aniqlanmagan; darslik sahifasi deb tasdiqlanmagan): 12
\end{tcolorbox}
\vspace{2mm}\hrule height 0.2pt\vspace{3mm}

% DB_ID d9138398-2cd7-9c21-ce31-0e1734b7fdd5
\needspace{5\baselineskip}
{\bfseries\sffamily\color{ink} 146.} {scriptsizecolor{muted} Y1 | review | d9138398}
\par
Ikkilamchi saqlash qurilmalari mavzusidagi vazifada ketma-ket murojaatli, arxiv uchun qulay magnit vosita talab qilindi. Qaysi tushuncha mos?
\begin{enumerate}[label=\Alph*),leftmargin=9mm]
\item magnit lenta
\item HDD
\item SSD
\item optik disk
\end{enumerate}
\begin{tcolorbox}[enhanced,breakable,colback=teal!4,colframe=teal!60!black,boxrule=0.45pt,arc=1.2mm,left=2.7mm,right=2.7mm,top=1.8mm,bottom=1.8mm,before skip=2mm,after skip=2mm]
{\bfseries\sffamily\color{teal}To'g'ri javob:} \textbf{A) magnit lenta}
\par\smallskip
{\bfseries Izoh:} magnit lenta. Berilgan vazifa magnit lenta tushunchasiga mos: ketmaket murojaatli, arxiv uchun qulay magnit vosita. Manba: ICT7, PDF 23-38; Qo'shimcha sahifa raqami (qaysi PDFga tegishli ekani aniqlanmagan; darslik sahifasi deb tasdiqlanmagan): 33
\end{tcolorbox}
\vspace{2mm}\hrule height 0.2pt\vspace{3mm}

% DB_ID e70a7e9c-5992-6f81-6fad-2cea679ee149
\needspace{5\baselineskip}
{\bfseries\sffamily\color{ink} 147.} {scriptsizecolor{muted} Y1 | review | e70a7e9c}
\par
Chiqarish va ikki vazifali qurilmalar mavzusida "lazer printer" qatnashgan atama-ta'rif juftliklaridan to'g'risini toping.
\begin{enumerate}[label=\Alph*),leftmargin=9mm]
\item lazer printer - toner va elektrofotografik jarayon bilan tez chop etuvchi printer
\item lazer printer - ekranda yoki audio orqali vaqtinchalik taqdim etilgan natija
\item lazer printer - qog'oz yoki boshqa fizik muhitdagi natija
\item lazer printer - mayda siyoh tomchilarini purkab tasvir hosil qiluvchi printer
\end{enumerate}
\begin{tcolorbox}[enhanced,breakable,colback=teal!4,colframe=teal!60!black,boxrule=0.45pt,arc=1.2mm,left=2.7mm,right=2.7mm,top=1.8mm,bottom=1.8mm,before skip=2mm,after skip=2mm]
{\bfseries\sffamily\color{teal}To'g'ri javob:} \textbf{A) lazer printer - toner va elektrofotografik jarayon bilan tez chop etuvchi printer}
\par\smallskip
{\bfseries Izoh:} lazer printer - toner va elektrofotografik jarayon bilan tez chop etuvchi printer. To'g'ri juftlik: lazer printer - toner va elektrofotografik jarayon bilan tez chop etuvchi printer. Manba: ICT5, PDF 29-36; Qo'shimcha sahifa raqami (qaysi PDFga tegishli ekani aniqlanmagan; darslik sahifasi deb tasdiqlanmagan): 44
\end{tcolorbox}
\vspace{2mm}\hrule height 0.2pt\vspace{3mm}

% DB_ID e7b3ed04-8692-ab23-5bfc-43ac02bb75f5
\needspace{5\baselineskip}
{\bfseries\sffamily\color{ink} 148.} {scriptsizecolor{muted} Y1 | review | e7b3ed04}
\par
Markaziy protsessor va bajarish sikli mavzusidagi vazifada CPU va RAM o'rtasidagi tezkor oraliq xotira talab qilindi. Qaysi tushuncha mos?
\begin{enumerate}[label=\Alph*),leftmargin=9mm]
\item CPU
\item ALU
\item kesh
\item CU
\end{enumerate}
\begin{tcolorbox}[enhanced,breakable,colback=teal!4,colframe=teal!60!black,boxrule=0.45pt,arc=1.2mm,left=2.7mm,right=2.7mm,top=1.8mm,bottom=1.8mm,before skip=2mm,after skip=2mm]
{\bfseries\sffamily\color{teal}To'g'ri javob:} \textbf{C) kesh}
\par\smallskip
{\bfseries Izoh:} kesh. Berilgan vazifa kesh tushunchasiga mos: CPU va RAM o'rtasidagi tezkor oraliq xotira. Manba: CAM10/11, PDF 32-34; Qo'shimcha sahifa raqami (qaysi PDFga tegishli ekani aniqlanmagan; darslik sahifasi deb tasdiqlanmagan): 22
\end{tcolorbox}
\vspace{2mm}\hrule height 0.2pt\vspace{3mm}

% DB_ID ef5398a7-01ff-83c6-9d16-2e11ee19328d
\needspace{5\baselineskip}
{\bfseries\sffamily\color{ink} 149.} {scriptsizecolor{muted} Y1 | review | ef5398a7}
\par
Chiqarish va ikki vazifali qurilmalar mavzusida "aktuator" qatnashgan atama-ta'rif juftliklaridan to'g'risini toping.
\begin{enumerate}[label=\Alph*),leftmargin=9mm]
\item aktuator - ekranda yoki audio orqali vaqtinchalik taqdim etilgan natija
\item aktuator - qog'oz yoki boshqa fizik muhitdagi natija
\item aktuator - boshqaruv signalini fizik harakat yoki ta'sirga aylantiruvchi qurilma
\item aktuator - toner va elektrofotografik jarayon bilan tez chop etuvchi printer
\end{enumerate}
\begin{tcolorbox}[enhanced,breakable,colback=teal!4,colframe=teal!60!black,boxrule=0.45pt,arc=1.2mm,left=2.7mm,right=2.7mm,top=1.8mm,bottom=1.8mm,before skip=2mm,after skip=2mm]
{\bfseries\sffamily\color{teal}To'g'ri javob:} \textbf{C) aktuator - boshqaruv signalini fizik harakat yoki ta'sirga aylantiruvchi qurilma}
\par\smallskip
{\bfseries Izoh:} aktuator - boshqaruv signalini fizik harakat yoki ta'sirga aylantiruvchi qurilma. To'g'ri juftlik: aktuator - boshqaruv signalini fizik harakat yoki ta'sirga aylantiruvchi qurilma. Manba: MAR, PDF 8-10; Qo'shimcha sahifa raqami (qaysi PDFga tegishli ekani aniqlanmagan; darslik sahifasi deb tasdiqlanmagan): 47
\end{tcolorbox}
\vspace{2mm}\hrule height 0.2pt\vspace{3mm}

% DB_ID ff50a630-454a-2484-3e48-5f7e95f69624
\needspace{5\baselineskip}
{\bfseries\sffamily\color{ink} 150.} {scriptsizecolor{muted} Y1 | review | ff50a630}
\par
Chiqarish va ikki vazifali qurilmalar mavzusidagi vazifada toner va elektrofotografik jarayon bilan tez chop etuvchi printer talab qilindi. Qaysi tushuncha mos?
\begin{enumerate}[label=\Alph*),leftmargin=9mm]
\item yumshoq nusxa
\item qattiq nusxa
\item siyohpurkagich
\item lazer printer
\end{enumerate}
\begin{tcolorbox}[enhanced,breakable,colback=teal!4,colframe=teal!60!black,boxrule=0.45pt,arc=1.2mm,left=2.7mm,right=2.7mm,top=1.8mm,bottom=1.8mm,before skip=2mm,after skip=2mm]
{\bfseries\sffamily\color{teal}To'g'ri javob:} \textbf{D) lazer printer}
\par\smallskip
{\bfseries Izoh:} lazer printer. Berilgan vazifa lazer printer tushunchasiga mos: toner va elektrofotografik jarayon bilan tez chop etuvchi printer. Manba: ICT5, PDF 29-36; Qo'shimcha sahifa raqami (qaysi PDFga tegishli ekani aniqlanmagan; darslik sahifasi deb tasdiqlanmagan): 44
\end{tcolorbox}
\vspace{2mm}\hrule height 0.2pt\vspace{3mm}

% DB_ID 1010df06-5ecf-cbd2-3adc-ecf683d7cd31
\needspace{5\baselineskip}
{\bfseries\sffamily\color{ink} 151.} {scriptsizecolor{muted} Y1 | review | 1010df06}
\par
Birlamchi xotira: registr, kesh, RAM va ROM bo'yicha qatorlar ichidan "firmware" noto'g'ri izohlanganini aniqlang.
\begin{enumerate}[label=\Alph*),leftmargin=9mm]
\item ROM - elektrsiz ham ma'lumotni saqlovchi o'zgarmas xotira
\item firmware - faol dastur va ma'lumotlar uchun o'zgaruvchan tezkor xotira
\item kesh - tez-tez kerak bo'ladigan ma'lumotni CPUga yaqin saqlovchi xotira
\item virtual xotira - RAM yetishmaganda ikkilamchi xotiradan ajratilgan sekinroq maydon
\end{enumerate}
\begin{tcolorbox}[enhanced,breakable,colback=teal!4,colframe=teal!60!black,boxrule=0.45pt,arc=1.2mm,left=2.7mm,right=2.7mm,top=1.8mm,bottom=1.8mm,before skip=2mm,after skip=2mm]
{\bfseries\sffamily\color{teal}To'g'ri javob:} \textbf{B) firmware - faol dastur va ma'lumotlar uchun o'zgaruvchan tezkor xotira}
\par\smallskip
{\bfseries Izoh:} firmware - faol dastur va ma'lumotlar uchun o'zgaruvchan tezkor xotira. firmware uchun to'g'ri ta'rif: qurilma ichiga o'rnatilgan past darajali boshqaruv dasturi. Manba: CAM9, PDF 11-14; Qo'shimcha sahifa raqami (qaysi PDFga tegishli ekani aniqlanmagan; darslik sahifasi deb tasdiqlanmagan): 28
\end{tcolorbox}
\vspace{2mm}\hrule height 0.2pt\vspace{3mm}

% DB_ID 11593eb6-bd28-cf54-d26f-03db58e03cf0
\needspace{5\baselineskip}
{\bfseries\sffamily\color{ink} 152.} {scriptsizecolor{muted} Y1 | review | 11593eb6}
\par
Birlamchi xotira: registr, kesh, RAM va ROM mavzusida "UEFI" qatnashgan atama-ta'rif juftliklaridan to'g'risini toping.
\begin{enumerate}[label=\Alph*),leftmargin=9mm]
\item UEFI - faol dastur va ma'lumotlar uchun o'zgaruvchan tezkor xotira
\item UEFI - qurilmani boshlang'ich sozlab, OT yuklanishini boshlovchi firmware muhiti
\item UEFI - elektrsiz ham ma'lumotni saqlovchi o'zgarmas xotira
\item UEFI - tez-tez kerak bo'ladigan ma'lumotni CPUga yaqin saqlovchi xotira
\end{enumerate}
\begin{tcolorbox}[enhanced,breakable,colback=teal!4,colframe=teal!60!black,boxrule=0.45pt,arc=1.2mm,left=2.7mm,right=2.7mm,top=1.8mm,bottom=1.8mm,before skip=2mm,after skip=2mm]
{\bfseries\sffamily\color{teal}To'g'ri javob:} \textbf{B) UEFI - qurilmani boshlang'ich sozlab, OT yuklanishini boshlovchi firmware muhiti}
\par\smallskip
{\bfseries Izoh:} UEFI - qurilmani boshlang'ich sozlab, OT yuklanishini boshlovchi firmware muhiti. To'g'ri juftlik: UEFI - qurilmani boshlang'ich sozlab, OT yuklanishini boshlovchi firmware muhiti. Manba: CAM10/11, PDF 33-36; Qo'shimcha sahifa raqami (qaysi PDFga tegishli ekani aniqlanmagan; darslik sahifasi deb tasdiqlanmagan): 29
\end{tcolorbox}
\vspace{2mm}\hrule height 0.2pt\vspace{3mm}

% DB_ID 15374c34-c0a9-93ac-4b60-943f85346eb6
\needspace{5\baselineskip}
{\bfseries\sffamily\color{ink} 153.} {scriptsizecolor{muted} Y1 | review | 15374c34}
\par
Kiritish qurilmalari mavzusida "OMR" qatnashgan atama-ta'rif juftliklaridan to'g'risini toping.
\begin{enumerate}[label=\Alph*),leftmargin=9mm]
\item OMR - tasvirdagi bosma matnni tahrirlanadigan belgilar sifatida tanish
\item OMR - oldindan belgilangan joylardagi belgilarni aniqlash
\item OMR - magnit siyoh bilan bosilgan maxsus belgilarni o'qish
\item OMR - chiziqli koddagi identifikatorni avtomatik kiritish
\end{enumerate}
\begin{tcolorbox}[enhanced,breakable,colback=teal!4,colframe=teal!60!black,boxrule=0.45pt,arc=1.2mm,left=2.7mm,right=2.7mm,top=1.8mm,bottom=1.8mm,before skip=2mm,after skip=2mm]
{\bfseries\sffamily\color{teal}To'g'ri javob:} \textbf{B) OMR - oldindan belgilangan joylardagi belgilarni aniqlash}
\par\smallskip
{\bfseries Izoh:} OMR - oldindan belgilangan joylardagi belgilarni aniqlash. To'g'ri juftlik: OMR - oldindan belgilangan joylardagi belgilarni aniqlash. Manba: CAM10/11, PDF 37-45; Qo'shimcha sahifa raqami (qaysi PDFga tegishli ekani aniqlanmagan; darslik sahifasi deb tasdiqlanmagan): 37
\end{tcolorbox}
\vspace{2mm}\hrule height 0.2pt\vspace{3mm}

% DB_ID 1749f1ee-350a-f5ef-75c9-a15e4b5141e0
\needspace{5\baselineskip}
{\bfseries\sffamily\color{ink} 154.} {scriptsizecolor{muted} Y1 | review | 1749f1ee}
\par
Attestatsiyadagi nozik farqlar va integratsion masalalar mavzusida "mezon" qatnashgan atama-ta'rif juftliklaridan to'g'risini toping.
\begin{enumerate}[label=\Alph*),leftmargin=9mm]
\item mezon - obyektni muhim belgilari bo'yicha guruhga ajratish
\item mezon - variantlarni baholash uchun qo'llanadigan belgi
\item mezon - yechim bajarishi shart bo'lgan chegara yoki talab
\item mezon - bir nechta mavzuni bitta vaziyatda bog'laydigan topshiriq
\end{enumerate}
\begin{tcolorbox}[enhanced,breakable,colback=teal!4,colframe=teal!60!black,boxrule=0.45pt,arc=1.2mm,left=2.7mm,right=2.7mm,top=1.8mm,bottom=1.8mm,before skip=2mm,after skip=2mm]
{\bfseries\sffamily\color{teal}To'g'ri javob:} \textbf{B) mezon - variantlarni baholash uchun qo'llanadigan belgi}
\par\smallskip
{\bfseries Izoh:} mezon - variantlarni baholash uchun qo'llanadigan belgi. To'g'ri juftlik: mezon - variantlarni baholash uchun qo'llanadigan belgi. Manba: CAM9, PDF 8-47; Qo'shimcha sahifa raqami (qaysi PDFga tegishli ekani aniqlanmagan; darslik sahifasi deb tasdiqlanmagan): 116
\end{tcolorbox}
\vspace{2mm}\hrule height 0.2pt\vspace{3mm}

% DB_ID 1b2ad97b-5ac0-5499-144d-1beca3bccffc
\needspace{5\baselineskip}
{\bfseries\sffamily\color{ink} 155.} {scriptsizecolor{muted} Y1 | review | 1b2ad97b}
\par
Kiritish qurilmalari mavzusida "shtrix-kod o'quvchi" qatnashgan atama-ta'rif juftliklaridan to'g'risini toping.
\begin{enumerate}[label=\Alph*),leftmargin=9mm]
\item shtrix-kod o'quvchi - tasvirdagi bosma matnni tahrirlanadigan belgilar sifatida tanish
\item shtrix-kod o'quvchi - oldindan belgilangan joylardagi belgilarni aniqlash
\item shtrix-kod o'quvchi - magnit siyoh bilan bosilgan maxsus belgilarni o'qish
\item shtrix-kod o'quvchi - chiziqli koddagi identifikatorni avtomatik kiritish
\end{enumerate}
\begin{tcolorbox}[enhanced,breakable,colback=teal!4,colframe=teal!60!black,boxrule=0.45pt,arc=1.2mm,left=2.7mm,right=2.7mm,top=1.8mm,bottom=1.8mm,before skip=2mm,after skip=2mm]
{\bfseries\sffamily\color{teal}To'g'ri javob:} \textbf{D) shtrix-kod o'quvchi - chiziqli koddagi identifikatorni avtomatik kiritish}
\par\smallskip
{\bfseries Izoh:} shtrix-kod o'quvchi - chiziqli koddagi identifikatorni avtomatik kiritish. To'g'ri juftlik: shtrix-kod o'quvchi - chiziqli koddagi identifikatorni avtomatik kiritish. Manba: ICT5, PDF 28-34; Qo'shimcha sahifa raqami (qaysi PDFga tegishli ekani aniqlanmagan; darslik sahifasi deb tasdiqlanmagan): 39
\end{tcolorbox}
\vspace{2mm}\hrule height 0.2pt\vspace{3mm}

% DB_ID 20e62b99-ad23-4fb7-a450-86162acc3151
\needspace{5\baselineskip}
{\bfseries\sffamily\color{ink} 156.} {scriptsizecolor{muted} Y1 | review | 20e62b99}
\par
Birlamchi xotira: registr, kesh, RAM va ROM bo'yicha qatorlar ichidan "kesh" noto'g'ri izohlanganini aniqlang.
\begin{enumerate}[label=\Alph*),leftmargin=9mm]
\item ROM - elektrsiz ham ma'lumotni saqlovchi o'zgarmas xotira
\item virtual xotira - RAM yetishmaganda ikkilamchi xotiradan ajratilgan sekinroq maydon
\item firmware - qurilma ichiga o'rnatilgan past darajali boshqaruv dasturi
\item kesh - faol dastur va ma'lumotlar uchun o'zgaruvchan tezkor xotira
\end{enumerate}
\begin{tcolorbox}[enhanced,breakable,colback=teal!4,colframe=teal!60!black,boxrule=0.45pt,arc=1.2mm,left=2.7mm,right=2.7mm,top=1.8mm,bottom=1.8mm,before skip=2mm,after skip=2mm]
{\bfseries\sffamily\color{teal}To'g'ri javob:} \textbf{D) kesh - faol dastur va ma'lumotlar uchun o'zgaruvchan tezkor xotira}
\par\smallskip
{\bfseries Izoh:} kesh - faol dastur va ma'lumotlar uchun o'zgaruvchan tezkor xotira. kesh uchun to'g'ri ta'rif: tez-tez kerak bo'ladigan ma'lumotni CPUga yaqin saqlovchi xotira. Manba: ICT10, PDF 15-20; Qo'shimcha sahifa raqami (qaysi PDFga tegishli ekani aniqlanmagan; darslik sahifasi deb tasdiqlanmagan): 26
\end{tcolorbox}
\vspace{2mm}\hrule height 0.2pt\vspace{3mm}

% DB_ID 2b2ccf8e-c609-0068-cdb2-12f20b5d6348
\needspace{5\baselineskip}
{\bfseries\sffamily\color{ink} 157.} {scriptsizecolor{muted} Y1 | review | 2b2ccf8e}
\par
Markaziy protsessor va bajarish sikli mavzusida "kesh" qatnashgan atama-ta'rif juftliklaridan to'g'risini toping.
\begin{enumerate}[label=\Alph*),leftmargin=9mm]
\item kesh - ko'rsatmalarni bajarib, ma'lumotlarni qayta ishlovchi markaziy qurilma
\item kesh - arifmetik va mantiqiy amallarni bajaruvchi qism
\item kesh - ko'rsatmalar oqimi va boshqa qismlar ishini boshqaruvchi qism
\item kesh - CPU va RAM o'rtasidagi tezkor oraliq xotira
\end{enumerate}
\begin{tcolorbox}[enhanced,breakable,colback=teal!4,colframe=teal!60!black,boxrule=0.45pt,arc=1.2mm,left=2.7mm,right=2.7mm,top=1.8mm,bottom=1.8mm,before skip=2mm,after skip=2mm]
{\bfseries\sffamily\color{teal}To'g'ri javob:} \textbf{D) kesh - CPU va RAM o'rtasidagi tezkor oraliq xotira}
\par\smallskip
{\bfseries Izoh:} kesh - CPU va RAM o'rtasidagi tezkor oraliq xotira. To'g'ri juftlik: kesh - CPU va RAM o'rtasidagi tezkor oraliq xotira. Manba: CAM10/11, PDF 32-34; Qo'shimcha sahifa raqami (qaysi PDFga tegishli ekani aniqlanmagan; darslik sahifasi deb tasdiqlanmagan): 22
\end{tcolorbox}
\vspace{2mm}\hrule height 0.2pt\vspace{3mm}

% DB_ID 2b748425-409b-6a78-5a6d-2abf37ebf002
\needspace{5\baselineskip}
{\bfseries\sffamily\color{ink} 158.} {scriptsizecolor{muted} Y1 | review | 2b748425}
\par
Birlamchi xotira: registr, kesh, RAM va ROM mavzusida "virtual xotira" qatnashgan atama-ta'rif juftliklaridan to'g'risini toping.
\begin{enumerate}[label=\Alph*),leftmargin=9mm]
\item virtual xotira - faol dastur va ma'lumotlar uchun o'zgaruvchan tezkor xotira
\item virtual xotira - elektrsiz ham ma'lumotni saqlovchi o'zgarmas xotira
\item virtual xotira - tez-tez kerak bo'ladigan ma'lumotni CPUga yaqin saqlovchi xotira
\item virtual xotira - RAM yetishmaganda ikkilamchi xotiradan ajratilgan sekinroq maydon
\end{enumerate}
\begin{tcolorbox}[enhanced,breakable,colback=teal!4,colframe=teal!60!black,boxrule=0.45pt,arc=1.2mm,left=2.7mm,right=2.7mm,top=1.8mm,bottom=1.8mm,before skip=2mm,after skip=2mm]
{\bfseries\sffamily\color{teal}To'g'ri javob:} \textbf{D) virtual xotira - RAM yetishmaganda ikkilamchi xotiradan ajratilgan sekinroq maydon}
\par\smallskip
{\bfseries Izoh:} virtual xotira - RAM yetishmaganda ikkilamchi xotiradan ajratilgan sekinroq maydon. To'g'ri juftlik: virtual xotira - RAM yetishmaganda ikkilamchi xotiradan ajratilgan sekinroq maydon. Manba: MAR, PDF 7-8; Qo'shimcha sahifa raqami (qaysi PDFga tegishli ekani aniqlanmagan; darslik sahifasi deb tasdiqlanmagan): 27
\end{tcolorbox}
\vspace{2mm}\hrule height 0.2pt\vspace{3mm}

% DB_ID 40533b7b-2115-0f91-ddc1-a8b5d462d802
\needspace{5\baselineskip}
{\bfseries\sffamily\color{ink} 159.} {scriptsizecolor{muted} Y1 | review | 40533b7b}
\par
Markaziy protsessor va bajarish sikli mavzusida "ALU" qatnashgan atama-ta'rif juftliklaridan to'g'risini toping.
\begin{enumerate}[label=\Alph*),leftmargin=9mm]
\item ALU - ko'rsatmalarni bajarib, ma'lumotlarni qayta ishlovchi markaziy qurilma
\item ALU - ko'rsatmalar oqimi va boshqa qismlar ishini boshqaruvchi qism
\item ALU - CPU ichidagi juda tez va kichik xotira
\item ALU - arifmetik va mantiqiy amallarni bajaruvchi qism
\end{enumerate}
\begin{tcolorbox}[enhanced,breakable,colback=teal!4,colframe=teal!60!black,boxrule=0.45pt,arc=1.2mm,left=2.7mm,right=2.7mm,top=1.8mm,bottom=1.8mm,before skip=2mm,after skip=2mm]
{\bfseries\sffamily\color{teal}To'g'ri javob:} \textbf{D) ALU - arifmetik va mantiqiy amallarni bajaruvchi qism}
\par\smallskip
{\bfseries Izoh:} ALU - arifmetik va mantiqiy amallarni bajaruvchi qism. To'g'ri juftlik: ALU - arifmetik va mantiqiy amallarni bajaruvchi qism. Manba: CAM10/11, PDF 32-34; Qo'shimcha sahifa raqami (qaysi PDFga tegishli ekani aniqlanmagan; darslik sahifasi deb tasdiqlanmagan): 18
\end{tcolorbox}
\vspace{2mm}\hrule height 0.2pt\vspace{3mm}

% DB_ID 405a3300-77d5-59bc-5dd0-04b87f4fbd35
\needspace{5\baselineskip}
{\bfseries\sffamily\color{ink} 160.} {scriptsizecolor{muted} Y1 | review | 405a3300}
\par
Unumdorlik, nosozliklarni aniqlash va xavfsiz xizmat mavzusida "profilaktika" qatnashgan atama-ta'rif juftliklaridan to'g'risini toping.
\begin{enumerate}[label=\Alph*),leftmargin=9mm]
\item profilaktika - tizim umumiy tezligini cheklayotgan eng sust resurs
\item profilaktika - RAM va ikkilamchi xotira o'rtasida sahifalarni almashish
\item profilaktika - qizish sabab protsessor tezligining avtomatik pasayishi
\item profilaktika - nosozlik yuz berish ehtimolini oldindan kamaytiruvchi xizmat
\end{enumerate}
\begin{tcolorbox}[enhanced,breakable,colback=teal!4,colframe=teal!60!black,boxrule=0.45pt,arc=1.2mm,left=2.7mm,right=2.7mm,top=1.8mm,bottom=1.8mm,before skip=2mm,after skip=2mm]
{\bfseries\sffamily\color{teal}To'g'ri javob:} \textbf{D) profilaktika - nosozlik yuz berish ehtimolini oldindan kamaytiruvchi xizmat}
\par\smallskip
{\bfseries Izoh:} profilaktika - nosozlik yuz berish ehtimolini oldindan kamaytiruvchi xizmat. To'g'ri juftlik: profilaktika - nosozlik yuz berish ehtimolini oldindan kamaytiruvchi xizmat. Manba: CAM10/11, PDF 24-61; Qo'shimcha sahifa raqami (qaysi PDFga tegishli ekani aniqlanmagan; darslik sahifasi deb tasdiqlanmagan): 114
\end{tcolorbox}
\vspace{2mm}\hrule height 0.2pt\vspace{3mm}

% DB_ID 4674acc7-2c33-4593-049f-2ef9e348d8a3
\needspace{5\baselineskip}
{\bfseries\sffamily\color{ink} 161.} {scriptsizecolor{muted} Y1 | review | 4674acc7}
\par
Tizimli blok, bosh plata va ulanishlar mavzusida "port" qatnashgan atama-ta'rif juftliklaridan to'g'risini toping.
\begin{enumerate}[label=\Alph*),leftmargin=9mm]
\item port - ichki komponentlarni bog'lovchi asosiy bosma plata
\item port - ma'lumot va boshqaruv signallarini tashuvchi aloqa yo'llari
\item port - ichki komponent o'rnatiladigan ulagich
\item port - tashqi qurilma ulanadigan nuqta
\end{enumerate}
\begin{tcolorbox}[enhanced,breakable,colback=teal!4,colframe=teal!60!black,boxrule=0.45pt,arc=1.2mm,left=2.7mm,right=2.7mm,top=1.8mm,bottom=1.8mm,before skip=2mm,after skip=2mm]
{\bfseries\sffamily\color{teal}To'g'ri javob:} \textbf{D) port - tashqi qurilma ulanadigan nuqta}
\par\smallskip
{\bfseries Izoh:} port - tashqi qurilma ulanadigan nuqta. To'g'ri juftlik: port - tashqi qurilma ulanadigan nuqta. Manba: MAR, PDF 5-8; Qo'shimcha sahifa raqami (qaysi PDFga tegishli ekani aniqlanmagan; darslik sahifasi deb tasdiqlanmagan): 13
\end{tcolorbox}
\vspace{2mm}\hrule height 0.2pt\vspace{3mm}

% DB_ID 53944749-0db5-5db7-5aa4-f78309844deb
\needspace{5\baselineskip}
{\bfseries\sffamily\color{ink} 162.} {scriptsizecolor{muted} Y1 | review | 53944749}
\par
Kompyuter tizimi va uning ishlash modeli mavzusida "umumiy maqsadli tizim" qatnashgan atama-ta'rif juftliklaridan to'g'risini toping.
\begin{enumerate}[label=\Alph*),leftmargin=9mm]
\item umumiy maqsadli tizim - bir maqsad uchun hamkorlik qiluvchi texnik vosita, dastur, ma'lumot va foydalanuvchi majmuasi
\item umumiy maqsadli tizim - kompyuter tizimining fizik qismlari
\item umumiy maqsadli tizim - kompyuterga vazifalarni bajarishni ko'rsatuvchi dasturlar
\item umumiy maqsadli tizim - turli dasturlar yordamida ko'p xil vazifani bajaradigan kompyuter tizimi
\end{enumerate}
\begin{tcolorbox}[enhanced,breakable,colback=teal!4,colframe=teal!60!black,boxrule=0.45pt,arc=1.2mm,left=2.7mm,right=2.7mm,top=1.8mm,bottom=1.8mm,before skip=2mm,after skip=2mm]
{\bfseries\sffamily\color{teal}To'g'ri javob:} \textbf{D) umumiy maqsadli tizim - turli dasturlar yordamida ko'p xil vazifani bajaradigan kompyuter tizimi}
\par\smallskip
{\bfseries Izoh:} umumiy maqsadli tizim - turli dasturlar yordamida ko'p xil vazifani bajaradigan kompyuter tizimi. To'g'ri juftlik: umumiy maqsadli tizim - turli dasturlar yordamida ko'p xil vazifani bajaradigan kompyuter tizimi. Manba: MAR, PDF 4-10; Qo'shimcha sahifa raqami (qaysi PDFga tegishli ekani aniqlanmagan; darslik sahifasi deb tasdiqlanmagan): 9
\end{tcolorbox}
\vspace{2mm}\hrule height 0.2pt\vspace{3mm}

% DB_ID 55f88ca1-0110-6f90-c34e-eb476221f582
\needspace{5\baselineskip}
{\bfseries\sffamily\color{ink} 163.} {scriptsizecolor{muted} Y1 | review | 55f88ca1}
\par
Chiqarish va ikki vazifali qurilmalar bo'yicha qatorlar ichidan "yumshoq nusxa" noto'g'ri izohlanganini aniqlang.
\begin{enumerate}[label=\Alph*),leftmargin=9mm]
\item lazer printer - toner va elektrofotografik jarayon bilan tez chop etuvchi printer
\item siyohpurkagich - mayda siyoh tomchilarini purkab tasvir hosil qiluvchi printer
\item aktuator - boshqaruv signalini fizik harakat yoki ta'sirga aylantiruvchi qurilma
\item yumshoq nusxa - qog'oz yoki boshqa fizik muhitdagi natija
\end{enumerate}
\begin{tcolorbox}[enhanced,breakable,colback=teal!4,colframe=teal!60!black,boxrule=0.45pt,arc=1.2mm,left=2.7mm,right=2.7mm,top=1.8mm,bottom=1.8mm,before skip=2mm,after skip=2mm]
{\bfseries\sffamily\color{teal}To'g'ri javob:} \textbf{D) yumshoq nusxa - qog'oz yoki boshqa fizik muhitdagi natija}
\par\smallskip
{\bfseries Izoh:} yumshoq nusxa - qog'oz yoki boshqa fizik muhitdagi natija. yumshoq nusxa uchun to'g'ri ta'rif: ekranda yoki audio orqali vaqtinchalik taqdim etilgan natija. Manba: CAM9, PDF 16-32; Qo'shimcha sahifa raqami (qaysi PDFga tegishli ekani aniqlanmagan; darslik sahifasi deb tasdiqlanmagan): 42
\end{tcolorbox}
\vspace{2mm}\hrule height 0.2pt\vspace{3mm}

% DB_ID 567a9b49-4b03-c7af-97b1-38f12e4fe0c1
\needspace{5\baselineskip}
{\bfseries\sffamily\color{ink} 164.} {scriptsizecolor{muted} Y1 | review | 567a9b49}
\par
Attestatsiyadagi nozik farqlar va integratsion masalalar bo'yicha qatorlar ichidan "cheklov" noto'g'ri izohlanganini aniqlang.
\begin{enumerate}[label=\Alph*),leftmargin=9mm]
\item mezon - variantlarni baholash uchun qo'llanadigan belgi
\item integratsion savol - bir nechta mavzuni bitta vaziyatda bog'laydigan topshiriq
\item distraktor - ishonarli, ammo noto'g'ri javob varianti
\item cheklov - obyektni muhim belgilari bo'yicha guruhga ajratish
\end{enumerate}
\begin{tcolorbox}[enhanced,breakable,colback=teal!4,colframe=teal!60!black,boxrule=0.45pt,arc=1.2mm,left=2.7mm,right=2.7mm,top=1.8mm,bottom=1.8mm,before skip=2mm,after skip=2mm]
{\bfseries\sffamily\color{teal}To'g'ri javob:} \textbf{D) cheklov - obyektni muhim belgilari bo'yicha guruhga ajratish}
\par\smallskip
{\bfseries Izoh:} cheklov - obyektni muhim belgilari bo'yicha guruhga ajratish. cheklov uchun to'g'ri ta'rif: yechim bajarishi shart bo'lgan chegara yoki talab. Manba: CAM10/11, PDF 23-61; Qo'shimcha sahifa raqami (qaysi PDFga tegishli ekani aniqlanmagan; darslik sahifasi deb tasdiqlanmagan): 117
\end{tcolorbox}
\vspace{2mm}\hrule height 0.2pt\vspace{3mm}

% DB_ID 569ae8d2-40e7-d430-5fd3-7dd4c5952b99
\needspace{5\baselineskip}
{\bfseries\sffamily\color{ink} 165.} {scriptsizecolor{muted} Y1 | review | 569ae8d2}
\par
Tizimli blok, bosh plata va ulanishlar mavzusida "shina" qatnashgan atama-ta'rif juftliklaridan to'g'risini toping.
\begin{enumerate}[label=\Alph*),leftmargin=9mm]
\item shina - ichki komponentlarni bog'lovchi asosiy bosma plata
\item shina - ma'lumot va boshqaruv signallarini tashuvchi aloqa yo'llari
\item shina - ichki komponent o'rnatiladigan ulagich
\item shina - tashqi qurilma ulanadigan nuqta
\end{enumerate}
\begin{tcolorbox}[enhanced,breakable,colback=teal!4,colframe=teal!60!black,boxrule=0.45pt,arc=1.2mm,left=2.7mm,right=2.7mm,top=1.8mm,bottom=1.8mm,before skip=2mm,after skip=2mm]
{\bfseries\sffamily\color{teal}To'g'ri javob:} \textbf{B) shina - ma'lumot va boshqaruv signallarini tashuvchi aloqa yo'llari}
\par\smallskip
{\bfseries Izoh:} shina - ma'lumot va boshqaruv signallarini tashuvchi aloqa yo'llari. To'g'ri juftlik: shina - ma'lumot va boshqaruv signallarini tashuvchi aloqa yo'llari. Manba: CAM10/11, PDF 32-36; Qo'shimcha sahifa raqami (qaysi PDFga tegishli ekani aniqlanmagan; darslik sahifasi deb tasdiqlanmagan): 11
\end{tcolorbox}
\vspace{2mm}\hrule height 0.2pt\vspace{3mm}

% DB_ID 5e4a380f-32c6-a3e9-f47c-2de81b87f222
\needspace{5\baselineskip}
{\bfseries\sffamily\color{ink} 166.} {scriptsizecolor{muted} Y1 | review | 5e4a380f}
\par
1 TB = 1024 GB deb olinsin. 256 GBli to'rtta SSDning jami sig'imi qancha?
\begin{enumerate}[label=\Alph*),leftmargin=9mm]
\item 1 TB
\item 512 GB
\item 2 TB
\item 1024 TB
\end{enumerate}
\begin{tcolorbox}[enhanced,breakable,colback=teal!4,colframe=teal!60!black,boxrule=0.45pt,arc=1.2mm,left=2.7mm,right=2.7mm,top=1.8mm,bottom=1.8mm,before skip=2mm,after skip=2mm]
{\bfseries\sffamily\color{teal}To'g'ri javob:} \textbf{A) 1 TB}
\par\smallskip
{\bfseries Izoh:} 1 TB. 4 \ensuremath{\times} 256 GB = 1024 GB = 1 TB. Manba: CAM9, PDF 33-47; Qo'shimcha sahifa raqami (qaysi PDFga tegishli ekani aniqlanmagan; darslik sahifasi deb tasdiqlanmagan): 35
\end{tcolorbox}
\vspace{2mm}\hrule height 0.2pt\vspace{3mm}

% DB_ID 604db54e-ead6-d7a6-b2a2-5e142e2d5647
\needspace{5\baselineskip}
{\bfseries\sffamily\color{ink} 167.} {scriptsizecolor{muted} Y1 | review | 604db54e}
\par
Unumdorlik, nosozliklarni aniqlash va xavfsiz xizmat mavzusida "diagnostika" qatnashgan atama-ta'rif juftliklaridan to'g'risini toping.
\begin{enumerate}[label=\Alph*),leftmargin=9mm]
\item diagnostika - tizim umumiy tezligini cheklayotgan eng sust resurs
\item diagnostika - belgi va dalillar asosida sababni tizimli aniqlash
\item diagnostika - RAM va ikkilamchi xotira o'rtasida sahifalarni almashish
\item diagnostika - qizish sabab protsessor tezligining avtomatik pasayishi
\end{enumerate}
\begin{tcolorbox}[enhanced,breakable,colback=teal!4,colframe=teal!60!black,boxrule=0.45pt,arc=1.2mm,left=2.7mm,right=2.7mm,top=1.8mm,bottom=1.8mm,before skip=2mm,after skip=2mm]
{\bfseries\sffamily\color{teal}To'g'ri javob:} \textbf{B) diagnostika - belgi va dalillar asosida sababni tizimli aniqlash}
\par\smallskip
{\bfseries Izoh:} diagnostika - belgi va dalillar asosida sababni tizimli aniqlash. To'g'ri juftlik: diagnostika - belgi va dalillar asosida sababni tizimli aniqlash. Manba: MAR, PDF 10-20; Qo'shimcha sahifa raqami (qaysi PDFga tegishli ekani aniqlanmagan; darslik sahifasi deb tasdiqlanmagan): 112
\end{tcolorbox}
\vspace{2mm}\hrule height 0.2pt\vspace{3mm}

% DB_ID 6076a5b5-bf2c-ac9a-7aae-5b699bc10589
\needspace{5\baselineskip}
{\bfseries\sffamily\color{ink} 168.} {scriptsizecolor{muted} Y1 | review | 6076a5b5}
\par
Kompyuter tizimi va uning ishlash modeli bo'yicha qatorlar ichidan "kompyuter tizimi" noto'g'ri izohlanganini aniqlang.
\begin{enumerate}[label=\Alph*),leftmargin=9mm]
\item dasturiy ta'minot - kompyuterga vazifalarni bajarishni ko'rsatuvchi dasturlar
\item IPOS - kiritish, qayta ishlash, chiqarish va saqlash modeli
\item o'rnatilgan tizim - kattaroq qurilma ichida muayyan vazifani bajaruvchi kompyuter tizimi
\item kompyuter tizimi - kompyuter tizimining fizik qismlari
\end{enumerate}
\begin{tcolorbox}[enhanced,breakable,colback=teal!4,colframe=teal!60!black,boxrule=0.45pt,arc=1.2mm,left=2.7mm,right=2.7mm,top=1.8mm,bottom=1.8mm,before skip=2mm,after skip=2mm]
{\bfseries\sffamily\color{teal}To'g'ri javob:} \textbf{D) kompyuter tizimi - kompyuter tizimining fizik qismlari}
\par\smallskip
{\bfseries Izoh:} kompyuter tizimi - kompyuter tizimining fizik qismlari. kompyuter tizimi uchun to'g'ri ta'rif: bir maqsad uchun hamkorlik qiluvchi texnik vosita, dastur, ma'lumot va foydalanuvchi majmuasi. Manba: CAM9, PDF 8-15; Qo'shimcha sahifa raqami (qaysi PDFga tegishli ekani aniqlanmagan; darslik sahifasi deb tasdiqlanmagan): 3
\end{tcolorbox}
\vspace{2mm}\hrule height 0.2pt\vspace{3mm}

% DB_ID 6140479b-9ba8-3de6-7078-b0adf20d0eb4
\needspace{5\baselineskip}
{\bfseries\sffamily\color{ink} 169.} {scriptsizecolor{muted} Y1 | review | 6140479b}
\par
Tizimli blok, bosh plata va ulanishlar mavzusida "kengaytirish kartasi" qatnashgan atama-ta'rif juftliklaridan to'g'risini toping.
\begin{enumerate}[label=\Alph*),leftmargin=9mm]
\item kengaytirish kartasi - ichki komponentlarni bog'lovchi asosiy bosma plata
\item kengaytirish kartasi - kompyuterga qo'shimcha funksiya berish uchun slotga o'rnatiladigan plata
\item kengaytirish kartasi - ma'lumot va boshqaruv signallarini tashuvchi aloqa yo'llari
\item kengaytirish kartasi - ichki komponent o'rnatiladigan ulagich
\end{enumerate}
\begin{tcolorbox}[enhanced,breakable,colback=teal!4,colframe=teal!60!black,boxrule=0.45pt,arc=1.2mm,left=2.7mm,right=2.7mm,top=1.8mm,bottom=1.8mm,before skip=2mm,after skip=2mm]
{\bfseries\sffamily\color{teal}To'g'ri javob:} \textbf{B) kengaytirish kartasi - kompyuterga qo'shimcha funksiya berish uchun slotga o'rnatiladigan plata}
\par\smallskip
{\bfseries Izoh:} kengaytirish kartasi - kompyuterga qo'shimcha funksiya berish uchun slotga o'rnatiladigan plata. To'g'ri juftlik: kengaytirish kartasi - kompyuterga qo'shimcha funksiya berish uchun slotga o'rnatiladigan plata. Manba: CAM10/11, PDF 32-36; Qo'shimcha sahifa raqami (qaysi PDFga tegishli ekani aniqlanmagan; darslik sahifasi deb tasdiqlanmagan): 15
\end{tcolorbox}
\vspace{2mm}\hrule height 0.2pt\vspace{3mm}

% DB_ID 65524f73-461c-97a5-16a0-e4cc7a91c8c4
\needspace{5\baselineskip}
{\bfseries\sffamily\color{ink} 170.} {scriptsizecolor{muted} Y1 | review | 65524f73}
\par
Kiritish qurilmalari bo'yicha qatorlar ichidan "sensor" noto'g'ri izohlanganini aniqlang.
\begin{enumerate}[label=\Alph*),leftmargin=9mm]
\item OMR - oldindan belgilangan joylardagi belgilarni aniqlash
\item sensor - tasvirdagi bosma matnni tahrirlanadigan belgilar sifatida tanish
\item MICR - magnit siyoh bilan bosilgan maxsus belgilarni o'qish
\item shtrix-kod o'quvchi - chiziqli koddagi identifikatorni avtomatik kiritish
\end{enumerate}
\begin{tcolorbox}[enhanced,breakable,colback=teal!4,colframe=teal!60!black,boxrule=0.45pt,arc=1.2mm,left=2.7mm,right=2.7mm,top=1.8mm,bottom=1.8mm,before skip=2mm,after skip=2mm]
{\bfseries\sffamily\color{teal}To'g'ri javob:} \textbf{B) sensor - tasvirdagi bosma matnni tahrirlanadigan belgilar sifatida tanish}
\par\smallskip
{\bfseries Izoh:} sensor - tasvirdagi bosma matnni tahrirlanadigan belgilar sifatida tanish. sensor uchun to'g'ri ta'rif: fizik kattalikni o'lchab signalga aylantiruvchi qurilma. Manba: ICT10, PDF 8-12; Qo'shimcha sahifa raqami (qaysi PDFga tegishli ekani aniqlanmagan; darslik sahifasi deb tasdiqlanmagan): 40
\end{tcolorbox}
\vspace{2mm}\hrule height 0.2pt\vspace{3mm}

% DB_ID 662e3716-bcae-bdd6-622e-06b38904deec
\needspace{5\baselineskip}
{\bfseries\sffamily\color{ink} 171.} {scriptsizecolor{muted} Y1 | review | 662e3716}
\par
Attestatsiyadagi nozik farqlar va integratsion masalalar bo'yicha qatorlar ichidan "distraktor" noto'g'ri izohlanganini aniqlang.
\begin{enumerate}[label=\Alph*),leftmargin=9mm]
\item mezon - variantlarni baholash uchun qo'llanadigan belgi
\item distraktor - obyektni muhim belgilari bo'yicha guruhga ajratish
\item cheklov - yechim bajarishi shart bo'lgan chegara yoki talab
\item integratsion savol - bir nechta mavzuni bitta vaziyatda bog'laydigan topshiriq
\end{enumerate}
\begin{tcolorbox}[enhanced,breakable,colback=teal!4,colframe=teal!60!black,boxrule=0.45pt,arc=1.2mm,left=2.7mm,right=2.7mm,top=1.8mm,bottom=1.8mm,before skip=2mm,after skip=2mm]
{\bfseries\sffamily\color{teal}To'g'ri javob:} \textbf{B) distraktor - obyektni muhim belgilari bo'yicha guruhga ajratish}
\par\smallskip
{\bfseries Izoh:} distraktor - obyektni muhim belgilari bo'yicha guruhga ajratish. distraktor uchun to'g'ri ta'rif: ishonarli, ammo noto'g'ri javob varianti. Manba: MAR, PDF 4-20; Qo'shimcha sahifa raqami (qaysi PDFga tegishli ekani aniqlanmagan; darslik sahifasi deb tasdiqlanmagan): 120
\end{tcolorbox}
\vspace{2mm}\hrule height 0.2pt\vspace{3mm}

% DB_ID 6cc799ee-2e14-3b2f-5af0-a164e9e6d4ca
\needspace{5\baselineskip}
{\bfseries\sffamily\color{ink} 172.} {scriptsizecolor{muted} Y1 | review | 6cc799ee}
\par
Tizimli blok, bosh plata va ulanishlar bo'yicha qatorlar ichidan "slot" noto'g'ri izohlanganini aniqlang.
\begin{enumerate}[label=\Alph*),leftmargin=9mm]
\item shina - ma'lumot va boshqaruv signallarini tashuvchi aloqa yo'llari
\item port - tashqi qurilma ulanadigan nuqta
\item quvvat bloki - tarmoq elektrini komponentlarga mos doimiy tokka aylantiruvchi qism
\item slot - ichki komponentlarni bog'lovchi asosiy bosma plata
\end{enumerate}
\begin{tcolorbox}[enhanced,breakable,colback=teal!4,colframe=teal!60!black,boxrule=0.45pt,arc=1.2mm,left=2.7mm,right=2.7mm,top=1.8mm,bottom=1.8mm,before skip=2mm,after skip=2mm]
{\bfseries\sffamily\color{teal}To'g'ri javob:} \textbf{D) slot - ichki komponentlarni bog'lovchi asosiy bosma plata}
\par\smallskip
{\bfseries Izoh:} slot - ichki komponentlarni bog'lovchi asosiy bosma plata. slot uchun to'g'ri ta'rif: ichki komponent o'rnatiladigan ulagich. Manba: ICT10, PDF 12-18; Qo'shimcha sahifa raqami (qaysi PDFga tegishli ekani aniqlanmagan; darslik sahifasi deb tasdiqlanmagan): 12
\end{tcolorbox}
\vspace{2mm}\hrule height 0.2pt\vspace{3mm}

% DB_ID 75acf185-1bf7-da8c-7ff0-f7e848a9550c
\needspace{5\baselineskip}
{\bfseries\sffamily\color{ink} 173.} {scriptsizecolor{muted} Y1 | review | 75acf185}
\par
Ikkilamchi saqlash qurilmalari bo'yicha qatorlar ichidan "optik disk" noto'g'ri izohlanganini aniqlang.
\begin{enumerate}[label=\Alph*),leftmargin=9mm]
\item SSD - flesh hujayralarga asoslangan harakatsiz ikkilamchi xotira
\item optik disk - magnit plastinkali, katta sig'imli ikkilamchi xotira
\item magnit lenta - ketma-ket murojaatli, arxiv uchun qulay magnit vosita
\item USB flesh - portativ solid-state saqlash qurilmasi
\end{enumerate}
\begin{tcolorbox}[enhanced,breakable,colback=teal!4,colframe=teal!60!black,boxrule=0.45pt,arc=1.2mm,left=2.7mm,right=2.7mm,top=1.8mm,bottom=1.8mm,before skip=2mm,after skip=2mm]
{\bfseries\sffamily\color{teal}To'g'ri javob:} \textbf{B) optik disk - magnit plastinkali, katta sig'imli ikkilamchi xotira}
\par\smallskip
{\bfseries Izoh:} optik disk - magnit plastinkali, katta sig'imli ikkilamchi xotira. optik disk uchun to'g'ri ta'rif: lazer orqali o'qiladigan va yoziladigan disk. Manba: ICT10, PDF 14-18; Qo'shimcha sahifa raqami (qaysi PDFga tegishli ekani aniqlanmagan; darslik sahifasi deb tasdiqlanmagan): 32
\end{tcolorbox}
\vspace{2mm}\hrule height 0.2pt\vspace{3mm}

% DB_ID 7961fbec-7df0-abf7-d55e-6ab032a3f22d
\needspace{5\baselineskip}
{\bfseries\sffamily\color{ink} 174.} {scriptsizecolor{muted} Y1 | review | 7961fbec}
\par
Attestatsiyadagi nozik farqlar va integratsion masalalar bo'yicha qatorlar ichidan "tasnif" noto'g'ri izohlanganini aniqlang.
\begin{enumerate}[label=\Alph*),leftmargin=9mm]
\item cheklov - yechim bajarishi shart bo'lgan chegara yoki talab
\item tasnif - variantlarni baholash uchun qo'llanadigan belgi
\item integratsion savol - bir nechta mavzuni bitta vaziyatda bog'laydigan topshiriq
\item distraktor - ishonarli, ammo noto'g'ri javob varianti
\end{enumerate}
\begin{tcolorbox}[enhanced,breakable,colback=teal!4,colframe=teal!60!black,boxrule=0.45pt,arc=1.2mm,left=2.7mm,right=2.7mm,top=1.8mm,bottom=1.8mm,before skip=2mm,after skip=2mm]
{\bfseries\sffamily\color{teal}To'g'ri javob:} \textbf{B) tasnif - variantlarni baholash uchun qo'llanadigan belgi}
\par\smallskip
{\bfseries Izoh:} tasnif - variantlarni baholash uchun qo'llanadigan belgi. tasnif uchun to'g'ri ta'rif: obyektni muhim belgilari bo'yicha guruhga ajratish. Manba: SPEC, PDF 3; Qo'shimcha sahifa raqami (qaysi PDFga tegishli ekani aniqlanmagan; darslik sahifasi deb tasdiqlanmagan): 115
\end{tcolorbox}
\vspace{2mm}\hrule height 0.2pt\vspace{3mm}

% DB_ID 7988093b-91ca-99e3-d573-e6d1298ae73f
\needspace{5\baselineskip}
{\bfseries\sffamily\color{ink} 175.} {scriptsizecolor{muted} Y1 | review | 7988093b}
\par
Ikkilamchi saqlash qurilmalari bo'yicha qatorlar ichidan "USB flesh" noto'g'ri izohlanganini aniqlang.
\begin{enumerate}[label=\Alph*),leftmargin=9mm]
\item SSD - flesh hujayralarga asoslangan harakatsiz ikkilamchi xotira
\item optik disk - lazer orqali o'qiladigan va yoziladigan disk
\item magnit lenta - ketma-ket murojaatli, arxiv uchun qulay magnit vosita
\item USB flesh - magnit plastinkali, katta sig'imli ikkilamchi xotira
\end{enumerate}
\begin{tcolorbox}[enhanced,breakable,colback=teal!4,colframe=teal!60!black,boxrule=0.45pt,arc=1.2mm,left=2.7mm,right=2.7mm,top=1.8mm,bottom=1.8mm,before skip=2mm,after skip=2mm]
{\bfseries\sffamily\color{teal}To'g'ri javob:} \textbf{D) USB flesh - magnit plastinkali, katta sig'imli ikkilamchi xotira}
\par\smallskip
{\bfseries Izoh:} USB flesh - magnit plastinkali, katta sig'imli ikkilamchi xotira. USB flesh uchun to'g'ri ta'rif: portativ solid-state saqlash qurilmasi. Manba: MAR, PDF 7-9; Qo'shimcha sahifa raqami (qaysi PDFga tegishli ekani aniqlanmagan; darslik sahifasi deb tasdiqlanmagan): 34
\end{tcolorbox}
\vspace{2mm}\hrule height 0.2pt\vspace{3mm}

% DB_ID 8059ca4b-be2a-a9fe-f3ee-1e4d4c787014
\needspace{5\baselineskip}
{\bfseries\sffamily\color{ink} 176.} {scriptsizecolor{muted} Y1 | review | 8059ca4b}
\par
Kompyuter yuklama ostida qizib, avtomatik sekinlashmoqda. Eng bevosita tekshiruv qaysi?
\begin{enumerate}[label=\Alph*),leftmargin=9mm]
\item fayllarning kengaytmalari
\item papkalarning nomlanishi
\item printer siyohi darajasi
\item radiator, ventilyator va havo oqimi
\end{enumerate}
\begin{tcolorbox}[enhanced,breakable,colback=teal!4,colframe=teal!60!black,boxrule=0.45pt,arc=1.2mm,left=2.7mm,right=2.7mm,top=1.8mm,bottom=1.8mm,before skip=2mm,after skip=2mm]
{\bfseries\sffamily\color{teal}To'g'ri javob:} \textbf{D) radiator, ventilyator va havo oqimi}
\par\smallskip
{\bfseries Izoh:} radiator, ventilyator va havo oqimi. Qizish va termal sekinlashish avvalo sovitish tizimi bilan bog'lanadi. Manba: CAM9, PDF 9-14; Qo'shimcha sahifa raqami (qaysi PDFga tegishli ekani aniqlanmagan; darslik sahifasi deb tasdiqlanmagan): 16
\end{tcolorbox}
\vspace{2mm}\hrule height 0.2pt\vspace{3mm}

% DB_ID 83ce8eff-6d8f-2678-485f-de74d24ab94a
\needspace{5\baselineskip}
{\bfseries\sffamily\color{ink} 177.} {scriptsizecolor{muted} Y1 | review | 83ce8eff}
\par
Ikkilamchi saqlash qurilmalari mavzusida "SSD" qatnashgan atama-ta'rif juftliklaridan to'g'risini toping.
\begin{enumerate}[label=\Alph*),leftmargin=9mm]
\item SSD - magnit plastinkali, katta sig'imli ikkilamchi xotira
\item SSD - lazer orqali o'qiladigan va yoziladigan disk
\item SSD - ketma-ket murojaatli, arxiv uchun qulay magnit vosita
\item SSD - flesh hujayralarga asoslangan harakatsiz ikkilamchi xotira
\end{enumerate}
\begin{tcolorbox}[enhanced,breakable,colback=teal!4,colframe=teal!60!black,boxrule=0.45pt,arc=1.2mm,left=2.7mm,right=2.7mm,top=1.8mm,bottom=1.8mm,before skip=2mm,after skip=2mm]
{\bfseries\sffamily\color{teal}To'g'ri javob:} \textbf{D) SSD - flesh hujayralarga asoslangan harakatsiz ikkilamchi xotira}
\par\smallskip
{\bfseries Izoh:} SSD - flesh hujayralarga asoslangan harakatsiz ikkilamchi xotira. To'g'ri juftlik: SSD - flesh hujayralarga asoslangan harakatsiz ikkilamchi xotira. Manba: CAM10/11, PDF 34-37; Qo'shimcha sahifa raqami (qaysi PDFga tegishli ekani aniqlanmagan; darslik sahifasi deb tasdiqlanmagan): 31
\end{tcolorbox}
\vspace{2mm}\hrule height 0.2pt\vspace{3mm}

% DB_ID 84a6bd97-7e3e-9f05-fdc1-5533bedc1909
\needspace{5\baselineskip}
{\bfseries\sffamily\color{ink} 178.} {scriptsizecolor{muted} Y1 | review | 84a6bd97}
\par
Chiqarish va ikki vazifali qurilmalar mavzusida "qattiq nusxa" qatnashgan atama-ta'rif juftliklaridan to'g'risini toping.
\begin{enumerate}[label=\Alph*),leftmargin=9mm]
\item qattiq nusxa - ekranda yoki audio orqali vaqtinchalik taqdim etilgan natija
\item qattiq nusxa - toner va elektrofotografik jarayon bilan tez chop etuvchi printer
\item qattiq nusxa - mayda siyoh tomchilarini purkab tasvir hosil qiluvchi printer
\item qattiq nusxa - qog'oz yoki boshqa fizik muhitdagi natija
\end{enumerate}
\begin{tcolorbox}[enhanced,breakable,colback=teal!4,colframe=teal!60!black,boxrule=0.45pt,arc=1.2mm,left=2.7mm,right=2.7mm,top=1.8mm,bottom=1.8mm,before skip=2mm,after skip=2mm]
{\bfseries\sffamily\color{teal}To'g'ri javob:} \textbf{D) qattiq nusxa - qog'oz yoki boshqa fizik muhitdagi natija}
\par\smallskip
{\bfseries Izoh:} qattiq nusxa - qog'oz yoki boshqa fizik muhitdagi natija. To'g'ri juftlik: qattiq nusxa - qog'oz yoki boshqa fizik muhitdagi natija. Manba: CAM10/11, PDF 37-45; Qo'shimcha sahifa raqami (qaysi PDFga tegishli ekani aniqlanmagan; darslik sahifasi deb tasdiqlanmagan): 43
\end{tcolorbox}
\vspace{2mm}\hrule height 0.2pt\vspace{3mm}

% DB_ID 86019e8f-c539-117b-320f-67438f6c78cb
\needspace{5\baselineskip}
{\bfseries\sffamily\color{ink} 179.} {scriptsizecolor{muted} Y1 | review | 86019e8f}
\par
Kompyuter tizimi va uning ishlash modeli bo'yicha qatorlar ichidan "dasturiy ta'minot" noto'g'ri izohlanganini aniqlang.
\begin{enumerate}[label=\Alph*),leftmargin=9mm]
\item texnik ta'minot - kompyuter tizimining fizik qismlari
\item dasturiy ta'minot - bir maqsad uchun hamkorlik qiluvchi texnik vosita, dastur, ma'lumot va foydalanuvchi majmuasi
\item IPOS - kiritish, qayta ishlash, chiqarish va saqlash modeli
\item o'rnatilgan tizim - kattaroq qurilma ichida muayyan vazifani bajaruvchi kompyuter tizimi
\end{enumerate}
\begin{tcolorbox}[enhanced,breakable,colback=teal!4,colframe=teal!60!black,boxrule=0.45pt,arc=1.2mm,left=2.7mm,right=2.7mm,top=1.8mm,bottom=1.8mm,before skip=2mm,after skip=2mm]
{\bfseries\sffamily\color{teal}To'g'ri javob:} \textbf{B) dasturiy ta'minot - bir maqsad uchun hamkorlik qiluvchi texnik vosita, dastur, ma'lumot va foydalanuvchi majmuasi}
\par\smallskip
{\bfseries Izoh:} dasturiy ta'minot - bir maqsad uchun hamkorlik qiluvchi texnik vosita, dastur, ma'lumot va foydalanuvchi majmuasi. dasturiy ta'minot uchun to'g'ri ta'rif: kompyuterga vazifalarni bajarishni ko'rsatuvchi dasturlar. Manba: ICT5, PDF 25-32; Qo'shimcha sahifa raqami (qaysi PDFga tegishli ekani aniqlanmagan; darslik sahifasi deb tasdiqlanmagan): 6
\end{tcolorbox}
\vspace{2mm}\hrule height 0.2pt\vspace{3mm}

% DB_ID 8b723c66-676f-46c5-07a6-ea20861ee4f6
\needspace{5\baselineskip}
{\bfseries\sffamily\color{ink} 180.} {scriptsizecolor{muted} Y1 | review | 8b723c66}
\par
Birlamchi xotira: registr, kesh, RAM va ROM mavzusida "ROM" qatnashgan atama-ta'rif juftliklaridan to'g'risini toping.
\begin{enumerate}[label=\Alph*),leftmargin=9mm]
\item ROM - faol dastur va ma'lumotlar uchun o'zgaruvchan tezkor xotira
\item ROM - elektrsiz ham ma'lumotni saqlovchi o'zgarmas xotira
\item ROM - tez-tez kerak bo'ladigan ma'lumotni CPUga yaqin saqlovchi xotira
\item ROM - RAM yetishmaganda ikkilamchi xotiradan ajratilgan sekinroq maydon
\end{enumerate}
\begin{tcolorbox}[enhanced,breakable,colback=teal!4,colframe=teal!60!black,boxrule=0.45pt,arc=1.2mm,left=2.7mm,right=2.7mm,top=1.8mm,bottom=1.8mm,before skip=2mm,after skip=2mm]
{\bfseries\sffamily\color{teal}To'g'ri javob:} \textbf{B) ROM - elektrsiz ham ma'lumotni saqlovchi o'zgarmas xotira}
\par\smallskip
{\bfseries Izoh:} ROM - elektrsiz ham ma'lumotni saqlovchi o'zgarmas xotira. To'g'ri juftlik: ROM - elektrsiz ham ma'lumotni saqlovchi o'zgarmas xotira. Manba: CAM10/11, PDF 33-36; Qo'shimcha sahifa raqami (qaysi PDFga tegishli ekani aniqlanmagan; darslik sahifasi deb tasdiqlanmagan): 24
\end{tcolorbox}
\vspace{2mm}\hrule height 0.2pt\vspace{3mm}

% DB_ID 8c15f6d7-2b64-c715-8378-d8c90317f2db
\needspace{5\baselineskip}
{\bfseries\sffamily\color{ink} 181.} {scriptsizecolor{muted} Y1 | review | 8c15f6d7}
\par
Unumdorlik, nosozliklarni aniqlash va xavfsiz xizmat bo'yicha qatorlar ichidan "bottleneck" noto'g'ri izohlanganini aniqlang.
\begin{enumerate}[label=\Alph*),leftmargin=9mm]
\item thermal throttling - qizish sabab protsessor tezligining avtomatik pasayishi
\item diagnostika - belgi va dalillar asosida sababni tizimli aniqlash
\item log - tizim hodisalari va xatolari qaydi
\item bottleneck - RAM va ikkilamchi xotira o'rtasida sahifalarni almashish
\end{enumerate}
\begin{tcolorbox}[enhanced,breakable,colback=teal!4,colframe=teal!60!black,boxrule=0.45pt,arc=1.2mm,left=2.7mm,right=2.7mm,top=1.8mm,bottom=1.8mm,before skip=2mm,after skip=2mm]
{\bfseries\sffamily\color{teal}To'g'ri javob:} \textbf{D) bottleneck - RAM va ikkilamchi xotira o'rtasida sahifalarni almashish}
\par\smallskip
{\bfseries Izoh:} bottleneck - RAM va ikkilamchi xotira o'rtasida sahifalarni almashish. bottleneck uchun to'g'ri ta'rif: tizim umumiy tezligini cheklayotgan eng sust resurs. Manba: CAM9, PDF 12-47; Qo'shimcha sahifa raqami (qaysi PDFga tegishli ekani aniqlanmagan; darslik sahifasi deb tasdiqlanmagan): 108
\end{tcolorbox}
\vspace{2mm}\hrule height 0.2pt\vspace{3mm}

% DB_ID 9288de8a-523b-1028-067f-664ca2d76464
\needspace{5\baselineskip}
{\bfseries\sffamily\color{ink} 182.} {scriptsizecolor{muted} Y1 | review | 9288de8a}
\par
Ikkilamchi saqlash qurilmalari bo'yicha qatorlar ichidan "HDD" noto'g'ri izohlanganini aniqlang.
\begin{enumerate}[label=\Alph*),leftmargin=9mm]
\item optik disk - lazer orqali o'qiladigan va yoziladigan disk
\item magnit lenta - ketma-ket murojaatli, arxiv uchun qulay magnit vosita
\item USB flesh - portativ solid-state saqlash qurilmasi
\item HDD - flesh hujayralarga asoslangan harakatsiz ikkilamchi xotira
\end{enumerate}
\begin{tcolorbox}[enhanced,breakable,colback=teal!4,colframe=teal!60!black,boxrule=0.45pt,arc=1.2mm,left=2.7mm,right=2.7mm,top=1.8mm,bottom=1.8mm,before skip=2mm,after skip=2mm]
{\bfseries\sffamily\color{teal}To'g'ri javob:} \textbf{D) HDD - flesh hujayralarga asoslangan harakatsiz ikkilamchi xotira}
\par\smallskip
{\bfseries Izoh:} HDD - flesh hujayralarga asoslangan harakatsiz ikkilamchi xotira. HDD uchun to'g'ri ta'rif: magnit plastinkali, katta sig'imli ikkilamchi xotira. Manba: CAM9, PDF 33-47; Qo'shimcha sahifa raqami (qaysi PDFga tegishli ekani aniqlanmagan; darslik sahifasi deb tasdiqlanmagan): 30
\end{tcolorbox}
\vspace{2mm}\hrule height 0.2pt\vspace{3mm}

% DB_ID 94741b31-5484-e641-ce93-f0ff32bc24a6
\needspace{5\baselineskip}
{\bfseries\sffamily\color{ink} 183.} {scriptsizecolor{muted} Y1 | review | 94741b31}
\par
Markaziy protsessor va bajarish sikli bo'yicha qatorlar ichidan "takt chastotasi" noto'g'ri izohlanganini aniqlang.
\begin{enumerate}[label=\Alph*),leftmargin=9mm]
\item ALU - arifmetik va mantiqiy amallarni bajaruvchi qism
\item CU - ko'rsatmalar oqimi va boshqa qismlar ishini boshqaruvchi qism
\item registr - CPU ichidagi juda tez va kichik xotira
\item takt chastotasi - ko'rsatmalarni bajarib, ma'lumotlarni qayta ishlovchi markaziy qurilma
\end{enumerate}
\begin{tcolorbox}[enhanced,breakable,colback=teal!4,colframe=teal!60!black,boxrule=0.45pt,arc=1.2mm,left=2.7mm,right=2.7mm,top=1.8mm,bottom=1.8mm,before skip=2mm,after skip=2mm]
{\bfseries\sffamily\color{teal}To'g'ri javob:} \textbf{D) takt chastotasi - ko'rsatmalarni bajarib, ma'lumotlarni qayta ishlovchi markaziy qurilma}
\par\smallskip
{\bfseries Izoh:} takt chastotasi - ko'rsatmalarni bajarib, ma'lumotlarni qayta ishlovchi markaziy qurilma. takt chastotasi uchun to'g'ri ta'rif: sekunddagi takt impulslar soni. Manba: CAM9, PDF 10-14; Qo'shimcha sahifa raqami (qaysi PDFga tegishli ekani aniqlanmagan; darslik sahifasi deb tasdiqlanmagan): 21
\end{tcolorbox}
\vspace{2mm}\hrule height 0.2pt\vspace{3mm}

% DB_ID 95cf2e64-5575-7197-b3b5-ac1ee667f980
\needspace{5\baselineskip}
{\bfseries\sffamily\color{ink} 184.} {scriptsizecolor{muted} Y1 | review | 95cf2e64}
\par
Markaziy protsessor va bajarish sikli bo'yicha qatorlar ichidan "CU" noto'g'ri izohlanganini aniqlang.
\begin{enumerate}[label=\Alph*),leftmargin=9mm]
\item ALU - arifmetik va mantiqiy amallarni bajaruvchi qism
\item CU - ko'rsatmalarni bajarib, ma'lumotlarni qayta ishlovchi markaziy qurilma
\item registr - CPU ichidagi juda tez va kichik xotira
\item takt chastotasi - sekunddagi takt impulslar soni
\end{enumerate}
\begin{tcolorbox}[enhanced,breakable,colback=teal!4,colframe=teal!60!black,boxrule=0.45pt,arc=1.2mm,left=2.7mm,right=2.7mm,top=1.8mm,bottom=1.8mm,before skip=2mm,after skip=2mm]
{\bfseries\sffamily\color{teal}To'g'ri javob:} \textbf{B) CU - ko'rsatmalarni bajarib, ma'lumotlarni qayta ishlovchi markaziy qurilma}
\par\smallskip
{\bfseries Izoh:} CU - ko'rsatmalarni bajarib, ma'lumotlarni qayta ishlovchi markaziy qurilma. CU uchun to'g'ri ta'rif: ko'rsatmalar oqimi va boshqa qismlar ishini boshqaruvchi qism. Manba: ICT10, PDF 17-20; Qo'shimcha sahifa raqami (qaysi PDFga tegishli ekani aniqlanmagan; darslik sahifasi deb tasdiqlanmagan): 19
\end{tcolorbox}
\vspace{2mm}\hrule height 0.2pt\vspace{3mm}

% DB_ID 98df209e-5664-5e73-3476-18a55291d51f
\needspace{5\baselineskip}
{\bfseries\sffamily\color{ink} 185.} {scriptsizecolor{muted} Y1 | review | 98df209e}
\par
Kompyuter tizimi va uning ishlash modeli bo'yicha qatorlar ichidan "o'rnatilgan tizim" noto'g'ri izohlanganini aniqlang.
\begin{enumerate}[label=\Alph*),leftmargin=9mm]
\item texnik ta'minot - kompyuter tizimining fizik qismlari
\item dasturiy ta'minot - kompyuterga vazifalarni bajarishni ko'rsatuvchi dasturlar
\item IPOS - kiritish, qayta ishlash, chiqarish va saqlash modeli
\item o'rnatilgan tizim - bir maqsad uchun hamkorlik qiluvchi texnik vosita, dastur, ma'lumot va foydalanuvchi majmuasi
\end{enumerate}
\begin{tcolorbox}[enhanced,breakable,colback=teal!4,colframe=teal!60!black,boxrule=0.45pt,arc=1.2mm,left=2.7mm,right=2.7mm,top=1.8mm,bottom=1.8mm,before skip=2mm,after skip=2mm]
{\bfseries\sffamily\color{teal}To'g'ri javob:} \textbf{D) o'rnatilgan tizim - bir maqsad uchun hamkorlik qiluvchi texnik vosita, dastur, ma'lumot va foydalanuvchi majmuasi}
\par\smallskip
{\bfseries Izoh:} o'rnatilgan tizim - bir maqsad uchun hamkorlik qiluvchi texnik vosita, dastur, ma'lumot va foydalanuvchi majmuasi. o'rnatilgan tizim uchun to'g'ri ta'rif: kattaroq qurilma ichida muayyan vazifani bajaruvchi kompyuter tizimi. Manba: ICT9, PDF 14; Qo'shimcha sahifa raqami (qaysi PDFga tegishli ekani aniqlanmagan; darslik sahifasi deb tasdiqlanmagan): 8
\end{tcolorbox}
\vspace{2mm}\hrule height 0.2pt\vspace{3mm}

% DB_ID 98f44105-6ad0-473c-3564-0bdeeb67fc55
\needspace{5\baselineskip}
{\bfseries\sffamily\color{ink} 186.} {scriptsizecolor{muted} Y1 | review | 98f44105}
\par
Chiqarish va ikki vazifali qurilmalar bo'yicha qatorlar ichidan "lazer printer" noto'g'ri izohlanganini aniqlang.
\begin{enumerate}[label=\Alph*),leftmargin=9mm]
\item qattiq nusxa - qog'oz yoki boshqa fizik muhitdagi natija
\item lazer printer - ekranda yoki audio orqali vaqtinchalik taqdim etilgan natija
\item siyohpurkagich - mayda siyoh tomchilarini purkab tasvir hosil qiluvchi printer
\item aktuator - boshqaruv signalini fizik harakat yoki ta'sirga aylantiruvchi qurilma
\end{enumerate}
\begin{tcolorbox}[enhanced,breakable,colback=teal!4,colframe=teal!60!black,boxrule=0.45pt,arc=1.2mm,left=2.7mm,right=2.7mm,top=1.8mm,bottom=1.8mm,before skip=2mm,after skip=2mm]
{\bfseries\sffamily\color{teal}To'g'ri javob:} \textbf{B) lazer printer - ekranda yoki audio orqali vaqtinchalik taqdim etilgan natija}
\par\smallskip
{\bfseries Izoh:} lazer printer - ekranda yoki audio orqali vaqtinchalik taqdim etilgan natija. lazer printer uchun to'g'ri ta'rif: toner va elektrofotografik jarayon bilan tez chop etuvchi printer. Manba: ICT5, PDF 29-36; Qo'shimcha sahifa raqami (qaysi PDFga tegishli ekani aniqlanmagan; darslik sahifasi deb tasdiqlanmagan): 44
\end{tcolorbox}
\vspace{2mm}\hrule height 0.2pt\vspace{3mm}

% DB_ID 9b80745a-446a-6613-85e8-ece4ffc7874d
\needspace{5\baselineskip}
{\bfseries\sffamily\color{ink} 187.} {scriptsizecolor{muted} Y1 | review | 9b80745a}
\par
Attestatsiyadagi nozik farqlar va integratsion masalalar mavzusida "integratsion savol" qatnashgan atama-ta'rif juftliklaridan to'g'risini toping.
\begin{enumerate}[label=\Alph*),leftmargin=9mm]
\item integratsion savol - obyektni muhim belgilari bo'yicha guruhga ajratish
\item integratsion savol - variantlarni baholash uchun qo'llanadigan belgi
\item integratsion savol - yechim bajarishi shart bo'lgan chegara yoki talab
\item integratsion savol - bir nechta mavzuni bitta vaziyatda bog'laydigan topshiriq
\end{enumerate}
\begin{tcolorbox}[enhanced,breakable,colback=teal!4,colframe=teal!60!black,boxrule=0.45pt,arc=1.2mm,left=2.7mm,right=2.7mm,top=1.8mm,bottom=1.8mm,before skip=2mm,after skip=2mm]
{\bfseries\sffamily\color{teal}To'g'ri javob:} \textbf{D) integratsion savol - bir nechta mavzuni bitta vaziyatda bog'laydigan topshiriq}
\par\smallskip
{\bfseries Izoh:} integratsion savol - bir nechta mavzuni bitta vaziyatda bog'laydigan topshiriq. To'g'ri juftlik: integratsion savol - bir nechta mavzuni bitta vaziyatda bog'laydigan topshiriq. Manba: ICT5, PDF 25-49; Qo'shimcha sahifa raqami (qaysi PDFga tegishli ekani aniqlanmagan; darslik sahifasi deb tasdiqlanmagan): 118
\end{tcolorbox}
\vspace{2mm}\hrule height 0.2pt\vspace{3mm}

% DB_ID 9c3d41ed-8b1c-7000-e1f1-94162de488e0
\needspace{5\baselineskip}
{\bfseries\sffamily\color{ink} 188.} {scriptsizecolor{muted} Y1 | review | 9c3d41ed}
\par
Kiritish qurilmalari mavzusida "biometrik kiritish" qatnashgan atama-ta'rif juftliklaridan to'g'risini toping.
\begin{enumerate}[label=\Alph*),leftmargin=9mm]
\item biometrik kiritish - tasvirdagi bosma matnni tahrirlanadigan belgilar sifatida tanish
\item biometrik kiritish - insonning noyob jismoniy yoki xulqiy belgisini kiritish
\item biometrik kiritish - oldindan belgilangan joylardagi belgilarni aniqlash
\item biometrik kiritish - magnit siyoh bilan bosilgan maxsus belgilarni o'qish
\end{enumerate}
\begin{tcolorbox}[enhanced,breakable,colback=teal!4,colframe=teal!60!black,boxrule=0.45pt,arc=1.2mm,left=2.7mm,right=2.7mm,top=1.8mm,bottom=1.8mm,before skip=2mm,after skip=2mm]
{\bfseries\sffamily\color{teal}To'g'ri javob:} \textbf{B) biometrik kiritish - insonning noyob jismoniy yoki xulqiy belgisini kiritish}
\par\smallskip
{\bfseries Izoh:} biometrik kiritish - insonning noyob jismoniy yoki xulqiy belgisini kiritish. To'g'ri juftlik: biometrik kiritish - insonning noyob jismoniy yoki xulqiy belgisini kiritish. Manba: MAR, PDF 8-9; Qo'shimcha sahifa raqami (qaysi PDFga tegishli ekani aniqlanmagan; darslik sahifasi deb tasdiqlanmagan): 41
\end{tcolorbox}
\vspace{2mm}\hrule height 0.2pt\vspace{3mm}

% DB_ID 9c7f85ba-9d33-6baa-22dc-6bba1788b29a
\needspace{5\baselineskip}
{\bfseries\sffamily\color{ink} 189.} {scriptsizecolor{muted} Y1 | review | 9c7f85ba}
\par
Kompyuter tizimi va uning ishlash modeli mavzusida "IPOS" qatnashgan atama-ta'rif juftliklaridan to'g'risini toping.
\begin{enumerate}[label=\Alph*),leftmargin=9mm]
\item IPOS - bir maqsad uchun hamkorlik qiluvchi texnik vosita, dastur, ma'lumot va foydalanuvchi majmuasi
\item IPOS - kiritish, qayta ishlash, chiqarish va saqlash modeli
\item IPOS - kompyuter tizimining fizik qismlari
\item IPOS - kompyuterga vazifalarni bajarishni ko'rsatuvchi dasturlar
\end{enumerate}
\begin{tcolorbox}[enhanced,breakable,colback=teal!4,colframe=teal!60!black,boxrule=0.45pt,arc=1.2mm,left=2.7mm,right=2.7mm,top=1.8mm,bottom=1.8mm,before skip=2mm,after skip=2mm]
{\bfseries\sffamily\color{teal}To'g'ri javob:} \textbf{B) IPOS - kiritish, qayta ishlash, chiqarish va saqlash modeli}
\par\smallskip
{\bfseries Izoh:} IPOS - kiritish, qayta ishlash, chiqarish va saqlash modeli. To'g'ri juftlik: IPOS - kiritish, qayta ishlash, chiqarish va saqlash modeli. Manba: ICT10, PDF 6-12; Qo'shimcha sahifa raqami (qaysi PDFga tegishli ekani aniqlanmagan; darslik sahifasi deb tasdiqlanmagan): 6
\end{tcolorbox}
\vspace{2mm}\hrule height 0.2pt\vspace{3mm}

% DB_ID 9fdd2a7c-7c34-65ba-9660-dfa1a3f35c1f
\needspace{5\baselineskip}
{\bfseries\sffamily\color{ink} 190.} {scriptsizecolor{muted} Y1 | review | 9fdd2a7c}
\par
Chiqarish va ikki vazifali qurilmalar mavzusida "siyohpurkagich" qatnashgan atama-ta'rif juftliklaridan to'g'risini toping.
\begin{enumerate}[label=\Alph*),leftmargin=9mm]
\item siyohpurkagich - ekranda yoki audio orqali vaqtinchalik taqdim etilgan natija
\item siyohpurkagich - mayda siyoh tomchilarini purkab tasvir hosil qiluvchi printer
\item siyohpurkagich - qog'oz yoki boshqa fizik muhitdagi natija
\item siyohpurkagich - toner va elektrofotografik jarayon bilan tez chop etuvchi printer
\end{enumerate}
\begin{tcolorbox}[enhanced,breakable,colback=teal!4,colframe=teal!60!black,boxrule=0.45pt,arc=1.2mm,left=2.7mm,right=2.7mm,top=1.8mm,bottom=1.8mm,before skip=2mm,after skip=2mm]
{\bfseries\sffamily\color{teal}To'g'ri javob:} \textbf{B) siyohpurkagich - mayda siyoh tomchilarini purkab tasvir hosil qiluvchi printer}
\par\smallskip
{\bfseries Izoh:} siyohpurkagich - mayda siyoh tomchilarini purkab tasvir hosil qiluvchi printer. To'g'ri juftlik: siyohpurkagich - mayda siyoh tomchilarini purkab tasvir hosil qiluvchi printer. Manba: ICT10, PDF 9-13; Qo'shimcha sahifa raqami (qaysi PDFga tegishli ekani aniqlanmagan; darslik sahifasi deb tasdiqlanmagan): 45
\end{tcolorbox}
\vspace{2mm}\hrule height 0.2pt\vspace{3mm}

% DB_ID a636efb8-bfa0-edda-f41c-f3889196ddc5
\needspace{5\baselineskip}
{\bfseries\sffamily\color{ink} 191.} {scriptsizecolor{muted} Y1 | review | a636efb8}
\par
Yangi grafik karta fizik jihatdan korpusga sig'adi, ammo bosh platadagi ulagich turi mos emas. Qaysi moslik tekshirilmagan?
\begin{enumerate}[label=\Alph*),leftmargin=9mm]
\item monitorning ekran o'lchami
\item fayl tizimining nomi
\item kengaytirish sloti interfeysi
\item klaviatura til tartibi
\end{enumerate}
\begin{tcolorbox}[enhanced,breakable,colback=teal!4,colframe=teal!60!black,boxrule=0.45pt,arc=1.2mm,left=2.7mm,right=2.7mm,top=1.8mm,bottom=1.8mm,before skip=2mm,after skip=2mm]
{\bfseries\sffamily\color{teal}To'g'ri javob:} \textbf{C) kengaytirish sloti interfeysi}
\par\smallskip
{\bfseries Izoh:} kengaytirish sloti interfeysi. Ichki karta mos PCIe kabi slot va interfeysni talab qiladi. Manba: CAM9, PDF 9-14; Qo'shimcha sahifa raqami (qaysi PDFga tegishli ekani aniqlanmagan; darslik sahifasi deb tasdiqlanmagan): 16
\end{tcolorbox}
\vspace{2mm}\hrule height 0.2pt\vspace{3mm}

% DB_ID aa7389d5-2a0c-8806-056d-a2f6b4e54627
\needspace{5\baselineskip}
{\bfseries\sffamily\color{ink} 192.} {scriptsizecolor{muted} Y1 | review | aa7389d5}
\par
Elektr o'chdi. Qaysi ma'lumot yo'qolishi ehtimoli eng yuqori?
\begin{enumerate}[label=\Alph*),leftmargin=9mm]
\item SSDga avval yozilgan loyiha
\item ROMdagi yuklash kodi
\item fleshga ko'chirilgan arxiv
\item saqlanmagan hujjatning RAMdagi o'zgarishlari
\end{enumerate}
\begin{tcolorbox}[enhanced,breakable,colback=teal!4,colframe=teal!60!black,boxrule=0.45pt,arc=1.2mm,left=2.7mm,right=2.7mm,top=1.8mm,bottom=1.8mm,before skip=2mm,after skip=2mm]
{\bfseries\sffamily\color{teal}To'g'ri javob:} \textbf{D) saqlanmagan hujjatning RAMdagi o'zgarishlari}
\par\smallskip
{\bfseries Izoh:} saqlanmagan hujjatning RAMdagi o'zgarishlari. RAM o'zgaruvchan; doimiy saqlashdagi yozilgan ma'lumot saqlanadi. Manba: CAM9, PDF 11-14; Qo'shimcha sahifa raqami (qaysi PDFga tegishli ekani aniqlanmagan; darslik sahifasi deb tasdiqlanmagan): 29
\end{tcolorbox}
\vspace{2mm}\hrule height 0.2pt\vspace{3mm}

% DB_ID b00f2bf3-ef12-cc77-93e3-461eed210f0f
\needspace{5\baselineskip}
{\bfseries\sffamily\color{ink} 193.} {scriptsizecolor{muted} Y1 | review | b00f2bf3}
\par
Ikkilamchi saqlash qurilmalari mavzusida "bulutli saqlash" qatnashgan atama-ta'rif juftliklaridan to'g'risini toping.
\begin{enumerate}[label=\Alph*),leftmargin=9mm]
\item bulutli saqlash - magnit plastinkali, katta sig'imli ikkilamchi xotira
\item bulutli saqlash - flesh hujayralarga asoslangan harakatsiz ikkilamchi xotira
\item bulutli saqlash - lazer orqali o'qiladigan va yoziladigan disk
\item bulutli saqlash - tarmoq orqali masofadagi server xotirasidan foydalanish xizmati
\end{enumerate}
\begin{tcolorbox}[enhanced,breakable,colback=teal!4,colframe=teal!60!black,boxrule=0.45pt,arc=1.2mm,left=2.7mm,right=2.7mm,top=1.8mm,bottom=1.8mm,before skip=2mm,after skip=2mm]
{\bfseries\sffamily\color{teal}To'g'ri javob:} \textbf{D) bulutli saqlash - tarmoq orqali masofadagi server xotirasidan foydalanish xizmati}
\par\smallskip
{\bfseries Izoh:} bulutli saqlash - tarmoq orqali masofadagi server xotirasidan foydalanish xizmati. To'g'ri juftlik: bulutli saqlash - tarmoq orqali masofadagi server xotirasidan foydalanish xizmati. Manba: CAM9, PDF 33-47; Qo'shimcha sahifa raqami (qaysi PDFga tegishli ekani aniqlanmagan; darslik sahifasi deb tasdiqlanmagan): 35
\end{tcolorbox}
\vspace{2mm}\hrule height 0.2pt\vspace{3mm}

% DB_ID b1d3c81c-ca2f-bab3-b6e0-84e4bc702674
\needspace{5\baselineskip}
{\bfseries\sffamily\color{ink} 194.} {scriptsizecolor{muted} Y1 | review | b1d3c81c}
\par
Chiqarish va ikki vazifali qurilmalar bo'yicha qatorlar ichidan "aktuator" noto'g'ri izohlanganini aniqlang.
\begin{enumerate}[label=\Alph*),leftmargin=9mm]
\item qattiq nusxa - qog'oz yoki boshqa fizik muhitdagi natija
\item lazer printer - toner va elektrofotografik jarayon bilan tez chop etuvchi printer
\item siyohpurkagich - mayda siyoh tomchilarini purkab tasvir hosil qiluvchi printer
\item aktuator - ekranda yoki audio orqali vaqtinchalik taqdim etilgan natija
\end{enumerate}
\begin{tcolorbox}[enhanced,breakable,colback=teal!4,colframe=teal!60!black,boxrule=0.45pt,arc=1.2mm,left=2.7mm,right=2.7mm,top=1.8mm,bottom=1.8mm,before skip=2mm,after skip=2mm]
{\bfseries\sffamily\color{teal}To'g'ri javob:} \textbf{D) aktuator - ekranda yoki audio orqali vaqtinchalik taqdim etilgan natija}
\par\smallskip
{\bfseries Izoh:} aktuator - ekranda yoki audio orqali vaqtinchalik taqdim etilgan natija. aktuator uchun to'g'ri ta'rif: boshqaruv signalini fizik harakat yoki ta'sirga aylantiruvchi qurilma. Manba: MAR, PDF 8-10; Qo'shimcha sahifa raqami (qaysi PDFga tegishli ekani aniqlanmagan; darslik sahifasi deb tasdiqlanmagan): 47
\end{tcolorbox}
\vspace{2mm}\hrule height 0.2pt\vspace{3mm}

% DB_ID b764ed10-5174-ef27-c083-370bc0eca5a5
\needspace{5\baselineskip}
{\bfseries\sffamily\color{ink} 195.} {scriptsizecolor{muted} Y1 | review | b764ed10}
\par
Chiqarish va ikki vazifali qurilmalar mavzusida "haptik chiqish" qatnashgan atama-ta'rif juftliklaridan to'g'risini toping.
\begin{enumerate}[label=\Alph*),leftmargin=9mm]
\item haptik chiqish - ekranda yoki audio orqali vaqtinchalik taqdim etilgan natija
\item haptik chiqish - qog'oz yoki boshqa fizik muhitdagi natija
\item haptik chiqish - toner va elektrofotografik jarayon bilan tez chop etuvchi printer
\item haptik chiqish - tebranish yoki kuch orqali beriladigan sezgi chiqishi
\end{enumerate}
\begin{tcolorbox}[enhanced,breakable,colback=teal!4,colframe=teal!60!black,boxrule=0.45pt,arc=1.2mm,left=2.7mm,right=2.7mm,top=1.8mm,bottom=1.8mm,before skip=2mm,after skip=2mm]
{\bfseries\sffamily\color{teal}To'g'ri javob:} \textbf{D) haptik chiqish - tebranish yoki kuch orqali beriladigan sezgi chiqishi}
\par\smallskip
{\bfseries Izoh:} haptik chiqish - tebranish yoki kuch orqali beriladigan sezgi chiqishi. To'g'ri juftlik: haptik chiqish - tebranish yoki kuch orqali beriladigan sezgi chiqishi. Manba: CAM9, PDF 16-32; Qo'shimcha sahifa raqami (qaysi PDFga tegishli ekani aniqlanmagan; darslik sahifasi deb tasdiqlanmagan): 48
\end{tcolorbox}
\vspace{2mm}\hrule height 0.2pt\vspace{3mm}

% DB_ID b76e294f-af6e-1116-fc36-29d4889b6641
\needspace{5\baselineskip}
{\bfseries\sffamily\color{ink} 196.} {scriptsizecolor{muted} Y1 | review | b76e294f}
\par
Tizimli blok, bosh plata va ulanishlar bo'yicha qatorlar ichidan "quvvat bloki" noto'g'ri izohlanganini aniqlang.
\begin{enumerate}[label=\Alph*),leftmargin=9mm]
\item shina - ma'lumot va boshqaruv signallarini tashuvchi aloqa yo'llari
\item quvvat bloki - ichki komponentlarni bog'lovchi asosiy bosma plata
\item slot - ichki komponent o'rnatiladigan ulagich
\item port - tashqi qurilma ulanadigan nuqta
\end{enumerate}
\begin{tcolorbox}[enhanced,breakable,colback=teal!4,colframe=teal!60!black,boxrule=0.45pt,arc=1.2mm,left=2.7mm,right=2.7mm,top=1.8mm,bottom=1.8mm,before skip=2mm,after skip=2mm]
{\bfseries\sffamily\color{teal}To'g'ri javob:} \textbf{B) quvvat bloki - ichki komponentlarni bog'lovchi asosiy bosma plata}
\par\smallskip
{\bfseries Izoh:} quvvat bloki - ichki komponentlarni bog'lovchi asosiy bosma plata. quvvat bloki uchun to'g'ri ta'rif: tarmoq elektrini komponentlarga mos doimiy tokka aylantiruvchi qism. Manba: CAM9, PDF 9-14; Qo'shimcha sahifa raqami (qaysi PDFga tegishli ekani aniqlanmagan; darslik sahifasi deb tasdiqlanmagan): 14
\end{tcolorbox}
\vspace{2mm}\hrule height 0.2pt\vspace{3mm}

% DB_ID bc1916cc-b6b2-ca17-c00d-9dcb98e3f4ec
\needspace{5\baselineskip}
{\bfseries\sffamily\color{ink} 197.} {scriptsizecolor{muted} Y1 | review | bc1916cc}
\par
Kiritish qurilmalari bo'yicha qatorlar ichidan "MICR" noto'g'ri izohlanganini aniqlang.
\begin{enumerate}[label=\Alph*),leftmargin=9mm]
\item OMR - oldindan belgilangan joylardagi belgilarni aniqlash
\item shtrix-kod o'quvchi - chiziqli koddagi identifikatorni avtomatik kiritish
\item sensor - fizik kattalikni o'lchab signalga aylantiruvchi qurilma
\item MICR - tasvirdagi bosma matnni tahrirlanadigan belgilar sifatida tanish
\end{enumerate}
\begin{tcolorbox}[enhanced,breakable,colback=teal!4,colframe=teal!60!black,boxrule=0.45pt,arc=1.2mm,left=2.7mm,right=2.7mm,top=1.8mm,bottom=1.8mm,before skip=2mm,after skip=2mm]
{\bfseries\sffamily\color{teal}To'g'ri javob:} \textbf{D) MICR - tasvirdagi bosma matnni tahrirlanadigan belgilar sifatida tanish}
\par\smallskip
{\bfseries Izoh:} MICR - tasvirdagi bosma matnni tahrirlanadigan belgilar sifatida tanish. MICR uchun to'g'ri ta'rif: magnit siyoh bilan bosilgan maxsus belgilarni o'qish. Manba: CAM5, PDF 60-70; Qo'shimcha sahifa raqami (qaysi PDFga tegishli ekani aniqlanmagan; darslik sahifasi deb tasdiqlanmagan): 38
\end{tcolorbox}
\vspace{2mm}\hrule height 0.2pt\vspace{3mm}

% DB_ID be56bb2e-2533-06a8-ed50-d65711d3fb9a
\needspace{5\baselineskip}
{\bfseries\sffamily\color{ink} 198.} {scriptsizecolor{muted} Y1 | review | be56bb2e}
\par
Avtomatik sug'orish tizimi klapanni ochdi. Klapan mexanizmi qaysi turga kiradi?
\begin{enumerate}[label=\Alph*),leftmargin=9mm]
\item sensor
\item aktuator
\item saqlash vositasi
\item matn kiritish qurilmasi
\end{enumerate}
\begin{tcolorbox}[enhanced,breakable,colback=teal!4,colframe=teal!60!black,boxrule=0.45pt,arc=1.2mm,left=2.7mm,right=2.7mm,top=1.8mm,bottom=1.8mm,before skip=2mm,after skip=2mm]
{\bfseries\sffamily\color{teal}To'g'ri javob:} \textbf{B) aktuator}
\par\smallskip
{\bfseries Izoh:} aktuator. Aktuator boshqaruv signalini fizik harakatga aylantiradi. Manba: CAM9, PDF 16-32; Qo'shimcha sahifa raqami (qaysi PDFga tegishli ekani aniqlanmagan; darslik sahifasi deb tasdiqlanmagan): 48
\end{tcolorbox}
\vspace{2mm}\hrule height 0.2pt\vspace{3mm}

% DB_ID c3c9416c-9edd-f53a-dc2b-a0b9be4a50c8
\needspace{5\baselineskip}
{\bfseries\sffamily\color{ink} 199.} {scriptsizecolor{muted} Y1 | review | c3c9416c}
\par
Markaziy protsessor va bajarish sikli mavzusida "registr" qatnashgan atama-ta'rif juftliklaridan to'g'risini toping.
\begin{enumerate}[label=\Alph*),leftmargin=9mm]
\item registr - ko'rsatmalarni bajarib, ma'lumotlarni qayta ishlovchi markaziy qurilma
\item registr - CPU ichidagi juda tez va kichik xotira
\item registr - arifmetik va mantiqiy amallarni bajaruvchi qism
\item registr - ko'rsatmalar oqimi va boshqa qismlar ishini boshqaruvchi qism
\end{enumerate}
\begin{tcolorbox}[enhanced,breakable,colback=teal!4,colframe=teal!60!black,boxrule=0.45pt,arc=1.2mm,left=2.7mm,right=2.7mm,top=1.8mm,bottom=1.8mm,before skip=2mm,after skip=2mm]
{\bfseries\sffamily\color{teal}To'g'ri javob:} \textbf{B) registr - CPU ichidagi juda tez va kichik xotira}
\par\smallskip
{\bfseries Izoh:} registr - CPU ichidagi juda tez va kichik xotira. To'g'ri juftlik: registr - CPU ichidagi juda tez va kichik xotira. Manba: MAR, PDF 6-7; Qo'shimcha sahifa raqami (qaysi PDFga tegishli ekani aniqlanmagan; darslik sahifasi deb tasdiqlanmagan): 20
\end{tcolorbox}
\vspace{2mm}\hrule height 0.2pt\vspace{3mm}

% DB_ID c441fd8a-de69-9328-523f-5dfbefef8253
\needspace{5\baselineskip}
{\bfseries\sffamily\color{ink} 200.} {scriptsizecolor{muted} Y1 | review | c441fd8a}
\par
Birlamchi xotira: registr, kesh, RAM va ROM bo'yicha qatorlar ichidan "RAM" noto'g'ri izohlanganini aniqlang.
\begin{enumerate}[label=\Alph*),leftmargin=9mm]
\item kesh - tez-tez kerak bo'ladigan ma'lumotni CPUga yaqin saqlovchi xotira
\item RAM - elektrsiz ham ma'lumotni saqlovchi o'zgarmas xotira
\item virtual xotira - RAM yetishmaganda ikkilamchi xotiradan ajratilgan sekinroq maydon
\item firmware - qurilma ichiga o'rnatilgan past darajali boshqaruv dasturi
\end{enumerate}
\begin{tcolorbox}[enhanced,breakable,colback=teal!4,colframe=teal!60!black,boxrule=0.45pt,arc=1.2mm,left=2.7mm,right=2.7mm,top=1.8mm,bottom=1.8mm,before skip=2mm,after skip=2mm]
{\bfseries\sffamily\color{teal}To'g'ri javob:} \textbf{B) RAM - elektrsiz ham ma'lumotni saqlovchi o'zgarmas xotira}
\par\smallskip
{\bfseries Izoh:} RAM - elektrsiz ham ma'lumotni saqlovchi o'zgarmas xotira. RAM uchun to'g'ri ta'rif: faol dastur va ma'lumotlar uchun o'zgaruvchan tezkor xotira. Manba: CAM9, PDF 11-14; Qo'shimcha sahifa raqami (qaysi PDFga tegishli ekani aniqlanmagan; darslik sahifasi deb tasdiqlanmagan): 23
\end{tcolorbox}
\vspace{2mm}\hrule height 0.2pt\vspace{3mm}

% DB_ID ce8f92fb-d482-a6a4-b0ad-0e8e46efad19
\needspace{5\baselineskip}
{\bfseries\sffamily\color{ink} 201.} {scriptsizecolor{muted} Y1 | review | ce8f92fb}
\par
Qaysi holat maxsus maqsadli o'rnatilgan tizimga eng yaxshi misol?
\begin{enumerate}[label=\Alph*),leftmargin=9mm]
\item turli fanlarda ishlatiladigan noutbuk
\item kir yuvish mashinasining sikl boshqaruvi
\item ofisdagi umumiy ish stansiyasi
\item dars uchun planshetdagi brauzer
\end{enumerate}
\begin{tcolorbox}[enhanced,breakable,colback=teal!4,colframe=teal!60!black,boxrule=0.45pt,arc=1.2mm,left=2.7mm,right=2.7mm,top=1.8mm,bottom=1.8mm,before skip=2mm,after skip=2mm]
{\bfseries\sffamily\color{teal}To'g'ri javob:} \textbf{B) kir yuvish mashinasining sikl boshqaruvi}
\par\smallskip
{\bfseries Izoh:} kir yuvish mashinasining sikl boshqaruvi. Kir yuvish mashinasi boshqaruvi bitta aniq vazifaga moslashtirilgan. Manba: CAM9, PDF 8-15; Qo'shimcha sahifa raqami (qaysi PDFga tegishli ekani aniqlanmagan; darslik sahifasi deb tasdiqlanmagan): 9
\end{tcolorbox}
\vspace{2mm}\hrule height 0.2pt\vspace{3mm}

% DB_ID cee0c4b5-eaae-fe58-ca5e-01b216134bd6
\needspace{5\baselineskip}
{\bfseries\sffamily\color{ink} 202.} {scriptsizecolor{muted} Y1 | review | cee0c4b5}
\par
Attestatsiyadagi nozik farqlar va integratsion masalalar mavzusida "asoslash" qatnashgan atama-ta'rif juftliklaridan to'g'risini toping.
\begin{enumerate}[label=\Alph*),leftmargin=9mm]
\item asoslash - obyektni muhim belgilari bo'yicha guruhga ajratish
\item asoslash - tanlovni vaziyatdagi dalil bilan bog'lash
\item asoslash - variantlarni baholash uchun qo'llanadigan belgi
\item asoslash - yechim bajarishi shart bo'lgan chegara yoki talab
\end{enumerate}
\begin{tcolorbox}[enhanced,breakable,colback=teal!4,colframe=teal!60!black,boxrule=0.45pt,arc=1.2mm,left=2.7mm,right=2.7mm,top=1.8mm,bottom=1.8mm,before skip=2mm,after skip=2mm]
{\bfseries\sffamily\color{teal}To'g'ri javob:} \textbf{B) asoslash - tanlovni vaziyatdagi dalil bilan bog'lash}
\par\smallskip
{\bfseries Izoh:} asoslash - tanlovni vaziyatdagi dalil bilan bog'lash. To'g'ri juftlik: asoslash - tanlovni vaziyatdagi dalil bilan bog'lash. Manba: SPEC, PDF 3; Qo'shimcha sahifa raqami (qaysi PDFga tegishli ekani aniqlanmagan; darslik sahifasi deb tasdiqlanmagan): 120
\end{tcolorbox}
\vspace{2mm}\hrule height 0.2pt\vspace{3mm}

% DB_ID cf3323a0-3423-a4c5-af9a-77fe0492ce5d
\needspace{5\baselineskip}
{\bfseries\sffamily\color{ink} 203.} {scriptsizecolor{muted} Y1 | review | cf3323a0}
\par
Unumdorlik, nosozliklarni aniqlash va xavfsiz xizmat bo'yicha qatorlar ichidan "log" noto'g'ri izohlanganini aniqlang.
\begin{enumerate}[label=\Alph*),leftmargin=9mm]
\item paging - RAM va ikkilamchi xotira o'rtasida sahifalarni almashish
\item thermal throttling - qizish sabab protsessor tezligining avtomatik pasayishi
\item diagnostika - belgi va dalillar asosida sababni tizimli aniqlash
\item log - tizim umumiy tezligini cheklayotgan eng sust resurs
\end{enumerate}
\begin{tcolorbox}[enhanced,breakable,colback=teal!4,colframe=teal!60!black,boxrule=0.45pt,arc=1.2mm,left=2.7mm,right=2.7mm,top=1.8mm,bottom=1.8mm,before skip=2mm,after skip=2mm]
{\bfseries\sffamily\color{teal}To'g'ri javob:} \textbf{D) log - tizim umumiy tezligini cheklayotgan eng sust resurs}
\par\smallskip
{\bfseries Izoh:} log - tizim umumiy tezligini cheklayotgan eng sust resurs. log uchun to'g'ri ta'rif: tizim hodisalari va xatolari qaydi. Manba: CAM9, PDF 12-47; Qo'shimcha sahifa raqami (qaysi PDFga tegishli ekani aniqlanmagan; darslik sahifasi deb tasdiqlanmagan): 113
\end{tcolorbox}
\vspace{2mm}\hrule height 0.2pt\vspace{3mm}

% DB_ID e2298b8a-1c46-f00d-4314-e9dcdc57532f
\needspace{5\baselineskip}
{\bfseries\sffamily\color{ink} 204.} {scriptsizecolor{muted} Y1 | review | e2298b8a}
\par
Unumdorlik, nosozliklarni aniqlash va xavfsiz xizmat mavzusida "paging" qatnashgan atama-ta'rif juftliklaridan to'g'risini toping.
\begin{enumerate}[label=\Alph*),leftmargin=9mm]
\item paging - tizim umumiy tezligini cheklayotgan eng sust resurs
\item paging - qizish sabab protsessor tezligining avtomatik pasayishi
\item paging - belgi va dalillar asosida sababni tizimli aniqlash
\item paging - RAM va ikkilamchi xotira o'rtasida sahifalarni almashish
\end{enumerate}
\begin{tcolorbox}[enhanced,breakable,colback=teal!4,colframe=teal!60!black,boxrule=0.45pt,arc=1.2mm,left=2.7mm,right=2.7mm,top=1.8mm,bottom=1.8mm,before skip=2mm,after skip=2mm]
{\bfseries\sffamily\color{teal}To'g'ri javob:} \textbf{D) paging - RAM va ikkilamchi xotira o'rtasida sahifalarni almashish}
\par\smallskip
{\bfseries Izoh:} paging - RAM va ikkilamchi xotira o'rtasida sahifalarni almashish. To'g'ri juftlik: paging - RAM va ikkilamchi xotira o'rtasida sahifalarni almashish. Manba: CAM10/11, PDF 24-61; Qo'shimcha sahifa raqami (qaysi PDFga tegishli ekani aniqlanmagan; darslik sahifasi deb tasdiqlanmagan): 109
\end{tcolorbox}
\vspace{2mm}\hrule height 0.2pt\vspace{3mm}

% DB_ID e2b94f72-8c4f-5a1d-c233-9f94e177e8f0
\needspace{5\baselineskip}
{\bfseries\sffamily\color{ink} 205.} {scriptsizecolor{muted} Y1 | review | e2b94f72}
\par
Tizimli blok, bosh plata va ulanishlar bo'yicha qatorlar ichidan "bosh plata" noto'g'ri izohlanganini aniqlang.
\begin{enumerate}[label=\Alph*),leftmargin=9mm]
\item slot - ichki komponent o'rnatiladigan ulagich
\item bosh plata - ma'lumot va boshqaruv signallarini tashuvchi aloqa yo'llari
\item port - tashqi qurilma ulanadigan nuqta
\item quvvat bloki - tarmoq elektrini komponentlarga mos doimiy tokka aylantiruvchi qism
\end{enumerate}
\begin{tcolorbox}[enhanced,breakable,colback=teal!4,colframe=teal!60!black,boxrule=0.45pt,arc=1.2mm,left=2.7mm,right=2.7mm,top=1.8mm,bottom=1.8mm,before skip=2mm,after skip=2mm]
{\bfseries\sffamily\color{teal}To'g'ri javob:} \textbf{B) bosh plata - ma'lumot va boshqaruv signallarini tashuvchi aloqa yo'llari}
\par\smallskip
{\bfseries Izoh:} bosh plata - ma'lumot va boshqaruv signallarini tashuvchi aloqa yo'llari. bosh plata uchun to'g'ri ta'rif: ichki komponentlarni bog'lovchi asosiy bosma plata. Manba: CAM9, PDF 9-14; Qo'shimcha sahifa raqami (qaysi PDFga tegishli ekani aniqlanmagan; darslik sahifasi deb tasdiqlanmagan): 10
\end{tcolorbox}
\vspace{2mm}\hrule height 0.2pt\vspace{3mm}

% DB_ID e683114e-1d73-d52e-c199-4c673f1946a6
\needspace{5\baselineskip}
{\bfseries\sffamily\color{ink} 206.} {scriptsizecolor{muted} Y1 | review | e683114e}
\par
Unumdorlik, nosozliklarni aniqlash va xavfsiz xizmat bo'yicha qatorlar ichidan "thermal throttling" noto'g'ri izohlanganini aniqlang.
\begin{enumerate}[label=\Alph*),leftmargin=9mm]
\item paging - RAM va ikkilamchi xotira o'rtasida sahifalarni almashish
\item thermal throttling - tizim umumiy tezligini cheklayotgan eng sust resurs
\item diagnostika - belgi va dalillar asosida sababni tizimli aniqlash
\item log - tizim hodisalari va xatolari qaydi
\end{enumerate}
\begin{tcolorbox}[enhanced,breakable,colback=teal!4,colframe=teal!60!black,boxrule=0.45pt,arc=1.2mm,left=2.7mm,right=2.7mm,top=1.8mm,bottom=1.8mm,before skip=2mm,after skip=2mm]
{\bfseries\sffamily\color{teal}To'g'ri javob:} \textbf{B) thermal throttling - tizim umumiy tezligini cheklayotgan eng sust resurs}
\par\smallskip
{\bfseries Izoh:} thermal throttling - tizim umumiy tezligini cheklayotgan eng sust resurs. thermal throttling uchun to'g'ri ta'rif: qizish sabab protsessor tezligining avtomatik pasayishi. Manba: ICT10, PDF 12-29; Qo'shimcha sahifa raqami (qaysi PDFga tegishli ekani aniqlanmagan; darslik sahifasi deb tasdiqlanmagan): 111
\end{tcolorbox}
\vspace{2mm}\hrule height 0.2pt\vspace{3mm}

% DB_ID e94e0c4e-f66d-d0b7-33c4-1d5ba35b94f6
\needspace{5\baselineskip}
{\bfseries\sffamily\color{ink} 207.} {scriptsizecolor{muted} Y1 | review | e94e0c4e}
\par
Ikkilamchi saqlash qurilmalari mavzusida "magnit lenta" qatnashgan atama-ta'rif juftliklaridan to'g'risini toping.
\begin{enumerate}[label=\Alph*),leftmargin=9mm]
\item magnit lenta - magnit plastinkali, katta sig'imli ikkilamchi xotira
\item magnit lenta - ketma-ket murojaatli, arxiv uchun qulay magnit vosita
\item magnit lenta - flesh hujayralarga asoslangan harakatsiz ikkilamchi xotira
\item magnit lenta - lazer orqali o'qiladigan va yoziladigan disk
\end{enumerate}
\begin{tcolorbox}[enhanced,breakable,colback=teal!4,colframe=teal!60!black,boxrule=0.45pt,arc=1.2mm,left=2.7mm,right=2.7mm,top=1.8mm,bottom=1.8mm,before skip=2mm,after skip=2mm]
{\bfseries\sffamily\color{teal}To'g'ri javob:} \textbf{B) magnit lenta - ketma-ket murojaatli, arxiv uchun qulay magnit vosita}
\par\smallskip
{\bfseries Izoh:} magnit lenta - ketmaket murojaatli, arxiv uchun qulay magnit vosita. To'g'ri juftlik: magnit lenta - ketma-ket murojaatli, arxiv uchun qulay magnit vosita. Manba: ICT7, PDF 23-38; Qo'shimcha sahifa raqami (qaysi PDFga tegishli ekani aniqlanmagan; darslik sahifasi deb tasdiqlanmagan): 33
\end{tcolorbox}
\vspace{2mm}\hrule height 0.2pt\vspace{3mm}

% DB_ID f1199953-1699-3e5b-a21d-f6074ff6021c
\needspace{5\baselineskip}
{\bfseries\sffamily\color{ink} 208.} {scriptsizecolor{muted} Y1 | review | f1199953}
\par
Qog'ozdagi matn skanerlandi, lekin natija rasm bo'lib qoldi. Uni tahrirlanadigan matnga aylantirish uchun nima kerak?
\begin{enumerate}[label=\Alph*),leftmargin=9mm]
\item OMR sensori
\item grafik karta
\item plotter drayveri
\item OCR dasturi
\end{enumerate}
\begin{tcolorbox}[enhanced,breakable,colback=teal!4,colframe=teal!60!black,boxrule=0.45pt,arc=1.2mm,left=2.7mm,right=2.7mm,top=1.8mm,bottom=1.8mm,before skip=2mm,after skip=2mm]
{\bfseries\sffamily\color{teal}To'g'ri javob:} \textbf{D) OCR dasturi}
\par\smallskip
{\bfseries Izoh:} OCR dasturi. Skaner tasvir yaratadi; OCR belgilarni matn sifatida taniydi. Manba: CAM9, PDF 16-32; Qo'shimcha sahifa raqami (qaysi PDFga tegishli ekani aniqlanmagan; darslik sahifasi deb tasdiqlanmagan): 41
\end{tcolorbox}
\vspace{2mm}\hrule height 0.2pt\vspace{3mm}

% DB_ID fb10bf3c-8adb-4924-8cd3-b1c01a8e0a41
\needspace{5\baselineskip}
{\bfseries\sffamily\color{ink} 209.} {scriptsizecolor{muted} Y1 | review | fb10bf3c}
\par
Kiritish qurilmalari bo'yicha qatorlar ichidan "OCR" noto'g'ri izohlanganini aniqlang.
\begin{enumerate}[label=\Alph*),leftmargin=9mm]
\item MICR - magnit siyoh bilan bosilgan maxsus belgilarni o'qish
\item OCR - oldindan belgilangan joylardagi belgilarni aniqlash
\item shtrix-kod o'quvchi - chiziqli koddagi identifikatorni avtomatik kiritish
\item sensor - fizik kattalikni o'lchab signalga aylantiruvchi qurilma
\end{enumerate}
\begin{tcolorbox}[enhanced,breakable,colback=teal!4,colframe=teal!60!black,boxrule=0.45pt,arc=1.2mm,left=2.7mm,right=2.7mm,top=1.8mm,bottom=1.8mm,before skip=2mm,after skip=2mm]
{\bfseries\sffamily\color{teal}To'g'ri javob:} \textbf{B) OCR - oldindan belgilangan joylardagi belgilarni aniqlash}
\par\smallskip
{\bfseries Izoh:} OCR - oldindan belgilangan joylardagi belgilarni aniqlash. OCR uchun to'g'ri ta'rif: tasvirdagi bosma matnni tahrirlanadigan belgilar sifatida tanish. Manba: CAM9, PDF 16-32; Qo'shimcha sahifa raqami (qaysi PDFga tegishli ekani aniqlanmagan; darslik sahifasi deb tasdiqlanmagan): 36
\end{tcolorbox}
\vspace{2mm}\hrule height 0.2pt\vspace{3mm}

% DB_ID fe6b58ae-9b3c-e41e-f20e-fb97f9610b91
\needspace{5\baselineskip}
{\bfseries\sffamily\color{ink} 210.} {scriptsizecolor{muted} Y1 | review | fe6b58ae}
\par
Kompyuter tizimi va uning ishlash modeli mavzusida "texnik ta'minot" qatnashgan atama-ta'rif juftliklaridan to'g'risini toping.
\begin{enumerate}[label=\Alph*),leftmargin=9mm]
\item texnik ta'minot - bir maqsad uchun hamkorlik qiluvchi texnik vosita, dastur, ma'lumot va foydalanuvchi majmuasi
\item texnik ta'minot - kompyuterga vazifalarni bajarishni ko'rsatuvchi dasturlar
\item texnik ta'minot - kiritish, qayta ishlash, chiqarish va saqlash modeli
\item texnik ta'minot - kompyuter tizimining fizik qismlari
\end{enumerate}
\begin{tcolorbox}[enhanced,breakable,colback=teal!4,colframe=teal!60!black,boxrule=0.45pt,arc=1.2mm,left=2.7mm,right=2.7mm,top=1.8mm,bottom=1.8mm,before skip=2mm,after skip=2mm]
{\bfseries\sffamily\color{teal}To'g'ri javob:} \textbf{D) texnik ta'minot - kompyuter tizimining fizik qismlari}
\par\smallskip
{\bfseries Izoh:} texnik ta'minot - kompyuter tizimining fizik qismlari. To'g'ri juftlik: texnik ta'minot - kompyuter tizimining fizik qismlari. Manba: CAM10/11, PDF 31-35; Qo'shimcha sahifa raqami (qaysi PDFga tegishli ekani aniqlanmagan; darslik sahifasi deb tasdiqlanmagan): 4
\end{tcolorbox}
\vspace{2mm}\hrule height 0.2pt\vspace{3mm}

% DB_ID fedde812-0a34-a7d8-534c-d7d77c02f0f6
\needspace{5\baselineskip}
{\bfseries\sffamily\color{ink} 211.} {scriptsizecolor{muted} Y1 | review | fedde812}
\par
Markaziy protsessor va bajarish sikli bo'yicha qatorlar ichidan "CPU" noto'g'ri izohlanganini aniqlang.
\begin{enumerate}[label=\Alph*),leftmargin=9mm]
\item CU - ko'rsatmalar oqimi va boshqa qismlar ishini boshqaruvchi qism
\item registr - CPU ichidagi juda tez va kichik xotira
\item takt chastotasi - sekunddagi takt impulslar soni
\item CPU - arifmetik va mantiqiy amallarni bajaruvchi qism
\end{enumerate}
\begin{tcolorbox}[enhanced,breakable,colback=teal!4,colframe=teal!60!black,boxrule=0.45pt,arc=1.2mm,left=2.7mm,right=2.7mm,top=1.8mm,bottom=1.8mm,before skip=2mm,after skip=2mm]
{\bfseries\sffamily\color{teal}To'g'ri javob:} \textbf{D) CPU - arifmetik va mantiqiy amallarni bajaruvchi qism}
\par\smallskip
{\bfseries Izoh:} CPU - arifmetik va mantiqiy amallarni bajaruvchi qism. CPU uchun to'g'ri ta'rif: ko'rsatmalarni bajarib, ma'lumotlarni qayta ishlovchi markaziy qurilma. Manba: CAM9, PDF 10-14; Qo'shimcha sahifa raqami (qaysi PDFga tegishli ekani aniqlanmagan; darslik sahifasi deb tasdiqlanmagan): 17
\end{tcolorbox}
\vspace{2mm}\hrule height 0.2pt\vspace{3mm}

% DB_ID 02a1048c-56be-51f7-5fac-d9d18d268b8f
\needspace{5\baselineskip}
{\bfseries\sffamily\color{ink} 212.} {scriptsizecolor{muted} Y1 | review | 02a1048c}
\par
Kiritish qurilmalari mavzusida: I. sensor - fizik kattalikni o'lchab signalga aylantiruvchi qurilma. II. OCR - tasvirdagi bosma matnni tahrirlanadigan belgilar sifatida tanish. Mulohazalarni baholang.
\begin{enumerate}[label=\Alph*),leftmargin=9mm]
\item faqat I
\item faqat II
\item I va II
\item I ham, II ham emas
\end{enumerate}
\begin{tcolorbox}[enhanced,breakable,colback=teal!4,colframe=teal!60!black,boxrule=0.45pt,arc=1.2mm,left=2.7mm,right=2.7mm,top=1.8mm,bottom=1.8mm,before skip=2mm,after skip=2mm]
{\bfseries\sffamily\color{teal}To'g'ri javob:} \textbf{C) I va II}
\par\smallskip
{\bfseries Izoh:} I va II. Birinchi va ikkinchi mulohazaning ta'riflari alohida tekshirildi; natija: I va II. Manba: ICT10, PDF 8-12; Qo'shimcha sahifa raqami (qaysi PDFga tegishli ekani aniqlanmagan; darslik sahifasi deb tasdiqlanmagan): 41
\end{tcolorbox}
\vspace{2mm}\hrule height 0.2pt\vspace{3mm}

% DB_ID 0a71de96-1617-ef6e-a429-f49205a58afc
\needspace{5\baselineskip}
{\bfseries\sffamily\color{ink} 213.} {scriptsizecolor{muted} Y1 | review | 0a71de96}
\par
Arxiv markazi magnit lentaga yozilgan katta faylning oxiridagi ma'lumotni topishga qiynalmoqda. Qaysi xususiyat buni izohlaydi?
\begin{enumerate}[label=\Alph*),leftmargin=9mm]
\item HDD faqat ketma-ket o'qishga mos bo'lgani uchun lenta bilan bir xil ishlaydi
\item Magnit lenta doimiy xotira bo'lmagani uchun barcha fayllar faqat RAMda turadi
\item Magnit lenta odatda ketma-ket murojaat qiladi; kerakli joyga yetish vaqt oladi
\item SSD tasodifiy murojaatni qo'llamagani uchun lenta undan har doim tezroq ishlaydi
\end{enumerate}
\begin{tcolorbox}[enhanced,breakable,colback=teal!4,colframe=teal!60!black,boxrule=0.45pt,arc=1.2mm,left=2.7mm,right=2.7mm,top=1.8mm,bottom=1.8mm,before skip=2mm,after skip=2mm]
{\bfseries\sffamily\color{teal}To'g'ri javob:} \textbf{C) Magnit lenta odatda ketma-ket murojaat qiladi; kerakli joyga yetish vaqt oladi}
\par\smallskip
{\bfseries Izoh:} Darslik qattiq diskdagi bevosita murojaat va magnit lentadagi ketma-ket o'qishni farqlaydi. Lentada kerakli joyga yetguncha material bo'ylab harakatlanish talab etilishi mumkin. [Manba: Cambridge+ 9-sinf, PDF 42-43.]
\end{tcolorbox}
\vspace{2mm}\hrule height 0.2pt\vspace{3mm}

% DB_ID 14d19622-8ef9-c248-58b6-f1499a6d1932
\needspace{5\baselineskip}
{\bfseries\sffamily\color{ink} 214.} {scriptsizecolor{muted} Y1 | review | 14d19622}
\par
Arxitektor yangi uyning katta formatdagi rangli texnik chizmasini chop etishi kerak. Qaysi qurilmani tanlash maqsadga muvofiq?
\begin{enumerate}[label=\Alph*),leftmargin=9mm]
\item Keng formatli plotter, chunki u texnik chizmalarni chiqarishga mos
\item Kvitansiya printeri, chunki u har bir chizmani bir xil kenglikda chiqaradi
\item Shtrix-kod printeri, chunki u katta grafik xaritalarni chop etishga mos
\item Oddiy matritsali printer, chunki u turli qog'ozda har qanday rasmni chiqaradi
\end{enumerate}
\begin{tcolorbox}[enhanced,breakable,colback=teal!4,colframe=teal!60!black,boxrule=0.45pt,arc=1.2mm,left=2.7mm,right=2.7mm,top=1.8mm,bottom=1.8mm,before skip=2mm,after skip=2mm]
{\bfseries\sffamily\color{teal}To'g'ri javob:} \textbf{A) Keng formatli plotter, chunki u texnik chizmalarni chiqarishga mos}
\par\smallskip
{\bfseries Izoh:} Cambridge+ darsligida arxitektorning katta formatdagi rangli uy chizmalarini chiqarishi uchun plotter qurilmasidan foydalanish misoli keltiriladi. [Manba: Cambridge+ 10-11-sinf, PDF 44-bet.]
\end{tcolorbox}
\vspace{2mm}\hrule height 0.2pt\vspace{3mm}

% DB_ID 15a061e3-c587-9fea-7720-a6279fda0ab6
\needspace{5\baselineskip}
{\bfseries\sffamily\color{ink} 215.} {scriptsizecolor{muted} Y1 | review | 15a061e3}
\par
Ko'p ko'chiriladigan noutbuk uchun harakatlanuvchi mexanik qismlari bo'lmagan saqlash vositasi kerak. Qaysi yechim mos?
\begin{enumerate}[label=\Alph*),leftmargin=9mm]
\item SSD flesh xotira orqali ishlaydi va unda harakatlanuvchi magnit disk plastinkalari yo'q
\item HDD magnit plastinkalarsiz ishlaydi, uni faqat elektron xotira hujayralari bilan ishlab chiqarishadi
\item Magnit lenta barcha sektorlar bo'ylab ixtiyoriy tez kirishga imkon beradigan yuqori tezlikli xotiradir
\item Optik disk ma'lumotni RAMga nisbatan tezroq o'qiydi va doimiy saqlash uchun hech qanday qurilma talab qilmaydi
\end{enumerate}
\begin{tcolorbox}[enhanced,breakable,colback=teal!4,colframe=teal!60!black,boxrule=0.45pt,arc=1.2mm,left=2.7mm,right=2.7mm,top=1.8mm,bottom=1.8mm,before skip=2mm,after skip=2mm]
{\bfseries\sffamily\color{teal}To'g'ri javob:} \textbf{A) SSD flesh xotira orqali ishlaydi va unda harakatlanuvchi magnit disk plastinkalari yo'q}
\par\smallskip
{\bfseries Izoh:} SSDda mexanik harakatlanuvchi qismlar yo'q, u flesh xotiraga asoslanadi. [Manba: Cambridge+ 10-11-sinf PDF 35, 41-43.]
\end{tcolorbox}
\vspace{2mm}\hrule height 0.2pt\vspace{3mm}

% DB_ID 16e9ccee-c729-a81c-9958-23c01693c4b6
\needspace{5\baselineskip}
{\bfseries\sffamily\color{ink} 216.} {scriptsizecolor{muted} Y1 | review | 16e9ccee}
\par
Protsessor juda qisqa muddatga joriy amal bilan bog'liq kichik ikkilik qiymatlarni saqlashi kerak. Bu vazifaga mos ichki elementni tanlang.
\begin{enumerate}[label=\Alph*),leftmargin=9mm]
\item Optik disk, chunki u protsessor ichida tez ishlaydigan vaqtinchalik xotiradir
\item Printer buferi, chunki u protsessordagi barcha hisoblash qiymatlarini saqlaydi
\item Registr, chunki u protsessor amallarida qiymatlarni vaqtincha saqlaydi
\item Tarmoq kabeli, chunki u protsessor ichidagi kichik qiymatlarni xotiraga yozadi
\end{enumerate}
\begin{tcolorbox}[enhanced,breakable,colback=teal!4,colframe=teal!60!black,boxrule=0.45pt,arc=1.2mm,left=2.7mm,right=2.7mm,top=1.8mm,bottom=1.8mm,before skip=2mm,after skip=2mm]
{\bfseries\sffamily\color{teal}To'g'ri javob:} \textbf{C) Registr, chunki u protsessor amallarida qiymatlarni vaqtincha saqlaydi}
\par\smallskip
{\bfseries Izoh:} Qo'llanmada registrlar ma'lumotni ikkilik shaklda saqlab turishga xizmat qilishi bayon qilingan. Protsessor ichidagi registrni optik disk yoki periferiya bilan adashtirmaslik kerak. [Manba: Ma'ruza (Tayyor), PDF 5-7.]
\end{tcolorbox}
\vspace{2mm}\hrule height 0.2pt\vspace{3mm}

% DB_ID 1a703cc8-668b-cd31-e1f7-2d7b690d4cab
\needspace{5\baselineskip}
{\bfseries\sffamily\color{ink} 217.} {scriptsizecolor{muted} Y1 | review | 1a703cc8}
\par
Tizimli blok, bosh plata va ulanishlar bo'yicha qatorlar ichidan "shina" noto'g'ri izohlanganini aniqlang.
\begin{enumerate}[label=\Alph*),leftmargin=9mm]
\item slot - ichki komponent o'rnatiladigan ulagich
\item port - tashqi qurilma ulanadigan nuqta
\item shina - ichki komponentlarni bog'lovchi asosiy bosma plata
\item quvvat bloki - tarmoq elektrini komponentlarga mos doimiy tokka aylantiruvchi qism
\end{enumerate}
\begin{tcolorbox}[enhanced,breakable,colback=teal!4,colframe=teal!60!black,boxrule=0.45pt,arc=1.2mm,left=2.7mm,right=2.7mm,top=1.8mm,bottom=1.8mm,before skip=2mm,after skip=2mm]
{\bfseries\sffamily\color{teal}To'g'ri javob:} \textbf{C) shina - ichki komponentlarni bog'lovchi asosiy bosma plata}
\par\smallskip
{\bfseries Izoh:} shina - ichki komponentlarni bog'lovchi asosiy bosma plata. shina uchun to'g'ri ta'rif: ma'lumot va boshqaruv signallarini tashuvchi aloqa yo'llari. Manba: CAM10/11, PDF 32-36; Qo'shimcha sahifa raqami (qaysi PDFga tegishli ekani aniqlanmagan; darslik sahifasi deb tasdiqlanmagan): 11
\end{tcolorbox}
\vspace{2mm}\hrule height 0.2pt\vspace{3mm}

% DB_ID 1b3bcfcf-99db-1665-74d9-a9ce744f3bd3
\needspace{5\baselineskip}
{\bfseries\sffamily\color{ink} 218.} {scriptsizecolor{muted} Y1 | review | 1b3bcfcf}
\par
Birlamchi xotira: registr, kesh, RAM va ROM mavzusida: I. kesh - tez-tez kerak bo'ladigan ma'lumotni CPUga yaqin saqlovchi xotira. II. RAM - tez-tez kerak bo'ladigan ma'lumotni CPUga yaqin saqlovchi xotira. Mulohazalarni baholang.
\begin{enumerate}[label=\Alph*),leftmargin=9mm]
\item faqat I
\item faqat II
\item I va II
\item I ham, II ham emas
\end{enumerate}
\begin{tcolorbox}[enhanced,breakable,colback=teal!4,colframe=teal!60!black,boxrule=0.45pt,arc=1.2mm,left=2.7mm,right=2.7mm,top=1.8mm,bottom=1.8mm,before skip=2mm,after skip=2mm]
{\bfseries\sffamily\color{teal}To'g'ri javob:} \textbf{A) faqat I}
\par\smallskip
{\bfseries Izoh:} faqat I. Birinchi va ikkinchi mulohazaning ta'riflari alohida tekshirildi; natija: faqat I. Manba: ICT10, PDF 15-20; Qo'shimcha sahifa raqami (qaysi PDFga tegishli ekani aniqlanmagan; darslik sahifasi deb tasdiqlanmagan): 26
\end{tcolorbox}
\vspace{2mm}\hrule height 0.2pt\vspace{3mm}

% DB_ID 1fe6460b-291e-3c75-a65f-bd4072df5899
\needspace{5\baselineskip}
{\bfseries\sffamily\color{ink} 219.} {scriptsizecolor{muted} Y1 | review | 1fe6460b}
\par
Birlamchi xotira: registr, kesh, RAM va ROM mavzusida: I. firmware - qurilma ichiga o'rnatilgan past darajali boshqaruv dasturi. II. RAM - faol dastur va ma'lumotlar uchun o'zgaruvchan tezkor xotira. Mulohazalarni baholang.
\begin{enumerate}[label=\Alph*),leftmargin=9mm]
\item faqat I
\item faqat II
\item I va II
\item I ham, II ham emas
\end{enumerate}
\begin{tcolorbox}[enhanced,breakable,colback=teal!4,colframe=teal!60!black,boxrule=0.45pt,arc=1.2mm,left=2.7mm,right=2.7mm,top=1.8mm,bottom=1.8mm,before skip=2mm,after skip=2mm]
{\bfseries\sffamily\color{teal}To'g'ri javob:} \textbf{C) I va II}
\par\smallskip
{\bfseries Izoh:} I va II. Birinchi va ikkinchi mulohazaning ta'riflari alohida tekshirildi; natija: I va II. Manba: CAM9, PDF 11-14; Qo'shimcha sahifa raqami (qaysi PDFga tegishli ekani aniqlanmagan; darslik sahifasi deb tasdiqlanmagan): 28
\end{tcolorbox}
\vspace{2mm}\hrule height 0.2pt\vspace{3mm}

% DB_ID 26ea2636-a8d3-0a60-b20c-0bd1e34d5aed
\needspace{5\baselineskip}
{\bfseries\sffamily\color{ink} 220.} {scriptsizecolor{muted} Y1 | review | 26ea2636}
\par
Ikkilamchi saqlash qurilmalari mavzusida: I. HDD - magnit plastinkali, katta sig'imli ikkilamchi xotira. II. SSD - flesh hujayralarga asoslangan harakatsiz ikkilamchi xotira. Mulohazalarni baholang.
\begin{enumerate}[label=\Alph*),leftmargin=9mm]
\item I va II
\item faqat I
\item faqat II
\item I ham, II ham emas
\end{enumerate}
\begin{tcolorbox}[enhanced,breakable,colback=teal!4,colframe=teal!60!black,boxrule=0.45pt,arc=1.2mm,left=2.7mm,right=2.7mm,top=1.8mm,bottom=1.8mm,before skip=2mm,after skip=2mm]
{\bfseries\sffamily\color{teal}To'g'ri javob:} \textbf{A) I va II}
\par\smallskip
{\bfseries Izoh:} I va II. Birinchi va ikkinchi mulohazaning ta'riflari alohida tekshirildi; natija: I va II. Manba: CAM9, PDF 33-47; Qo'shimcha sahifa raqami (qaysi PDFga tegishli ekani aniqlanmagan; darslik sahifasi deb tasdiqlanmagan): 31
\end{tcolorbox}
\vspace{2mm}\hrule height 0.2pt\vspace{3mm}

% DB_ID 2a56eb13-ff2a-70c0-d0ad-125a15801b7b
\needspace{5\baselineskip}
{\bfseries\sffamily\color{ink} 221.} {scriptsizecolor{muted} Y1 | review | 2a56eb13}
\par
Attestatsiyadagi nozik farqlar va integratsion masalalar mavzusida: I. distraktor - ishonarli, ammo noto'g'ri javob varianti. II. tasnif - obyektni muhim belgilari bo'yicha guruhga ajratish. Mulohazalarni baholang.
\begin{enumerate}[label=\Alph*),leftmargin=9mm]
\item faqat I
\item faqat II
\item I va II
\item I ham, II ham emas
\end{enumerate}
\begin{tcolorbox}[enhanced,breakable,colback=teal!4,colframe=teal!60!black,boxrule=0.45pt,arc=1.2mm,left=2.7mm,right=2.7mm,top=1.8mm,bottom=1.8mm,before skip=2mm,after skip=2mm]
{\bfseries\sffamily\color{teal}To'g'ri javob:} \textbf{C) I va II}
\par\smallskip
{\bfseries Izoh:} I va II. Birinchi va ikkinchi mulohazaning ta'riflari alohida tekshirildi; natija: I va II. Manba: MAR, PDF 4-20; Qo'shimcha sahifa raqami (qaysi PDFga tegishli ekani aniqlanmagan; darslik sahifasi deb tasdiqlanmagan): 120
\end{tcolorbox}
\vspace{2mm}\hrule height 0.2pt\vspace{3mm}

% DB_ID 2bdbf9a4-13e0-9f82-9ea4-33ad59eea4c2
\needspace{5\baselineskip}
{\bfseries\sffamily\color{ink} 222.} {scriptsizecolor{muted} Y1 | review | 2bdbf9a4}
\par
CPU buyruq bajarish jarayonida ikki sonni qo'shib, natija ma'lum chegaradan katta-kichikligini solishtirdi. Qaysi ichki qism shu amallarni bajaradi?
\begin{enumerate}[label=\Alph*),leftmargin=9mm]
\item ALU arifmetik va mantiqiy amallarni bajaruvchi birlik sifatida ishlaydi
\item CU xotiradagi har bir faylni doimiy saqlovchi qurilma sifatida ishlaydi
\item Registr tashqi qurilmadan kelgan analog signalni tasvirga aylantiradi
\item Kesh kompyuterning elektr quvvatini doimiy kuchlanishga aylantiradi
\end{enumerate}
\begin{tcolorbox}[enhanced,breakable,colback=teal!4,colframe=teal!60!black,boxrule=0.45pt,arc=1.2mm,left=2.7mm,right=2.7mm,top=1.8mm,bottom=1.8mm,before skip=2mm,after skip=2mm]
{\bfseries\sffamily\color{teal}To'g'ri javob:} \textbf{A) ALU arifmetik va mantiqiy amallarni bajaruvchi birlik sifatida ishlaydi}
\par\smallskip
{\bfseries Izoh:} ALU hisob-kitob va mantiqiy operatsiyalarni bajaradi. CU ko'rsatmalar oqimini muvofiqlashtiradi; registr ma'lumot saqlashga, kesh esa tezkor murojaatga xizmat qiladi. [Manba: Cambridge+ 10-11-sinf, PDF 33.]
\end{tcolorbox}
\vspace{2mm}\hrule height 0.2pt\vspace{3mm}

% DB_ID 2c3dd5e6-0285-81a8-d472-feaef0c252c1
\needspace{5\baselineskip}
{\bfseries\sffamily\color{ink} 223.} {scriptsizecolor{muted} Y1 | review | 2c3dd5e6}
\par
Ofis oyiga o'n minglab oq-qora matn sahifa chiqaradi. Eng mos tanlov qaysi?
\begin{enumerate}[label=\Alph*),leftmargin=9mm]
\item dupleks lazer printer
\item kichik foto siyohpurkagich
\item 3D printer
\item grafik planshet
\end{enumerate}
\begin{tcolorbox}[enhanced,breakable,colback=teal!4,colframe=teal!60!black,boxrule=0.45pt,arc=1.2mm,left=2.7mm,right=2.7mm,top=1.8mm,bottom=1.8mm,before skip=2mm,after skip=2mm]
{\bfseries\sffamily\color{teal}To'g'ri javob:} \textbf{A) dupleks lazer printer}
\par\smallskip
{\bfseries Izoh:} dupleks lazer printer. Lazer printer katta hajmli matnni tez va past sahifa tannarxida chiqarishga mos. Manba: CAM9, PDF 16-32; Qo'shimcha sahifa raqami (qaysi PDFga tegishli ekani aniqlanmagan; darslik sahifasi deb tasdiqlanmagan): 48
\end{tcolorbox}
\vspace{2mm}\hrule height 0.2pt\vspace{3mm}

% DB_ID 313c9b46-c6f9-7457-1a62-a907b1a9975d
\needspace{5\baselineskip}
{\bfseries\sffamily\color{ink} 224.} {scriptsizecolor{muted} Y1 | review | 313c9b46}
\par
Kompyuter tizimi va uning ishlash modeli mavzusida: I. dasturiy ta'minot - kompyuterga vazifalarni bajarishni ko'rsatuvchi dasturlar. II. kompyuter tizimi - kompyuterga vazifalarni bajarishni ko'rsatuvchi dasturlar. Mulohazalarni baholang.
\begin{enumerate}[label=\Alph*),leftmargin=9mm]
\item faqat II
\item I va II
\item faqat I
\item I ham, II ham emas
\end{enumerate}
\begin{tcolorbox}[enhanced,breakable,colback=teal!4,colframe=teal!60!black,boxrule=0.45pt,arc=1.2mm,left=2.7mm,right=2.7mm,top=1.8mm,bottom=1.8mm,before skip=2mm,after skip=2mm]
{\bfseries\sffamily\color{teal}To'g'ri javob:} \textbf{C) faqat I}
\par\smallskip
{\bfseries Izoh:} faqat I. Birinchi va ikkinchi mulohazaning ta'riflari alohida tekshirildi; natija: faqat I. Manba: ICT5, PDF 25-32; Qo'shimcha sahifa raqami (qaysi PDFga tegishli ekani aniqlanmagan; darslik sahifasi deb tasdiqlanmagan): 6
\end{tcolorbox}
\vspace{2mm}\hrule height 0.2pt\vspace{3mm}

% DB_ID 332c942e-0502-1926-0e30-f2fd96bd809b
\needspace{5\baselineskip}
{\bfseries\sffamily\color{ink} 225.} {scriptsizecolor{muted} Y1 | review | 332c942e}
\par
Supermarketdagi sotuvchi mahsulot ustidagi parallel qora-oq chiziqlardan mahsulot identifikatorini tez kiritishi kerak. Qaysi uskuna mos?
\begin{enumerate}[label=\Alph*),leftmargin=9mm]
\item Shtrix-kod o'quvchi chiziqlar orqali kodlangan identifikatorni kiritadi
\item OCR dasturi mahsulotdagi har bir rasmdan to'liq elektron jadval yaratadi
\item OMR qurilmasi koddagi raqamlarni magnitli siyoh yordamida belgilaydi
\item MICR o'quvchi chiziq uzunligidan mahsulot vaznini bevosita aniqlaydi
\end{enumerate}
\begin{tcolorbox}[enhanced,breakable,colback=teal!4,colframe=teal!60!black,boxrule=0.45pt,arc=1.2mm,left=2.7mm,right=2.7mm,top=1.8mm,bottom=1.8mm,before skip=2mm,after skip=2mm]
{\bfseries\sffamily\color{teal}To'g'ri javob:} \textbf{A) Shtrix-kod o'quvchi chiziqlar orqali kodlangan identifikatorni kiritadi}
\par\smallskip
{\bfseries Izoh:} Shtrix-kod rider turli kenglikdagi parallel chiziqli kodlardan identifikatorni avtomatik o'qish uchun mos; OCR/OMR/MICR boshqa maqsadlarda ishlatiladi. [Manba: Cambridge+ 9-sinf, PDF 33.]
\end{tcolorbox}
\vspace{2mm}\hrule height 0.2pt\vspace{3mm}

% DB_ID 3d96448e-6ae9-5fa0-9bbe-c242411501ab
\needspace{5\baselineskip}
{\bfseries\sffamily\color{ink} 226.} {scriptsizecolor{muted} Y1 | review | 3d96448e}
\par
Chiqarish va ikki vazifali qurilmalar mavzusida: I. aktuator - boshqaruv signalini fizik harakat yoki ta'sirga aylantiruvchi qurilma. II. yumshoq nusxa - ekranda yoki audio orqali vaqtinchalik taqdim etilgan natija. Mulohazalarni baholang.
\begin{enumerate}[label=\Alph*),leftmargin=9mm]
\item I va II
\item faqat I
\item faqat II
\item I ham, II ham emas
\end{enumerate}
\begin{tcolorbox}[enhanced,breakable,colback=teal!4,colframe=teal!60!black,boxrule=0.45pt,arc=1.2mm,left=2.7mm,right=2.7mm,top=1.8mm,bottom=1.8mm,before skip=2mm,after skip=2mm]
{\bfseries\sffamily\color{teal}To'g'ri javob:} \textbf{A) I va II}
\par\smallskip
{\bfseries Izoh:} I va II. Birinchi va ikkinchi mulohazaning ta'riflari alohida tekshirildi; natija: I va II. Manba: MAR, PDF 8-10; Qo'shimcha sahifa raqami (qaysi PDFga tegishli ekani aniqlanmagan; darslik sahifasi deb tasdiqlanmagan): 47
\end{tcolorbox}
\vspace{2mm}\hrule height 0.2pt\vspace{3mm}

% DB_ID 493862f7-8759-d2d6-eb5c-2ebf2d5f0b81
\needspace{5\baselineskip}
{\bfseries\sffamily\color{ink} 227.} {scriptsizecolor{muted} Y1 | review | 493862f7}
\par
8 GB RAM to'lganidan keyin kompyuterda SSDga murojaat ko'paydi va ilovalar sekinlashdi. Bu holatning eng mantiqiy izohi qaysi?
\begin{enumerate}[label=\Alph*),leftmargin=9mm]
\item OT SSDdagi virtual xotiradan foydalanmoqda; unga murojaat RAMdan sekinroq
\item OT ROMdagi yuklash ko'rsatmalarini SSDga o'tkazmoqda; bu RAMni oshiradi
\item OT keshni doimiy saqlash diskiga aylantirmoqda; uning tezligi RAMdan yuqori
\item OT barcha ilovalarni grafik xotiraga ko'chirmoqda; diskdan foydalanish pasayadi
\end{enumerate}
\begin{tcolorbox}[enhanced,breakable,colback=teal!4,colframe=teal!60!black,boxrule=0.45pt,arc=1.2mm,left=2.7mm,right=2.7mm,top=1.8mm,bottom=1.8mm,before skip=2mm,after skip=2mm]
{\bfseries\sffamily\color{teal}To'g'ri javob:} \textbf{A) OT SSDdagi virtual xotiradan foydalanmoqda; unga murojaat RAMdan sekinroq}
\par\smallskip
{\bfseries Izoh:} Darslik RAM to'lganda HDD yoki SSDda virtual xotira maydonidan foydalanilishi, ammo ikkilamchi xotiraga murojaat RAMdan ancha sekin ekanini tushuntiradi. [Manba: Cambridge+ 10-11-sinf, PDF 34.]
\end{tcolorbox}
\vspace{2mm}\hrule height 0.2pt\vspace{3mm}

% DB_ID 4d4351ef-6827-1d63-6110-1660e6d062f9
\needspace{5\baselineskip}
{\bfseries\sffamily\color{ink} 228.} {scriptsizecolor{muted} Y1 | review | 4d4351ef}
\par
Birlamchi xotira: registr, kesh, RAM va ROM bo'yicha qatorlar ichidan "ROM" noto'g'ri izohlanganini aniqlang.
\begin{enumerate}[label=\Alph*),leftmargin=9mm]
\item kesh - tez-tez kerak bo'ladigan ma'lumotni CPUga yaqin saqlovchi xotira
\item virtual xotira - RAM yetishmaganda ikkilamchi xotiradan ajratilgan sekinroq maydon
\item ROM - faol dastur va ma'lumotlar uchun o'zgaruvchan tezkor xotira
\item firmware - qurilma ichiga o'rnatilgan past darajali boshqaruv dasturi
\end{enumerate}
\begin{tcolorbox}[enhanced,breakable,colback=teal!4,colframe=teal!60!black,boxrule=0.45pt,arc=1.2mm,left=2.7mm,right=2.7mm,top=1.8mm,bottom=1.8mm,before skip=2mm,after skip=2mm]
{\bfseries\sffamily\color{teal}To'g'ri javob:} \textbf{C) ROM - faol dastur va ma'lumotlar uchun o'zgaruvchan tezkor xotira}
\par\smallskip
{\bfseries Izoh:} ROM - faol dastur va ma'lumotlar uchun o'zgaruvchan tezkor xotira. ROM uchun to'g'ri ta'rif: elektrsiz ham ma'lumotni saqlovchi o'zgarmas xotira. Manba: CAM10/11, PDF 33-36; Qo'shimcha sahifa raqami (qaysi PDFga tegishli ekani aniqlanmagan; darslik sahifasi deb tasdiqlanmagan): 24
\end{tcolorbox}
\vspace{2mm}\hrule height 0.2pt\vspace{3mm}

% DB_ID 53cbbf53-9abf-1f21-feea-5696a2a7f727
\needspace{5\baselineskip}
{\bfseries\sffamily\color{ink} 229.} {scriptsizecolor{muted} Y1 | review | 53cbbf53}
\par
RAM to'lib, tizim SSDdagi virtual xotiradan faol foydalanmoqda. Eng kutiladigan natija qaysi?
\begin{enumerate}[label=\Alph*),leftmargin=9mm]
\item barcha fayllar ROMga ko'chadi
\item CPU takt chastotasi doim oshadi
\item dasturlar ishlaydi, ammo sezilarli sekinlashadi
\item monitor aniqligi avtomatik kamayadi
\end{enumerate}
\begin{tcolorbox}[enhanced,breakable,colback=teal!4,colframe=teal!60!black,boxrule=0.45pt,arc=1.2mm,left=2.7mm,right=2.7mm,top=1.8mm,bottom=1.8mm,before skip=2mm,after skip=2mm]
{\bfseries\sffamily\color{teal}To'g'ri javob:} \textbf{C) dasturlar ishlaydi, ammo sezilarli sekinlashadi}
\par\smallskip
{\bfseries Izoh:} dasturlar ishlaydi, ammo sezilarli sekinlashadi. Virtual xotira RAMdan sekin ikkilamchi xotirada joylashadi. Manba: CAM9, PDF 11-14; Qo'shimcha sahifa raqami (qaysi PDFga tegishli ekani aniqlanmagan; darslik sahifasi deb tasdiqlanmagan): 29
\end{tcolorbox}
\vspace{2mm}\hrule height 0.2pt\vspace{3mm}

% DB_ID 54fb8a44-a8d3-7899-6753-91584d1b3679
\needspace{5\baselineskip}
{\bfseries\sffamily\color{ink} 230.} {scriptsizecolor{muted} Y1 | review | 54fb8a44}
\par
Unumdorlik, nosozliklarni aniqlash va xavfsiz xizmat mavzusida: I. bottleneck - tizim umumiy tezligini cheklayotgan eng sust resurs. II. paging - RAM va ikkilamchi xotira o'rtasida sahifalarni almashish. Mulohazalarni baholang.
\begin{enumerate}[label=\Alph*),leftmargin=9mm]
\item I va II
\item faqat I
\item faqat II
\item I ham, II ham emas
\end{enumerate}
\begin{tcolorbox}[enhanced,breakable,colback=teal!4,colframe=teal!60!black,boxrule=0.45pt,arc=1.2mm,left=2.7mm,right=2.7mm,top=1.8mm,bottom=1.8mm,before skip=2mm,after skip=2mm]
{\bfseries\sffamily\color{teal}To'g'ri javob:} \textbf{A) I va II}
\par\smallskip
{\bfseries Izoh:} I va II. Birinchi va ikkinchi mulohazaning ta'riflari alohida tekshirildi; natija: I va II. Manba: CAM9, PDF 12-47; Qo'shimcha sahifa raqami (qaysi PDFga tegishli ekani aniqlanmagan; darslik sahifasi deb tasdiqlanmagan): 109
\end{tcolorbox}
\vspace{2mm}\hrule height 0.2pt\vspace{3mm}

% DB_ID 55484bed-33ff-28c1-ac7a-4b49326028e2
\needspace{5\baselineskip}
{\bfseries\sffamily\color{ink} 231.} {scriptsizecolor{muted} Y1 | review | 55484bed}
\par
QR-kodlar shtrix-kodlardan farqli ravishda kengroq turdagi ma'lumotni o'z ichiga olishi mumkin. Restoran menyusining veb manzilini tarqatish uchun qaysi yondashuv ma'qul?
\begin{enumerate}[label=\Alph*),leftmargin=9mm]
\item QR-kodni kamera orqali o'qib, unda berilgan veb manzilni ochish
\item Oddiy narx yorlig'ini bosib, uning rangidan veb manzilni aniqlash
\item Magnitli siyoh raqamini tanib, undan menyu rasmini tiklash
\item Shtrix-kod o'rniga printer drayverini tanlab, unda menyuni kodlash
\end{enumerate}
\begin{tcolorbox}[enhanced,breakable,colback=teal!4,colframe=teal!60!black,boxrule=0.45pt,arc=1.2mm,left=2.7mm,right=2.7mm,top=1.8mm,bottom=1.8mm,before skip=2mm,after skip=2mm]
{\bfseries\sffamily\color{teal}To'g'ri javob:} \textbf{A) QR-kodni kamera orqali o'qib, unda berilgan veb manzilni ochish}
\par\smallskip
{\bfseries Izoh:} Darslik QR-kodni shtrix-koddan ko'proq turdagi axborot elementlarini taqdim eta oluvchi kod sifatida muhokama qiladi. Unda veb manzil kodlash amaliy misoldir. [Manba: Cambridge+ 10-11-sinf, PDF 39-bet.]
\end{tcolorbox}
\vspace{2mm}\hrule height 0.2pt\vspace{3mm}

% DB_ID 6645ada2-d94b-ace4-d841-c28fc146f074
\needspace{5\baselineskip}
{\bfseries\sffamily\color{ink} 232.} {scriptsizecolor{muted} Y1 | review | 6645ada2}
\par
Kompyuter yig'ilayotganda mutaxassis qaysi komponentlar apparat (texnik) ta'minotga kirishini ajratishi kerak. Qaysi tanlovda faqat fizik qismlar bor?
\begin{enumerate}[label=\Alph*),leftmargin=9mm]
\item Protsessor, operativ xotira, tizim platasi va klaviatura
\item Operatsion tizim, xotira, matn muharriri va klaviatura
\item Tizim platasi, antivirus, videokarta va drayver dasturi
\item Protsessor, brauzer, hisoblash jadvali va printer drayveri
\end{enumerate}
\begin{tcolorbox}[enhanced,breakable,colback=teal!4,colframe=teal!60!black,boxrule=0.45pt,arc=1.2mm,left=2.7mm,right=2.7mm,top=1.8mm,bottom=1.8mm,before skip=2mm,after skip=2mm]
{\bfseries\sffamily\color{teal}To'g'ri javob:} \textbf{A) Protsessor, operativ xotira, tizim platasi va klaviatura}
\par\smallskip
{\bfseries Izoh:} Darslik texnik ta'minotni kompyuterning fizik komponentlari deb ta'riflaydi va protsessor, xotira, plata hamda tashqi qurilmalarni misol keltiradi. Dastur yoki drayver qatnashgan guruhlar faqat hardware emas. [Manba: Cambridge+ 10-11-sinf, PDF 33.]
\end{tcolorbox}
\vspace{2mm}\hrule height 0.2pt\vspace{3mm}

% DB_ID 6baa2fc2-a604-21cb-b857-95a8b62d5e51
\needspace{5\baselineskip}
{\bfseries\sffamily\color{ink} 233.} {scriptsizecolor{muted} Y1 | review | 6baa2fc2}
\par
Imtihon varaqasida javoblar oldindan bosilgan A-D kataklarini bo'yash orqali belgilanadi. Tizim matnni emas, aynan belgining joyini topishi kerak. Qaysi texnologiya mos?
\begin{enumerate}[label=\Alph*),leftmargin=9mm]
\item OCR, chunki u hujjatdagi bosma harflarni tahrirlanadigan matnga aylantiradi
\item MICR, chunki u chekdagi magnitli belgilarni o'qish uchun ishlatiladi
\item OMR, chunki u blankdagi belgilangan kataklarning joylashuvini tekshiradi
\item Shtrix-kod o'quvchi, chunki u parallel chiziqlar orqali kodni o'qiydi
\end{enumerate}
\begin{tcolorbox}[enhanced,breakable,colback=teal!4,colframe=teal!60!black,boxrule=0.45pt,arc=1.2mm,left=2.7mm,right=2.7mm,top=1.8mm,bottom=1.8mm,before skip=2mm,after skip=2mm]
{\bfseries\sffamily\color{teal}To'g'ri javob:} \textbf{C) OMR, chunki u blankdagi belgilangan kataklarning joylashuvini tekshiradi}
\par\smallskip
{\bfseries Izoh:} OMR maxsus blankdagi belgi qo'yilgan joylarni aniqlaydi, belgining harf shaklini OCR kabi tanimaydi. Darslik aynan kataklar va ularning joylashuvini muhokama qiladi. [Manba: Cambridge+ 9-sinf, PDF 30-32.]
\end{tcolorbox}
\vspace{2mm}\hrule height 0.2pt\vspace{3mm}

% DB_ID 6db2dc63-7c0d-984a-1904-45a1fcf9a107
\needspace{5\baselineskip}
{\bfseries\sffamily\color{ink} 234.} {scriptsizecolor{muted} Y1 | review | 6db2dc63}
\par
Ikkilamchi saqlash qurilmalari mavzusida: I. USB flesh - portativ solid-state saqlash qurilmasi. II. HDD - magnit plastinkali, katta sig'imli ikkilamchi xotira. Mulohazalarni baholang.
\begin{enumerate}[label=\Alph*),leftmargin=9mm]
\item I va II
\item faqat I
\item faqat II
\item I ham, II ham emas
\end{enumerate}
\begin{tcolorbox}[enhanced,breakable,colback=teal!4,colframe=teal!60!black,boxrule=0.45pt,arc=1.2mm,left=2.7mm,right=2.7mm,top=1.8mm,bottom=1.8mm,before skip=2mm,after skip=2mm]
{\bfseries\sffamily\color{teal}To'g'ri javob:} \textbf{A) I va II}
\par\smallskip
{\bfseries Izoh:} I va II. Birinchi va ikkinchi mulohazaning ta'riflari alohida tekshirildi; natija: I va II. Manba: MAR, PDF 7-9; Qo'shimcha sahifa raqami (qaysi PDFga tegishli ekani aniqlanmagan; darslik sahifasi deb tasdiqlanmagan): 35
\end{tcolorbox}
\vspace{2mm}\hrule height 0.2pt\vspace{3mm}

% DB_ID 722c3d12-c163-930b-9824-a44d126b96a0
\needspace{5\baselineskip}
{\bfseries\sffamily\color{ink} 235.} {scriptsizecolor{muted} Y1 | review | 722c3d12}
\par
Tizimli blok, bosh plata va ulanishlar mavzusida: I. quvvat bloki - tarmoq elektrini komponentlarga mos doimiy tokka aylantiruvchi qism. II. bosh plata - ichki komponentlarni bog'lovchi asosiy bosma plata. Mulohazalarni baholang.
\begin{enumerate}[label=\Alph*),leftmargin=9mm]
\item faqat I
\item faqat II
\item I va II
\item I ham, II ham emas
\end{enumerate}
\begin{tcolorbox}[enhanced,breakable,colback=teal!4,colframe=teal!60!black,boxrule=0.45pt,arc=1.2mm,left=2.7mm,right=2.7mm,top=1.8mm,bottom=1.8mm,before skip=2mm,after skip=2mm]
{\bfseries\sffamily\color{teal}To'g'ri javob:} \textbf{C) I va II}
\par\smallskip
{\bfseries Izoh:} I va II. Birinchi va ikkinchi mulohazaning ta'riflari alohida tekshirildi; natija: I va II. Manba: CAM9, PDF 9-14; Qo'shimcha sahifa raqami (qaysi PDFga tegishli ekani aniqlanmagan; darslik sahifasi deb tasdiqlanmagan): 15
\end{tcolorbox}
\vspace{2mm}\hrule height 0.2pt\vspace{3mm}

% DB_ID 79d06fcd-b3c5-6efc-67c2-837f1b7e0187
\needspace{5\baselineskip}
{\bfseries\sffamily\color{ink} 236.} {scriptsizecolor{muted} Y1 | review | 79d06fcd}
\par
Tizimli blok, bosh plata va ulanishlar bo'yicha qatorlar ichidan "port" noto'g'ri izohlanganini aniqlang.
\begin{enumerate}[label=\Alph*),leftmargin=9mm]
\item port - ichki komponentlarni bog'lovchi asosiy bosma plata
\item shina - ma'lumot va boshqaruv signallarini tashuvchi aloqa yo'llari
\item slot - ichki komponent o'rnatiladigan ulagich
\item quvvat bloki - tarmoq elektrini komponentlarga mos doimiy tokka aylantiruvchi qism
\end{enumerate}
\begin{tcolorbox}[enhanced,breakable,colback=teal!4,colframe=teal!60!black,boxrule=0.45pt,arc=1.2mm,left=2.7mm,right=2.7mm,top=1.8mm,bottom=1.8mm,before skip=2mm,after skip=2mm]
{\bfseries\sffamily\color{teal}To'g'ri javob:} \textbf{A) port - ichki komponentlarni bog'lovchi asosiy bosma plata}
\par\smallskip
{\bfseries Izoh:} port - ichki komponentlarni bog'lovchi asosiy bosma plata. port uchun to'g'ri ta'rif: tashqi qurilma ulanadigan nuqta. Manba: MAR, PDF 5-8; Qo'shimcha sahifa raqami (qaysi PDFga tegishli ekani aniqlanmagan; darslik sahifasi deb tasdiqlanmagan): 13
\end{tcolorbox}
\vspace{2mm}\hrule height 0.2pt\vspace{3mm}

% DB_ID 7b1b8f11-ef8c-4c75-1436-5b65686fbfb3
\needspace{5\baselineskip}
{\bfseries\sffamily\color{ink} 237.} {scriptsizecolor{muted} Y1 | review | 7b1b8f11}
\par
Markaziy protsessor va bajarish sikli mavzusida: I. takt chastotasi - sekunddagi takt impulslar soni. II. CPU - ko'rsatmalarni bajarib, ma'lumotlarni qayta ishlovchi markaziy qurilma. Mulohazalarni baholang.
\begin{enumerate}[label=\Alph*),leftmargin=9mm]
\item I va II
\item faqat I
\item faqat II
\item I ham, II ham emas
\end{enumerate}
\begin{tcolorbox}[enhanced,breakable,colback=teal!4,colframe=teal!60!black,boxrule=0.45pt,arc=1.2mm,left=2.7mm,right=2.7mm,top=1.8mm,bottom=1.8mm,before skip=2mm,after skip=2mm]
{\bfseries\sffamily\color{teal}To'g'ri javob:} \textbf{A) I va II}
\par\smallskip
{\bfseries Izoh:} I va II. Birinchi va ikkinchi mulohazaning ta'riflari alohida tekshirildi; natija: I va II. Manba: CAM9, PDF 10-14; Qo'shimcha sahifa raqami (qaysi PDFga tegishli ekani aniqlanmagan; darslik sahifasi deb tasdiqlanmagan): 22
\end{tcolorbox}
\vspace{2mm}\hrule height 0.2pt\vspace{3mm}

% DB_ID 7b22530a-d0ac-e1c1-37ad-ceb954d4b431
\needspace{5\baselineskip}
{\bfseries\sffamily\color{ink} 238.} {scriptsizecolor{muted} Y1 | review | 7b22530a}
\par
Aqlli issiqxona haroratni o'lchadi, qiymatni belgilangan chegara bilan taqqosladi va ventilyatorni yoqdi. To'g'ri IPOS ketma-ketligini toping.
\begin{enumerate}[label=\Alph*),leftmargin=9mm]
\item sensor \ensuremath{\rightarrow} protsessor \ensuremath{\rightarrow} aktuator
\item aktuator \ensuremath{\rightarrow} protsessor \ensuremath{\rightarrow} sensor
\item protsessor \ensuremath{\rightarrow} sensor \ensuremath{\rightarrow} saqlash
\item sensor \ensuremath{\rightarrow} monitor \ensuremath{\rightarrow} protsessor
\end{enumerate}
\begin{tcolorbox}[enhanced,breakable,colback=teal!4,colframe=teal!60!black,boxrule=0.45pt,arc=1.2mm,left=2.7mm,right=2.7mm,top=1.8mm,bottom=1.8mm,before skip=2mm,after skip=2mm]
{\bfseries\sffamily\color{teal}To'g'ri javob:} \textbf{A) sensor \ensuremath{\rightarrow} protsessor \ensuremath{\rightarrow} aktuator}
\par\smallskip
{\bfseries Izoh:} sensor \ensuremath{\rightarrow} protsessor \ensuremath{\rightarrow} aktuator. Sensor kiritadi, protsessor qaror qiladi, aktuator fizik chiqishni bajaradi. Manba: CAM9, PDF 8-15; Qo'shimcha sahifa raqami (qaysi PDFga tegishli ekani aniqlanmagan; darslik sahifasi deb tasdiqlanmagan): 9
\end{tcolorbox}
\vspace{2mm}\hrule height 0.2pt\vspace{3mm}

% DB_ID 7c0a84ae-12b8-eb74-0f73-c53d9b94ac6f
\needspace{5\baselineskip}
{\bfseries\sffamily\color{ink} 239.} {scriptsizecolor{muted} Y1 | review | 7c0a84ae}
\par
Talaba chizmani monitor orqali ko'rib chiqdi va uni qog'ozga printerda chop etdi. Qaysi tasnif to'g'ri?
\begin{enumerate}[label=\Alph*),leftmargin=9mm]
\item Qog'ozga bosilgan natija qattiq nusxa, monitor orqali ko'rsatilgani elektron nusxa
\item Monitor ekranidagi raqamli tasvir qattiq nusxa hisoblanadi, printer varag'i esa vaqtinchalik tezkor xotira nusxasidir
\item Printerdan chiqqan qog'oz ekran orqali ko'rsatiladigan yumshoq nusxa, monitor tasviri esa doimiy qog'oz nusxadir
\item Monitor axborotni uzoq muddat saqlovchi mexanik disk, printer esa faqat grafik ma'lumotlarni kirituvchi qurilma
\end{enumerate}
\begin{tcolorbox}[enhanced,breakable,colback=teal!4,colframe=teal!60!black,boxrule=0.45pt,arc=1.2mm,left=2.7mm,right=2.7mm,top=1.8mm,bottom=1.8mm,before skip=2mm,after skip=2mm]
{\bfseries\sffamily\color{teal}To'g'ri javob:} \textbf{A) Qog'ozga bosilgan natija qattiq nusxa, monitor orqali ko'rsatilgani elektron nusxa}
\par\smallskip
{\bfseries Izoh:} Printerning qog'ozga chiqaradigan fizik natijasi qattiq nusxa, monitor esa elektron tasvir beradi. [Manba: Cambridge+ 10-11-sinf PDF 36-37.]
\end{tcolorbox}
\vspace{2mm}\hrule height 0.2pt\vspace{3mm}

% DB_ID 7c2991e0-39e1-e156-5179-989c6ff85834
\needspace{5\baselineskip}
{\bfseries\sffamily\color{ink} 240.} {scriptsizecolor{muted} Y1 | review | 7c2991e0}
\par
Chiqarish va ikki vazifali qurilmalar mavzusida: I. yumshoq nusxa - ekranda yoki audio orqali vaqtinchalik taqdim etilgan natija. II. qattiq nusxa - qog'oz yoki boshqa fizik muhitdagi natija. Mulohazalarni baholang.
\begin{enumerate}[label=\Alph*),leftmargin=9mm]
\item I va II
\item faqat I
\item faqat II
\item I ham, II ham emas
\end{enumerate}
\begin{tcolorbox}[enhanced,breakable,colback=teal!4,colframe=teal!60!black,boxrule=0.45pt,arc=1.2mm,left=2.7mm,right=2.7mm,top=1.8mm,bottom=1.8mm,before skip=2mm,after skip=2mm]
{\bfseries\sffamily\color{teal}To'g'ri javob:} \textbf{A) I va II}
\par\smallskip
{\bfseries Izoh:} I va II. Birinchi va ikkinchi mulohazaning ta'riflari alohida tekshirildi; natija: I va II. Manba: CAM9, PDF 16-32; Qo'shimcha sahifa raqami (qaysi PDFga tegishli ekani aniqlanmagan; darslik sahifasi deb tasdiqlanmagan): 43
\end{tcolorbox}
\vspace{2mm}\hrule height 0.2pt\vspace{3mm}

% DB_ID 9193a62d-701c-789c-965a-28bf59514ab1
\needspace{5\baselineskip}
{\bfseries\sffamily\color{ink} 241.} {scriptsizecolor{muted} Y1 | review | 9193a62d}
\par
2.4 GHz takt chastotali protsessor 0.5 sekundda nechta takt impulsiga ega bo'ladi?
\begin{enumerate}[label=\Alph*),leftmargin=9mm]
\item 1.2 milliard
\item 4.8 milliard
\item 2.4 million
\item 1.2 million
\end{enumerate}
\begin{tcolorbox}[enhanced,breakable,colback=teal!4,colframe=teal!60!black,boxrule=0.45pt,arc=1.2mm,left=2.7mm,right=2.7mm,top=1.8mm,bottom=1.8mm,before skip=2mm,after skip=2mm]
{\bfseries\sffamily\color{teal}To'g'ri javob:} \textbf{A) 1.2 milliard}
\par\smallskip
{\bfseries Izoh:} 1.2 milliard. 2.4 milliard \ensuremath{\times} 0.5 = 1.2 milliard takt. Bu ko'rsatmalar soni bilan aynan teng bo'lishi shart emas. Manba: CAM9, PDF 10-14; Qo'shimcha sahifa raqami (qaysi PDFga tegishli ekani aniqlanmagan; darslik sahifasi deb tasdiqlanmagan): 22
\end{tcolorbox}
\vspace{2mm}\hrule height 0.2pt\vspace{3mm}

% DB_ID 92b41a5c-ff8c-f1b8-a000-9b2e20cd3fd1
\needspace{5\baselineskip}
{\bfseries\sffamily\color{ink} 242.} {scriptsizecolor{muted} Y1 | review | 92b41a5c}
\par
Kiritish qurilmalari mavzusida: I. OCR - tasvirdagi bosma matnni tahrirlanadigan belgilar sifatida tanish. II. OMR - oldindan belgilangan joylardagi belgilarni aniqlash. Mulohazalarni baholang.
\begin{enumerate}[label=\Alph*),leftmargin=9mm]
\item faqat I
\item faqat II
\item I va II
\item I ham, II ham emas
\end{enumerate}
\begin{tcolorbox}[enhanced,breakable,colback=teal!4,colframe=teal!60!black,boxrule=0.45pt,arc=1.2mm,left=2.7mm,right=2.7mm,top=1.8mm,bottom=1.8mm,before skip=2mm,after skip=2mm]
{\bfseries\sffamily\color{teal}To'g'ri javob:} \textbf{C) I va II}
\par\smallskip
{\bfseries Izoh:} I va II. Birinchi va ikkinchi mulohazaning ta'riflari alohida tekshirildi; natija: I va II. Manba: CAM9, PDF 16-32; Qo'shimcha sahifa raqami (qaysi PDFga tegishli ekani aniqlanmagan; darslik sahifasi deb tasdiqlanmagan): 37
\end{tcolorbox}
\vspace{2mm}\hrule height 0.2pt\vspace{3mm}

% DB_ID 94ce06bc-47d3-d366-d5ca-30e1a5750496
\needspace{5\baselineskip}
{\bfseries\sffamily\color{ink} 243.} {scriptsizecolor{muted} Y1 | review | 94ce06bc}
\par
Tizimli blok, bosh plata va ulanishlar mavzusida: I. bosh plata - ichki komponentlarni bog'lovchi asosiy bosma plata. II. shina - ma'lumot va boshqaruv signallarini tashuvchi aloqa yo'llari. Mulohazalarni baholang.
\begin{enumerate}[label=\Alph*),leftmargin=9mm]
\item faqat I
\item faqat II
\item I va II
\item I ham, II ham emas
\end{enumerate}
\begin{tcolorbox}[enhanced,breakable,colback=teal!4,colframe=teal!60!black,boxrule=0.45pt,arc=1.2mm,left=2.7mm,right=2.7mm,top=1.8mm,bottom=1.8mm,before skip=2mm,after skip=2mm]
{\bfseries\sffamily\color{teal}To'g'ri javob:} \textbf{C) I va II}
\par\smallskip
{\bfseries Izoh:} I va II. Birinchi va ikkinchi mulohazaning ta'riflari alohida tekshirildi; natija: I va II. Manba: CAM9, PDF 9-14; Qo'shimcha sahifa raqami (qaysi PDFga tegishli ekani aniqlanmagan; darslik sahifasi deb tasdiqlanmagan): 11
\end{tcolorbox}
\vspace{2mm}\hrule height 0.2pt\vspace{3mm}

% DB_ID 99e91a46-842c-691c-d823-816d0b77f2ef
\needspace{5\baselineskip}
{\bfseries\sffamily\color{ink} 244.} {scriptsizecolor{muted} Y1 | review | 99e91a46}
\par
Unumdorlik, nosozliklarni aniqlash va xavfsiz xizmat bo'yicha qatorlar ichidan "paging" noto'g'ri izohlanganini aniqlang.
\begin{enumerate}[label=\Alph*),leftmargin=9mm]
\item paging - tizim umumiy tezligini cheklayotgan eng sust resurs
\item thermal throttling - qizish sabab protsessor tezligining avtomatik pasayishi
\item diagnostika - belgi va dalillar asosida sababni tizimli aniqlash
\item log - tizim hodisalari va xatolari qaydi
\end{enumerate}
\begin{tcolorbox}[enhanced,breakable,colback=teal!4,colframe=teal!60!black,boxrule=0.45pt,arc=1.2mm,left=2.7mm,right=2.7mm,top=1.8mm,bottom=1.8mm,before skip=2mm,after skip=2mm]
{\bfseries\sffamily\color{teal}To'g'ri javob:} \textbf{A) paging - tizim umumiy tezligini cheklayotgan eng sust resurs}
\par\smallskip
{\bfseries Izoh:} paging - tizim umumiy tezligini cheklayotgan eng sust resurs. paging uchun to'g'ri ta'rif: RAM va ikkilamchi xotira o'rtasida sahifalarni almashish. Manba: CAM10/11, PDF 24-61; Qo'shimcha sahifa raqami (qaysi PDFga tegishli ekani aniqlanmagan; darslik sahifasi deb tasdiqlanmagan): 109
\end{tcolorbox}
\vspace{2mm}\hrule height 0.2pt\vspace{3mm}

% DB_ID a03119d9-de17-0f08-3405-89a283e8f0a3
\needspace{5\baselineskip}
{\bfseries\sffamily\color{ink} 245.} {scriptsizecolor{muted} Y1 | review | a03119d9}
\par
Attestatsiyadagi nozik farqlar va integratsion masalalar bo'yicha qatorlar ichidan "integratsion savol" noto'g'ri izohlanganini aniqlang.
\begin{enumerate}[label=\Alph*),leftmargin=9mm]
\item integratsion savol - obyektni muhim belgilari bo'yicha guruhga ajratish
\item mezon - variantlarni baholash uchun qo'llanadigan belgi
\item cheklov - yechim bajarishi shart bo'lgan chegara yoki talab
\item distraktor - ishonarli, ammo noto'g'ri javob varianti
\end{enumerate}
\begin{tcolorbox}[enhanced,breakable,colback=teal!4,colframe=teal!60!black,boxrule=0.45pt,arc=1.2mm,left=2.7mm,right=2.7mm,top=1.8mm,bottom=1.8mm,before skip=2mm,after skip=2mm]
{\bfseries\sffamily\color{teal}To'g'ri javob:} \textbf{A) integratsion savol - obyektni muhim belgilari bo'yicha guruhga ajratish}
\par\smallskip
{\bfseries Izoh:} integratsion savol - obyektni muhim belgilari bo'yicha guruhga ajratish. integratsion savol uchun to'g'ri ta'rif: bir nechta mavzuni bitta vaziyatda bog'laydigan topshiriq. Manba: ICT5, PDF 25-49; Qo'shimcha sahifa raqami (qaysi PDFga tegishli ekani aniqlanmagan; darslik sahifasi deb tasdiqlanmagan): 119
\end{tcolorbox}
\vspace{2mm}\hrule height 0.2pt\vspace{3mm}

% DB_ID a2cc4025-633d-c5fd-44d2-0ca88415a853
\needspace{5\baselineskip}
{\bfseries\sffamily\color{ink} 246.} {scriptsizecolor{muted} Y1 | review | a2cc4025}
\par
Unumdorlik, nosozliklarni aniqlash va xavfsiz xizmat mavzusida: I. log - tizim hodisalari va xatolari qaydi. II. bottleneck - tizim umumiy tezligini cheklayotgan eng sust resurs. Mulohazalarni baholang.
\begin{enumerate}[label=\Alph*),leftmargin=9mm]
\item I va II
\item faqat I
\item faqat II
\item I ham, II ham emas
\end{enumerate}
\begin{tcolorbox}[enhanced,breakable,colback=teal!4,colframe=teal!60!black,boxrule=0.45pt,arc=1.2mm,left=2.7mm,right=2.7mm,top=1.8mm,bottom=1.8mm,before skip=2mm,after skip=2mm]
{\bfseries\sffamily\color{teal}To'g'ri javob:} \textbf{A) I va II}
\par\smallskip
{\bfseries Izoh:} I va II. Birinchi va ikkinchi mulohazaning ta'riflari alohida tekshirildi; natija: I va II. Manba: CAM9, PDF 12-47; Qo'shimcha sahifa raqami (qaysi PDFga tegishli ekani aniqlanmagan; darslik sahifasi deb tasdiqlanmagan): 113
\end{tcolorbox}
\vspace{2mm}\hrule height 0.2pt\vspace{3mm}

% DB_ID aea7d1d5-16be-2123-8b35-7ac5e15cb0b1
\needspace{5\baselineskip}
{\bfseries\sffamily\color{ink} 247.} {scriptsizecolor{muted} Y1 | review | aea7d1d5}
\par
Attestatsiyadagi nozik farqlar va integratsion masalalar mavzusida: I. cheklov - yechim bajarishi shart bo'lgan chegara yoki talab. II. tasnif - yechim bajarishi shart bo'lgan chegara yoki talab. Mulohazalarni baholang.
\begin{enumerate}[label=\Alph*),leftmargin=9mm]
\item faqat I
\item faqat II
\item I va II
\item I ham, II ham emas
\end{enumerate}
\begin{tcolorbox}[enhanced,breakable,colback=teal!4,colframe=teal!60!black,boxrule=0.45pt,arc=1.2mm,left=2.7mm,right=2.7mm,top=1.8mm,bottom=1.8mm,before skip=2mm,after skip=2mm]
{\bfseries\sffamily\color{teal}To'g'ri javob:} \textbf{A) faqat I}
\par\smallskip
{\bfseries Izoh:} faqat I. Birinchi va ikkinchi mulohazaning ta'riflari alohida tekshirildi; natija: faqat I. Manba: CAM10/11, PDF 23-61; Qo'shimcha sahifa raqami (qaysi PDFga tegishli ekani aniqlanmagan; darslik sahifasi deb tasdiqlanmagan): 118
\end{tcolorbox}
\vspace{2mm}\hrule height 0.2pt\vspace{3mm}

% DB_ID b45539fc-84e2-f9ee-b8c3-9f2a004a4f81
\needspace{5\baselineskip}
{\bfseries\sffamily\color{ink} 248.} {scriptsizecolor{muted} Y1 | review | b45539fc}
\par
Kompyuter tizimi va uning ishlash modeli mavzusida: I. o'rnatilgan tizim - kattaroq qurilma ichida muayyan vazifani bajaruvchi kompyuter tizimi. II. kompyuter tizimi - bir maqsad uchun hamkorlik qiluvchi texnik vosita, dastur, ma'lumot va foydalanuvchi majmuasi. Mulohazalarni baholang.
\begin{enumerate}[label=\Alph*),leftmargin=9mm]
\item I va II
\item faqat I
\item faqat II
\item I ham, II ham emas
\end{enumerate}
\begin{tcolorbox}[enhanced,breakable,colback=teal!4,colframe=teal!60!black,boxrule=0.45pt,arc=1.2mm,left=2.7mm,right=2.7mm,top=1.8mm,bottom=1.8mm,before skip=2mm,after skip=2mm]
{\bfseries\sffamily\color{teal}To'g'ri javob:} \textbf{A) I va II}
\par\smallskip
{\bfseries Izoh:} I va II. Birinchi va ikkinchi mulohazaning ta'riflari alohida tekshirildi; natija: I va II. Manba: ICT9, PDF 14; Qo'shimcha sahifa raqami (qaysi PDFga tegishli ekani aniqlanmagan; darslik sahifasi deb tasdiqlanmagan): 8
\end{tcolorbox}
\vspace{2mm}\hrule height 0.2pt\vspace{3mm}

% DB_ID b588bd3d-1a55-35a7-0fbe-f48f6e86746f
\needspace{5\baselineskip}
{\bfseries\sffamily\color{ink} 249.} {scriptsizecolor{muted} Y1 | review | b588bd3d}
\par
Chiqarish va ikki vazifali qurilmalar mavzusida: I. lazer printer - toner va elektrofotografik jarayon bilan tez chop etuvchi printer. II. yumshoq nusxa - toner va elektrofotografik jarayon bilan tez chop etuvchi printer. Mulohazalarni baholang.
\begin{enumerate}[label=\Alph*),leftmargin=9mm]
\item faqat II
\item I va II
\item faqat I
\item I ham, II ham emas
\end{enumerate}
\begin{tcolorbox}[enhanced,breakable,colback=teal!4,colframe=teal!60!black,boxrule=0.45pt,arc=1.2mm,left=2.7mm,right=2.7mm,top=1.8mm,bottom=1.8mm,before skip=2mm,after skip=2mm]
{\bfseries\sffamily\color{teal}To'g'ri javob:} \textbf{C) faqat I}
\par\smallskip
{\bfseries Izoh:} faqat I. Birinchi va ikkinchi mulohazaning ta'riflari alohida tekshirildi; natija: faqat I. Manba: ICT5, PDF 29-36; Qo'shimcha sahifa raqami (qaysi PDFga tegishli ekani aniqlanmagan; darslik sahifasi deb tasdiqlanmagan): 45
\end{tcolorbox}
\vspace{2mm}\hrule height 0.2pt\vspace{3mm}

% DB_ID b701f6c6-a3f6-b9d3-2b75-225a30a44420
\needspace{5\baselineskip}
{\bfseries\sffamily\color{ink} 250.} {scriptsizecolor{muted} Y1 | review | b701f6c6}
\par
Attestatsiyadagi nozik farqlar va integratsion masalalar bo'yicha qatorlar ichidan "mezon" noto'g'ri izohlanganini aniqlang.
\begin{enumerate}[label=\Alph*),leftmargin=9mm]
\item cheklov - yechim bajarishi shart bo'lgan chegara yoki talab
\item integratsion savol - bir nechta mavzuni bitta vaziyatda bog'laydigan topshiriq
\item mezon - obyektni muhim belgilari bo'yicha guruhga ajratish
\item distraktor - ishonarli, ammo noto'g'ri javob varianti
\end{enumerate}
\begin{tcolorbox}[enhanced,breakable,colback=teal!4,colframe=teal!60!black,boxrule=0.45pt,arc=1.2mm,left=2.7mm,right=2.7mm,top=1.8mm,bottom=1.8mm,before skip=2mm,after skip=2mm]
{\bfseries\sffamily\color{teal}To'g'ri javob:} \textbf{C) mezon - obyektni muhim belgilari bo'yicha guruhga ajratish}
\par\smallskip
{\bfseries Izoh:} mezon - obyektni muhim belgilari bo'yicha guruhga ajratish. mezon uchun to'g'ri ta'rif: variantlarni baholash uchun qo'llanadigan belgi. Manba: CAM9, PDF 8-47; Qo'shimcha sahifa raqami (qaysi PDFga tegishli ekani aniqlanmagan; darslik sahifasi deb tasdiqlanmagan): 116
\end{tcolorbox}
\vspace{2mm}\hrule height 0.2pt\vspace{3mm}

% DB_ID c845f735-952a-c38a-9263-9760803eb673
\needspace{5\baselineskip}
{\bfseries\sffamily\color{ink} 251.} {scriptsizecolor{muted} Y1 | review | c845f735}
\par
Unumdorlik, nosozliklarni aniqlash va xavfsiz xizmat mavzusida: I. thermal throttling - qizish sabab protsessor tezligining avtomatik pasayishi. II. bottleneck - qizish sabab protsessor tezligining avtomatik pasayishi. Mulohazalarni baholang.
\begin{enumerate}[label=\Alph*),leftmargin=9mm]
\item faqat II
\item I va II
\item faqat I
\item I ham, II ham emas
\end{enumerate}
\begin{tcolorbox}[enhanced,breakable,colback=teal!4,colframe=teal!60!black,boxrule=0.45pt,arc=1.2mm,left=2.7mm,right=2.7mm,top=1.8mm,bottom=1.8mm,before skip=2mm,after skip=2mm]
{\bfseries\sffamily\color{teal}To'g'ri javob:} \textbf{C) faqat I}
\par\smallskip
{\bfseries Izoh:} faqat I. Birinchi va ikkinchi mulohazaning ta'riflari alohida tekshirildi; natija: faqat I. Manba: ICT10, PDF 12-29; Qo'shimcha sahifa raqami (qaysi PDFga tegishli ekani aniqlanmagan; darslik sahifasi deb tasdiqlanmagan): 111
\end{tcolorbox}
\vspace{2mm}\hrule height 0.2pt\vspace{3mm}

% DB_ID ca09a37e-409e-cee8-eca6-7a00c63196e9
\needspace{5\baselineskip}
{\bfseries\sffamily\color{ink} 252.} {scriptsizecolor{muted} Y1 | review | ca09a37e}
\par
Kompyuter tizimi va uning ishlash modeli mavzusida: I. kompyuter tizimi - bir maqsad uchun hamkorlik qiluvchi texnik vosita, dastur, ma'lumot va foydalanuvchi majmuasi. II. texnik ta'minot - kompyuter tizimining fizik qismlari. Mulohazalarni baholang.
\begin{enumerate}[label=\Alph*),leftmargin=9mm]
\item I va II
\item faqat I
\item faqat II
\item I ham, II ham emas
\end{enumerate}
\begin{tcolorbox}[enhanced,breakable,colback=teal!4,colframe=teal!60!black,boxrule=0.45pt,arc=1.2mm,left=2.7mm,right=2.7mm,top=1.8mm,bottom=1.8mm,before skip=2mm,after skip=2mm]
{\bfseries\sffamily\color{teal}To'g'ri javob:} \textbf{A) I va II}
\par\smallskip
{\bfseries Izoh:} I va II. Birinchi va ikkinchi mulohazaning ta'riflari alohida tekshirildi; natija: I va II. Manba: CAM9, PDF 8-15; Qo'shimcha sahifa raqami (qaysi PDFga tegishli ekani aniqlanmagan; darslik sahifasi deb tasdiqlanmagan): 4
\end{tcolorbox}
\vspace{2mm}\hrule height 0.2pt\vspace{3mm}

% DB_ID caf4b6da-1d02-e8b1-f275-10f341f2327b
\needspace{5\baselineskip}
{\bfseries\sffamily\color{ink} 253.} {scriptsizecolor{muted} Y1 | review | caf4b6da}
\par
Ikkilamchi saqlash qurilmalari mavzusida: I. optik disk - lazer orqali o'qiladigan va yoziladigan disk. II. HDD - lazer orqali o'qiladigan va yoziladigan disk. Mulohazalarni baholang.
\begin{enumerate}[label=\Alph*),leftmargin=9mm]
\item faqat II
\item I va II
\item faqat I
\item I ham, II ham emas
\end{enumerate}
\begin{tcolorbox}[enhanced,breakable,colback=teal!4,colframe=teal!60!black,boxrule=0.45pt,arc=1.2mm,left=2.7mm,right=2.7mm,top=1.8mm,bottom=1.8mm,before skip=2mm,after skip=2mm]
{\bfseries\sffamily\color{teal}To'g'ri javob:} \textbf{C) faqat I}
\par\smallskip
{\bfseries Izoh:} faqat I. Birinchi va ikkinchi mulohazaning ta'riflari alohida tekshirildi; natija: faqat I. Manba: ICT10, PDF 14-18; Qo'shimcha sahifa raqami (qaysi PDFga tegishli ekani aniqlanmagan; darslik sahifasi deb tasdiqlanmagan): 33
\end{tcolorbox}
\vspace{2mm}\hrule height 0.2pt\vspace{3mm}

% DB_ID ce312cd8-b7ae-33a2-7e29-d9c1d1b71608
\needspace{5\baselineskip}
{\bfseries\sffamily\color{ink} 254.} {scriptsizecolor{muted} Y1 | review | ce312cd8}
\par
Tizimli blok, bosh plata va ulanishlar mavzusida: I. slot - ichki komponent o'rnatiladigan ulagich. II. bosh plata - ichki komponent o'rnatiladigan ulagich. Mulohazalarni baholang.
\begin{enumerate}[label=\Alph*),leftmargin=9mm]
\item faqat I
\item faqat II
\item I va II
\item I ham, II ham emas
\end{enumerate}
\begin{tcolorbox}[enhanced,breakable,colback=teal!4,colframe=teal!60!black,boxrule=0.45pt,arc=1.2mm,left=2.7mm,right=2.7mm,top=1.8mm,bottom=1.8mm,before skip=2mm,after skip=2mm]
{\bfseries\sffamily\color{teal}To'g'ri javob:} \textbf{A) faqat I}
\par\smallskip
{\bfseries Izoh:} faqat I. Birinchi va ikkinchi mulohazaning ta'riflari alohida tekshirildi; natija: faqat I. Manba: ICT10, PDF 12-18; Qo'shimcha sahifa raqami (qaysi PDFga tegishli ekani aniqlanmagan; darslik sahifasi deb tasdiqlanmagan): 12
\end{tcolorbox}
\vspace{2mm}\hrule height 0.2pt\vspace{3mm}

% DB_ID d4491e55-d451-9701-e541-47fc66b82819
\needspace{5\baselineskip}
{\bfseries\sffamily\color{ink} 255.} {scriptsizecolor{muted} Y1 | review | d4491e55}
\par
Unumdorlik, nosozliklarni aniqlash va xavfsiz xizmat bo'yicha qatorlar ichidan "diagnostika" noto'g'ri izohlanganini aniqlang.
\begin{enumerate}[label=\Alph*),leftmargin=9mm]
\item paging - RAM va ikkilamchi xotira o'rtasida sahifalarni almashish
\item thermal throttling - qizish sabab protsessor tezligining avtomatik pasayishi
\item diagnostika - tizim umumiy tezligini cheklayotgan eng sust resurs
\item log - tizim hodisalari va xatolari qaydi
\end{enumerate}
\begin{tcolorbox}[enhanced,breakable,colback=teal!4,colframe=teal!60!black,boxrule=0.45pt,arc=1.2mm,left=2.7mm,right=2.7mm,top=1.8mm,bottom=1.8mm,before skip=2mm,after skip=2mm]
{\bfseries\sffamily\color{teal}To'g'ri javob:} \textbf{C) diagnostika - tizim umumiy tezligini cheklayotgan eng sust resurs}
\par\smallskip
{\bfseries Izoh:} diagnostika - tizim umumiy tezligini cheklayotgan eng sust resurs. diagnostika uchun to'g'ri ta'rif: belgi va dalillar asosida sababni tizimli aniqlash. Manba: MAR, PDF 10-20; Qo'shimcha sahifa raqami (qaysi PDFga tegishli ekani aniqlanmagan; darslik sahifasi deb tasdiqlanmagan): 112
\end{tcolorbox}
\vspace{2mm}\hrule height 0.2pt\vspace{3mm}

% DB_ID d9bd79ba-ec0b-7b02-30e4-734e8bc7474c
\needspace{5\baselineskip}
{\bfseries\sffamily\color{ink} 256.} {scriptsizecolor{muted} Y1 | review | d9bd79ba}
\par
Markaziy protsessor va bajarish sikli mavzusida: I. CPU - ko'rsatmalarni bajarib, ma'lumotlarni qayta ishlovchi markaziy qurilma. II. ALU - arifmetik va mantiqiy amallarni bajaruvchi qism. Mulohazalarni baholang.
\begin{enumerate}[label=\Alph*),leftmargin=9mm]
\item I va II
\item faqat I
\item faqat II
\item I ham, II ham emas
\end{enumerate}
\begin{tcolorbox}[enhanced,breakable,colback=teal!4,colframe=teal!60!black,boxrule=0.45pt,arc=1.2mm,left=2.7mm,right=2.7mm,top=1.8mm,bottom=1.8mm,before skip=2mm,after skip=2mm]
{\bfseries\sffamily\color{teal}To'g'ri javob:} \textbf{A) I va II}
\par\smallskip
{\bfseries Izoh:} I va II. Birinchi va ikkinchi mulohazaning ta'riflari alohida tekshirildi; natija: I va II. Manba: CAM9, PDF 10-14; Qo'shimcha sahifa raqami (qaysi PDFga tegishli ekani aniqlanmagan; darslik sahifasi deb tasdiqlanmagan): 18
\end{tcolorbox}
\vspace{2mm}\hrule height 0.2pt\vspace{3mm}

% DB_ID da1f2605-f3f8-bbdf-81be-301a505a894a
\needspace{5\baselineskip}
{\bfseries\sffamily\color{ink} 257.} {scriptsizecolor{muted} Y1 | review | da1f2605}
\par
Qaysi qatorda ikki tushuncha bir-biriga bog'liq, ammo aynan bir xil emas?
\begin{enumerate}[label=\Alph*),leftmargin=9mm]
\item RAM va tezkor xotira
\item CPU va markaziy protsessor
\item skaner va OCR
\item papka va katalog
\end{enumerate}
\begin{tcolorbox}[enhanced,breakable,colback=teal!4,colframe=teal!60!black,boxrule=0.45pt,arc=1.2mm,left=2.7mm,right=2.7mm,top=1.8mm,bottom=1.8mm,before skip=2mm,after skip=2mm]
{\bfseries\sffamily\color{teal}To'g'ri javob:} \textbf{C) skaner va OCR}
\par\smallskip
{\bfseries Izoh:} skaner va OCR. Skaner tasvirni kiritadi, OCR esa tasvirdagi belgilarni matn sifatida taniydi. Manba: SPEC, PDF 3; Qo'shimcha sahifa raqami (qaysi PDFga tegishli ekani aniqlanmagan; darslik sahifasi deb tasdiqlanmagan): 121
\end{tcolorbox}
\vspace{2mm}\hrule height 0.2pt\vspace{3mm}

% DB_ID dccb1b07-6b38-dc77-b38b-bae3931afb9b
\needspace{5\baselineskip}
{\bfseries\sffamily\color{ink} 258.} {scriptsizecolor{muted} Y1 | review | dccb1b07}
\par
Do'kon kassasida skaner mahsulot kodini o'qidi, tizim chegirmani hisobladi, ekranda yakuniy narx chiqdi va savdo tarixi diskka yozildi. IPOS ketma-ketligi qaysi?
\begin{enumerate}[label=\Alph*),leftmargin=9mm]
\item Ekranga chiqarish \ensuremath{\rightarrow} hisoblash \ensuremath{\rightarrow} kodni o'qish \ensuremath{\rightarrow} diskka yozish
\item Kodni o'qish \ensuremath{\rightarrow} ekranga chiqarish \ensuremath{\rightarrow} diskka yozish \ensuremath{\rightarrow} hisoblash
\item Kodni o'qish \ensuremath{\rightarrow} hisoblash \ensuremath{\rightarrow} ekranga chiqarish \ensuremath{\rightarrow} diskka yozish
\item Diskka yozish \ensuremath{\rightarrow} kodni o'qish \ensuremath{\rightarrow} hisoblash \ensuremath{\rightarrow} ekranga chiqarish
\end{enumerate}
\begin{tcolorbox}[enhanced,breakable,colback=teal!4,colframe=teal!60!black,boxrule=0.45pt,arc=1.2mm,left=2.7mm,right=2.7mm,top=1.8mm,bottom=1.8mm,before skip=2mm,after skip=2mm]
{\bfseries\sffamily\color{teal}To'g'ri javob:} \textbf{C) Kodni o'qish \ensuremath{\rightarrow} hisoblash \ensuremath{\rightarrow} ekranga chiqarish \ensuremath{\rightarrow} diskka yozish}
\par\smallskip
{\bfseries Izoh:} Kodni skaner yordamida kiritish, chegirmani hisoblash, narxni ekranga chiqarish va tranzaksiyani doimiy saqlash mos ravishda input, process, output, storage bosqichlaridir. [Manba: Cambridge+ 9-sinf, PDF 9-10.]
\end{tcolorbox}
\vspace{2mm}\hrule height 0.2pt\vspace{3mm}

% DB_ID e1b6d698-8680-9a16-7f88-9ab811b6a75f
\needspace{5\baselineskip}
{\bfseries\sffamily\color{ink} 259.} {scriptsizecolor{muted} Y1 | review | e1b6d698}
\par
Markaziy protsessor va bajarish sikli mavzusida: I. CU - ko'rsatmalar oqimi va boshqa qismlar ishini boshqaruvchi qism. II. CPU - ko'rsatmalar oqimi va boshqa qismlar ishini boshqaruvchi qism. Mulohazalarni baholang.
\begin{enumerate}[label=\Alph*),leftmargin=9mm]
\item faqat II
\item I va II
\item faqat I
\item I ham, II ham emas
\end{enumerate}
\begin{tcolorbox}[enhanced,breakable,colback=teal!4,colframe=teal!60!black,boxrule=0.45pt,arc=1.2mm,left=2.7mm,right=2.7mm,top=1.8mm,bottom=1.8mm,before skip=2mm,after skip=2mm]
{\bfseries\sffamily\color{teal}To'g'ri javob:} \textbf{C) faqat I}
\par\smallskip
{\bfseries Izoh:} faqat I. Birinchi va ikkinchi mulohazaning ta'riflari alohida tekshirildi; natija: faqat I. Manba: ICT10, PDF 17-20; Qo'shimcha sahifa raqami (qaysi PDFga tegishli ekani aniqlanmagan; darslik sahifasi deb tasdiqlanmagan): 20
\end{tcolorbox}
\vspace{2mm}\hrule height 0.2pt\vspace{3mm}

% DB_ID e95d4cd4-a9ca-2d96-6e25-5b4b9624ca57
\needspace{5\baselineskip}
{\bfseries\sffamily\color{ink} 260.} {scriptsizecolor{muted} Y1 | review | e95d4cd4}
\par
Har kuni 800 GB zaxira olinadi. 4 TB lenta sig'imida 1 TB = 1000 GB deb olinsa, nazariy jihatdan nechta to'liq kun saqlanadi?
\begin{enumerate}[label=\Alph*),leftmargin=9mm]
\item 4 kun
\item 5 kun
\item 6 kun
\item 8 kun
\end{enumerate}
\begin{tcolorbox}[enhanced,breakable,colback=teal!4,colframe=teal!60!black,boxrule=0.45pt,arc=1.2mm,left=2.7mm,right=2.7mm,top=1.8mm,bottom=1.8mm,before skip=2mm,after skip=2mm]
{\bfseries\sffamily\color{teal}To'g'ri javob:} \textbf{B) 5 kun}
\par\smallskip
{\bfseries Izoh:} 5 kun. 4000 \ensuremath{\div} 800 = 5 ta to'liq zaxira. Manba: CAM9, PDF 33-47; Qo'shimcha sahifa raqami (qaysi PDFga tegishli ekani aniqlanmagan; darslik sahifasi deb tasdiqlanmagan): 35
\end{tcolorbox}
\vspace{2mm}\hrule height 0.2pt\vspace{3mm}

% DB_ID eeb9ce03-e639-297a-8e8e-feaa97dcf238
\needspace{5\baselineskip}
{\bfseries\sffamily\color{ink} 261.} {scriptsizecolor{muted} Y1 | review | eeb9ce03}
\par
Birlamchi xotira: registr, kesh, RAM va ROM mavzusida: I. RAM - faol dastur va ma'lumotlar uchun o'zgaruvchan tezkor xotira. II. ROM - elektrsiz ham ma'lumotni saqlovchi o'zgarmas xotira. Mulohazalarni baholang.
\begin{enumerate}[label=\Alph*),leftmargin=9mm]
\item faqat I
\item faqat II
\item I va II
\item I ham, II ham emas
\end{enumerate}
\begin{tcolorbox}[enhanced,breakable,colback=teal!4,colframe=teal!60!black,boxrule=0.45pt,arc=1.2mm,left=2.7mm,right=2.7mm,top=1.8mm,bottom=1.8mm,before skip=2mm,after skip=2mm]
{\bfseries\sffamily\color{teal}To'g'ri javob:} \textbf{C) I va II}
\par\smallskip
{\bfseries Izoh:} I va II. Birinchi va ikkinchi mulohazaning ta'riflari alohida tekshirildi; natija: I va II. Manba: CAM9, PDF 11-14; Qo'shimcha sahifa raqami (qaysi PDFga tegishli ekani aniqlanmagan; darslik sahifasi deb tasdiqlanmagan): 24
\end{tcolorbox}
\vspace{2mm}\hrule height 0.2pt\vspace{3mm}

% DB_ID f0e6ba2c-41f7-5be8-f520-c5be9c53f799
\needspace{5\baselineskip}
{\bfseries\sffamily\color{ink} 262.} {scriptsizecolor{muted} Y1 | review | f0e6ba2c}
\par
Chiqarish va ikki vazifali qurilmalar bo'yicha qatorlar ichidan "siyohpurkagich" noto'g'ri izohlanganini aniqlang.
\begin{enumerate}[label=\Alph*),leftmargin=9mm]
\item qattiq nusxa - qog'oz yoki boshqa fizik muhitdagi natija
\item lazer printer - toner va elektrofotografik jarayon bilan tez chop etuvchi printer
\item siyohpurkagich - ekranda yoki audio orqali vaqtinchalik taqdim etilgan natija
\item aktuator - boshqaruv signalini fizik harakat yoki ta'sirga aylantiruvchi qurilma
\end{enumerate}
\begin{tcolorbox}[enhanced,breakable,colback=teal!4,colframe=teal!60!black,boxrule=0.45pt,arc=1.2mm,left=2.7mm,right=2.7mm,top=1.8mm,bottom=1.8mm,before skip=2mm,after skip=2mm]
{\bfseries\sffamily\color{teal}To'g'ri javob:} \textbf{C) siyohpurkagich - ekranda yoki audio orqali vaqtinchalik taqdim etilgan natija}
\par\smallskip
{\bfseries Izoh:} siyohpurkagich - ekranda yoki audio orqali vaqtinchalik taqdim etilgan natija. siyohpurkagich uchun to'g'ri ta'rif: mayda siyoh tomchilarini purkab tasvir hosil qiluvchi printer. Manba: ICT10, PDF 9-13; Qo'shimcha sahifa raqami (qaysi PDFga tegishli ekani aniqlanmagan; darslik sahifasi deb tasdiqlanmagan): 46
\end{tcolorbox}
\vspace{2mm}\hrule height 0.2pt\vspace{3mm}

% DB_ID f279e292-f23d-fb2f-f588-fdb11e54acda
\needspace{5\baselineskip}
{\bfseries\sffamily\color{ink} 263.} {scriptsizecolor{muted} Y1 | review | f279e292}
\par
Kiritish qurilmalari mavzusida: I. MICR - magnit siyoh bilan bosilgan maxsus belgilarni o'qish. II. OCR - magnit siyoh bilan bosilgan maxsus belgilarni o'qish. Mulohazalarni baholang.
\begin{enumerate}[label=\Alph*),leftmargin=9mm]
\item faqat I
\item faqat II
\item I va II
\item I ham, II ham emas
\end{enumerate}
\begin{tcolorbox}[enhanced,breakable,colback=teal!4,colframe=teal!60!black,boxrule=0.45pt,arc=1.2mm,left=2.7mm,right=2.7mm,top=1.8mm,bottom=1.8mm,before skip=2mm,after skip=2mm]
{\bfseries\sffamily\color{teal}To'g'ri javob:} \textbf{A) faqat I}
\par\smallskip
{\bfseries Izoh:} faqat I. Birinchi va ikkinchi mulohazaning ta'riflari alohida tekshirildi; natija: faqat I. Manba: CAM5, PDF 60-70; Qo'shimcha sahifa raqami (qaysi PDFga tegishli ekani aniqlanmagan; darslik sahifasi deb tasdiqlanmagan): 39
\end{tcolorbox}
\vspace{2mm}\hrule height 0.2pt\vspace{3mm}

% DB_ID f36ede51-1bb2-2036-6899-5a0ac2b39a41
\needspace{5\baselineskip}
{\bfseries\sffamily\color{ink} 264.} {scriptsizecolor{muted} Y1 | review | f36ede51}
\par
Attestatsiyadagi nozik farqlar va integratsion masalalar mavzusida: I. tasnif - obyektni muhim belgilari bo'yicha guruhga ajratish. II. mezon - variantlarni baholash uchun qo'llanadigan belgi. Mulohazalarni baholang.
\begin{enumerate}[label=\Alph*),leftmargin=9mm]
\item faqat I
\item faqat II
\item I va II
\item I ham, II ham emas
\end{enumerate}
\begin{tcolorbox}[enhanced,breakable,colback=teal!4,colframe=teal!60!black,boxrule=0.45pt,arc=1.2mm,left=2.7mm,right=2.7mm,top=1.8mm,bottom=1.8mm,before skip=2mm,after skip=2mm]
{\bfseries\sffamily\color{teal}To'g'ri javob:} \textbf{C) I va II}
\par\smallskip
{\bfseries Izoh:} I va II. Birinchi va ikkinchi mulohazaning ta'riflari alohida tekshirildi; natija: I va II. Manba: SPEC, PDF 3; Qo'shimcha sahifa raqami (qaysi PDFga tegishli ekani aniqlanmagan; darslik sahifasi deb tasdiqlanmagan): 116
\end{tcolorbox}
\vspace{2mm}\hrule height 0.2pt\vspace{3mm}

% DB_ID 0b7132bf-6692-799c-085a-473c765c621b
\needspace{5\baselineskip}
{\bfseries\sffamily\color{ink} 265.} {scriptsizecolor{muted} Y1 | review | 0b7132bf}
\par
Maktabda 900 ta to'rtta javob katakchali test blankini avtomatik baholash rejalashtirilgan. Har bir katakchadagi belgi mavjudligini qanday aniqlash maqsadga muvofiq?
\begin{enumerate}[label=\Alph*),leftmargin=9mm]
\item OCR orqali javob varaqasidagi bosma harflarni tahrirlanadigan matnga aylantirish
\item OMR orqali oldindan belgilangan javob katakchalaridagi belgilarni aniqlash
\item MICR orqali blankda magnit siyoh bilan yozilgan bank identifikatorini aniqlash
\item Shtrix-kod o'quvchi orqali blankning mahsulot identifikatori raqamini kiritish
\end{enumerate}
\begin{tcolorbox}[enhanced,breakable,colback=teal!4,colframe=teal!60!black,boxrule=0.45pt,arc=1.2mm,left=2.7mm,right=2.7mm,top=1.8mm,bottom=1.8mm,before skip=2mm,after skip=2mm]
{\bfseries\sffamily\color{teal}To'g'ri javob:} \textbf{B) OMR orqali oldindan belgilangan javob katakchalaridagi belgilarni aniqlash}
\par\smallskip
{\bfseries Izoh:} OMR standart shakldagi to'ldirilgan belgi katakchalarini aniqlash uchun mos. OCR esa rasm ichidagi belgilarni matnga o'tkazadi. Manba: Cambridge+ 9-sinf, PDF 30-31
\end{tcolorbox}
\vspace{2mm}\hrule height 0.2pt\vspace{3mm}

% DB_ID 0f74ad4b-1b0b-ba7e-5432-4bea7f9a04df
\needspace{5\baselineskip}
{\bfseries\sffamily\color{ink} 266.} {scriptsizecolor{muted} Y1 | review | 0f74ad4b}
\par
Noutbukda RAM bandligi 98\%, disk faoliyati juda yuqori, CPU esa 25\% band. Tashxisni qaysi tekshiruvdan boshlash eng asosli?
\begin{enumerate}[label=\Alph*),leftmargin=9mm]
\item Protsessor 25\% bo'lsa ham uning takt chastotasini asosiy sabab deb qabul qilib, boshqa resurslarni tekshirmaslik
\item Xotira sarflayotgan jarayonlarni va RAM yetishmagani uchun diskka sahifalash mavjudligini tekshirish
\item Disk yuklamasi yuqori ekan, RAM ko'rsatkichini hisobga olmay diskning jismoniy nosozligini yakuniy sabab deyish
\item Foydalanuvchi dasturlaridan qat'i nazar faqat videokarta drayverini qayta o'rnatish bilan cheklanish
\end{enumerate}
\begin{tcolorbox}[enhanced,breakable,colback=teal!4,colframe=teal!60!black,boxrule=0.45pt,arc=1.2mm,left=2.7mm,right=2.7mm,top=1.8mm,bottom=1.8mm,before skip=2mm,after skip=2mm]
{\bfseries\sffamily\color{teal}To'g'ri javob:} \textbf{B) Xotira sarflayotgan jarayonlarni va RAM yetishmagani uchun diskka sahifalash mavjudligini tekshirish}
\par\smallskip
{\bfseries Izoh:} RAM yetishmovchiligi virtual xotiraga murojaatni kuchaytirib, disk bandligini va sekinlashishni oshirishi mumkin. 98\% RAM va disk faolligi pagingni tekshirish uchun asosdir; bu yakuniy tashxis emas. Manba: Cambridge+ 10-11-sinf, PDF 34
\end{tcolorbox}
\vspace{2mm}\hrule height 0.2pt\vspace{3mm}

% DB_ID 130ab2ec-342e-d51c-fdee-dc668ed7a2f8
\needspace{5\baselineskip}
{\bfseries\sffamily\color{ink} 267.} {scriptsizecolor{muted} Y1 | review | 130ab2ec}
\par
Dastur ikkita sonning yig'indisini hisoblab, keyingi buyruqni bajarishga o'tadi. Protsessorning ALU va boshqaruv qurilmasi (CU) vazifalarini qaysi izoh aniq farqlaydi?
\begin{enumerate}[label=\Alph*),leftmargin=9mm]
\item ALU barcha fayllarni diskda saqlaydi, CU esa monitorda tasvir chiqaradi
\item ALU hisoblashni bajaradi, CU esa ko'rsatmalar bajarilishini boshqaradi
\item ALU printer drayverini o'rnatadi, CU esa fayl kengaytmasini o'zgartiradi
\item ALU tarmoq kabelini ulaydi, CU esa faqat qog'oz chop etishga xizmat qiladi
\end{enumerate}
\begin{tcolorbox}[enhanced,breakable,colback=teal!4,colframe=teal!60!black,boxrule=0.45pt,arc=1.2mm,left=2.7mm,right=2.7mm,top=1.8mm,bottom=1.8mm,before skip=2mm,after skip=2mm]
{\bfseries\sffamily\color{teal}To'g'ri javob:} \textbf{B) ALU hisoblashni bajaradi, CU esa ko'rsatmalar bajarilishini boshqaradi}
\par\smallskip
{\bfseries Izoh:} Arifmetik-mantiqiy qurilma (ALU) hisoblash va mantiqiy amallarga, boshqaruv qurilmasi (CU) esa buyruqlarning bajarilishi va qismlar faoliyatini muvofiqlashtirishga xizmat qiladi. [Manba: Cambridge+ 10-11-sinf, PDF 33]
\end{tcolorbox}
\vspace{2mm}\hrule height 0.2pt\vspace{3mm}

% DB_ID 163a0227-5516-8689-4142-591d8d1ec61a
\needspace{5\baselineskip}
{\bfseries\sffamily\color{ink} 268.} {scriptsizecolor{muted} Y1 | review | 163a0227}
\par
Unumdorlik, nosozliklarni aniqlash va xavfsiz xizmat mavzusida: I. paging - qizish sabab protsessor tezligining avtomatik pasayishi. II. bottleneck - tizim umumiy tezligini cheklayotgan eng sust resurs. Mulohazalarni baholang.
\begin{enumerate}[label=\Alph*),leftmargin=9mm]
\item faqat I
\item faqat II
\item I va II
\item I ham, II ham emas
\end{enumerate}
\begin{tcolorbox}[enhanced,breakable,colback=teal!4,colframe=teal!60!black,boxrule=0.45pt,arc=1.2mm,left=2.7mm,right=2.7mm,top=1.8mm,bottom=1.8mm,before skip=2mm,after skip=2mm]
{\bfseries\sffamily\color{teal}To'g'ri javob:} \textbf{B) faqat II}
\par\smallskip
{\bfseries Izoh:} faqat II. Birinchi va ikkinchi mulohazaning ta'riflari alohida tekshirildi; natija: faqat II. Manba: CAM10/11, PDF 24-61; Qo'shimcha sahifa raqami (qaysi PDFga tegishli ekani aniqlanmagan; darslik sahifasi deb tasdiqlanmagan): 110
\end{tcolorbox}
\vspace{2mm}\hrule height 0.2pt\vspace{3mm}

% DB_ID 1919181a-7978-3018-45b6-122e6e53da7e
\needspace{5\baselineskip}
{\bfseries\sffamily\color{ink} 269.} {scriptsizecolor{muted} Y1 | review | 1919181a}
\par
Protsessor ikki sonni taqqoslash natijasini qisqa vaqt saqlashi kerak. ALU va registr vazifalari qaysi variantda to'g'ri berilgan?
\begin{enumerate}[label=\Alph*),leftmargin=9mm]
\item ALU doimiy fayllarni saqlaydi; registr esa monitorga tasvir yuboradi
\item ALU buyruqlarni diskka ko'chiradi; registr faqat printer tezligini boshqaradi
\item ALU faqat matnli fayllarni ochadi; registr esa tarmoq qurilmasiga ulanadi
\item ALU arifmetik va mantiqiy amalni bajaradi; registr natijani qisqa saqlaydi
\end{enumerate}
\begin{tcolorbox}[enhanced,breakable,colback=teal!4,colframe=teal!60!black,boxrule=0.45pt,arc=1.2mm,left=2.7mm,right=2.7mm,top=1.8mm,bottom=1.8mm,before skip=2mm,after skip=2mm]
{\bfseries\sffamily\color{teal}To'g'ri javob:} \textbf{D) ALU arifmetik va mantiqiy amalni bajaradi; registr natijani qisqa saqlaydi}
\par\smallskip
{\bfseries Izoh:} ALU arifmetik va mantiqiy amallarni bajaradi, registr esa protsessorning ichki tezkor kichik xotira elementidir. [Manba: Cambridge+ 10-11-sinf, PDF 33.]
\end{tcolorbox}
\vspace{2mm}\hrule height 0.2pt\vspace{3mm}

% DB_ID 4c23be68-69c2-f1e6-cf71-d748a955f214
\needspace{5\baselineskip}
{\bfseries\sffamily\color{ink} 270.} {scriptsizecolor{muted} Y1 | review | 4c23be68}
\par
Dastur faqat bitta ketma-ket oqimda ishlaydi. 4 yadroli CPU o'rnatilganda qaysi xulosa eng to'g'ri?
\begin{enumerate}[label=\Alph*),leftmargin=9mm]
\item har bir ko'rsatma to'rt nusxada bajariladi
\item tezlik to'rt marta oshishi kafolatlanmaydi
\item RAM hajmi avtomatik to'rt baravar bo'ladi
\item takt chastotasi doim to'rt baravar oshadi
\end{enumerate}
\begin{tcolorbox}[enhanced,breakable,colback=teal!4,colframe=teal!60!black,boxrule=0.45pt,arc=1.2mm,left=2.7mm,right=2.7mm,top=1.8mm,bottom=1.8mm,before skip=2mm,after skip=2mm]
{\bfseries\sffamily\color{teal}To'g'ri javob:} \textbf{B) tezlik to'rt marta oshishi kafolatlanmaydi}
\par\smallskip
{\bfseries Izoh:} tezlik to'rt marta oshishi kafolatlanmaydi. Ketma-ket vazifa barcha yadrolardan parallel foydalana olmasligi mumkin. Manba: CAM9, PDF 10-14; Qo'shimcha sahifa raqami (qaysi PDFga tegishli ekani aniqlanmagan; darslik sahifasi deb tasdiqlanmagan): 22
\end{tcolorbox}
\vspace{2mm}\hrule height 0.2pt\vspace{3mm}

% DB_ID 53bf26f5-c4b6-7932-73ab-49fb3a016ec0
\needspace{5\baselineskip}
{\bfseries\sffamily\color{ink} 271.} {scriptsizecolor{muted} Y1 | review | 53bf26f5}
\par
Ikkilamchi saqlash qurilmalari mavzusida: I. SSD - lazer orqali o'qiladigan va yoziladigan disk. II. HDD - magnit plastinkali, katta sig'imli ikkilamchi xotira. Mulohazalarni baholang.
\begin{enumerate}[label=\Alph*),leftmargin=9mm]
\item faqat I
\item faqat II
\item I va II
\item I ham, II ham emas
\end{enumerate}
\begin{tcolorbox}[enhanced,breakable,colback=teal!4,colframe=teal!60!black,boxrule=0.45pt,arc=1.2mm,left=2.7mm,right=2.7mm,top=1.8mm,bottom=1.8mm,before skip=2mm,after skip=2mm]
{\bfseries\sffamily\color{teal}To'g'ri javob:} \textbf{B) faqat II}
\par\smallskip
{\bfseries Izoh:} faqat II. Birinchi va ikkinchi mulohazaning ta'riflari alohida tekshirildi; natija: faqat II. Manba: CAM10/11, PDF 34-37; Qo'shimcha sahifa raqami (qaysi PDFga tegishli ekani aniqlanmagan; darslik sahifasi deb tasdiqlanmagan): 32
\end{tcolorbox}
\vspace{2mm}\hrule height 0.2pt\vspace{3mm}

% DB_ID 55acbab0-d8bd-be92-eda9-7fc60d075346
\needspace{5\baselineskip}
{\bfseries\sffamily\color{ink} 272.} {scriptsizecolor{muted} Y1 | review | 55acbab0}
\par
Attestatsiyadagi nozik farqlar va integratsion masalalar mavzusida: I. integratsion savol - variantlarni baholash uchun qo'llanadigan belgi. II. tasnif - bir nechta mavzuni bitta vaziyatda bog'laydigan topshiriq. Mulohazalarni baholang.
\begin{enumerate}[label=\Alph*),leftmargin=9mm]
\item faqat I
\item I ham, II ham emas
\item faqat II
\item I va II
\end{enumerate}
\begin{tcolorbox}[enhanced,breakable,colback=teal!4,colframe=teal!60!black,boxrule=0.45pt,arc=1.2mm,left=2.7mm,right=2.7mm,top=1.8mm,bottom=1.8mm,before skip=2mm,after skip=2mm]
{\bfseries\sffamily\color{teal}To'g'ri javob:} \textbf{B) I ham, II ham emas}
\par\smallskip
{\bfseries Izoh:} I ham, II ham emas. Birinchi va ikkinchi mulohazaning ta'riflari alohida tekshirildi; natija: I ham, II ham emas. Manba: ICT5, PDF 25-49; Qo'shimcha sahifa raqami (qaysi PDFga tegishli ekani aniqlanmagan; darslik sahifasi deb tasdiqlanmagan): 119
\end{tcolorbox}
\vspace{2mm}\hrule height 0.2pt\vspace{3mm}

% DB_ID 5c7ab97f-18e6-a93e-fe74-09fa40fc20a1
\needspace{5\baselineskip}
{\bfseries\sffamily\color{ink} 273.} {scriptsizecolor{muted} Y1 | review | 5c7ab97f}
\par
Arxiv markazi eski ma'lumotni magnit lentada saqlaydi, ishchi fayllarga esa HDD orqali tez-tez murojaat qiladi. Qaysi taqqoslash saqlash vositalariga mos?
\begin{enumerate}[label=\Alph*),leftmargin=9mm]
\item Magnit lenta barcha joyga bevosita murojaat qiladi; HDD esa faqat ketma-ket bitlarni oxirigacha o'qiydi
\item Magnit lenta flesh hujayralarga asoslanadi; HDD esa doimiy xotira uchun ishlatilmaydi
\item Magnit lenta istalgan sektorni darhol ochadi; HDD arxivdan ma'lumot qidirishga yaroqsiz
\item Magnit lenta ketma-ket arxivlashga, HDD esa bevosita murojaat qilishga mos
\end{enumerate}
\begin{tcolorbox}[enhanced,breakable,colback=teal!4,colframe=teal!60!black,boxrule=0.45pt,arc=1.2mm,left=2.7mm,right=2.7mm,top=1.8mm,bottom=1.8mm,before skip=2mm,after skip=2mm]
{\bfseries\sffamily\color{teal}To'g'ri javob:} \textbf{D) Magnit lenta ketma-ket arxivlashga, HDD esa bevosita murojaat qilishga mos}
\par\smallskip
{\bfseries Izoh:} Darslikda HDD ma'lumotga to'g'ridan-to'g'ri murojaat qilishi va magnit lenta ma'lumotni ketma-ketlik bilan saqlab o'qishi farqlanadi. [Manba: Cambridge+ 9-sinf, PDF 41-43.]
\end{tcolorbox}
\vspace{2mm}\hrule height 0.2pt\vspace{3mm}

% DB_ID 66f696a9-868a-878f-4d29-a12001d5d34d
\needspace{5\baselineskip}
{\bfseries\sffamily\color{ink} 274.} {scriptsizecolor{muted} Y1 | review | 66f696a9}
\par
Mutaxassis yangi grafik kartani tizim blokiga o'rnatib, keyin tashqi monitorni ulayapti. Bosh plata va ulanish portining vazifalari qaysi variantda to'g'ri farqlangan?
\begin{enumerate}[label=\Alph*),leftmargin=9mm]
\item Bosh plata faqat tashqi qurilmalarni ro'yxatlaydi, port esa protsessor hisoblash tezligini o'zgartirib turadi
\item Bosh plata ichki komponentlarni bog'laydi; port qurilmalarni ulash nuqtasini ta'minlaydi
\item Bosh plata faqat saqlash qurilmasi sifatida ishlaydi, port esa dasturlarning buyruqlarini tarjima qiladi
\item Bosh plata tashqi tarmoq hisoblanadi, port esa CPU ichidagi arifmetik birlikni boshqaradi
\end{enumerate}
\begin{tcolorbox}[enhanced,breakable,colback=teal!4,colframe=teal!60!black,boxrule=0.45pt,arc=1.2mm,left=2.7mm,right=2.7mm,top=1.8mm,bottom=1.8mm,before skip=2mm,after skip=2mm]
{\bfseries\sffamily\color{teal}To'g'ri javob:} \textbf{B) Bosh plata ichki komponentlarni bog'laydi; port qurilmalarni ulash nuqtasini ta'minlaydi}
\par\smallskip
{\bfseries Izoh:} Bosh plata protsessor, xotira, kengaytma kabi ichki komponentlarni bog'laydi. Port tashqi yoki boshqa mos qurilmalarni ulash uchun interfeys vazifasini bajaradi. [Manba: Cambridge+ 10-11-sinf, PDF 33-34.]
\end{tcolorbox}
\vspace{2mm}\hrule height 0.2pt\vspace{3mm}

% DB_ID 78dddebe-40e5-92ab-5672-ab91be3539b5
\needspace{5\baselineskip}
{\bfseries\sffamily\color{ink} 275.} {scriptsizecolor{muted} Y1 | review | 78dddebe}
\par
Attestatsiyadagi nozik farqlar va integratsion masalalar mavzusida: I. mezon - yechim bajarishi shart bo'lgan chegara yoki talab. II. tasnif - obyektni muhim belgilari bo'yicha guruhga ajratish. Mulohazalarni baholang.
\begin{enumerate}[label=\Alph*),leftmargin=9mm]
\item faqat I
\item I va II
\item I ham, II ham emas
\item faqat II
\end{enumerate}
\begin{tcolorbox}[enhanced,breakable,colback=teal!4,colframe=teal!60!black,boxrule=0.45pt,arc=1.2mm,left=2.7mm,right=2.7mm,top=1.8mm,bottom=1.8mm,before skip=2mm,after skip=2mm]
{\bfseries\sffamily\color{teal}To'g'ri javob:} \textbf{D) faqat II}
\par\smallskip
{\bfseries Izoh:} faqat II. Birinchi va ikkinchi mulohazaning ta'riflari alohida tekshirildi; natija: faqat II. Manba: CAM9, PDF 8-47; Qo'shimcha sahifa raqami (qaysi PDFga tegishli ekani aniqlanmagan; darslik sahifasi deb tasdiqlanmagan): 117
\end{tcolorbox}
\vspace{2mm}\hrule height 0.2pt\vspace{3mm}

% DB_ID 8f6e2519-6ce8-8654-ebf2-53d689214fa5
\needspace{5\baselineskip}
{\bfseries\sffamily\color{ink} 276.} {scriptsizecolor{muted} Y1 | review | 8f6e2519}
\par
Tadqiqotchi qog'ozdagi maqolani avval tasvir sifatida raqamlashtirdi, keyin undagi yozuvni ajratib olib matn protsessorida tahrirladi. Vazifalar to'g'ri taqsimoti qaysi?
\begin{enumerate}[label=\Alph*),leftmargin=9mm]
\item Skaner - tasvirni kiritadi; OCR - belgini matn sifatida tanidi; matn protsessori - tahrirlaydi
\item Skaner - yozuvni tahrirlaydi; OCR - chop etadi; matn protsessori - tasvirni sensor bilan o'qiydi
\item Skaner - matnni kompilyatsiya qiladi; OCR - RAMni boshqaradi; matn protsessori - diskni formatlaydi
\item Skaner - faylni saqlash formatini tanlaydi; OCR - rasmni proyeksiyalaydi; matn protsessori - sensorni sozlaydi
\end{enumerate}
\begin{tcolorbox}[enhanced,breakable,colback=teal!4,colframe=teal!60!black,boxrule=0.45pt,arc=1.2mm,left=2.7mm,right=2.7mm,top=1.8mm,bottom=1.8mm,before skip=2mm,after skip=2mm]
{\bfseries\sffamily\color{teal}To'g'ri javob:} \textbf{A) Skaner - tasvirni kiritadi; OCR - belgini matn sifatida tanidi; matn protsessori - tahrirlaydi}
\par\smallskip
{\bfseries Izoh:} Skaner qog'oz tasvirini kiritadi; OCR belgilarni matnga aylantiradi; matn protsessori hosil bo'lgan matnni tahrirlaydi. Manba: Cambridge+ 9-sinf, PDF 29-31
\end{tcolorbox}
\vspace{2mm}\hrule height 0.2pt\vspace{3mm}

% DB_ID 96e1a754-f913-909f-c8fa-83b1eea624ce
\needspace{5\baselineskip}
{\bfseries\sffamily\color{ink} 277.} {scriptsizecolor{muted} Y1 | review | 96e1a754}
\par
Maktab barcha zaxira fayllarni asosiy server bilan bir xonada saqlamoqda. Xonadagi yong'in xavfini hisobga olganda qaysi chorani ustuvor tanlash kerak?
\begin{enumerate}[label=\Alph*),leftmargin=9mm]
\item Serverdagi xotirani kattalashtirib, barcha nusxalarni shu serverda saqlash
\item Har bir zaxirani faqat vaqtinchalik RAMga joylab qo'yish va test qilish
\item Nusxalar sonini oshirib, ularni bir xil xonada turli papkalarda saqlash
\item Zaxiraning kamida bir mustaqil nusxasini boshqa jismoniy joyda saqlash
\end{enumerate}
\begin{tcolorbox}[enhanced,breakable,colback=teal!4,colframe=teal!60!black,boxrule=0.45pt,arc=1.2mm,left=2.7mm,right=2.7mm,top=1.8mm,bottom=1.8mm,before skip=2mm,after skip=2mm]
{\bfseries\sffamily\color{teal}To'g'ri javob:} \textbf{D) Zaxiraning kamida bir mustaqil nusxasini boshqa jismoniy joyda saqlash}
\par\smallskip
{\bfseries Izoh:} Darslikda maktab serveridagi ma'lumotning har kunlik nusxasini maktabdan tashqarida saqlash vaziyati beriladi. Bir joydagi halokatdan himoya qilish uchun uzoqdagi nusxa kerak. [Manba: Cambridge+ 10-11-sinf, PDF 43-bet.]
\end{tcolorbox}
\vspace{2mm}\hrule height 0.2pt\vspace{3mm}

% DB_ID 97017dc9-1e55-dec7-da4c-5927ed0a5ff0
\needspace{5\baselineskip}
{\bfseries\sffamily\color{ink} 278.} {scriptsizecolor{muted} Y1 | review | 97017dc9}
\par
Birlamchi xotira: registr, kesh, RAM va ROM mavzusida: I. ROM - tez-tez kerak bo'ladigan ma'lumotni CPUga yaqin saqlovchi xotira. II. RAM - faol dastur va ma'lumotlar uchun o'zgaruvchan tezkor xotira. Mulohazalarni baholang.
\begin{enumerate}[label=\Alph*),leftmargin=9mm]
\item faqat I
\item I va II
\item I ham, II ham emas
\item faqat II
\end{enumerate}
\begin{tcolorbox}[enhanced,breakable,colback=teal!4,colframe=teal!60!black,boxrule=0.45pt,arc=1.2mm,left=2.7mm,right=2.7mm,top=1.8mm,bottom=1.8mm,before skip=2mm,after skip=2mm]
{\bfseries\sffamily\color{teal}To'g'ri javob:} \textbf{D) faqat II}
\par\smallskip
{\bfseries Izoh:} faqat II. Birinchi va ikkinchi mulohazaning ta'riflari alohida tekshirildi; natija: faqat II. Manba: CAM10/11, PDF 33-36; Qo'shimcha sahifa raqami (qaysi PDFga tegishli ekani aniqlanmagan; darslik sahifasi deb tasdiqlanmagan): 25
\end{tcolorbox}
\vspace{2mm}\hrule height 0.2pt\vspace{3mm}

% DB_ID 98b7fe27-5e4a-8ce5-d119-fa9be5764425
\needspace{5\baselineskip}
{\bfseries\sffamily\color{ink} 279.} {scriptsizecolor{muted} Y1 | review | 98b7fe27}
\par
Do'konda skaner mahsulot kodini o'qidi, dastur narxni hisoblab ekranga chiqardi va sotuvni diskka saqladi. Bu jarayonning IPOS modeliga mos ketma-ketligini toping.
\begin{enumerate}[label=\Alph*),leftmargin=9mm]
\item Protsessor \ensuremath{\rightarrow} skaner \ensuremath{\rightarrow} sotuv fayli \ensuremath{\rightarrow} ekran: hisoblashdan so'ng kiritish bosqichi bajariladi
\item Skaner \ensuremath{\rightarrow} ekran \ensuremath{\rightarrow} protsessor \ensuremath{\rightarrow} disk: natija hisoblash boshlanishidan avval chiqariladi
\item Ekran \ensuremath{\rightarrow} skaner \ensuremath{\rightarrow} disk \ensuremath{\rightarrow} protsessor: foydalanuvchi ekranda ko'ringan natijani avval kiritadi, so'ng hisoblaydi
\item Skaner \ensuremath{\rightarrow} narxni hisoblash \ensuremath{\rightarrow} ekran \ensuremath{\rightarrow} sotuv yozuvi: kiritish, qayta ishlash, chiqarish, saqlash
\end{enumerate}
\begin{tcolorbox}[enhanced,breakable,colback=teal!4,colframe=teal!60!black,boxrule=0.45pt,arc=1.2mm,left=2.7mm,right=2.7mm,top=1.8mm,bottom=1.8mm,before skip=2mm,after skip=2mm]
{\bfseries\sffamily\color{teal}To'g'ri javob:} \textbf{D) Skaner \ensuremath{\rightarrow} narxni hisoblash \ensuremath{\rightarrow} ekran \ensuremath{\rightarrow} sotuv yozuvi: kiritish, qayta ishlash, chiqarish, saqlash}
\par\smallskip
{\bfseries Izoh:} Skaner axborotni kiritadi, protsessor hisoblaydi, ekran natijani chiqaradi, disk esa ma'lumotni saqlaydi. Bu kiritish-qayta ishlash-chiqarish-saqlash modelidir. Manba: Cambridge+ 9-sinf, PDF 9-10
\end{tcolorbox}
\vspace{2mm}\hrule height 0.2pt\vspace{3mm}

% DB_ID a9db982f-d2b7-7107-f1f9-30be287831af
\needspace{5\baselineskip}
{\bfseries\sffamily\color{ink} 280.} {scriptsizecolor{muted} Y1 | review | a9db982f}
\par
8 GB RAMga ega kompyuterda ko'plab dasturlar ochilganidan keyin SSDga murojaat keskin ko'paydi va tizim sekinlashdi. Buni qanday izohlash mumkin?
\begin{enumerate}[label=\Alph*),leftmargin=9mm]
\item SSDga kirish RAMga kirishdan tezroq bo'lgani uchun dasturlar ma'lumotni boshqa joyga o'tkazmoqda
\item RAM yetishmayotgani sababli SSD maydonidan virtual xotira sifatida foydalanilmoqda
\item ROMdagi boshlang'ich dastur ko'payib, foydalanuvchi ilovalari uchun ajratilgan joyni egallamoqda
\item CPU keshining hajmi ortgani uchun faol jarayonlar barcha saqlash qurilmalariga bir vaqtda yozmoqda
\end{enumerate}
\begin{tcolorbox}[enhanced,breakable,colback=teal!4,colframe=teal!60!black,boxrule=0.45pt,arc=1.2mm,left=2.7mm,right=2.7mm,top=1.8mm,bottom=1.8mm,before skip=2mm,after skip=2mm]
{\bfseries\sffamily\color{teal}To'g'ri javob:} \textbf{B) RAM yetishmayotgani sababli SSD maydonidan virtual xotira sifatida foydalanilmoqda}
\par\smallskip
{\bfseries Izoh:} RAM to'lganda tizim HDD yoki SSDdagi ikkilamchi xotira maydonidan virtual xotira sifatida foydalanishi mumkin. Bunday murojaat RAMga qaraganda sekinroq. Manba: Cambridge+ 10-11-sinf, PDF 34
\end{tcolorbox}
\vspace{2mm}\hrule height 0.2pt\vspace{3mm}

% DB_ID ad232a83-50cb-7e74-8f0a-5445dbcb4ad6
\needspace{5\baselineskip}
{\bfseries\sffamily\color{ink} 281.} {scriptsizecolor{muted} Y1 | review | ad232a83}
\par
Chiqarish va ikki vazifali qurilmalar mavzusida: I. qattiq nusxa - toner va elektrofotografik jarayon bilan tez chop etuvchi printer. II. yumshoq nusxa - ekranda yoki audio orqali vaqtinchalik taqdim etilgan natija. Mulohazalarni baholang.
\begin{enumerate}[label=\Alph*),leftmargin=9mm]
\item faqat I
\item faqat II
\item I va II
\item I ham, II ham emas
\end{enumerate}
\begin{tcolorbox}[enhanced,breakable,colback=teal!4,colframe=teal!60!black,boxrule=0.45pt,arc=1.2mm,left=2.7mm,right=2.7mm,top=1.8mm,bottom=1.8mm,before skip=2mm,after skip=2mm]
{\bfseries\sffamily\color{teal}To'g'ri javob:} \textbf{B) faqat II}
\par\smallskip
{\bfseries Izoh:} faqat II. Birinchi va ikkinchi mulohazaning ta'riflari alohida tekshirildi; natija: faqat II. Manba: CAM10/11, PDF 37-45; Qo'shimcha sahifa raqami (qaysi PDFga tegishli ekani aniqlanmagan; darslik sahifasi deb tasdiqlanmagan): 44
\end{tcolorbox}
\vspace{2mm}\hrule height 0.2pt\vspace{3mm}

% DB_ID b244f8c5-067c-5da6-af64-0298f4499266
\needspace{5\baselineskip}
{\bfseries\sffamily\color{ink} 282.} {scriptsizecolor{muted} Y1 | review | b244f8c5}
\par
Kompyuter tizimi va uning ishlash modeli mavzusida: I. texnik ta'minot - kompyuterga vazifalarni bajarishni ko'rsatuvchi dasturlar. II. kompyuter tizimi - bir maqsad uchun hamkorlik qiluvchi texnik vosita, dastur, ma'lumot va foydalanuvchi majmuasi. Mulohazalarni baholang.
\begin{enumerate}[label=\Alph*),leftmargin=9mm]
\item faqat I
\item faqat II
\item I va II
\item I ham, II ham emas
\end{enumerate}
\begin{tcolorbox}[enhanced,breakable,colback=teal!4,colframe=teal!60!black,boxrule=0.45pt,arc=1.2mm,left=2.7mm,right=2.7mm,top=1.8mm,bottom=1.8mm,before skip=2mm,after skip=2mm]
{\bfseries\sffamily\color{teal}To'g'ri javob:} \textbf{B) faqat II}
\par\smallskip
{\bfseries Izoh:} faqat II. Birinchi va ikkinchi mulohazaning ta'riflari alohida tekshirildi; natija: faqat II. Manba: CAM10/11, PDF 31-35; Qo'shimcha sahifa raqami (qaysi PDFga tegishli ekani aniqlanmagan; darslik sahifasi deb tasdiqlanmagan): 5
\end{tcolorbox}
\vspace{2mm}\hrule height 0.2pt\vspace{3mm}

% DB_ID ba27923a-0096-4fd8-0cea-298bc0df1c7f
\needspace{5\baselineskip}
{\bfseries\sffamily\color{ink} 283.} {scriptsizecolor{muted} Y1 | review | ba27923a}
\par
Muhandis tez-tez joydan joyga ko'chib, katta ma'lumot fayllari bilan ishlaydi. Qurilma tez bo'lishi va tushib ketishga chidamliroq bo'lishi kerak. Qaysi tanlov asosli?
\begin{enumerate}[label=\Alph*),leftmargin=9mm]
\item Mexanik diski mavjud HDDni, harakatlanuvchi qismlarning zarbaga sezgirligini hisobga olmasdan tanlash
\item Optik diskni, qayta-qayta katta fayl yozishdagi sig'im va tezlik cheklovlarini tekshirmasdan tanlash
\item USB flesh qurilmani, interfeysning ko'chirish tezligi va yozish resursini baholamasdan tanlash
\item Harakatlanuvchi qismi bo'lmagan, talabga mos sig'im hamda tezlikdagi tashqi SSDni tanlash
\end{enumerate}
\begin{tcolorbox}[enhanced,breakable,colback=teal!4,colframe=teal!60!black,boxrule=0.45pt,arc=1.2mm,left=2.7mm,right=2.7mm,top=1.8mm,bottom=1.8mm,before skip=2mm,after skip=2mm]
{\bfseries\sffamily\color{teal}To'g'ri javob:} \textbf{D) Harakatlanuvchi qismi bo'lmagan, talabga mos sig'im hamda tezlikdagi tashqi SSDni tanlash}
\par\smallskip
{\bfseries Izoh:} Darslik SSDni HDDga qaraganda tezroq hamda mexanik harakatlanuvchi qismlarga ega bo'lmagan yechim sifatida taqqoslaydi. Ko'chma va zarba xavfi bo'lgan vaziyatda mos xususiyatli tashqi SSD ma'qul. Manba: Cambridge+ 10-11-sinf, PDF 39-41
\end{tcolorbox}
\vspace{2mm}\hrule height 0.2pt\vspace{3mm}

% DB_ID c4ed559c-a4d3-43ee-f10d-5461d250ce9e
\needspace{5\baselineskip}
{\bfseries\sffamily\color{ink} 284.} {scriptsizecolor{muted} Y1 | review | c4ed559c}
\par
Kiritish qurilmalari mavzusida: I. OMR - magnit siyoh bilan bosilgan maxsus belgilarni o'qish. II. OCR - tasvirdagi bosma matnni tahrirlanadigan belgilar sifatida tanish. Mulohazalarni baholang.
\begin{enumerate}[label=\Alph*),leftmargin=9mm]
\item faqat I
\item I va II
\item I ham, II ham emas
\item faqat II
\end{enumerate}
\begin{tcolorbox}[enhanced,breakable,colback=teal!4,colframe=teal!60!black,boxrule=0.45pt,arc=1.2mm,left=2.7mm,right=2.7mm,top=1.8mm,bottom=1.8mm,before skip=2mm,after skip=2mm]
{\bfseries\sffamily\color{teal}To'g'ri javob:} \textbf{D) faqat II}
\par\smallskip
{\bfseries Izoh:} faqat II. Birinchi va ikkinchi mulohazaning ta'riflari alohida tekshirildi; natija: faqat II. Manba: CAM10/11, PDF 37-45; Qo'shimcha sahifa raqami (qaysi PDFga tegishli ekani aniqlanmagan; darslik sahifasi deb tasdiqlanmagan): 38
\end{tcolorbox}
\vspace{2mm}\hrule height 0.2pt\vspace{3mm}

% DB_ID cd48e53d-c566-22e3-ed6c-c004d21a99b6
\needspace{5\baselineskip}
{\bfseries\sffamily\color{ink} 285.} {scriptsizecolor{muted} Y1 | review | cd48e53d}
\par
Unumdorlik, nosozliklarni aniqlash va xavfsiz xizmat mavzusida: I. diagnostika - RAM va ikkilamchi xotira o'rtasida sahifalarni almashish. II. bottleneck - belgi va dalillar asosida sababni tizimli aniqlash. Mulohazalarni baholang.
\begin{enumerate}[label=\Alph*),leftmargin=9mm]
\item faqat I
\item faqat II
\item I va II
\item I ham, II ham emas
\end{enumerate}
\begin{tcolorbox}[enhanced,breakable,colback=teal!4,colframe=teal!60!black,boxrule=0.45pt,arc=1.2mm,left=2.7mm,right=2.7mm,top=1.8mm,bottom=1.8mm,before skip=2mm,after skip=2mm]
{\bfseries\sffamily\color{teal}To'g'ri javob:} \textbf{D) I ham, II ham emas}
\par\smallskip
{\bfseries Izoh:} I ham, II ham emas. Birinchi va ikkinchi mulohazaning ta'riflari alohida tekshirildi; natija: I ham, II ham emas. Manba: MAR, PDF 10-20; Qo'shimcha sahifa raqami (qaysi PDFga tegishli ekani aniqlanmagan; darslik sahifasi deb tasdiqlanmagan): 112
\end{tcolorbox}
\vspace{2mm}\hrule height 0.2pt\vspace{3mm}

% DB_ID e1d8d26f-f44a-8453-2f34-bf5cdecc3080
\needspace{5\baselineskip}
{\bfseries\sffamily\color{ink} 286.} {scriptsizecolor{muted} Y1 | review | e1d8d26f}
\par
O'quvchi hujjatni ekranda ko'rib chiqdi, keyin uni qog'ozga siyohpurkagich printerda chiqardi. Qaysi taqqoslash to'g'ri?
\begin{enumerate}[label=\Alph*),leftmargin=9mm]
\item Monitor qog'ozdagi hard copy beradi, printer esa faqat ekrandagi soft copy beradi
\item Monitor tasvirni saqlash qurilmasi, printer esa matnni kiritish qurilmasi hisoblanadi
\item Ekrandagi tasvir hard copy, printerdagi fizik varaq esa vaqtinchalik RAM nusxasidir
\item Ekrandagi ko'rinish soft copy, chop etilgan qog'oz esa hard copy hisoblanadi
\end{enumerate}
\begin{tcolorbox}[enhanced,breakable,colback=teal!4,colframe=teal!60!black,boxrule=0.45pt,arc=1.2mm,left=2.7mm,right=2.7mm,top=1.8mm,bottom=1.8mm,before skip=2mm,after skip=2mm]
{\bfseries\sffamily\color{teal}To'g'ri javob:} \textbf{D) Ekrandagi ko'rinish soft copy, chop etilgan qog'oz esa hard copy hisoblanadi}
\par\smallskip
{\bfseries Izoh:} Ekranda ko'rsatilgan ma'lumot soft copy, printerda qog'ozga chiqarilgan nusxa hard copy sifatida farqlanadi. [Manba: Cambridge+ 9-sinf, PDF 35-38.]
\end{tcolorbox}
\vspace{2mm}\hrule height 0.2pt\vspace{3mm}

% DB_ID e4d100a9-7104-1bfc-fd12-31d90fce34aa
\needspace{5\baselineskip}
{\bfseries\sffamily\color{ink} 287.} {scriptsizecolor{muted} Y1 | review | e4d100a9}
\par
Kompyuterga videokarta va operativ xotira o'rnatilishi kerak. Bosh plata, shina va kengaytirish slotining vazifalarini qaysi izoh eng to'g'ri farqlaydi?
\begin{enumerate}[label=\Alph*),leftmargin=9mm]
\item Bosh plata fayl tizimini boshqaradi; shina esa faqat dastur kodini tarjima qiladi
\item Shina tashqi printer turini belgilaydi; bosh plata esa faqat SSD hajmini oshiradi
\item Kengaytirish sloti protsessordagi arifmetik amallarni bajaradi, shina esa printerga qog'oz beradi
\item Bosh plata qismlarni birlashtiradi, shina signal tashiydi, slot esa kartani ulaydi
\end{enumerate}
\begin{tcolorbox}[enhanced,breakable,colback=teal!4,colframe=teal!60!black,boxrule=0.45pt,arc=1.2mm,left=2.7mm,right=2.7mm,top=1.8mm,bottom=1.8mm,before skip=2mm,after skip=2mm]
{\bfseries\sffamily\color{teal}To'g'ri javob:} \textbf{D) Bosh plata qismlarni birlashtiradi, shina signal tashiydi, slot esa kartani ulaydi}
\par\smallskip
{\bfseries Izoh:} Bosh plata ichki komponentlarni birlashtiradi. Shina ularning o'zaro signal almashish yo'li, kengaytirish sloti esa mos qo'shimcha kartani ulash joyidir. [Manba: Cambridge+ 10-11-sinf, PDF 33-34]
\end{tcolorbox}
\vspace{2mm}\hrule height 0.2pt\vspace{3mm}
\clearpage
\section{Operatsion tizimlar va ularning imkoniyatlari}
\subsection{Nashr qilingan testlar}

% DB_ID 02a00696-d826-4a76-be6a-7ce735778bfe
\needspace{5\baselineskip}
{\bfseries\sffamily\color{ink} 288.} {scriptsizecolor{muted} Y1 | published | 02a00696}
\par
CLI bilan ishlashning asosiy belgisi qaysi?
\begin{enumerate}[label=\Alph*),leftmargin=9mm]
\item Buyruqlarni matn ko'rinishida kiritish
\item Faqat sichqoncha bilan ishlash
\item Faqat rasmlar bilan ishlash
\item Faqat audio buyruqlar
\end{enumerate}
\begin{tcolorbox}[enhanced,breakable,colback=teal!4,colframe=teal!60!black,boxrule=0.45pt,arc=1.2mm,left=2.7mm,right=2.7mm,top=1.8mm,bottom=1.8mm,before skip=2mm,after skip=2mm]
{\bfseries\sffamily\color{teal}To'g'ri javob:} \textbf{A) Buyruqlarni matn ko'rinishida kiritish}
\par\smallskip
{\bfseries Izoh:} CLI foydalanuvchidan matnli buyruqlarni kiritishni talab qiladi.
\end{tcolorbox}
\vspace{2mm}\hrule height 0.2pt\vspace{3mm}

% DB_ID 2bfa394d-6f8a-4a26-88a2-466e17d1f415
\needspace{5\baselineskip}
{\bfseries\sffamily\color{ink} 289.} {scriptsizecolor{muted} Y1 | published | 2bfa394d}
\par
GUI nimani anglatadi?
\begin{enumerate}[label=\Alph*),leftmargin=9mm]
\item Grafik foydalanuvchi interfeysi
\item Faqat matnli buyruqlar oynasi
\item Doimiy xotira turi
\item Fayl siqish algoritmi
\end{enumerate}
\begin{tcolorbox}[enhanced,breakable,colback=teal!4,colframe=teal!60!black,boxrule=0.45pt,arc=1.2mm,left=2.7mm,right=2.7mm,top=1.8mm,bottom=1.8mm,before skip=2mm,after skip=2mm]
{\bfseries\sffamily\color{teal}To'g'ri javob:} \textbf{A) Grafik foydalanuvchi interfeysi}
\par\smallskip
{\bfseries Izoh:} GUI oynalar, menyular, tugmalar va piktogrammalar kabi grafik elementlar orqali muloqotni ta'minlaydi.
\end{tcolorbox}
\vspace{2mm}\hrule height 0.2pt\vspace{3mm}

% DB_ID 513ae5c6-452e-4bc3-b78e-e4eb014e52dd
\needspace{5\baselineskip}
{\bfseries\sffamily\color{ink} 290.} {scriptsizecolor{muted} Y1 | published | 513ae5c6}
\par
Kompyuterda bir nechta dastur ishlaganda xotira, protsessor va qurilma resurslaridan foydalanishni muvofiqlashtiradigan asosiy tizim dasturining vazifasi nima?
\begin{enumerate}[label=\Alph*),leftmargin=9mm]
\item Dasturlar va apparat resurslarini boshqarish
\item Fayllarni arxivlash va siqishni boshqarish
\item Veb-sahifalar va brauzer tarixini boshqarish
\item Bitta faylni tahrirlash va chop etish
\end{enumerate}
\begin{tcolorbox}[enhanced,breakable,colback=teal!4,colframe=teal!60!black,boxrule=0.45pt,arc=1.2mm,left=2.7mm,right=2.7mm,top=1.8mm,bottom=1.8mm,before skip=2mm,after skip=2mm]
{\bfseries\sffamily\color{teal}To'g'ri javob:} \textbf{A) Dasturlar va apparat resurslarini boshqarish}
\par\smallskip
{\bfseries Izoh:} Operatsion tizim jarayonlarni, tezkor xotirani, qurilmalarga murojaatni va boshqa tizim resurslarini muvofiqlashtiradi. Fayllarni siqish yoki matnni tahrirlash alohida ilovalarning vazifasidir.
\end{tcolorbox}
\vspace{2mm}\hrule height 0.2pt\vspace{3mm}

% DB_ID b3eb84b6-e239-4d47-99ec-f0af1609a749
\needspace{5\baselineskip}
{\bfseries\sffamily\color{ink} 291.} {scriptsizecolor{muted} Y1 | published | b3eb84b6}
\par
Quyidagilardan qaysi biri mobil operatsion tizimga misol?
\begin{enumerate}[label=\Alph*),leftmargin=9mm]
\item Android
\item Microsoft Word
\item Adobe Photoshop
\item WinRAR
\end{enumerate}
\begin{tcolorbox}[enhanced,breakable,colback=teal!4,colframe=teal!60!black,boxrule=0.45pt,arc=1.2mm,left=2.7mm,right=2.7mm,top=1.8mm,bottom=1.8mm,before skip=2mm,after skip=2mm]
{\bfseries\sffamily\color{teal}To'g'ri javob:} \textbf{A) Android}
\par\smallskip
{\bfseries Izoh:} Android mobil qurilmalarda ishlatiladigan operatsion tizimdir.
\end{tcolorbox}
\vspace{2mm}\hrule height 0.2pt\vspace{3mm}

% DB_ID c069be13-d3f1-5b0b-8649-17384910a151
\needspace{5\baselineskip}
{\bfseries\sffamily\color{ink} 292.} {scriptsizecolor{muted} Y1 | published | c069be13}
\par
Grafik foydalanuvchi interfeysida ikonka nima?
\begin{enumerate}[label=\Alph*),leftmargin=9mm]
\item Matn kiritiladigan maydon
\item Obyekt, fayl, dastur yoki amalni ifodalovchi kichik grafik belgi
\item Klaviatura tugmalari majmuasi
\item Printerdan chiqarilgan hujjat
\end{enumerate}
\begin{tcolorbox}[enhanced,breakable,colback=teal!4,colframe=teal!60!black,boxrule=0.45pt,arc=1.2mm,left=2.7mm,right=2.7mm,top=1.8mm,bottom=1.8mm,before skip=2mm,after skip=2mm]
{\bfseries\sffamily\color{teal}To'g'ri javob:} \textbf{B) Obyekt, fayl, dastur yoki amalni ifodalovchi kichik grafik belgi}
\par\smallskip
{\bfseries Izoh:} Ikonka - obyekt yoki amalni ko'rgazmali ifodalovchi kichik grafik belgi.
\end{tcolorbox}
\vspace{2mm}\hrule height 0.2pt\vspace{3mm}

% DB_ID 10724a45-d2ff-487c-8626-2814ea16c998
\needspace{5\baselineskip}
{\bfseries\sffamily\color{ink} 293.} {scriptsizecolor{muted} Y1 | published | 10724a45}
\par
Kompyuterda bir nechta dastur ishlayotganda xotira va protsessor vaqtini taqsimlashni qaysi dasturiy qatlam boshqaradi?
\begin{enumerate}[label=\Alph*),leftmargin=9mm]
\item Operatsion tizim
\item Matn protsessori
\item Rasm muharriri
\item Brauzer kengaytmasi
\end{enumerate}
\begin{tcolorbox}[enhanced,breakable,colback=teal!4,colframe=teal!60!black,boxrule=0.45pt,arc=1.2mm,left=2.7mm,right=2.7mm,top=1.8mm,bottom=1.8mm,before skip=2mm,after skip=2mm]
{\bfseries\sffamily\color{teal}To'g'ri javob:} \textbf{A) Operatsion tizim}
\par\smallskip
{\bfseries Izoh:} Resurslarni dasturlar o'rtasida taqsimlash operatsion tizim vazifasidir.
\end{tcolorbox}
\vspace{2mm}\hrule height 0.2pt\vspace{3mm}

% DB_ID 10d6a6fd-8454-40ad-8825-00c3c0895b3a
\needspace{5\baselineskip}
{\bfseries\sffamily\color{ink} 294.} {scriptsizecolor{muted} Y1 | published | 10d6a6fd}
\par
Telefon ilovalarni ishga tushiradi, fayllarni boshqaradi va sensorlar bilan ishlaydi. Bu vazifalarni muvofiqlashtiruvchi qatlam qaysi?
\begin{enumerate}[label=\Alph*),leftmargin=9mm]
\item Mobil operatsion tizim
\item Faqat brauzer
\item Matn muharriri
\item Arxivator
\end{enumerate}
\begin{tcolorbox}[enhanced,breakable,colback=teal!4,colframe=teal!60!black,boxrule=0.45pt,arc=1.2mm,left=2.7mm,right=2.7mm,top=1.8mm,bottom=1.8mm,before skip=2mm,after skip=2mm]
{\bfseries\sffamily\color{teal}To'g'ri javob:} \textbf{A) Mobil operatsion tizim}
\par\smallskip
{\bfseries Izoh:} Mobil operatsion tizim qurilma resurslari va ilovalar ishini boshqaradi.
\end{tcolorbox}
\vspace{2mm}\hrule height 0.2pt\vspace{3mm}

% DB_ID 13167a4f-fd6f-492e-b953-bb8400b35bdf
\needspace{5\baselineskip}
{\bfseries\sffamily\color{ink} 295.} {scriptsizecolor{muted} Y1 | published | 13167a4f}
\par
Yangi printer ulandi, lekin operatsion tizim uni taniy olmadi. Eng mantiqiy keyingi qadam qaysi?
\begin{enumerate}[label=\Alph*),leftmargin=9mm]
\item Mos qurilma drayverini o'rnatish
\item Word faylini qayta nomlash
\item RAMni formatlash
\item Brauzer tarixini tozalash
\end{enumerate}
\begin{tcolorbox}[enhanced,breakable,colback=teal!4,colframe=teal!60!black,boxrule=0.45pt,arc=1.2mm,left=2.7mm,right=2.7mm,top=1.8mm,bottom=1.8mm,before skip=2mm,after skip=2mm]
{\bfseries\sffamily\color{teal}To'g'ri javob:} \textbf{A) Mos qurilma drayverini o'rnatish}
\par\smallskip
{\bfseries Izoh:} Operatsion tizim qurilma bilan ishlashi uchun mos drayver kerak bo'lishi mumkin.
\end{tcolorbox}
\vspace{2mm}\hrule height 0.2pt\vspace{3mm}

% DB_ID 5f9fb965-e051-41c1-8f19-33fe2626f188
\needspace{5\baselineskip}
{\bfseries\sffamily\color{ink} 296.} {scriptsizecolor{muted} Y1 | published | 5f9fb965}
\par
Administrator serverni grafik oynalarsiz tezkor matnli buyruqlar orqali boshqarmoqchi. Qaysi interfeys mos?
\begin{enumerate}[label=\Alph*),leftmargin=9mm]
\item CLI
\item GUI
\item Slayd-shou
\item WYSIWYG muharriri
\end{enumerate}
\begin{tcolorbox}[enhanced,breakable,colback=teal!4,colframe=teal!60!black,boxrule=0.45pt,arc=1.2mm,left=2.7mm,right=2.7mm,top=1.8mm,bottom=1.8mm,before skip=2mm,after skip=2mm]
{\bfseries\sffamily\color{teal}To'g'ri javob:} \textbf{A) CLI}
\par\smallskip
{\bfseries Izoh:} Matnli buyruqlar bilan boshqarish CLI ga xos.
\end{tcolorbox}
\vspace{2mm}\hrule height 0.2pt\vspace{3mm}

% DB_ID b9f11db4-faec-461e-af5a-8f1c1d08c9ca
\needspace{5\baselineskip}
{\bfseries\sffamily\color{ink} 297.} {scriptsizecolor{muted} Y1 | published | b9f11db4}
\par
Faqat operatsion tizimlar keltirilgan qatorni belgilang.
\begin{enumerate}[label=\Alph*),leftmargin=9mm]
\item IBM, Windows
\item Linux, CMS
\item iOS, Moodle
\item Windows, Android
\end{enumerate}
\begin{tcolorbox}[enhanced,breakable,colback=teal!4,colframe=teal!60!black,boxrule=0.45pt,arc=1.2mm,left=2.7mm,right=2.7mm,top=1.8mm,bottom=1.8mm,before skip=2mm,after skip=2mm]
{\bfseries\sffamily\color{teal}To'g'ri javob:} \textbf{D) Windows, Android}
\par\smallskip
{\bfseries Izoh:} Windows va Android ikkalasi operatsion tizim. [2025 kuzgi demo YECHIMLARI, sahifa 10.]
\end{tcolorbox}
\vspace{2mm}\hrule height 0.2pt\vspace{3mm}
\subsection{Tekshiruvdagi testlar}

% DB_ID 0638433b-07f6-6355-34cd-a65ccaa069d6
\needspace{5\baselineskip}
{\bfseries\sffamily\color{ink} 298.} {scriptsizecolor{muted} Y1 | review | 0638433b}
\par
Operatsion tizim turlari va interfeyslar doirasida "real vaqt OT" atamasining eng aniq ta'rifini toping.
\begin{enumerate}[label=\Alph*),leftmargin=9mm]
\item oyna, ikona, menyu va ko'rsatkichga tayangan grafik interfeys
\item matnli buyruqlar orqali boshqariladigan interfeys
\item tarmoq xizmatlari va ko'p mijoz resurslarini boshqarishga mo'ljallangan OT
\item javobni belgilangan vaqt chegarasida berishni kafolatlashga yo'naltirilgan OT
\end{enumerate}
\begin{tcolorbox}[enhanced,breakable,colback=teal!4,colframe=teal!60!black,boxrule=0.45pt,arc=1.2mm,left=2.7mm,right=2.7mm,top=1.8mm,bottom=1.8mm,before skip=2mm,after skip=2mm]
{\bfseries\sffamily\color{teal}To'g'ri javob:} \textbf{D) javobni belgilangan vaqt chegarasida berishni kafolatlashga yo'naltirilgan OT}
\par\smallskip
{\bfseries Izoh:} javobni belgilangan vaqt chegarasida berishni kafolatlashga yo'naltirilgan OT. real vaqt OT - javobni belgilangan vaqt chegarasida berishni kafolatlashga yo'naltirilgan OT. Manba: MAR, PDF 12-15; Qo'shimcha sahifa raqami (qaysi PDFga tegishli ekani aniqlanmagan; darslik sahifasi deb tasdiqlanmagan): 67
\end{tcolorbox}
\vspace{2mm}\hrule height 0.2pt\vspace{3mm}

% DB_ID 2e127eeb-dbfa-1d06-8c0c-1c92ea9d465f
\needspace{5\baselineskip}
{\bfseries\sffamily\color{ink} 299.} {scriptsizecolor{muted} Y1 | review | 2e127eeb}
\par
Operatsion tizim turlari va interfeyslar doirasida "CLI" atamasining eng aniq ta'rifini toping.
\begin{enumerate}[label=\Alph*),leftmargin=9mm]
\item oyna, ikona, menyu va ko'rsatkichga tayangan grafik interfeys
\item tarmoq xizmatlari va ko'p mijoz resurslarini boshqarishga mo'ljallangan OT
\item sensorli, portativ va energiya cheklangan qurilmaga mos OT
\item matnli buyruqlar orqali boshqariladigan interfeys
\end{enumerate}
\begin{tcolorbox}[enhanced,breakable,colback=teal!4,colframe=teal!60!black,boxrule=0.45pt,arc=1.2mm,left=2.7mm,right=2.7mm,top=1.8mm,bottom=1.8mm,before skip=2mm,after skip=2mm]
{\bfseries\sffamily\color{teal}To'g'ri javob:} \textbf{D) matnli buyruqlar orqali boshqariladigan interfeys}
\par\smallskip
{\bfseries Izoh:} matnli buyruqlar orqali boshqariladigan interfeys. CLI - matnli buyruqlar orqali boshqariladigan interfeys. Manba: CAM10/11, PDF 44-48; Qo'shimcha sahifa raqami (qaysi PDFga tegishli ekani aniqlanmagan; darslik sahifasi deb tasdiqlanmagan): 63
\end{tcolorbox}
\vspace{2mm}\hrule height 0.2pt\vspace{3mm}

% DB_ID 629a1742-99c0-00f7-d680-26e28ad6616b
\needspace{5\baselineskip}
{\bfseries\sffamily\color{ink} 300.} {scriptsizecolor{muted} Y1 | review | 629a1742}
\par
Operatsion tizimning vazifalari doirasida "yadro" atamasining eng aniq ta'rifini toping.
\begin{enumerate}[label=\Alph*),leftmargin=9mm]
\item OTning resurslar va apparatni boshqaruvchi markaziy qismi
\item resurslarni boshqarib, foydalanuvchi va ilovalarga xizmat ko'rsatuvchi tizimli dastur
\item bajarilayotgan dastur nusxasi
\item bir nechta vazifani navbatlashtirib yoki parallel bajarish
\end{enumerate}
\begin{tcolorbox}[enhanced,breakable,colback=teal!4,colframe=teal!60!black,boxrule=0.45pt,arc=1.2mm,left=2.7mm,right=2.7mm,top=1.8mm,bottom=1.8mm,before skip=2mm,after skip=2mm]
{\bfseries\sffamily\color{teal}To'g'ri javob:} \textbf{A) OTning resurslar va apparatni boshqaruvchi markaziy qismi}
\par\smallskip
{\bfseries Izoh:} OTning resurslar va apparatni boshqaruvchi markaziy qismi. yadro - OTning resurslar va apparatni boshqaruvchi markaziy qismi. Manba: MAR, PDF 11-16; Qo'shimcha sahifa raqami (qaysi PDFga tegishli ekani aniqlanmagan; darslik sahifasi deb tasdiqlanmagan): 60
\end{tcolorbox}
\vspace{2mm}\hrule height 0.2pt\vspace{3mm}

% DB_ID 7cf166e4-7402-effc-40c5-7cc6214f4c83
\needspace{5\baselineskip}
{\bfseries\sffamily\color{ink} 301.} {scriptsizecolor{muted} Y1 | review | 7cf166e4}
\par
Operatsion tizimning vazifalari doirasida "operatsion tizim" atamasining eng aniq ta'rifini toping.
\begin{enumerate}[label=\Alph*),leftmargin=9mm]
\item resurslarni boshqarib, foydalanuvchi va ilovalarga xizmat ko'rsatuvchi tizimli dastur
\item bajarilayotgan dastur nusxasi
\item bir nechta vazifani navbatlashtirib yoki parallel bajarish
\item fayllarni nomlash, joylashtirish va boshqarish qoidalari
\end{enumerate}
\begin{tcolorbox}[enhanced,breakable,colback=teal!4,colframe=teal!60!black,boxrule=0.45pt,arc=1.2mm,left=2.7mm,right=2.7mm,top=1.8mm,bottom=1.8mm,before skip=2mm,after skip=2mm]
{\bfseries\sffamily\color{teal}To'g'ri javob:} \textbf{A) resurslarni boshqarib, foydalanuvchi va ilovalarga xizmat ko'rsatuvchi tizimli dastur}
\par\smallskip
{\bfseries Izoh:} resurslarni boshqarib, foydalanuvchi va ilovalarga xizmat ko'rsatuvchi tizimli dastur. operatsion tizim - resurslarni boshqarib, foydalanuvchi va ilovalarga xizmat ko'rsatuvchi tizimli dastur. Manba: CAM9, PDF 12-15; Qo'shimcha sahifa raqami (qaysi PDFga tegishli ekani aniqlanmagan; darslik sahifasi deb tasdiqlanmagan): 56
\end{tcolorbox}
\vspace{2mm}\hrule height 0.2pt\vspace{3mm}

% DB_ID 7dbaa4c2-7ba4-118e-2869-9c0c2368d6f7
\needspace{5\baselineskip}
{\bfseries\sffamily\color{ink} 302.} {scriptsizecolor{muted} Y1 | review | 7dbaa4c2}
\par
Operatsion tizimning vazifalari doirasida "multitasking" atamasining eng aniq ta'rifini toping.
\begin{enumerate}[label=\Alph*),leftmargin=9mm]
\item resurslarni boshqarib, foydalanuvchi va ilovalarga xizmat ko'rsatuvchi tizimli dastur
\item bajarilayotgan dastur nusxasi
\item bir nechta vazifani navbatlashtirib yoki parallel bajarish
\item fayllarni nomlash, joylashtirish va boshqarish qoidalari
\end{enumerate}
\begin{tcolorbox}[enhanced,breakable,colback=teal!4,colframe=teal!60!black,boxrule=0.45pt,arc=1.2mm,left=2.7mm,right=2.7mm,top=1.8mm,bottom=1.8mm,before skip=2mm,after skip=2mm]
{\bfseries\sffamily\color{teal}To'g'ri javob:} \textbf{C) bir nechta vazifani navbatlashtirib yoki parallel bajarish}
\par\smallskip
{\bfseries Izoh:} bir nechta vazifani navbatlashtirib yoki parallel bajarish. multitasking - bir nechta vazifani navbatlashtirib yoki parallel bajarish. Manba: ICT5, PDF 38-41; Qo'shimcha sahifa raqami (qaysi PDFga tegishli ekani aniqlanmagan; darslik sahifasi deb tasdiqlanmagan): 58
\end{tcolorbox}
\vspace{2mm}\hrule height 0.2pt\vspace{3mm}

% DB_ID 8f2139f8-1446-da5b-bd83-cce5a125d522
\needspace{5\baselineskip}
{\bfseries\sffamily\color{ink} 303.} {scriptsizecolor{muted} Y1 | review | 8f2139f8}
\par
Operatsion tizimning vazifalari doirasida "jarayon" atamasining eng aniq ta'rifini toping.
\begin{enumerate}[label=\Alph*),leftmargin=9mm]
\item resurslarni boshqarib, foydalanuvchi va ilovalarga xizmat ko'rsatuvchi tizimli dastur
\item bajarilayotgan dastur nusxasi
\item bir nechta vazifani navbatlashtirib yoki parallel bajarish
\item fayllarni nomlash, joylashtirish va boshqarish qoidalari
\end{enumerate}
\begin{tcolorbox}[enhanced,breakable,colback=teal!4,colframe=teal!60!black,boxrule=0.45pt,arc=1.2mm,left=2.7mm,right=2.7mm,top=1.8mm,bottom=1.8mm,before skip=2mm,after skip=2mm]
{\bfseries\sffamily\color{teal}To'g'ri javob:} \textbf{B) bajarilayotgan dastur nusxasi}
\par\smallskip
{\bfseries Izoh:} bajarilayotgan dastur nusxasi. jarayon - bajarilayotgan dastur nusxasi. Manba: CAM10/11, PDF 43-47; Qo'shimcha sahifa raqami (qaysi PDFga tegishli ekani aniqlanmagan; darslik sahifasi deb tasdiqlanmagan): 57
\end{tcolorbox}
\vspace{2mm}\hrule height 0.2pt\vspace{3mm}

% DB_ID 8f950b26-4536-b06a-7c9e-527aa40bbbf1
\needspace{5\baselineskip}
{\bfseries\sffamily\color{ink} 304.} {scriptsizecolor{muted} Y1 | review | 8f950b26}
\par
Operatsion tizimning vazifalari doirasida "fayl tizimi" atamasining eng aniq ta'rifini toping.
\begin{enumerate}[label=\Alph*),leftmargin=9mm]
\item resurslarni boshqarib, foydalanuvchi va ilovalarga xizmat ko'rsatuvchi tizimli dastur
\item bajarilayotgan dastur nusxasi
\item bir nechta vazifani navbatlashtirib yoki parallel bajarish
\item fayllarni nomlash, joylashtirish va boshqarish qoidalari
\end{enumerate}
\begin{tcolorbox}[enhanced,breakable,colback=teal!4,colframe=teal!60!black,boxrule=0.45pt,arc=1.2mm,left=2.7mm,right=2.7mm,top=1.8mm,bottom=1.8mm,before skip=2mm,after skip=2mm]
{\bfseries\sffamily\color{teal}To'g'ri javob:} \textbf{D) fayllarni nomlash, joylashtirish va boshqarish qoidalari}
\par\smallskip
{\bfseries Izoh:} fayllarni nomlash, joylashtirish va boshqarish qoidalari. fayl tizimi - fayllarni nomlash, joylashtirish va boshqarish qoidalari. Manba: CAM8, PDF 99-106; Qo'shimcha sahifa raqami (qaysi PDFga tegishli ekani aniqlanmagan; darslik sahifasi deb tasdiqlanmagan): 59
\end{tcolorbox}
\vspace{2mm}\hrule height 0.2pt\vspace{3mm}

% DB_ID 9ace4589-cfb7-a191-e2fb-8d8b7ba547b7
\needspace{5\baselineskip}
{\bfseries\sffamily\color{ink} 305.} {scriptsizecolor{muted} Y1 | review | 9ace4589}
\par
Operatsion tizim turlari va interfeyslar doirasida "server OT" atamasining eng aniq ta'rifini toping.
\begin{enumerate}[label=\Alph*),leftmargin=9mm]
\item tarmoq xizmatlari va ko'p mijoz resurslarini boshqarishga mo'ljallangan OT
\item oyna, ikona, menyu va ko'rsatkichga tayangan grafik interfeys
\item matnli buyruqlar orqali boshqariladigan interfeys
\item sensorli, portativ va energiya cheklangan qurilmaga mos OT
\end{enumerate}
\begin{tcolorbox}[enhanced,breakable,colback=teal!4,colframe=teal!60!black,boxrule=0.45pt,arc=1.2mm,left=2.7mm,right=2.7mm,top=1.8mm,bottom=1.8mm,before skip=2mm,after skip=2mm]
{\bfseries\sffamily\color{teal}To'g'ri javob:} \textbf{A) tarmoq xizmatlari va ko'p mijoz resurslarini boshqarishga mo'ljallangan OT}
\par\smallskip
{\bfseries Izoh:} tarmoq xizmatlari va ko'p mijoz resurslarini boshqarishga mo'ljallangan OT. server OT - tarmoq xizmatlari va ko'p mijoz resurslarini boshqarishga mo'ljallangan OT. Manba: ICT5, PDF 39-41; Qo'shimcha sahifa raqami (qaysi PDFga tegishli ekani aniqlanmagan; darslik sahifasi deb tasdiqlanmagan): 64
\end{tcolorbox}
\vspace{2mm}\hrule height 0.2pt\vspace{3mm}

% DB_ID beaae02d-a6de-f144-173c-67dbe54a4558
\needspace{5\baselineskip}
{\bfseries\sffamily\color{ink} 306.} {scriptsizecolor{muted} Y1 | review | beaae02d}
\par
Operatsion tizim turlari va interfeyslar doirasida "mobil OT" atamasining eng aniq ta'rifini toping.
\begin{enumerate}[label=\Alph*),leftmargin=9mm]
\item oyna, ikona, menyu va ko'rsatkichga tayangan grafik interfeys
\item sensorli, portativ va energiya cheklangan qurilmaga mos OT
\item matnli buyruqlar orqali boshqariladigan interfeys
\item tarmoq xizmatlari va ko'p mijoz resurslarini boshqarishga mo'ljallangan OT
\end{enumerate}
\begin{tcolorbox}[enhanced,breakable,colback=teal!4,colframe=teal!60!black,boxrule=0.45pt,arc=1.2mm,left=2.7mm,right=2.7mm,top=1.8mm,bottom=1.8mm,before skip=2mm,after skip=2mm]
{\bfseries\sffamily\color{teal}To'g'ri javob:} \textbf{B) sensorli, portativ va energiya cheklangan qurilmaga mos OT}
\par\smallskip
{\bfseries Izoh:} sensorli, portativ va energiya cheklangan qurilmaga mos OT. mobil OT - sensorli, portativ va energiya cheklangan qurilmaga mos OT. Manba: CAM7, PDF 140-146; Qo'shimcha sahifa raqami (qaysi PDFga tegishli ekani aniqlanmagan; darslik sahifasi deb tasdiqlanmagan): 65
\end{tcolorbox}
\vspace{2mm}\hrule height 0.2pt\vspace{3mm}

% DB_ID d9353796-f25e-64ae-909b-6a009e429eb8
\needspace{5\baselineskip}
{\bfseries\sffamily\color{ink} 307.} {scriptsizecolor{muted} Y1 | review | d9353796}
\par
Operatsion tizim turlari va interfeyslar doirasida "embedded OT" atamasining eng aniq ta'rifini toping.
\begin{enumerate}[label=\Alph*),leftmargin=9mm]
\item oyna, ikona, menyu va ko'rsatkichga tayangan grafik interfeys
\item matnli buyruqlar orqali boshqariladigan interfeys
\item maxsus qurilma vazifasiga moslashtirilgan OT
\item tarmoq xizmatlari va ko'p mijoz resurslarini boshqarishga mo'ljallangan OT
\end{enumerate}
\begin{tcolorbox}[enhanced,breakable,colback=teal!4,colframe=teal!60!black,boxrule=0.45pt,arc=1.2mm,left=2.7mm,right=2.7mm,top=1.8mm,bottom=1.8mm,before skip=2mm,after skip=2mm]
{\bfseries\sffamily\color{teal}To'g'ri javob:} \textbf{C) maxsus qurilma vazifasiga moslashtirilgan OT}
\par\smallskip
{\bfseries Izoh:} maxsus qurilma vazifasiga moslashtirilgan OT. embedded OT - maxsus qurilma vazifasiga moslashtirilgan OT. Manba: ICT11, PDF 244-245; Qo'shimcha sahifa raqami (qaysi PDFga tegishli ekani aniqlanmagan; darslik sahifasi deb tasdiqlanmagan): 66
\end{tcolorbox}
\vspace{2mm}\hrule height 0.2pt\vspace{3mm}

% DB_ID e8db1091-40ba-769a-a709-9c3fa0b63820
\needspace{5\baselineskip}
{\bfseries\sffamily\color{ink} 308.} {scriptsizecolor{muted} Y1 | review | e8db1091}
\par
Operatsion tizimning vazifalari doirasida "bootloader" atamasining eng aniq ta'rifini toping.
\begin{enumerate}[label=\Alph*),leftmargin=9mm]
\item resurslarni boshqarib, foydalanuvchi va ilovalarga xizmat ko'rsatuvchi tizimli dastur
\item operatsion tizim yadrosini xotiraga yuklovchi dastur
\item bajarilayotgan dastur nusxasi
\item bir nechta vazifani navbatlashtirib yoki parallel bajarish
\end{enumerate}
\begin{tcolorbox}[enhanced,breakable,colback=teal!4,colframe=teal!60!black,boxrule=0.45pt,arc=1.2mm,left=2.7mm,right=2.7mm,top=1.8mm,bottom=1.8mm,before skip=2mm,after skip=2mm]
{\bfseries\sffamily\color{teal}To'g'ri javob:} \textbf{B) operatsion tizim yadrosini xotiraga yuklovchi dastur}
\par\smallskip
{\bfseries Izoh:} operatsion tizim yadrosini xotiraga yuklovchi dastur. bootloader - operatsion tizim yadrosini xotiraga yuklovchi dastur. Manba: CAM9, PDF 12-15; Qo'shimcha sahifa raqami (qaysi PDFga tegishli ekani aniqlanmagan; darslik sahifasi deb tasdiqlanmagan): 61
\end{tcolorbox}
\vspace{2mm}\hrule height 0.2pt\vspace{3mm}

% DB_ID edc13eec-f4a5-6700-9e7d-2e6f815bacc9
\needspace{5\baselineskip}
{\bfseries\sffamily\color{ink} 309.} {scriptsizecolor{muted} Y1 | review | edc13eec}
\par
Operatsion tizim turlari va interfeyslar doirasida "GUI" atamasining eng aniq ta'rifini toping.
\begin{enumerate}[label=\Alph*),leftmargin=9mm]
\item matnli buyruqlar orqali boshqariladigan interfeys
\item tarmoq xizmatlari va ko'p mijoz resurslarini boshqarishga mo'ljallangan OT
\item oyna, ikona, menyu va ko'rsatkichga tayangan grafik interfeys
\item sensorli, portativ va energiya cheklangan qurilmaga mos OT
\end{enumerate}
\begin{tcolorbox}[enhanced,breakable,colback=teal!4,colframe=teal!60!black,boxrule=0.45pt,arc=1.2mm,left=2.7mm,right=2.7mm,top=1.8mm,bottom=1.8mm,before skip=2mm,after skip=2mm]
{\bfseries\sffamily\color{teal}To'g'ri javob:} \textbf{C) oyna, ikona, menyu va ko'rsatkichga tayangan grafik interfeys}
\par\smallskip
{\bfseries Izoh:} oyna, ikona, menyu va ko'rsatkichga tayangan grafik interfeys. GUI - oyna, ikona, menyu va ko'rsatkichga tayangan grafik interfeys. Manba: CAM9, PDF 13-15; Qo'shimcha sahifa raqami (qaysi PDFga tegishli ekani aniqlanmagan; darslik sahifasi deb tasdiqlanmagan): 62
\end{tcolorbox}
\vspace{2mm}\hrule height 0.2pt\vspace{3mm}

% DB_ID 0e3b2815-b2e3-e8b6-f8ea-64d216b94a35
\needspace{5\baselineskip}
{\bfseries\sffamily\color{ink} 310.} {scriptsizecolor{muted} Y1 | review | 0e3b2815}
\par
Operatsion tizim turlari va interfeyslar mavzusidagi vazifada javobni belgilangan vaqt chegarasida berishni kafolatlashga yo'naltirilgan OT talab qilindi. Qaysi tushuncha mos?
\begin{enumerate}[label=\Alph*),leftmargin=9mm]
\item real vaqt OT
\item GUI
\item CLI
\item server OT
\end{enumerate}
\begin{tcolorbox}[enhanced,breakable,colback=teal!4,colframe=teal!60!black,boxrule=0.45pt,arc=1.2mm,left=2.7mm,right=2.7mm,top=1.8mm,bottom=1.8mm,before skip=2mm,after skip=2mm]
{\bfseries\sffamily\color{teal}To'g'ri javob:} \textbf{A) real vaqt OT}
\par\smallskip
{\bfseries Izoh:} real vaqt OT. Berilgan vazifa real vaqt OT tushunchasiga mos: javobni belgilangan vaqt chegarasida berishni kafolatlashga yo'naltirilgan OT. Manba: MAR, PDF 12-15; Qo'shimcha sahifa raqami (qaysi PDFga tegishli ekani aniqlanmagan; darslik sahifasi deb tasdiqlanmagan): 68
\end{tcolorbox}
\vspace{2mm}\hrule height 0.2pt\vspace{3mm}

% DB_ID 1a9586c0-243f-1f6c-ed72-57348c0e913c
\needspace{5\baselineskip}
{\bfseries\sffamily\color{ink} 311.} {scriptsizecolor{muted} Y1 | review | 1a9586c0}
\par
Operatsion tizimning vazifalari mavzusida "yadro" qatnashgan atama-ta'rif juftliklaridan to'g'risini toping.
\begin{enumerate}[label=\Alph*),leftmargin=9mm]
\item yadro - resurslarni boshqarib, foydalanuvchi va ilovalarga xizmat ko'rsatuvchi tizimli dastur
\item yadro - bajarilayotgan dastur nusxasi
\item yadro - OTning resurslar va apparatni boshqaruvchi markaziy qismi
\item yadro - bir nechta vazifani navbatlashtirib yoki parallel bajarish
\end{enumerate}
\begin{tcolorbox}[enhanced,breakable,colback=teal!4,colframe=teal!60!black,boxrule=0.45pt,arc=1.2mm,left=2.7mm,right=2.7mm,top=1.8mm,bottom=1.8mm,before skip=2mm,after skip=2mm]
{\bfseries\sffamily\color{teal}To'g'ri javob:} \textbf{C) yadro - OTning resurslar va apparatni boshqaruvchi markaziy qismi}
\par\smallskip
{\bfseries Izoh:} yadro - OTning resurslar va apparatni boshqaruvchi markaziy qismi. To'g'ri juftlik: yadro - OTning resurslar va apparatni boshqaruvchi markaziy qismi. Manba: MAR, PDF 11-16; Qo'shimcha sahifa raqami (qaysi PDFga tegishli ekani aniqlanmagan; darslik sahifasi deb tasdiqlanmagan): 60
\end{tcolorbox}
\vspace{2mm}\hrule height 0.2pt\vspace{3mm}

% DB_ID 24d1e209-f1ce-9967-b56a-e8c8572d00d9
\needspace{5\baselineskip}
{\bfseries\sffamily\color{ink} 312.} {scriptsizecolor{muted} Y1 | review | 24d1e209}
\par
Operatsion tizim turlari va interfeyslar mavzusidagi vazifada oyna, ikona, menyu va ko'rsatkichga tayangan grafik interfeys talab qilindi. Qaysi tushuncha mos?
\begin{enumerate}[label=\Alph*),leftmargin=9mm]
\item CLI
\item server OT
\item mobil OT
\item GUI
\end{enumerate}
\begin{tcolorbox}[enhanced,breakable,colback=teal!4,colframe=teal!60!black,boxrule=0.45pt,arc=1.2mm,left=2.7mm,right=2.7mm,top=1.8mm,bottom=1.8mm,before skip=2mm,after skip=2mm]
{\bfseries\sffamily\color{teal}To'g'ri javob:} \textbf{D) GUI}
\par\smallskip
{\bfseries Izoh:} GUI. Berilgan vazifa GUI tushunchasiga mos: oyna, ikona, menyu va ko'rsatkichga tayangan grafik interfeys. Manba: CAM9, PDF 13-15; Qo'shimcha sahifa raqami (qaysi PDFga tegishli ekani aniqlanmagan; darslik sahifasi deb tasdiqlanmagan): 62
\end{tcolorbox}
\vspace{2mm}\hrule height 0.2pt\vspace{3mm}

% DB_ID 25cf632e-59b0-4d85-77f7-9ca0ace50f3c
\needspace{5\baselineskip}
{\bfseries\sffamily\color{ink} 313.} {scriptsizecolor{muted} Y1 | review | 25cf632e}
\par
Operatsion tizim turlari va interfeyslar mavzusida "embedded OT" qatnashgan atama-ta'rif juftliklaridan to'g'risini toping.
\begin{enumerate}[label=\Alph*),leftmargin=9mm]
\item embedded OT - maxsus qurilma vazifasiga moslashtirilgan OT
\item embedded OT - oyna, ikona, menyu va ko'rsatkichga tayangan grafik interfeys
\item embedded OT - matnli buyruqlar orqali boshqariladigan interfeys
\item embedded OT - tarmoq xizmatlari va ko'p mijoz resurslarini boshqarishga mo'ljallangan OT
\end{enumerate}
\begin{tcolorbox}[enhanced,breakable,colback=teal!4,colframe=teal!60!black,boxrule=0.45pt,arc=1.2mm,left=2.7mm,right=2.7mm,top=1.8mm,bottom=1.8mm,before skip=2mm,after skip=2mm]
{\bfseries\sffamily\color{teal}To'g'ri javob:} \textbf{A) embedded OT - maxsus qurilma vazifasiga moslashtirilgan OT}
\par\smallskip
{\bfseries Izoh:} embedded OT - maxsus qurilma vazifasiga moslashtirilgan OT. To'g'ri juftlik: embedded OT - maxsus qurilma vazifasiga moslashtirilgan OT. Manba: ICT11, PDF 244-245; Qo'shimcha sahifa raqami (qaysi PDFga tegishli ekani aniqlanmagan; darslik sahifasi deb tasdiqlanmagan): 67
\end{tcolorbox}
\vspace{2mm}\hrule height 0.2pt\vspace{3mm}

% DB_ID 2cb799ba-d0f6-1aec-5c4b-81ec3b813937
\needspace{5\baselineskip}
{\bfseries\sffamily\color{ink} 314.} {scriptsizecolor{muted} Y1 | review | 2cb799ba}
\par
Operatsion tizim turlari va interfeyslar mavzusidagi vazifada maxsus qurilma vazifasiga moslashtirilgan OT talab qilindi. Qaysi tushuncha mos?
\begin{enumerate}[label=\Alph*),leftmargin=9mm]
\item GUI
\item CLI
\item server OT
\item embedded OT
\end{enumerate}
\begin{tcolorbox}[enhanced,breakable,colback=teal!4,colframe=teal!60!black,boxrule=0.45pt,arc=1.2mm,left=2.7mm,right=2.7mm,top=1.8mm,bottom=1.8mm,before skip=2mm,after skip=2mm]
{\bfseries\sffamily\color{teal}To'g'ri javob:} \textbf{D) embedded OT}
\par\smallskip
{\bfseries Izoh:} embedded OT. Berilgan vazifa embedded OT tushunchasiga mos: maxsus qurilma vazifasiga moslashtirilgan OT. Manba: ICT11, PDF 244-245; Qo'shimcha sahifa raqami (qaysi PDFga tegishli ekani aniqlanmagan; darslik sahifasi deb tasdiqlanmagan): 66
\end{tcolorbox}
\vspace{2mm}\hrule height 0.2pt\vspace{3mm}

% DB_ID 33186aee-3e1e-bbbe-e0ad-52fc03b9ba7e
\needspace{5\baselineskip}
{\bfseries\sffamily\color{ink} 315.} {scriptsizecolor{muted} Y1 | review | 33186aee}
\par
Operatsion tizim turlari va interfeyslar mavzusidagi vazifada sensorli, portativ va energiya cheklangan qurilmaga mos OT talab qilindi. Qaysi tushuncha mos?
\begin{enumerate}[label=\Alph*),leftmargin=9mm]
\item GUI
\item CLI
\item mobil OT
\item server OT
\end{enumerate}
\begin{tcolorbox}[enhanced,breakable,colback=teal!4,colframe=teal!60!black,boxrule=0.45pt,arc=1.2mm,left=2.7mm,right=2.7mm,top=1.8mm,bottom=1.8mm,before skip=2mm,after skip=2mm]
{\bfseries\sffamily\color{teal}To'g'ri javob:} \textbf{C) mobil OT}
\par\smallskip
{\bfseries Izoh:} mobil OT. Berilgan vazifa mobil OT tushunchasiga mos: sensorli, portativ va energiya cheklangan qurilmaga mos OT. Manba: CAM7, PDF 140-146; Qo'shimcha sahifa raqami (qaysi PDFga tegishli ekani aniqlanmagan; darslik sahifasi deb tasdiqlanmagan): 65
\end{tcolorbox}
\vspace{2mm}\hrule height 0.2pt\vspace{3mm}

% DB_ID 4084ade4-3814-ffcc-171a-f629e92baef3
\needspace{5\baselineskip}
{\bfseries\sffamily\color{ink} 316.} {scriptsizecolor{muted} Y1 | review | 4084ade4}
\par
Operatsion tizimning vazifalari mavzusidagi vazifada resurslarni boshqarib, foydalanuvchi va ilovalarga xizmat ko'rsatuvchi tizimli dastur talab qilindi. Qaysi tushuncha mos?
\begin{enumerate}[label=\Alph*),leftmargin=9mm]
\item jarayon
\item operatsion tizim
\item multitasking
\item fayl tizimi
\end{enumerate}
\begin{tcolorbox}[enhanced,breakable,colback=teal!4,colframe=teal!60!black,boxrule=0.45pt,arc=1.2mm,left=2.7mm,right=2.7mm,top=1.8mm,bottom=1.8mm,before skip=2mm,after skip=2mm]
{\bfseries\sffamily\color{teal}To'g'ri javob:} \textbf{B) operatsion tizim}
\par\smallskip
{\bfseries Izoh:} operatsion tizim. Berilgan vazifa operatsion tizim tushunchasiga mos: resurslarni boshqarib, foydalanuvchi va ilovalarga xizmat ko'rsatuvchi tizimli dastur. Manba: CAM9, PDF 12-15; Qo'shimcha sahifa raqami (qaysi PDFga tegishli ekani aniqlanmagan; darslik sahifasi deb tasdiqlanmagan): 56
\end{tcolorbox}
\vspace{2mm}\hrule height 0.2pt\vspace{3mm}

% DB_ID 6f18921e-902f-6d36-78b0-c333d326fcf1
\needspace{5\baselineskip}
{\bfseries\sffamily\color{ink} 317.} {scriptsizecolor{muted} Y1 | review | 6f18921e}
\par
Operatsion tizimning vazifalari mavzusidagi vazifada fayllarni nomlash, joylashtirish va boshqarish qoidalari talab qilindi. Qaysi tushuncha mos?
\begin{enumerate}[label=\Alph*),leftmargin=9mm]
\item fayl tizimi
\item operatsion tizim
\item jarayon
\item multitasking
\end{enumerate}
\begin{tcolorbox}[enhanced,breakable,colback=teal!4,colframe=teal!60!black,boxrule=0.45pt,arc=1.2mm,left=2.7mm,right=2.7mm,top=1.8mm,bottom=1.8mm,before skip=2mm,after skip=2mm]
{\bfseries\sffamily\color{teal}To'g'ri javob:} \textbf{A) fayl tizimi}
\par\smallskip
{\bfseries Izoh:} fayl tizimi. Berilgan vazifa fayl tizimi tushunchasiga mos: fayllarni nomlash, joylashtirish va boshqarish qoidalari. Manba: CAM8, PDF 99-106; Qo'shimcha sahifa raqami (qaysi PDFga tegishli ekani aniqlanmagan; darslik sahifasi deb tasdiqlanmagan): 59
\end{tcolorbox}
\vspace{2mm}\hrule height 0.2pt\vspace{3mm}

% DB_ID 7d4102e1-8acc-bc65-137a-c49bcf1425a6
\needspace{5\baselineskip}
{\bfseries\sffamily\color{ink} 318.} {scriptsizecolor{muted} Y1 | review | 7d4102e1}
\par
Operatsion tizim turlari va interfeyslar mavzusidagi vazifada matnli buyruqlar orqali boshqariladigan interfeys talab qilindi. Qaysi tushuncha mos?
\begin{enumerate}[label=\Alph*),leftmargin=9mm]
\item CLI
\item GUI
\item server OT
\item mobil OT
\end{enumerate}
\begin{tcolorbox}[enhanced,breakable,colback=teal!4,colframe=teal!60!black,boxrule=0.45pt,arc=1.2mm,left=2.7mm,right=2.7mm,top=1.8mm,bottom=1.8mm,before skip=2mm,after skip=2mm]
{\bfseries\sffamily\color{teal}To'g'ri javob:} \textbf{A) CLI}
\par\smallskip
{\bfseries Izoh:} CLI. Berilgan vazifa CLI tushunchasiga mos: matnli buyruqlar orqali boshqariladigan interfeys. Manba: CAM10/11, PDF 44-48; Qo'shimcha sahifa raqami (qaysi PDFga tegishli ekani aniqlanmagan; darslik sahifasi deb tasdiqlanmagan): 63
\end{tcolorbox}
\vspace{2mm}\hrule height 0.2pt\vspace{3mm}

% DB_ID 7d9ccf88-02ed-0a84-4631-001a5b733a0c
\needspace{5\baselineskip}
{\bfseries\sffamily\color{ink} 319.} {scriptsizecolor{muted} Y1 | review | 7d9ccf88}
\par
Operatsion tizimning vazifalari mavzusida "operatsion tizim" qatnashgan atama-ta'rif juftliklaridan to'g'risini toping.
\begin{enumerate}[label=\Alph*),leftmargin=9mm]
\item operatsion tizim - bajarilayotgan dastur nusxasi
\item operatsion tizim - bir nechta vazifani navbatlashtirib yoki parallel bajarish
\item operatsion tizim - resurslarni boshqarib, foydalanuvchi va ilovalarga xizmat ko'rsatuvchi tizimli dastur
\item operatsion tizim - fayllarni nomlash, joylashtirish va boshqarish qoidalari
\end{enumerate}
\begin{tcolorbox}[enhanced,breakable,colback=teal!4,colframe=teal!60!black,boxrule=0.45pt,arc=1.2mm,left=2.7mm,right=2.7mm,top=1.8mm,bottom=1.8mm,before skip=2mm,after skip=2mm]
{\bfseries\sffamily\color{teal}To'g'ri javob:} \textbf{C) operatsion tizim - resurslarni boshqarib, foydalanuvchi va ilovalarga xizmat ko'rsatuvchi tizimli dastur}
\par\smallskip
{\bfseries Izoh:} operatsion tizim - resurslarni boshqarib, foydalanuvchi va ilovalarga xizmat ko'rsatuvchi tizimli dastur. To'g'ri juftlik: operatsion tizim - resurslarni boshqarib, foydalanuvchi va ilovalarga xizmat ko'rsatuvchi tizimli dastur. Manba: CAM9, PDF 12-15; Qo'shimcha sahifa raqami (qaysi PDFga tegishli ekani aniqlanmagan; darslik sahifasi deb tasdiqlanmagan): 56
\end{tcolorbox}
\vspace{2mm}\hrule height 0.2pt\vspace{3mm}

% DB_ID 82d926d4-40bf-6be0-47b2-70c4760efb17
\needspace{5\baselineskip}
{\bfseries\sffamily\color{ink} 320.} {scriptsizecolor{muted} Y1 | review | 82d926d4}
\par
Operatsion tizim turlari va interfeyslar mavzusida "server OT" qatnashgan atama-ta'rif juftliklaridan to'g'risini toping.
\begin{enumerate}[label=\Alph*),leftmargin=9mm]
\item server OT - oyna, ikona, menyu va ko'rsatkichga tayangan grafik interfeys
\item server OT - matnli buyruqlar orqali boshqariladigan interfeys
\item server OT - tarmoq xizmatlari va ko'p mijoz resurslarini boshqarishga mo'ljallangan OT
\item server OT - sensorli, portativ va energiya cheklangan qurilmaga mos OT
\end{enumerate}
\begin{tcolorbox}[enhanced,breakable,colback=teal!4,colframe=teal!60!black,boxrule=0.45pt,arc=1.2mm,left=2.7mm,right=2.7mm,top=1.8mm,bottom=1.8mm,before skip=2mm,after skip=2mm]
{\bfseries\sffamily\color{teal}To'g'ri javob:} \textbf{C) server OT - tarmoq xizmatlari va ko'p mijoz resurslarini boshqarishga mo'ljallangan OT}
\par\smallskip
{\bfseries Izoh:} server OT - tarmoq xizmatlari va ko'p mijoz resurslarini boshqarishga mo'ljallangan OT. To'g'ri juftlik: server OT - tarmoq xizmatlari va ko'p mijoz resurslarini boshqarishga mo'ljallangan OT. Manba: ICT5, PDF 39-41; Qo'shimcha sahifa raqami (qaysi PDFga tegishli ekani aniqlanmagan; darslik sahifasi deb tasdiqlanmagan): 64
\end{tcolorbox}
\vspace{2mm}\hrule height 0.2pt\vspace{3mm}

% DB_ID 9714d742-a62e-49e3-92d0-d7e4ef8d87ae
\needspace{5\baselineskip}
{\bfseries\sffamily\color{ink} 321.} {scriptsizecolor{muted} Y1 | review | 9714d742}
\par
Operatsion tizimning vazifalari mavzusida "multitasking" qatnashgan atama-ta'rif juftliklaridan to'g'risini toping.
\begin{enumerate}[label=\Alph*),leftmargin=9mm]
\item multitasking - bir nechta vazifani navbatlashtirib yoki parallel bajarish
\item multitasking - resurslarni boshqarib, foydalanuvchi va ilovalarga xizmat ko'rsatuvchi tizimli dastur
\item multitasking - bajarilayotgan dastur nusxasi
\item multitasking - fayllarni nomlash, joylashtirish va boshqarish qoidalari
\end{enumerate}
\begin{tcolorbox}[enhanced,breakable,colback=teal!4,colframe=teal!60!black,boxrule=0.45pt,arc=1.2mm,left=2.7mm,right=2.7mm,top=1.8mm,bottom=1.8mm,before skip=2mm,after skip=2mm]
{\bfseries\sffamily\color{teal}To'g'ri javob:} \textbf{A) multitasking - bir nechta vazifani navbatlashtirib yoki parallel bajarish}
\par\smallskip
{\bfseries Izoh:} multitasking - bir nechta vazifani navbatlashtirib yoki parallel bajarish. To'g'ri juftlik: multitasking - bir nechta vazifani navbatlashtirib yoki parallel bajarish. Manba: ICT5, PDF 38-41; Qo'shimcha sahifa raqami (qaysi PDFga tegishli ekani aniqlanmagan; darslik sahifasi deb tasdiqlanmagan): 58
\end{tcolorbox}
\vspace{2mm}\hrule height 0.2pt\vspace{3mm}

% DB_ID 9cd73c56-1cbe-22aa-fa6a-105bd4209cbb
\needspace{5\baselineskip}
{\bfseries\sffamily\color{ink} 322.} {scriptsizecolor{muted} Y1 | review | 9cd73c56}
\par
Operatsion tizimning vazifalari mavzusidagi vazifada bir nechta vazifani navbatlashtirib yoki parallel bajarish talab qilindi. Qaysi tushuncha mos?
\begin{enumerate}[label=\Alph*),leftmargin=9mm]
\item operatsion tizim
\item jarayon
\item fayl tizimi
\item multitasking
\end{enumerate}
\begin{tcolorbox}[enhanced,breakable,colback=teal!4,colframe=teal!60!black,boxrule=0.45pt,arc=1.2mm,left=2.7mm,right=2.7mm,top=1.8mm,bottom=1.8mm,before skip=2mm,after skip=2mm]
{\bfseries\sffamily\color{teal}To'g'ri javob:} \textbf{D) multitasking}
\par\smallskip
{\bfseries Izoh:} multitasking. Berilgan vazifa multitasking tushunchasiga mos: bir nechta vazifani navbatlashtirib yoki parallel bajarish. Manba: ICT5, PDF 38-41; Qo'shimcha sahifa raqami (qaysi PDFga tegishli ekani aniqlanmagan; darslik sahifasi deb tasdiqlanmagan): 58
\end{tcolorbox}
\vspace{2mm}\hrule height 0.2pt\vspace{3mm}

% DB_ID a0dfc5f8-5166-9790-92c1-416cee2d6c8f
\needspace{5\baselineskip}
{\bfseries\sffamily\color{ink} 323.} {scriptsizecolor{muted} Y1 | review | a0dfc5f8}
\par
Operatsion tizim turlari va interfeyslar mavzusida "GUI" qatnashgan atama-ta'rif juftliklaridan to'g'risini toping.
\begin{enumerate}[label=\Alph*),leftmargin=9mm]
\item GUI - oyna, ikona, menyu va ko'rsatkichga tayangan grafik interfeys
\item GUI - matnli buyruqlar orqali boshqariladigan interfeys
\item GUI - tarmoq xizmatlari va ko'p mijoz resurslarini boshqarishga mo'ljallangan OT
\item GUI - sensorli, portativ va energiya cheklangan qurilmaga mos OT
\end{enumerate}
\begin{tcolorbox}[enhanced,breakable,colback=teal!4,colframe=teal!60!black,boxrule=0.45pt,arc=1.2mm,left=2.7mm,right=2.7mm,top=1.8mm,bottom=1.8mm,before skip=2mm,after skip=2mm]
{\bfseries\sffamily\color{teal}To'g'ri javob:} \textbf{A) GUI - oyna, ikona, menyu va ko'rsatkichga tayangan grafik interfeys}
\par\smallskip
{\bfseries Izoh:} GUI - oyna, ikona, menyu va ko'rsatkichga tayangan grafik interfeys. To'g'ri juftlik: GUI - oyna, ikona, menyu va ko'rsatkichga tayangan grafik interfeys. Manba: CAM9, PDF 13-15; Qo'shimcha sahifa raqami (qaysi PDFga tegishli ekani aniqlanmagan; darslik sahifasi deb tasdiqlanmagan): 62
\end{tcolorbox}
\vspace{2mm}\hrule height 0.2pt\vspace{3mm}

% DB_ID b30d4f68-b2d9-7d5f-6dcd-65140b4e9442
\needspace{5\baselineskip}
{\bfseries\sffamily\color{ink} 324.} {scriptsizecolor{muted} Y1 | review | b30d4f68}
\par
Operatsion tizim turlari va interfeyslar mavzusidagi vazifada tarmoq xizmatlari va ko'p mijoz resurslarini boshqarishga mo'ljallangan OT talab qilindi. Qaysi tushuncha mos?
\begin{enumerate}[label=\Alph*),leftmargin=9mm]
\item GUI
\item server OT
\item CLI
\item mobil OT
\end{enumerate}
\begin{tcolorbox}[enhanced,breakable,colback=teal!4,colframe=teal!60!black,boxrule=0.45pt,arc=1.2mm,left=2.7mm,right=2.7mm,top=1.8mm,bottom=1.8mm,before skip=2mm,after skip=2mm]
{\bfseries\sffamily\color{teal}To'g'ri javob:} \textbf{B) server OT}
\par\smallskip
{\bfseries Izoh:} server OT. Berilgan vazifa server OT tushunchasiga mos: tarmoq xizmatlari va ko'p mijoz resurslarini boshqarishga mo'ljallangan OT. Manba: ICT5, PDF 39-41; Qo'shimcha sahifa raqami (qaysi PDFga tegishli ekani aniqlanmagan; darslik sahifasi deb tasdiqlanmagan): 64
\end{tcolorbox}
\vspace{2mm}\hrule height 0.2pt\vspace{3mm}

% DB_ID cd3e145c-c916-4f3e-8a23-8ed819a51c9c
\needspace{5\baselineskip}
{\bfseries\sffamily\color{ink} 325.} {scriptsizecolor{muted} Y1 | review | cd3e145c}
\par
Operatsion tizimning vazifalari mavzusidagi vazifada bajarilayotgan dastur nusxasi talab qilindi. Qaysi tushuncha mos?
\begin{enumerate}[label=\Alph*),leftmargin=9mm]
\item operatsion tizim
\item multitasking
\item jarayon
\item fayl tizimi
\end{enumerate}
\begin{tcolorbox}[enhanced,breakable,colback=teal!4,colframe=teal!60!black,boxrule=0.45pt,arc=1.2mm,left=2.7mm,right=2.7mm,top=1.8mm,bottom=1.8mm,before skip=2mm,after skip=2mm]
{\bfseries\sffamily\color{teal}To'g'ri javob:} \textbf{C) jarayon}
\par\smallskip
{\bfseries Izoh:} jarayon. Berilgan vazifa jarayon tushunchasiga mos: bajarilayotgan dastur nusxasi. Manba: CAM10/11, PDF 43-47; Qo'shimcha sahifa raqami (qaysi PDFga tegishli ekani aniqlanmagan; darslik sahifasi deb tasdiqlanmagan): 57
\end{tcolorbox}
\vspace{2mm}\hrule height 0.2pt\vspace{3mm}

% DB_ID e9d77918-af9e-6a2d-4a5e-1439330d06d5
\needspace{5\baselineskip}
{\bfseries\sffamily\color{ink} 326.} {scriptsizecolor{muted} Y1 | review | e9d77918}
\par
Operatsion tizimning vazifalari mavzusidagi vazifada OTning resurslar va apparatni boshqaruvchi markaziy qismi talab qilindi. Qaysi tushuncha mos?
\begin{enumerate}[label=\Alph*),leftmargin=9mm]
\item operatsion tizim
\item yadro
\item jarayon
\item multitasking
\end{enumerate}
\begin{tcolorbox}[enhanced,breakable,colback=teal!4,colframe=teal!60!black,boxrule=0.45pt,arc=1.2mm,left=2.7mm,right=2.7mm,top=1.8mm,bottom=1.8mm,before skip=2mm,after skip=2mm]
{\bfseries\sffamily\color{teal}To'g'ri javob:} \textbf{B) yadro}
\par\smallskip
{\bfseries Izoh:} yadro. Berilgan vazifa yadro tushunchasiga mos: OTning resurslar va apparatni boshqaruvchi markaziy qismi. Manba: MAR, PDF 11-16; Qo'shimcha sahifa raqami (qaysi PDFga tegishli ekani aniqlanmagan; darslik sahifasi deb tasdiqlanmagan): 60
\end{tcolorbox}
\vspace{2mm}\hrule height 0.2pt\vspace{3mm}

% DB_ID f77e1b4d-f4fc-556c-919b-dbda9ea40046
\needspace{5\baselineskip}
{\bfseries\sffamily\color{ink} 327.} {scriptsizecolor{muted} Y1 | review | f77e1b4d}
\par
Operatsion tizimning vazifalari mavzusidagi vazifada operatsion tizim yadrosini xotiraga yuklovchi dastur talab qilindi. Qaysi tushuncha mos?
\begin{enumerate}[label=\Alph*),leftmargin=9mm]
\item operatsion tizim
\item jarayon
\item bootloader
\item multitasking
\end{enumerate}
\begin{tcolorbox}[enhanced,breakable,colback=teal!4,colframe=teal!60!black,boxrule=0.45pt,arc=1.2mm,left=2.7mm,right=2.7mm,top=1.8mm,bottom=1.8mm,before skip=2mm,after skip=2mm]
{\bfseries\sffamily\color{teal}To'g'ri javob:} \textbf{C) bootloader}
\par\smallskip
{\bfseries Izoh:} bootloader. Berilgan vazifa bootloader tushunchasiga mos: operatsion tizim yadrosini xotiraga yuklovchi dastur. Manba: CAM9, PDF 12-15; Qo'shimcha sahifa raqami (qaysi PDFga tegishli ekani aniqlanmagan; darslik sahifasi deb tasdiqlanmagan): 61
\end{tcolorbox}
\vspace{2mm}\hrule height 0.2pt\vspace{3mm}

% DB_ID 05d9fe80-7e64-11b7-ebd6-7884a663091f
\needspace{5\baselineskip}
{\bfseries\sffamily\color{ink} 328.} {scriptsizecolor{muted} Y1 | review | 05d9fe80}
\par
Operatsion tizimning vazifalari bo'yicha qatorlar ichidan "operatsion tizim" noto'g'ri izohlanganini aniqlang.
\begin{enumerate}[label=\Alph*),leftmargin=9mm]
\item multitasking - bir nechta vazifani navbatlashtirib yoki parallel bajarish
\item fayl tizimi - fayllarni nomlash, joylashtirish va boshqarish qoidalari
\item yadro - OTning resurslar va apparatni boshqaruvchi markaziy qismi
\item operatsion tizim - bajarilayotgan dastur nusxasi
\end{enumerate}
\begin{tcolorbox}[enhanced,breakable,colback=teal!4,colframe=teal!60!black,boxrule=0.45pt,arc=1.2mm,left=2.7mm,right=2.7mm,top=1.8mm,bottom=1.8mm,before skip=2mm,after skip=2mm]
{\bfseries\sffamily\color{teal}To'g'ri javob:} \textbf{D) operatsion tizim - bajarilayotgan dastur nusxasi}
\par\smallskip
{\bfseries Izoh:} operatsion tizim - bajarilayotgan dastur nusxasi. operatsion tizim uchun to'g'ri ta'rif: resurslarni boshqarib, foydalanuvchi va ilovalarga xizmat ko'rsatuvchi tizimli dastur. Manba: CAM9, PDF 12-15; Qo'shimcha sahifa raqami (qaysi PDFga tegishli ekani aniqlanmagan; darslik sahifasi deb tasdiqlanmagan): 56
\end{tcolorbox}
\vspace{2mm}\hrule height 0.2pt\vspace{3mm}

% DB_ID 0ba2a3e6-f369-5246-9e1a-8689ccc702ec
\needspace{5\baselineskip}
{\bfseries\sffamily\color{ink} 329.} {scriptsizecolor{muted} Y1 | review | 0ba2a3e6}
\par
Operatsion tizim turlari va interfeyslar mavzusida "mobil OT" qatnashgan atama-ta'rif juftliklaridan to'g'risini toping.
\begin{enumerate}[label=\Alph*),leftmargin=9mm]
\item mobil OT - oyna, ikona, menyu va ko'rsatkichga tayangan grafik interfeys
\item mobil OT - matnli buyruqlar orqali boshqariladigan interfeys
\item mobil OT - tarmoq xizmatlari va ko'p mijoz resurslarini boshqarishga mo'ljallangan OT
\item mobil OT - sensorli, portativ va energiya cheklangan qurilmaga mos OT
\end{enumerate}
\begin{tcolorbox}[enhanced,breakable,colback=teal!4,colframe=teal!60!black,boxrule=0.45pt,arc=1.2mm,left=2.7mm,right=2.7mm,top=1.8mm,bottom=1.8mm,before skip=2mm,after skip=2mm]
{\bfseries\sffamily\color{teal}To'g'ri javob:} \textbf{D) mobil OT - sensorli, portativ va energiya cheklangan qurilmaga mos OT}
\par\smallskip
{\bfseries Izoh:} mobil OT - sensorli, portativ va energiya cheklangan qurilmaga mos OT. To'g'ri juftlik: mobil OT - sensorli, portativ va energiya cheklangan qurilmaga mos OT. Manba: CAM7, PDF 140-146; Qo'shimcha sahifa raqami (qaysi PDFga tegishli ekani aniqlanmagan; darslik sahifasi deb tasdiqlanmagan): 66
\end{tcolorbox}
\vspace{2mm}\hrule height 0.2pt\vspace{3mm}

% DB_ID 288cd318-833d-19a6-7aae-1177cf670934
\needspace{5\baselineskip}
{\bfseries\sffamily\color{ink} 330.} {scriptsizecolor{muted} Y1 | review | 288cd318}
\par
Operatsion tizimning vazifalari bo'yicha qatorlar ichidan "multitasking" noto'g'ri izohlanganini aniqlang.
\begin{enumerate}[label=\Alph*),leftmargin=9mm]
\item jarayon - bajarilayotgan dastur nusxasi
\item multitasking - resurslarni boshqarib, foydalanuvchi va ilovalarga xizmat ko'rsatuvchi tizimli dastur
\item fayl tizimi - fayllarni nomlash, joylashtirish va boshqarish qoidalari
\item yadro - OTning resurslar va apparatni boshqaruvchi markaziy qismi
\end{enumerate}
\begin{tcolorbox}[enhanced,breakable,colback=teal!4,colframe=teal!60!black,boxrule=0.45pt,arc=1.2mm,left=2.7mm,right=2.7mm,top=1.8mm,bottom=1.8mm,before skip=2mm,after skip=2mm]
{\bfseries\sffamily\color{teal}To'g'ri javob:} \textbf{B) multitasking - resurslarni boshqarib, foydalanuvchi va ilovalarga xizmat ko'rsatuvchi tizimli dastur}
\par\smallskip
{\bfseries Izoh:} multitasking - resurslarni boshqarib, foydalanuvchi va ilovalarga xizmat ko'rsatuvchi tizimli dastur. multitasking uchun to'g'ri ta'rif: bir nechta vazifani navbatlashtirib yoki parallel bajarish. Manba: ICT5, PDF 38-41; Qo'shimcha sahifa raqami (qaysi PDFga tegishli ekani aniqlanmagan; darslik sahifasi deb tasdiqlanmagan): 58
\end{tcolorbox}
\vspace{2mm}\hrule height 0.2pt\vspace{3mm}

% DB_ID 4d12a18d-bcff-a308-d0ee-0acb9975a74a
\needspace{5\baselineskip}
{\bfseries\sffamily\color{ink} 331.} {scriptsizecolor{muted} Y1 | review | 4d12a18d}
\par
Operatsion tizim turlari va interfeyslar bo'yicha qatorlar ichidan "GUI" noto'g'ri izohlanganini aniqlang.
\begin{enumerate}[label=\Alph*),leftmargin=9mm]
\item server OT - tarmoq xizmatlari va ko'p mijoz resurslarini boshqarishga mo'ljallangan OT
\item GUI - matnli buyruqlar orqali boshqariladigan interfeys
\item mobil OT - sensorli, portativ va energiya cheklangan qurilmaga mos OT
\item embedded OT - maxsus qurilma vazifasiga moslashtirilgan OT
\end{enumerate}
\begin{tcolorbox}[enhanced,breakable,colback=teal!4,colframe=teal!60!black,boxrule=0.45pt,arc=1.2mm,left=2.7mm,right=2.7mm,top=1.8mm,bottom=1.8mm,before skip=2mm,after skip=2mm]
{\bfseries\sffamily\color{teal}To'g'ri javob:} \textbf{B) GUI - matnli buyruqlar orqali boshqariladigan interfeys}
\par\smallskip
{\bfseries Izoh:} GUI - matnli buyruqlar orqali boshqariladigan interfeys. GUI uchun to'g'ri ta'rif: oyna, ikona, menyu va ko'rsatkichga tayangan grafik interfeys. Manba: CAM9, PDF 13-15; Qo'shimcha sahifa raqami (qaysi PDFga tegishli ekani aniqlanmagan; darslik sahifasi deb tasdiqlanmagan): 62
\end{tcolorbox}
\vspace{2mm}\hrule height 0.2pt\vspace{3mm}

% DB_ID 5170cdbd-a9c6-6bd0-9b32-8a00b7a2c884
\needspace{5\baselineskip}
{\bfseries\sffamily\color{ink} 332.} {scriptsizecolor{muted} Y1 | review | 5170cdbd}
\par
Operatsion tizim turlari va interfeyslar bo'yicha qatorlar ichidan "server OT" noto'g'ri izohlanganini aniqlang.
\begin{enumerate}[label=\Alph*),leftmargin=9mm]
\item CLI - matnli buyruqlar orqali boshqariladigan interfeys
\item mobil OT - sensorli, portativ va energiya cheklangan qurilmaga mos OT
\item embedded OT - maxsus qurilma vazifasiga moslashtirilgan OT
\item server OT - oyna, ikona, menyu va ko'rsatkichga tayangan grafik interfeys
\end{enumerate}
\begin{tcolorbox}[enhanced,breakable,colback=teal!4,colframe=teal!60!black,boxrule=0.45pt,arc=1.2mm,left=2.7mm,right=2.7mm,top=1.8mm,bottom=1.8mm,before skip=2mm,after skip=2mm]
{\bfseries\sffamily\color{teal}To'g'ri javob:} \textbf{D) server OT - oyna, ikona, menyu va ko'rsatkichga tayangan grafik interfeys}
\par\smallskip
{\bfseries Izoh:} server OT - oyna, ikona, menyu va ko'rsatkichga tayangan grafik interfeys. server OT uchun to'g'ri ta'rif: tarmoq xizmatlari va ko'p mijoz resurslarini boshqarishga mo'ljallangan OT. Manba: ICT5, PDF 39-41; Qo'shimcha sahifa raqami (qaysi PDFga tegishli ekani aniqlanmagan; darslik sahifasi deb tasdiqlanmagan): 65
\end{tcolorbox}
\vspace{2mm}\hrule height 0.2pt\vspace{3mm}

% DB_ID 577b6236-7967-0681-09d2-63b2175a7ef3
\needspace{5\baselineskip}
{\bfseries\sffamily\color{ink} 333.} {scriptsizecolor{muted} Y1 | review | 577b6236}
\par
Operatsion tizimning vazifalari mavzusida "jarayon" qatnashgan atama-ta'rif juftliklaridan to'g'risini toping.
\begin{enumerate}[label=\Alph*),leftmargin=9mm]
\item jarayon - resurslarni boshqarib, foydalanuvchi va ilovalarga xizmat ko'rsatuvchi tizimli dastur
\item jarayon - bir nechta vazifani navbatlashtirib yoki parallel bajarish
\item jarayon - fayllarni nomlash, joylashtirish va boshqarish qoidalari
\item jarayon - bajarilayotgan dastur nusxasi
\end{enumerate}
\begin{tcolorbox}[enhanced,breakable,colback=teal!4,colframe=teal!60!black,boxrule=0.45pt,arc=1.2mm,left=2.7mm,right=2.7mm,top=1.8mm,bottom=1.8mm,before skip=2mm,after skip=2mm]
{\bfseries\sffamily\color{teal}To'g'ri javob:} \textbf{D) jarayon - bajarilayotgan dastur nusxasi}
\par\smallskip
{\bfseries Izoh:} jarayon - bajarilayotgan dastur nusxasi. To'g'ri juftlik: jarayon - bajarilayotgan dastur nusxasi. Manba: CAM10/11, PDF 43-47; Qo'shimcha sahifa raqami (qaysi PDFga tegishli ekani aniqlanmagan; darslik sahifasi deb tasdiqlanmagan): 57
\end{tcolorbox}
\vspace{2mm}\hrule height 0.2pt\vspace{3mm}

% DB_ID 625eae6e-be79-5f84-485b-037d85772579
\needspace{5\baselineskip}
{\bfseries\sffamily\color{ink} 334.} {scriptsizecolor{muted} Y1 | review | 625eae6e}
\par
Operatsion tizim turlari va interfeyslar mavzusida "real vaqt OT" qatnashgan atama-ta'rif juftliklaridan to'g'risini toping.
\begin{enumerate}[label=\Alph*),leftmargin=9mm]
\item real vaqt OT - oyna, ikona, menyu va ko'rsatkichga tayangan grafik interfeys
\item real vaqt OT - javobni belgilangan vaqt chegarasida berishni kafolatlashga yo'naltirilgan OT
\item real vaqt OT - matnli buyruqlar orqali boshqariladigan interfeys
\item real vaqt OT - tarmoq xizmatlari va ko'p mijoz resurslarini boshqarishga mo'ljallangan OT
\end{enumerate}
\begin{tcolorbox}[enhanced,breakable,colback=teal!4,colframe=teal!60!black,boxrule=0.45pt,arc=1.2mm,left=2.7mm,right=2.7mm,top=1.8mm,bottom=1.8mm,before skip=2mm,after skip=2mm]
{\bfseries\sffamily\color{teal}To'g'ri javob:} \textbf{B) real vaqt OT - javobni belgilangan vaqt chegarasida berishni kafolatlashga yo'naltirilgan OT}
\par\smallskip
{\bfseries Izoh:} real vaqt OT - javobni belgilangan vaqt chegarasida berishni kafolatlashga yo'naltirilgan OT. To'g'ri juftlik: real vaqt OT - javobni belgilangan vaqt chegarasida berishni kafolatlashga yo'naltirilgan OT. Manba: MAR, PDF 12-15; Qo'shimcha sahifa raqami (qaysi PDFga tegishli ekani aniqlanmagan; darslik sahifasi deb tasdiqlanmagan): 68
\end{tcolorbox}
\vspace{2mm}\hrule height 0.2pt\vspace{3mm}

% DB_ID 65e62fa9-8f65-e60f-2861-661895bd3f85
\needspace{5\baselineskip}
{\bfseries\sffamily\color{ink} 335.} {scriptsizecolor{muted} Y1 | review | 65e62fa9}
\par
Operatsion tizim turlari va interfeyslar bo'yicha qatorlar ichidan "embedded OT" noto'g'ri izohlanganini aniqlang.
\begin{enumerate}[label=\Alph*),leftmargin=9mm]
\item CLI - matnli buyruqlar orqali boshqariladigan interfeys
\item embedded OT - oyna, ikona, menyu va ko'rsatkichga tayangan grafik interfeys
\item server OT - tarmoq xizmatlari va ko'p mijoz resurslarini boshqarishga mo'ljallangan OT
\item mobil OT - sensorli, portativ va energiya cheklangan qurilmaga mos OT
\end{enumerate}
\begin{tcolorbox}[enhanced,breakable,colback=teal!4,colframe=teal!60!black,boxrule=0.45pt,arc=1.2mm,left=2.7mm,right=2.7mm,top=1.8mm,bottom=1.8mm,before skip=2mm,after skip=2mm]
{\bfseries\sffamily\color{teal}To'g'ri javob:} \textbf{B) embedded OT - oyna, ikona, menyu va ko'rsatkichga tayangan grafik interfeys}
\par\smallskip
{\bfseries Izoh:} embedded OT - oyna, ikona, menyu va ko'rsatkichga tayangan grafik interfeys. embedded OT uchun to'g'ri ta'rif: maxsus qurilma vazifasiga moslashtirilgan OT. Manba: ICT11, PDF 244-245; Qo'shimcha sahifa raqami (qaysi PDFga tegishli ekani aniqlanmagan; darslik sahifasi deb tasdiqlanmagan): 67
\end{tcolorbox}
\vspace{2mm}\hrule height 0.2pt\vspace{3mm}

% DB_ID 855e1aa8-5561-8d9e-85ba-16f4f0e16657
\needspace{5\baselineskip}
{\bfseries\sffamily\color{ink} 336.} {scriptsizecolor{muted} Y1 | review | 855e1aa8}
\par
Ikki dastur bir vaqtda ishlayotgandek ko'rinadi. CPU vaqtini ularga taqsimlayotgan qism qaysi?
\begin{enumerate}[label=\Alph*),leftmargin=9mm]
\item operatsion tizim
\item matn protsessori
\item monitor firmware'i
\item fayl kengaytmasi
\end{enumerate}
\begin{tcolorbox}[enhanced,breakable,colback=teal!4,colframe=teal!60!black,boxrule=0.45pt,arc=1.2mm,left=2.7mm,right=2.7mm,top=1.8mm,bottom=1.8mm,before skip=2mm,after skip=2mm]
{\bfseries\sffamily\color{teal}To'g'ri javob:} \textbf{A) operatsion tizim}
\par\smallskip
{\bfseries Izoh:} operatsion tizim. Jarayonlarga CPU vaqtini rejalashtirish OT vazifasidir. Manba: CAM9, PDF 12-15; Qo'shimcha sahifa raqami (qaysi PDFga tegishli ekani aniqlanmagan; darslik sahifasi deb tasdiqlanmagan): 61
\end{tcolorbox}
\vspace{2mm}\hrule height 0.2pt\vspace{3mm}

% DB_ID aaaa869c-5762-e63a-8a3e-824b56731e86
\needspace{5\baselineskip}
{\bfseries\sffamily\color{ink} 337.} {scriptsizecolor{muted} Y1 | review | aaaa869c}
\par
Operatsion tizim turlari va interfeyslar mavzusida "CLI" qatnashgan atama-ta'rif juftliklaridan to'g'risini toping.
\begin{enumerate}[label=\Alph*),leftmargin=9mm]
\item CLI - oyna, ikona, menyu va ko'rsatkichga tayangan grafik interfeys
\item CLI - matnli buyruqlar orqali boshqariladigan interfeys
\item CLI - tarmoq xizmatlari va ko'p mijoz resurslarini boshqarishga mo'ljallangan OT
\item CLI - sensorli, portativ va energiya cheklangan qurilmaga mos OT
\end{enumerate}
\begin{tcolorbox}[enhanced,breakable,colback=teal!4,colframe=teal!60!black,boxrule=0.45pt,arc=1.2mm,left=2.7mm,right=2.7mm,top=1.8mm,bottom=1.8mm,before skip=2mm,after skip=2mm]
{\bfseries\sffamily\color{teal}To'g'ri javob:} \textbf{B) CLI - matnli buyruqlar orqali boshqariladigan interfeys}
\par\smallskip
{\bfseries Izoh:} CLI - matnli buyruqlar orqali boshqariladigan interfeys. To'g'ri juftlik: CLI - matnli buyruqlar orqali boshqariladigan interfeys. Manba: CAM10/11, PDF 44-48; Qo'shimcha sahifa raqami (qaysi PDFga tegishli ekani aniqlanmagan; darslik sahifasi deb tasdiqlanmagan): 63
\end{tcolorbox}
\vspace{2mm}\hrule height 0.2pt\vspace{3mm}

% DB_ID b03daae7-4636-7c10-e51f-080bd6d6b9e4
\needspace{5\baselineskip}
{\bfseries\sffamily\color{ink} 338.} {scriptsizecolor{muted} Y1 | review | b03daae7}
\par
Operatsion tizimning vazifalari mavzusida "fayl tizimi" qatnashgan atama-ta'rif juftliklaridan to'g'risini toping.
\begin{enumerate}[label=\Alph*),leftmargin=9mm]
\item fayl tizimi - resurslarni boshqarib, foydalanuvchi va ilovalarga xizmat ko'rsatuvchi tizimli dastur
\item fayl tizimi - fayllarni nomlash, joylashtirish va boshqarish qoidalari
\item fayl tizimi - bajarilayotgan dastur nusxasi
\item fayl tizimi - bir nechta vazifani navbatlashtirib yoki parallel bajarish
\end{enumerate}
\begin{tcolorbox}[enhanced,breakable,colback=teal!4,colframe=teal!60!black,boxrule=0.45pt,arc=1.2mm,left=2.7mm,right=2.7mm,top=1.8mm,bottom=1.8mm,before skip=2mm,after skip=2mm]
{\bfseries\sffamily\color{teal}To'g'ri javob:} \textbf{B) fayl tizimi - fayllarni nomlash, joylashtirish va boshqarish qoidalari}
\par\smallskip
{\bfseries Izoh:} fayl tizimi - fayllarni nomlash, joylashtirish va boshqarish qoidalari. To'g'ri juftlik: fayl tizimi - fayllarni nomlash, joylashtirish va boshqarish qoidalari. Manba: CAM8, PDF 99-106; Qo'shimcha sahifa raqami (qaysi PDFga tegishli ekani aniqlanmagan; darslik sahifasi deb tasdiqlanmagan): 59
\end{tcolorbox}
\vspace{2mm}\hrule height 0.2pt\vspace{3mm}

% DB_ID c3698c6c-da44-fc0b-37d1-c65d3eb03930
\needspace{5\baselineskip}
{\bfseries\sffamily\color{ink} 339.} {scriptsizecolor{muted} Y1 | review | c3698c6c}
\par
Operatsion tizimning vazifalari bo'yicha qatorlar ichidan "yadro" noto'g'ri izohlanganini aniqlang.
\begin{enumerate}[label=\Alph*),leftmargin=9mm]
\item jarayon - bajarilayotgan dastur nusxasi
\item multitasking - bir nechta vazifani navbatlashtirib yoki parallel bajarish
\item fayl tizimi - fayllarni nomlash, joylashtirish va boshqarish qoidalari
\item yadro - resurslarni boshqarib, foydalanuvchi va ilovalarga xizmat ko'rsatuvchi tizimli dastur
\end{enumerate}
\begin{tcolorbox}[enhanced,breakable,colback=teal!4,colframe=teal!60!black,boxrule=0.45pt,arc=1.2mm,left=2.7mm,right=2.7mm,top=1.8mm,bottom=1.8mm,before skip=2mm,after skip=2mm]
{\bfseries\sffamily\color{teal}To'g'ri javob:} \textbf{D) yadro - resurslarni boshqarib, foydalanuvchi va ilovalarga xizmat ko'rsatuvchi tizimli dastur}
\par\smallskip
{\bfseries Izoh:} yadro - resurslarni boshqarib, foydalanuvchi va ilovalarga xizmat ko'rsatuvchi tizimli dastur. yadro uchun to'g'ri ta'rif: OTning resurslar va apparatni boshqaruvchi markaziy qismi. Manba: MAR, PDF 11-16; Qo'shimcha sahifa raqami (qaysi PDFga tegishli ekani aniqlanmagan; darslik sahifasi deb tasdiqlanmagan): 60
\end{tcolorbox}
\vspace{2mm}\hrule height 0.2pt\vspace{3mm}

% DB_ID ddf66256-219f-7e33-b050-4d573e399b6c
\needspace{5\baselineskip}
{\bfseries\sffamily\color{ink} 340.} {scriptsizecolor{muted} Y1 | review | ddf66256}
\par
Operatsion tizimning vazifalari mavzusida "bootloader" qatnashgan atama-ta'rif juftliklaridan to'g'risini toping.
\begin{enumerate}[label=\Alph*),leftmargin=9mm]
\item bootloader - resurslarni boshqarib, foydalanuvchi va ilovalarga xizmat ko'rsatuvchi tizimli dastur
\item bootloader - bajarilayotgan dastur nusxasi
\item bootloader - bir nechta vazifani navbatlashtirib yoki parallel bajarish
\item bootloader - operatsion tizim yadrosini xotiraga yuklovchi dastur
\end{enumerate}
\begin{tcolorbox}[enhanced,breakable,colback=teal!4,colframe=teal!60!black,boxrule=0.45pt,arc=1.2mm,left=2.7mm,right=2.7mm,top=1.8mm,bottom=1.8mm,before skip=2mm,after skip=2mm]
{\bfseries\sffamily\color{teal}To'g'ri javob:} \textbf{D) bootloader - operatsion tizim yadrosini xotiraga yuklovchi dastur}
\par\smallskip
{\bfseries Izoh:} bootloader - operatsion tizim yadrosini xotiraga yuklovchi dastur. To'g'ri juftlik: bootloader - operatsion tizim yadrosini xotiraga yuklovchi dastur. Manba: CAM9, PDF 12-15; Qo'shimcha sahifa raqami (qaysi PDFga tegishli ekani aniqlanmagan; darslik sahifasi deb tasdiqlanmagan): 61
\end{tcolorbox}
\vspace{2mm}\hrule height 0.2pt\vspace{3mm}

% DB_ID ece58b14-f6a6-2a48-cd81-b1f5b9b21590
\needspace{5\baselineskip}
{\bfseries\sffamily\color{ink} 341.} {scriptsizecolor{muted} Y1 | review | ece58b14}
\par
Kompyuterda matn muharriri, brauzer va taqdimot dasturi birga ishlayapti. Ular protsessor va xotiradan foydalanayotganda asosiy boshqaruv qatlami qaysi?
\begin{enumerate}[label=\Alph*),leftmargin=9mm]
\item Matn muharriri, chunki u barcha dasturlar uchun ishga tushirish navbatini va resurslarni to'liq taqsimlaydi
\item Taqdimot dasturi, chunki u barcha dasturlar o'rtasidagi apparat resurslarini avtomatik boshqaradi
\item Operatsion tizim, chunki u jarayonlar va apparat resurslarini muvofiqlashtiradi
\item Brauzer, chunki u barcha xotira manzillarini boshqarib, protsessor vaqtini ham taqsimlab beradi
\end{enumerate}
\begin{tcolorbox}[enhanced,breakable,colback=teal!4,colframe=teal!60!black,boxrule=0.45pt,arc=1.2mm,left=2.7mm,right=2.7mm,top=1.8mm,bottom=1.8mm,before skip=2mm,after skip=2mm]
{\bfseries\sffamily\color{teal}To'g'ri javob:} \textbf{C) Operatsion tizim, chunki u jarayonlar va apparat resurslarini muvofiqlashtiradi}
\par\smallskip
{\bfseries Izoh:} Operatsion tizim kompyuterning qurilma va dasturlari, shuningdek ular orasidagi aloqa hamda resurslarini boshqaradi. [Manba: ICT 10-sinf (2021), PDF 24-bet.]
\end{tcolorbox}
\vspace{2mm}\hrule height 0.2pt\vspace{3mm}

% DB_ID 0322568b-48f1-2251-e265-20d23d576297
\needspace{5\baselineskip}
{\bfseries\sffamily\color{ink} 342.} {scriptsizecolor{muted} Y1 | review | 0322568b}
\par
Operatsion tizimning vazifalari mavzusida: I. yadro - OTning resurslar va apparatni boshqaruvchi markaziy qismi. II. operatsion tizim - resurslarni boshqarib, foydalanuvchi va ilovalarga xizmat ko'rsatuvchi tizimli dastur. Mulohazalarni baholang.
\begin{enumerate}[label=\Alph*),leftmargin=9mm]
\item I va II
\item faqat I
\item faqat II
\item I ham, II ham emas
\end{enumerate}
\begin{tcolorbox}[enhanced,breakable,colback=teal!4,colframe=teal!60!black,boxrule=0.45pt,arc=1.2mm,left=2.7mm,right=2.7mm,top=1.8mm,bottom=1.8mm,before skip=2mm,after skip=2mm]
{\bfseries\sffamily\color{teal}To'g'ri javob:} \textbf{A) I va II}
\par\smallskip
{\bfseries Izoh:} I va II. Birinchi va ikkinchi mulohazaning ta'riflari alohida tekshirildi; natija: I va II. Manba: MAR, PDF 11-16; Qo'shimcha sahifa raqami (qaysi PDFga tegishli ekani aniqlanmagan; darslik sahifasi deb tasdiqlanmagan): 61
\end{tcolorbox}
\vspace{2mm}\hrule height 0.2pt\vspace{3mm}

% DB_ID 10c62a73-52b1-672a-ce3c-603bcc548e74
\needspace{5\baselineskip}
{\bfseries\sffamily\color{ink} 343.} {scriptsizecolor{muted} Y1 | review | 10c62a73}
\par
Administrator yuzlab serverda bir xil buyruqlar ketma-ketligini avtomatlashtirmoqchi. Qaysi interfeys afzal?
\begin{enumerate}[label=\Alph*),leftmargin=9mm]
\item faqat sensorli menyu
\item faqat grafik ikonalar
\item CLI va skript
\item ovozli diktant
\end{enumerate}
\begin{tcolorbox}[enhanced,breakable,colback=teal!4,colframe=teal!60!black,boxrule=0.45pt,arc=1.2mm,left=2.7mm,right=2.7mm,top=1.8mm,bottom=1.8mm,before skip=2mm,after skip=2mm]
{\bfseries\sffamily\color{teal}To'g'ri javob:} \textbf{C) CLI va skript}
\par\smallskip
{\bfseries Izoh:} CLI va skript. CLI buyruqlarni skriptga birlashtirib takroriy boshqaruvni avtomatlashtiradi. Manba: CAM9, PDF 13-15; Qo'shimcha sahifa raqami (qaysi PDFga tegishli ekani aniqlanmagan; darslik sahifasi deb tasdiqlanmagan): 68
\end{tcolorbox}
\vspace{2mm}\hrule height 0.2pt\vspace{3mm}

% DB_ID 18faa2c0-baab-f029-8b39-305edc9f27b9
\needspace{5\baselineskip}
{\bfseries\sffamily\color{ink} 344.} {scriptsizecolor{muted} Y1 | review | 18faa2c0}
\par
IT mutaxassisi kompyuterda Windows va NTFS atamalarini uchratdi. Ularning vazifasi bo'yicha qaysi izoh aniq?
\begin{enumerate}[label=\Alph*),leftmargin=9mm]
\item Windows fayllarni faqat siqadi, NTFS esa faqat foydalanuvchi oynalarini chizadi
\item Windows printer drayveri, NTFS esa matn hujjatlarini ko'rsatadigan ilovadir
\item Windows operatsion tizim, NTFS esa fayllarni tashkil etuvchi fayl tizimidir
\item Windows fayl kengaytmasi, NTFS esa operativ xotirada saqlanuvchi jarayondir
\end{enumerate}
\begin{tcolorbox}[enhanced,breakable,colback=teal!4,colframe=teal!60!black,boxrule=0.45pt,arc=1.2mm,left=2.7mm,right=2.7mm,top=1.8mm,bottom=1.8mm,before skip=2mm,after skip=2mm]
{\bfseries\sffamily\color{teal}To'g'ri javob:} \textbf{C) Windows operatsion tizim, NTFS esa fayllarni tashkil etuvchi fayl tizimidir}
\par\smallskip
{\bfseries Izoh:} ICT 10-sinfda Windows operatsion tizim, NTFS esa saqlashni tashkil etuvchi fayl tizimiga misol sifatida keltirilgan. [Manba: ICT 10-sinf, PDF 25-26.]
\end{tcolorbox}
\vspace{2mm}\hrule height 0.2pt\vspace{3mm}

% DB_ID 23d34be9-4eec-322a-124d-ee469b902f68
\needspace{5\baselineskip}
{\bfseries\sffamily\color{ink} 345.} {scriptsizecolor{muted} Y1 | review | 23d34be9}
\par
Operatsion tizim turlari va interfeyslar mavzusida: I. server OT - tarmoq xizmatlari va ko'p mijoz resurslarini boshqarishga mo'ljallangan OT. II. GUI - tarmoq xizmatlari va ko'p mijoz resurslarini boshqarishga mo'ljallangan OT. Mulohazalarni baholang.
\begin{enumerate}[label=\Alph*),leftmargin=9mm]
\item faqat I
\item faqat II
\item I va II
\item I ham, II ham emas
\end{enumerate}
\begin{tcolorbox}[enhanced,breakable,colback=teal!4,colframe=teal!60!black,boxrule=0.45pt,arc=1.2mm,left=2.7mm,right=2.7mm,top=1.8mm,bottom=1.8mm,before skip=2mm,after skip=2mm]
{\bfseries\sffamily\color{teal}To'g'ri javob:} \textbf{A) faqat I}
\par\smallskip
{\bfseries Izoh:} faqat I. Birinchi va ikkinchi mulohazaning ta'riflari alohida tekshirildi; natija: faqat I. Manba: ICT5, PDF 39-41; Qo'shimcha sahifa raqami (qaysi PDFga tegishli ekani aniqlanmagan; darslik sahifasi deb tasdiqlanmagan): 65
\end{tcolorbox}
\vspace{2mm}\hrule height 0.2pt\vspace{3mm}

% DB_ID 47632ac0-0e69-b202-7b67-439fd6e8d5e7
\needspace{5\baselineskip}
{\bfseries\sffamily\color{ink} 346.} {scriptsizecolor{muted} Y1 | review | 47632ac0}
\par
Operatsion tizimning vazifalari mavzusida: I. multitasking - bir nechta vazifani navbatlashtirib yoki parallel bajarish. II. operatsion tizim - bir nechta vazifani navbatlashtirib yoki parallel bajarish. Mulohazalarni baholang.
\begin{enumerate}[label=\Alph*),leftmargin=9mm]
\item faqat II
\item I va II
\item faqat I
\item I ham, II ham emas
\end{enumerate}
\begin{tcolorbox}[enhanced,breakable,colback=teal!4,colframe=teal!60!black,boxrule=0.45pt,arc=1.2mm,left=2.7mm,right=2.7mm,top=1.8mm,bottom=1.8mm,before skip=2mm,after skip=2mm]
{\bfseries\sffamily\color{teal}To'g'ri javob:} \textbf{C) faqat I}
\par\smallskip
{\bfseries Izoh:} faqat I. Birinchi va ikkinchi mulohazaning ta'riflari alohida tekshirildi; natija: faqat I. Manba: ICT5, PDF 38-41; Qo'shimcha sahifa raqami (qaysi PDFga tegishli ekani aniqlanmagan; darslik sahifasi deb tasdiqlanmagan): 59
\end{tcolorbox}
\vspace{2mm}\hrule height 0.2pt\vspace{3mm}

% DB_ID 70742eb5-54ce-c80d-c124-4361692deaa6
\needspace{5\baselineskip}
{\bfseries\sffamily\color{ink} 347.} {scriptsizecolor{muted} Y1 | review | 70742eb5}
\par
Administrator bir xil kompyuter sozlamasini ellik marta takrorlashi kerak. Qaysi usul vaqt va bir xillik nuqtayi nazaridan eng qulay?
\begin{enumerate}[label=\Alph*),leftmargin=9mm]
\item Har safar sichqonchani boshqa oynaga ko'chirib, amallarni qo'lda qaytarish
\item CLI buyruqlarini skriptga yozib, bir xil ketma-ketlikni takrorlash
\item Kompyuter fon rasmini yangilab, barcha sozlamalarni qayta boshlash
\item Monitorning tasvir aniqligini o'zgartirib, buyruqlarni tezroq bosish
\end{enumerate}
\begin{tcolorbox}[enhanced,breakable,colback=teal!4,colframe=teal!60!black,boxrule=0.45pt,arc=1.2mm,left=2.7mm,right=2.7mm,top=1.8mm,bottom=1.8mm,before skip=2mm,after skip=2mm]
{\bfseries\sffamily\color{teal}To'g'ri javob:} \textbf{B) CLI buyruqlarini skriptga yozib, bir xil ketma-ketlikni takrorlash}
\par\smallskip
{\bfseries Izoh:} ICT 10-sinf CLI ni matnli buyruq kiritish usuli sifatida tavsiflaydi. Bir xil buyruqlarni skript bilan takrorlash umumiy hisoblash amaliyotidagi qo'llash xulosasidir. [Manba: ICT 10-sinf (2021), PDF 24-bet.]
\end{tcolorbox}
\vspace{2mm}\hrule height 0.2pt\vspace{3mm}

% DB_ID 7f4ce9b7-ebe0-14ac-1dcd-de6fce8b0465
\needspace{5\baselineskip}
{\bfseries\sffamily\color{ink} 348.} {scriptsizecolor{muted} Y1 | review | 7f4ce9b7}
\par
Operatsion tizimning vazifalari mavzusida: I. operatsion tizim - resurslarni boshqarib, foydalanuvchi va ilovalarga xizmat ko'rsatuvchi tizimli dastur. II. jarayon - bajarilayotgan dastur nusxasi. Mulohazalarni baholang.
\begin{enumerate}[label=\Alph*),leftmargin=9mm]
\item I va II
\item faqat I
\item faqat II
\item I ham, II ham emas
\end{enumerate}
\begin{tcolorbox}[enhanced,breakable,colback=teal!4,colframe=teal!60!black,boxrule=0.45pt,arc=1.2mm,left=2.7mm,right=2.7mm,top=1.8mm,bottom=1.8mm,before skip=2mm,after skip=2mm]
{\bfseries\sffamily\color{teal}To'g'ri javob:} \textbf{A) I va II}
\par\smallskip
{\bfseries Izoh:} I va II. Birinchi va ikkinchi mulohazaning ta'riflari alohida tekshirildi; natija: I va II. Manba: CAM9, PDF 12-15; Qo'shimcha sahifa raqami (qaysi PDFga tegishli ekani aniqlanmagan; darslik sahifasi deb tasdiqlanmagan): 57
\end{tcolorbox}
\vspace{2mm}\hrule height 0.2pt\vspace{3mm}

% DB_ID 855378da-a098-af1b-356d-f990bb9e6bdc
\needspace{5\baselineskip}
{\bfseries\sffamily\color{ink} 349.} {scriptsizecolor{muted} Y1 | review | 855378da}
\par
Kompyuter yoqilgach firmware qaysi ishni keyingi bosqichga topshiradi?
\begin{enumerate}[label=\Alph*),leftmargin=9mm]
\item foydalanuvchi hujjatini chop etishni
\item bootloader orqali OT yadrosini yuklashni
\item elektron jadval formulasini hisoblashni
\item rasm formatini siqishni
\end{enumerate}
\begin{tcolorbox}[enhanced,breakable,colback=teal!4,colframe=teal!60!black,boxrule=0.45pt,arc=1.2mm,left=2.7mm,right=2.7mm,top=1.8mm,bottom=1.8mm,before skip=2mm,after skip=2mm]
{\bfseries\sffamily\color{teal}To'g'ri javob:} \textbf{B) bootloader orqali OT yadrosini yuklashni}
\par\smallskip
{\bfseries Izoh:} bootloader orqali OT yadrosini yuklashni. Firmware boshlang'ich tekshiruvdan so'ng yuklash zanjirini ishga tushiradi. Manba: CAM9, PDF 12-15; Qo'shimcha sahifa raqami (qaysi PDFga tegishli ekani aniqlanmagan; darslik sahifasi deb tasdiqlanmagan): 61
\end{tcolorbox}
\vspace{2mm}\hrule height 0.2pt\vspace{3mm}

% DB_ID 9a96695c-1b37-e640-442b-df5d18a0d457
\needspace{5\baselineskip}
{\bfseries\sffamily\color{ink} 350.} {scriptsizecolor{muted} Y1 | review | 9a96695c}
\par
Administrator buyruqlar qatorida, boshqa xodim esa oynali menyuda ishlamoqda. Bu ikki boshqaruv usulini ajrating.
\begin{enumerate}[label=\Alph*),leftmargin=9mm]
\item Matnli buyruq va grafik menyu protsessorning ikkita mustaqil apparat qismini ifodalaydi
\item Matnli interfeys GUI, grafik oynalar esa CLI deb yuritiladi va ikkalasi ham alohida disk turidir
\item Matnli buyruqlar orqali boshqarish CLI, grafik oynalar orqali boshqarish GUI deb ataladi
\item Buyruq qatori ROM bilan, oynali menyu esa faqat elektr quvvati bloki bilan bog'liq
\end{enumerate}
\begin{tcolorbox}[enhanced,breakable,colback=teal!4,colframe=teal!60!black,boxrule=0.45pt,arc=1.2mm,left=2.7mm,right=2.7mm,top=1.8mm,bottom=1.8mm,before skip=2mm,after skip=2mm]
{\bfseries\sffamily\color{teal}To'g'ri javob:} \textbf{C) Matnli buyruqlar orqali boshqarish CLI, grafik oynalar orqali boshqarish GUI deb ataladi}
\par\smallskip
{\bfseries Izoh:} CLI matnli buyruq, GUI esa grafik elementlar orqali foydalanuvchi muloqotini ta'minlaydi. [Manba: ICT 10-sinf PDF 24.]
\end{tcolorbox}
\vspace{2mm}\hrule height 0.2pt\vspace{3mm}

% DB_ID 9de45a73-7fa4-4926-761a-e9f3a511553b
\needspace{5\baselineskip}
{\bfseries\sffamily\color{ink} 351.} {scriptsizecolor{muted} Y1 | review | 9de45a73}
\par
Operatsion tizim turlari va interfeyslar mavzusida: I. embedded OT - maxsus qurilma vazifasiga moslashtirilgan OT. II. GUI - oyna, ikona, menyu va ko'rsatkichga tayangan grafik interfeys. Mulohazalarni baholang.
\begin{enumerate}[label=\Alph*),leftmargin=9mm]
\item faqat I
\item faqat II
\item I va II
\item I ham, II ham emas
\end{enumerate}
\begin{tcolorbox}[enhanced,breakable,colback=teal!4,colframe=teal!60!black,boxrule=0.45pt,arc=1.2mm,left=2.7mm,right=2.7mm,top=1.8mm,bottom=1.8mm,before skip=2mm,after skip=2mm]
{\bfseries\sffamily\color{teal}To'g'ri javob:} \textbf{C) I va II}
\par\smallskip
{\bfseries Izoh:} I va II. Birinchi va ikkinchi mulohazaning ta'riflari alohida tekshirildi; natija: I va II. Manba: ICT11, PDF 244-245; Qo'shimcha sahifa raqami (qaysi PDFga tegishli ekani aniqlanmagan; darslik sahifasi deb tasdiqlanmagan): 67
\end{tcolorbox}
\vspace{2mm}\hrule height 0.2pt\vspace{3mm}

% DB_ID b664ec0c-8023-5f4d-a135-a186a96d9663
\needspace{5\baselineskip}
{\bfseries\sffamily\color{ink} 352.} {scriptsizecolor{muted} Y1 | review | b664ec0c}
\par
Operatsion tizim turlari va interfeyslar mavzusida: I. GUI - oyna, ikona, menyu va ko'rsatkichga tayangan grafik interfeys. II. CLI - matnli buyruqlar orqali boshqariladigan interfeys. Mulohazalarni baholang.
\begin{enumerate}[label=\Alph*),leftmargin=9mm]
\item faqat I
\item faqat II
\item I va II
\item I ham, II ham emas
\end{enumerate}
\begin{tcolorbox}[enhanced,breakable,colback=teal!4,colframe=teal!60!black,boxrule=0.45pt,arc=1.2mm,left=2.7mm,right=2.7mm,top=1.8mm,bottom=1.8mm,before skip=2mm,after skip=2mm]
{\bfseries\sffamily\color{teal}To'g'ri javob:} \textbf{C) I va II}
\par\smallskip
{\bfseries Izoh:} I va II. Birinchi va ikkinchi mulohazaning ta'riflari alohida tekshirildi; natija: I va II. Manba: CAM9, PDF 13-15; Qo'shimcha sahifa raqami (qaysi PDFga tegishli ekani aniqlanmagan; darslik sahifasi deb tasdiqlanmagan): 63
\end{tcolorbox}
\vspace{2mm}\hrule height 0.2pt\vspace{3mm}

% DB_ID bc494e7f-f30f-261b-d9dd-48472af46776
\needspace{5\baselineskip}
{\bfseries\sffamily\color{ink} 353.} {scriptsizecolor{muted} Y1 | review | bc494e7f}
\par
Operatsion tizim turlari va interfeyslar bo'yicha qatorlar ichidan "mobil OT" noto'g'ri izohlanganini aniqlang.
\begin{enumerate}[label=\Alph*),leftmargin=9mm]
\item mobil OT - oyna, ikona, menyu va ko'rsatkichga tayangan grafik interfeys
\item CLI - matnli buyruqlar orqali boshqariladigan interfeys
\item server OT - tarmoq xizmatlari va ko'p mijoz resurslarini boshqarishga mo'ljallangan OT
\item embedded OT - maxsus qurilma vazifasiga moslashtirilgan OT
\end{enumerate}
\begin{tcolorbox}[enhanced,breakable,colback=teal!4,colframe=teal!60!black,boxrule=0.45pt,arc=1.2mm,left=2.7mm,right=2.7mm,top=1.8mm,bottom=1.8mm,before skip=2mm,after skip=2mm]
{\bfseries\sffamily\color{teal}To'g'ri javob:} \textbf{A) mobil OT - oyna, ikona, menyu va ko'rsatkichga tayangan grafik interfeys}
\par\smallskip
{\bfseries Izoh:} mobil OT - oyna, ikona, menyu va ko'rsatkichga tayangan grafik interfeys. mobil OT uchun to'g'ri ta'rif: sensorli, portativ va energiya cheklangan qurilmaga mos OT. Manba: CAM7, PDF 140-146; Qo'shimcha sahifa raqami (qaysi PDFga tegishli ekani aniqlanmagan; darslik sahifasi deb tasdiqlanmagan): 66
\end{tcolorbox}
\vspace{2mm}\hrule height 0.2pt\vspace{3mm}

% DB_ID e945642f-43e1-d39e-5507-3328b1f1843e
\needspace{5\baselineskip}
{\bfseries\sffamily\color{ink} 354.} {scriptsizecolor{muted} Y1 | review | e945642f}
\par
Bir vaqtning o'zida brauzer, jadval va audio dasturi ochilgan. Protsessor hamda xotira resurslari foydalanishini boshqarishda qaysi dastur asosiy?
\begin{enumerate}[label=\Alph*),leftmargin=9mm]
\item Operatsion tizim jarayonlar va kompyuter resurslarini muvofiqlashtiradi
\item Brauzer barcha jarayonlarning apparat darajasidagi boshqaruvini yagona xizmat sifatida amalga oshiradi
\item Elektron jadval dasturi kompyuterning barcha resurslarini operatsion tizimsiz o'zi boshqarib boradi
\item Audio pleyer boshqa dasturlar ishlash tartibini majburiy belgilaydi va barcha drayverlarni yangilaydi
\end{enumerate}
\begin{tcolorbox}[enhanced,breakable,colback=teal!4,colframe=teal!60!black,boxrule=0.45pt,arc=1.2mm,left=2.7mm,right=2.7mm,top=1.8mm,bottom=1.8mm,before skip=2mm,after skip=2mm]
{\bfseries\sffamily\color{teal}To'g'ri javob:} \textbf{A) Operatsion tizim jarayonlar va kompyuter resurslarini muvofiqlashtiradi}
\par\smallskip
{\bfseries Izoh:} OT dastur va qurilma resurslarini boshqaradigan tizimli dasturiy qatlamdir. [Manba: ICT 10-sinf PDF 24-25.]
\end{tcolorbox}
\vspace{2mm}\hrule height 0.2pt\vspace{3mm}

% DB_ID 0652bb4a-3d1b-b965-ae82-6d631ded5cf7
\needspace{5\baselineskip}
{\bfseries\sffamily\color{ink} 355.} {scriptsizecolor{muted} Y1 | review | 0652bb4a}
\par
Operatsion tizimning vazifalari mavzusida: I. jarayon - bir nechta vazifani navbatlashtirib yoki parallel bajarish. II. operatsion tizim - resurslarni boshqarib, foydalanuvchi va ilovalarga xizmat ko'rsatuvchi tizimli dastur. Mulohazalarni baholang.
\begin{enumerate}[label=\Alph*),leftmargin=9mm]
\item faqat I
\item faqat II
\item I va II
\item I ham, II ham emas
\end{enumerate}
\begin{tcolorbox}[enhanced,breakable,colback=teal!4,colframe=teal!60!black,boxrule=0.45pt,arc=1.2mm,left=2.7mm,right=2.7mm,top=1.8mm,bottom=1.8mm,before skip=2mm,after skip=2mm]
{\bfseries\sffamily\color{teal}To'g'ri javob:} \textbf{B) faqat II}
\par\smallskip
{\bfseries Izoh:} faqat II. Birinchi va ikkinchi mulohazaning ta'riflari alohida tekshirildi; natija: faqat II. Manba: CAM10/11, PDF 43-47; Qo'shimcha sahifa raqami (qaysi PDFga tegishli ekani aniqlanmagan; darslik sahifasi deb tasdiqlanmagan): 58
\end{tcolorbox}
\vspace{2mm}\hrule height 0.2pt\vspace{3mm}

% DB_ID 1cea25ea-3cf1-e7a4-a222-761bd0a19e13
\needspace{5\baselineskip}
{\bfseries\sffamily\color{ink} 356.} {scriptsizecolor{muted} Y1 | review | 1cea25ea}
\par
Operatsion tizim turlari va interfeyslar mavzusida: I. CLI - tarmoq xizmatlari va ko'p mijoz resurslarini boshqarishga mo'ljallangan OT. II. GUI - oyna, ikona, menyu va ko'rsatkichga tayangan grafik interfeys. Mulohazalarni baholang.
\begin{enumerate}[label=\Alph*),leftmargin=9mm]
\item faqat I
\item I va II
\item I ham, II ham emas
\item faqat II
\end{enumerate}
\begin{tcolorbox}[enhanced,breakable,colback=teal!4,colframe=teal!60!black,boxrule=0.45pt,arc=1.2mm,left=2.7mm,right=2.7mm,top=1.8mm,bottom=1.8mm,before skip=2mm,after skip=2mm]
{\bfseries\sffamily\color{teal}To'g'ri javob:} \textbf{D) faqat II}
\par\smallskip
{\bfseries Izoh:} faqat II. Birinchi va ikkinchi mulohazaning ta'riflari alohida tekshirildi; natija: faqat II. Manba: CAM10/11, PDF 44-48; Qo'shimcha sahifa raqami (qaysi PDFga tegishli ekani aniqlanmagan; darslik sahifasi deb tasdiqlanmagan): 64
\end{tcolorbox}
\vspace{2mm}\hrule height 0.2pt\vspace{3mm}

% DB_ID 429c83be-fd4d-ce72-6fe0-98f66fe91a89
\needspace{5\baselineskip}
{\bfseries\sffamily\color{ink} 357.} {scriptsizecolor{muted} Y1 | review | 429c83be}
\par
Bir xodim operatsion tizim bilan matnli buyruqlar yordamida ishlaydi, boshqasi esa oynalar va menyulardan foydalanadi. Bu farqni qanday izohlash to'g'ri?
\begin{enumerate}[label=\Alph*),leftmargin=9mm]
\item CLI va GUI kompyuterning bir-biriga mos kelmaydigan ikki apparat arxitekturasini bildiradi
\item CLI matnli buyruqlar, GUI esa grafik elementlar orqali boshqarish interfeysidir
\item GUI faqat ko'p foydalanuvchili serverlarga, CLI esa shaxsiy kompyuterlarga mo'ljallangan
\item CLI grafikli menyularga asoslanadi, GUI esa maxsus buyruq satrini talab qiladi
\end{enumerate}
\begin{tcolorbox}[enhanced,breakable,colback=teal!4,colframe=teal!60!black,boxrule=0.45pt,arc=1.2mm,left=2.7mm,right=2.7mm,top=1.8mm,bottom=1.8mm,before skip=2mm,after skip=2mm]
{\bfseries\sffamily\color{teal}To'g'ri javob:} \textbf{B) CLI matnli buyruqlar, GUI esa grafik elementlar orqali boshqarish interfeysidir}
\par\smallskip
{\bfseries Izoh:} CLI - buyruqlar qatori interfeysi, GUI esa grafik foydalanuvchi interfeysi. Ular apparat turi yoki alohida OT mahsuloti emas, boshqaruv usullaridir. Manba: ICT 10-sinf, PDF 25
\end{tcolorbox}
\vspace{2mm}\hrule height 0.2pt\vspace{3mm}

% DB_ID 42fbdd00-1821-5769-a361-a6adc0ae8ebc
\needspace{5\baselineskip}
{\bfseries\sffamily\color{ink} 358.} {scriptsizecolor{muted} Y1 | review | 42fbdd00}
\par
Texnik xodim fayllarni tartibga solish qoidasi va kompyuter resurslarini boshqarish dasturini ajratmoqchi. Qaysi fikrlar juftligi to'g'ri?
\begin{enumerate}[label=\Alph*),leftmargin=9mm]
\item Fayl tizimi - foydalanuvchi ilovasi; OT - printerdan chiqariladigan sahifa, faqat qog'oz bilan ishlaydi
\item Fayl tizimi - xotira mikrosxemasi; OT - ma'lumot saqlovchi magnit lenta
\item Fayl tizimi - alohida protsessor; OT - faqat shu protsessorda ishlaydigan dastur
\item Fayl tizimi - saqlashni tashkil etish qoidasi; OT - tizim resurslarini boshqaradi
\end{enumerate}
\begin{tcolorbox}[enhanced,breakable,colback=teal!4,colframe=teal!60!black,boxrule=0.45pt,arc=1.2mm,left=2.7mm,right=2.7mm,top=1.8mm,bottom=1.8mm,before skip=2mm,after skip=2mm]
{\bfseries\sffamily\color{teal}To'g'ri javob:} \textbf{D) Fayl tizimi - saqlashni tashkil etish qoidasi; OT - tizim resurslarini boshqaradi}
\par\smallskip
{\bfseries Izoh:} OT resurslarni boshqaradi, fayl tizimi esa fayl va papkalarning saqlash va tuzilish usulini belgilaydi. Ikkalasi bir xil tushuncha emas. [Manba: ICT 10-sinf (2021), PDF 24-26.]
\end{tcolorbox}
\vspace{2mm}\hrule height 0.2pt\vspace{3mm}

% DB_ID ab7f4ee2-25e6-0816-46e8-25fa23bafdda
\needspace{5\baselineskip}
{\bfseries\sffamily\color{ink} 359.} {scriptsizecolor{muted} Y1 | review | ab7f4ee2}
\par
Printer o'zining sinov sahifasini chop etadi va USB kabeli ishlaydi, lekin kompyuter "mos drayver topilmadi" xabarini beradi. Tekshiruvni qaysi yo'nalishdan boshlash maqsadga muvofiq?
\begin{enumerate}[label=\Alph*),leftmargin=9mm]
\item Kompyuterning displey drayveri va ekran rezolyutsiyasi sozlamalarini avval tekshirish
\item Printerdagi qog'oz turini va chop etish sifatini tizim drayveriga aloqasiz qayta sozlash
\item Hujjat yaratish dasturidagi barcha sahifa chegaralarini o'zgartirib, qayta yuborish
\item Printer modeliga mos qurilma drayverini o'rnatish va OT bilan mosligini tekshirish
\end{enumerate}
\begin{tcolorbox}[enhanced,breakable,colback=teal!4,colframe=teal!60!black,boxrule=0.45pt,arc=1.2mm,left=2.7mm,right=2.7mm,top=1.8mm,bottom=1.8mm,before skip=2mm,after skip=2mm]
{\bfseries\sffamily\color{teal}To'g'ri javob:} \textbf{D) Printer modeliga mos qurilma drayverini o'rnatish va OT bilan mosligini tekshirish}
\par\smallskip
{\bfseries Izoh:} Darslik qurilma drayverini printer kabi yordamchi qurilmalarni boshqaruvchi dastur deb izohlaydi. Drayver topilmaganida mos drayverni tekshirish oqilona birinchi qadamdir. Manba: Cambridge+ 9-sinf, PDF 12
\end{tcolorbox}
\vspace{2mm}\hrule height 0.2pt\vspace{3mm}
\clearpage
\section{Fayl, papka va fayl tizimi bilan ishlash}
\subsection{Nashr qilingan testlar}

% DB_ID 34704322-9832-58a2-8f08-cd9521527331
\needspace{5\baselineskip}
{\bfseries\sffamily\color{ink} 360.} {scriptsizecolor{muted} Y1 | published | 34704322}
\par
\texttt{Informatika} so'zi oddiy matn fayliga UTF-8 formatida BOM va satr oxiri belgisisiz, shuningdek MS Word \texttt{.doc} hujjatiga saqlandi. Fayllar hajmi haqidagi qaysi fikr to'g'ri?
\begin{enumerate}[label=\Alph*),leftmargin=9mm]
\item Oddiy matn fayli 11 bayt, MS Word fayli bir necha KB
\item Ikkala fayl ham 11 bayt
\item Oddiy matn fayli bir necha KB, MS Word fayli 11 bayt
\item Ikkala fayl ham bir necha KB
\end{enumerate}
\begin{tcolorbox}[enhanced,breakable,colback=teal!4,colframe=teal!60!black,boxrule=0.45pt,arc=1.2mm,left=2.7mm,right=2.7mm,top=1.8mm,bottom=1.8mm,before skip=2mm,after skip=2mm]
{\bfseries\sffamily\color{teal}To'g'ri javob:} \textbf{A) Oddiy matn fayli 11 bayt, MS Word fayli bir necha KB}
\par\smallskip
{\bfseries Izoh:} \texttt{Informatika} 11 ta ASCII belgidan iborat va UTF-8da har biri 1 bayt: jami 11 bayt. Word hujjati matndan tashqari format va tuzilma metama'lumotlarini ham saqlagani uchun hajmi bir necha KB bo'ladi.
\end{tcolorbox}
\vspace{2mm}\hrule height 0.2pt\vspace{3mm}

% DB_ID 39174c74-b13c-48e1-b9d9-b029d48e5480
\needspace{5\baselineskip}
{\bfseries\sffamily\color{ink} 361.} {scriptsizecolor{muted} Y1 | published | 39174c74}
\par
Papka (folder)ning vazifasi qaysi?
\begin{enumerate}[label=\Alph*),leftmargin=9mm]
\item Fayl va boshqa papkalarni tartibli saqlash
\item CPU tezligini oshirish
\item Internet tezligini o'lchash
\item Printer siyohini boshqarish
\end{enumerate}
\begin{tcolorbox}[enhanced,breakable,colback=teal!4,colframe=teal!60!black,boxrule=0.45pt,arc=1.2mm,left=2.7mm,right=2.7mm,top=1.8mm,bottom=1.8mm,before skip=2mm,after skip=2mm]
{\bfseries\sffamily\color{teal}To'g'ri javob:} \textbf{A) Fayl va boshqa papkalarni tartibli saqlash}
\par\smallskip
{\bfseries Izoh:} Papka fayl va ichki papkalarni mantiqiy guruhlash uchun ishlatiladi.
\end{tcolorbox}
\vspace{2mm}\hrule height 0.2pt\vspace{3mm}

% DB_ID 5e07f6ec-1fb7-4c38-a20c-45f0150fab27
\needspace{5\baselineskip}
{\bfseries\sffamily\color{ink} 362.} {scriptsizecolor{muted} Y1 | published | 5e07f6ec}
\par
Fayl kengaytmasi odatda nimani ko'rsatishga yordam beradi?
\begin{enumerate}[label=\Alph*),leftmargin=9mm]
\item Faylning turi yoki formati
\item Kompyuterning IP manzilini
\item RAM hajmini
\item Protsessor chastotasini
\end{enumerate}
\begin{tcolorbox}[enhanced,breakable,colback=teal!4,colframe=teal!60!black,boxrule=0.45pt,arc=1.2mm,left=2.7mm,right=2.7mm,top=1.8mm,bottom=1.8mm,before skip=2mm,after skip=2mm]
{\bfseries\sffamily\color{teal}To'g'ri javob:} \textbf{A) Faylning turi yoki formati}
\par\smallskip
{\bfseries Izoh:} .docx, .png, .pdf kabi kengaytmalar fayl turini aniqlashga yordam beradi.
\end{tcolorbox}
\vspace{2mm}\hrule height 0.2pt\vspace{3mm}

% DB_ID 7009f45d-d659-4873-9392-5495df97168b
\needspace{5\baselineskip}
{\bfseries\sffamily\color{ink} 363.} {scriptsizecolor{muted} Y1 | published | 7009f45d}
\par
Fayl turlarini mos kengaytmalar bilan juftlang: 1) audio; 2) video; 3) rastr rasm. a) \texttt{.jpg}, \texttt{.png}, \texttt{.bmp}; b) \texttt{.mp4}, \texttt{.avi}, \texttt{.mov}; c) \texttt{.aac}, \texttt{.flac}, \texttt{.gif}; d) \texttt{.mp3}, \texttt{.wav}, \texttt{.ogg}.
\begin{enumerate}[label=\Alph*),leftmargin=9mm]
\item 1-d; 2-b; 3-a
\item 1-c; 2-b; 3-a
\item 1-d; 2-b; 3-c
\item 1-b; 2-d; 3-a
\end{enumerate}
\begin{tcolorbox}[enhanced,breakable,colback=teal!4,colframe=teal!60!black,boxrule=0.45pt,arc=1.2mm,left=2.7mm,right=2.7mm,top=1.8mm,bottom=1.8mm,before skip=2mm,after skip=2mm]
{\bfseries\sffamily\color{teal}To'g'ri javob:} \textbf{A) 1-d; 2-b; 3-a}
\par\smallskip
{\bfseries Izoh:} Audio uchun d, video uchun b, rastr rasm uchun a mos. c to'plamida \texttt{.gif} audio formati emas.
\end{tcolorbox}
\vspace{2mm}\hrule height 0.2pt\vspace{3mm}

% DB_ID 9331e65e-7666-4637-bbca-a583dcce3ac4
\needspace{5\baselineskip}
{\bfseries\sffamily\color{ink} 364.} {scriptsizecolor{muted} Y1 | published | 9331e65e}
\par
\texttt{ege.htm} fayli yangi yaratilgan podkatalogga ko'chirilgach, uning to'liq manzili \texttt{A:\textbackslash{}Exam\textbackslash{}Information\textbackslash{}ECDL\textbackslash{}EXAM\textbackslash{}ege.htm} bo'ldi. Yangi podkatalog nomi qaysi?
\begin{enumerate}[label=\Alph*),leftmargin=9mm]
\item EGE
\item Information
\item ECDL
\item EXAM
\end{enumerate}
\begin{tcolorbox}[enhanced,breakable,colback=teal!4,colframe=teal!60!black,boxrule=0.45pt,arc=1.2mm,left=2.7mm,right=2.7mm,top=1.8mm,bottom=1.8mm,before skip=2mm,after skip=2mm]
{\bfseries\sffamily\color{teal}To'g'ri javob:} \textbf{D) EXAM}
\par\smallskip
{\bfseries Izoh:} Fayl nomidan bevosita oldingi yo'l qismi uning ota papkasi: EXAM.
\end{tcolorbox}
\vspace{2mm}\hrule height 0.2pt\vspace{3mm}

% DB_ID d808c2bb-68d0-5626-b80a-dcf322c00672
\needspace{5\baselineskip}
{\bfseries\sffamily\color{ink} 365.} {scriptsizecolor{muted} Y1 | published | d808c2bb}
\par
Dasturdagi ma'lumotni boshqa dastur foydalanishi mumkin bo'lgan formatga chiqarish jarayoni nima deyiladi?
\begin{enumerate}[label=\Alph*),leftmargin=9mm]
\item Import
\item Eksport
\item Copy
\item More
\end{enumerate}
\begin{tcolorbox}[enhanced,breakable,colback=teal!4,colframe=teal!60!black,boxrule=0.45pt,arc=1.2mm,left=2.7mm,right=2.7mm,top=1.8mm,bottom=1.8mm,before skip=2mm,after skip=2mm]
{\bfseries\sffamily\color{teal}To'g'ri javob:} \textbf{B) Eksport}
\par\smallskip
{\bfseries Izoh:} Eksport ma'lumotni tashqi fayl yoki boshqa dastur qabul qiladigan formatga chiqaradi. Import esa aksincha, tashqi ma'lumotni dastur ichiga olib kiradi.
\end{tcolorbox}
\vspace{2mm}\hrule height 0.2pt\vspace{3mm}

% DB_ID 23d73e3b-1b5d-5262-9e2e-a7b02b6ddd61
\needspace{5\baselineskip}
{\bfseries\sffamily\color{ink} 366.} {scriptsizecolor{muted} Y1 | published | 23d73e3b}
\par
Windowsda \texttt{?} aynan bitta belgini, \texttt{*} esa istalgancha (hatto nol) belgini bildiradi. \texttt{template.sys}, \texttt{jumpland.xls} va \texttt{campmap.acs} fayllarining uchalasiga ham mos keladigan maskani toping.
\begin{enumerate}[label=\Alph*),leftmargin=9mm]
\item \texttt{??mp?a*.?*}
\item \texttt{*m?p*a??.??s}
\item \texttt{????p*a?.?s*}
\item \texttt{?*m*p?a??.*s}
\end{enumerate}
\begin{tcolorbox}[enhanced,breakable,colback=teal!4,colframe=teal!60!black,boxrule=0.45pt,arc=1.2mm,left=2.7mm,right=2.7mm,top=1.8mm,bottom=1.8mm,before skip=2mm,after skip=2mm]
{\bfseries\sffamily\color{teal}To'g'ri javob:} \textbf{A) \texttt{??mp?a*.?*}}
\par\smallskip
{\bfseries Izoh:} \texttt{??mp?a*.?*} maskasida dastlabki ikki belgi mos ravishda \texttt{te}, \texttt{ju}, \texttt{ca}; keyin \texttt{mp}, bitta belgi va \texttt{a} keladi. Qolgan nom hamda kengaytma qismlarini \texttt{*} qoplaydi. Qolgan maskalardan kamida bittasi \texttt{campmap.acs} ga mos kelmaydi.
\end{tcolorbox}
\vspace{2mm}\hrule height 0.2pt\vspace{3mm}

% DB_ID 4563968e-0256-489c-8a7d-16546483371a
\needspace{5\baselineskip}
{\bfseries\sffamily\color{ink} 367.} {scriptsizecolor{muted} Y1 | published | 4563968e}
\par
Windows fayl maskasida \texttt{?} - aynan bitta belgi, \texttt{*} - istalgan uzunlikdagi belgilar ketma-ketligi. \texttt{word.doc}, \texttt{load.dxf}, \texttt{password.doc}, \texttt{pound.dxf} fayllarining barchasiga mos maskani toping.
\begin{enumerate}[label=\Alph*),leftmargin=9mm]
\item \texttt{*o*.d??}
\item \texttt{?o?.d**}
\item \texttt{*o?d.d*}
\item \texttt{?o*d?.d?}
\end{enumerate}
\begin{tcolorbox}[enhanced,breakable,colback=teal!4,colframe=teal!60!black,boxrule=0.45pt,arc=1.2mm,left=2.7mm,right=2.7mm,top=1.8mm,bottom=1.8mm,before skip=2mm,after skip=2mm]
{\bfseries\sffamily\color{teal}To'g'ri javob:} \textbf{A) \texttt{*o*.d??}}
\par\smallskip
{\bfseries Izoh:} Barcha asosiy nomlarda \texttt{o} bor; barcha kengaytmalar \texttt{d} bilan boshlanib, undan keyin aynan ikki belgi keladi. Shuning uchun \texttt{*o*.d??} mos.
\end{tcolorbox}
\vspace{2mm}\hrule height 0.2pt\vspace{3mm}

% DB_ID ad1b1ab5-732e-4096-93c0-8ff3bd14fb1e
\needspace{5\baselineskip}
{\bfseries\sffamily\color{ink} 368.} {scriptsizecolor{muted} Y1 | published | ad1b1ab5}
\par
Windowsda \texttt{C:\textbackslash{}Maktab\textbackslash{}9A\textbackslash{}hisobot.docx} yo'li ko'rsatilgan. Ushbu yo'ldagi \texttt{hisobot.docx} qismi nimani ifodalaydi?
\begin{enumerate}[label=\Alph*),leftmargin=9mm]
\item Fayl nomi va kengaytmasi
\item Disk nomi va harfi
\item Papka nomi va manzili
\item Katalog tuzilishi va yo'li
\end{enumerate}
\begin{tcolorbox}[enhanced,breakable,colback=teal!4,colframe=teal!60!black,boxrule=0.45pt,arc=1.2mm,left=2.7mm,right=2.7mm,top=1.8mm,bottom=1.8mm,before skip=2mm,after skip=2mm]
{\bfseries\sffamily\color{teal}To'g'ri javob:} \textbf{A) Fayl nomi va kengaytmasi}
\par\smallskip
{\bfseries Izoh:} \texttt{hisobot} fayl nomi, \texttt{.docx} kengaytmasi; undan oldingi \texttt{Maktab} va \texttt{9A} katalog nomlaridir. \texttt{C:} esa disk belgisini bildiradi.
\end{tcolorbox}
\vspace{2mm}\hrule height 0.2pt\vspace{3mm}

% DB_ID c1e99d0c-890f-5bb0-a162-4ddf19d8494e
\needspace{5\baselineskip}
{\bfseries\sffamily\color{ink} 369.} {scriptsizecolor{muted} Y1 | published | c1e99d0c}
\par
Papkada \texttt{print1.pdf}, \texttt{sprint.pdf}, \texttt{print22.pdf}, \texttt{reprint.pdf} va \texttt{picture.jpg} fayllari bor. Nomi ichida \texttt{print} qatnashgan va kengaytmasi aynan \texttt{.pdf} bo'lgan fayllarni tanlaydigan maskani toping.
\begin{enumerate}[label=\Alph*),leftmargin=9mm]
\item \texttt{*print*.pdf}
\item \texttt{?print.pdf}
\item \texttt{?print*.pdf}
\item \texttt{print*.pdf}
\end{enumerate}
\begin{tcolorbox}[enhanced,breakable,colback=teal!4,colframe=teal!60!black,boxrule=0.45pt,arc=1.2mm,left=2.7mm,right=2.7mm,top=1.8mm,bottom=1.8mm,before skip=2mm,after skip=2mm]
{\bfseries\sffamily\color{teal}To'g'ri javob:} \textbf{A) \texttt{*print*.pdf}}
\par\smallskip
{\bfseries Izoh:} Birinchi \texttt{*} \texttt{print} oldidagi, ikkinchisi undan keyingi istalgancha belgini qamraydi; \texttt{.pdf} kengaytma aynan PDF bo'lishini talab qiladi.
\end{tcolorbox}
\vspace{2mm}\hrule height 0.2pt\vspace{3mm}

% DB_ID 6172077b-bef1-4f42-b871-d80834fad1e7
\needspace{5\baselineskip}
{\bfseries\sffamily\color{ink} 370.} {scriptsizecolor{muted} Y1 | published | 6172077b}
\par
\texttt{D:\textbackslash{}Foto\textbackslash{}sayohat\textbackslash{}rasm.png} faylini topish uchun qaysi ma'lumot eng foydali?
\begin{enumerate}[label=\Alph*),leftmargin=9mm]
\item To'liq fayl yo'li
\item CPU modeli
\item Monitor o'lchami
\item Internet tezligi
\end{enumerate}
\begin{tcolorbox}[enhanced,breakable,colback=teal!4,colframe=teal!60!black,boxrule=0.45pt,arc=1.2mm,left=2.7mm,right=2.7mm,top=1.8mm,bottom=1.8mm,before skip=2mm,after skip=2mm]
{\bfseries\sffamily\color{teal}To'g'ri javob:} \textbf{A) To'liq fayl yo'li}
\par\smallskip
{\bfseries Izoh:} To'liq yo'l disk, papkalar va fayl nomini ko'rsatadi.
\end{tcolorbox}
\vspace{2mm}\hrule height 0.2pt\vspace{3mm}

% DB_ID 7ea17a0a-ced2-482f-8856-824a002a16e9
\needspace{5\baselineskip}
{\bfseries\sffamily\color{ink} 371.} {scriptsizecolor{muted} Y1 | published | 7ea17a0a}
\par
Windowsda \texttt{*e?i*e*.?m?} fayl qidirish maskasiga qaysi nom mos keladi?
\begin{enumerate}[label=\Alph*),leftmargin=9mm]
\item decline.cma
\item hemisphere.bmp
\item environment.tmp
\item incentive.cmt
\end{enumerate}
\begin{tcolorbox}[enhanced,breakable,colback=teal!4,colframe=teal!60!black,boxrule=0.45pt,arc=1.2mm,left=2.7mm,right=2.7mm,top=1.8mm,bottom=1.8mm,before skip=2mm,after skip=2mm]
{\bfseries\sffamily\color{teal}To'g'ri javob:} \textbf{B) hemisphere.bmp}
\par\smallskip
{\bfseries Izoh:} hemisphere nomi e-m-i ketma-ketligiga va .bmp kengaytmasi ?m?ga mos.
\end{tcolorbox}
\vspace{2mm}\hrule height 0.2pt\vspace{3mm}

% DB_ID a28ef2f7-cfb8-4760-9cc1-20d557ae7245
\needspace{5\baselineskip}
{\bfseries\sffamily\color{ink} 372.} {scriptsizecolor{muted} Y1 | published | a28ef2f7}
\par
\texttt{C:\textbackslash{}Program Files\textbackslash{}PostgreSQL\textbackslash{}school.sql} manzilida moslang: 1) C: 2) Program Files 3) PostgreSQL 4) school.sql. a) joriy katalog b) diskning bosh katalogi c) fayl nomi d) bosh katalogning katalogosti e) faylning xususiy nomi.
\begin{enumerate}[label=\Alph*),leftmargin=9mm]
\item 1-b,2-a,3-d,4-e
\item 1-b,2-d,3-a,4-c
\item 1-b,2-d,3-e,4-c
\item 1-a,2-b,3-d,4-e
\end{enumerate}
\begin{tcolorbox}[enhanced,breakable,colback=teal!4,colframe=teal!60!black,boxrule=0.45pt,arc=1.2mm,left=2.7mm,right=2.7mm,top=1.8mm,bottom=1.8mm,before skip=2mm,after skip=2mm]
{\bfseries\sffamily\color{teal}To'g'ri javob:} \textbf{B) 1-b,2-d,3-a,4-c}
\par\smallskip
{\bfseries Izoh:} C: disk ildizi, Program Files katalogosti, PostgreSQL joriy katalog, school.sql fayl nomi.
\end{tcolorbox}
\vspace{2mm}\hrule height 0.2pt\vspace{3mm}

% DB_ID a9e5c25f-a539-43fa-aba4-ff62d38efd78
\needspace{5\baselineskip}
{\bfseries\sffamily\color{ink} 373.} {scriptsizecolor{muted} Y1 | published | a9e5c25f}
\par
O'quvchi \texttt{referat.docx} faylini \texttt{9A} papkasidan \texttt{Arxiv} papkasiga ko'chirdi. Faylning mazmuni o'zgarmadi. Qaysi amal bajarildi?
\begin{enumerate}[label=\Alph*),leftmargin=9mm]
\item Faylning joylashuvi o'zgartirildi
\item Fayl formati o'zgartirildi
\item Fayl shifrlanib qoldi
\item Fayl RAMga aylantirildi
\end{enumerate}
\begin{tcolorbox}[enhanced,breakable,colback=teal!4,colframe=teal!60!black,boxrule=0.45pt,arc=1.2mm,left=2.7mm,right=2.7mm,top=1.8mm,bottom=1.8mm,before skip=2mm,after skip=2mm]
{\bfseries\sffamily\color{teal}To'g'ri javob:} \textbf{A) Faylning joylashuvi o'zgartirildi}
\par\smallskip
{\bfseries Izoh:} Papka orasida ko'chirish fayl joylashuvini o'zgartiradi, mazmun va formatni emas.
\end{tcolorbox}
\vspace{2mm}\hrule height 0.2pt\vspace{3mm}

% DB_ID b2550fc2-d13e-4e02-9330-91e7766a33c9
\needspace{5\baselineskip}
{\bfseries\sffamily\color{ink} 374.} {scriptsizecolor{muted} Y3 | published | b2550fc2}
\par
Yangi loyiha fayllarini tartibli joylashtirish bosqichlarini tartiblang.
\begin{enumerate}[label=\Alph*),leftmargin=9mm]
\item Fayllarni mos ichki papkalarga saqlash yoki ko'chirish
\item Loyiha uchun asosiy papka yaratish
\item Yo'l va fayl nomlarini tekshirish
\item Kerakli ichki papkalarni nomlash
\end{enumerate}
\begin{tcolorbox}[enhanced,breakable,colback=teal!4,colframe=teal!60!black,boxrule=0.45pt,arc=1.2mm,left=2.7mm,right=2.7mm,top=1.8mm,bottom=1.8mm,before skip=2mm,after skip=2mm]
{\bfseries\sffamily\color{teal}To'g'ri javob:} \textbf{B - D - A - C}
\par\smallskip
{\bfseries Izoh:} Avval tuzilma yaratiladi, so'ng fayllar joylashtiriladi va natija tekshiriladi.
\end{tcolorbox}
\vspace{2mm}\hrule height 0.2pt\vspace{3mm}

% DB_ID c13a23c3-ad03-47b9-8435-d27061d17330
\needspace{5\baselineskip}
{\bfseries\sffamily\color{ink} 375.} {scriptsizecolor{muted} Y1 | published | c13a23c3}
\par
Foydalanuvchi \texttt{ma'ruza.pdf} faylining nomini \texttt{dars.pdf} ga o'zgartirdi. Qaysi xulosa to'g'ri?
\begin{enumerate}[label=\Alph*),leftmargin=9mm]
\item Nom o'zgardi, format .pdf bo'lib qoldi
\item Fayl Word formatiga aylandi
\item Fayl o'chirildi
\item Fayl papkaga aylandi
\end{enumerate}
\begin{tcolorbox}[enhanced,breakable,colback=teal!4,colframe=teal!60!black,boxrule=0.45pt,arc=1.2mm,left=2.7mm,right=2.7mm,top=1.8mm,bottom=1.8mm,before skip=2mm,after skip=2mm]
{\bfseries\sffamily\color{teal}To'g'ri javob:} \textbf{A) Nom o'zgardi, format .pdf bo'lib qoldi}
\par\smallskip
{\bfseries Izoh:} Faqat bazaviy nom o'zgargan, .pdf kengaytmasi saqlangan.
\end{tcolorbox}
\vspace{2mm}\hrule height 0.2pt\vspace{3mm}

% DB_ID db57760c-c003-484e-bd00-ad5803a32e61
\needspace{5\baselineskip}
{\bfseries\sffamily\color{ink} 376.} {scriptsizecolor{muted} Y1 | published | db57760c}
\par
Bir papkada \texttt{hisobot.docx} nomli fayl bor. Yana bir faylni aynan shu nom bilan shu papkada saqlashga urinishda tizim ogohlantirdi. Sabab nima?
\begin{enumerate}[label=\Alph*),leftmargin=9mm]
\item Bir papkada bir xil nomli ikki faylga odatda ruxsat berilmaydi
\item DOCX fayl faqat bitta marta yaratiladi
\item Papka faqat bitta fayl saqlaydi
\item RAM yetishmayapti
\end{enumerate}
\begin{tcolorbox}[enhanced,breakable,colback=teal!4,colframe=teal!60!black,boxrule=0.45pt,arc=1.2mm,left=2.7mm,right=2.7mm,top=1.8mm,bottom=1.8mm,before skip=2mm,after skip=2mm]
{\bfseries\sffamily\color{teal}To'g'ri javob:} \textbf{A) Bir papkada bir xil nomli ikki faylga odatda ruxsat berilmaydi}
\par\smallskip
{\bfseries Izoh:} Bir katalog ichida fayl nomi noyob bo'lishi kerak.
\end{tcolorbox}
\vspace{2mm}\hrule height 0.2pt\vspace{3mm}

% DB_ID ea8a9712-7718-421c-b089-0535a148ec60
\needspace{5\baselineskip}
{\bfseries\sffamily\color{ink} 377.} {scriptsizecolor{muted} Y1 | published | ea8a9712}
\par
Quyidagi juftliklarning qaysi ikkalasi ham grafik fayl kengaytmasi?
\begin{enumerate}[label=\Alph*),leftmargin=9mm]
\item .gif, .dll
\item .jpg, .png
\item .xls, .zip
\item .rar, .exe
\end{enumerate}
\begin{tcolorbox}[enhanced,breakable,colback=teal!4,colframe=teal!60!black,boxrule=0.45pt,arc=1.2mm,left=2.7mm,right=2.7mm,top=1.8mm,bottom=1.8mm,before skip=2mm,after skip=2mm]
{\bfseries\sffamily\color{teal}To'g'ri javob:} \textbf{B) .jpg, .png}
\par\smallskip
{\bfseries Izoh:} JPEG va PNG rasm formatlari.
\end{tcolorbox}
\vspace{2mm}\hrule height 0.2pt\vspace{3mm}
\subsection{Tekshiruvdagi testlar}

% DB_ID 0915997c-aa7f-2ad7-5a71-69e851b89e24
\needspace{5\baselineskip}
{\bfseries\sffamily\color{ink} 378.} {scriptsizecolor{muted} Y1 | review | 0915997c}
\par
Papka, katalog, yo'l va ierarxiya doirasida "mutlaq yo'l" atamasining eng aniq ta'rifini toping.
\begin{enumerate}[label=\Alph*),leftmargin=9mm]
\item fayl va quyi papkalarni guruhlovchi konteyner
\item fayl tizimi daraxtining eng yuqori katalogi
\item ildizdan boshlanadigan to'liq joylashuv
\item joriy obyektni bevosita o'z ichiga olgan katalog
\end{enumerate}
\begin{tcolorbox}[enhanced,breakable,colback=teal!4,colframe=teal!60!black,boxrule=0.45pt,arc=1.2mm,left=2.7mm,right=2.7mm,top=1.8mm,bottom=1.8mm,before skip=2mm,after skip=2mm]
{\bfseries\sffamily\color{teal}To'g'ri javob:} \textbf{C) ildizdan boshlanadigan to'liq joylashuv}
\par\smallskip
{\bfseries Izoh:} ildizdan boshlanadigan to'liq joylashuv. mutlaq yo'l - ildizdan boshlanadigan to'liq joylashuv. Manba: ICT6, PDF 116-124; Qo'shimcha sahifa raqami (qaysi PDFga tegishli ekani aniqlanmagan; darslik sahifasi deb tasdiqlanmagan): 93
\end{tcolorbox}
\vspace{2mm}\hrule height 0.2pt\vspace{3mm}

% DB_ID 0fefaa8f-4858-2355-bd28-3aaf20dcc342
\needspace{5\baselineskip}
{\bfseries\sffamily\color{ink} 379.} {scriptsizecolor{muted} Y1 | review | 0fefaa8f}
\par
Fayl, kengaytma, format va atribut doirasida "fayl" atamasining eng aniq ta'rifini toping.
\begin{enumerate}[label=\Alph*),leftmargin=9mm]
\item nom bilan saqlanadigan ma'lumotlar birligi
\item fayl nomining formatni ko'rsatuvchi oxirgi qismi
\item fayl ichidagi ma'lumot tuzilishi va kodlash qoidasi
\item faylning read-only yoki hidden kabi boshqaruv xususiyati
\end{enumerate}
\begin{tcolorbox}[enhanced,breakable,colback=teal!4,colframe=teal!60!black,boxrule=0.45pt,arc=1.2mm,left=2.7mm,right=2.7mm,top=1.8mm,bottom=1.8mm,before skip=2mm,after skip=2mm]
{\bfseries\sffamily\color{teal}To'g'ri javob:} \textbf{A) nom bilan saqlanadigan ma'lumotlar birligi}
\par\smallskip
{\bfseries Izoh:} nom bilan saqlanadigan ma'lumotlar birligi. fayl - nom bilan saqlanadigan ma'lumotlar birligi. Manba: ICT5, PDF 42-45; Qo'shimcha sahifa raqami (qaysi PDFga tegishli ekani aniqlanmagan; darslik sahifasi deb tasdiqlanmagan): 83
\end{tcolorbox}
\vspace{2mm}\hrule height 0.2pt\vspace{3mm}

% DB_ID 20d61082-2524-91d5-6a3e-47d0731e511c
\needspace{5\baselineskip}
{\bfseries\sffamily\color{ink} 380.} {scriptsizecolor{muted} Y1 | review | 20d61082}
\par
Fayl, kengaytma, format va atribut doirasida "kengaytma" atamasining eng aniq ta'rifini toping.
\begin{enumerate}[label=\Alph*),leftmargin=9mm]
\item nom bilan saqlanadigan ma'lumotlar birligi
\item fayl nomining formatni ko'rsatuvchi oxirgi qismi
\item fayl ichidagi ma'lumot tuzilishi va kodlash qoidasi
\item faylning read-only yoki hidden kabi boshqaruv xususiyati
\end{enumerate}
\begin{tcolorbox}[enhanced,breakable,colback=teal!4,colframe=teal!60!black,boxrule=0.45pt,arc=1.2mm,left=2.7mm,right=2.7mm,top=1.8mm,bottom=1.8mm,before skip=2mm,after skip=2mm]
{\bfseries\sffamily\color{teal}To'g'ri javob:} \textbf{B) fayl nomining formatni ko'rsatuvchi oxirgi qismi}
\par\smallskip
{\bfseries Izoh:} fayl nomining formatni ko'rsatuvchi oxirgi qismi. kengaytma - fayl nomining formatni ko'rsatuvchi oxirgi qismi. Manba: MAR, PDF 16-20; Qo'shimcha sahifa raqami (qaysi PDFga tegishli ekani aniqlanmagan; darslik sahifasi deb tasdiqlanmagan): 84
\end{tcolorbox}
\vspace{2mm}\hrule height 0.2pt\vspace{3mm}

% DB_ID 434c3653-b18c-6abf-e583-c6e2e8b473b6
\needspace{5\baselineskip}
{\bfseries\sffamily\color{ink} 381.} {scriptsizecolor{muted} Y1 | review | 434c3653}
\par
Fayl, kengaytma, format va atribut doirasida "metama'lumot" atamasining eng aniq ta'rifini toping.
\begin{enumerate}[label=\Alph*),leftmargin=9mm]
\item fayl mazmuni haqida tavsiflovchi ma'lumot
\item nom bilan saqlanadigan ma'lumotlar birligi
\item fayl nomining formatni ko'rsatuvchi oxirgi qismi
\item fayl ichidagi ma'lumot tuzilishi va kodlash qoidasi
\end{enumerate}
\begin{tcolorbox}[enhanced,breakable,colback=teal!4,colframe=teal!60!black,boxrule=0.45pt,arc=1.2mm,left=2.7mm,right=2.7mm,top=1.8mm,bottom=1.8mm,before skip=2mm,after skip=2mm]
{\bfseries\sffamily\color{teal}To'g'ri javob:} \textbf{A) fayl mazmuni haqida tavsiflovchi ma'lumot}
\par\smallskip
{\bfseries Izoh:} fayl mazmuni haqida tavsiflovchi ma'lumot. metama'lumot - fayl mazmuni haqida tavsiflovchi ma'lumot. Manba: CAM7, PDF 41-54; Qo'shimcha sahifa raqami (qaysi PDFga tegishli ekani aniqlanmagan; darslik sahifasi deb tasdiqlanmagan): 87
\end{tcolorbox}
\vspace{2mm}\hrule height 0.2pt\vspace{3mm}

% DB_ID 59092d24-a05d-942c-d210-9bc80f070818
\needspace{5\baselineskip}
{\bfseries\sffamily\color{ink} 382.} {scriptsizecolor{muted} Y1 | review | 59092d24}
\par
Fayl, kengaytma, format va atribut doirasida "konvertatsiya" atamasining eng aniq ta'rifini toping.
\begin{enumerate}[label=\Alph*),leftmargin=9mm]
\item nom bilan saqlanadigan ma'lumotlar birligi
\item ma'lumotni bir format tuzilishidan boshqasiga o'tkazish
\item fayl nomining formatni ko'rsatuvchi oxirgi qismi
\item fayl ichidagi ma'lumot tuzilishi va kodlash qoidasi
\end{enumerate}
\begin{tcolorbox}[enhanced,breakable,colback=teal!4,colframe=teal!60!black,boxrule=0.45pt,arc=1.2mm,left=2.7mm,right=2.7mm,top=1.8mm,bottom=1.8mm,before skip=2mm,after skip=2mm]
{\bfseries\sffamily\color{teal}To'g'ri javob:} \textbf{B) ma'lumotni bir format tuzilishidan boshqasiga o'tkazish}
\par\smallskip
{\bfseries Izoh:} ma'lumotni bir format tuzilishidan boshqasiga o'tkazish. konvertatsiya - ma'lumotni bir format tuzilishidan boshqasiga o'tkazish. Manba: ICT5, PDF 42-45; Qo'shimcha sahifa raqami (qaysi PDFga tegishli ekani aniqlanmagan; darslik sahifasi deb tasdiqlanmagan): 88
\end{tcolorbox}
\vspace{2mm}\hrule height 0.2pt\vspace{3mm}

% DB_ID 5e1ca296-e6b5-b32f-ae28-37bd873101fd
\needspace{5\baselineskip}
{\bfseries\sffamily\color{ink} 383.} {scriptsizecolor{muted} Y1 | review | 5e1ca296}
\par
Papka, katalog, yo'l va ierarxiya doirasida "papka" atamasining eng aniq ta'rifini toping.
\begin{enumerate}[label=\Alph*),leftmargin=9mm]
\item fayl tizimi daraxtining eng yuqori katalogi
\item joriy obyektni bevosita o'z ichiga olgan katalog
\item fayl va quyi papkalarni guruhlovchi konteyner
\item boshqa katalog ichida joylashgan katalog
\end{enumerate}
\begin{tcolorbox}[enhanced,breakable,colback=teal!4,colframe=teal!60!black,boxrule=0.45pt,arc=1.2mm,left=2.7mm,right=2.7mm,top=1.8mm,bottom=1.8mm,before skip=2mm,after skip=2mm]
{\bfseries\sffamily\color{teal}To'g'ri javob:} \textbf{C) fayl va quyi papkalarni guruhlovchi konteyner}
\par\smallskip
{\bfseries Izoh:} fayl va quyi papkalarni guruhlovchi konteyner. papka - fayl va quyi papkalarni guruhlovchi konteyner. Manba: ICT5, PDF 42-45; Qo'shimcha sahifa raqami (qaysi PDFga tegishli ekani aniqlanmagan; darslik sahifasi deb tasdiqlanmagan): 89
\end{tcolorbox}
\vspace{2mm}\hrule height 0.2pt\vspace{3mm}

% DB_ID 69e7ba49-808b-c14b-9ef4-5cd395933e1c
\needspace{5\baselineskip}
{\bfseries\sffamily\color{ink} 384.} {scriptsizecolor{muted} Y1 | review | 69e7ba49}
\par
Papka, katalog, yo'l va ierarxiya doirasida "quyi katalog" atamasining eng aniq ta'rifini toping.
\begin{enumerate}[label=\Alph*),leftmargin=9mm]
\item fayl va quyi papkalarni guruhlovchi konteyner
\item boshqa katalog ichida joylashgan katalog
\item fayl tizimi daraxtining eng yuqori katalogi
\item joriy obyektni bevosita o'z ichiga olgan katalog
\end{enumerate}
\begin{tcolorbox}[enhanced,breakable,colback=teal!4,colframe=teal!60!black,boxrule=0.45pt,arc=1.2mm,left=2.7mm,right=2.7mm,top=1.8mm,bottom=1.8mm,before skip=2mm,after skip=2mm]
{\bfseries\sffamily\color{teal}To'g'ri javob:} \textbf{B) boshqa katalog ichida joylashgan katalog}
\par\smallskip
{\bfseries Izoh:} boshqa katalog ichida joylashgan katalog. quyi katalog - boshqa katalog ichida joylashgan katalog. Manba: CAM9, PDF 147-149; Qo'shimcha sahifa raqami (qaysi PDFga tegishli ekani aniqlanmagan; darslik sahifasi deb tasdiqlanmagan): 92
\end{tcolorbox}
\vspace{2mm}\hrule height 0.2pt\vspace{3mm}

% DB_ID 7c8f770f-51c9-0512-5ed8-48843f6d04ef
\needspace{5\baselineskip}
{\bfseries\sffamily\color{ink} 385.} {scriptsizecolor{muted} Y1 | review | 7c8f770f}
\par
Fayl va papkalar ustida amallar doirasida "filter" atamasining eng aniq ta'rifini toping.
\begin{enumerate}[label=\Alph*),leftmargin=9mm]
\item manbani saqlagan holda yangi nusxa yaratish
\item obyektning joylashuvini o'zgartirish
\item obyekt nomini o'zgartirish
\item shartga mos obyektlargina ko'rinadigan qilish
\end{enumerate}
\begin{tcolorbox}[enhanced,breakable,colback=teal!4,colframe=teal!60!black,boxrule=0.45pt,arc=1.2mm,left=2.7mm,right=2.7mm,top=1.8mm,bottom=1.8mm,before skip=2mm,after skip=2mm]
{\bfseries\sffamily\color{teal}To'g'ri javob:} \textbf{D) shartga mos obyektlargina ko'rinadigan qilish}
\par\smallskip
{\bfseries Izoh:} shartga mos obyektlargina ko'rinadigan qilish. filter - shartga mos obyektlargina ko'rinadigan qilish. Manba: CAM9, PDF 147-149; Qo'shimcha sahifa raqami (qaysi PDFga tegishli ekani aniqlanmagan; darslik sahifasi deb tasdiqlanmagan): 97
\end{tcolorbox}
\vspace{2mm}\hrule height 0.2pt\vspace{3mm}

% DB_ID 84f77474-eb41-a038-bb8e-b8fde9d7392d
\needspace{5\baselineskip}
{\bfseries\sffamily\color{ink} 386.} {scriptsizecolor{muted} Y1 | review | 84f77474}
\par
Papka, katalog, yo'l va ierarxiya doirasida "ildiz katalog" atamasining eng aniq ta'rifini toping.
\begin{enumerate}[label=\Alph*),leftmargin=9mm]
\item fayl va quyi papkalarni guruhlovchi konteyner
\item joriy obyektni bevosita o'z ichiga olgan katalog
\item boshqa katalog ichida joylashgan katalog
\item fayl tizimi daraxtining eng yuqori katalogi
\end{enumerate}
\begin{tcolorbox}[enhanced,breakable,colback=teal!4,colframe=teal!60!black,boxrule=0.45pt,arc=1.2mm,left=2.7mm,right=2.7mm,top=1.8mm,bottom=1.8mm,before skip=2mm,after skip=2mm]
{\bfseries\sffamily\color{teal}To'g'ri javob:} \textbf{D) fayl tizimi daraxtining eng yuqori katalogi}
\par\smallskip
{\bfseries Izoh:} fayl tizimi daraxtining eng yuqori katalogi. ildiz katalog - fayl tizimi daraxtining eng yuqori katalogi. Manba: MAR, PDF 16-20; Qo'shimcha sahifa raqami (qaysi PDFga tegishli ekani aniqlanmagan; darslik sahifasi deb tasdiqlanmagan): 90
\end{tcolorbox}
\vspace{2mm}\hrule height 0.2pt\vspace{3mm}

% DB_ID 89af10b0-ee78-4c11-b544-e267aada3bd3
\needspace{5\baselineskip}
{\bfseries\sffamily\color{ink} 387.} {scriptsizecolor{muted} Y1 | review | 89af10b0}
\par
Fayl, kengaytma, format va atribut doirasida "format" atamasining eng aniq ta'rifini toping.
\begin{enumerate}[label=\Alph*),leftmargin=9mm]
\item nom bilan saqlanadigan ma'lumotlar birligi
\item fayl nomining formatni ko'rsatuvchi oxirgi qismi
\item fayl ichidagi ma'lumot tuzilishi va kodlash qoidasi
\item faylning read-only yoki hidden kabi boshqaruv xususiyati
\end{enumerate}
\begin{tcolorbox}[enhanced,breakable,colback=teal!4,colframe=teal!60!black,boxrule=0.45pt,arc=1.2mm,left=2.7mm,right=2.7mm,top=1.8mm,bottom=1.8mm,before skip=2mm,after skip=2mm]
{\bfseries\sffamily\color{teal}To'g'ri javob:} \textbf{C) fayl ichidagi ma'lumot tuzilishi va kodlash qoidasi}
\par\smallskip
{\bfseries Izoh:} fayl ichidagi ma'lumot tuzilishi va kodlash qoidasi. format - fayl ichidagi ma'lumot tuzilishi va kodlash qoidasi. Manba: CAM9, PDF 147-149; Qo'shimcha sahifa raqami (qaysi PDFga tegishli ekani aniqlanmagan; darslik sahifasi deb tasdiqlanmagan): 85
\end{tcolorbox}
\vspace{2mm}\hrule height 0.2pt\vspace{3mm}

% DB_ID 8b48358a-9716-acbd-ce27-aabb8bbb0844
\needspace{5\baselineskip}
{\bfseries\sffamily\color{ink} 388.} {scriptsizecolor{muted} Y1 | review | 8b48358a}
\par
Fayl va papkalar ustida amallar doirasida "backup" atamasining eng aniq ta'rifini toping.
\begin{enumerate}[label=\Alph*),leftmargin=9mm]
\item manbani saqlagan holda yangi nusxa yaratish
\item yo'qolishdan tiklash uchun alohida joydagi qo'shimcha nusxa
\item obyektning joylashuvini o'zgartirish
\item obyekt nomini o'zgartirish
\end{enumerate}
\begin{tcolorbox}[enhanced,breakable,colback=teal!4,colframe=teal!60!black,boxrule=0.45pt,arc=1.2mm,left=2.7mm,right=2.7mm,top=1.8mm,bottom=1.8mm,before skip=2mm,after skip=2mm]
{\bfseries\sffamily\color{teal}To'g'ri javob:} \textbf{B) yo'qolishdan tiklash uchun alohida joydagi qo'shimcha nusxa}
\par\smallskip
{\bfseries Izoh:} yo'qolishdan tiklash uchun alohida joydagi qo'shimcha nusxa. backup - yo'qolishdan tiklash uchun alohida joydagi qo'shimcha nusxa. Manba: ICT5, PDF 43-49; Qo'shimcha sahifa raqami (qaysi PDFga tegishli ekani aniqlanmagan; darslik sahifasi deb tasdiqlanmagan): 99
\end{tcolorbox}
\vspace{2mm}\hrule height 0.2pt\vspace{3mm}

% DB_ID 91701120-5b74-ab0e-595b-b353e97e746b
\needspace{5\baselineskip}
{\bfseries\sffamily\color{ink} 389.} {scriptsizecolor{muted} Y1 | review | 91701120}
\par
Fayl va papkalar ustida amallar doirasida "archive" atamasining eng aniq ta'rifini toping.
\begin{enumerate}[label=\Alph*),leftmargin=9mm]
\item bir nechta faylni konteynerga yig'ish va ko'pincha siqish
\item manbani saqlagan holda yangi nusxa yaratish
\item obyektning joylashuvini o'zgartirish
\item obyekt nomini o'zgartirish
\end{enumerate}
\begin{tcolorbox}[enhanced,breakable,colback=teal!4,colframe=teal!60!black,boxrule=0.45pt,arc=1.2mm,left=2.7mm,right=2.7mm,top=1.8mm,bottom=1.8mm,before skip=2mm,after skip=2mm]
{\bfseries\sffamily\color{teal}To'g'ri javob:} \textbf{A) bir nechta faylni konteynerga yig'ish va ko'pincha siqish}
\par\smallskip
{\bfseries Izoh:} bir nechta faylni konteynerga yig'ish va ko'pincha siqish. archive - bir nechta faylni konteynerga yig'ish va ko'pincha siqish. Manba: ICT6, PDF 116-126; Qo'shimcha sahifa raqami (qaysi PDFga tegishli ekani aniqlanmagan; darslik sahifasi deb tasdiqlanmagan): 98
\end{tcolorbox}
\vspace{2mm}\hrule height 0.2pt\vspace{3mm}

% DB_ID a26f95d8-efac-b22d-fa39-b9fed6e3f813
\needspace{5\baselineskip}
{\bfseries\sffamily\color{ink} 390.} {scriptsizecolor{muted} Y1 | review | a26f95d8}
\par
Fayl, kengaytma, format va atribut doirasida "atribut" atamasining eng aniq ta'rifini toping.
\begin{enumerate}[label=\Alph*),leftmargin=9mm]
\item nom bilan saqlanadigan ma'lumotlar birligi
\item fayl nomining formatni ko'rsatuvchi oxirgi qismi
\item fayl ichidagi ma'lumot tuzilishi va kodlash qoidasi
\item faylning read-only yoki hidden kabi boshqaruv xususiyati
\end{enumerate}
\begin{tcolorbox}[enhanced,breakable,colback=teal!4,colframe=teal!60!black,boxrule=0.45pt,arc=1.2mm,left=2.7mm,right=2.7mm,top=1.8mm,bottom=1.8mm,before skip=2mm,after skip=2mm]
{\bfseries\sffamily\color{teal}To'g'ri javob:} \textbf{D) faylning read-only yoki hidden kabi boshqaruv xususiyati}
\par\smallskip
{\bfseries Izoh:} faylning read-only yoki hidden kabi boshqaruv xususiyati. atribut - faylning read-only yoki hidden kabi boshqaruv xususiyati. Manba: CAM6, PDF 178-205; Qo'shimcha sahifa raqami (qaysi PDFga tegishli ekani aniqlanmagan; darslik sahifasi deb tasdiqlanmagan): 86
\end{tcolorbox}
\vspace{2mm}\hrule height 0.2pt\vspace{3mm}

% DB_ID d94caafa-0df1-98b3-60e7-fced2b2dd317
\needspace{5\baselineskip}
{\bfseries\sffamily\color{ink} 391.} {scriptsizecolor{muted} Y1 | review | d94caafa}
\par
Fayl va papkalar ustida amallar doirasida "rename" atamasining eng aniq ta'rifini toping.
\begin{enumerate}[label=\Alph*),leftmargin=9mm]
\item manbani saqlagan holda yangi nusxa yaratish
\item obyektning joylashuvini o'zgartirish
\item obyekt nomini o'zgartirish
\item shartga mos obyektlargina ko'rinadigan qilish
\end{enumerate}
\begin{tcolorbox}[enhanced,breakable,colback=teal!4,colframe=teal!60!black,boxrule=0.45pt,arc=1.2mm,left=2.7mm,right=2.7mm,top=1.8mm,bottom=1.8mm,before skip=2mm,after skip=2mm]
{\bfseries\sffamily\color{teal}To'g'ri javob:} \textbf{C) obyekt nomini o'zgartirish}
\par\smallskip
{\bfseries Izoh:} obyekt nomini o'zgartirish. rename - obyekt nomini o'zgartirish. Manba: CAM6, PDF 196-205; Qo'shimcha sahifa raqami (qaysi PDFga tegishli ekani aniqlanmagan; darslik sahifasi deb tasdiqlanmagan): 97
\end{tcolorbox}
\vspace{2mm}\hrule height 0.2pt\vspace{3mm}

% DB_ID e87aef15-4d24-39a8-e730-e2a86cd48954
\needspace{5\baselineskip}
{\bfseries\sffamily\color{ink} 392.} {scriptsizecolor{muted} Y1 | review | e87aef15}
\par
Papka, katalog, yo'l va ierarxiya doirasida "ona katalog" atamasining eng aniq ta'rifini toping.
\begin{enumerate}[label=\Alph*),leftmargin=9mm]
\item joriy obyektni bevosita o'z ichiga olgan katalog
\item fayl va quyi papkalarni guruhlovchi konteyner
\item fayl tizimi daraxtining eng yuqori katalogi
\item boshqa katalog ichida joylashgan katalog
\end{enumerate}
\begin{tcolorbox}[enhanced,breakable,colback=teal!4,colframe=teal!60!black,boxrule=0.45pt,arc=1.2mm,left=2.7mm,right=2.7mm,top=1.8mm,bottom=1.8mm,before skip=2mm,after skip=2mm]
{\bfseries\sffamily\color{teal}To'g'ri javob:} \textbf{A) joriy obyektni bevosita o'z ichiga olgan katalog}
\par\smallskip
{\bfseries Izoh:} joriy obyektni bevosita o'z ichiga olgan katalog. ona katalog - joriy obyektni bevosita o'z ichiga olgan katalog. Manba: CAM6, PDF 178-205; Qo'shimcha sahifa raqami (qaysi PDFga tegishli ekani aniqlanmagan; darslik sahifasi deb tasdiqlanmagan): 91
\end{tcolorbox}
\vspace{2mm}\hrule height 0.2pt\vspace{3mm}

% DB_ID fd2db667-4f8e-c544-c961-6d4659f977f1
\needspace{5\baselineskip}
{\bfseries\sffamily\color{ink} 393.} {scriptsizecolor{muted} Y1 | review | fd2db667}
\par
Fayl va papkalar ustida amallar doirasida "copy" atamasining eng aniq ta'rifini toping.
\begin{enumerate}[label=\Alph*),leftmargin=9mm]
\item manbani saqlagan holda yangi nusxa yaratish
\item obyektning joylashuvini o'zgartirish
\item obyekt nomini o'zgartirish
\item shartga mos obyektlargina ko'rinadigan qilish
\end{enumerate}
\begin{tcolorbox}[enhanced,breakable,colback=teal!4,colframe=teal!60!black,boxrule=0.45pt,arc=1.2mm,left=2.7mm,right=2.7mm,top=1.8mm,bottom=1.8mm,before skip=2mm,after skip=2mm]
{\bfseries\sffamily\color{teal}To'g'ri javob:} \textbf{A) manbani saqlagan holda yangi nusxa yaratish}
\par\smallskip
{\bfseries Izoh:} manbani saqlagan holda yangi nusxa yaratish. copy - manbani saqlagan holda yangi nusxa yaratish. Manba: ICT5, PDF 43-49; Qo'shimcha sahifa raqami (qaysi PDFga tegishli ekani aniqlanmagan; darslik sahifasi deb tasdiqlanmagan): 95
\end{tcolorbox}
\vspace{2mm}\hrule height 0.2pt\vspace{3mm}

% DB_ID fd9e92d8-e762-203e-b637-50a3ffc28dfc
\needspace{5\baselineskip}
{\bfseries\sffamily\color{ink} 394.} {scriptsizecolor{muted} Y1 | review | fd9e92d8}
\par
Fayl va papkalar ustida amallar doirasida "move" atamasining eng aniq ta'rifini toping.
\begin{enumerate}[label=\Alph*),leftmargin=9mm]
\item manbani saqlagan holda yangi nusxa yaratish
\item obyektning joylashuvini o'zgartirish
\item obyekt nomini o'zgartirish
\item shartga mos obyektlargina ko'rinadigan qilish
\end{enumerate}
\begin{tcolorbox}[enhanced,breakable,colback=teal!4,colframe=teal!60!black,boxrule=0.45pt,arc=1.2mm,left=2.7mm,right=2.7mm,top=1.8mm,bottom=1.8mm,before skip=2mm,after skip=2mm]
{\bfseries\sffamily\color{teal}To'g'ri javob:} \textbf{B) obyektning joylashuvini o'zgartirish}
\par\smallskip
{\bfseries Izoh:} obyektning joylashuvini o'zgartirish. move - obyektning joylashuvini o'zgartirish. Manba: MAR, PDF 17-20; Qo'shimcha sahifa raqami (qaysi PDFga tegishli ekani aniqlanmagan; darslik sahifasi deb tasdiqlanmagan): 96
\end{tcolorbox}
\vspace{2mm}\hrule height 0.2pt\vspace{3mm}

% DB_ID fe57d38a-eb7b-54ec-3a37-f29b636e7381
\needspace{5\baselineskip}
{\bfseries\sffamily\color{ink} 395.} {scriptsizecolor{muted} Y1 | review | fe57d38a}
\par
Papka, katalog, yo'l va ierarxiya doirasida "nisbiy yo'l" atamasining eng aniq ta'rifini toping.
\begin{enumerate}[label=\Alph*),leftmargin=9mm]
\item fayl va quyi papkalarni guruhlovchi konteyner
\item fayl tizimi daraxtining eng yuqori katalogi
\item joriy obyektni bevosita o'z ichiga olgan katalog
\item joriy katalogga nisbatan berilgan joylashuv
\end{enumerate}
\begin{tcolorbox}[enhanced,breakable,colback=teal!4,colframe=teal!60!black,boxrule=0.45pt,arc=1.2mm,left=2.7mm,right=2.7mm,top=1.8mm,bottom=1.8mm,before skip=2mm,after skip=2mm]
{\bfseries\sffamily\color{teal}To'g'ri javob:} \textbf{D) joriy katalogga nisbatan berilgan joylashuv}
\par\smallskip
{\bfseries Izoh:} joriy katalogga nisbatan berilgan joylashuv. nisbiy yo'l - joriy katalogga nisbatan berilgan joylashuv. Manba: ICT5, PDF 42-45; Qo'shimcha sahifa raqami (qaysi PDFga tegishli ekani aniqlanmagan; darslik sahifasi deb tasdiqlanmagan): 94
\end{tcolorbox}
\vspace{2mm}\hrule height 0.2pt\vspace{3mm}

% DB_ID 02c505b4-95b3-bda6-43a3-95bb0ec3b761
\needspace{5\baselineskip}
{\bfseries\sffamily\color{ink} 396.} {scriptsizecolor{muted} Y1 | review | 02c505b4}
\par
Fayl, kengaytma, format va atribut mavzusida "fayl" qatnashgan atama-ta'rif juftliklaridan to'g'risini toping.
\begin{enumerate}[label=\Alph*),leftmargin=9mm]
\item fayl - fayl nomining formatni ko'rsatuvchi oxirgi qismi
\item fayl - fayl ichidagi ma'lumot tuzilishi va kodlash qoidasi
\item fayl - nom bilan saqlanadigan ma'lumotlar birligi
\item fayl - faylning read-only yoki hidden kabi boshqaruv xususiyati
\end{enumerate}
\begin{tcolorbox}[enhanced,breakable,colback=teal!4,colframe=teal!60!black,boxrule=0.45pt,arc=1.2mm,left=2.7mm,right=2.7mm,top=1.8mm,bottom=1.8mm,before skip=2mm,after skip=2mm]
{\bfseries\sffamily\color{teal}To'g'ri javob:} \textbf{C) fayl - nom bilan saqlanadigan ma'lumotlar birligi}
\par\smallskip
{\bfseries Izoh:} fayl - nom bilan saqlanadigan ma'lumotlar birligi. To'g'ri juftlik: fayl - nom bilan saqlanadigan ma'lumotlar birligi. Manba: ICT5, PDF 42-45; Qo'shimcha sahifa raqami (qaysi PDFga tegishli ekani aniqlanmagan; darslik sahifasi deb tasdiqlanmagan): 83
\end{tcolorbox}
\vspace{2mm}\hrule height 0.2pt\vspace{3mm}

% DB_ID 11f52a78-649d-268d-8553-11911f42cade
\needspace{5\baselineskip}
{\bfseries\sffamily\color{ink} 397.} {scriptsizecolor{muted} Y1 | review | 11f52a78}
\par
Fayl, kengaytma, format va atribut mavzusidagi vazifada fayl ichidagi ma'lumot tuzilishi va kodlash qoidasi talab qilindi. Qaysi tushuncha mos?
\begin{enumerate}[label=\Alph*),leftmargin=9mm]
\item fayl
\item kengaytma
\item atribut
\item format
\end{enumerate}
\begin{tcolorbox}[enhanced,breakable,colback=teal!4,colframe=teal!60!black,boxrule=0.45pt,arc=1.2mm,left=2.7mm,right=2.7mm,top=1.8mm,bottom=1.8mm,before skip=2mm,after skip=2mm]
{\bfseries\sffamily\color{teal}To'g'ri javob:} \textbf{D) format}
\par\smallskip
{\bfseries Izoh:} format. Berilgan vazifa format tushunchasiga mos: fayl ichidagi ma'lumot tuzilishi va kodlash qoidasi. Manba: CAM9, PDF 147-149; Qo'shimcha sahifa raqami (qaysi PDFga tegishli ekani aniqlanmagan; darslik sahifasi deb tasdiqlanmagan): 85
\end{tcolorbox}
\vspace{2mm}\hrule height 0.2pt\vspace{3mm}

% DB_ID 3bf607d2-944e-f06a-f491-f79ae09a4a49
\needspace{5\baselineskip}
{\bfseries\sffamily\color{ink} 398.} {scriptsizecolor{muted} Y1 | review | 3bf607d2}
\par
Fayl va papkalar ustida amallar mavzusida "rename" qatnashgan atama-ta'rif juftliklaridan to'g'risini toping.
\begin{enumerate}[label=\Alph*),leftmargin=9mm]
\item rename - obyekt nomini o'zgartirish
\item rename - manbani saqlagan holda yangi nusxa yaratish
\item rename - obyektning joylashuvini o'zgartirish
\item rename - shartga mos obyektlargina ko'rinadigan qilish
\end{enumerate}
\begin{tcolorbox}[enhanced,breakable,colback=teal!4,colframe=teal!60!black,boxrule=0.45pt,arc=1.2mm,left=2.7mm,right=2.7mm,top=1.8mm,bottom=1.8mm,before skip=2mm,after skip=2mm]
{\bfseries\sffamily\color{teal}To'g'ri javob:} \textbf{A) rename - obyekt nomini o'zgartirish}
\par\smallskip
{\bfseries Izoh:} rename - obyekt nomini o'zgartirish. To'g'ri juftlik: rename - obyekt nomini o'zgartirish. Manba: CAM6, PDF 196-205; Qo'shimcha sahifa raqami (qaysi PDFga tegishli ekani aniqlanmagan; darslik sahifasi deb tasdiqlanmagan): 97
\end{tcolorbox}
\vspace{2mm}\hrule height 0.2pt\vspace{3mm}

% DB_ID 46df4099-b009-8748-5e31-5b9abd7c1c89
\needspace{5\baselineskip}
{\bfseries\sffamily\color{ink} 399.} {scriptsizecolor{muted} Y1 | review | 46df4099}
\par
Papka, katalog, yo'l va ierarxiya mavzusidagi vazifada ildizdan boshlanadigan to'liq joylashuv talab qilindi. Qaysi tushuncha mos?
\begin{enumerate}[label=\Alph*),leftmargin=9mm]
\item papka
\item ildiz katalog
\item ona katalog
\item mutlaq yo'l
\end{enumerate}
\begin{tcolorbox}[enhanced,breakable,colback=teal!4,colframe=teal!60!black,boxrule=0.45pt,arc=1.2mm,left=2.7mm,right=2.7mm,top=1.8mm,bottom=1.8mm,before skip=2mm,after skip=2mm]
{\bfseries\sffamily\color{teal}To'g'ri javob:} \textbf{D) mutlaq yo'l}
\par\smallskip
{\bfseries Izoh:} mutlaq yo'l. Berilgan vazifa mutlaq yo'l tushunchasiga mos: ildizdan boshlanadigan to'liq joylashuv. Manba: ICT6, PDF 116-124; Qo'shimcha sahifa raqami (qaysi PDFga tegishli ekani aniqlanmagan; darslik sahifasi deb tasdiqlanmagan): 93
\end{tcolorbox}
\vspace{2mm}\hrule height 0.2pt\vspace{3mm}

% DB_ID 4e20d53a-efb5-cb78-a6df-d7fdbe83119b
\needspace{5\baselineskip}
{\bfseries\sffamily\color{ink} 400.} {scriptsizecolor{muted} Y1 | review | 4e20d53a}
\par
Papka, katalog, yo'l va ierarxiya mavzusida "papka" qatnashgan atama-ta'rif juftliklaridan to'g'risini toping.
\begin{enumerate}[label=\Alph*),leftmargin=9mm]
\item papka - fayl va quyi papkalarni guruhlovchi konteyner
\item papka - fayl tizimi daraxtining eng yuqori katalogi
\item papka - joriy obyektni bevosita o'z ichiga olgan katalog
\item papka - boshqa katalog ichida joylashgan katalog
\end{enumerate}
\begin{tcolorbox}[enhanced,breakable,colback=teal!4,colframe=teal!60!black,boxrule=0.45pt,arc=1.2mm,left=2.7mm,right=2.7mm,top=1.8mm,bottom=1.8mm,before skip=2mm,after skip=2mm]
{\bfseries\sffamily\color{teal}To'g'ri javob:} \textbf{A) papka - fayl va quyi papkalarni guruhlovchi konteyner}
\par\smallskip
{\bfseries Izoh:} papka - fayl va quyi papkalarni guruhlovchi konteyner. To'g'ri juftlik: papka - fayl va quyi papkalarni guruhlovchi konteyner. Manba: ICT5, PDF 42-45; Qo'shimcha sahifa raqami (qaysi PDFga tegishli ekani aniqlanmagan; darslik sahifasi deb tasdiqlanmagan): 89
\end{tcolorbox}
\vspace{2mm}\hrule height 0.2pt\vspace{3mm}

% DB_ID 5dbecd23-c334-14d2-e340-521c61c38672
\needspace{5\baselineskip}
{\bfseries\sffamily\color{ink} 401.} {scriptsizecolor{muted} Y1 | review | 5dbecd23}
\par
Fayl va papkalar ustida amallar mavzusidagi vazifada shartga mos obyektlargina ko'rinadigan qilish talab qilindi. Qaysi tushuncha mos?
\begin{enumerate}[label=\Alph*),leftmargin=9mm]
\item filter
\item copy
\item move
\item rename
\end{enumerate}
\begin{tcolorbox}[enhanced,breakable,colback=teal!4,colframe=teal!60!black,boxrule=0.45pt,arc=1.2mm,left=2.7mm,right=2.7mm,top=1.8mm,bottom=1.8mm,before skip=2mm,after skip=2mm]
{\bfseries\sffamily\color{teal}To'g'ri javob:} \textbf{A) filter}
\par\smallskip
{\bfseries Izoh:} filter. Berilgan vazifa filter tushunchasiga mos: shartga mos obyektlargina ko'rinadigan qilish. Manba: CAM9, PDF 147-149; Qo'shimcha sahifa raqami (qaysi PDFga tegishli ekani aniqlanmagan; darslik sahifasi deb tasdiqlanmagan): 98
\end{tcolorbox}
\vspace{2mm}\hrule height 0.2pt\vspace{3mm}

% DB_ID 6827d8c6-ff1e-e8bb-e60e-14797ca13c0e
\needspace{5\baselineskip}
{\bfseries\sffamily\color{ink} 402.} {scriptsizecolor{muted} Y1 | review | 6827d8c6}
\par
Fayl, kengaytma, format va atribut mavzusidagi vazifada nom bilan saqlanadigan ma'lumotlar birligi talab qilindi. Qaysi tushuncha mos?
\begin{enumerate}[label=\Alph*),leftmargin=9mm]
\item kengaytma
\item fayl
\item format
\item atribut
\end{enumerate}
\begin{tcolorbox}[enhanced,breakable,colback=teal!4,colframe=teal!60!black,boxrule=0.45pt,arc=1.2mm,left=2.7mm,right=2.7mm,top=1.8mm,bottom=1.8mm,before skip=2mm,after skip=2mm]
{\bfseries\sffamily\color{teal}To'g'ri javob:} \textbf{B) fayl}
\par\smallskip
{\bfseries Izoh:} fayl. Berilgan vazifa fayl tushunchasiga mos: nom bilan saqlanadigan ma'lumotlar birligi. Manba: ICT5, PDF 42-45; Qo'shimcha sahifa raqami (qaysi PDFga tegishli ekani aniqlanmagan; darslik sahifasi deb tasdiqlanmagan): 83
\end{tcolorbox}
\vspace{2mm}\hrule height 0.2pt\vspace{3mm}

% DB_ID 694d7547-002c-c5ab-e3d1-034c4543c21d
\needspace{5\baselineskip}
{\bfseries\sffamily\color{ink} 403.} {scriptsizecolor{muted} Y1 | review | 694d7547}
\par
Papka, katalog, yo'l va ierarxiya mavzusidagi vazifada joriy obyektni bevosita o'z ichiga olgan katalog talab qilindi. Qaysi tushuncha mos?
\begin{enumerate}[label=\Alph*),leftmargin=9mm]
\item papka
\item ona katalog
\item ildiz katalog
\item quyi katalog
\end{enumerate}
\begin{tcolorbox}[enhanced,breakable,colback=teal!4,colframe=teal!60!black,boxrule=0.45pt,arc=1.2mm,left=2.7mm,right=2.7mm,top=1.8mm,bottom=1.8mm,before skip=2mm,after skip=2mm]
{\bfseries\sffamily\color{teal}To'g'ri javob:} \textbf{B) ona katalog}
\par\smallskip
{\bfseries Izoh:} ona katalog. Berilgan vazifa ona katalog tushunchasiga mos: joriy obyektni bevosita o'z ichiga olgan katalog. Manba: CAM6, PDF 178-205; Qo'shimcha sahifa raqami (qaysi PDFga tegishli ekani aniqlanmagan; darslik sahifasi deb tasdiqlanmagan): 91
\end{tcolorbox}
\vspace{2mm}\hrule height 0.2pt\vspace{3mm}

% DB_ID 6c9f6d48-dae5-aaf6-dadb-6a91a5c5376d
\needspace{5\baselineskip}
{\bfseries\sffamily\color{ink} 404.} {scriptsizecolor{muted} Y1 | review | 6c9f6d48}
\par
Fayl, kengaytma, format va atribut mavzusidagi vazifada ma'lumotni bir format tuzilishidan boshqasiga o'tkazish talab qilindi. Qaysi tushuncha mos?
\begin{enumerate}[label=\Alph*),leftmargin=9mm]
\item fayl
\item kengaytma
\item konvertatsiya
\item format
\end{enumerate}
\begin{tcolorbox}[enhanced,breakable,colback=teal!4,colframe=teal!60!black,boxrule=0.45pt,arc=1.2mm,left=2.7mm,right=2.7mm,top=1.8mm,bottom=1.8mm,before skip=2mm,after skip=2mm]
{\bfseries\sffamily\color{teal}To'g'ri javob:} \textbf{C) konvertatsiya}
\par\smallskip
{\bfseries Izoh:} konvertatsiya. Berilgan vazifa konvertatsiya tushunchasiga mos: ma'lumotni bir format tuzilishidan boshqasiga o'tkazish. Manba: ICT5, PDF 42-45; Qo'shimcha sahifa raqami (qaysi PDFga tegishli ekani aniqlanmagan; darslik sahifasi deb tasdiqlanmagan): 88
\end{tcolorbox}
\vspace{2mm}\hrule height 0.2pt\vspace{3mm}

% DB_ID 6f2bc84f-9711-11f2-a3f6-aebb04ee791e
\needspace{5\baselineskip}
{\bfseries\sffamily\color{ink} 405.} {scriptsizecolor{muted} Y1 | review | 6f2bc84f}
\par
Fayl, kengaytma, format va atribut mavzusida "metama'lumot" qatnashgan atama-ta'rif juftliklaridan to'g'risini toping.
\begin{enumerate}[label=\Alph*),leftmargin=9mm]
\item metama'lumot - nom bilan saqlanadigan ma'lumotlar birligi
\item metama'lumot - fayl nomining formatni ko'rsatuvchi oxirgi qismi
\item metama'lumot - fayl mazmuni haqida tavsiflovchi ma'lumot
\item metama'lumot - fayl ichidagi ma'lumot tuzilishi va kodlash qoidasi
\end{enumerate}
\begin{tcolorbox}[enhanced,breakable,colback=teal!4,colframe=teal!60!black,boxrule=0.45pt,arc=1.2mm,left=2.7mm,right=2.7mm,top=1.8mm,bottom=1.8mm,before skip=2mm,after skip=2mm]
{\bfseries\sffamily\color{teal}To'g'ri javob:} \textbf{C) metama'lumot - fayl mazmuni haqida tavsiflovchi ma'lumot}
\par\smallskip
{\bfseries Izoh:} metama'lumot - fayl mazmuni haqida tavsiflovchi ma'lumot. To'g'ri juftlik: metama'lumot - fayl mazmuni haqida tavsiflovchi ma'lumot. Manba: CAM7, PDF 41-54; Qo'shimcha sahifa raqami (qaysi PDFga tegishli ekani aniqlanmagan; darslik sahifasi deb tasdiqlanmagan): 87
\end{tcolorbox}
\vspace{2mm}\hrule height 0.2pt\vspace{3mm}

% DB_ID 72f5d561-9b20-d263-eac5-416dfba5f86f
\needspace{5\baselineskip}
{\bfseries\sffamily\color{ink} 406.} {scriptsizecolor{muted} Y1 | review | 72f5d561}
\par
Papka, katalog, yo'l va ierarxiya mavzusida "ona katalog" qatnashgan atama-ta'rif juftliklaridan to'g'risini toping.
\begin{enumerate}[label=\Alph*),leftmargin=9mm]
\item ona katalog - fayl va quyi papkalarni guruhlovchi konteyner
\item ona katalog - fayl tizimi daraxtining eng yuqori katalogi
\item ona katalog - joriy obyektni bevosita o'z ichiga olgan katalog
\item ona katalog - boshqa katalog ichida joylashgan katalog
\end{enumerate}
\begin{tcolorbox}[enhanced,breakable,colback=teal!4,colframe=teal!60!black,boxrule=0.45pt,arc=1.2mm,left=2.7mm,right=2.7mm,top=1.8mm,bottom=1.8mm,before skip=2mm,after skip=2mm]
{\bfseries\sffamily\color{teal}To'g'ri javob:} \textbf{C) ona katalog - joriy obyektni bevosita o'z ichiga olgan katalog}
\par\smallskip
{\bfseries Izoh:} ona katalog - joriy obyektni bevosita o'z ichiga olgan katalog. To'g'ri juftlik: ona katalog - joriy obyektni bevosita o'z ichiga olgan katalog. Manba: CAM6, PDF 178-205; Qo'shimcha sahifa raqami (qaysi PDFga tegishli ekani aniqlanmagan; darslik sahifasi deb tasdiqlanmagan): 91
\end{tcolorbox}
\vspace{2mm}\hrule height 0.2pt\vspace{3mm}

% DB_ID 790b4ac1-3b56-ff60-6ee1-12ec0573cb19
\needspace{5\baselineskip}
{\bfseries\sffamily\color{ink} 407.} {scriptsizecolor{muted} Y1 | review | 790b4ac1}
\par
Fayl va papkalar ustida amallar mavzusidagi vazifada bir nechta faylni konteynerga yig'ish va ko'pincha siqish talab qilindi. Qaysi tushuncha mos?
\begin{enumerate}[label=\Alph*),leftmargin=9mm]
\item copy
\item archive
\item move
\item rename
\end{enumerate}
\begin{tcolorbox}[enhanced,breakable,colback=teal!4,colframe=teal!60!black,boxrule=0.45pt,arc=1.2mm,left=2.7mm,right=2.7mm,top=1.8mm,bottom=1.8mm,before skip=2mm,after skip=2mm]
{\bfseries\sffamily\color{teal}To'g'ri javob:} \textbf{B) archive}
\par\smallskip
{\bfseries Izoh:} archive. Berilgan vazifa archive tushunchasiga mos: bir nechta faylni konteynerga yig'ish va ko'pincha siqish. Manba: ICT6, PDF 116-126; Qo'shimcha sahifa raqami (qaysi PDFga tegishli ekani aniqlanmagan; darslik sahifasi deb tasdiqlanmagan): 98
\end{tcolorbox}
\vspace{2mm}\hrule height 0.2pt\vspace{3mm}

% DB_ID 83495fe6-5c13-1f57-83e2-dfe246ee22c2
\needspace{5\baselineskip}
{\bfseries\sffamily\color{ink} 408.} {scriptsizecolor{muted} Y1 | review | 83495fe6}
\par
Fayl, kengaytma, format va atribut mavzusidagi vazifada faylning read-only yoki hidden kabi boshqaruv xususiyati talab qilindi. Qaysi tushuncha mos?
\begin{enumerate}[label=\Alph*),leftmargin=9mm]
\item atribut
\item fayl
\item kengaytma
\item format
\end{enumerate}
\begin{tcolorbox}[enhanced,breakable,colback=teal!4,colframe=teal!60!black,boxrule=0.45pt,arc=1.2mm,left=2.7mm,right=2.7mm,top=1.8mm,bottom=1.8mm,before skip=2mm,after skip=2mm]
{\bfseries\sffamily\color{teal}To'g'ri javob:} \textbf{A) atribut}
\par\smallskip
{\bfseries Izoh:} atribut. Berilgan vazifa atribut tushunchasiga mos: faylning readonly yoki hidden kabi boshqaruv xususiyati. Manba: CAM6, PDF 178-205; Qo'shimcha sahifa raqami (qaysi PDFga tegishli ekani aniqlanmagan; darslik sahifasi deb tasdiqlanmagan): 86
\end{tcolorbox}
\vspace{2mm}\hrule height 0.2pt\vspace{3mm}

% DB_ID 86d9c5d0-8dc4-0553-bd17-b413502607c4
\needspace{5\baselineskip}
{\bfseries\sffamily\color{ink} 409.} {scriptsizecolor{muted} Y1 | review | 86d9c5d0}
\par
Fayl va papkalar ustida amallar mavzusidagi vazifada obyekt nomini o'zgartirish talab qilindi. Qaysi tushuncha mos?
\begin{enumerate}[label=\Alph*),leftmargin=9mm]
\item copy
\item move
\item filter
\item rename
\end{enumerate}
\begin{tcolorbox}[enhanced,breakable,colback=teal!4,colframe=teal!60!black,boxrule=0.45pt,arc=1.2mm,left=2.7mm,right=2.7mm,top=1.8mm,bottom=1.8mm,before skip=2mm,after skip=2mm]
{\bfseries\sffamily\color{teal}To'g'ri javob:} \textbf{D) rename}
\par\smallskip
{\bfseries Izoh:} rename. Berilgan vazifa rename tushunchasiga mos: obyekt nomini o'zgartirish. Manba: CAM6, PDF 196-205; Qo'shimcha sahifa raqami (qaysi PDFga tegishli ekani aniqlanmagan; darslik sahifasi deb tasdiqlanmagan): 97
\end{tcolorbox}
\vspace{2mm}\hrule height 0.2pt\vspace{3mm}

% DB_ID 88f844d4-b4c8-de76-b5fc-bdf4ccc45920
\needspace{5\baselineskip}
{\bfseries\sffamily\color{ink} 410.} {scriptsizecolor{muted} Y1 | review | 88f844d4}
\par
Fayl va papkalar ustida amallar mavzusidagi vazifada manbani saqlagan holda yangi nusxa yaratish talab qilindi. Qaysi tushuncha mos?
\begin{enumerate}[label=\Alph*),leftmargin=9mm]
\item move
\item copy
\item rename
\item filter
\end{enumerate}
\begin{tcolorbox}[enhanced,breakable,colback=teal!4,colframe=teal!60!black,boxrule=0.45pt,arc=1.2mm,left=2.7mm,right=2.7mm,top=1.8mm,bottom=1.8mm,before skip=2mm,after skip=2mm]
{\bfseries\sffamily\color{teal}To'g'ri javob:} \textbf{B) copy}
\par\smallskip
{\bfseries Izoh:} copy. Berilgan vazifa copy tushunchasiga mos: manbani saqlagan holda yangi nusxa yaratish. Manba: ICT5, PDF 43-49; Qo'shimcha sahifa raqami (qaysi PDFga tegishli ekani aniqlanmagan; darslik sahifasi deb tasdiqlanmagan): 95
\end{tcolorbox}
\vspace{2mm}\hrule height 0.2pt\vspace{3mm}

% DB_ID 9db82814-1e31-28a5-d70e-c9f2bfaf30ed
\needspace{5\baselineskip}
{\bfseries\sffamily\color{ink} 411.} {scriptsizecolor{muted} Y1 | review | 9db82814}
\par
Fayl va papkalar ustida amallar mavzusidagi vazifada yo'qolishdan tiklash uchun alohida joydagi qo'shimcha nusxa talab qilindi. Qaysi tushuncha mos?
\begin{enumerate}[label=\Alph*),leftmargin=9mm]
\item copy
\item move
\item backup
\item rename
\end{enumerate}
\begin{tcolorbox}[enhanced,breakable,colback=teal!4,colframe=teal!60!black,boxrule=0.45pt,arc=1.2mm,left=2.7mm,right=2.7mm,top=1.8mm,bottom=1.8mm,before skip=2mm,after skip=2mm]
{\bfseries\sffamily\color{teal}To'g'ri javob:} \textbf{C) backup}
\par\smallskip
{\bfseries Izoh:} backup. Berilgan vazifa backup tushunchasiga mos: yo'qolishdan tiklash uchun alohida joydagi qo'shimcha nusxa. Manba: ICT5, PDF 43-49; Qo'shimcha sahifa raqami (qaysi PDFga tegishli ekani aniqlanmagan; darslik sahifasi deb tasdiqlanmagan): 99
\end{tcolorbox}
\vspace{2mm}\hrule height 0.2pt\vspace{3mm}

% DB_ID 9fded869-debb-2910-25fc-93ca571e6013
\needspace{5\baselineskip}
{\bfseries\sffamily\color{ink} 412.} {scriptsizecolor{muted} Y1 | review | 9fded869}
\par
Fayl va papkalar ustida amallar mavzusidagi vazifada obyektning joylashuvini o'zgartirish talab qilindi. Qaysi tushuncha mos?
\begin{enumerate}[label=\Alph*),leftmargin=9mm]
\item copy
\item rename
\item move
\item filter
\end{enumerate}
\begin{tcolorbox}[enhanced,breakable,colback=teal!4,colframe=teal!60!black,boxrule=0.45pt,arc=1.2mm,left=2.7mm,right=2.7mm,top=1.8mm,bottom=1.8mm,before skip=2mm,after skip=2mm]
{\bfseries\sffamily\color{teal}To'g'ri javob:} \textbf{C) move}
\par\smallskip
{\bfseries Izoh:} move. Berilgan vazifa move tushunchasiga mos: obyektning joylashuvini o'zgartirish. Manba: MAR, PDF 17-20; Qo'shimcha sahifa raqami (qaysi PDFga tegishli ekani aniqlanmagan; darslik sahifasi deb tasdiqlanmagan): 96
\end{tcolorbox}
\vspace{2mm}\hrule height 0.2pt\vspace{3mm}

% DB_ID aa7b744e-65fe-e8cb-37cd-805f19e5331e
\needspace{5\baselineskip}
{\bfseries\sffamily\color{ink} 413.} {scriptsizecolor{muted} Y1 | review | aa7b744e}
\par
Fayl va papkalar ustida amallar mavzusida "copy" qatnashgan atama-ta'rif juftliklaridan to'g'risini toping.
\begin{enumerate}[label=\Alph*),leftmargin=9mm]
\item copy - obyektning joylashuvini o'zgartirish
\item copy - obyekt nomini o'zgartirish
\item copy - manbani saqlagan holda yangi nusxa yaratish
\item copy - shartga mos obyektlargina ko'rinadigan qilish
\end{enumerate}
\begin{tcolorbox}[enhanced,breakable,colback=teal!4,colframe=teal!60!black,boxrule=0.45pt,arc=1.2mm,left=2.7mm,right=2.7mm,top=1.8mm,bottom=1.8mm,before skip=2mm,after skip=2mm]
{\bfseries\sffamily\color{teal}To'g'ri javob:} \textbf{C) copy - manbani saqlagan holda yangi nusxa yaratish}
\par\smallskip
{\bfseries Izoh:} copy - manbani saqlagan holda yangi nusxa yaratish. To'g'ri juftlik: copy - manbani saqlagan holda yangi nusxa yaratish. Manba: ICT5, PDF 43-49; Qo'shimcha sahifa raqami (qaysi PDFga tegishli ekani aniqlanmagan; darslik sahifasi deb tasdiqlanmagan): 95
\end{tcolorbox}
\vspace{2mm}\hrule height 0.2pt\vspace{3mm}

% DB_ID abc3eb99-3d45-d741-23b4-a2f294a43bc7
\needspace{5\baselineskip}
{\bfseries\sffamily\color{ink} 414.} {scriptsizecolor{muted} Y1 | review | abc3eb99}
\par
Fayl va papkalar ustida amallar mavzusida "archive" qatnashgan atama-ta'rif juftliklaridan to'g'risini toping.
\begin{enumerate}[label=\Alph*),leftmargin=9mm]
\item archive - manbani saqlagan holda yangi nusxa yaratish
\item archive - obyektning joylashuvini o'zgartirish
\item archive - bir nechta faylni konteynerga yig'ish va ko'pincha siqish
\item archive - obyekt nomini o'zgartirish
\end{enumerate}
\begin{tcolorbox}[enhanced,breakable,colback=teal!4,colframe=teal!60!black,boxrule=0.45pt,arc=1.2mm,left=2.7mm,right=2.7mm,top=1.8mm,bottom=1.8mm,before skip=2mm,after skip=2mm]
{\bfseries\sffamily\color{teal}To'g'ri javob:} \textbf{C) archive - bir nechta faylni konteynerga yig'ish va ko'pincha siqish}
\par\smallskip
{\bfseries Izoh:} archive - bir nechta faylni konteynerga yig'ish va ko'pincha siqish. To'g'ri juftlik: archive - bir nechta faylni konteynerga yig'ish va ko'pincha siqish. Manba: ICT6, PDF 116-126; Qo'shimcha sahifa raqami (qaysi PDFga tegishli ekani aniqlanmagan; darslik sahifasi deb tasdiqlanmagan): 99
\end{tcolorbox}
\vspace{2mm}\hrule height 0.2pt\vspace{3mm}

% DB_ID ad33e7e2-49df-50d3-50eb-31d2c71deec0
\needspace{5\baselineskip}
{\bfseries\sffamily\color{ink} 415.} {scriptsizecolor{muted} Y1 | review | ad33e7e2}
\par
Fayl, kengaytma, format va atribut mavzusidagi vazifada fayl nomining formatni ko'rsatuvchi oxirgi qismi talab qilindi. Qaysi tushuncha mos?
\begin{enumerate}[label=\Alph*),leftmargin=9mm]
\item fayl
\item format
\item kengaytma
\item atribut
\end{enumerate}
\begin{tcolorbox}[enhanced,breakable,colback=teal!4,colframe=teal!60!black,boxrule=0.45pt,arc=1.2mm,left=2.7mm,right=2.7mm,top=1.8mm,bottom=1.8mm,before skip=2mm,after skip=2mm]
{\bfseries\sffamily\color{teal}To'g'ri javob:} \textbf{C) kengaytma}
\par\smallskip
{\bfseries Izoh:} kengaytma. Berilgan vazifa kengaytma tushunchasiga mos: fayl nomining formatni ko'rsatuvchi oxirgi qismi. Manba: MAR, PDF 16-20; Qo'shimcha sahifa raqami (qaysi PDFga tegishli ekani aniqlanmagan; darslik sahifasi deb tasdiqlanmagan): 84
\end{tcolorbox}
\vspace{2mm}\hrule height 0.2pt\vspace{3mm}

% DB_ID ba8d55e6-dcf2-041a-d956-1f1736f8a7d5
\needspace{5\baselineskip}
{\bfseries\sffamily\color{ink} 416.} {scriptsizecolor{muted} Y1 | review | ba8d55e6}
\par
Papka, katalog, yo'l va ierarxiya mavzusidagi vazifada boshqa katalog ichida joylashgan katalog talab qilindi. Qaysi tushuncha mos?
\begin{enumerate}[label=\Alph*),leftmargin=9mm]
\item papka
\item ildiz katalog
\item quyi katalog
\item ona katalog
\end{enumerate}
\begin{tcolorbox}[enhanced,breakable,colback=teal!4,colframe=teal!60!black,boxrule=0.45pt,arc=1.2mm,left=2.7mm,right=2.7mm,top=1.8mm,bottom=1.8mm,before skip=2mm,after skip=2mm]
{\bfseries\sffamily\color{teal}To'g'ri javob:} \textbf{C) quyi katalog}
\par\smallskip
{\bfseries Izoh:} quyi katalog. Berilgan vazifa quyi katalog tushunchasiga mos: boshqa katalog ichida joylashgan katalog. Manba: CAM9, PDF 147-149; Qo'shimcha sahifa raqami (qaysi PDFga tegishli ekani aniqlanmagan; darslik sahifasi deb tasdiqlanmagan): 92
\end{tcolorbox}
\vspace{2mm}\hrule height 0.2pt\vspace{3mm}

% DB_ID cb274639-ec6d-e24a-62b3-604aee372532
\needspace{5\baselineskip}
{\bfseries\sffamily\color{ink} 417.} {scriptsizecolor{muted} Y1 | review | cb274639}
\par
Papka, katalog, yo'l va ierarxiya mavzusidagi vazifada fayl va quyi papkalarni guruhlovchi konteyner talab qilindi. Qaysi tushuncha mos?
\begin{enumerate}[label=\Alph*),leftmargin=9mm]
\item ildiz katalog
\item ona katalog
\item quyi katalog
\item papka
\end{enumerate}
\begin{tcolorbox}[enhanced,breakable,colback=teal!4,colframe=teal!60!black,boxrule=0.45pt,arc=1.2mm,left=2.7mm,right=2.7mm,top=1.8mm,bottom=1.8mm,before skip=2mm,after skip=2mm]
{\bfseries\sffamily\color{teal}To'g'ri javob:} \textbf{D) papka}
\par\smallskip
{\bfseries Izoh:} papka. Berilgan vazifa papka tushunchasiga mos: fayl va quyi papkalarni guruhlovchi konteyner. Manba: ICT5, PDF 42-45; Qo'shimcha sahifa raqami (qaysi PDFga tegishli ekani aniqlanmagan; darslik sahifasi deb tasdiqlanmagan): 89
\end{tcolorbox}
\vspace{2mm}\hrule height 0.2pt\vspace{3mm}

% DB_ID cce4e6fc-aa97-f1ed-6561-efd004a282ff
\needspace{5\baselineskip}
{\bfseries\sffamily\color{ink} 418.} {scriptsizecolor{muted} Y1 | review | cce4e6fc}
\par
Papka, katalog, yo'l va ierarxiya mavzusidagi vazifada joriy katalogga nisbatan berilgan joylashuv talab qilindi. Qaysi tushuncha mos?
\begin{enumerate}[label=\Alph*),leftmargin=9mm]
\item nisbiy yo'l
\item papka
\item ildiz katalog
\item ona katalog
\end{enumerate}
\begin{tcolorbox}[enhanced,breakable,colback=teal!4,colframe=teal!60!black,boxrule=0.45pt,arc=1.2mm,left=2.7mm,right=2.7mm,top=1.8mm,bottom=1.8mm,before skip=2mm,after skip=2mm]
{\bfseries\sffamily\color{teal}To'g'ri javob:} \textbf{A) nisbiy yo'l}
\par\smallskip
{\bfseries Izoh:} nisbiy yo'l. Berilgan vazifa nisbiy yo'l tushunchasiga mos: joriy katalogga nisbatan berilgan joylashuv. Manba: ICT5, PDF 42-45; Qo'shimcha sahifa raqami (qaysi PDFga tegishli ekani aniqlanmagan; darslik sahifasi deb tasdiqlanmagan): 94
\end{tcolorbox}
\vspace{2mm}\hrule height 0.2pt\vspace{3mm}

% DB_ID d04e5193-85bc-90e6-7518-dd31235b21ab
\needspace{5\baselineskip}
{\bfseries\sffamily\color{ink} 419.} {scriptsizecolor{muted} Y1 | review | d04e5193}
\par
Papka, katalog, yo'l va ierarxiya mavzusida "mutlaq yo'l" qatnashgan atama-ta'rif juftliklaridan to'g'risini toping.
\begin{enumerate}[label=\Alph*),leftmargin=9mm]
\item mutlaq yo'l - ildizdan boshlanadigan to'liq joylashuv
\item mutlaq yo'l - fayl va quyi papkalarni guruhlovchi konteyner
\item mutlaq yo'l - fayl tizimi daraxtining eng yuqori katalogi
\item mutlaq yo'l - joriy obyektni bevosita o'z ichiga olgan katalog
\end{enumerate}
\begin{tcolorbox}[enhanced,breakable,colback=teal!4,colframe=teal!60!black,boxrule=0.45pt,arc=1.2mm,left=2.7mm,right=2.7mm,top=1.8mm,bottom=1.8mm,before skip=2mm,after skip=2mm]
{\bfseries\sffamily\color{teal}To'g'ri javob:} \textbf{A) mutlaq yo'l - ildizdan boshlanadigan to'liq joylashuv}
\par\smallskip
{\bfseries Izoh:} mutlaq yo'l - ildizdan boshlanadigan to'liq joylashuv. To'g'ri juftlik: mutlaq yo'l - ildizdan boshlanadigan to'liq joylashuv. Manba: ICT6, PDF 116-124; Qo'shimcha sahifa raqami (qaysi PDFga tegishli ekani aniqlanmagan; darslik sahifasi deb tasdiqlanmagan): 93
\end{tcolorbox}
\vspace{2mm}\hrule height 0.2pt\vspace{3mm}

% DB_ID e0ab763d-4f98-a12a-323d-2bbbbaab5c8d
\needspace{5\baselineskip}
{\bfseries\sffamily\color{ink} 420.} {scriptsizecolor{muted} Y1 | review | e0ab763d}
\par
Fayl, kengaytma, format va atribut mavzusidagi vazifada fayl mazmuni haqida tavsiflovchi ma'lumot talab qilindi. Qaysi tushuncha mos?
\begin{enumerate}[label=\Alph*),leftmargin=9mm]
\item fayl
\item metama'lumot
\item kengaytma
\item format
\end{enumerate}
\begin{tcolorbox}[enhanced,breakable,colback=teal!4,colframe=teal!60!black,boxrule=0.45pt,arc=1.2mm,left=2.7mm,right=2.7mm,top=1.8mm,bottom=1.8mm,before skip=2mm,after skip=2mm]
{\bfseries\sffamily\color{teal}To'g'ri javob:} \textbf{B) metama'lumot}
\par\smallskip
{\bfseries Izoh:} metama'lumot. Berilgan vazifa metama'lumot tushunchasiga mos: fayl mazmuni haqida tavsiflovchi ma'lumot. Manba: CAM7, PDF 41-54; Qo'shimcha sahifa raqami (qaysi PDFga tegishli ekani aniqlanmagan; darslik sahifasi deb tasdiqlanmagan): 87
\end{tcolorbox}
\vspace{2mm}\hrule height 0.2pt\vspace{3mm}

% DB_ID f11003ad-bb83-686a-2616-9dab37576b6f
\needspace{5\baselineskip}
{\bfseries\sffamily\color{ink} 421.} {scriptsizecolor{muted} Y1 | review | f11003ad}
\par
Papka, katalog, yo'l va ierarxiya mavzusidagi vazifada fayl tizimi daraxtining eng yuqori katalogi talab qilindi. Qaysi tushuncha mos?
\begin{enumerate}[label=\Alph*),leftmargin=9mm]
\item ildiz katalog
\item papka
\item ona katalog
\item quyi katalog
\end{enumerate}
\begin{tcolorbox}[enhanced,breakable,colback=teal!4,colframe=teal!60!black,boxrule=0.45pt,arc=1.2mm,left=2.7mm,right=2.7mm,top=1.8mm,bottom=1.8mm,before skip=2mm,after skip=2mm]
{\bfseries\sffamily\color{teal}To'g'ri javob:} \textbf{A) ildiz katalog}
\par\smallskip
{\bfseries Izoh:} ildiz katalog. Berilgan vazifa ildiz katalog tushunchasiga mos: fayl tizimi daraxtining eng yuqori katalogi. Manba: MAR, PDF 16-20; Qo'shimcha sahifa raqami (qaysi PDFga tegishli ekani aniqlanmagan; darslik sahifasi deb tasdiqlanmagan): 90
\end{tcolorbox}
\vspace{2mm}\hrule height 0.2pt\vspace{3mm}

% DB_ID ffc17a90-f1b3-866d-f417-d85c11e0f4b1
\needspace{5\baselineskip}
{\bfseries\sffamily\color{ink} 422.} {scriptsizecolor{muted} Y1 | review | ffc17a90}
\par
Fayl, kengaytma, format va atribut mavzusida "format" qatnashgan atama-ta'rif juftliklaridan to'g'risini toping.
\begin{enumerate}[label=\Alph*),leftmargin=9mm]
\item format - fayl ichidagi ma'lumot tuzilishi va kodlash qoidasi
\item format - nom bilan saqlanadigan ma'lumotlar birligi
\item format - fayl nomining formatni ko'rsatuvchi oxirgi qismi
\item format - faylning read-only yoki hidden kabi boshqaruv xususiyati
\end{enumerate}
\begin{tcolorbox}[enhanced,breakable,colback=teal!4,colframe=teal!60!black,boxrule=0.45pt,arc=1.2mm,left=2.7mm,right=2.7mm,top=1.8mm,bottom=1.8mm,before skip=2mm,after skip=2mm]
{\bfseries\sffamily\color{teal}To'g'ri javob:} \textbf{A) format - fayl ichidagi ma'lumot tuzilishi va kodlash qoidasi}
\par\smallskip
{\bfseries Izoh:} format - fayl ichidagi ma'lumot tuzilishi va kodlash qoidasi. To'g'ri juftlik: format - fayl ichidagi ma'lumot tuzilishi va kodlash qoidasi. Manba: CAM9, PDF 147-149; Qo'shimcha sahifa raqami (qaysi PDFga tegishli ekani aniqlanmagan; darslik sahifasi deb tasdiqlanmagan): 85
\end{tcolorbox}
\vspace{2mm}\hrule height 0.2pt\vspace{3mm}

% DB_ID 296625f0-1efe-4253-96ec-1d86454b28f2
\needspace{5\baselineskip}
{\bfseries\sffamily\color{ink} 423.} {scriptsizecolor{muted} Y1 | review | 296625f0}
\par
C:\textbackslash{}Users\textbackslash{}Ali\textbackslash{}Documents\textbackslash{}reja.docx yo'lida faylning bevosita ona papkasi qaysi?
\begin{enumerate}[label=\Alph*),leftmargin=9mm]
\item Users
\item Ali
\item Documents
\item C:
\end{enumerate}
\begin{tcolorbox}[enhanced,breakable,colback=teal!4,colframe=teal!60!black,boxrule=0.45pt,arc=1.2mm,left=2.7mm,right=2.7mm,top=1.8mm,bottom=1.8mm,before skip=2mm,after skip=2mm]
{\bfseries\sffamily\color{teal}To'g'ri javob:} \textbf{C) Documents}
\par\smallskip
{\bfseries Izoh:} Documents. Faylni bevosita o'z ichiga olgan katalog Documents. Manba: ICT5, PDF 42-45; Qo'shimcha sahifa raqami (qaysi PDFga tegishli ekani aniqlanmagan; darslik sahifasi deb tasdiqlanmagan): 94
\end{tcolorbox}
\vspace{2mm}\hrule height 0.2pt\vspace{3mm}

% DB_ID 3ae200a3-ac30-adff-e85c-802ce52f11ca
\needspace{5\baselineskip}
{\bfseries\sffamily\color{ink} 424.} {scriptsizecolor{muted} Y1 | review | 3ae200a3}
\par
Papka, katalog, yo'l va ierarxiya mavzusida "ildiz katalog" qatnashgan atama-ta'rif juftliklaridan to'g'risini toping.
\begin{enumerate}[label=\Alph*),leftmargin=9mm]
\item ildiz katalog - fayl va quyi papkalarni guruhlovchi konteyner
\item ildiz katalog - fayl tizimi daraxtining eng yuqori katalogi
\item ildiz katalog - joriy obyektni bevosita o'z ichiga olgan katalog
\item ildiz katalog - boshqa katalog ichida joylashgan katalog
\end{enumerate}
\begin{tcolorbox}[enhanced,breakable,colback=teal!4,colframe=teal!60!black,boxrule=0.45pt,arc=1.2mm,left=2.7mm,right=2.7mm,top=1.8mm,bottom=1.8mm,before skip=2mm,after skip=2mm]
{\bfseries\sffamily\color{teal}To'g'ri javob:} \textbf{B) ildiz katalog - fayl tizimi daraxtining eng yuqori katalogi}
\par\smallskip
{\bfseries Izoh:} ildiz katalog - fayl tizimi daraxtining eng yuqori katalogi. To'g'ri juftlik: ildiz katalog - fayl tizimi daraxtining eng yuqori katalogi. Manba: MAR, PDF 16-20; Qo'shimcha sahifa raqami (qaysi PDFga tegishli ekani aniqlanmagan; darslik sahifasi deb tasdiqlanmagan): 90
\end{tcolorbox}
\vspace{2mm}\hrule height 0.2pt\vspace{3mm}

% DB_ID 3bdb3932-6fe3-0405-951e-29d125919968
\needspace{5\baselineskip}
{\bfseries\sffamily\color{ink} 425.} {scriptsizecolor{muted} Y1 | review | 3bdb3932}
\par
Fayl va papkalar ustida amallar mavzusida "move" qatnashgan atama-ta'rif juftliklaridan to'g'risini toping.
\begin{enumerate}[label=\Alph*),leftmargin=9mm]
\item move - manbani saqlagan holda yangi nusxa yaratish
\item move - obyekt nomini o'zgartirish
\item move - shartga mos obyektlargina ko'rinadigan qilish
\item move - obyektning joylashuvini o'zgartirish
\end{enumerate}
\begin{tcolorbox}[enhanced,breakable,colback=teal!4,colframe=teal!60!black,boxrule=0.45pt,arc=1.2mm,left=2.7mm,right=2.7mm,top=1.8mm,bottom=1.8mm,before skip=2mm,after skip=2mm]
{\bfseries\sffamily\color{teal}To'g'ri javob:} \textbf{D) move - obyektning joylashuvini o'zgartirish}
\par\smallskip
{\bfseries Izoh:} move - obyektning joylashuvini o'zgartirish. To'g'ri juftlik: move - obyektning joylashuvini o'zgartirish. Manba: MAR, PDF 17-20; Qo'shimcha sahifa raqami (qaysi PDFga tegishli ekani aniqlanmagan; darslik sahifasi deb tasdiqlanmagan): 96
\end{tcolorbox}
\vspace{2mm}\hrule height 0.2pt\vspace{3mm}

% DB_ID 44cc06e4-3c57-8516-8ce5-7a97278e90a1
\needspace{5\baselineskip}
{\bfseries\sffamily\color{ink} 426.} {scriptsizecolor{muted} Y1 | review | 44cc06e4}
\par
Papka, katalog, yo'l va ierarxiya bo'yicha qatorlar ichidan "papka" noto'g'ri izohlanganini aniqlang.
\begin{enumerate}[label=\Alph*),leftmargin=9mm]
\item ona katalog - joriy obyektni bevosita o'z ichiga olgan katalog
\item papka - fayl tizimi daraxtining eng yuqori katalogi
\item quyi katalog - boshqa katalog ichida joylashgan katalog
\item mutlaq yo'l - ildizdan boshlanadigan to'liq joylashuv
\end{enumerate}
\begin{tcolorbox}[enhanced,breakable,colback=teal!4,colframe=teal!60!black,boxrule=0.45pt,arc=1.2mm,left=2.7mm,right=2.7mm,top=1.8mm,bottom=1.8mm,before skip=2mm,after skip=2mm]
{\bfseries\sffamily\color{teal}To'g'ri javob:} \textbf{B) papka - fayl tizimi daraxtining eng yuqori katalogi}
\par\smallskip
{\bfseries Izoh:} papka - fayl tizimi daraxtining eng yuqori katalogi. papka uchun to'g'ri ta'rif: fayl va quyi papkalarni guruhlovchi konteyner. Manba: ICT5, PDF 42-45; Qo'shimcha sahifa raqami (qaysi PDFga tegishli ekani aniqlanmagan; darslik sahifasi deb tasdiqlanmagan): 89
\end{tcolorbox}
\vspace{2mm}\hrule height 0.2pt\vspace{3mm}

% DB_ID 4c4abd25-e521-d258-e53d-9f50a9d5f601
\needspace{5\baselineskip}
{\bfseries\sffamily\color{ink} 427.} {scriptsizecolor{muted} Y1 | review | 4c4abd25}
\par
Fayl, kengaytma, format va atribut bo'yicha qatorlar ichidan "metama'lumot" noto'g'ri izohlanganini aniqlang.
\begin{enumerate}[label=\Alph*),leftmargin=9mm]
\item kengaytma - fayl nomining formatni ko'rsatuvchi oxirgi qismi
\item format - fayl ichidagi ma'lumot tuzilishi va kodlash qoidasi
\item atribut - faylning read-only yoki hidden kabi boshqaruv xususiyati
\item metama'lumot - nom bilan saqlanadigan ma'lumotlar birligi
\end{enumerate}
\begin{tcolorbox}[enhanced,breakable,colback=teal!4,colframe=teal!60!black,boxrule=0.45pt,arc=1.2mm,left=2.7mm,right=2.7mm,top=1.8mm,bottom=1.8mm,before skip=2mm,after skip=2mm]
{\bfseries\sffamily\color{teal}To'g'ri javob:} \textbf{D) metama'lumot - nom bilan saqlanadigan ma'lumotlar birligi}
\par\smallskip
{\bfseries Izoh:} metama'lumot - nom bilan saqlanadigan ma'lumotlar birligi. metama'lumot uchun to'g'ri ta'rif: fayl mazmuni haqida tavsiflovchi ma'lumot. Manba: CAM7, PDF 41-54; Qo'shimcha sahifa raqami (qaysi PDFga tegishli ekani aniqlanmagan; darslik sahifasi deb tasdiqlanmagan): 87
\end{tcolorbox}
\vspace{2mm}\hrule height 0.2pt\vspace{3mm}

% DB_ID 523dee48-77ad-b372-495d-10037e962abc
\needspace{5\baselineskip}
{\bfseries\sffamily\color{ink} 428.} {scriptsizecolor{muted} Y1 | review | 523dee48}
\par
Fayl va papkalar ustida amallar mavzusida "backup" qatnashgan atama-ta'rif juftliklaridan to'g'risini toping.
\begin{enumerate}[label=\Alph*),leftmargin=9mm]
\item backup - manbani saqlagan holda yangi nusxa yaratish
\item backup - obyektning joylashuvini o'zgartirish
\item backup - obyekt nomini o'zgartirish
\item backup - yo'qolishdan tiklash uchun alohida joydagi qo'shimcha nusxa
\end{enumerate}
\begin{tcolorbox}[enhanced,breakable,colback=teal!4,colframe=teal!60!black,boxrule=0.45pt,arc=1.2mm,left=2.7mm,right=2.7mm,top=1.8mm,bottom=1.8mm,before skip=2mm,after skip=2mm]
{\bfseries\sffamily\color{teal}To'g'ri javob:} \textbf{D) backup - yo'qolishdan tiklash uchun alohida joydagi qo'shimcha nusxa}
\par\smallskip
{\bfseries Izoh:} backup - yo'qolishdan tiklash uchun alohida joydagi qo'shimcha nusxa. To'g'ri juftlik: backup - yo'qolishdan tiklash uchun alohida joydagi qo'shimcha nusxa. Manba: ICT5, PDF 43-49; Qo'shimcha sahifa raqami (qaysi PDFga tegishli ekani aniqlanmagan; darslik sahifasi deb tasdiqlanmagan): 100
\end{tcolorbox}
\vspace{2mm}\hrule height 0.2pt\vspace{3mm}

% DB_ID 5b70ccb6-0c3a-e207-4d8d-3d1e7825bd0d
\needspace{5\baselineskip}
{\bfseries\sffamily\color{ink} 429.} {scriptsizecolor{muted} Y1 | review | 5b70ccb6}
\par
Soddalashtirilgan fayl niqobi qoidasida savol belgisi (?) aynan bitta belgiga mos keladi deb olinsin. hisob\_??.txt niqobiga qaysi fayl nomi mos?
\begin{enumerate}[label=\Alph*),leftmargin=9mm]
\item hisob\_7.txt
\item hisob\_07.txt
\item hisob\_123.txt
\item hisob\_07.docx
\end{enumerate}
\begin{tcolorbox}[enhanced,breakable,colback=teal!4,colframe=teal!60!black,boxrule=0.45pt,arc=1.2mm,left=2.7mm,right=2.7mm,top=1.8mm,bottom=1.8mm,before skip=2mm,after skip=2mm]
{\bfseries\sffamily\color{teal}To'g'ri javob:} \textbf{B) hisob\_07.txt}
\par\smallskip
{\bfseries Izoh:} Ushbu savolda belgilangan soddalashtirilgan qoida bo'yicha har bir ? aynan bitta belgini bildiradi; hisob\_07.txt mos. Haqiqiy Windows/DOS dir niqobida ? ayrim vaziyatlarda kamroq belgiga ham mos kelishi mumkin; shu sabab oldingi umumiy ifoda noaniq edi. Microsoft: https://learn.microsoft.com/en-us/windows-server/administration/windows-commands/dir; https://learn.microsoft.com/en-us/openspecs/windows\_protocols/ms-cifs/dc92d939-ec45-40c8-96e5-4c4091e4ab43. AI texnik audit, mustaqil inson eksperti tasdig'i emas.
\end{tcolorbox}
\vspace{2mm}\hrule height 0.2pt\vspace{3mm}

% DB_ID 62e1caff-e68a-73bb-4fcf-d7833f2848b5
\needspace{5\baselineskip}
{\bfseries\sffamily\color{ink} 430.} {scriptsizecolor{muted} Y1 | review | 62e1caff}
\par
Fayl, kengaytma, format va atribut bo'yicha qatorlar ichidan "format" noto'g'ri izohlanganini aniqlang.
\begin{enumerate}[label=\Alph*),leftmargin=9mm]
\item kengaytma - fayl nomining formatni ko'rsatuvchi oxirgi qismi
\item format - nom bilan saqlanadigan ma'lumotlar birligi
\item atribut - faylning read-only yoki hidden kabi boshqaruv xususiyati
\item metama'lumot - fayl mazmuni haqida tavsiflovchi ma'lumot
\end{enumerate}
\begin{tcolorbox}[enhanced,breakable,colback=teal!4,colframe=teal!60!black,boxrule=0.45pt,arc=1.2mm,left=2.7mm,right=2.7mm,top=1.8mm,bottom=1.8mm,before skip=2mm,after skip=2mm]
{\bfseries\sffamily\color{teal}To'g'ri javob:} \textbf{B) format - nom bilan saqlanadigan ma'lumotlar birligi}
\par\smallskip
{\bfseries Izoh:} format - nom bilan saqlanadigan ma'lumotlar birligi. format uchun to'g'ri ta'rif: fayl ichidagi ma'lumot tuzilishi va kodlash qoidasi. Manba: CAM9, PDF 147-149; Qo'shimcha sahifa raqami (qaysi PDFga tegishli ekani aniqlanmagan; darslik sahifasi deb tasdiqlanmagan): 85
\end{tcolorbox}
\vspace{2mm}\hrule height 0.2pt\vspace{3mm}

% DB_ID 6e7cb003-f334-2ee9-12cd-581bd261ff19
\needspace{5\baselineskip}
{\bfseries\sffamily\color{ink} 431.} {scriptsizecolor{muted} Y1 | review | 6e7cb003}
\par
Fayl va papkalar ustida amallar bo'yicha qatorlar ichidan "archive" noto'g'ri izohlanganini aniqlang.
\begin{enumerate}[label=\Alph*),leftmargin=9mm]
\item move - obyektning joylashuvini o'zgartirish
\item rename - obyekt nomini o'zgartirish
\item filter - shartga mos obyektlargina ko'rinadigan qilish
\item archive - manbani saqlagan holda yangi nusxa yaratish
\end{enumerate}
\begin{tcolorbox}[enhanced,breakable,colback=teal!4,colframe=teal!60!black,boxrule=0.45pt,arc=1.2mm,left=2.7mm,right=2.7mm,top=1.8mm,bottom=1.8mm,before skip=2mm,after skip=2mm]
{\bfseries\sffamily\color{teal}To'g'ri javob:} \textbf{D) archive - manbani saqlagan holda yangi nusxa yaratish}
\par\smallskip
{\bfseries Izoh:} archive - manbani saqlagan holda yangi nusxa yaratish. archive uchun to'g'ri ta'rif: bir nechta faylni konteynerga yig'ish va ko'pincha siqish. Manba: ICT6, PDF 116-126; Qo'shimcha sahifa raqami (qaysi PDFga tegishli ekani aniqlanmagan; darslik sahifasi deb tasdiqlanmagan): 99
\end{tcolorbox}
\vspace{2mm}\hrule height 0.2pt\vspace{3mm}

% DB_ID 71e9f33f-9e77-4f74-2467-9244287aec37
\needspace{5\baselineskip}
{\bfseries\sffamily\color{ink} 432.} {scriptsizecolor{muted} Y1 | review | 71e9f33f}
\par
Fayl, kengaytma, format va atribut mavzusida "konvertatsiya" qatnashgan atama-ta'rif juftliklaridan to'g'risini toping.
\begin{enumerate}[label=\Alph*),leftmargin=9mm]
\item konvertatsiya - nom bilan saqlanadigan ma'lumotlar birligi
\item konvertatsiya - fayl nomining formatni ko'rsatuvchi oxirgi qismi
\item konvertatsiya - fayl ichidagi ma'lumot tuzilishi va kodlash qoidasi
\item konvertatsiya - ma'lumotni bir format tuzilishidan boshqasiga o'tkazish
\end{enumerate}
\begin{tcolorbox}[enhanced,breakable,colback=teal!4,colframe=teal!60!black,boxrule=0.45pt,arc=1.2mm,left=2.7mm,right=2.7mm,top=1.8mm,bottom=1.8mm,before skip=2mm,after skip=2mm]
{\bfseries\sffamily\color{teal}To'g'ri javob:} \textbf{D) konvertatsiya - ma'lumotni bir format tuzilishidan boshqasiga o'tkazish}
\par\smallskip
{\bfseries Izoh:} konvertatsiya - ma'lumotni bir format tuzilishidan boshqasiga o'tkazish. To'g'ri juftlik: konvertatsiya - ma'lumotni bir format tuzilishidan boshqasiga o'tkazish. Manba: ICT5, PDF 42-45; Qo'shimcha sahifa raqami (qaysi PDFga tegishli ekani aniqlanmagan; darslik sahifasi deb tasdiqlanmagan): 88
\end{tcolorbox}
\vspace{2mm}\hrule height 0.2pt\vspace{3mm}

% DB_ID 749fae2a-d7a7-6240-2dc3-666554c85f6d
\needspace{5\baselineskip}
{\bfseries\sffamily\color{ink} 433.} {scriptsizecolor{muted} Y1 | review | 749fae2a}
\par
Joriy papka C:\textbackslash{}Users\textbackslash{}Ali bo'lsa, Documents\textbackslash{}reja.docx qanday yo'l?
\begin{enumerate}[label=\Alph*),leftmargin=9mm]
\item ildiz yo'l
\item URL manzil
\item fayl kengaytmasi
\item nisbiy yo'l
\end{enumerate}
\begin{tcolorbox}[enhanced,breakable,colback=teal!4,colframe=teal!60!black,boxrule=0.45pt,arc=1.2mm,left=2.7mm,right=2.7mm,top=1.8mm,bottom=1.8mm,before skip=2mm,after skip=2mm]
{\bfseries\sffamily\color{teal}To'g'ri javob:} \textbf{D) nisbiy yo'l}
\par\smallskip
{\bfseries Izoh:} nisbiy yo'l. Yo'l joriy papkaga nisbatan berilgan va disk ildizidan boshlanmagan. Manba: ICT5, PDF 42-45; Qo'shimcha sahifa raqami (qaysi PDFga tegishli ekani aniqlanmagan; darslik sahifasi deb tasdiqlanmagan): 94
\end{tcolorbox}
\vspace{2mm}\hrule height 0.2pt\vspace{3mm}

% DB_ID 80d065d9-5610-6317-7729-0713512343b5
\needspace{5\baselineskip}
{\bfseries\sffamily\color{ink} 434.} {scriptsizecolor{muted} Y1 | review | 80d065d9}
\par
Papka, katalog, yo'l va ierarxiya bo'yicha qatorlar ichidan "ona katalog" noto'g'ri izohlanganini aniqlang.
\begin{enumerate}[label=\Alph*),leftmargin=9mm]
\item ildiz katalog - fayl tizimi daraxtining eng yuqori katalogi
\item quyi katalog - boshqa katalog ichida joylashgan katalog
\item mutlaq yo'l - ildizdan boshlanadigan to'liq joylashuv
\item ona katalog - fayl va quyi papkalarni guruhlovchi konteyner
\end{enumerate}
\begin{tcolorbox}[enhanced,breakable,colback=teal!4,colframe=teal!60!black,boxrule=0.45pt,arc=1.2mm,left=2.7mm,right=2.7mm,top=1.8mm,bottom=1.8mm,before skip=2mm,after skip=2mm]
{\bfseries\sffamily\color{teal}To'g'ri javob:} \textbf{D) ona katalog - fayl va quyi papkalarni guruhlovchi konteyner}
\par\smallskip
{\bfseries Izoh:} ona katalog - fayl va quyi papkalarni guruhlovchi konteyner. ona katalog uchun to'g'ri ta'rif: joriy obyektni bevosita o'z ichiga olgan katalog. Manba: CAM6, PDF 178-205; Qo'shimcha sahifa raqami (qaysi PDFga tegishli ekani aniqlanmagan; darslik sahifasi deb tasdiqlanmagan): 91
\end{tcolorbox}
\vspace{2mm}\hrule height 0.2pt\vspace{3mm}

% DB_ID 864adaae-214d-0e99-2e36-04bb9caf4d8d
\needspace{5\baselineskip}
{\bfseries\sffamily\color{ink} 435.} {scriptsizecolor{muted} Y1 | review | 864adaae}
\par
Fayl va papkalar ustida amallar bo'yicha qatorlar ichidan "rename" noto'g'ri izohlanganini aniqlang.
\begin{enumerate}[label=\Alph*),leftmargin=9mm]
\item move - obyektning joylashuvini o'zgartirish
\item rename - manbani saqlagan holda yangi nusxa yaratish
\item filter - shartga mos obyektlargina ko'rinadigan qilish
\item archive - bir nechta faylni konteynerga yig'ish va ko'pincha siqish
\end{enumerate}
\begin{tcolorbox}[enhanced,breakable,colback=teal!4,colframe=teal!60!black,boxrule=0.45pt,arc=1.2mm,left=2.7mm,right=2.7mm,top=1.8mm,bottom=1.8mm,before skip=2mm,after skip=2mm]
{\bfseries\sffamily\color{teal}To'g'ri javob:} \textbf{B) rename - manbani saqlagan holda yangi nusxa yaratish}
\par\smallskip
{\bfseries Izoh:} rename - manbani saqlagan holda yangi nusxa yaratish. rename uchun to'g'ri ta'rif: obyekt nomini o'zgartirish. Manba: CAM6, PDF 196-205; Qo'shimcha sahifa raqami (qaysi PDFga tegishli ekani aniqlanmagan; darslik sahifasi deb tasdiqlanmagan): 97
\end{tcolorbox}
\vspace{2mm}\hrule height 0.2pt\vspace{3mm}

% DB_ID 8fd083b7-9d45-07c0-3d89-5357adbb8532
\needspace{5\baselineskip}
{\bfseries\sffamily\color{ink} 436.} {scriptsizecolor{muted} Y1 | review | 8fd083b7}
\par
Papka, katalog, yo'l va ierarxiya mavzusida "nisbiy yo'l" qatnashgan atama-ta'rif juftliklaridan to'g'risini toping.
\begin{enumerate}[label=\Alph*),leftmargin=9mm]
\item nisbiy yo'l - fayl va quyi papkalarni guruhlovchi konteyner
\item nisbiy yo'l - joriy katalogga nisbatan berilgan joylashuv
\item nisbiy yo'l - fayl tizimi daraxtining eng yuqori katalogi
\item nisbiy yo'l - joriy obyektni bevosita o'z ichiga olgan katalog
\end{enumerate}
\begin{tcolorbox}[enhanced,breakable,colback=teal!4,colframe=teal!60!black,boxrule=0.45pt,arc=1.2mm,left=2.7mm,right=2.7mm,top=1.8mm,bottom=1.8mm,before skip=2mm,after skip=2mm]
{\bfseries\sffamily\color{teal}To'g'ri javob:} \textbf{B) nisbiy yo'l - joriy katalogga nisbatan berilgan joylashuv}
\par\smallskip
{\bfseries Izoh:} nisbiy yo'l - joriy katalogga nisbatan berilgan joylashuv. To'g'ri juftlik: nisbiy yo'l - joriy katalogga nisbatan berilgan joylashuv. Manba: ICT5, PDF 42-45; Qo'shimcha sahifa raqami (qaysi PDFga tegishli ekani aniqlanmagan; darslik sahifasi deb tasdiqlanmagan): 94
\end{tcolorbox}
\vspace{2mm}\hrule height 0.2pt\vspace{3mm}

% DB_ID 93c6f207-f3d0-c31c-0930-9428582b0398
\needspace{5\baselineskip}
{\bfseries\sffamily\color{ink} 437.} {scriptsizecolor{muted} Y1 | review | 93c6f207}
\par
Fayl, kengaytma, format va atribut mavzusida "kengaytma" qatnashgan atama-ta'rif juftliklaridan to'g'risini toping.
\begin{enumerate}[label=\Alph*),leftmargin=9mm]
\item kengaytma - nom bilan saqlanadigan ma'lumotlar birligi
\item kengaytma - fayl ichidagi ma'lumot tuzilishi va kodlash qoidasi
\item kengaytma - faylning read-only yoki hidden kabi boshqaruv xususiyati
\item kengaytma - fayl nomining formatni ko'rsatuvchi oxirgi qismi
\end{enumerate}
\begin{tcolorbox}[enhanced,breakable,colback=teal!4,colframe=teal!60!black,boxrule=0.45pt,arc=1.2mm,left=2.7mm,right=2.7mm,top=1.8mm,bottom=1.8mm,before skip=2mm,after skip=2mm]
{\bfseries\sffamily\color{teal}To'g'ri javob:} \textbf{D) kengaytma - fayl nomining formatni ko'rsatuvchi oxirgi qismi}
\par\smallskip
{\bfseries Izoh:} kengaytma - fayl nomining formatni ko'rsatuvchi oxirgi qismi. To'g'ri juftlik: kengaytma - fayl nomining formatni ko'rsatuvchi oxirgi qismi. Manba: MAR, PDF 16-20; Qo'shimcha sahifa raqami (qaysi PDFga tegishli ekani aniqlanmagan; darslik sahifasi deb tasdiqlanmagan): 84
\end{tcolorbox}
\vspace{2mm}\hrule height 0.2pt\vspace{3mm}

% DB_ID a0c53f9f-ec12-4f71-e68a-d859edfde223
\needspace{5\baselineskip}
{\bfseries\sffamily\color{ink} 438.} {scriptsizecolor{muted} Y1 | review | a0c53f9f}
\par
Fayl va papkalar ustida amallar bo'yicha qatorlar ichidan "copy" noto'g'ri izohlanganini aniqlang.
\begin{enumerate}[label=\Alph*),leftmargin=9mm]
\item rename - obyekt nomini o'zgartirish
\item filter - shartga mos obyektlargina ko'rinadigan qilish
\item archive - bir nechta faylni konteynerga yig'ish va ko'pincha siqish
\item copy - obyektning joylashuvini o'zgartirish
\end{enumerate}
\begin{tcolorbox}[enhanced,breakable,colback=teal!4,colframe=teal!60!black,boxrule=0.45pt,arc=1.2mm,left=2.7mm,right=2.7mm,top=1.8mm,bottom=1.8mm,before skip=2mm,after skip=2mm]
{\bfseries\sffamily\color{teal}To'g'ri javob:} \textbf{D) copy - obyektning joylashuvini o'zgartirish}
\par\smallskip
{\bfseries Izoh:} copy - obyektning joylashuvini o'zgartirish. copy uchun to'g'ri ta'rif: manbani saqlagan holda yangi nusxa yaratish. Manba: ICT5, PDF 43-49; Qo'shimcha sahifa raqami (qaysi PDFga tegishli ekani aniqlanmagan; darslik sahifasi deb tasdiqlanmagan): 95
\end{tcolorbox}
\vspace{2mm}\hrule height 0.2pt\vspace{3mm}

% DB_ID a94acb81-9c6a-2c25-24a4-b53a43416558
\needspace{5\baselineskip}
{\bfseries\sffamily\color{ink} 439.} {scriptsizecolor{muted} Y1 | review | a94acb81}
\par
Fayl, kengaytma, format va atribut mavzusida "atribut" qatnashgan atama-ta'rif juftliklaridan to'g'risini toping.
\begin{enumerate}[label=\Alph*),leftmargin=9mm]
\item atribut - nom bilan saqlanadigan ma'lumotlar birligi
\item atribut - faylning read-only yoki hidden kabi boshqaruv xususiyati
\item atribut - fayl nomining formatni ko'rsatuvchi oxirgi qismi
\item atribut - fayl ichidagi ma'lumot tuzilishi va kodlash qoidasi
\end{enumerate}
\begin{tcolorbox}[enhanced,breakable,colback=teal!4,colframe=teal!60!black,boxrule=0.45pt,arc=1.2mm,left=2.7mm,right=2.7mm,top=1.8mm,bottom=1.8mm,before skip=2mm,after skip=2mm]
{\bfseries\sffamily\color{teal}To'g'ri javob:} \textbf{B) atribut - faylning read-only yoki hidden kabi boshqaruv xususiyati}
\par\smallskip
{\bfseries Izoh:} atribut - faylning readonly yoki hidden kabi boshqaruv xususiyati. To'g'ri juftlik: atribut - faylning read-only yoki hidden kabi boshqaruv xususiyati. Manba: CAM6, PDF 178-205; Qo'shimcha sahifa raqami (qaysi PDFga tegishli ekani aniqlanmagan; darslik sahifasi deb tasdiqlanmagan): 86
\end{tcolorbox}
\vspace{2mm}\hrule height 0.2pt\vspace{3mm}

% DB_ID c2b28edd-64b3-67c5-e093-2d1526fc7f7a
\needspace{5\baselineskip}
{\bfseries\sffamily\color{ink} 440.} {scriptsizecolor{muted} Y1 | review | c2b28edd}
\par
Fayl va papkalar ustida amallar mavzusida "filter" qatnashgan atama-ta'rif juftliklaridan to'g'risini toping.
\begin{enumerate}[label=\Alph*),leftmargin=9mm]
\item filter - manbani saqlagan holda yangi nusxa yaratish
\item filter - shartga mos obyektlargina ko'rinadigan qilish
\item filter - obyektning joylashuvini o'zgartirish
\item filter - obyekt nomini o'zgartirish
\end{enumerate}
\begin{tcolorbox}[enhanced,breakable,colback=teal!4,colframe=teal!60!black,boxrule=0.45pt,arc=1.2mm,left=2.7mm,right=2.7mm,top=1.8mm,bottom=1.8mm,before skip=2mm,after skip=2mm]
{\bfseries\sffamily\color{teal}To'g'ri javob:} \textbf{B) filter - shartga mos obyektlargina ko'rinadigan qilish}
\par\smallskip
{\bfseries Izoh:} filter - shartga mos obyektlargina ko'rinadigan qilish. To'g'ri juftlik: filter - shartga mos obyektlargina ko'rinadigan qilish. Manba: CAM9, PDF 147-149; Qo'shimcha sahifa raqami (qaysi PDFga tegishli ekani aniqlanmagan; darslik sahifasi deb tasdiqlanmagan): 98
\end{tcolorbox}
\vspace{2mm}\hrule height 0.2pt\vspace{3mm}

% DB_ID c33ff590-e02d-7d13-588f-bfba189e5e3c
\needspace{5\baselineskip}
{\bfseries\sffamily\color{ink} 441.} {scriptsizecolor{muted} Y1 | review | c33ff590}
\par
Foydalanuvchi faqat \texttt{grafik\_} bilan boshlanuvchi va \texttt{.png} bilan tugaydigan fayllarni tanlaydi. \texttt{grafik\_*.png} niqobiga qaysi nom mos?
\begin{enumerate}[label=\Alph*),leftmargin=9mm]
\item grafik\_reja.jpg
\item grafik.png
\item reja\_grafik.png
\item grafik\_reja.png
\end{enumerate}
\begin{tcolorbox}[enhanced,breakable,colback=teal!4,colframe=teal!60!black,boxrule=0.45pt,arc=1.2mm,left=2.7mm,right=2.7mm,top=1.8mm,bottom=1.8mm,before skip=2mm,after skip=2mm]
{\bfseries\sffamily\color{teal}To'g'ri javob:} \textbf{D) grafik\_reja.png}
\par\smallskip
{\bfseries Izoh:} \texttt{*} belgi ketma-ketligiga mos, ammo qat'iy \texttt{grafik\_} prefiksi va \texttt{.png} kengaytmasi saqlanishi kerak. [Manba: Demo 2026 test 10 000 so'mlik, PDF 4 (uslub); ICT 5-sinf, PDF 44-45 (fayl kengaytmasi).]
\end{tcolorbox}
\vspace{2mm}\hrule height 0.2pt\vspace{3mm}

% DB_ID c35b4ba5-c016-73a5-7368-92970b30f08a
\needspace{5\baselineskip}
{\bfseries\sffamily\color{ink} 442.} {scriptsizecolor{muted} Y1 | review | c35b4ba5}
\par
Fayl hidden atributiga ega. Bu nimani anglatadi?
\begin{enumerate}[label=\Alph*),leftmargin=9mm]
\item mazmuni kuchli shifrlangan
\item odatdagi ko'rinishda yashirilishi mumkin
\item fayl internetga yuborilgan
\item faylning zaxirasi yaratilgan
\end{enumerate}
\begin{tcolorbox}[enhanced,breakable,colback=teal!4,colframe=teal!60!black,boxrule=0.45pt,arc=1.2mm,left=2.7mm,right=2.7mm,top=1.8mm,bottom=1.8mm,before skip=2mm,after skip=2mm]
{\bfseries\sffamily\color{teal}To'g'ri javob:} \textbf{B) odatdagi ko'rinishda yashirilishi mumkin}
\par\smallskip
{\bfseries Izoh:} odatdagi ko'rinishda yashirilishi mumkin. Hidden ko'rinish atributi, xavfsiz shifrlash emas. Manba: ICT5, PDF 42-45; Qo'shimcha sahifa raqami (qaysi PDFga tegishli ekani aniqlanmagan; darslik sahifasi deb tasdiqlanmagan): 88
\end{tcolorbox}
\vspace{2mm}\hrule height 0.2pt\vspace{3mm}

% DB_ID cee0222c-08ec-9758-5969-61792ade0291
\needspace{5\baselineskip}
{\bfseries\sffamily\color{ink} 443.} {scriptsizecolor{muted} Y1 | review | cee0222c}
\par
Papka, katalog, yo'l va ierarxiya mavzusida "quyi katalog" qatnashgan atama-ta'rif juftliklaridan to'g'risini toping.
\begin{enumerate}[label=\Alph*),leftmargin=9mm]
\item quyi katalog - fayl va quyi papkalarni guruhlovchi konteyner
\item quyi katalog - fayl tizimi daraxtining eng yuqori katalogi
\item quyi katalog - joriy obyektni bevosita o'z ichiga olgan katalog
\item quyi katalog - boshqa katalog ichida joylashgan katalog
\end{enumerate}
\begin{tcolorbox}[enhanced,breakable,colback=teal!4,colframe=teal!60!black,boxrule=0.45pt,arc=1.2mm,left=2.7mm,right=2.7mm,top=1.8mm,bottom=1.8mm,before skip=2mm,after skip=2mm]
{\bfseries\sffamily\color{teal}To'g'ri javob:} \textbf{D) quyi katalog - boshqa katalog ichida joylashgan katalog}
\par\smallskip
{\bfseries Izoh:} quyi katalog - boshqa katalog ichida joylashgan katalog. To'g'ri juftlik: quyi katalog - boshqa katalog ichida joylashgan katalog. Manba: CAM9, PDF 147-149; Qo'shimcha sahifa raqami (qaysi PDFga tegishli ekani aniqlanmagan; darslik sahifasi deb tasdiqlanmagan): 92
\end{tcolorbox}
\vspace{2mm}\hrule height 0.2pt\vspace{3mm}

% DB_ID d3860fa7-7601-58b2-1af8-f22f0cfc1b09
\needspace{5\baselineskip}
{\bfseries\sffamily\color{ink} 444.} {scriptsizecolor{muted} Y1 | review | d3860fa7}
\par
Fayl, kengaytma, format va atribut bo'yicha qatorlar ichidan "fayl" noto'g'ri izohlanganini aniqlang.
\begin{enumerate}[label=\Alph*),leftmargin=9mm]
\item format - fayl ichidagi ma'lumot tuzilishi va kodlash qoidasi
\item atribut - faylning read-only yoki hidden kabi boshqaruv xususiyati
\item metama'lumot - fayl mazmuni haqida tavsiflovchi ma'lumot
\item fayl - fayl nomining formatni ko'rsatuvchi oxirgi qismi
\end{enumerate}
\begin{tcolorbox}[enhanced,breakable,colback=teal!4,colframe=teal!60!black,boxrule=0.45pt,arc=1.2mm,left=2.7mm,right=2.7mm,top=1.8mm,bottom=1.8mm,before skip=2mm,after skip=2mm]
{\bfseries\sffamily\color{teal}To'g'ri javob:} \textbf{D) fayl - fayl nomining formatni ko'rsatuvchi oxirgi qismi}
\par\smallskip
{\bfseries Izoh:} fayl - fayl nomining formatni ko'rsatuvchi oxirgi qismi. fayl uchun to'g'ri ta'rif: nom bilan saqlanadigan ma'lumotlar birligi. Manba: ICT5, PDF 42-45; Qo'shimcha sahifa raqami (qaysi PDFga tegishli ekani aniqlanmagan; darslik sahifasi deb tasdiqlanmagan): 83
\end{tcolorbox}
\vspace{2mm}\hrule height 0.2pt\vspace{3mm}

% DB_ID f2a535cd-0179-eec3-334d-8aa1747eb2cb
\needspace{5\baselineskip}
{\bfseries\sffamily\color{ink} 445.} {scriptsizecolor{muted} Y1 | review | f2a535cd}
\par
Papka, katalog, yo'l va ierarxiya bo'yicha qatorlar ichidan "mutlaq yo'l" noto'g'ri izohlanganini aniqlang.
\begin{enumerate}[label=\Alph*),leftmargin=9mm]
\item ildiz katalog - fayl tizimi daraxtining eng yuqori katalogi
\item mutlaq yo'l - fayl va quyi papkalarni guruhlovchi konteyner
\item ona katalog - joriy obyektni bevosita o'z ichiga olgan katalog
\item quyi katalog - boshqa katalog ichida joylashgan katalog
\end{enumerate}
\begin{tcolorbox}[enhanced,breakable,colback=teal!4,colframe=teal!60!black,boxrule=0.45pt,arc=1.2mm,left=2.7mm,right=2.7mm,top=1.8mm,bottom=1.8mm,before skip=2mm,after skip=2mm]
{\bfseries\sffamily\color{teal}To'g'ri javob:} \textbf{B) mutlaq yo'l - fayl va quyi papkalarni guruhlovchi konteyner}
\par\smallskip
{\bfseries Izoh:} mutlaq yo'l - fayl va quyi papkalarni guruhlovchi konteyner. mutlaq yo'l uchun to'g'ri ta'rif: ildizdan boshlanadigan to'liq joylashuv. Manba: ICT6, PDF 116-124; Qo'shimcha sahifa raqami (qaysi PDFga tegishli ekani aniqlanmagan; darslik sahifasi deb tasdiqlanmagan): 93
\end{tcolorbox}
\vspace{2mm}\hrule height 0.2pt\vspace{3mm}

% DB_ID 00d36d7a-d681-2f61-3e0a-ae7a37e0a5e4
\needspace{5\baselineskip}
{\bfseries\sffamily\color{ink} 446.} {scriptsizecolor{muted} Y1 | review | 00d36d7a}
\par
Fayl va papkalar ustida amallar bo'yicha qatorlar ichidan "filter" noto'g'ri izohlanganini aniqlang.
\begin{enumerate}[label=\Alph*),leftmargin=9mm]
\item move - obyektning joylashuvini o'zgartirish
\item rename - obyekt nomini o'zgartirish
\item filter - manbani saqlagan holda yangi nusxa yaratish
\item archive - bir nechta faylni konteynerga yig'ish va ko'pincha siqish
\end{enumerate}
\begin{tcolorbox}[enhanced,breakable,colback=teal!4,colframe=teal!60!black,boxrule=0.45pt,arc=1.2mm,left=2.7mm,right=2.7mm,top=1.8mm,bottom=1.8mm,before skip=2mm,after skip=2mm]
{\bfseries\sffamily\color{teal}To'g'ri javob:} \textbf{C) filter - manbani saqlagan holda yangi nusxa yaratish}
\par\smallskip
{\bfseries Izoh:} filter - manbani saqlagan holda yangi nusxa yaratish. filter uchun to'g'ri ta'rif: shartga mos obyektlargina ko'rinadigan qilish. Manba: CAM9, PDF 147-149; Qo'shimcha sahifa raqami (qaysi PDFga tegishli ekani aniqlanmagan; darslik sahifasi deb tasdiqlanmagan): 98
\end{tcolorbox}
\vspace{2mm}\hrule height 0.2pt\vspace{3mm}

% DB_ID 02d76683-068a-e78e-7e0c-d7eda44f8c1a
\needspace{5\baselineskip}
{\bfseries\sffamily\color{ink} 447.} {scriptsizecolor{muted} Y1 | review | 02d76683}
\par
foto.jpg fayli foto.docx deb qayta nomlandi. Eng to'g'ri xulosa qaysi?
\begin{enumerate}[label=\Alph*),leftmargin=9mm]
\item ichki format o'zgarmadi va fayl ochilmasligi mumkin
\item rasm Word hujjatiga to'liq aylandi
\item fayl avtomatik shifrlanadi
\item rasm sifati albatta oshadi
\end{enumerate}
\begin{tcolorbox}[enhanced,breakable,colback=teal!4,colframe=teal!60!black,boxrule=0.45pt,arc=1.2mm,left=2.7mm,right=2.7mm,top=1.8mm,bottom=1.8mm,before skip=2mm,after skip=2mm]
{\bfseries\sffamily\color{teal}To'g'ri javob:} \textbf{A) ichki format o'zgarmadi va fayl ochilmasligi mumkin}
\par\smallskip
{\bfseries Izoh:} ichki format o'zgarmadi va fayl ochilmasligi mumkin. Nomdagi kengaytmani almashtirish ichki formatni konvertatsiya qilmaydi. Manba: ICT5, PDF 42-45; Qo'shimcha sahifa raqami (qaysi PDFga tegishli ekani aniqlanmagan; darslik sahifasi deb tasdiqlanmagan): 88
\end{tcolorbox}
\vspace{2mm}\hrule height 0.2pt\vspace{3mm}

% DB_ID 039ab23a-9197-4c44-8398-402e9562e357
\needspace{5\baselineskip}
{\bfseries\sffamily\color{ink} 448.} {scriptsizecolor{muted} Y1 | review | 039ab23a}
\par
<<hisobot.txt>> faylining kengaytmasi nomlash orqali <<hisobot.jpg>>ga o'zgardi. Konvertatsiya dasturi ishlatilmagan. Qaysi fikr to'g'ri?
\begin{enumerate}[label=\Alph*),leftmargin=9mm]
\item Fayl endi rastrli tasvir formatiga to'liq o'tkazilgan deb hisoblanadi
\item Faylning ichki matn baytlari nom o'zgarishi tufayli albatta yo'qoladi
\item Fayl endi ko'proq xotira talab qiladi, chunki yangi kengaytma uzunroq
\item Faylning nomi o'zgargan, lekin ichki ma'lumot formati saqlangan
\end{enumerate}
\begin{tcolorbox}[enhanced,breakable,colback=teal!4,colframe=teal!60!black,boxrule=0.45pt,arc=1.2mm,left=2.7mm,right=2.7mm,top=1.8mm,bottom=1.8mm,before skip=2mm,after skip=2mm]
{\bfseries\sffamily\color{teal}To'g'ri javob:} \textbf{D) Faylning nomi o'zgargan, lekin ichki ma'lumot formati saqlangan}
\par\smallskip
{\bfseries Izoh:} Fayl nomi va kengaytmasi uning ichki mazmunidan farq qiladi. Kengaytmani qayta nomlash fayl baytlarini yangi formatga konvertatsiya qilmaydi. [Manba: ICT 5-sinf (2020), PDF 45-bet.]
\end{tcolorbox}
\vspace{2mm}\hrule height 0.2pt\vspace{3mm}

% DB_ID 0b9175b7-adb8-41d2-cbfa-057423a35b4a
\needspace{5\baselineskip}
{\bfseries\sffamily\color{ink} 449.} {scriptsizecolor{muted} Y1 | review | 0b9175b7}
\par
Ikki papkada bir xil <<nazorat.docx>> nomli fayllar mavjud. Bitta faylni aynan ko'rsatish uchun qaysi ma'lumot yetarlicha aniq?
\begin{enumerate}[label=\Alph*),leftmargin=9mm]
\item Faylning qaysi dastur yordamida oxirgi marta ochilganini va ekranidagi piktogrammasini ko'rsatish
\item Faylning papkalar ketma-ketligi bilan birga to'liq manzilini ko'rsatish
\item Fayl nomining harflari bosh yoki kichik yozilganini va uning fayl kengaytmasini alohida ko'rsatish
\item Faylning ekran ikonkasidagi rangi, shakli hamda oxirgi tahrir etilgan dasturini yozib qo'yish
\end{enumerate}
\begin{tcolorbox}[enhanced,breakable,colback=teal!4,colframe=teal!60!black,boxrule=0.45pt,arc=1.2mm,left=2.7mm,right=2.7mm,top=1.8mm,bottom=1.8mm,before skip=2mm,after skip=2mm]
{\bfseries\sffamily\color{teal}To'g'ri javob:} \textbf{B) Faylning papkalar ketma-ketligi bilan birga to'liq manzilini ko'rsatish}
\par\smallskip
{\bfseries Izoh:} Bir xil nomdagi fayllar turli papkalarda bo'lishi mumkin. Papkalar ierarxiyasi orqali tuzilgan to'liq yo'l faylning saqlanish joyini ajratib beradi. [Manba: ICT 5-sinf (2020), PDF 43-bet.]
\end{tcolorbox}
\vspace{2mm}\hrule height 0.2pt\vspace{3mm}

% DB_ID 23b9ca58-50f9-4e92-3a90-28a185ab9b25
\needspace{5\baselineskip}
{\bfseries\sffamily\color{ink} 450.} {scriptsizecolor{muted} Y1 | review | 23b9ca58}
\par
A papkadan B papkaga fayl nusxalandi, so'ng B dagi nusxa qayta nomlandi. Yakunda nechta fayl nusxasi mavjud?
\begin{enumerate}[label=\Alph*),leftmargin=9mm]
\item ikkita, nomlari turlicha
\item bitta, eski nom bilan
\item bitta, yangi nom bilan
\item uchta, nomlari bir xil
\end{enumerate}
\begin{tcolorbox}[enhanced,breakable,colback=teal!4,colframe=teal!60!black,boxrule=0.45pt,arc=1.2mm,left=2.7mm,right=2.7mm,top=1.8mm,bottom=1.8mm,before skip=2mm,after skip=2mm]
{\bfseries\sffamily\color{teal}To'g'ri javob:} \textbf{A) ikkita, nomlari turlicha}
\par\smallskip
{\bfseries Izoh:} ikkita, nomlari turlicha. Copy ikki nusxa hosil qiladi; rename faqat B dagi nusxa nomini o'zgartiradi. Manba: ICT5, PDF 43-49; Qo'shimcha sahifa raqami (qaysi PDFga tegishli ekani aniqlanmagan; darslik sahifasi deb tasdiqlanmagan): 100
\end{tcolorbox}
\vspace{2mm}\hrule height 0.2pt\vspace{3mm}

% DB_ID 2e1bae84-407e-df7b-e5bb-d8134a2e1132
\needspace{5\baselineskip}
{\bfseries\sffamily\color{ink} 451.} {scriptsizecolor{muted} Y1 | review | 2e1bae84}
\par
<<natija.txt>> faylining nomi <<natija.jpg>>ga almashtirildi, lekin konvertatsiya dasturi ishlatilmadi. Nima o'zgardi?
\begin{enumerate}[label=\Alph*),leftmargin=9mm]
\item Fayl kengaytmasi o'zgargan, uning ichki matn ma'lumotlari o'z holicha saqlanib qolgan
\item Fayl ichidagi matn barcha grafik piksellarga avtomatik aylangan va har bir belgi yangi rastrli qatlamga joylashtirilgan
\item Barcha baytlar o'chirilib, tizim tomonidan yangi JPEG tasvir fayli sifatida to'liq qayta qurilgan
\item Faqat faylni qayta nomlash operatsiyasi bajarilgan bo'lsa ham, barcha formatlar avtomatik bir-biriga konvertatsiya qilinadi
\end{enumerate}
\begin{tcolorbox}[enhanced,breakable,colback=teal!4,colframe=teal!60!black,boxrule=0.45pt,arc=1.2mm,left=2.7mm,right=2.7mm,top=1.8mm,bottom=1.8mm,before skip=2mm,after skip=2mm]
{\bfseries\sffamily\color{teal}To'g'ri javob:} \textbf{A) Fayl kengaytmasi o'zgargan, uning ichki matn ma'lumotlari o'z holicha saqlanib qolgan}
\par\smallskip
{\bfseries Izoh:} Faylni qayta nomlash uning ichki ma'lumot tuzilishini yangi grafik formatga avtomatik o'tkazmaydi. [Manba: ICT 5-sinf PDF 44-45.]
\end{tcolorbox}
\vspace{2mm}\hrule height 0.2pt\vspace{3mm}

% DB_ID 3e6bf43c-ccf2-b2af-c5df-321d33618681
\needspace{5\baselineskip}
{\bfseries\sffamily\color{ink} 452.} {scriptsizecolor{muted} Y1 | review | 3e6bf43c}
\par
Papka, katalog, yo'l va ierarxiya mavzusida: I. ona katalog - joriy obyektni bevosita o'z ichiga olgan katalog. II. papka - joriy obyektni bevosita o'z ichiga olgan katalog. Mulohazalarni baholang.
\begin{enumerate}[label=\Alph*),leftmargin=9mm]
\item faqat I
\item faqat II
\item I va II
\item I ham, II ham emas
\end{enumerate}
\begin{tcolorbox}[enhanced,breakable,colback=teal!4,colframe=teal!60!black,boxrule=0.45pt,arc=1.2mm,left=2.7mm,right=2.7mm,top=1.8mm,bottom=1.8mm,before skip=2mm,after skip=2mm]
{\bfseries\sffamily\color{teal}To'g'ri javob:} \textbf{A) faqat I}
\par\smallskip
{\bfseries Izoh:} faqat I. Birinchi va ikkinchi mulohazaning ta'riflari alohida tekshirildi; natija: faqat I. Manba: CAM6, PDF 178-205; Qo'shimcha sahifa raqami (qaysi PDFga tegishli ekani aniqlanmagan; darslik sahifasi deb tasdiqlanmagan): 92
\end{tcolorbox}
\vspace{2mm}\hrule height 0.2pt\vspace{3mm}

% DB_ID 520e0da1-86e7-e131-9981-dd3258efa46a
\needspace{5\baselineskip}
{\bfseries\sffamily\color{ink} 453.} {scriptsizecolor{muted} Y1 | review | 520e0da1}
\par
Fayl, kengaytma, format va atribut bo'yicha qatorlar ichidan "atribut" noto'g'ri izohlanganini aniqlang.
\begin{enumerate}[label=\Alph*),leftmargin=9mm]
\item kengaytma - fayl nomining formatni ko'rsatuvchi oxirgi qismi
\item format - fayl ichidagi ma'lumot tuzilishi va kodlash qoidasi
\item atribut - nom bilan saqlanadigan ma'lumotlar birligi
\item metama'lumot - fayl mazmuni haqida tavsiflovchi ma'lumot
\end{enumerate}
\begin{tcolorbox}[enhanced,breakable,colback=teal!4,colframe=teal!60!black,boxrule=0.45pt,arc=1.2mm,left=2.7mm,right=2.7mm,top=1.8mm,bottom=1.8mm,before skip=2mm,after skip=2mm]
{\bfseries\sffamily\color{teal}To'g'ri javob:} \textbf{C) atribut - nom bilan saqlanadigan ma'lumotlar birligi}
\par\smallskip
{\bfseries Izoh:} atribut - nom bilan saqlanadigan ma'lumotlar birligi. atribut uchun to'g'ri ta'rif: faylning read-only yoki hidden kabi boshqaruv xususiyati. Manba: CAM6, PDF 178-205; Qo'shimcha sahifa raqami (qaysi PDFga tegishli ekani aniqlanmagan; darslik sahifasi deb tasdiqlanmagan): 86
\end{tcolorbox}
\vspace{2mm}\hrule height 0.2pt\vspace{3mm}

% DB_ID 544b4ff7-ad2a-c1d3-8607-42103c4fdb96
\needspace{5\baselineskip}
{\bfseries\sffamily\color{ink} 454.} {scriptsizecolor{muted} Y1 | review | 544b4ff7}
\par
Fayl va papkalar ustida amallar mavzusida: I. rename - obyekt nomini o'zgartirish. II. copy - obyekt nomini o'zgartirish. Mulohazalarni baholang.
\begin{enumerate}[label=\Alph*),leftmargin=9mm]
\item faqat II
\item I va II
\item faqat I
\item I ham, II ham emas
\end{enumerate}
\begin{tcolorbox}[enhanced,breakable,colback=teal!4,colframe=teal!60!black,boxrule=0.45pt,arc=1.2mm,left=2.7mm,right=2.7mm,top=1.8mm,bottom=1.8mm,before skip=2mm,after skip=2mm]
{\bfseries\sffamily\color{teal}To'g'ri javob:} \textbf{C) faqat I}
\par\smallskip
{\bfseries Izoh:} faqat I. Birinchi va ikkinchi mulohazaning ta'riflari alohida tekshirildi; natija: faqat I. Manba: CAM6, PDF 196-205; Qo'shimcha sahifa raqami (qaysi PDFga tegishli ekani aniqlanmagan; darslik sahifasi deb tasdiqlanmagan): 97
\end{tcolorbox}
\vspace{2mm}\hrule height 0.2pt\vspace{3mm}

% DB_ID 54ccaedb-be87-2565-aa8b-5f322c63d86c
\needspace{5\baselineskip}
{\bfseries\sffamily\color{ink} 455.} {scriptsizecolor{muted} Y1 | review | 54ccaedb}
\par
Fayl, kengaytma, format va atribut mavzusida: I. metama'lumot - fayl mazmuni haqida tavsiflovchi ma'lumot. II. fayl - nom bilan saqlanadigan ma'lumotlar birligi. Mulohazalarni baholang.
\begin{enumerate}[label=\Alph*),leftmargin=9mm]
\item I va II
\item faqat I
\item faqat II
\item I ham, II ham emas
\end{enumerate}
\begin{tcolorbox}[enhanced,breakable,colback=teal!4,colframe=teal!60!black,boxrule=0.45pt,arc=1.2mm,left=2.7mm,right=2.7mm,top=1.8mm,bottom=1.8mm,before skip=2mm,after skip=2mm]
{\bfseries\sffamily\color{teal}To'g'ri javob:} \textbf{A) I va II}
\par\smallskip
{\bfseries Izoh:} I va II. Birinchi va ikkinchi mulohazaning ta'riflari alohida tekshirildi; natija: I va II. Manba: CAM7, PDF 41-54; Qo'shimcha sahifa raqami (qaysi PDFga tegishli ekani aniqlanmagan; darslik sahifasi deb tasdiqlanmagan): 88
\end{tcolorbox}
\vspace{2mm}\hrule height 0.2pt\vspace{3mm}

% DB_ID 557cfc3a-a07e-2a10-8444-41a26946c682
\needspace{5\baselineskip}
{\bfseries\sffamily\color{ink} 456.} {scriptsizecolor{muted} Y1 | review | 557cfc3a}
\par
Fayl va papkalar ustida amallar mavzusida: I. archive - bir nechta faylni konteynerga yig'ish va ko'pincha siqish. II. copy - manbani saqlagan holda yangi nusxa yaratish. Mulohazalarni baholang.
\begin{enumerate}[label=\Alph*),leftmargin=9mm]
\item I va II
\item faqat I
\item faqat II
\item I ham, II ham emas
\end{enumerate}
\begin{tcolorbox}[enhanced,breakable,colback=teal!4,colframe=teal!60!black,boxrule=0.45pt,arc=1.2mm,left=2.7mm,right=2.7mm,top=1.8mm,bottom=1.8mm,before skip=2mm,after skip=2mm]
{\bfseries\sffamily\color{teal}To'g'ri javob:} \textbf{A) I va II}
\par\smallskip
{\bfseries Izoh:} I va II. Birinchi va ikkinchi mulohazaning ta'riflari alohida tekshirildi; natija: I va II. Manba: ICT6, PDF 116-126; Qo'shimcha sahifa raqami (qaysi PDFga tegishli ekani aniqlanmagan; darslik sahifasi deb tasdiqlanmagan): 99
\end{tcolorbox}
\vspace{2mm}\hrule height 0.2pt\vspace{3mm}

% DB_ID 63a3cb9d-916d-2af2-0056-185c4c92b26d
\needspace{5\baselineskip}
{\bfseries\sffamily\color{ink} 457.} {scriptsizecolor{muted} Y1 | review | 63a3cb9d}
\par
Kompyuterda NTFS atamasi va fayllarni ko'rish oynasi uchradi. Ularni bir-biridan qanday ajratish mumkin?
\begin{enumerate}[label=\Alph*),leftmargin=9mm]
\item NTFS fayl tizimi, oynadagi fayl menejeri esa foydalanuvchi interfeysi
\item NTFS matn muharriri, fayl menejeri esa tasvirni qayta ishlash protsessori
\item NTFS qurilma drayveri, fayl menejeri esa ma'lumotlar bazasi dvigateli
\item NTFS operatsion tizim oilasi, fayl menejeri esa magnit disk mexanizmi
\end{enumerate}
\begin{tcolorbox}[enhanced,breakable,colback=teal!4,colframe=teal!60!black,boxrule=0.45pt,arc=1.2mm,left=2.7mm,right=2.7mm,top=1.8mm,bottom=1.8mm,before skip=2mm,after skip=2mm]
{\bfseries\sffamily\color{teal}To'g'ri javob:} \textbf{A) NTFS fayl tizimi, oynadagi fayl menejeri esa foydalanuvchi interfeysi}
\par\smallskip
{\bfseries Izoh:} ICT 10-sinfda NTFS fayl tizimi sifatida keltiriladi. Fayl menejeri foydalanuvchining papka va fayllar bilan ishlash interfeysidir. [Manba: ICT 10-sinf (2021), PDF 26-bet.]
\end{tcolorbox}
\vspace{2mm}\hrule height 0.2pt\vspace{3mm}

% DB_ID 6be5b25b-cfbc-7c17-59ee-c9e6b62dc4c7
\needspace{5\baselineskip}
{\bfseries\sffamily\color{ink} 458.} {scriptsizecolor{muted} Y1 | review | 6be5b25b}
\par
Papka, katalog, yo'l va ierarxiya mavzusida: I. mutlaq yo'l - ildizdan boshlanadigan to'liq joylashuv. II. papka - fayl va quyi papkalarni guruhlovchi konteyner. Mulohazalarni baholang.
\begin{enumerate}[label=\Alph*),leftmargin=9mm]
\item faqat I
\item faqat II
\item I va II
\item I ham, II ham emas
\end{enumerate}
\begin{tcolorbox}[enhanced,breakable,colback=teal!4,colframe=teal!60!black,boxrule=0.45pt,arc=1.2mm,left=2.7mm,right=2.7mm,top=1.8mm,bottom=1.8mm,before skip=2mm,after skip=2mm]
{\bfseries\sffamily\color{teal}To'g'ri javob:} \textbf{C) I va II}
\par\smallskip
{\bfseries Izoh:} I va II. Birinchi va ikkinchi mulohazaning ta'riflari alohida tekshirildi; natija: I va II. Manba: ICT6, PDF 116-124; Qo'shimcha sahifa raqami (qaysi PDFga tegishli ekani aniqlanmagan; darslik sahifasi deb tasdiqlanmagan): 94
\end{tcolorbox}
\vspace{2mm}\hrule height 0.2pt\vspace{3mm}

% DB_ID 8d26a33f-b239-db0f-3bd4-34b613fba2a4
\needspace{5\baselineskip}
{\bfseries\sffamily\color{ink} 459.} {scriptsizecolor{muted} Y1 | review | 8d26a33f}
\par
Papka, katalog, yo'l va ierarxiya mavzusida: I. papka - fayl va quyi papkalarni guruhlovchi konteyner. II. ildiz katalog - fayl tizimi daraxtining eng yuqori katalogi. Mulohazalarni baholang.
\begin{enumerate}[label=\Alph*),leftmargin=9mm]
\item faqat I
\item faqat II
\item I va II
\item I ham, II ham emas
\end{enumerate}
\begin{tcolorbox}[enhanced,breakable,colback=teal!4,colframe=teal!60!black,boxrule=0.45pt,arc=1.2mm,left=2.7mm,right=2.7mm,top=1.8mm,bottom=1.8mm,before skip=2mm,after skip=2mm]
{\bfseries\sffamily\color{teal}To'g'ri javob:} \textbf{C) I va II}
\par\smallskip
{\bfseries Izoh:} I va II. Birinchi va ikkinchi mulohazaning ta'riflari alohida tekshirildi; natija: I va II. Manba: ICT5, PDF 42-45; Qo'shimcha sahifa raqami (qaysi PDFga tegishli ekani aniqlanmagan; darslik sahifasi deb tasdiqlanmagan): 90
\end{tcolorbox}
\vspace{2mm}\hrule height 0.2pt\vspace{3mm}

% DB_ID 8f4af1fc-66c3-76f9-7c2c-317a3fe65132
\needspace{5\baselineskip}
{\bfseries\sffamily\color{ink} 460.} {scriptsizecolor{muted} Y1 | review | 8f4af1fc}
\par
<<Reja.docx>> faylining asl nusxasi ishchi papkada qolib, arxiv papkada ikkinchi nusxasi paydo bo'lishi kerak. Qaysi amal mos?
\begin{enumerate}[label=\Alph*),leftmargin=9mm]
\item Nusxalash asl faylni saqlagan holda boshqa papkada ikkinchi nusxani yaratadi
\item Ko'chirish asosiy faylni manba papkadan chiqarib, yangi joyga olib o'tadi; shunga qaramasdan, har ikkala papkada ham mustaqil nusxa qoladi
\item Qayta nomlash papka yo'lini o'zgartirmasdan ikki mustaqil jismoniy nusxa yaratib, ularni ikki xil diskka taqsimlaydi
\item O'chirish faylni boshqa papkaga avtomatik ko'chirib, uning asl nusxasini ham to'liq saqlab qoladi
\end{enumerate}
\begin{tcolorbox}[enhanced,breakable,colback=teal!4,colframe=teal!60!black,boxrule=0.45pt,arc=1.2mm,left=2.7mm,right=2.7mm,top=1.8mm,bottom=1.8mm,before skip=2mm,after skip=2mm]
{\bfseries\sffamily\color{teal}To'g'ri javob:} \textbf{A) Nusxalash asl faylni saqlagan holda boshqa papkada ikkinchi nusxani yaratadi}
\par\smallskip
{\bfseries Izoh:} Nusxalash asl faylni saqlab, ikkinchi nusxani boshqa papkada hosil qiladi. [Manba: ICT 5-sinf PDF 45.]
\end{tcolorbox}
\vspace{2mm}\hrule height 0.2pt\vspace{3mm}

% DB_ID d36bdde6-2a17-e101-7a79-79df03508160
\needspace{5\baselineskip}
{\bfseries\sffamily\color{ink} 461.} {scriptsizecolor{muted} Y1 | review | d36bdde6}
\par
<<Darslar>> ichida <<Informatika>>, uning ichida esa <<Testlar>> papkasi joylashgan. <<Testlar>>ning joylashuvini aniq qanday yozish mumkin?
\begin{enumerate}[label=\Alph*),leftmargin=9mm]
\item Papkalarning faqat alifbo bo'yicha tartiblangan umumiy nomlar ro'yxatini berish
\item Papka nomlarining ekrandagi belgilari, yasalgan sana va dastur piktogrammasining rangini birlashtirib tavsiflash
\item Darslar/Informatika/Testlar ichma-ich papkalar ketma-ketligini ko'rsatish
\item Saqlash qurilmasidagi barcha papkalar nomini yozish va har birining ichki fayllarini ham ko'rsatish
\end{enumerate}
\begin{tcolorbox}[enhanced,breakable,colback=teal!4,colframe=teal!60!black,boxrule=0.45pt,arc=1.2mm,left=2.7mm,right=2.7mm,top=1.8mm,bottom=1.8mm,before skip=2mm,after skip=2mm]
{\bfseries\sffamily\color{teal}To'g'ri javob:} \textbf{C) Darslar/Informatika/Testlar ichma-ich papkalar ketma-ketligini ko'rsatish}
\par\smallskip
{\bfseries Izoh:} Papka daraxtidagi ichma-ich yo'l ketma-ket joylashuvni aniqlaydi. [Manba: ICT 5-sinf PDF 43-45.]
\end{tcolorbox}
\vspace{2mm}\hrule height 0.2pt\vspace{3mm}

% DB_ID d8da987b-9f4c-547a-2866-61a7d6980f74
\needspace{5\baselineskip}
{\bfseries\sffamily\color{ink} 462.} {scriptsizecolor{muted} Y1 | review | d8da987b}
\par
Papka, katalog, yo'l va ierarxiya bo'yicha qatorlar ichidan "quyi katalog" noto'g'ri izohlanganini aniqlang.
\begin{enumerate}[label=\Alph*),leftmargin=9mm]
\item quyi katalog - fayl va quyi papkalarni guruhlovchi konteyner
\item ildiz katalog - fayl tizimi daraxtining eng yuqori katalogi
\item ona katalog - joriy obyektni bevosita o'z ichiga olgan katalog
\item mutlaq yo'l - ildizdan boshlanadigan to'liq joylashuv
\end{enumerate}
\begin{tcolorbox}[enhanced,breakable,colback=teal!4,colframe=teal!60!black,boxrule=0.45pt,arc=1.2mm,left=2.7mm,right=2.7mm,top=1.8mm,bottom=1.8mm,before skip=2mm,after skip=2mm]
{\bfseries\sffamily\color{teal}To'g'ri javob:} \textbf{A) quyi katalog - fayl va quyi papkalarni guruhlovchi konteyner}
\par\smallskip
{\bfseries Izoh:} quyi katalog - fayl va quyi papkalarni guruhlovchi konteyner. quyi katalog uchun to'g'ri ta'rif: boshqa katalog ichida joylashgan katalog. Manba: CAM9, PDF 147-149; Qo'shimcha sahifa raqami (qaysi PDFga tegishli ekani aniqlanmagan; darslik sahifasi deb tasdiqlanmagan): 92
\end{tcolorbox}
\vspace{2mm}\hrule height 0.2pt\vspace{3mm}

% DB_ID e433496c-e50e-8ccb-5982-54630ce3d6cf
\needspace{5\baselineskip}
{\bfseries\sffamily\color{ink} 463.} {scriptsizecolor{muted} Y1 | review | e433496c}
\par
Fayl va papkalar ustida amallar mavzusida: I. copy - manbani saqlagan holda yangi nusxa yaratish. II. move - obyektning joylashuvini o'zgartirish. Mulohazalarni baholang.
\begin{enumerate}[label=\Alph*),leftmargin=9mm]
\item I va II
\item faqat I
\item faqat II
\item I ham, II ham emas
\end{enumerate}
\begin{tcolorbox}[enhanced,breakable,colback=teal!4,colframe=teal!60!black,boxrule=0.45pt,arc=1.2mm,left=2.7mm,right=2.7mm,top=1.8mm,bottom=1.8mm,before skip=2mm,after skip=2mm]
{\bfseries\sffamily\color{teal}To'g'ri javob:} \textbf{A) I va II}
\par\smallskip
{\bfseries Izoh:} I va II. Birinchi va ikkinchi mulohazaning ta'riflari alohida tekshirildi; natija: I va II. Manba: ICT5, PDF 43-49; Qo'shimcha sahifa raqami (qaysi PDFga tegishli ekani aniqlanmagan; darslik sahifasi deb tasdiqlanmagan): 96
\end{tcolorbox}
\vspace{2mm}\hrule height 0.2pt\vspace{3mm}

% DB_ID f131740f-e496-7626-af74-078b6e0c9157
\needspace{5\baselineskip}
{\bfseries\sffamily\color{ink} 464.} {scriptsizecolor{muted} Y1 | review | f131740f}
\par
Fayl, kengaytma, format va atribut mavzusida: I. fayl - nom bilan saqlanadigan ma'lumotlar birligi. II. kengaytma - fayl nomining formatni ko'rsatuvchi oxirgi qismi. Mulohazalarni baholang.
\begin{enumerate}[label=\Alph*),leftmargin=9mm]
\item I va II
\item faqat I
\item faqat II
\item I ham, II ham emas
\end{enumerate}
\begin{tcolorbox}[enhanced,breakable,colback=teal!4,colframe=teal!60!black,boxrule=0.45pt,arc=1.2mm,left=2.7mm,right=2.7mm,top=1.8mm,bottom=1.8mm,before skip=2mm,after skip=2mm]
{\bfseries\sffamily\color{teal}To'g'ri javob:} \textbf{A) I va II}
\par\smallskip
{\bfseries Izoh:} I va II. Birinchi va ikkinchi mulohazaning ta'riflari alohida tekshirildi; natija: I va II. Manba: ICT5, PDF 42-45; Qo'shimcha sahifa raqami (qaysi PDFga tegishli ekani aniqlanmagan; darslik sahifasi deb tasdiqlanmagan): 84
\end{tcolorbox}
\vspace{2mm}\hrule height 0.2pt\vspace{3mm}

% DB_ID fbbe5377-355d-17a3-dbaf-68907f6dd208
\needspace{5\baselineskip}
{\bfseries\sffamily\color{ink} 465.} {scriptsizecolor{muted} Y1 | review | fbbe5377}
\par
Fayl, kengaytma, format va atribut mavzusida: I. format - fayl ichidagi ma'lumot tuzilishi va kodlash qoidasi. II. fayl - fayl ichidagi ma'lumot tuzilishi va kodlash qoidasi. Mulohazalarni baholang.
\begin{enumerate}[label=\Alph*),leftmargin=9mm]
\item faqat II
\item I va II
\item faqat I
\item I ham, II ham emas
\end{enumerate}
\begin{tcolorbox}[enhanced,breakable,colback=teal!4,colframe=teal!60!black,boxrule=0.45pt,arc=1.2mm,left=2.7mm,right=2.7mm,top=1.8mm,bottom=1.8mm,before skip=2mm,after skip=2mm]
{\bfseries\sffamily\color{teal}To'g'ri javob:} \textbf{C) faqat I}
\par\smallskip
{\bfseries Izoh:} faqat I. Birinchi va ikkinchi mulohazaning ta'riflari alohida tekshirildi; natija: faqat I. Manba: CAM9, PDF 147-149; Qo'shimcha sahifa raqami (qaysi PDFga tegishli ekani aniqlanmagan; darslik sahifasi deb tasdiqlanmagan): 86
\end{tcolorbox}
\vspace{2mm}\hrule height 0.2pt\vspace{3mm}

% DB_ID 1786fc81-2830-fbbc-d895-81aaa94c6641
\needspace{5\baselineskip}
{\bfseries\sffamily\color{ink} 466.} {scriptsizecolor{muted} Y1 | review | 1786fc81}
\par
Fayl, kengaytma, format va atribut mavzusida: I. atribut - fayl nomining formatni ko'rsatuvchi oxirgi qismi. II. fayl - faylning read-only yoki hidden kabi boshqaruv xususiyati. Mulohazalarni baholang.
\begin{enumerate}[label=\Alph*),leftmargin=9mm]
\item faqat I
\item faqat II
\item I va II
\item I ham, II ham emas
\end{enumerate}
\begin{tcolorbox}[enhanced,breakable,colback=teal!4,colframe=teal!60!black,boxrule=0.45pt,arc=1.2mm,left=2.7mm,right=2.7mm,top=1.8mm,bottom=1.8mm,before skip=2mm,after skip=2mm]
{\bfseries\sffamily\color{teal}To'g'ri javob:} \textbf{D) I ham, II ham emas}
\par\smallskip
{\bfseries Izoh:} I ham, II ham emas. Birinchi va ikkinchi mulohazaning ta'riflari alohida tekshirildi; natija: I ham, II ham emas. Manba: CAM6, PDF 178-205; Qo'shimcha sahifa raqami (qaysi PDFga tegishli ekani aniqlanmagan; darslik sahifasi deb tasdiqlanmagan): 87
\end{tcolorbox}
\vspace{2mm}\hrule height 0.2pt\vspace{3mm}

% DB_ID 1deca5c0-5b9c-e024-1d30-86bac027c6dc
\needspace{5\baselineskip}
{\bfseries\sffamily\color{ink} 467.} {scriptsizecolor{muted} Y1 | review | 1deca5c0}
\par
Fayl va papkalar ustida amallar mavzusida: I. filter - obyektning joylashuvini o'zgartirish. II. copy - shartga mos obyektlargina ko'rinadigan qilish. Mulohazalarni baholang.
\begin{enumerate}[label=\Alph*),leftmargin=9mm]
\item faqat I
\item faqat II
\item I va II
\item I ham, II ham emas
\end{enumerate}
\begin{tcolorbox}[enhanced,breakable,colback=teal!4,colframe=teal!60!black,boxrule=0.45pt,arc=1.2mm,left=2.7mm,right=2.7mm,top=1.8mm,bottom=1.8mm,before skip=2mm,after skip=2mm]
{\bfseries\sffamily\color{teal}To'g'ri javob:} \textbf{D) I ham, II ham emas}
\par\smallskip
{\bfseries Izoh:} I ham, II ham emas. Birinchi va ikkinchi mulohazaning ta'riflari alohida tekshirildi; natija: I ham, II ham emas. Manba: CAM9, PDF 147-149; Qo'shimcha sahifa raqami (qaysi PDFga tegishli ekani aniqlanmagan; darslik sahifasi deb tasdiqlanmagan): 98
\end{tcolorbox}
\vspace{2mm}\hrule height 0.2pt\vspace{3mm}

% DB_ID 3d1b5b9c-48cc-e671-87e3-e2dc8eab1d7b
\needspace{5\baselineskip}
{\bfseries\sffamily\color{ink} 468.} {scriptsizecolor{muted} Y1 | review | 3d1b5b9c}
\par
Fayl va papkalar ustida amallar mavzusida: I. move - obyekt nomini o'zgartirish. II. copy - manbani saqlagan holda yangi nusxa yaratish. Mulohazalarni baholang.
\begin{enumerate}[label=\Alph*),leftmargin=9mm]
\item faqat I
\item faqat II
\item I va II
\item I ham, II ham emas
\end{enumerate}
\begin{tcolorbox}[enhanced,breakable,colback=teal!4,colframe=teal!60!black,boxrule=0.45pt,arc=1.2mm,left=2.7mm,right=2.7mm,top=1.8mm,bottom=1.8mm,before skip=2mm,after skip=2mm]
{\bfseries\sffamily\color{teal}To'g'ri javob:} \textbf{B) faqat II}
\par\smallskip
{\bfseries Izoh:} faqat II. Birinchi va ikkinchi mulohazaning ta'riflari alohida tekshirildi; natija: faqat II. Manba: MAR, PDF 17-20; Qo'shimcha sahifa raqami (qaysi PDFga tegishli ekani aniqlanmagan; darslik sahifasi deb tasdiqlanmagan): 96
\end{tcolorbox}
\vspace{2mm}\hrule height 0.2pt\vspace{3mm}

% DB_ID 6ec2c040-efa7-a4f3-084f-aa22f894656e
\needspace{5\baselineskip}
{\bfseries\sffamily\color{ink} 469.} {scriptsizecolor{muted} Y1 | review | 6ec2c040}
\par
Fayl bulut bilan sinxron papkadan o'chirildi va o'chirish barcha qurilmalarga tarqaldi. Qaysi himoya yetishmagan?
\begin{enumerate}[label=\Alph*),leftmargin=9mm]
\item yana bir shortcut
\item mustaqil versiyali zaxira nusxa
\item fayl kengaytmasini yashirish
\item papkani nom bo'yicha saralash
\end{enumerate}
\begin{tcolorbox}[enhanced,breakable,colback=teal!4,colframe=teal!60!black,boxrule=0.45pt,arc=1.2mm,left=2.7mm,right=2.7mm,top=1.8mm,bottom=1.8mm,before skip=2mm,after skip=2mm]
{\bfseries\sffamily\color{teal}To'g'ri javob:} \textbf{B) mustaqil versiyali zaxira nusxa}
\par\smallskip
{\bfseries Izoh:} mustaqil versiyali zaxira nusxa. Sinxronlash o'chirishni tarqatadi; alohida backup tiklash imkonini beradi. Manba: ICT5, PDF 43-49; Qo'shimcha sahifa raqami (qaysi PDFga tegishli ekani aniqlanmagan; darslik sahifasi deb tasdiqlanmagan): 100
\end{tcolorbox}
\vspace{2mm}\hrule height 0.2pt\vspace{3mm}

% DB_ID 720edfa7-5786-1378-ff65-b248389ab24a
\needspace{5\baselineskip}
{\bfseries\sffamily\color{ink} 470.} {scriptsizecolor{muted} Y1 | review | 720edfa7}
\par
Kompyuterda "Darslar" papkasi ichida "Informatika", uning ichida esa "Testlar" papkasi joylashgan. Bu tuzilmaning to'g'ri tavsifi qaysi?
\begin{enumerate}[label=\Alph*),leftmargin=9mm]
\item Har bir quyi papka alohida fizik disk bo'lib, yuqori papkaga bog'lanmaydi
\item Papka ichida quyi papkalar bo'lishi va kataloglar daraxt ko'rinishida tuzilishi mumkin
\item Ichki papka yaratilganda yuqori papkada turgan mavjud fayllar avtomatik tarzda shu ichki papkaga ko'chadi
\item Fayl kengaytmalari papkalarning qaysi qurilmaga tegishli ekanini belgilaydi
\end{enumerate}
\begin{tcolorbox}[enhanced,breakable,colback=teal!4,colframe=teal!60!black,boxrule=0.45pt,arc=1.2mm,left=2.7mm,right=2.7mm,top=1.8mm,bottom=1.8mm,before skip=2mm,after skip=2mm]
{\bfseries\sffamily\color{teal}To'g'ri javob:} \textbf{B) Papka ichida quyi papkalar bo'lishi va kataloglar daraxt ko'rinishida tuzilishi mumkin}
\par\smallskip
{\bfseries Izoh:} Fayllar tartib uchun papkalar bo'yicha guruhlanadi; papka ichida boshqa papka joylashishi mumkin. Bu ierarxik katalog tuzilmasidir. Manba: ICT 5-sinf, PDF 43-44
\end{tcolorbox}
\vspace{2mm}\hrule height 0.2pt\vspace{3mm}

% DB_ID 78da937b-2ded-adab-1bc0-be956db30ac2
\needspace{5\baselineskip}
{\bfseries\sffamily\color{ink} 471.} {scriptsizecolor{muted} Y1 | review | 78da937b}
\par
Fayl, kengaytma, format va atribut mavzusida: I. kengaytma - fayl ichidagi ma'lumot tuzilishi va kodlash qoidasi. II. fayl - nom bilan saqlanadigan ma'lumotlar birligi. Mulohazalarni baholang.
\begin{enumerate}[label=\Alph*),leftmargin=9mm]
\item faqat I
\item faqat II
\item I va II
\item I ham, II ham emas
\end{enumerate}
\begin{tcolorbox}[enhanced,breakable,colback=teal!4,colframe=teal!60!black,boxrule=0.45pt,arc=1.2mm,left=2.7mm,right=2.7mm,top=1.8mm,bottom=1.8mm,before skip=2mm,after skip=2mm]
{\bfseries\sffamily\color{teal}To'g'ri javob:} \textbf{B) faqat II}
\par\smallskip
{\bfseries Izoh:} faqat II. Birinchi va ikkinchi mulohazaning ta'riflari alohida tekshirildi; natija: faqat II. Manba: MAR, PDF 16-20; Qo'shimcha sahifa raqami (qaysi PDFga tegishli ekani aniqlanmagan; darslik sahifasi deb tasdiqlanmagan): 85
\end{tcolorbox}
\vspace{2mm}\hrule height 0.2pt\vspace{3mm}

% DB_ID e775e662-509a-f7a1-5743-997e5f831b07
\needspace{5\baselineskip}
{\bfseries\sffamily\color{ink} 472.} {scriptsizecolor{muted} Y1 | review | e775e662}
\par
Papka, katalog, yo'l va ierarxiya mavzusida: I. ildiz katalog - joriy obyektni bevosita o'z ichiga olgan katalog. II. papka - fayl va quyi papkalarni guruhlovchi konteyner. Mulohazalarni baholang.
\begin{enumerate}[label=\Alph*),leftmargin=9mm]
\item faqat I
\item I va II
\item I ham, II ham emas
\item faqat II
\end{enumerate}
\begin{tcolorbox}[enhanced,breakable,colback=teal!4,colframe=teal!60!black,boxrule=0.45pt,arc=1.2mm,left=2.7mm,right=2.7mm,top=1.8mm,bottom=1.8mm,before skip=2mm,after skip=2mm]
{\bfseries\sffamily\color{teal}To'g'ri javob:} \textbf{D) faqat II}
\par\smallskip
{\bfseries Izoh:} faqat II. Birinchi va ikkinchi mulohazaning ta'riflari alohida tekshirildi; natija: faqat II. Manba: MAR, PDF 16-20; Qo'shimcha sahifa raqami (qaysi PDFga tegishli ekani aniqlanmagan; darslik sahifasi deb tasdiqlanmagan): 91
\end{tcolorbox}
\vspace{2mm}\hrule height 0.2pt\vspace{3mm}
\clearpage
\section{Tizimli, amaliy va xizmat ko'rsatuvchi dasturlar}
\subsection{Nashr qilingan testlar}

% DB_ID 182fa3d2-d1fe-4e17-a535-d4842f6f1a5a
\needspace{5\baselineskip}
{\bfseries\sffamily\color{ink} 473.} {scriptsizecolor{muted} Y1 | published | 182fa3d2}
\par
Printer yangi kompyuterga ulandi, ammo operatsion tizim uni boshqarish uchun mos dasturiy interfeysga ega emas. Bunday vaziyatda drayverning vazifasi qaysi?
\begin{enumerate}[label=\Alph*),leftmargin=9mm]
\item Operatsion tizimning muayyan apparat qurilmasi bilan ishlashiga yordam beradi
\item Kompyuterdagi barcha fayllarni avtomatik ravishda diskdan diskka ko'chiradi
\item Kompyuterning barcha tarmoq paketlarini tashqi serverlar bo'ylab yo'naltiradi
\item Matn faylidagi imlo xatolarini avtomatik topib o'rniga to'g'risini yozadi
\end{enumerate}
\begin{tcolorbox}[enhanced,breakable,colback=teal!4,colframe=teal!60!black,boxrule=0.45pt,arc=1.2mm,left=2.7mm,right=2.7mm,top=1.8mm,bottom=1.8mm,before skip=2mm,after skip=2mm]
{\bfseries\sffamily\color{teal}To'g'ri javob:} \textbf{A) Operatsion tizimning muayyan apparat qurilmasi bilan ishlashiga yordam beradi}
\par\smallskip
{\bfseries Izoh:} Qurilma drayveri operatsion tizim bilan muayyan apparat qurilmasi o'rtasida ishlash interfeysini ta'minlaydi. Drayver alohida imlo tekshiruvchi yoki router emas.
\end{tcolorbox}
\vspace{2mm}\hrule height 0.2pt\vspace{3mm}

% DB_ID 2cf941fb-d49e-41ed-bba0-26b4c207ca7a
\needspace{5\baselineskip}
{\bfseries\sffamily\color{ink} 474.} {scriptsizecolor{muted} Y1 | published | 2cf941fb}
\par
Moslang: 1) veb-sahifa ishlab chiqish vositalari; 2) ma'lumotlar bazasini boshqarish tizimlari; 3) grafika bilan ishlash dasturlari. a) CorelDRAW, IrfanView; b) PowerPoint, Access; c) HTML, CSS, JavaScript; d) Access, MySQL.
\begin{enumerate}[label=\Alph*),leftmargin=9mm]
\item 1-a; 2-d; 3-c
\item 1-b; 2-d; 3-a
\item 1-c; 2-d; 3-a
\item 1-c; 2-d; 3-b
\end{enumerate}
\begin{tcolorbox}[enhanced,breakable,colback=teal!4,colframe=teal!60!black,boxrule=0.45pt,arc=1.2mm,left=2.7mm,right=2.7mm,top=1.8mm,bottom=1.8mm,before skip=2mm,after skip=2mm]
{\bfseries\sffamily\color{teal}To'g'ri javob:} \textbf{C) 1-c; 2-d; 3-a}
\par\smallskip
{\bfseries Izoh:} Veb ishlab chiqish uchun c, MBBT uchun d, grafika uchun a mos: 1-c; 2-d; 3-a.
\end{tcolorbox}
\vspace{2mm}\hrule height 0.2pt\vspace{3mm}

% DB_ID 3b693d17-3ffb-4ce5-8cd7-194876f417fb
\needspace{5\baselineskip}
{\bfseries\sffamily\color{ink} 475.} {scriptsizecolor{muted} Y1 | published | 3b693d17}
\par
Videoklipga harakatsiz tasvir joylashtirish odatda qanday bajariladi?
\begin{enumerate}[label=\Alph*),leftmargin=9mm]
\item Tasvir import qilinib, \texttt{Add image} buyrug'i bilan vaqt shkalasiga joylanadi
\item Tasvir eksport qilinib, sudrab tashlash orqali vaqt shkalasiga joylanadi
\item Tasvir import qilinib, sudrab tashlash orqali vaqt shkalasiga joylanadi
\item Tasvir eksport qilinib, \texttt{Add image} buyrug'i bilan vaqt shkalasiga joylanadi
\end{enumerate}
\begin{tcolorbox}[enhanced,breakable,colback=teal!4,colframe=teal!60!black,boxrule=0.45pt,arc=1.2mm,left=2.7mm,right=2.7mm,top=1.8mm,bottom=1.8mm,before skip=2mm,after skip=2mm]
{\bfseries\sffamily\color{teal}To'g'ri javob:} \textbf{C) Tasvir import qilinib, sudrab tashlash orqali vaqt shkalasiga joylanadi}
\par\smallskip
{\bfseries Izoh:} Darslikda statik tasvir avval tahrirlash dasturiga import qilinishi, so'ng vaqt shkalasiga tortib joylashtirilishi ko'rsatilgan.
\end{tcolorbox}
\vspace{2mm}\hrule height 0.2pt\vspace{3mm}

% DB_ID 46a0c9b9-963e-4ec4-a5ac-f98cc65ffbaf
\needspace{5\baselineskip}
{\bfseries\sffamily\color{ink} 476.} {scriptsizecolor{muted} Y1 | published | 46a0c9b9}
\par
Antivirus kompyuterni tekshirib, xavfli fayllarni aniqlash va karantinga olish kabi texnik xizmat vazifalarini bajaradi. Vazifasiga ko'ra u qaysi dastur toifasiga kiradi?
\begin{enumerate}[label=\Alph*),leftmargin=9mm]
\item Xizmat ko'rsatuvchi (utility) dastur
\item Amaliy matn muharriri dasturi
\item Operatsion tizimning asosiy yadrosi
\item Qurilmaga xizmat qiluvchi drayver
\end{enumerate}
\begin{tcolorbox}[enhanced,breakable,colback=teal!4,colframe=teal!60!black,boxrule=0.45pt,arc=1.2mm,left=2.7mm,right=2.7mm,top=1.8mm,bottom=1.8mm,before skip=2mm,after skip=2mm]
{\bfseries\sffamily\color{teal}To'g'ri javob:} \textbf{A) Xizmat ko'rsatuvchi (utility) dastur}
\par\smallskip
{\bfseries Izoh:} Antivirus foydalanuvchi tizimiga texnik xizmat ko'rsatadigan utility dasturlarga kiradi. U matn muharriri yoki bevosita operatsion tizim yadrosi emas.
\end{tcolorbox}
\vspace{2mm}\hrule height 0.2pt\vspace{3mm}

% DB_ID 9ee1bdbd-e382-4626-9ecb-23fa439655b5
\needspace{5\baselineskip}
{\bfseries\sffamily\color{ink} 477.} {scriptsizecolor{muted} Y1 | published | 9ee1bdbd}
\par
Microsoft Word qaysi turdagi dasturiy ta'minotga kiradi?
\begin{enumerate}[label=\Alph*),leftmargin=9mm]
\item Amaliy dastur
\item Operatsion tizim
\item Qurilma drayveri
\item BIOS
\end{enumerate}
\begin{tcolorbox}[enhanced,breakable,colback=teal!4,colframe=teal!60!black,boxrule=0.45pt,arc=1.2mm,left=2.7mm,right=2.7mm,top=1.8mm,bottom=1.8mm,before skip=2mm,after skip=2mm]
{\bfseries\sffamily\color{teal}To'g'ri javob:} \textbf{A) Amaliy dastur}
\par\smallskip
{\bfseries Izoh:} Word foydalanuvchining hujjat yaratish vazifasini bajaradigan amaliy dasturdir.
\end{tcolorbox}
\vspace{2mm}\hrule height 0.2pt\vspace{3mm}

% DB_ID beae8d23-aac8-59f3-9055-6f9757f8e20f
\needspace{5\baselineskip}
{\bfseries\sffamily\color{ink} 478.} {scriptsizecolor{muted} Y1 | published | beae8d23}
\par
Dasturlarni eng aniq vazifaviy turiga moslang: 1) CorelDRAW; 2) Disk Cleanup.
\begin{enumerate}[label=\Alph*),leftmargin=9mm]
\item 1 - rastr grafik muharriri; 2 - operatsion tizim
\item 1 - vektor grafik muharriri; 2 - xizmat ko'rsatuvchi (utility) dastur
\item 1 - xizmat ko'rsatuvchi dastur; 2 - vektor grafik muharriri
\item 1 - matn muharriri; 2 - drayver
\end{enumerate}
\begin{tcolorbox}[enhanced,breakable,colback=teal!4,colframe=teal!60!black,boxrule=0.45pt,arc=1.2mm,left=2.7mm,right=2.7mm,top=1.8mm,bottom=1.8mm,before skip=2mm,after skip=2mm]
{\bfseries\sffamily\color{teal}To'g'ri javob:} \textbf{B) 1 - vektor grafik muharriri; 2 - xizmat ko'rsatuvchi (utility) dastur}
\par\smallskip
{\bfseries Izoh:} CorelDRAW vektor grafik muharriri, Disk Cleanup esa diskni keraksiz fayllardan tozalovchi utility dastur.
\end{tcolorbox}
\vspace{2mm}\hrule height 0.2pt\vspace{3mm}

% DB_ID ce5c388e-9913-43b5-8f30-a4d4df7b09c3
\needspace{5\baselineskip}
{\bfseries\sffamily\color{ink} 479.} {scriptsizecolor{muted} Y1 | published | ce5c388e}
\par
"Dasturiy ta'minot" atamasiga eng aniq ta'rifni toping.
\begin{enumerate}[label=\Alph*),leftmargin=9mm]
\item tizimni yaratadigan, boshqaradigan va undan foydalanadigan insonlar hamda ularning bilimlari
\item qayd etilgan belgi, qiymat yoki kuzatuv bo'lib, kontekstsiz mazmuni yetarli bo'lmasligi mumkin
\item kompyuter bajaradigan dasturlar, operatsion tizimlar va algoritmik vositalar majmui
\item qabul qiluvchi uchun mazmunli bo'lib, noaniqlikni kamaytiradigan ma'lumot
\end{enumerate}
\begin{tcolorbox}[enhanced,breakable,colback=teal!4,colframe=teal!60!black,boxrule=0.45pt,arc=1.2mm,left=2.7mm,right=2.7mm,top=1.8mm,bottom=1.8mm,before skip=2mm,after skip=2mm]
{\bfseries\sffamily\color{teal}To'g'ri javob:} \textbf{C) kompyuter bajaradigan dasturlar, operatsion tizimlar va algoritmik vositalar majmui}
\par\smallskip
{\bfseries Izoh:} Dasturiy ta'minot - kompyuter bajaradigan dasturlar, operatsion tizimlar va algoritmik vositalar majmui.
\end{tcolorbox}
\vspace{2mm}\hrule height 0.2pt\vspace{3mm}

% DB_ID 27e0f930-ea29-4c86-a340-a21879d948ce
\needspace{5\baselineskip}
{\bfseries\sffamily\color{ink} 480.} {scriptsizecolor{muted} Y1 | published | 27e0f930}
\par
Quyidagi ta'rif qaysi atamaga tegishli? Kompyuter bajaradigan dasturlar, operatsion tizimlar va algoritmik vositalar majmui.
\begin{enumerate}[label=\Alph*),leftmargin=9mm]
\item Dasturiy ta'minot
\item Insoniy ta'minot
\item Ma'lumot
\item Axborot
\end{enumerate}
\begin{tcolorbox}[enhanced,breakable,colback=teal!4,colframe=teal!60!black,boxrule=0.45pt,arc=1.2mm,left=2.7mm,right=2.7mm,top=1.8mm,bottom=1.8mm,before skip=2mm,after skip=2mm]
{\bfseries\sffamily\color{teal}To'g'ri javob:} \textbf{A) Dasturiy ta'minot}
\par\smallskip
{\bfseries Izoh:} Ta'rif Dasturiy ta'minot tushunchasiga tegishli.
\end{tcolorbox}
\vspace{2mm}\hrule height 0.2pt\vspace{3mm}

% DB_ID 16f0e147-b66d-4c35-abc9-1d3fe8853b23
\needspace{5\baselineskip}
{\bfseries\sffamily\color{ink} 481.} {scriptsizecolor{muted} Y1 | published | 16f0e147}
\par
Kompyuterga yangi videokarta o'rnatildi. Operatsion tizim u bilan to'g'ri ishlashi uchun qaysi turdagi dastur kerak bo'lishi mumkin?
\begin{enumerate}[label=\Alph*),leftmargin=9mm]
\item Qurilma drayveri
\item Matn muharriri
\item Taqdimot dasturi
\item Arxivator
\end{enumerate}
\begin{tcolorbox}[enhanced,breakable,colback=teal!4,colframe=teal!60!black,boxrule=0.45pt,arc=1.2mm,left=2.7mm,right=2.7mm,top=1.8mm,bottom=1.8mm,before skip=2mm,after skip=2mm]
{\bfseries\sffamily\color{teal}To'g'ri javob:} \textbf{A) Qurilma drayveri}
\par\smallskip
{\bfseries Izoh:} Videokarta drayveri tizim va apparat o'rtasidagi ishlashni ta'minlaydi.
\end{tcolorbox}
\vspace{2mm}\hrule height 0.2pt\vspace{3mm}

% DB_ID 8c04d3fc-4a03-4a70-a602-db8625affa55
\needspace{5\baselineskip}
{\bfseries\sffamily\color{ink} 482.} {scriptsizecolor{muted} Y1 | published | 8c04d3fc}
\par
Diskdagi fayllarni bitta siqilgan paketga jamlash kerak. Qaysi dastur turi mos?
\begin{enumerate}[label=\Alph*),leftmargin=9mm]
\item Arxivator utility
\item Elektron jadval
\item Operatsion tizim yadrosi
\item Grafik drayver
\end{enumerate}
\begin{tcolorbox}[enhanced,breakable,colback=teal!4,colframe=teal!60!black,boxrule=0.45pt,arc=1.2mm,left=2.7mm,right=2.7mm,top=1.8mm,bottom=1.8mm,before skip=2mm,after skip=2mm]
{\bfseries\sffamily\color{teal}To'g'ri javob:} \textbf{A) Arxivator utility}
\par\smallskip
{\bfseries Izoh:} Arxivator fayllarni siqish va paketlashga xizmat qiladigan utility dasturdir.
\end{tcolorbox}
\vspace{2mm}\hrule height 0.2pt\vspace{3mm}

% DB_ID 91d3ae5b-8616-4825-8d6d-37d87444921d
\needspace{5\baselineskip}
{\bfseries\sffamily\color{ink} 483.} {scriptsizecolor{muted} Y2 | published | 91d3ae5b}
\par
Dasturiy vositalarni ularning roliga moslang.
\textbf{Chap ustun:}
\begin{enumerate}[label=\arabic*),leftmargin=9mm]
\item Windows
\item Microsoft Word
\item Antivirus
\item Qurilma drayveri
\end{enumerate}
\textbf{O'ng ustun:}
\begin{enumerate}[label=\Alph*),leftmargin=9mm]
\item Xizmat ko'rsatuvchi himoya dasturi
\item Operatsion tizim
\item Operatsion tizimning apparat bilan ishlashini ta'minlovchi tizimli komponent
\item Amaliy matn protsessori
\end{enumerate}
\begin{tcolorbox}[enhanced,breakable,colback=teal!4,colframe=teal!60!black,boxrule=0.45pt,arc=1.2mm,left=2.7mm,right=2.7mm,top=1.8mm,bottom=1.8mm,before skip=2mm,after skip=2mm]
{\bfseries\sffamily\color{teal}To'g'ri javob:} \textbf{1-B, 2-D, 3-A, 4-C}
\par\smallskip
{\bfseries Izoh:} Windows - OT; Word - amaliy dastur; antivirus - utility; drayver - apparat bilan ishlashga yordam beruvchi tizimli komponent.
\end{tcolorbox}
\vspace{2mm}\hrule height 0.2pt\vspace{3mm}

% DB_ID 9c120793-e765-4a2a-a109-08c46f65b524
\needspace{5\baselineskip}
{\bfseries\sffamily\color{ink} 484.} {scriptsizecolor{muted} Y1 | published | 9c120793}
\par
Foydalanuvchi matnli hisobot tayyorlashi kerak. Qaysi dastur turi eng mos?
\begin{enumerate}[label=\Alph*),leftmargin=9mm]
\item Amaliy matn protsessori
\item Qurilma drayveri
\item Operatsion tizim yadrosi
\item BIOS utilitasi
\end{enumerate}
\begin{tcolorbox}[enhanced,breakable,colback=teal!4,colframe=teal!60!black,boxrule=0.45pt,arc=1.2mm,left=2.7mm,right=2.7mm,top=1.8mm,bottom=1.8mm,before skip=2mm,after skip=2mm]
{\bfseries\sffamily\color{teal}To'g'ri javob:} \textbf{A) Amaliy matn protsessori}
\par\smallskip
{\bfseries Izoh:} Hisobot yozish foydalanuvchi vazifasi bo'lib, matn protsessori amaliy dastur hisoblanadi.
\end{tcolorbox}
\vspace{2mm}\hrule height 0.2pt\vspace{3mm}

% DB_ID c1fe006b-dc3c-4894-874c-2990271d822d
\needspace{5\baselineskip}
{\bfseries\sffamily\color{ink} 485.} {scriptsizecolor{muted} Y1 | published | c1fe006b}
\par
Foydalanuvchi zararli fayllarni aniqlab karantinga olishni xohlaydi. Qaysi dastur eng mos?
\begin{enumerate}[label=\Alph*),leftmargin=9mm]
\item Antivirus utility
\item Matn protsessori
\item Grafik muharrir
\item Taqdimot dasturi
\end{enumerate}
\begin{tcolorbox}[enhanced,breakable,colback=teal!4,colframe=teal!60!black,boxrule=0.45pt,arc=1.2mm,left=2.7mm,right=2.7mm,top=1.8mm,bottom=1.8mm,before skip=2mm,after skip=2mm]
{\bfseries\sffamily\color{teal}To'g'ri javob:} \textbf{A) Antivirus utility}
\par\smallskip
{\bfseries Izoh:} Antivirus zararli dasturlarni aniqlash va zararsizlantirish uchun mo'ljallangan utilitydir.
\end{tcolorbox}
\vspace{2mm}\hrule height 0.2pt\vspace{3mm}

% DB_ID e9f1db9f-7259-4349-ae1b-5c189a70f467
\needspace{5\baselineskip}
{\bfseries\sffamily\color{ink} 486.} {scriptsizecolor{muted} Y1 | published | e9f1db9f}
\par
Quyidagi dasturlardan qaysilari tizimli dastur (system software) hisoblanadi va qaysilari amaliy dastur (application software)? Antivirus, MS Excel, drayverlar dasturi, brauzer, Photoshop.
\begin{enumerate}[label=\Alph*),leftmargin=9mm]
\item Tizimli: MS Excel, Photoshop; amaliy: antivirus, drayverlar
\item Tizimli: antivirus, drayverlar; amaliy: MS Excel, brauzer, Photoshop
\item Brauzer va antivirus faqat tizimli dastur, qolganlari amaliy emas
\item Barchasi amaliy dasturlar hisoblanadi
\end{enumerate}
\begin{tcolorbox}[enhanced,breakable,colback=teal!4,colframe=teal!60!black,boxrule=0.45pt,arc=1.2mm,left=2.7mm,right=2.7mm,top=1.8mm,bottom=1.8mm,before skip=2mm,after skip=2mm]
{\bfseries\sffamily\color{teal}To'g'ri javob:} \textbf{B) Tizimli: antivirus, drayverlar; amaliy: MS Excel, brauzer, Photoshop}
\par\smallskip
{\bfseries Izoh:} Antivirus va drayverlar xizmat/tizimli, Excel, brauzer va Photoshop foydalanuvchi vazifalarini bajaruvchi amaliy dasturlardir.
\end{tcolorbox}
\vspace{2mm}\hrule height 0.2pt\vspace{3mm}

% DB_ID eadcb13e-a515-4faf-958c-8250040bf76f
\needspace{5\baselineskip}
{\bfseries\sffamily\color{ink} 487.} {scriptsizecolor{muted} Y1 | published | eadcb13e}
\par
Yangi operatsion tizim va matn muharriri kompyuterlarga o'rnatildi. Bu vaziyatdagi asosiy tushunchani aniqlang.
\begin{enumerate}[label=\Alph*),leftmargin=9mm]
\item Insoniy ta'minot
\item Ma'lumot
\item Dasturiy ta'minot
\item Axborot
\end{enumerate}
\begin{tcolorbox}[enhanced,breakable,colback=teal!4,colframe=teal!60!black,boxrule=0.45pt,arc=1.2mm,left=2.7mm,right=2.7mm,top=1.8mm,bottom=1.8mm,before skip=2mm,after skip=2mm]
{\bfseries\sffamily\color{teal}To'g'ri javob:} \textbf{C) Dasturiy ta'minot}
\par\smallskip
{\bfseries Izoh:} Vaziyat Dasturiy ta'minotga misol: kompyuter bajaradigan dasturlar, operatsion tizimlar va algoritmik vositalar majmui.
\end{tcolorbox}
\vspace{2mm}\hrule height 0.2pt\vspace{3mm}
\subsection{Tekshiruvdagi testlar}

% DB_ID 10dac0c2-0e20-cb06-64d5-f39758f91886
\needspace{5\baselineskip}
{\bfseries\sffamily\color{ink} 488.} {scriptsizecolor{muted} Y1 | review | 10dac0c2}
\par
Drayver, utilita, firmware va translatorlar doirasida "kompilyator" atamasining eng aniq ta'rifini toping.
\begin{enumerate}[label=\Alph*),leftmargin=9mm]
\item OT bilan aniq qurilma modeli o'rtasida buyruqlarni moslashtiruvchi dastur
\item tizimga xizmat ko'rsatuvchi yordamchi dastur
\item qurilma ichidagi past darajali boshqaruv dasturi
\item butun manba kodini obyekt yoki bajariladigan kodga o'giruvchi translator
\end{enumerate}
\begin{tcolorbox}[enhanced,breakable,colback=teal!4,colframe=teal!60!black,boxrule=0.45pt,arc=1.2mm,left=2.7mm,right=2.7mm,top=1.8mm,bottom=1.8mm,before skip=2mm,after skip=2mm]
{\bfseries\sffamily\color{teal}To'g'ri javob:} \textbf{D) butun manba kodini obyekt yoki bajariladigan kodga o'giruvchi translator}
\par\smallskip
{\bfseries Izoh:} butun manba kodini obyekt yoki bajariladigan kodga o'giruvchi translator. kompilyator - butun manba kodini obyekt yoki bajariladigan kodga o'giruvchi translator. Manba: ICT8, PDF 106-110; Qo'shimcha sahifa raqami (qaysi PDFga tegishli ekani aniqlanmagan; darslik sahifasi deb tasdiqlanmagan): 72
\end{tcolorbox}
\vspace{2mm}\hrule height 0.2pt\vspace{3mm}

% DB_ID 173ceab3-f91e-568b-7f86-35861e6b49cd
\needspace{5\baselineskip}
{\bfseries\sffamily\color{ink} 489.} {scriptsizecolor{muted} Y1 | review | 173ceab3}
\par
Drayver, utilita, firmware va translatorlar doirasida "drayver" atamasining eng aniq ta'rifini toping.
\begin{enumerate}[label=\Alph*),leftmargin=9mm]
\item OT bilan aniq qurilma modeli o'rtasida buyruqlarni moslashtiruvchi dastur
\item tizimga xizmat ko'rsatuvchi yordamchi dastur
\item qurilma ichidagi past darajali boshqaruv dasturi
\item butun manba kodini obyekt yoki bajariladigan kodga o'giruvchi translator
\end{enumerate}
\begin{tcolorbox}[enhanced,breakable,colback=teal!4,colframe=teal!60!black,boxrule=0.45pt,arc=1.2mm,left=2.7mm,right=2.7mm,top=1.8mm,bottom=1.8mm,before skip=2mm,after skip=2mm]
{\bfseries\sffamily\color{teal}To'g'ri javob:} \textbf{A) OT bilan aniq qurilma modeli o'rtasida buyruqlarni moslashtiruvchi dastur}
\par\smallskip
{\bfseries Izoh:} OT bilan aniq qurilma modeli o'rtasida buyruqlarni moslashtiruvchi dastur. drayver - OT bilan aniq qurilma modeli o'rtasida buyruqlarni moslashtiruvchi dastur. Manba: CAM9, PDF 14-15; Qo'shimcha sahifa raqami (qaysi PDFga tegishli ekani aniqlanmagan; darslik sahifasi deb tasdiqlanmagan): 69
\end{tcolorbox}
\vspace{2mm}\hrule height 0.2pt\vspace{3mm}

% DB_ID 30b03e82-7727-f1d4-9542-4f4b634c7ca1
\needspace{5\baselineskip}
{\bfseries\sffamily\color{ink} 490.} {scriptsizecolor{muted} Y1 | review | 30b03e82}
\par
Amaliy dasturlar va vazifaga mos tanlov doirasida "elektron jadval" atamasining eng aniq ta'rifini toping.
\begin{enumerate}[label=\Alph*),leftmargin=9mm]
\item tuzilmali matnli hujjat yaratish va formatlash dasturi
\item tuzilmali ma'lumotlar bazasini yaratish, so'rash va boshqarish tizimi
\item professional sahifa maketi va nashr tayyorlash dasturi
\item formulalar asosida hisoblash, model va diagramma yaratish dasturi
\end{enumerate}
\begin{tcolorbox}[enhanced,breakable,colback=teal!4,colframe=teal!60!black,boxrule=0.45pt,arc=1.2mm,left=2.7mm,right=2.7mm,top=1.8mm,bottom=1.8mm,before skip=2mm,after skip=2mm]
{\bfseries\sffamily\color{teal}To'g'ri javob:} \textbf{D) formulalar asosida hisoblash, model va diagramma yaratish dasturi}
\par\smallskip
{\bfseries Izoh:} formulalar asosida hisoblash, model va diagramma yaratish dasturi. elektron jadval - formulalar asosida hisoblash, model va diagramma yaratish dasturi. Manba: CAM6, PDF 34-205; Qo'shimcha sahifa raqami (qaysi PDFga tegishli ekani aniqlanmagan; darslik sahifasi deb tasdiqlanmagan): 77
\end{tcolorbox}
\vspace{2mm}\hrule height 0.2pt\vspace{3mm}

% DB_ID 3208c1fa-851b-0ffa-ec9d-45864bf09a19
\needspace{5\baselineskip}
{\bfseries\sffamily\color{ink} 491.} {scriptsizecolor{muted} Y1 | review | 3208c1fa}
\par
Dastur o'rnatish, yangilash, litsenziya va moslik doirasida "installation" atamasining eng aniq ta'rifini toping.
\begin{enumerate}[label=\Alph*),leftmargin=9mm]
\item dastur va ro'yxatdan o'tgan komponentlarini nazoratli olib tashlash
\item mavjud versiyaga tuzatish yoki kichik yaxshilanish o'rnatish
\item dastur komponentlarini tizimda ishlashga tayyorlab joylashtirish
\item qurilma, OT, dastur va formatlarning birgalikda ishlay olish xususiyati
\end{enumerate}
\begin{tcolorbox}[enhanced,breakable,colback=teal!4,colframe=teal!60!black,boxrule=0.45pt,arc=1.2mm,left=2.7mm,right=2.7mm,top=1.8mm,bottom=1.8mm,before skip=2mm,after skip=2mm]
{\bfseries\sffamily\color{teal}To'g'ri javob:} \textbf{C) dastur komponentlarini tizimda ishlashga tayyorlab joylashtirish}
\par\smallskip
{\bfseries Izoh:} dastur komponentlarini tizimda ishlashga tayyorlab joylashtirish. installation - dastur komponentlarini tizimda ishlashga tayyorlab joylashtirish. Manba: MAR, PDF 10-16; Qo'shimcha sahifa raqami (qaysi PDFga tegishli ekani aniqlanmagan; darslik sahifasi deb tasdiqlanmagan): 101
\end{tcolorbox}
\vspace{2mm}\hrule height 0.2pt\vspace{3mm}

% DB_ID 377fa4a5-fdf3-b8fc-71f7-a11920a87196
\needspace{5\baselineskip}
{\bfseries\sffamily\color{ink} 492.} {scriptsizecolor{muted} Y1 | review | 377fa4a5}
\par
Dastur o'rnatish, yangilash, litsenziya va moslik doirasida "open-source" atamasining eng aniq ta'rifini toping.
\begin{enumerate}[label=\Alph*),leftmargin=9mm]
\item dastur komponentlarini tizimda ishlashga tayyorlab joylashtirish
\item dastur va ro'yxatdan o'tgan komponentlarini nazoratli olib tashlash
\item manba kodi litsenziya asosida ochiq dastur
\item mavjud versiyaga tuzatish yoki kichik yaxshilanish o'rnatish
\end{enumerate}
\begin{tcolorbox}[enhanced,breakable,colback=teal!4,colframe=teal!60!black,boxrule=0.45pt,arc=1.2mm,left=2.7mm,right=2.7mm,top=1.8mm,bottom=1.8mm,before skip=2mm,after skip=2mm]
{\bfseries\sffamily\color{teal}To'g'ri javob:} \textbf{C) manba kodi litsenziya asosida ochiq dastur}
\par\smallskip
{\bfseries Izoh:} manba kodi litsenziya asosida ochiq dastur. open-source - manba kodi litsenziya asosida ochiq dastur. Manba: MAR, PDF 10-16; Qo'shimcha sahifa raqami (qaysi PDFga tegishli ekani aniqlanmagan; darslik sahifasi deb tasdiqlanmagan): 106
\end{tcolorbox}
\vspace{2mm}\hrule height 0.2pt\vspace{3mm}

% DB_ID 430834e1-fe29-83c0-a2e0-0251b0441893
\needspace{5\baselineskip}
{\bfseries\sffamily\color{ink} 493.} {scriptsizecolor{muted} Y1 | review | 430834e1}
\par
Dastur o'rnatish, yangilash, litsenziya va moslik doirasida "update" atamasining eng aniq ta'rifini toping.
\begin{enumerate}[label=\Alph*),leftmargin=9mm]
\item mavjud versiyaga tuzatish yoki kichik yaxshilanish o'rnatish
\item dastur komponentlarini tizimda ishlashga tayyorlab joylashtirish
\item dastur va ro'yxatdan o'tgan komponentlarini nazoratli olib tashlash
\item qurilma, OT, dastur va formatlarning birgalikda ishlay olish xususiyati
\end{enumerate}
\begin{tcolorbox}[enhanced,breakable,colback=teal!4,colframe=teal!60!black,boxrule=0.45pt,arc=1.2mm,left=2.7mm,right=2.7mm,top=1.8mm,bottom=1.8mm,before skip=2mm,after skip=2mm]
{\bfseries\sffamily\color{teal}To'g'ri javob:} \textbf{A) mavjud versiyaga tuzatish yoki kichik yaxshilanish o'rnatish}
\par\smallskip
{\bfseries Izoh:} mavjud versiyaga tuzatish yoki kichik yaxshilanish o'rnatish. update - mavjud versiyaga tuzatish yoki kichik yaxshilanish o'rnatish. Manba: ICT8, PDF 106-110; Qo'shimcha sahifa raqami (qaysi PDFga tegishli ekani aniqlanmagan; darslik sahifasi deb tasdiqlanmagan): 103
\end{tcolorbox}
\vspace{2mm}\hrule height 0.2pt\vspace{3mm}

% DB_ID 470263e8-74ea-61c2-b737-5de348af15ca
\needspace{5\baselineskip}
{\bfseries\sffamily\color{ink} 494.} {scriptsizecolor{muted} Y1 | review | 470263e8}
\par
Dasturiy ta'minotning tasnifi doirasida "tizimli dastur" atamasining eng aniq ta'rifini toping.
\begin{enumerate}[label=\Alph*),leftmargin=9mm]
\item foydalanuvchining aniq vazifasini bajaruvchi dastur
\item tizimga xizmat ko'rsatuvchi yordamchi dastur
\item kompyuter resurslarini boshqarib, boshqa dasturlarga muhit yaratuvchi dastur
\item dastur kodini boshqa, odatda mashina bajaradigan shaklga o'giruvchi vosita
\end{enumerate}
\begin{tcolorbox}[enhanced,breakable,colback=teal!4,colframe=teal!60!black,boxrule=0.45pt,arc=1.2mm,left=2.7mm,right=2.7mm,top=1.8mm,bottom=1.8mm,before skip=2mm,after skip=2mm]
{\bfseries\sffamily\color{teal}To'g'ri javob:} \textbf{C) kompyuter resurslarini boshqarib, boshqa dasturlarga muhit yaratuvchi dastur}
\par\smallskip
{\bfseries Izoh:} kompyuter resurslarini boshqarib, boshqa dasturlarga muhit yaratuvchi dastur. tizimli dastur - kompyuter resurslarini boshqarib, boshqa dasturlarga muhit yaratuvchi dastur. Manba: CAM9, PDF 11-15; Qo'shimcha sahifa raqami (qaysi PDFga tegishli ekani aniqlanmagan; darslik sahifasi deb tasdiqlanmagan): 49
\end{tcolorbox}
\vspace{2mm}\hrule height 0.2pt\vspace{3mm}

% DB_ID 551ff514-4960-3775-09a3-90b5df8abda2
\needspace{5\baselineskip}
{\bfseries\sffamily\color{ink} 495.} {scriptsizecolor{muted} Y1 | review | 551ff514}
\par
Drayver, utilita, firmware va translatorlar doirasida "linker" atamasining eng aniq ta'rifini toping.
\begin{enumerate}[label=\Alph*),leftmargin=9mm]
\item OT bilan aniq qurilma modeli o'rtasida buyruqlarni moslashtiruvchi dastur
\item obyekt modullari va kutubxonalarni bajariladigan faylga bog'lovchi vosita
\item tizimga xizmat ko'rsatuvchi yordamchi dastur
\item qurilma ichidagi past darajali boshqaruv dasturi
\end{enumerate}
\begin{tcolorbox}[enhanced,breakable,colback=teal!4,colframe=teal!60!black,boxrule=0.45pt,arc=1.2mm,left=2.7mm,right=2.7mm,top=1.8mm,bottom=1.8mm,before skip=2mm,after skip=2mm]
{\bfseries\sffamily\color{teal}To'g'ri javob:} \textbf{B) obyekt modullari va kutubxonalarni bajariladigan faylga bog'lovchi vosita}
\par\smallskip
{\bfseries Izoh:} obyekt modullari va kutubxonalarni bajariladigan faylga bog'lovchi vosita. linker - obyekt modullari va kutubxonalarni bajariladigan faylga bog'lovchi vosita. Manba: CAM10/11, PDF 44-49; Qo'shimcha sahifa raqami (qaysi PDFga tegishli ekani aniqlanmagan; darslik sahifasi deb tasdiqlanmagan): 74
\end{tcolorbox}
\vspace{2mm}\hrule height 0.2pt\vspace{3mm}

% DB_ID 5bc0a2f3-1318-f109-9b75-99ac1b348965
\needspace{5\baselineskip}
{\bfseries\sffamily\color{ink} 496.} {scriptsizecolor{muted} Y1 | review | 5bc0a2f3}
\par
Dastur o'rnatish, yangilash, litsenziya va moslik doirasida "uninstall" atamasining eng aniq ta'rifini toping.
\begin{enumerate}[label=\Alph*),leftmargin=9mm]
\item dastur komponentlarini tizimda ishlashga tayyorlab joylashtirish
\item mavjud versiyaga tuzatish yoki kichik yaxshilanish o'rnatish
\item qurilma, OT, dastur va formatlarning birgalikda ishlay olish xususiyati
\item dastur va ro'yxatdan o'tgan komponentlarini nazoratli olib tashlash
\end{enumerate}
\begin{tcolorbox}[enhanced,breakable,colback=teal!4,colframe=teal!60!black,boxrule=0.45pt,arc=1.2mm,left=2.7mm,right=2.7mm,top=1.8mm,bottom=1.8mm,before skip=2mm,after skip=2mm]
{\bfseries\sffamily\color{teal}To'g'ri javob:} \textbf{D) dastur va ro'yxatdan o'tgan komponentlarini nazoratli olib tashlash}
\par\smallskip
{\bfseries Izoh:} dastur va ro'yxatdan o'tgan komponentlarini nazoratli olib tashlash. uninstall - dastur va ro'yxatdan o'tgan komponentlarini nazoratli olib tashlash. Manba: CAM8, PDF 99-110; Qo'shimcha sahifa raqami (qaysi PDFga tegishli ekani aniqlanmagan; darslik sahifasi deb tasdiqlanmagan): 102
\end{tcolorbox}
\vspace{2mm}\hrule height 0.2pt\vspace{3mm}

% DB_ID 6db806b7-98f5-c865-1a46-3ab7877b2216
\needspace{5\baselineskip}
{\bfseries\sffamily\color{ink} 497.} {scriptsizecolor{muted} Y1 | review | 6db806b7}
\par
Drayver, utilita, firmware va translatorlar doirasida "interpretator" atamasining eng aniq ta'rifini toping.
\begin{enumerate}[label=\Alph*),leftmargin=9mm]
\item kodni bajarish paytida bosqichma-bosqich tarjima qiluvchi translator
\item OT bilan aniq qurilma modeli o'rtasida buyruqlarni moslashtiruvchi dastur
\item tizimga xizmat ko'rsatuvchi yordamchi dastur
\item qurilma ichidagi past darajali boshqaruv dasturi
\end{enumerate}
\begin{tcolorbox}[enhanced,breakable,colback=teal!4,colframe=teal!60!black,boxrule=0.45pt,arc=1.2mm,left=2.7mm,right=2.7mm,top=1.8mm,bottom=1.8mm,before skip=2mm,after skip=2mm]
{\bfseries\sffamily\color{teal}To'g'ri javob:} \textbf{A) kodni bajarish paytida bosqichma-bosqich tarjima qiluvchi translator}
\par\smallskip
{\bfseries Izoh:} kodni bajarish paytida bosqichma-bosqich tarjima qiluvchi translator. interpretator - kodni bajarish paytida bosqichma-bosqich tarjima qiluvchi translator. Manba: CAM9, PDF 14-15; Qo'shimcha sahifa raqami (qaysi PDFga tegishli ekani aniqlanmagan; darslik sahifasi deb tasdiqlanmagan): 73
\end{tcolorbox}
\vspace{2mm}\hrule height 0.2pt\vspace{3mm}

% DB_ID 73088236-3574-aabf-13ee-a6bc036c2d4d
\needspace{5\baselineskip}
{\bfseries\sffamily\color{ink} 498.} {scriptsizecolor{muted} Y1 | review | 73088236}
\par
Amaliy dasturlar va vazifaga mos tanlov doirasida "DTP" atamasining eng aniq ta'rifini toping.
\begin{enumerate}[label=\Alph*),leftmargin=9mm]
\item tuzilmali matnli hujjat yaratish va formatlash dasturi
\item professional sahifa maketi va nashr tayyorlash dasturi
\item formulalar asosida hisoblash, model va diagramma yaratish dasturi
\item tuzilmali ma'lumotlar bazasini yaratish, so'rash va boshqarish tizimi
\end{enumerate}
\begin{tcolorbox}[enhanced,breakable,colback=teal!4,colframe=teal!60!black,boxrule=0.45pt,arc=1.2mm,left=2.7mm,right=2.7mm,top=1.8mm,bottom=1.8mm,before skip=2mm,after skip=2mm]
{\bfseries\sffamily\color{teal}To'g'ri javob:} \textbf{B) professional sahifa maketi va nashr tayyorlash dasturi}
\par\smallskip
{\bfseries Izoh:} professional sahifa maketi va nashr tayyorlash dasturi. DTP - professional sahifa maketi va nashr tayyorlash dasturi. Manba: CAM5, PDF 45-138; Qo'shimcha sahifa raqami (qaysi PDFga tegishli ekani aniqlanmagan; darslik sahifasi deb tasdiqlanmagan): 79
\end{tcolorbox}
\vspace{2mm}\hrule height 0.2pt\vspace{3mm}

% DB_ID 79410c6c-24c0-72ee-514a-6c8d9f1118b3
\needspace{5\baselineskip}
{\bfseries\sffamily\color{ink} 499.} {scriptsizecolor{muted} Y1 | review | 79410c6c}
\par
Dasturiy ta'minotning tasnifi doirasida "maxsus dastur" atamasining eng aniq ta'rifini toping.
\begin{enumerate}[label=\Alph*),leftmargin=9mm]
\item kompyuter resurslarini boshqarib, boshqa dasturlarga muhit yaratuvchi dastur
\item foydalanuvchining aniq vazifasini bajaruvchi dastur
\item tizimga xizmat ko'rsatuvchi yordamchi dastur
\item muayyan buyurtmachi yoki jarayon uchun yaratilgan dastur
\end{enumerate}
\begin{tcolorbox}[enhanced,breakable,colback=teal!4,colframe=teal!60!black,boxrule=0.45pt,arc=1.2mm,left=2.7mm,right=2.7mm,top=1.8mm,bottom=1.8mm,before skip=2mm,after skip=2mm]
{\bfseries\sffamily\color{teal}To'g'ri javob:} \textbf{D) muayyan buyurtmachi yoki jarayon uchun yaratilgan dastur}
\par\smallskip
{\bfseries Izoh:} muayyan buyurtmachi yoki jarayon uchun yaratilgan dastur. maxsus dastur - muayyan buyurtmachi yoki jarayon uchun yaratilgan dastur. Manba: CAM10/11, PDF 44-58; Qo'shimcha sahifa raqami (qaysi PDFga tegishli ekani aniqlanmagan; darslik sahifasi deb tasdiqlanmagan): 54
\end{tcolorbox}
\vspace{2mm}\hrule height 0.2pt\vspace{3mm}

% DB_ID 794f884d-567e-ce86-7b13-a0a631da9241
\needspace{5\baselineskip}
{\bfseries\sffamily\color{ink} 500.} {scriptsizecolor{muted} Y1 | review | 794f884d}
\par
Dasturiy ta'minotning tasnifi doirasida "umumiy maqsadli dastur" atamasining eng aniq ta'rifini toping.
\begin{enumerate}[label=\Alph*),leftmargin=9mm]
\item kompyuter resurslarini boshqarib, boshqa dasturlarga muhit yaratuvchi dastur
\item foydalanuvchining aniq vazifasini bajaruvchi dastur
\item ko'p foydalanuvchi va vazifalarga mo'ljallangan tayyor dastur
\item tizimga xizmat ko'rsatuvchi yordamchi dastur
\end{enumerate}
\begin{tcolorbox}[enhanced,breakable,colback=teal!4,colframe=teal!60!black,boxrule=0.45pt,arc=1.2mm,left=2.7mm,right=2.7mm,top=1.8mm,bottom=1.8mm,before skip=2mm,after skip=2mm]
{\bfseries\sffamily\color{teal}To'g'ri javob:} \textbf{C) ko'p foydalanuvchi va vazifalarga mo'ljallangan tayyor dastur}
\par\smallskip
{\bfseries Izoh:} ko'p foydalanuvchi va vazifalarga mo'ljallangan tayyor dastur. umumiy maqsadli dastur - ko'p foydalanuvchi va vazifalarga mo'ljallangan tayyor dastur. Manba: CAM9, PDF 11-15; Qo'shimcha sahifa raqami (qaysi PDFga tegishli ekani aniqlanmagan; darslik sahifasi deb tasdiqlanmagan): 53
\end{tcolorbox}
\vspace{2mm}\hrule height 0.2pt\vspace{3mm}

% DB_ID 8df30060-61b3-53bb-c0ec-60a7a60ffc9c
\needspace{5\baselineskip}
{\bfseries\sffamily\color{ink} 501.} {scriptsizecolor{muted} Y1 | review | 8df30060}
\par
Amaliy dasturlar va vazifaga mos tanlov doirasida "CAD" atamasining eng aniq ta'rifini toping.
\begin{enumerate}[label=\Alph*),leftmargin=9mm]
\item tuzilmali matnli hujjat yaratish va formatlash dasturi
\item formulalar asosida hisoblash, model va diagramma yaratish dasturi
\item aniq texnik chizma va model yaratish dasturi
\item tuzilmali ma'lumotlar bazasini yaratish, so'rash va boshqarish tizimi
\end{enumerate}
\begin{tcolorbox}[enhanced,breakable,colback=teal!4,colframe=teal!60!black,boxrule=0.45pt,arc=1.2mm,left=2.7mm,right=2.7mm,top=1.8mm,bottom=1.8mm,before skip=2mm,after skip=2mm]
{\bfseries\sffamily\color{teal}To'g'ri javob:} \textbf{C) aniq texnik chizma va model yaratish dasturi}
\par\smallskip
{\bfseries Izoh:} aniq texnik chizma va model yaratish dasturi. CAD - aniq texnik chizma va model yaratish dasturi. Manba: ICT6, PDF 62-126; Qo'shimcha sahifa raqami (qaysi PDFga tegishli ekani aniqlanmagan; darslik sahifasi deb tasdiqlanmagan): 80
\end{tcolorbox}
\vspace{2mm}\hrule height 0.2pt\vspace{3mm}

% DB_ID 9cf48587-541e-3408-e027-44a1ebf8814f
\needspace{5\baselineskip}
{\bfseries\sffamily\color{ink} 502.} {scriptsizecolor{muted} Y1 | review | 9cf48587}
\par
Dasturiy ta'minotning tasnifi doirasida "amaliy dastur" atamasining eng aniq ta'rifini toping.
\begin{enumerate}[label=\Alph*),leftmargin=9mm]
\item kompyuter resurslarini boshqarib, boshqa dasturlarga muhit yaratuvchi dastur
\item tizimga xizmat ko'rsatuvchi yordamchi dastur
\item dastur kodini boshqa, odatda mashina bajaradigan shaklga o'giruvchi vosita
\item foydalanuvchining aniq vazifasini bajaruvchi dastur
\end{enumerate}
\begin{tcolorbox}[enhanced,breakable,colback=teal!4,colframe=teal!60!black,boxrule=0.45pt,arc=1.2mm,left=2.7mm,right=2.7mm,top=1.8mm,bottom=1.8mm,before skip=2mm,after skip=2mm]
{\bfseries\sffamily\color{teal}To'g'ri javob:} \textbf{D) foydalanuvchining aniq vazifasini bajaruvchi dastur}
\par\smallskip
{\bfseries Izoh:} foydalanuvchining aniq vazifasini bajaruvchi dastur. amaliy dastur - foydalanuvchining aniq vazifasini bajaruvchi dastur. Manba: CAM10/11, PDF 44-58; Qo'shimcha sahifa raqami (qaysi PDFga tegishli ekani aniqlanmagan; darslik sahifasi deb tasdiqlanmagan): 50
\end{tcolorbox}
\vspace{2mm}\hrule height 0.2pt\vspace{3mm}

% DB_ID bbf026fa-37df-c595-38ff-a338f961cbfe
\needspace{5\baselineskip}
{\bfseries\sffamily\color{ink} 503.} {scriptsizecolor{muted} Y1 | review | bbf026fa}
\par
Dasturiy ta'minotning tasnifi doirasida "translator" atamasining eng aniq ta'rifini toping.
\begin{enumerate}[label=\Alph*),leftmargin=9mm]
\item kompyuter resurslarini boshqarib, boshqa dasturlarga muhit yaratuvchi dastur
\item dastur kodini boshqa, odatda mashina bajaradigan shaklga o'giruvchi vosita
\item foydalanuvchining aniq vazifasini bajaruvchi dastur
\item tizimga xizmat ko'rsatuvchi yordamchi dastur
\end{enumerate}
\begin{tcolorbox}[enhanced,breakable,colback=teal!4,colframe=teal!60!black,boxrule=0.45pt,arc=1.2mm,left=2.7mm,right=2.7mm,top=1.8mm,bottom=1.8mm,before skip=2mm,after skip=2mm]
{\bfseries\sffamily\color{teal}To'g'ri javob:} \textbf{B) dastur kodini boshqa, odatda mashina bajaradigan shaklga o'giruvchi vosita}
\par\smallskip
{\bfseries Izoh:} dastur kodini boshqa, odatda mashina bajaradigan shaklga o'giruvchi vosita. translator - dastur kodini boshqa, odatda mashina bajaradigan shaklga o'giruvchi vosita. Manba: MAR, PDF 10-12; Qo'shimcha sahifa raqami (qaysi PDFga tegishli ekani aniqlanmagan; darslik sahifasi deb tasdiqlanmagan): 52
\end{tcolorbox}
\vspace{2mm}\hrule height 0.2pt\vspace{3mm}

% DB_ID c000f7ec-5b94-6cd5-48e4-5d89099ece91
\needspace{5\baselineskip}
{\bfseries\sffamily\color{ink} 504.} {scriptsizecolor{muted} Y1 | review | c000f7ec}
\par
Amaliy dasturlar va vazifaga mos tanlov doirasida "matn protsessori" atamasining eng aniq ta'rifini toping.
\begin{enumerate}[label=\Alph*),leftmargin=9mm]
\item formulalar asosida hisoblash, model va diagramma yaratish dasturi
\item tuzilmali ma'lumotlar bazasini yaratish, so'rash va boshqarish tizimi
\item tuzilmali matnli hujjat yaratish va formatlash dasturi
\item professional sahifa maketi va nashr tayyorlash dasturi
\end{enumerate}
\begin{tcolorbox}[enhanced,breakable,colback=teal!4,colframe=teal!60!black,boxrule=0.45pt,arc=1.2mm,left=2.7mm,right=2.7mm,top=1.8mm,bottom=1.8mm,before skip=2mm,after skip=2mm]
{\bfseries\sffamily\color{teal}To'g'ri javob:} \textbf{C) tuzilmali matnli hujjat yaratish va formatlash dasturi}
\par\smallskip
{\bfseries Izoh:} tuzilmali matnli hujjat yaratish va formatlash dasturi. matn protsessori - tuzilmali matnli hujjat yaratish va formatlash dasturi. Manba: CAM10/11, PDF 45-58; Qo'shimcha sahifa raqami (qaysi PDFga tegishli ekani aniqlanmagan; darslik sahifasi deb tasdiqlanmagan): 76
\end{tcolorbox}
\vspace{2mm}\hrule height 0.2pt\vspace{3mm}

% DB_ID ca5979d5-5ad7-0031-bd62-0d5c143389b7
\needspace{5\baselineskip}
{\bfseries\sffamily\color{ink} 505.} {scriptsizecolor{muted} Y1 | review | ca5979d5}
\par
Dastur o'rnatish, yangilash, litsenziya va moslik doirasida "freeware" atamasining eng aniq ta'rifini toping.
\begin{enumerate}[label=\Alph*),leftmargin=9mm]
\item dastur komponentlarini tizimda ishlashga tayyorlab joylashtirish
\item dastur va ro'yxatdan o'tgan komponentlarini nazoratli olib tashlash
\item mavjud versiyaga tuzatish yoki kichik yaxshilanish o'rnatish
\item odatda bepul foydalaniladigan, ammo kodi shart emas ochiq dastur
\end{enumerate}
\begin{tcolorbox}[enhanced,breakable,colback=teal!4,colframe=teal!60!black,boxrule=0.45pt,arc=1.2mm,left=2.7mm,right=2.7mm,top=1.8mm,bottom=1.8mm,before skip=2mm,after skip=2mm]
{\bfseries\sffamily\color{teal}To'g'ri javob:} \textbf{D) odatda bepul foydalaniladigan, ammo kodi shart emas ochiq dastur}
\par\smallskip
{\bfseries Izoh:} odatda bepul foydalaniladigan, ammo kodi shart emas ochiq dastur. freeware - odatda bepul foydalaniladigan, ammo kodi shart emas ochiq dastur. Manba: CAM8, PDF 99-110; Qo'shimcha sahifa raqami (qaysi PDFga tegishli ekani aniqlanmagan; darslik sahifasi deb tasdiqlanmagan): 107
\end{tcolorbox}
\vspace{2mm}\hrule height 0.2pt\vspace{3mm}

% DB_ID d6ac6c65-7ffd-1288-e599-2a8cbd2cd02d
\needspace{5\baselineskip}
{\bfseries\sffamily\color{ink} 506.} {scriptsizecolor{muted} Y1 | review | d6ac6c65}
\par
Drayver, utilita, firmware va translatorlar doirasida "utilita" atamasining eng aniq ta'rifini toping.
\begin{enumerate}[label=\Alph*),leftmargin=9mm]
\item OT bilan aniq qurilma modeli o'rtasida buyruqlarni moslashtiruvchi dastur
\item tizimga xizmat ko'rsatuvchi yordamchi dastur
\item qurilma ichidagi past darajali boshqaruv dasturi
\item butun manba kodini obyekt yoki bajariladigan kodga o'giruvchi translator
\end{enumerate}
\begin{tcolorbox}[enhanced,breakable,colback=teal!4,colframe=teal!60!black,boxrule=0.45pt,arc=1.2mm,left=2.7mm,right=2.7mm,top=1.8mm,bottom=1.8mm,before skip=2mm,after skip=2mm]
{\bfseries\sffamily\color{teal}To'g'ri javob:} \textbf{B) tizimga xizmat ko'rsatuvchi yordamchi dastur}
\par\smallskip
{\bfseries Izoh:} tizimga xizmat ko'rsatuvchi yordamchi dastur. utilita - tizimga xizmat ko'rsatuvchi yordamchi dastur. Manba: CAM10/11, PDF 44-49; Qo'shimcha sahifa raqami (qaysi PDFga tegishli ekani aniqlanmagan; darslik sahifasi deb tasdiqlanmagan): 70
\end{tcolorbox}
\vspace{2mm}\hrule height 0.2pt\vspace{3mm}

% DB_ID dc1b3117-44d6-d505-f23b-6429d633fc69
\needspace{5\baselineskip}
{\bfseries\sffamily\color{ink} 507.} {scriptsizecolor{muted} Y1 | review | dc1b3117}
\par
Dastur o'rnatish, yangilash, litsenziya va moslik doirasida "moslik" atamasining eng aniq ta'rifini toping.
\begin{enumerate}[label=\Alph*),leftmargin=9mm]
\item dastur komponentlarini tizimda ishlashga tayyorlab joylashtirish
\item qurilma, OT, dastur va formatlarning birgalikda ishlay olish xususiyati
\item dastur va ro'yxatdan o'tgan komponentlarini nazoratli olib tashlash
\item mavjud versiyaga tuzatish yoki kichik yaxshilanish o'rnatish
\end{enumerate}
\begin{tcolorbox}[enhanced,breakable,colback=teal!4,colframe=teal!60!black,boxrule=0.45pt,arc=1.2mm,left=2.7mm,right=2.7mm,top=1.8mm,bottom=1.8mm,before skip=2mm,after skip=2mm]
{\bfseries\sffamily\color{teal}To'g'ri javob:} \textbf{B) qurilma, OT, dastur va formatlarning birgalikda ishlay olish xususiyati}
\par\smallskip
{\bfseries Izoh:} qurilma, OT, dastur va formatlarning birgalikda ishlay olish xususiyati. moslik - qurilma, OT, dastur va formatlarning birgalikda ishlay olish xususiyati. Manba: CAM10/11, PDF 46-61; Qo'shimcha sahifa raqami (qaysi PDFga tegishli ekani aniqlanmagan; darslik sahifasi deb tasdiqlanmagan): 104
\end{tcolorbox}
\vspace{2mm}\hrule height 0.2pt\vspace{3mm}

% DB_ID ecb1a5be-bdad-f1d4-e899-0cbe3f04d9bc
\needspace{5\baselineskip}
{\bfseries\sffamily\color{ink} 508.} {scriptsizecolor{muted} Y1 | review | ecb1a5be}
\par
Drayver, utilita, firmware va translatorlar doirasida "firmware" atamasining eng aniq ta'rifini toping.
\begin{enumerate}[label=\Alph*),leftmargin=9mm]
\item OT bilan aniq qurilma modeli o'rtasida buyruqlarni moslashtiruvchi dastur
\item tizimga xizmat ko'rsatuvchi yordamchi dastur
\item qurilma ichidagi past darajali boshqaruv dasturi
\item butun manba kodini obyekt yoki bajariladigan kodga o'giruvchi translator
\end{enumerate}
\begin{tcolorbox}[enhanced,breakable,colback=teal!4,colframe=teal!60!black,boxrule=0.45pt,arc=1.2mm,left=2.7mm,right=2.7mm,top=1.8mm,bottom=1.8mm,before skip=2mm,after skip=2mm]
{\bfseries\sffamily\color{teal}To'g'ri javob:} \textbf{C) qurilma ichidagi past darajali boshqaruv dasturi}
\par\smallskip
{\bfseries Izoh:} qurilma ichidagi past darajali boshqaruv dasturi. firmware - qurilma ichidagi past darajali boshqaruv dasturi. Manba: MAR, PDF 10-16; Qo'shimcha sahifa raqami (qaysi PDFga tegishli ekani aniqlanmagan; darslik sahifasi deb tasdiqlanmagan): 71
\end{tcolorbox}
\vspace{2mm}\hrule height 0.2pt\vspace{3mm}

% DB_ID ef19e9f9-7180-7846-bb6f-b727ba54a2c6
\needspace{5\baselineskip}
{\bfseries\sffamily\color{ink} 509.} {scriptsizecolor{muted} Y1 | review | ef19e9f9}
\par
Dasturiy ta'minotning tasnifi doirasida "utilita" atamasining eng aniq ta'rifini toping.
\begin{enumerate}[label=\Alph*),leftmargin=9mm]
\item tizimga xizmat ko'rsatuvchi yordamchi dastur
\item kompyuter resurslarini boshqarib, boshqa dasturlarga muhit yaratuvchi dastur
\item foydalanuvchining aniq vazifasini bajaruvchi dastur
\item dastur kodini boshqa, odatda mashina bajaradigan shaklga o'giruvchi vosita
\end{enumerate}
\begin{tcolorbox}[enhanced,breakable,colback=teal!4,colframe=teal!60!black,boxrule=0.45pt,arc=1.2mm,left=2.7mm,right=2.7mm,top=1.8mm,bottom=1.8mm,before skip=2mm,after skip=2mm]
{\bfseries\sffamily\color{teal}To'g'ri javob:} \textbf{A) tizimga xizmat ko'rsatuvchi yordamchi dastur}
\par\smallskip
{\bfseries Izoh:} tizimga xizmat ko'rsatuvchi yordamchi dastur. utilita - tizimga xizmat ko'rsatuvchi yordamchi dastur. Manba: ICT5, PDF 38-41; Qo'shimcha sahifa raqami (qaysi PDFga tegishli ekani aniqlanmagan; darslik sahifasi deb tasdiqlanmagan): 51
\end{tcolorbox}
\vspace{2mm}\hrule height 0.2pt\vspace{3mm}

% DB_ID f5395b09-0694-9d86-f8e7-5adad86ee713
\needspace{5\baselineskip}
{\bfseries\sffamily\color{ink} 510.} {scriptsizecolor{muted} Y1 | review | f5395b09}
\par
Amaliy dasturlar va vazifaga mos tanlov doirasida "maxsus dastur" atamasining eng aniq ta'rifini toping.
\begin{enumerate}[label=\Alph*),leftmargin=9mm]
\item tuzilmali matnli hujjat yaratish va formatlash dasturi
\item formulalar asosida hisoblash, model va diagramma yaratish dasturi
\item tuzilmali ma'lumotlar bazasini yaratish, so'rash va boshqarish tizimi
\item muayyan tashkilot ehtiyojiga mos ishlab chiqilgan dastur
\end{enumerate}
\begin{tcolorbox}[enhanced,breakable,colback=teal!4,colframe=teal!60!black,boxrule=0.45pt,arc=1.2mm,left=2.7mm,right=2.7mm,top=1.8mm,bottom=1.8mm,before skip=2mm,after skip=2mm]
{\bfseries\sffamily\color{teal}To'g'ri javob:} \textbf{D) muayyan tashkilot ehtiyojiga mos ishlab chiqilgan dastur}
\par\smallskip
{\bfseries Izoh:} muayyan tashkilot ehtiyojiga mos ishlab chiqilgan dastur. maxsus dastur - muayyan tashkilot ehtiyojiga mos ishlab chiqilgan dastur. Manba: ICT8, PDF 95-110; Qo'shimcha sahifa raqami (qaysi PDFga tegishli ekani aniqlanmagan; darslik sahifasi deb tasdiqlanmagan): 81
\end{tcolorbox}
\vspace{2mm}\hrule height 0.2pt\vspace{3mm}

% DB_ID fe9c5ee1-8103-046f-8e9a-fc413307d81a
\needspace{5\baselineskip}
{\bfseries\sffamily\color{ink} 511.} {scriptsizecolor{muted} Y1 | review | fe9c5ee1}
\par
Amaliy dasturlar va vazifaga mos tanlov doirasida "DBMS" atamasining eng aniq ta'rifini toping.
\begin{enumerate}[label=\Alph*),leftmargin=9mm]
\item tuzilmali ma'lumotlar bazasini yaratish, so'rash va boshqarish tizimi
\item tuzilmali matnli hujjat yaratish va formatlash dasturi
\item formulalar asosida hisoblash, model va diagramma yaratish dasturi
\item professional sahifa maketi va nashr tayyorlash dasturi
\end{enumerate}
\begin{tcolorbox}[enhanced,breakable,colback=teal!4,colframe=teal!60!black,boxrule=0.45pt,arc=1.2mm,left=2.7mm,right=2.7mm,top=1.8mm,bottom=1.8mm,before skip=2mm,after skip=2mm]
{\bfseries\sffamily\color{teal}To'g'ri javob:} \textbf{A) tuzilmali ma'lumotlar bazasini yaratish, so'rash va boshqarish tizimi}
\par\smallskip
{\bfseries Izoh:} tuzilmali ma'lumotlar bazasini yaratish, so'rash va boshqarish tizimi. DBMS - tuzilmali ma'lumotlar bazasini yaratish, so'rash va boshqarish tizimi. Manba: CAM8, PDF 118-150; Qo'shimcha sahifa raqami (qaysi PDFga tegishli ekani aniqlanmagan; darslik sahifasi deb tasdiqlanmagan): 78
\end{tcolorbox}
\vspace{2mm}\hrule height 0.2pt\vspace{3mm}

% DB_ID 0ef0f329-a249-2ce4-7312-54618506b137
\needspace{5\baselineskip}
{\bfseries\sffamily\color{ink} 512.} {scriptsizecolor{muted} Y1 | review | 0ef0f329}
\par
Dastur o'rnatish, yangilash, litsenziya va moslik mavzusida "update" qatnashgan atama-ta'rif juftliklaridan to'g'risini toping.
\begin{enumerate}[label=\Alph*),leftmargin=9mm]
\item update - dastur komponentlarini tizimda ishlashga tayyorlab joylashtirish
\item update - dastur va ro'yxatdan o'tgan komponentlarini nazoratli olib tashlash
\item update - mavjud versiyaga tuzatish yoki kichik yaxshilanish o'rnatish
\item update - qurilma, OT, dastur va formatlarning birgalikda ishlay olish xususiyati
\end{enumerate}
\begin{tcolorbox}[enhanced,breakable,colback=teal!4,colframe=teal!60!black,boxrule=0.45pt,arc=1.2mm,left=2.7mm,right=2.7mm,top=1.8mm,bottom=1.8mm,before skip=2mm,after skip=2mm]
{\bfseries\sffamily\color{teal}To'g'ri javob:} \textbf{C) update - mavjud versiyaga tuzatish yoki kichik yaxshilanish o'rnatish}
\par\smallskip
{\bfseries Izoh:} update - mavjud versiyaga tuzatish yoki kichik yaxshilanish o'rnatish. To'g'ri juftlik: update - mavjud versiyaga tuzatish yoki kichik yaxshilanish o'rnatish. Manba: ICT8, PDF 106-110; Qo'shimcha sahifa raqami (qaysi PDFga tegishli ekani aniqlanmagan; darslik sahifasi deb tasdiqlanmagan): 104
\end{tcolorbox}
\vspace{2mm}\hrule height 0.2pt\vspace{3mm}

% DB_ID 0fc43941-609e-6e79-6b03-3722115e02a2
\needspace{5\baselineskip}
{\bfseries\sffamily\color{ink} 513.} {scriptsizecolor{muted} Y1 | review | 0fc43941}
\par
Drayver, utilita, firmware va translatorlar mavzusidagi vazifada tizimga xizmat ko'rsatuvchi yordamchi dastur talab qilindi. Qaysi tushuncha mos?
\begin{enumerate}[label=\Alph*),leftmargin=9mm]
\item drayver
\item firmware
\item utilita
\item kompilyator
\end{enumerate}
\begin{tcolorbox}[enhanced,breakable,colback=teal!4,colframe=teal!60!black,boxrule=0.45pt,arc=1.2mm,left=2.7mm,right=2.7mm,top=1.8mm,bottom=1.8mm,before skip=2mm,after skip=2mm]
{\bfseries\sffamily\color{teal}To'g'ri javob:} \textbf{C) utilita}
\par\smallskip
{\bfseries Izoh:} utilita. Berilgan vazifa utilita tushunchasiga mos: tizimga xizmat ko'rsatuvchi yordamchi dastur. Manba: CAM10/11, PDF 44-49; Qo'shimcha sahifa raqami (qaysi PDFga tegishli ekani aniqlanmagan; darslik sahifasi deb tasdiqlanmagan): 70
\end{tcolorbox}
\vspace{2mm}\hrule height 0.2pt\vspace{3mm}

% DB_ID 247fcac6-f11b-daa1-8bf1-f52dc9783fcf
\needspace{5\baselineskip}
{\bfseries\sffamily\color{ink} 514.} {scriptsizecolor{muted} Y1 | review | 247fcac6}
\par
Drayver, utilita, firmware va translatorlar mavzusida "drayver" qatnashgan atama-ta'rif juftliklaridan to'g'risini toping.
\begin{enumerate}[label=\Alph*),leftmargin=9mm]
\item drayver - tizimga xizmat ko'rsatuvchi yordamchi dastur
\item drayver - qurilma ichidagi past darajali boshqaruv dasturi
\item drayver - OT bilan aniq qurilma modeli o'rtasida buyruqlarni moslashtiruvchi dastur
\item drayver - butun manba kodini obyekt yoki bajariladigan kodga o'giruvchi translator
\end{enumerate}
\begin{tcolorbox}[enhanced,breakable,colback=teal!4,colframe=teal!60!black,boxrule=0.45pt,arc=1.2mm,left=2.7mm,right=2.7mm,top=1.8mm,bottom=1.8mm,before skip=2mm,after skip=2mm]
{\bfseries\sffamily\color{teal}To'g'ri javob:} \textbf{C) drayver - OT bilan aniq qurilma modeli o'rtasida buyruqlarni moslashtiruvchi dastur}
\par\smallskip
{\bfseries Izoh:} drayver - OT bilan aniq qurilma modeli o'rtasida buyruqlarni moslashtiruvchi dastur. To'g'ri juftlik: drayver - OT bilan aniq qurilma modeli o'rtasida buyruqlarni moslashtiruvchi dastur. Manba: CAM9, PDF 14-15; Qo'shimcha sahifa raqami (qaysi PDFga tegishli ekani aniqlanmagan; darslik sahifasi deb tasdiqlanmagan): 69
\end{tcolorbox}
\vspace{2mm}\hrule height 0.2pt\vspace{3mm}

% DB_ID 312afb65-eb4e-7681-d5df-b87469a509a1
\needspace{5\baselineskip}
{\bfseries\sffamily\color{ink} 515.} {scriptsizecolor{muted} Y1 | review | 312afb65}
\par
Dasturiy ta'minotning tasnifi mavzusidagi vazifada kompyuter resurslarini boshqarib, boshqa dasturlarga muhit yaratuvchi dastur talab qilindi. Qaysi tushuncha mos?
\begin{enumerate}[label=\Alph*),leftmargin=9mm]
\item amaliy dastur
\item utilita
\item translator
\item tizimli dastur
\end{enumerate}
\begin{tcolorbox}[enhanced,breakable,colback=teal!4,colframe=teal!60!black,boxrule=0.45pt,arc=1.2mm,left=2.7mm,right=2.7mm,top=1.8mm,bottom=1.8mm,before skip=2mm,after skip=2mm]
{\bfseries\sffamily\color{teal}To'g'ri javob:} \textbf{D) tizimli dastur}
\par\smallskip
{\bfseries Izoh:} tizimli dastur. Berilgan vazifa tizimli dastur tushunchasiga mos: kompyuter resurslarini boshqarib, boshqa dasturlarga muhit yaratuvchi dastur. Manba: CAM9, PDF 11-15; Qo'shimcha sahifa raqami (qaysi PDFga tegishli ekani aniqlanmagan; darslik sahifasi deb tasdiqlanmagan): 49
\end{tcolorbox}
\vspace{2mm}\hrule height 0.2pt\vspace{3mm}

% DB_ID 3354fcc0-68fb-6b3e-94c0-925c80313de5
\needspace{5\baselineskip}
{\bfseries\sffamily\color{ink} 516.} {scriptsizecolor{muted} Y1 | review | 3354fcc0}
\par
Dastur o'rnatish, yangilash, litsenziya va moslik mavzusida "open-source" qatnashgan atama-ta'rif juftliklaridan to'g'risini toping.
\begin{enumerate}[label=\Alph*),leftmargin=9mm]
\item open-source - manba kodi litsenziya asosida ochiq dastur
\item open-source - dastur komponentlarini tizimda ishlashga tayyorlab joylashtirish
\item open-source - dastur va ro'yxatdan o'tgan komponentlarini nazoratli olib tashlash
\item open-source - mavjud versiyaga tuzatish yoki kichik yaxshilanish o'rnatish
\end{enumerate}
\begin{tcolorbox}[enhanced,breakable,colback=teal!4,colframe=teal!60!black,boxrule=0.45pt,arc=1.2mm,left=2.7mm,right=2.7mm,top=1.8mm,bottom=1.8mm,before skip=2mm,after skip=2mm]
{\bfseries\sffamily\color{teal}To'g'ri javob:} \textbf{A) open-source - manba kodi litsenziya asosida ochiq dastur}
\par\smallskip
{\bfseries Izoh:} open-source - manba kodi litsenziya asosida ochiq dastur. To'g'ri juftlik: open-source - manba kodi litsenziya asosida ochiq dastur. Manba: MAR, PDF 10-16; Qo'shimcha sahifa raqami (qaysi PDFga tegishli ekani aniqlanmagan; darslik sahifasi deb tasdiqlanmagan): 106
\end{tcolorbox}
\vspace{2mm}\hrule height 0.2pt\vspace{3mm}

% DB_ID 3f7097b8-5b83-7b5d-7562-5843f8566a8d
\needspace{5\baselineskip}
{\bfseries\sffamily\color{ink} 517.} {scriptsizecolor{muted} Y1 | review | 3f7097b8}
\par
Drayver, utilita, firmware va translatorlar mavzusidagi vazifada kodni bajarish paytida bosqichma-bosqich tarjima qiluvchi translator talab qilindi. Qaysi tushuncha mos?
\begin{enumerate}[label=\Alph*),leftmargin=9mm]
\item drayver
\item interpretator
\item utilita
\item firmware
\end{enumerate}
\begin{tcolorbox}[enhanced,breakable,colback=teal!4,colframe=teal!60!black,boxrule=0.45pt,arc=1.2mm,left=2.7mm,right=2.7mm,top=1.8mm,bottom=1.8mm,before skip=2mm,after skip=2mm]
{\bfseries\sffamily\color{teal}To'g'ri javob:} \textbf{B) interpretator}
\par\smallskip
{\bfseries Izoh:} interpretator. Berilgan vazifa interpretator tushunchasiga mos: kodni bajarish paytida bosqichma-bosqich tarjima qiluvchi translator. Manba: CAM9, PDF 14-15; Qo'shimcha sahifa raqami (qaysi PDFga tegishli ekani aniqlanmagan; darslik sahifasi deb tasdiqlanmagan): 73
\end{tcolorbox}
\vspace{2mm}\hrule height 0.2pt\vspace{3mm}

% DB_ID 40f600ed-2532-fb61-f40e-5022deaf1c8a
\needspace{5\baselineskip}
{\bfseries\sffamily\color{ink} 518.} {scriptsizecolor{muted} Y1 | review | 40f600ed}
\par
Dasturiy ta'minotning tasnifi mavzusidagi vazifada muayyan buyurtmachi yoki jarayon uchun yaratilgan dastur talab qilindi. Qaysi tushuncha mos?
\begin{enumerate}[label=\Alph*),leftmargin=9mm]
\item maxsus dastur
\item tizimli dastur
\item amaliy dastur
\item utilita
\end{enumerate}
\begin{tcolorbox}[enhanced,breakable,colback=teal!4,colframe=teal!60!black,boxrule=0.45pt,arc=1.2mm,left=2.7mm,right=2.7mm,top=1.8mm,bottom=1.8mm,before skip=2mm,after skip=2mm]
{\bfseries\sffamily\color{teal}To'g'ri javob:} \textbf{A) maxsus dastur}
\par\smallskip
{\bfseries Izoh:} maxsus dastur. Berilgan vazifa maxsus dastur tushunchasiga mos: muayyan buyurtmachi yoki jarayon uchun yaratilgan dastur. Manba: CAM10/11, PDF 44-58; Qo'shimcha sahifa raqami (qaysi PDFga tegishli ekani aniqlanmagan; darslik sahifasi deb tasdiqlanmagan): 54
\end{tcolorbox}
\vspace{2mm}\hrule height 0.2pt\vspace{3mm}

% DB_ID 5bb3252e-fe25-c8d6-17b5-aa41495f78ea
\needspace{5\baselineskip}
{\bfseries\sffamily\color{ink} 519.} {scriptsizecolor{muted} Y1 | review | 5bb3252e}
\par
Dasturiy ta'minotning tasnifi mavzusidagi vazifada foydalanuvchining aniq vazifasini bajaruvchi dastur talab qilindi. Qaysi tushuncha mos?
\begin{enumerate}[label=\Alph*),leftmargin=9mm]
\item amaliy dastur
\item tizimli dastur
\item utilita
\item translator
\end{enumerate}
\begin{tcolorbox}[enhanced,breakable,colback=teal!4,colframe=teal!60!black,boxrule=0.45pt,arc=1.2mm,left=2.7mm,right=2.7mm,top=1.8mm,bottom=1.8mm,before skip=2mm,after skip=2mm]
{\bfseries\sffamily\color{teal}To'g'ri javob:} \textbf{A) amaliy dastur}
\par\smallskip
{\bfseries Izoh:} amaliy dastur. Berilgan vazifa amaliy dastur tushunchasiga mos: foydalanuvchining aniq vazifasini bajaruvchi dastur. Manba: CAM10/11, PDF 44-58; Qo'shimcha sahifa raqami (qaysi PDFga tegishli ekani aniqlanmagan; darslik sahifasi deb tasdiqlanmagan): 50
\end{tcolorbox}
\vspace{2mm}\hrule height 0.2pt\vspace{3mm}

% DB_ID 6e4fe344-3bd5-6d49-560b-5212ab851999
\needspace{5\baselineskip}
{\bfseries\sffamily\color{ink} 520.} {scriptsizecolor{muted} Y1 | review | 6e4fe344}
\par
Dasturiy ta'minotning tasnifi mavzusida "umumiy maqsadli dastur" qatnashgan atama-ta'rif juftliklaridan to'g'risini toping.
\begin{enumerate}[label=\Alph*),leftmargin=9mm]
\item umumiy maqsadli dastur - ko'p foydalanuvchi va vazifalarga mo'ljallangan tayyor dastur
\item umumiy maqsadli dastur - kompyuter resurslarini boshqarib, boshqa dasturlarga muhit yaratuvchi dastur
\item umumiy maqsadli dastur - foydalanuvchining aniq vazifasini bajaruvchi dastur
\item umumiy maqsadli dastur - tizimga xizmat ko'rsatuvchi yordamchi dastur
\end{enumerate}
\begin{tcolorbox}[enhanced,breakable,colback=teal!4,colframe=teal!60!black,boxrule=0.45pt,arc=1.2mm,left=2.7mm,right=2.7mm,top=1.8mm,bottom=1.8mm,before skip=2mm,after skip=2mm]
{\bfseries\sffamily\color{teal}To'g'ri javob:} \textbf{A) umumiy maqsadli dastur - ko'p foydalanuvchi va vazifalarga mo'ljallangan tayyor dastur}
\par\smallskip
{\bfseries Izoh:} umumiy maqsadli dastur - ko'p foydalanuvchi va vazifalarga mo'ljallangan tayyor dastur. To'g'ri juftlik: umumiy maqsadli dastur - ko'p foydalanuvchi va vazifalarga mo'ljallangan tayyor dastur. Manba: CAM9, PDF 11-15; Qo'shimcha sahifa raqami (qaysi PDFga tegishli ekani aniqlanmagan; darslik sahifasi deb tasdiqlanmagan): 53
\end{tcolorbox}
\vspace{2mm}\hrule height 0.2pt\vspace{3mm}

% DB_ID 72b4b97f-923d-0b41-84a2-951eac49f31f
\needspace{5\baselineskip}
{\bfseries\sffamily\color{ink} 521.} {scriptsizecolor{muted} Y1 | review | 72b4b97f}
\par
Drayver, utilita, firmware va translatorlar mavzusida "firmware" qatnashgan atama-ta'rif juftliklaridan to'g'risini toping.
\begin{enumerate}[label=\Alph*),leftmargin=9mm]
\item firmware - qurilma ichidagi past darajali boshqaruv dasturi
\item firmware - OT bilan aniq qurilma modeli o'rtasida buyruqlarni moslashtiruvchi dastur
\item firmware - tizimga xizmat ko'rsatuvchi yordamchi dastur
\item firmware - butun manba kodini obyekt yoki bajariladigan kodga o'giruvchi translator
\end{enumerate}
\begin{tcolorbox}[enhanced,breakable,colback=teal!4,colframe=teal!60!black,boxrule=0.45pt,arc=1.2mm,left=2.7mm,right=2.7mm,top=1.8mm,bottom=1.8mm,before skip=2mm,after skip=2mm]
{\bfseries\sffamily\color{teal}To'g'ri javob:} \textbf{A) firmware - qurilma ichidagi past darajali boshqaruv dasturi}
\par\smallskip
{\bfseries Izoh:} firmware - qurilma ichidagi past darajali boshqaruv dasturi. To'g'ri juftlik: firmware - qurilma ichidagi past darajali boshqaruv dasturi. Manba: MAR, PDF 10-16; Qo'shimcha sahifa raqami (qaysi PDFga tegishli ekani aniqlanmagan; darslik sahifasi deb tasdiqlanmagan): 71
\end{tcolorbox}
\vspace{2mm}\hrule height 0.2pt\vspace{3mm}

% DB_ID 7a5a9e9f-bdea-2a44-14d1-ea73faed8b56
\needspace{5\baselineskip}
{\bfseries\sffamily\color{ink} 522.} {scriptsizecolor{muted} Y1 | review | 7a5a9e9f}
\par
Dastur o'rnatish, yangilash, litsenziya va moslik mavzusidagi vazifada odatda bepul foydalaniladigan, ammo kodi shart emas ochiq dastur talab qilindi. Qaysi tushuncha mos?
\begin{enumerate}[label=\Alph*),leftmargin=9mm]
\item freeware
\item installation
\item uninstall
\item update
\end{enumerate}
\begin{tcolorbox}[enhanced,breakable,colback=teal!4,colframe=teal!60!black,boxrule=0.45pt,arc=1.2mm,left=2.7mm,right=2.7mm,top=1.8mm,bottom=1.8mm,before skip=2mm,after skip=2mm]
{\bfseries\sffamily\color{teal}To'g'ri javob:} \textbf{A) freeware}
\par\smallskip
{\bfseries Izoh:} freeware. Berilgan vazifa freeware tushunchasiga mos: odatda bepul foydalaniladigan, ammo kodi shart emas ochiq dastur. Manba: CAM8, PDF 99-110; Qo'shimcha sahifa raqami (qaysi PDFga tegishli ekani aniqlanmagan; darslik sahifasi deb tasdiqlanmagan): 107
\end{tcolorbox}
\vspace{2mm}\hrule height 0.2pt\vspace{3mm}

% DB_ID 7e03d2ed-afcc-4eef-75d1-4e2f885ae6b3
\needspace{5\baselineskip}
{\bfseries\sffamily\color{ink} 523.} {scriptsizecolor{muted} Y1 | review | 7e03d2ed}
\par
Amaliy dasturlar va vazifaga mos tanlov mavzusida "CAD" qatnashgan atama-ta'rif juftliklaridan to'g'risini toping.
\begin{enumerate}[label=\Alph*),leftmargin=9mm]
\item CAD - aniq texnik chizma va model yaratish dasturi
\item CAD - tuzilmali matnli hujjat yaratish va formatlash dasturi
\item CAD - formulalar asosida hisoblash, model va diagramma yaratish dasturi
\item CAD - tuzilmali ma'lumotlar bazasini yaratish, so'rash va boshqarish tizimi
\end{enumerate}
\begin{tcolorbox}[enhanced,breakable,colback=teal!4,colframe=teal!60!black,boxrule=0.45pt,arc=1.2mm,left=2.7mm,right=2.7mm,top=1.8mm,bottom=1.8mm,before skip=2mm,after skip=2mm]
{\bfseries\sffamily\color{teal}To'g'ri javob:} \textbf{A) CAD - aniq texnik chizma va model yaratish dasturi}
\par\smallskip
{\bfseries Izoh:} CAD - aniq texnik chizma va model yaratish dasturi. To'g'ri juftlik: CAD - aniq texnik chizma va model yaratish dasturi. Manba: ICT6, PDF 62-126; Qo'shimcha sahifa raqami (qaysi PDFga tegishli ekani aniqlanmagan; darslik sahifasi deb tasdiqlanmagan): 81
\end{tcolorbox}
\vspace{2mm}\hrule height 0.2pt\vspace{3mm}

% DB_ID 9464da0a-7190-0e16-d0d9-508ee161c956
\needspace{5\baselineskip}
{\bfseries\sffamily\color{ink} 524.} {scriptsizecolor{muted} Y1 | review | 9464da0a}
\par
Amaliy dasturlar va vazifaga mos tanlov mavzusidagi vazifada aniq texnik chizma va model yaratish dasturi talab qilindi. Qaysi tushuncha mos?
\begin{enumerate}[label=\Alph*),leftmargin=9mm]
\item matn protsessori
\item elektron jadval
\item DBMS
\item CAD
\end{enumerate}
\begin{tcolorbox}[enhanced,breakable,colback=teal!4,colframe=teal!60!black,boxrule=0.45pt,arc=1.2mm,left=2.7mm,right=2.7mm,top=1.8mm,bottom=1.8mm,before skip=2mm,after skip=2mm]
{\bfseries\sffamily\color{teal}To'g'ri javob:} \textbf{D) CAD}
\par\smallskip
{\bfseries Izoh:} CAD. Berilgan vazifa CAD tushunchasiga mos: aniq texnik chizma va model yaratish dasturi. Manba: ICT6, PDF 62-126; Qo'shimcha sahifa raqami (qaysi PDFga tegishli ekani aniqlanmagan; darslik sahifasi deb tasdiqlanmagan): 80
\end{tcolorbox}
\vspace{2mm}\hrule height 0.2pt\vspace{3mm}

% DB_ID 9c98e11b-f292-c174-fa97-c0cdecec8cd2
\needspace{5\baselineskip}
{\bfseries\sffamily\color{ink} 525.} {scriptsizecolor{muted} Y1 | review | 9c98e11b}
\par
Dastur o'rnatish, yangilash, litsenziya va moslik mavzusidagi vazifada dastur va ro'yxatdan o'tgan komponentlarini nazoratli olib tashlash talab qilindi. Qaysi tushuncha mos?
\begin{enumerate}[label=\Alph*),leftmargin=9mm]
\item uninstall
\item installation
\item update
\item moslik
\end{enumerate}
\begin{tcolorbox}[enhanced,breakable,colback=teal!4,colframe=teal!60!black,boxrule=0.45pt,arc=1.2mm,left=2.7mm,right=2.7mm,top=1.8mm,bottom=1.8mm,before skip=2mm,after skip=2mm]
{\bfseries\sffamily\color{teal}To'g'ri javob:} \textbf{A) uninstall}
\par\smallskip
{\bfseries Izoh:} uninstall. Berilgan vazifa uninstall tushunchasiga mos: dastur va ro'yxatdan o'tgan komponentlarini nazoratli olib tashlash. Manba: CAM8, PDF 99-110; Qo'shimcha sahifa raqami (qaysi PDFga tegishli ekani aniqlanmagan; darslik sahifasi deb tasdiqlanmagan): 102
\end{tcolorbox}
\vspace{2mm}\hrule height 0.2pt\vspace{3mm}

% DB_ID a07be903-a74e-d9b8-b5fd-271a5b39a3ad
\needspace{5\baselineskip}
{\bfseries\sffamily\color{ink} 526.} {scriptsizecolor{muted} Y1 | review | a07be903}
\par
Amaliy dasturlar va vazifaga mos tanlov mavzusidagi vazifada tuzilmali ma'lumotlar bazasini yaratish, so'rash va boshqarish tizimi talab qilindi. Qaysi tushuncha mos?
\begin{enumerate}[label=\Alph*),leftmargin=9mm]
\item matn protsessori
\item DBMS
\item elektron jadval
\item DTP
\end{enumerate}
\begin{tcolorbox}[enhanced,breakable,colback=teal!4,colframe=teal!60!black,boxrule=0.45pt,arc=1.2mm,left=2.7mm,right=2.7mm,top=1.8mm,bottom=1.8mm,before skip=2mm,after skip=2mm]
{\bfseries\sffamily\color{teal}To'g'ri javob:} \textbf{B) DBMS}
\par\smallskip
{\bfseries Izoh:} DBMS. Berilgan vazifa DBMS tushunchasiga mos: tuzilmali ma'lumotlar bazasini yaratish, so'rash va boshqarish tizimi. Manba: CAM8, PDF 118-150; Qo'shimcha sahifa raqami (qaysi PDFga tegishli ekani aniqlanmagan; darslik sahifasi deb tasdiqlanmagan): 78
\end{tcolorbox}
\vspace{2mm}\hrule height 0.2pt\vspace{3mm}

% DB_ID a39074e9-68c2-f5d6-3f1e-462d0a97baec
\needspace{5\baselineskip}
{\bfseries\sffamily\color{ink} 527.} {scriptsizecolor{muted} Y1 | review | a39074e9}
\par
Drayver, utilita, firmware va translatorlar mavzusidagi vazifada obyekt modullari va kutubxonalarni bajariladigan faylga bog'lovchi vosita talab qilindi. Qaysi tushuncha mos?
\begin{enumerate}[label=\Alph*),leftmargin=9mm]
\item drayver
\item utilita
\item linker
\item firmware
\end{enumerate}
\begin{tcolorbox}[enhanced,breakable,colback=teal!4,colframe=teal!60!black,boxrule=0.45pt,arc=1.2mm,left=2.7mm,right=2.7mm,top=1.8mm,bottom=1.8mm,before skip=2mm,after skip=2mm]
{\bfseries\sffamily\color{teal}To'g'ri javob:} \textbf{C) linker}
\par\smallskip
{\bfseries Izoh:} linker. Berilgan vazifa linker tushunchasiga mos: obyekt modullari va kutubxonalarni bajariladigan faylga bog'lovchi vosita. Manba: CAM10/11, PDF 44-49; Qo'shimcha sahifa raqami (qaysi PDFga tegishli ekani aniqlanmagan; darslik sahifasi deb tasdiqlanmagan): 75
\end{tcolorbox}
\vspace{2mm}\hrule height 0.2pt\vspace{3mm}

% DB_ID a71c81e5-6edf-248b-0b94-1c1184f224a2
\needspace{5\baselineskip}
{\bfseries\sffamily\color{ink} 528.} {scriptsizecolor{muted} Y1 | review | a71c81e5}
\par
Amaliy dasturlar va vazifaga mos tanlov mavzusida "DBMS" qatnashgan atama-ta'rif juftliklaridan to'g'risini toping.
\begin{enumerate}[label=\Alph*),leftmargin=9mm]
\item DBMS - tuzilmali matnli hujjat yaratish va formatlash dasturi
\item DBMS - formulalar asosida hisoblash, model va diagramma yaratish dasturi
\item DBMS - tuzilmali ma'lumotlar bazasini yaratish, so'rash va boshqarish tizimi
\item DBMS - professional sahifa maketi va nashr tayyorlash dasturi
\end{enumerate}
\begin{tcolorbox}[enhanced,breakable,colback=teal!4,colframe=teal!60!black,boxrule=0.45pt,arc=1.2mm,left=2.7mm,right=2.7mm,top=1.8mm,bottom=1.8mm,before skip=2mm,after skip=2mm]
{\bfseries\sffamily\color{teal}To'g'ri javob:} \textbf{C) DBMS - tuzilmali ma'lumotlar bazasini yaratish, so'rash va boshqarish tizimi}
\par\smallskip
{\bfseries Izoh:} DBMS - tuzilmali ma'lumotlar bazasini yaratish, so'rash va boshqarish tizimi. To'g'ri juftlik: DBMS - tuzilmali ma'lumotlar bazasini yaratish, so'rash va boshqarish tizimi. Manba: CAM8, PDF 118-150; Qo'shimcha sahifa raqami (qaysi PDFga tegishli ekani aniqlanmagan; darslik sahifasi deb tasdiqlanmagan): 79
\end{tcolorbox}
\vspace{2mm}\hrule height 0.2pt\vspace{3mm}

% DB_ID aa51994f-92e7-117c-a191-a0b3c320f7ee
\needspace{5\baselineskip}
{\bfseries\sffamily\color{ink} 529.} {scriptsizecolor{muted} Y1 | review | aa51994f}
\par
Dastur o'rnatish, yangilash, litsenziya va moslik mavzusidagi vazifada mavjud versiyaga tuzatish yoki kichik yaxshilanish o'rnatish talab qilindi. Qaysi tushuncha mos?
\begin{enumerate}[label=\Alph*),leftmargin=9mm]
\item installation
\item update
\item uninstall
\item moslik
\end{enumerate}
\begin{tcolorbox}[enhanced,breakable,colback=teal!4,colframe=teal!60!black,boxrule=0.45pt,arc=1.2mm,left=2.7mm,right=2.7mm,top=1.8mm,bottom=1.8mm,before skip=2mm,after skip=2mm]
{\bfseries\sffamily\color{teal}To'g'ri javob:} \textbf{B) update}
\par\smallskip
{\bfseries Izoh:} update. Berilgan vazifa update tushunchasiga mos: mavjud versiyaga tuzatish yoki kichik yaxshilanish o'rnatish. Manba: ICT8, PDF 106-110; Qo'shimcha sahifa raqami (qaysi PDFga tegishli ekani aniqlanmagan; darslik sahifasi deb tasdiqlanmagan): 103
\end{tcolorbox}
\vspace{2mm}\hrule height 0.2pt\vspace{3mm}

% DB_ID aad989e8-e1e1-1746-e5d0-8886bb484e2b
\needspace{5\baselineskip}
{\bfseries\sffamily\color{ink} 530.} {scriptsizecolor{muted} Y1 | review | aad989e8}
\par
Drayver, utilita, firmware va translatorlar mavzusida "interpretator" qatnashgan atama-ta'rif juftliklaridan to'g'risini toping.
\begin{enumerate}[label=\Alph*),leftmargin=9mm]
\item interpretator - OT bilan aniq qurilma modeli o'rtasida buyruqlarni moslashtiruvchi dastur
\item interpretator - tizimga xizmat ko'rsatuvchi yordamchi dastur
\item interpretator - kodni bajarish paytida bosqichma-bosqich tarjima qiluvchi translator
\item interpretator - qurilma ichidagi past darajali boshqaruv dasturi
\end{enumerate}
\begin{tcolorbox}[enhanced,breakable,colback=teal!4,colframe=teal!60!black,boxrule=0.45pt,arc=1.2mm,left=2.7mm,right=2.7mm,top=1.8mm,bottom=1.8mm,before skip=2mm,after skip=2mm]
{\bfseries\sffamily\color{teal}To'g'ri javob:} \textbf{C) interpretator - kodni bajarish paytida bosqichma-bosqich tarjima qiluvchi translator}
\par\smallskip
{\bfseries Izoh:} interpretator - kodni bajarish paytida bosqichma-bosqich tarjima qiluvchi translator. To'g'ri juftlik: interpretator - kodni bajarish paytida bosqichmabosqich tarjima qiluvchi translator. Manba: CAM9, PDF 14-15; Qo'shimcha sahifa raqami (qaysi PDFga tegishli ekani aniqlanmagan; darslik sahifasi deb tasdiqlanmagan): 74
\end{tcolorbox}
\vspace{2mm}\hrule height 0.2pt\vspace{3mm}

% DB_ID af88a919-ceef-6586-07f0-6721bdd33074
\needspace{5\baselineskip}
{\bfseries\sffamily\color{ink} 531.} {scriptsizecolor{muted} Y1 | review | af88a919}
\par
Amaliy dasturlar va vazifaga mos tanlov mavzusidagi vazifada professional sahifa maketi va nashr tayyorlash dasturi talab qilindi. Qaysi tushuncha mos?
\begin{enumerate}[label=\Alph*),leftmargin=9mm]
\item matn protsessori
\item elektron jadval
\item DTP
\item DBMS
\end{enumerate}
\begin{tcolorbox}[enhanced,breakable,colback=teal!4,colframe=teal!60!black,boxrule=0.45pt,arc=1.2mm,left=2.7mm,right=2.7mm,top=1.8mm,bottom=1.8mm,before skip=2mm,after skip=2mm]
{\bfseries\sffamily\color{teal}To'g'ri javob:} \textbf{C) DTP}
\par\smallskip
{\bfseries Izoh:} DTP. Berilgan vazifa DTP tushunchasiga mos: professional sahifa maketi va nashr tayyorlash dasturi. Manba: CAM5, PDF 45-138; Qo'shimcha sahifa raqami (qaysi PDFga tegishli ekani aniqlanmagan; darslik sahifasi deb tasdiqlanmagan): 79
\end{tcolorbox}
\vspace{2mm}\hrule height 0.2pt\vspace{3mm}

% DB_ID b00d1f2a-7fd5-007e-3c6a-cc145a6f4e47
\needspace{5\baselineskip}
{\bfseries\sffamily\color{ink} 532.} {scriptsizecolor{muted} Y1 | review | b00d1f2a}
\par
Amaliy dasturlar va vazifaga mos tanlov mavzusida "matn protsessori" qatnashgan atama-ta'rif juftliklaridan to'g'risini toping.
\begin{enumerate}[label=\Alph*),leftmargin=9mm]
\item matn protsessori - tuzilmali matnli hujjat yaratish va formatlash dasturi
\item matn protsessori - formulalar asosida hisoblash, model va diagramma yaratish dasturi
\item matn protsessori - tuzilmali ma'lumotlar bazasini yaratish, so'rash va boshqarish tizimi
\item matn protsessori - professional sahifa maketi va nashr tayyorlash dasturi
\end{enumerate}
\begin{tcolorbox}[enhanced,breakable,colback=teal!4,colframe=teal!60!black,boxrule=0.45pt,arc=1.2mm,left=2.7mm,right=2.7mm,top=1.8mm,bottom=1.8mm,before skip=2mm,after skip=2mm]
{\bfseries\sffamily\color{teal}To'g'ri javob:} \textbf{A) matn protsessori - tuzilmali matnli hujjat yaratish va formatlash dasturi}
\par\smallskip
{\bfseries Izoh:} matn protsessori - tuzilmali matnli hujjat yaratish va formatlash dasturi. To'g'ri juftlik: matn protsessori - tuzilmali matnli hujjat yaratish va formatlash dasturi. Manba: CAM10/11, PDF 45-58; Qo'shimcha sahifa raqami (qaysi PDFga tegishli ekani aniqlanmagan; darslik sahifasi deb tasdiqlanmagan): 76
\end{tcolorbox}
\vspace{2mm}\hrule height 0.2pt\vspace{3mm}

% DB_ID b9795126-5040-1fea-ddaa-079a73a74aae
\needspace{5\baselineskip}
{\bfseries\sffamily\color{ink} 533.} {scriptsizecolor{muted} Y1 | review | b9795126}
\par
Drayver, utilita, firmware va translatorlar mavzusidagi vazifada OT bilan aniq qurilma modeli o'rtasida buyruqlarni moslashtiruvchi dastur talab qilindi. Qaysi tushuncha mos?
\begin{enumerate}[label=\Alph*),leftmargin=9mm]
\item utilita
\item drayver
\item firmware
\item kompilyator
\end{enumerate}
\begin{tcolorbox}[enhanced,breakable,colback=teal!4,colframe=teal!60!black,boxrule=0.45pt,arc=1.2mm,left=2.7mm,right=2.7mm,top=1.8mm,bottom=1.8mm,before skip=2mm,after skip=2mm]
{\bfseries\sffamily\color{teal}To'g'ri javob:} \textbf{B) drayver}
\par\smallskip
{\bfseries Izoh:} drayver. Berilgan vazifa drayver tushunchasiga mos: OT bilan aniq qurilma modeli o'rtasida buyruqlarni moslashtiruvchi dastur. Manba: CAM9, PDF 14-15; Qo'shimcha sahifa raqami (qaysi PDFga tegishli ekani aniqlanmagan; darslik sahifasi deb tasdiqlanmagan): 69
\end{tcolorbox}
\vspace{2mm}\hrule height 0.2pt\vspace{3mm}

% DB_ID baaea13e-b334-580a-bd74-64e73c9a6e80
\needspace{5\baselineskip}
{\bfseries\sffamily\color{ink} 534.} {scriptsizecolor{muted} Y1 | review | baaea13e}
\par
Dastur o'rnatish, yangilash, litsenziya va moslik mavzusidagi vazifada qurilma, OT, dastur va formatlarning birgalikda ishlay olish xususiyati talab qilindi. Qaysi tushuncha mos?
\begin{enumerate}[label=\Alph*),leftmargin=9mm]
\item installation
\item uninstall
\item moslik
\item update
\end{enumerate}
\begin{tcolorbox}[enhanced,breakable,colback=teal!4,colframe=teal!60!black,boxrule=0.45pt,arc=1.2mm,left=2.7mm,right=2.7mm,top=1.8mm,bottom=1.8mm,before skip=2mm,after skip=2mm]
{\bfseries\sffamily\color{teal}To'g'ri javob:} \textbf{C) moslik}
\par\smallskip
{\bfseries Izoh:} moslik. Berilgan vazifa moslik tushunchasiga mos: qurilma, OT, dastur va formatlarning birgalikda ishlay olish xususiyati. Manba: CAM10/11, PDF 46-61; Qo'shimcha sahifa raqami (qaysi PDFga tegishli ekani aniqlanmagan; darslik sahifasi deb tasdiqlanmagan): 105
\end{tcolorbox}
\vspace{2mm}\hrule height 0.2pt\vspace{3mm}

% DB_ID bf8bc2f8-6f9b-b772-70c6-fd783628a6f3
\needspace{5\baselineskip}
{\bfseries\sffamily\color{ink} 535.} {scriptsizecolor{muted} Y1 | review | bf8bc2f8}
\par
Dastur o'rnatish, yangilash, litsenziya va moslik mavzusida "installation" qatnashgan atama-ta'rif juftliklaridan to'g'risini toping.
\begin{enumerate}[label=\Alph*),leftmargin=9mm]
\item installation - dastur komponentlarini tizimda ishlashga tayyorlab joylashtirish
\item installation - dastur va ro'yxatdan o'tgan komponentlarini nazoratli olib tashlash
\item installation - mavjud versiyaga tuzatish yoki kichik yaxshilanish o'rnatish
\item installation - qurilma, OT, dastur va formatlarning birgalikda ishlay olish xususiyati
\end{enumerate}
\begin{tcolorbox}[enhanced,breakable,colback=teal!4,colframe=teal!60!black,boxrule=0.45pt,arc=1.2mm,left=2.7mm,right=2.7mm,top=1.8mm,bottom=1.8mm,before skip=2mm,after skip=2mm]
{\bfseries\sffamily\color{teal}To'g'ri javob:} \textbf{A) installation - dastur komponentlarini tizimda ishlashga tayyorlab joylashtirish}
\par\smallskip
{\bfseries Izoh:} installation - dastur komponentlarini tizimda ishlashga tayyorlab joylashtirish. To'g'ri juftlik: installation - dastur komponentlarini tizimda ishlashga tayyorlab joylashtirish. Manba: MAR, PDF 10-16; Qo'shimcha sahifa raqami (qaysi PDFga tegishli ekani aniqlanmagan; darslik sahifasi deb tasdiqlanmagan): 101
\end{tcolorbox}
\vspace{2mm}\hrule height 0.2pt\vspace{3mm}

% DB_ID c2682ded-d624-3ed3-9ea2-a3e58c92d629
\needspace{5\baselineskip}
{\bfseries\sffamily\color{ink} 536.} {scriptsizecolor{muted} Y1 | review | c2682ded}
\par
Drayver, utilita, firmware va translatorlar mavzusidagi vazifada qurilma ichidagi past darajali boshqaruv dasturi talab qilindi. Qaysi tushuncha mos?
\begin{enumerate}[label=\Alph*),leftmargin=9mm]
\item drayver
\item utilita
\item kompilyator
\item firmware
\end{enumerate}
\begin{tcolorbox}[enhanced,breakable,colback=teal!4,colframe=teal!60!black,boxrule=0.45pt,arc=1.2mm,left=2.7mm,right=2.7mm,top=1.8mm,bottom=1.8mm,before skip=2mm,after skip=2mm]
{\bfseries\sffamily\color{teal}To'g'ri javob:} \textbf{D) firmware}
\par\smallskip
{\bfseries Izoh:} firmware. Berilgan vazifa firmware tushunchasiga mos: qurilma ichidagi past darajali boshqaruv dasturi. Manba: MAR, PDF 10-16; Qo'shimcha sahifa raqami (qaysi PDFga tegishli ekani aniqlanmagan; darslik sahifasi deb tasdiqlanmagan): 71
\end{tcolorbox}
\vspace{2mm}\hrule height 0.2pt\vspace{3mm}

% DB_ID cc9ae83b-c13e-3f15-fff0-bc59ca2a7e22
\needspace{5\baselineskip}
{\bfseries\sffamily\color{ink} 537.} {scriptsizecolor{muted} Y1 | review | cc9ae83b}
\par
Dastur o'rnatish, yangilash, litsenziya va moslik mavzusidagi vazifada dastur komponentlarini tizimda ishlashga tayyorlab joylashtirish talab qilindi. Qaysi tushuncha mos?
\begin{enumerate}[label=\Alph*),leftmargin=9mm]
\item uninstall
\item update
\item moslik
\item installation
\end{enumerate}
\begin{tcolorbox}[enhanced,breakable,colback=teal!4,colframe=teal!60!black,boxrule=0.45pt,arc=1.2mm,left=2.7mm,right=2.7mm,top=1.8mm,bottom=1.8mm,before skip=2mm,after skip=2mm]
{\bfseries\sffamily\color{teal}To'g'ri javob:} \textbf{D) installation}
\par\smallskip
{\bfseries Izoh:} installation. Berilgan vazifa installation tushunchasiga mos: dastur komponentlarini tizimda ishlashga tayyorlab joylashtirish. Manba: MAR, PDF 10-16; Qo'shimcha sahifa raqami (qaysi PDFga tegishli ekani aniqlanmagan; darslik sahifasi deb tasdiqlanmagan): 101
\end{tcolorbox}
\vspace{2mm}\hrule height 0.2pt\vspace{3mm}

% DB_ID d2975ab8-e704-b76c-fbc2-e8aecc83379b
\needspace{5\baselineskip}
{\bfseries\sffamily\color{ink} 538.} {scriptsizecolor{muted} Y1 | review | d2975ab8}
\par
Dasturiy ta'minotning tasnifi mavzusidagi vazifada dastur kodini boshqa, odatda mashina bajaradigan shaklga o'giruvchi vosita talab qilindi. Qaysi tushuncha mos?
\begin{enumerate}[label=\Alph*),leftmargin=9mm]
\item tizimli dastur
\item amaliy dastur
\item translator
\item utilita
\end{enumerate}
\begin{tcolorbox}[enhanced,breakable,colback=teal!4,colframe=teal!60!black,boxrule=0.45pt,arc=1.2mm,left=2.7mm,right=2.7mm,top=1.8mm,bottom=1.8mm,before skip=2mm,after skip=2mm]
{\bfseries\sffamily\color{teal}To'g'ri javob:} \textbf{C) translator}
\par\smallskip
{\bfseries Izoh:} translator. Berilgan vazifa translator tushunchasiga mos: dastur kodini boshqa, odatda mashina bajaradigan shaklga o'giruvchi vosita. Manba: MAR, PDF 10-12; Qo'shimcha sahifa raqami (qaysi PDFga tegishli ekani aniqlanmagan; darslik sahifasi deb tasdiqlanmagan): 52
\end{tcolorbox}
\vspace{2mm}\hrule height 0.2pt\vspace{3mm}

% DB_ID d8165070-e009-8050-e603-8258f2edec3d
\needspace{5\baselineskip}
{\bfseries\sffamily\color{ink} 539.} {scriptsizecolor{muted} Y1 | review | d8165070}
\par
Dastur o'rnatish, yangilash, litsenziya va moslik mavzusidagi vazifada manba kodi litsenziya asosida ochiq dastur talab qilindi. Qaysi tushuncha mos?
\begin{enumerate}[label=\Alph*),leftmargin=9mm]
\item installation
\item uninstall
\item update
\item open-source
\end{enumerate}
\begin{tcolorbox}[enhanced,breakable,colback=teal!4,colframe=teal!60!black,boxrule=0.45pt,arc=1.2mm,left=2.7mm,right=2.7mm,top=1.8mm,bottom=1.8mm,before skip=2mm,after skip=2mm]
{\bfseries\sffamily\color{teal}To'g'ri javob:} \textbf{D) open-source}
\par\smallskip
{\bfseries Izoh:} open-source. Berilgan vazifa open-source tushunchasiga mos: manba kodi litsenziya asosida ochiq dastur. Manba: MAR, PDF 10-16; Qo'shimcha sahifa raqami (qaysi PDFga tegishli ekani aniqlanmagan; darslik sahifasi deb tasdiqlanmagan): 106
\end{tcolorbox}
\vspace{2mm}\hrule height 0.2pt\vspace{3mm}

% DB_ID d88fd672-fc8b-3b96-f8b1-bd2076492bdf
\needspace{5\baselineskip}
{\bfseries\sffamily\color{ink} 540.} {scriptsizecolor{muted} Y1 | review | d88fd672}
\par
Amaliy dasturlar va vazifaga mos tanlov mavzusidagi vazifada muayyan tashkilot ehtiyojiga mos ishlab chiqilgan dastur talab qilindi. Qaysi tushuncha mos?
\begin{enumerate}[label=\Alph*),leftmargin=9mm]
\item maxsus dastur
\item matn protsessori
\item elektron jadval
\item DBMS
\end{enumerate}
\begin{tcolorbox}[enhanced,breakable,colback=teal!4,colframe=teal!60!black,boxrule=0.45pt,arc=1.2mm,left=2.7mm,right=2.7mm,top=1.8mm,bottom=1.8mm,before skip=2mm,after skip=2mm]
{\bfseries\sffamily\color{teal}To'g'ri javob:} \textbf{A) maxsus dastur}
\par\smallskip
{\bfseries Izoh:} maxsus dastur. Berilgan vazifa maxsus dastur tushunchasiga mos: muayyan tashkilot ehtiyojiga mos ishlab chiqilgan dastur. Manba: ICT8, PDF 95-110; Qo'shimcha sahifa raqami (qaysi PDFga tegishli ekani aniqlanmagan; darslik sahifasi deb tasdiqlanmagan): 81
\end{tcolorbox}
\vspace{2mm}\hrule height 0.2pt\vspace{3mm}

% DB_ID da7971a8-7bf7-7338-3a53-f75fc42d88f4
\needspace{5\baselineskip}
{\bfseries\sffamily\color{ink} 541.} {scriptsizecolor{muted} Y1 | review | da7971a8}
\par
Dasturiy ta'minotning tasnifi mavzusidagi vazifada ko'p foydalanuvchi va vazifalarga mo'ljallangan tayyor dastur talab qilindi. Qaysi tushuncha mos?
\begin{enumerate}[label=\Alph*),leftmargin=9mm]
\item tizimli dastur
\item amaliy dastur
\item utilita
\item umumiy maqsadli dastur
\end{enumerate}
\begin{tcolorbox}[enhanced,breakable,colback=teal!4,colframe=teal!60!black,boxrule=0.45pt,arc=1.2mm,left=2.7mm,right=2.7mm,top=1.8mm,bottom=1.8mm,before skip=2mm,after skip=2mm]
{\bfseries\sffamily\color{teal}To'g'ri javob:} \textbf{D) umumiy maqsadli dastur}
\par\smallskip
{\bfseries Izoh:} umumiy maqsadli dastur. Berilgan vazifa umumiy maqsadli dastur tushunchasiga mos: ko'p foydalanuvchi va vazifalarga mo'ljallangan tayyor dastur. Manba: CAM9, PDF 11-15; Qo'shimcha sahifa raqami (qaysi PDFga tegishli ekani aniqlanmagan; darslik sahifasi deb tasdiqlanmagan): 53
\end{tcolorbox}
\vspace{2mm}\hrule height 0.2pt\vspace{3mm}

% DB_ID e797607a-0aaa-f760-7ed6-fc37d7de2f87
\needspace{5\baselineskip}
{\bfseries\sffamily\color{ink} 542.} {scriptsizecolor{muted} Y1 | review | e797607a}
\par
Dasturiy ta'minotning tasnifi mavzusida "utilita" qatnashgan atama-ta'rif juftliklaridan to'g'risini toping.
\begin{enumerate}[label=\Alph*),leftmargin=9mm]
\item utilita - kompyuter resurslarini boshqarib, boshqa dasturlarga muhit yaratuvchi dastur
\item utilita - foydalanuvchining aniq vazifasini bajaruvchi dastur
\item utilita - tizimga xizmat ko'rsatuvchi yordamchi dastur
\item utilita - dastur kodini boshqa, odatda mashina bajaradigan shaklga o'giruvchi vosita
\end{enumerate}
\begin{tcolorbox}[enhanced,breakable,colback=teal!4,colframe=teal!60!black,boxrule=0.45pt,arc=1.2mm,left=2.7mm,right=2.7mm,top=1.8mm,bottom=1.8mm,before skip=2mm,after skip=2mm]
{\bfseries\sffamily\color{teal}To'g'ri javob:} \textbf{C) utilita - tizimga xizmat ko'rsatuvchi yordamchi dastur}
\par\smallskip
{\bfseries Izoh:} utilita - tizimga xizmat ko'rsatuvchi yordamchi dastur. To'g'ri juftlik: utilita - tizimga xizmat ko'rsatuvchi yordamchi dastur. Manba: ICT5, PDF 38-41; Qo'shimcha sahifa raqami (qaysi PDFga tegishli ekani aniqlanmagan; darslik sahifasi deb tasdiqlanmagan): 51
\end{tcolorbox}
\vspace{2mm}\hrule height 0.2pt\vspace{3mm}

% DB_ID ea60e68b-8cae-2b0d-a4fd-3d4d641afb29
\needspace{5\baselineskip}
{\bfseries\sffamily\color{ink} 543.} {scriptsizecolor{muted} Y1 | review | ea60e68b}
\par
Dasturiy ta'minotning tasnifi mavzusidagi vazifada tizimga xizmat ko'rsatuvchi yordamchi dastur talab qilindi. Qaysi tushuncha mos?
\begin{enumerate}[label=\Alph*),leftmargin=9mm]
\item tizimli dastur
\item utilita
\item amaliy dastur
\item translator
\end{enumerate}
\begin{tcolorbox}[enhanced,breakable,colback=teal!4,colframe=teal!60!black,boxrule=0.45pt,arc=1.2mm,left=2.7mm,right=2.7mm,top=1.8mm,bottom=1.8mm,before skip=2mm,after skip=2mm]
{\bfseries\sffamily\color{teal}To'g'ri javob:} \textbf{B) utilita}
\par\smallskip
{\bfseries Izoh:} utilita. Berilgan vazifa utilita tushunchasiga mos: tizimga xizmat ko'rsatuvchi yordamchi dastur. Manba: ICT5, PDF 38-41; Qo'shimcha sahifa raqami (qaysi PDFga tegishli ekani aniqlanmagan; darslik sahifasi deb tasdiqlanmagan): 51
\end{tcolorbox}
\vspace{2mm}\hrule height 0.2pt\vspace{3mm}

% DB_ID eace7868-5256-9732-e0c0-c739a4a8e454
\needspace{5\baselineskip}
{\bfseries\sffamily\color{ink} 544.} {scriptsizecolor{muted} Y1 | review | eace7868}
\par
Drayver, utilita, firmware va translatorlar mavzusidagi vazifada butun manba kodini obyekt yoki bajariladigan kodga o'giruvchi translator talab qilindi. Qaysi tushuncha mos?
\begin{enumerate}[label=\Alph*),leftmargin=9mm]
\item kompilyator
\item drayver
\item utilita
\item firmware
\end{enumerate}
\begin{tcolorbox}[enhanced,breakable,colback=teal!4,colframe=teal!60!black,boxrule=0.45pt,arc=1.2mm,left=2.7mm,right=2.7mm,top=1.8mm,bottom=1.8mm,before skip=2mm,after skip=2mm]
{\bfseries\sffamily\color{teal}To'g'ri javob:} \textbf{A) kompilyator}
\par\smallskip
{\bfseries Izoh:} kompilyator. Berilgan vazifa kompilyator tushunchasiga mos: butun manba kodini obyekt yoki bajariladigan kodga o'giruvchi translator. Manba: ICT8, PDF 106-110; Qo'shimcha sahifa raqami (qaysi PDFga tegishli ekani aniqlanmagan; darslik sahifasi deb tasdiqlanmagan): 72
\end{tcolorbox}
\vspace{2mm}\hrule height 0.2pt\vspace{3mm}

% DB_ID ee642fe0-76a6-a481-25bd-a48a8e882bbd
\needspace{5\baselineskip}
{\bfseries\sffamily\color{ink} 545.} {scriptsizecolor{muted} Y1 | review | ee642fe0}
\par
Dasturiy ta'minotning tasnifi mavzusida "tizimli dastur" qatnashgan atama-ta'rif juftliklaridan to'g'risini toping.
\begin{enumerate}[label=\Alph*),leftmargin=9mm]
\item tizimli dastur - kompyuter resurslarini boshqarib, boshqa dasturlarga muhit yaratuvchi dastur
\item tizimli dastur - foydalanuvchining aniq vazifasini bajaruvchi dastur
\item tizimli dastur - tizimga xizmat ko'rsatuvchi yordamchi dastur
\item tizimli dastur - dastur kodini boshqa, odatda mashina bajaradigan shaklga o'giruvchi vosita
\end{enumerate}
\begin{tcolorbox}[enhanced,breakable,colback=teal!4,colframe=teal!60!black,boxrule=0.45pt,arc=1.2mm,left=2.7mm,right=2.7mm,top=1.8mm,bottom=1.8mm,before skip=2mm,after skip=2mm]
{\bfseries\sffamily\color{teal}To'g'ri javob:} \textbf{A) tizimli dastur - kompyuter resurslarini boshqarib, boshqa dasturlarga muhit yaratuvchi dastur}
\par\smallskip
{\bfseries Izoh:} tizimli dastur - kompyuter resurslarini boshqarib, boshqa dasturlarga muhit yaratuvchi dastur. To'g'ri juftlik: tizimli dastur - kompyuter resurslarini boshqarib, boshqa dasturlarga muhit yaratuvchi dastur. Manba: CAM9, PDF 11-15; Qo'shimcha sahifa raqami (qaysi PDFga tegishli ekani aniqlanmagan; darslik sahifasi deb tasdiqlanmagan): 49
\end{tcolorbox}
\vspace{2mm}\hrule height 0.2pt\vspace{3mm}

% DB_ID 0859e838-d502-b574-3989-c100a5578ef6
\needspace{5\baselineskip}
{\bfseries\sffamily\color{ink} 546.} {scriptsizecolor{muted} Y1 | review | 0859e838}
\par
Amaliy dasturlar va vazifaga mos tanlov mavzusida "maxsus dastur" qatnashgan atama-ta'rif juftliklaridan to'g'risini toping.
\begin{enumerate}[label=\Alph*),leftmargin=9mm]
\item maxsus dastur - tuzilmali matnli hujjat yaratish va formatlash dasturi
\item maxsus dastur - muayyan tashkilot ehtiyojiga mos ishlab chiqilgan dastur
\item maxsus dastur - formulalar asosida hisoblash, model va diagramma yaratish dasturi
\item maxsus dastur - tuzilmali ma'lumotlar bazasini yaratish, so'rash va boshqarish tizimi
\end{enumerate}
\begin{tcolorbox}[enhanced,breakable,colback=teal!4,colframe=teal!60!black,boxrule=0.45pt,arc=1.2mm,left=2.7mm,right=2.7mm,top=1.8mm,bottom=1.8mm,before skip=2mm,after skip=2mm]
{\bfseries\sffamily\color{teal}To'g'ri javob:} \textbf{B) maxsus dastur - muayyan tashkilot ehtiyojiga mos ishlab chiqilgan dastur}
\par\smallskip
{\bfseries Izoh:} maxsus dastur - muayyan tashkilot ehtiyojiga mos ishlab chiqilgan dastur. To'g'ri juftlik: maxsus dastur - muayyan tashkilot ehtiyojiga mos ishlab chiqilgan dastur. Manba: ICT8, PDF 95-110; Qo'shimcha sahifa raqami (qaysi PDFga tegishli ekani aniqlanmagan; darslik sahifasi deb tasdiqlanmagan): 82
\end{tcolorbox}
\vspace{2mm}\hrule height 0.2pt\vspace{3mm}

% DB_ID 11091fd3-b746-d62a-0b3b-2cba192922d4
\needspace{5\baselineskip}
{\bfseries\sffamily\color{ink} 547.} {scriptsizecolor{muted} Y1 | review | 11091fd3}
\par
Drayver, utilita, firmware va translatorlar bo'yicha qatorlar ichidan "drayver" noto'g'ri izohlanganini aniqlang.
\begin{enumerate}[label=\Alph*),leftmargin=9mm]
\item firmware - qurilma ichidagi past darajali boshqaruv dasturi
\item kompilyator - butun manba kodini obyekt yoki bajariladigan kodga o'giruvchi translator
\item interpretator - kodni bajarish paytida bosqichma-bosqich tarjima qiluvchi translator
\item drayver - tizimga xizmat ko'rsatuvchi yordamchi dastur
\end{enumerate}
\begin{tcolorbox}[enhanced,breakable,colback=teal!4,colframe=teal!60!black,boxrule=0.45pt,arc=1.2mm,left=2.7mm,right=2.7mm,top=1.8mm,bottom=1.8mm,before skip=2mm,after skip=2mm]
{\bfseries\sffamily\color{teal}To'g'ri javob:} \textbf{D) drayver - tizimga xizmat ko'rsatuvchi yordamchi dastur}
\par\smallskip
{\bfseries Izoh:} drayver - tizimga xizmat ko'rsatuvchi yordamchi dastur. drayver uchun to'g'ri ta'rif: OT bilan aniq qurilma modeli o'rtasida buyruqlarni moslashtiruvchi dastur. Manba: CAM9, PDF 14-15; Qo'shimcha sahifa raqami (qaysi PDFga tegishli ekani aniqlanmagan; darslik sahifasi deb tasdiqlanmagan): 69
\end{tcolorbox}
\vspace{2mm}\hrule height 0.2pt\vspace{3mm}

% DB_ID 1dec924a-c69d-5203-634a-9f87d3eb68ff
\needspace{5\baselineskip}
{\bfseries\sffamily\color{ink} 548.} {scriptsizecolor{muted} Y1 | review | 1dec924a}
\par
Dastur bepul yuklanadi, ammo manba kodi yopiq. Qaysi atama mos?
\begin{enumerate}[label=\Alph*),leftmargin=9mm]
\item open-source
\item firmware
\item public domain shart
\item freeware
\end{enumerate}
\begin{tcolorbox}[enhanced,breakable,colback=teal!4,colframe=teal!60!black,boxrule=0.45pt,arc=1.2mm,left=2.7mm,right=2.7mm,top=1.8mm,bottom=1.8mm,before skip=2mm,after skip=2mm]
{\bfseries\sffamily\color{teal}To'g'ri javob:} \textbf{D) freeware}
\par\smallskip
{\bfseries Izoh:} freeware. Bepul foydalanish manba kodi ochiqligini anglatmaydi. Manba: MAR, PDF 10-16; Qo'shimcha sahifa raqami (qaysi PDFga tegishli ekani aniqlanmagan; darslik sahifasi deb tasdiqlanmagan): 107
\end{tcolorbox}
\vspace{2mm}\hrule height 0.2pt\vspace{3mm}

% DB_ID 297f760c-3ca4-dfdc-1802-b4e54503a915
\needspace{5\baselineskip}
{\bfseries\sffamily\color{ink} 549.} {scriptsizecolor{muted} Y1 | review | 297f760c}
\par
Dastur o'rnatish, yangilash, litsenziya va moslik bo'yicha qatorlar ichidan "open-source" noto'g'ri izohlanganini aniqlang.
\begin{enumerate}[label=\Alph*),leftmargin=9mm]
\item uninstall - dastur va ro'yxatdan o'tgan komponentlarini nazoratli olib tashlash
\item open-source - dastur komponentlarini tizimda ishlashga tayyorlab joylashtirish
\item update - mavjud versiyaga tuzatish yoki kichik yaxshilanish o'rnatish
\item moslik - qurilma, OT, dastur va formatlarning birgalikda ishlay olish xususiyati
\end{enumerate}
\begin{tcolorbox}[enhanced,breakable,colback=teal!4,colframe=teal!60!black,boxrule=0.45pt,arc=1.2mm,left=2.7mm,right=2.7mm,top=1.8mm,bottom=1.8mm,before skip=2mm,after skip=2mm]
{\bfseries\sffamily\color{teal}To'g'ri javob:} \textbf{B) open-source - dastur komponentlarini tizimda ishlashga tayyorlab joylashtirish}
\par\smallskip
{\bfseries Izoh:} open-source - dastur komponentlarini tizimda ishlashga tayyorlab joylashtirish. open-source uchun to'g'ri ta'rif: manba kodi litsenziya asosida ochiq dastur. Manba: MAR, PDF 10-16; Qo'shimcha sahifa raqami (qaysi PDFga tegishli ekani aniqlanmagan; darslik sahifasi deb tasdiqlanmagan): 106
\end{tcolorbox}
\vspace{2mm}\hrule height 0.2pt\vspace{3mm}

% DB_ID 36e740ec-ae9a-ed70-c227-199b264618e7
\needspace{5\baselineskip}
{\bfseries\sffamily\color{ink} 550.} {scriptsizecolor{muted} Y1 | review | 36e740ec}
\par
Amaliy dasturlar va vazifaga mos tanlov mavzusida "DTP" qatnashgan atama-ta'rif juftliklaridan to'g'risini toping.
\begin{enumerate}[label=\Alph*),leftmargin=9mm]
\item DTP - tuzilmali matnli hujjat yaratish va formatlash dasturi
\item DTP - formulalar asosida hisoblash, model va diagramma yaratish dasturi
\item DTP - tuzilmali ma'lumotlar bazasini yaratish, so'rash va boshqarish tizimi
\item DTP - professional sahifa maketi va nashr tayyorlash dasturi
\end{enumerate}
\begin{tcolorbox}[enhanced,breakable,colback=teal!4,colframe=teal!60!black,boxrule=0.45pt,arc=1.2mm,left=2.7mm,right=2.7mm,top=1.8mm,bottom=1.8mm,before skip=2mm,after skip=2mm]
{\bfseries\sffamily\color{teal}To'g'ri javob:} \textbf{D) DTP - professional sahifa maketi va nashr tayyorlash dasturi}
\par\smallskip
{\bfseries Izoh:} DTP - professional sahifa maketi va nashr tayyorlash dasturi. To'g'ri juftlik: DTP - professional sahifa maketi va nashr tayyorlash dasturi. Manba: CAM5, PDF 45-138; Qo'shimcha sahifa raqami (qaysi PDFga tegishli ekani aniqlanmagan; darslik sahifasi deb tasdiqlanmagan): 80
\end{tcolorbox}
\vspace{2mm}\hrule height 0.2pt\vspace{3mm}

% DB_ID 38fcb52d-8add-fe7e-4019-cae37027a97b
\needspace{5\baselineskip}
{\bfseries\sffamily\color{ink} 551.} {scriptsizecolor{muted} Y1 | review | 38fcb52d}
\par
Dasturiy ta'minotning tasnifi mavzusida "translator" qatnashgan atama-ta'rif juftliklaridan to'g'risini toping.
\begin{enumerate}[label=\Alph*),leftmargin=9mm]
\item translator - kompyuter resurslarini boshqarib, boshqa dasturlarga muhit yaratuvchi dastur
\item translator - foydalanuvchining aniq vazifasini bajaruvchi dastur
\item translator - tizimga xizmat ko'rsatuvchi yordamchi dastur
\item translator - dastur kodini boshqa, odatda mashina bajaradigan shaklga o'giruvchi vosita
\end{enumerate}
\begin{tcolorbox}[enhanced,breakable,colback=teal!4,colframe=teal!60!black,boxrule=0.45pt,arc=1.2mm,left=2.7mm,right=2.7mm,top=1.8mm,bottom=1.8mm,before skip=2mm,after skip=2mm]
{\bfseries\sffamily\color{teal}To'g'ri javob:} \textbf{D) translator - dastur kodini boshqa, odatda mashina bajaradigan shaklga o'giruvchi vosita}
\par\smallskip
{\bfseries Izoh:} translator - dastur kodini boshqa, odatda mashina bajaradigan shaklga o'giruvchi vosita. To'g'ri juftlik: translator - dastur kodini boshqa, odatda mashina bajaradigan shaklga o'giruvchi vosita. Manba: MAR, PDF 10-12; Qo'shimcha sahifa raqami (qaysi PDFga tegishli ekani aniqlanmagan; darslik sahifasi deb tasdiqlanmagan): 52
\end{tcolorbox}
\vspace{2mm}\hrule height 0.2pt\vspace{3mm}

% DB_ID 4818c272-114a-4cc8-f579-6e8388e1a08a
\needspace{5\baselineskip}
{\bfseries\sffamily\color{ink} 552.} {scriptsizecolor{muted} Y1 | review | 4818c272}
\par
Drayver, utilita, firmware va translatorlar mavzusida "linker" qatnashgan atama-ta'rif juftliklaridan to'g'risini toping.
\begin{enumerate}[label=\Alph*),leftmargin=9mm]
\item linker - OT bilan aniq qurilma modeli o'rtasida buyruqlarni moslashtiruvchi dastur
\item linker - tizimga xizmat ko'rsatuvchi yordamchi dastur
\item linker - qurilma ichidagi past darajali boshqaruv dasturi
\item linker - obyekt modullari va kutubxonalarni bajariladigan faylga bog'lovchi vosita
\end{enumerate}
\begin{tcolorbox}[enhanced,breakable,colback=teal!4,colframe=teal!60!black,boxrule=0.45pt,arc=1.2mm,left=2.7mm,right=2.7mm,top=1.8mm,bottom=1.8mm,before skip=2mm,after skip=2mm]
{\bfseries\sffamily\color{teal}To'g'ri javob:} \textbf{D) linker - obyekt modullari va kutubxonalarni bajariladigan faylga bog'lovchi vosita}
\par\smallskip
{\bfseries Izoh:} linker - obyekt modullari va kutubxonalarni bajariladigan faylga bog'lovchi vosita. To'g'ri juftlik: linker - obyekt modullari va kutubxonalarni bajariladigan faylga bog'lovchi vosita. Manba: CAM10/11, PDF 44-49; Qo'shimcha sahifa raqami (qaysi PDFga tegishli ekani aniqlanmagan; darslik sahifasi deb tasdiqlanmagan): 75
\end{tcolorbox}
\vspace{2mm}\hrule height 0.2pt\vspace{3mm}

% DB_ID 4bad6be9-67d7-c112-5221-97d5c565570a
\needspace{5\baselineskip}
{\bfseries\sffamily\color{ink} 553.} {scriptsizecolor{muted} Y1 | review | 4bad6be9}
\par
Drayver, utilita, firmware va translatorlar mavzusida "utilita" qatnashgan atama-ta'rif juftliklaridan to'g'risini toping.
\begin{enumerate}[label=\Alph*),leftmargin=9mm]
\item utilita - OT bilan aniq qurilma modeli o'rtasida buyruqlarni moslashtiruvchi dastur
\item utilita - qurilma ichidagi past darajali boshqaruv dasturi
\item utilita - butun manba kodini obyekt yoki bajariladigan kodga o'giruvchi translator
\item utilita - tizimga xizmat ko'rsatuvchi yordamchi dastur
\end{enumerate}
\begin{tcolorbox}[enhanced,breakable,colback=teal!4,colframe=teal!60!black,boxrule=0.45pt,arc=1.2mm,left=2.7mm,right=2.7mm,top=1.8mm,bottom=1.8mm,before skip=2mm,after skip=2mm]
{\bfseries\sffamily\color{teal}To'g'ri javob:} \textbf{D) utilita - tizimga xizmat ko'rsatuvchi yordamchi dastur}
\par\smallskip
{\bfseries Izoh:} utilita - tizimga xizmat ko'rsatuvchi yordamchi dastur. To'g'ri juftlik: utilita - tizimga xizmat ko'rsatuvchi yordamchi dastur. Manba: CAM10/11, PDF 44-49; Qo'shimcha sahifa raqami (qaysi PDFga tegishli ekani aniqlanmagan; darslik sahifasi deb tasdiqlanmagan): 70
\end{tcolorbox}
\vspace{2mm}\hrule height 0.2pt\vspace{3mm}

% DB_ID 5909138d-24bc-6f95-39e8-8b6ffabe1fe6
\needspace{5\baselineskip}
{\bfseries\sffamily\color{ink} 554.} {scriptsizecolor{muted} Y1 | review | 5909138d}
\par
Dastur o'rnatish, yangilash, litsenziya va moslik mavzusida "uninstall" qatnashgan atama-ta'rif juftliklaridan to'g'risini toping.
\begin{enumerate}[label=\Alph*),leftmargin=9mm]
\item uninstall - dastur komponentlarini tizimda ishlashga tayyorlab joylashtirish
\item uninstall - dastur va ro'yxatdan o'tgan komponentlarini nazoratli olib tashlash
\item uninstall - mavjud versiyaga tuzatish yoki kichik yaxshilanish o'rnatish
\item uninstall - qurilma, OT, dastur va formatlarning birgalikda ishlay olish xususiyati
\end{enumerate}
\begin{tcolorbox}[enhanced,breakable,colback=teal!4,colframe=teal!60!black,boxrule=0.45pt,arc=1.2mm,left=2.7mm,right=2.7mm,top=1.8mm,bottom=1.8mm,before skip=2mm,after skip=2mm]
{\bfseries\sffamily\color{teal}To'g'ri javob:} \textbf{B) uninstall - dastur va ro'yxatdan o'tgan komponentlarini nazoratli olib tashlash}
\par\smallskip
{\bfseries Izoh:} uninstall - dastur va ro'yxatdan o'tgan komponentlarini nazoratli olib tashlash. To'g'ri juftlik: uninstall - dastur va ro'yxatdan o'tgan komponentlarini nazoratli olib tashlash. Manba: CAM8, PDF 99-110; Qo'shimcha sahifa raqami (qaysi PDFga tegishli ekani aniqlanmagan; darslik sahifasi deb tasdiqlanmagan): 102
\end{tcolorbox}
\vspace{2mm}\hrule height 0.2pt\vspace{3mm}

% DB_ID 70d9b85e-b2e6-9696-af72-236af93d7531
\needspace{5\baselineskip}
{\bfseries\sffamily\color{ink} 555.} {scriptsizecolor{muted} Y1 | review | 70d9b85e}
\par
Dasturiy ta'minotning tasnifi bo'yicha qatorlar ichidan "utilita" noto'g'ri izohlanganini aniqlang.
\begin{enumerate}[label=\Alph*),leftmargin=9mm]
\item amaliy dastur - foydalanuvchining aniq vazifasini bajaruvchi dastur
\item translator - dastur kodini boshqa, odatda mashina bajaradigan shaklga o'giruvchi vosita
\item umumiy maqsadli dastur - ko'p foydalanuvchi va vazifalarga mo'ljallangan tayyor dastur
\item utilita - kompyuter resurslarini boshqarib, boshqa dasturlarga muhit yaratuvchi dastur
\end{enumerate}
\begin{tcolorbox}[enhanced,breakable,colback=teal!4,colframe=teal!60!black,boxrule=0.45pt,arc=1.2mm,left=2.7mm,right=2.7mm,top=1.8mm,bottom=1.8mm,before skip=2mm,after skip=2mm]
{\bfseries\sffamily\color{teal}To'g'ri javob:} \textbf{D) utilita - kompyuter resurslarini boshqarib, boshqa dasturlarga muhit yaratuvchi dastur}
\par\smallskip
{\bfseries Izoh:} utilita - kompyuter resurslarini boshqarib, boshqa dasturlarga muhit yaratuvchi dastur. utilita uchun to'g'ri ta'rif: tizimga xizmat ko'rsatuvchi yordamchi dastur. Manba: ICT5, PDF 38-41; Qo'shimcha sahifa raqami (qaysi PDFga tegishli ekani aniqlanmagan; darslik sahifasi deb tasdiqlanmagan): 51
\end{tcolorbox}
\vspace{2mm}\hrule height 0.2pt\vspace{3mm}

% DB_ID 866f9e0b-ea3c-ce1b-151d-ab7d3a1e7e30
\needspace{5\baselineskip}
{\bfseries\sffamily\color{ink} 556.} {scriptsizecolor{muted} Y1 | review | 866f9e0b}
\par
Dastur o'rnatish, yangilash, litsenziya va moslik mavzusida "freeware" qatnashgan atama-ta'rif juftliklaridan to'g'risini toping.
\begin{enumerate}[label=\Alph*),leftmargin=9mm]
\item freeware - dastur komponentlarini tizimda ishlashga tayyorlab joylashtirish
\item freeware - odatda bepul foydalaniladigan, ammo kodi shart emas ochiq dastur
\item freeware - dastur va ro'yxatdan o'tgan komponentlarini nazoratli olib tashlash
\item freeware - mavjud versiyaga tuzatish yoki kichik yaxshilanish o'rnatish
\end{enumerate}
\begin{tcolorbox}[enhanced,breakable,colback=teal!4,colframe=teal!60!black,boxrule=0.45pt,arc=1.2mm,left=2.7mm,right=2.7mm,top=1.8mm,bottom=1.8mm,before skip=2mm,after skip=2mm]
{\bfseries\sffamily\color{teal}To'g'ri javob:} \textbf{B) freeware - odatda bepul foydalaniladigan, ammo kodi shart emas ochiq dastur}
\par\smallskip
{\bfseries Izoh:} freeware - odatda bepul foydalaniladigan, ammo kodi shart emas ochiq dastur. To'g'ri juftlik: freeware - odatda bepul foydalaniladigan, ammo kodi shart emas ochiq dastur. Manba: CAM8, PDF 99-110; Qo'shimcha sahifa raqami (qaysi PDFga tegishli ekani aniqlanmagan; darslik sahifasi deb tasdiqlanmagan): 107
\end{tcolorbox}
\vspace{2mm}\hrule height 0.2pt\vspace{3mm}

% DB_ID 879f3b71-cb4a-cb23-04b0-95c9c83f67ef
\needspace{5\baselineskip}
{\bfseries\sffamily\color{ink} 557.} {scriptsizecolor{muted} Y1 | review | 879f3b71}
\par
Brauzer operatsion tizim bilan birga o'rnatilgan. U odatda qaysi turga kiradi?
\begin{enumerate}[label=\Alph*),leftmargin=9mm]
\item qurilma drayveri
\item firmware
\item bootloader
\item amaliy dastur
\end{enumerate}
\begin{tcolorbox}[enhanced,breakable,colback=teal!4,colframe=teal!60!black,boxrule=0.45pt,arc=1.2mm,left=2.7mm,right=2.7mm,top=1.8mm,bottom=1.8mm,before skip=2mm,after skip=2mm]
{\bfseries\sffamily\color{teal}To'g'ri javob:} \textbf{D) amaliy dastur}
\par\smallskip
{\bfseries Izoh:} amaliy dastur. Oldindan o'rnatilganligi tasnifni o'zgartirmaydi; brauzer foydalanuvchi vazifasini bajaradi. Manba: CAM9, PDF 11-15; Qo'shimcha sahifa raqami (qaysi PDFga tegishli ekani aniqlanmagan; darslik sahifasi deb tasdiqlanmagan): 55
\end{tcolorbox}
\vspace{2mm}\hrule height 0.2pt\vspace{3mm}

% DB_ID 8a897b0d-c85d-8971-bfdc-12f448735ffd
\needspace{5\baselineskip}
{\bfseries\sffamily\color{ink} 558.} {scriptsizecolor{muted} Y1 | review | 8a897b0d}
\par
Amaliy dasturlar va vazifaga mos tanlov mavzusida "elektron jadval" qatnashgan atama-ta'rif juftliklaridan to'g'risini toping.
\begin{enumerate}[label=\Alph*),leftmargin=9mm]
\item elektron jadval - tuzilmali matnli hujjat yaratish va formatlash dasturi
\item elektron jadval - formulalar asosida hisoblash, model va diagramma yaratish dasturi
\item elektron jadval - tuzilmali ma'lumotlar bazasini yaratish, so'rash va boshqarish tizimi
\item elektron jadval - professional sahifa maketi va nashr tayyorlash dasturi
\end{enumerate}
\begin{tcolorbox}[enhanced,breakable,colback=teal!4,colframe=teal!60!black,boxrule=0.45pt,arc=1.2mm,left=2.7mm,right=2.7mm,top=1.8mm,bottom=1.8mm,before skip=2mm,after skip=2mm]
{\bfseries\sffamily\color{teal}To'g'ri javob:} \textbf{B) elektron jadval - formulalar asosida hisoblash, model va diagramma yaratish dasturi}
\par\smallskip
{\bfseries Izoh:} elektron jadval - formulalar asosida hisoblash, model va diagramma yaratish dasturi. To'g'ri juftlik: elektron jadval - formulalar asosida hisoblash, model va diagramma yaratish dasturi. Manba: CAM6, PDF 34-205; Qo'shimcha sahifa raqami (qaysi PDFga tegishli ekani aniqlanmagan; darslik sahifasi deb tasdiqlanmagan): 77
\end{tcolorbox}
\vspace{2mm}\hrule height 0.2pt\vspace{3mm}

% DB_ID 94995d81-6ff3-6063-1b38-8817a1127285
\needspace{5\baselineskip}
{\bfseries\sffamily\color{ink} 559.} {scriptsizecolor{muted} Y1 | review | 94995d81}
\par
Drayver, utilita, firmware va translatorlar bo'yicha qatorlar ichidan "firmware" noto'g'ri izohlanganini aniqlang.
\begin{enumerate}[label=\Alph*),leftmargin=9mm]
\item utilita - tizimga xizmat ko'rsatuvchi yordamchi dastur
\item firmware - OT bilan aniq qurilma modeli o'rtasida buyruqlarni moslashtiruvchi dastur
\item kompilyator - butun manba kodini obyekt yoki bajariladigan kodga o'giruvchi translator
\item interpretator - kodni bajarish paytida bosqichma-bosqich tarjima qiluvchi translator
\end{enumerate}
\begin{tcolorbox}[enhanced,breakable,colback=teal!4,colframe=teal!60!black,boxrule=0.45pt,arc=1.2mm,left=2.7mm,right=2.7mm,top=1.8mm,bottom=1.8mm,before skip=2mm,after skip=2mm]
{\bfseries\sffamily\color{teal}To'g'ri javob:} \textbf{B) firmware - OT bilan aniq qurilma modeli o'rtasida buyruqlarni moslashtiruvchi dastur}
\par\smallskip
{\bfseries Izoh:} firmware - OT bilan aniq qurilma modeli o'rtasida buyruqlarni moslashtiruvchi dastur. firmware uchun to'g'ri ta'rif: qurilma ichidagi past darajali boshqaruv dasturi. Manba: MAR, PDF 10-16; Qo'shimcha sahifa raqami (qaysi PDFga tegishli ekani aniqlanmagan; darslik sahifasi deb tasdiqlanmagan): 71
\end{tcolorbox}
\vspace{2mm}\hrule height 0.2pt\vspace{3mm}

% DB_ID 9ea3ed9a-c995-15bc-15f0-1682d91b66d6
\needspace{5\baselineskip}
{\bfseries\sffamily\color{ink} 560.} {scriptsizecolor{muted} Y1 | review | 9ea3ed9a}
\par
Dasturiy ta'minotning tasnifi bo'yicha qatorlar ichidan "tizimli dastur" noto'g'ri izohlanganini aniqlang.
\begin{enumerate}[label=\Alph*),leftmargin=9mm]
\item utilita - tizimga xizmat ko'rsatuvchi yordamchi dastur
\item tizimli dastur - foydalanuvchining aniq vazifasini bajaruvchi dastur
\item translator - dastur kodini boshqa, odatda mashina bajaradigan shaklga o'giruvchi vosita
\item umumiy maqsadli dastur - ko'p foydalanuvchi va vazifalarga mo'ljallangan tayyor dastur
\end{enumerate}
\begin{tcolorbox}[enhanced,breakable,colback=teal!4,colframe=teal!60!black,boxrule=0.45pt,arc=1.2mm,left=2.7mm,right=2.7mm,top=1.8mm,bottom=1.8mm,before skip=2mm,after skip=2mm]
{\bfseries\sffamily\color{teal}To'g'ri javob:} \textbf{B) tizimli dastur - foydalanuvchining aniq vazifasini bajaruvchi dastur}
\par\smallskip
{\bfseries Izoh:} tizimli dastur - foydalanuvchining aniq vazifasini bajaruvchi dastur. tizimli dastur uchun to'g'ri ta'rif: kompyuter resurslarini boshqarib, boshqa dasturlarga muhit yaratuvchi dastur. Manba: CAM9, PDF 11-15; Qo'shimcha sahifa raqami (qaysi PDFga tegishli ekani aniqlanmagan; darslik sahifasi deb tasdiqlanmagan): 49
\end{tcolorbox}
\vspace{2mm}\hrule height 0.2pt\vspace{3mm}

% DB_ID a0bda807-9d13-ffc9-7286-91c8fa8f7fe6
\needspace{5\baselineskip}
{\bfseries\sffamily\color{ink} 561.} {scriptsizecolor{muted} Y1 | review | a0bda807}
\par
Dasturiy ta'minotning tasnifi mavzusida "maxsus dastur" qatnashgan atama-ta'rif juftliklaridan to'g'risini toping.
\begin{enumerate}[label=\Alph*),leftmargin=9mm]
\item maxsus dastur - kompyuter resurslarini boshqarib, boshqa dasturlarga muhit yaratuvchi dastur
\item maxsus dastur - muayyan buyurtmachi yoki jarayon uchun yaratilgan dastur
\item maxsus dastur - foydalanuvchining aniq vazifasini bajaruvchi dastur
\item maxsus dastur - tizimga xizmat ko'rsatuvchi yordamchi dastur
\end{enumerate}
\begin{tcolorbox}[enhanced,breakable,colback=teal!4,colframe=teal!60!black,boxrule=0.45pt,arc=1.2mm,left=2.7mm,right=2.7mm,top=1.8mm,bottom=1.8mm,before skip=2mm,after skip=2mm]
{\bfseries\sffamily\color{teal}To'g'ri javob:} \textbf{B) maxsus dastur - muayyan buyurtmachi yoki jarayon uchun yaratilgan dastur}
\par\smallskip
{\bfseries Izoh:} maxsus dastur - muayyan buyurtmachi yoki jarayon uchun yaratilgan dastur. To'g'ri juftlik: maxsus dastur - muayyan buyurtmachi yoki jarayon uchun yaratilgan dastur. Manba: CAM10/11, PDF 44-58; Qo'shimcha sahifa raqami (qaysi PDFga tegishli ekani aniqlanmagan; darslik sahifasi deb tasdiqlanmagan): 54
\end{tcolorbox}
\vspace{2mm}\hrule height 0.2pt\vspace{3mm}

% DB_ID a29e14ad-8773-13cd-f253-41bf506fa3a6
\needspace{5\baselineskip}
{\bfseries\sffamily\color{ink} 562.} {scriptsizecolor{muted} Y1 | review | a29e14ad}
\par
Dasturiy ta'minotning tasnifi bo'yicha qatorlar ichidan "umumiy maqsadli dastur" noto'g'ri izohlanganini aniqlang.
\begin{enumerate}[label=\Alph*),leftmargin=9mm]
\item amaliy dastur - foydalanuvchining aniq vazifasini bajaruvchi dastur
\item umumiy maqsadli dastur - kompyuter resurslarini boshqarib, boshqa dasturlarga muhit yaratuvchi dastur
\item utilita - tizimga xizmat ko'rsatuvchi yordamchi dastur
\item translator - dastur kodini boshqa, odatda mashina bajaradigan shaklga o'giruvchi vosita
\end{enumerate}
\begin{tcolorbox}[enhanced,breakable,colback=teal!4,colframe=teal!60!black,boxrule=0.45pt,arc=1.2mm,left=2.7mm,right=2.7mm,top=1.8mm,bottom=1.8mm,before skip=2mm,after skip=2mm]
{\bfseries\sffamily\color{teal}To'g'ri javob:} \textbf{B) umumiy maqsadli dastur - kompyuter resurslarini boshqarib, boshqa dasturlarga muhit yaratuvchi dastur}
\par\smallskip
{\bfseries Izoh:} umumiy maqsadli dastur - kompyuter resurslarini boshqarib, boshqa dasturlarga muhit yaratuvchi dastur. umumiy maqsadli dastur uchun to'g'ri ta'rif: ko'p foydalanuvchi va vazifalarga mo'ljallangan tayyor dastur. Manba: CAM9, PDF 11-15; Qo'shimcha sahifa raqami (qaysi PDFga tegishli ekani aniqlanmagan; darslik sahifasi deb tasdiqlanmagan): 53
\end{tcolorbox}
\vspace{2mm}\hrule height 0.2pt\vspace{3mm}

% DB_ID bcc3f31a-5b14-bd64-6f56-ff5b01fbff16
\needspace{5\baselineskip}
{\bfseries\sffamily\color{ink} 563.} {scriptsizecolor{muted} Y1 | review | bcc3f31a}
\par
Amaliy dasturlar va vazifaga mos tanlov bo'yicha qatorlar ichidan "CAD" noto'g'ri izohlanganini aniqlang.
\begin{enumerate}[label=\Alph*),leftmargin=9mm]
\item elektron jadval - formulalar asosida hisoblash, model va diagramma yaratish dasturi
\item CAD - tuzilmali matnli hujjat yaratish va formatlash dasturi
\item DBMS - tuzilmali ma'lumotlar bazasini yaratish, so'rash va boshqarish tizimi
\item DTP - professional sahifa maketi va nashr tayyorlash dasturi
\end{enumerate}
\begin{tcolorbox}[enhanced,breakable,colback=teal!4,colframe=teal!60!black,boxrule=0.45pt,arc=1.2mm,left=2.7mm,right=2.7mm,top=1.8mm,bottom=1.8mm,before skip=2mm,after skip=2mm]
{\bfseries\sffamily\color{teal}To'g'ri javob:} \textbf{B) CAD - tuzilmali matnli hujjat yaratish va formatlash dasturi}
\par\smallskip
{\bfseries Izoh:} CAD - tuzilmali matnli hujjat yaratish va formatlash dasturi. CAD uchun to'g'ri ta'rif: aniq texnik chizma va model yaratish dasturi. Manba: ICT6, PDF 62-126; Qo'shimcha sahifa raqami (qaysi PDFga tegishli ekani aniqlanmagan; darslik sahifasi deb tasdiqlanmagan): 81
\end{tcolorbox}
\vspace{2mm}\hrule height 0.2pt\vspace{3mm}

% DB_ID be64f553-8fc5-2ea0-5cd6-8b54a18bb853
\needspace{5\baselineskip}
{\bfseries\sffamily\color{ink} 564.} {scriptsizecolor{muted} Y1 | review | be64f553}
\par
Dastur o'rnatish, yangilash, litsenziya va moslik bo'yicha qatorlar ichidan "update" noto'g'ri izohlanganini aniqlang.
\begin{enumerate}[label=\Alph*),leftmargin=9mm]
\item uninstall - dastur va ro'yxatdan o'tgan komponentlarini nazoratli olib tashlash
\item moslik - qurilma, OT, dastur va formatlarning birgalikda ishlay olish xususiyati
\item open-source - manba kodi litsenziya asosida ochiq dastur
\item update - dastur komponentlarini tizimda ishlashga tayyorlab joylashtirish
\end{enumerate}
\begin{tcolorbox}[enhanced,breakable,colback=teal!4,colframe=teal!60!black,boxrule=0.45pt,arc=1.2mm,left=2.7mm,right=2.7mm,top=1.8mm,bottom=1.8mm,before skip=2mm,after skip=2mm]
{\bfseries\sffamily\color{teal}To'g'ri javob:} \textbf{D) update - dastur komponentlarini tizimda ishlashga tayyorlab joylashtirish}
\par\smallskip
{\bfseries Izoh:} update - dastur komponentlarini tizimda ishlashga tayyorlab joylashtirish. update uchun to'g'ri ta'rif: mavjud versiyaga tuzatish yoki kichik yaxshilanish o'rnatish. Manba: ICT8, PDF 106-110; Qo'shimcha sahifa raqami (qaysi PDFga tegishli ekani aniqlanmagan; darslik sahifasi deb tasdiqlanmagan): 104
\end{tcolorbox}
\vspace{2mm}\hrule height 0.2pt\vspace{3mm}

% DB_ID c10d16e2-05d1-9908-7f59-104411e7aa07
\needspace{5\baselineskip}
{\bfseries\sffamily\color{ink} 565.} {scriptsizecolor{muted} Y1 | review | c10d16e2}
\par
Dastur o'rnatish, yangilash, litsenziya va moslik bo'yicha qatorlar ichidan "installation" noto'g'ri izohlanganini aniqlang.
\begin{enumerate}[label=\Alph*),leftmargin=9mm]
\item update - mavjud versiyaga tuzatish yoki kichik yaxshilanish o'rnatish
\item installation - dastur va ro'yxatdan o'tgan komponentlarini nazoratli olib tashlash
\item moslik - qurilma, OT, dastur va formatlarning birgalikda ishlay olish xususiyati
\item open-source - manba kodi litsenziya asosida ochiq dastur
\end{enumerate}
\begin{tcolorbox}[enhanced,breakable,colback=teal!4,colframe=teal!60!black,boxrule=0.45pt,arc=1.2mm,left=2.7mm,right=2.7mm,top=1.8mm,bottom=1.8mm,before skip=2mm,after skip=2mm]
{\bfseries\sffamily\color{teal}To'g'ri javob:} \textbf{B) installation - dastur va ro'yxatdan o'tgan komponentlarini nazoratli olib tashlash}
\par\smallskip
{\bfseries Izoh:} installation - dastur va ro'yxatdan o'tgan komponentlarini nazoratli olib tashlash. installation uchun to'g'ri ta'rif: dastur komponentlarini tizimda ishlashga tayyorlab joylashtirish. Manba: MAR, PDF 10-16; Qo'shimcha sahifa raqami (qaysi PDFga tegishli ekani aniqlanmagan; darslik sahifasi deb tasdiqlanmagan): 101
\end{tcolorbox}
\vspace{2mm}\hrule height 0.2pt\vspace{3mm}

% DB_ID d7d9ad46-9a4f-2383-4c3c-b958f443e15b
\needspace{5\baselineskip}
{\bfseries\sffamily\color{ink} 566.} {scriptsizecolor{muted} Y1 | review | d7d9ad46}
\par
Drayver, utilita, firmware va translatorlar mavzusida "kompilyator" qatnashgan atama-ta'rif juftliklaridan to'g'risini toping.
\begin{enumerate}[label=\Alph*),leftmargin=9mm]
\item kompilyator - OT bilan aniq qurilma modeli o'rtasida buyruqlarni moslashtiruvchi dastur
\item kompilyator - butun manba kodini obyekt yoki bajariladigan kodga o'giruvchi translator
\item kompilyator - tizimga xizmat ko'rsatuvchi yordamchi dastur
\item kompilyator - qurilma ichidagi past darajali boshqaruv dasturi
\end{enumerate}
\begin{tcolorbox}[enhanced,breakable,colback=teal!4,colframe=teal!60!black,boxrule=0.45pt,arc=1.2mm,left=2.7mm,right=2.7mm,top=1.8mm,bottom=1.8mm,before skip=2mm,after skip=2mm]
{\bfseries\sffamily\color{teal}To'g'ri javob:} \textbf{B) kompilyator - butun manba kodini obyekt yoki bajariladigan kodga o'giruvchi translator}
\par\smallskip
{\bfseries Izoh:} kompilyator - butun manba kodini obyekt yoki bajariladigan kodga o'giruvchi translator. To'g'ri juftlik: kompilyator - butun manba kodini obyekt yoki bajariladigan kodga o'giruvchi translator. Manba: ICT8, PDF 106-110; Qo'shimcha sahifa raqami (qaysi PDFga tegishli ekani aniqlanmagan; darslik sahifasi deb tasdiqlanmagan): 72
\end{tcolorbox}
\vspace{2mm}\hrule height 0.2pt\vspace{3mm}

% DB_ID d9dde41a-5b31-3218-a2ec-ecee782e5156
\needspace{5\baselineskip}
{\bfseries\sffamily\color{ink} 567.} {scriptsizecolor{muted} Y1 | review | d9dde41a}
\par
Dasturiy ta'minotning tasnifi mavzusida "amaliy dastur" qatnashgan atama-ta'rif juftliklaridan to'g'risini toping.
\begin{enumerate}[label=\Alph*),leftmargin=9mm]
\item amaliy dastur - kompyuter resurslarini boshqarib, boshqa dasturlarga muhit yaratuvchi dastur
\item amaliy dastur - foydalanuvchining aniq vazifasini bajaruvchi dastur
\item amaliy dastur - tizimga xizmat ko'rsatuvchi yordamchi dastur
\item amaliy dastur - dastur kodini boshqa, odatda mashina bajaradigan shaklga o'giruvchi vosita
\end{enumerate}
\begin{tcolorbox}[enhanced,breakable,colback=teal!4,colframe=teal!60!black,boxrule=0.45pt,arc=1.2mm,left=2.7mm,right=2.7mm,top=1.8mm,bottom=1.8mm,before skip=2mm,after skip=2mm]
{\bfseries\sffamily\color{teal}To'g'ri javob:} \textbf{B) amaliy dastur - foydalanuvchining aniq vazifasini bajaruvchi dastur}
\par\smallskip
{\bfseries Izoh:} amaliy dastur - foydalanuvchining aniq vazifasini bajaruvchi dastur. To'g'ri juftlik: amaliy dastur - foydalanuvchining aniq vazifasini bajaruvchi dastur. Manba: CAM10/11, PDF 44-58; Qo'shimcha sahifa raqami (qaysi PDFga tegishli ekani aniqlanmagan; darslik sahifasi deb tasdiqlanmagan): 50
\end{tcolorbox}
\vspace{2mm}\hrule height 0.2pt\vspace{3mm}

% DB_ID d9fc2806-73e1-0974-ee51-5c3a4beb2845
\needspace{5\baselineskip}
{\bfseries\sffamily\color{ink} 568.} {scriptsizecolor{muted} Y1 | review | d9fc2806}
\par
Qiymat o'zgarsa natija va diagramma avtomatik yangilanishi kerak. Qaysi dastur turi mos?
\begin{enumerate}[label=\Alph*),leftmargin=9mm]
\item matn protsessori
\item audio muharrir
\item fayl arxivatori
\item elektron jadval
\end{enumerate}
\begin{tcolorbox}[enhanced,breakable,colback=teal!4,colframe=teal!60!black,boxrule=0.45pt,arc=1.2mm,left=2.7mm,right=2.7mm,top=1.8mm,bottom=1.8mm,before skip=2mm,after skip=2mm]
{\bfseries\sffamily\color{teal}To'g'ri javob:} \textbf{D) elektron jadval}
\par\smallskip
{\bfseries Izoh:} elektron jadval. Elektron jadval formulalarni qayta hisoblab, bog'langan diagrammani yangilaydi. Manba: CAM10/11, PDF 45-58; Qo'shimcha sahifa raqami (qaysi PDFga tegishli ekani aniqlanmagan; darslik sahifasi deb tasdiqlanmagan): 82
\end{tcolorbox}
\vspace{2mm}\hrule height 0.2pt\vspace{3mm}

% DB_ID e828a724-f7fd-845e-44b6-1a369ad84ebe
\needspace{5\baselineskip}
{\bfseries\sffamily\color{ink} 569.} {scriptsizecolor{muted} Y1 | review | e828a724}
\par
Dastur o'rnatish, yangilash, litsenziya va moslik mavzusida "moslik" qatnashgan atama-ta'rif juftliklaridan to'g'risini toping.
\begin{enumerate}[label=\Alph*),leftmargin=9mm]
\item moslik - dastur komponentlarini tizimda ishlashga tayyorlab joylashtirish
\item moslik - dastur va ro'yxatdan o'tgan komponentlarini nazoratli olib tashlash
\item moslik - mavjud versiyaga tuzatish yoki kichik yaxshilanish o'rnatish
\item moslik - qurilma, OT, dastur va formatlarning birgalikda ishlay olish xususiyati
\end{enumerate}
\begin{tcolorbox}[enhanced,breakable,colback=teal!4,colframe=teal!60!black,boxrule=0.45pt,arc=1.2mm,left=2.7mm,right=2.7mm,top=1.8mm,bottom=1.8mm,before skip=2mm,after skip=2mm]
{\bfseries\sffamily\color{teal}To'g'ri javob:} \textbf{D) moslik - qurilma, OT, dastur va formatlarning birgalikda ishlay olish xususiyati}
\par\smallskip
{\bfseries Izoh:} moslik - qurilma, OT, dastur va formatlarning birgalikda ishlay olish xususiyati. To'g'ri juftlik: moslik - qurilma, OT, dastur va formatlarning birgalikda ishlay olish xususiyati. Manba: CAM10/11, PDF 46-61; Qo'shimcha sahifa raqami (qaysi PDFga tegishli ekani aniqlanmagan; darslik sahifasi deb tasdiqlanmagan): 105
\end{tcolorbox}
\vspace{2mm}\hrule height 0.2pt\vspace{3mm}

% DB_ID fa42c06a-8c0d-8737-f3ef-6caeafaba7f0
\needspace{5\baselineskip}
{\bfseries\sffamily\color{ink} 570.} {scriptsizecolor{muted} Y1 | review | fa42c06a}
\par
Amaliy dasturlar va vazifaga mos tanlov bo'yicha qatorlar ichidan "DBMS" noto'g'ri izohlanganini aniqlang.
\begin{enumerate}[label=\Alph*),leftmargin=9mm]
\item elektron jadval - formulalar asosida hisoblash, model va diagramma yaratish dasturi
\item DTP - professional sahifa maketi va nashr tayyorlash dasturi
\item CAD - aniq texnik chizma va model yaratish dasturi
\item DBMS - tuzilmali matnli hujjat yaratish va formatlash dasturi
\end{enumerate}
\begin{tcolorbox}[enhanced,breakable,colback=teal!4,colframe=teal!60!black,boxrule=0.45pt,arc=1.2mm,left=2.7mm,right=2.7mm,top=1.8mm,bottom=1.8mm,before skip=2mm,after skip=2mm]
{\bfseries\sffamily\color{teal}To'g'ri javob:} \textbf{D) DBMS - tuzilmali matnli hujjat yaratish va formatlash dasturi}
\par\smallskip
{\bfseries Izoh:} DBMS - tuzilmali matnli hujjat yaratish va formatlash dasturi. DBMS uchun to'g'ri ta'rif: tuzilmali ma'lumotlar bazasini yaratish, so'rash va boshqarish tizimi. Manba: CAM8, PDF 118-150; Qo'shimcha sahifa raqami (qaysi PDFga tegishli ekani aniqlanmagan; darslik sahifasi deb tasdiqlanmagan): 79
\end{tcolorbox}
\vspace{2mm}\hrule height 0.2pt\vspace{3mm}

% DB_ID fa976909-fb8d-2946-82f2-77608b302f60
\needspace{5\baselineskip}
{\bfseries\sffamily\color{ink} 571.} {scriptsizecolor{muted} Y1 | review | fa976909}
\par
Drayver, utilita, firmware va translatorlar bo'yicha qatorlar ichidan "interpretator" noto'g'ri izohlanganini aniqlang.
\begin{enumerate}[label=\Alph*),leftmargin=9mm]
\item utilita - tizimga xizmat ko'rsatuvchi yordamchi dastur
\item firmware - qurilma ichidagi past darajali boshqaruv dasturi
\item kompilyator - butun manba kodini obyekt yoki bajariladigan kodga o'giruvchi translator
\item interpretator - OT bilan aniq qurilma modeli o'rtasida buyruqlarni moslashtiruvchi dastur
\end{enumerate}
\begin{tcolorbox}[enhanced,breakable,colback=teal!4,colframe=teal!60!black,boxrule=0.45pt,arc=1.2mm,left=2.7mm,right=2.7mm,top=1.8mm,bottom=1.8mm,before skip=2mm,after skip=2mm]
{\bfseries\sffamily\color{teal}To'g'ri javob:} \textbf{D) interpretator - OT bilan aniq qurilma modeli o'rtasida buyruqlarni moslashtiruvchi dastur}
\par\smallskip
{\bfseries Izoh:} interpretator - OT bilan aniq qurilma modeli o'rtasida buyruqlarni moslashtiruvchi dastur. interpretator uchun to'g'ri ta'rif: kodni bajarish paytida bosqichma-bosqich tarjima qiluvchi translator. Manba: CAM9, PDF 14-15; Qo'shimcha sahifa raqami (qaysi PDFga tegishli ekani aniqlanmagan; darslik sahifasi deb tasdiqlanmagan): 74
\end{tcolorbox}
\vspace{2mm}\hrule height 0.2pt\vspace{3mm}

% DB_ID faf1fa17-92f1-5eec-0546-8ae2c800a120
\needspace{5\baselineskip}
{\bfseries\sffamily\color{ink} 572.} {scriptsizecolor{muted} Y1 | review | faf1fa17}
\par
Amaliy dasturlar va vazifaga mos tanlov bo'yicha qatorlar ichidan "matn protsessori" noto'g'ri izohlanganini aniqlang.
\begin{enumerate}[label=\Alph*),leftmargin=9mm]
\item DBMS - tuzilmali ma'lumotlar bazasini yaratish, so'rash va boshqarish tizimi
\item matn protsessori - formulalar asosida hisoblash, model va diagramma yaratish dasturi
\item DTP - professional sahifa maketi va nashr tayyorlash dasturi
\item CAD - aniq texnik chizma va model yaratish dasturi
\end{enumerate}
\begin{tcolorbox}[enhanced,breakable,colback=teal!4,colframe=teal!60!black,boxrule=0.45pt,arc=1.2mm,left=2.7mm,right=2.7mm,top=1.8mm,bottom=1.8mm,before skip=2mm,after skip=2mm]
{\bfseries\sffamily\color{teal}To'g'ri javob:} \textbf{B) matn protsessori - formulalar asosida hisoblash, model va diagramma yaratish dasturi}
\par\smallskip
{\bfseries Izoh:} matn protsessori - formulalar asosida hisoblash, model va diagramma yaratish dasturi. matn protsessori uchun to'g'ri ta'rif: tuzilmali matnli hujjat yaratish va formatlash dasturi. Manba: CAM10/11, PDF 45-58; Qo'shimcha sahifa raqami (qaysi PDFga tegishli ekani aniqlanmagan; darslik sahifasi deb tasdiqlanmagan): 76
\end{tcolorbox}
\vspace{2mm}\hrule height 0.2pt\vspace{3mm}

% DB_ID 029b8f34-401d-2b1c-89d4-784f1581a060
\needspace{5\baselineskip}
{\bfseries\sffamily\color{ink} 573.} {scriptsizecolor{muted} Y1 | review | 029b8f34}
\par
Dasturiy ta'minotning tasnifi mavzusida: I. utilita - tizimga xizmat ko'rsatuvchi yordamchi dastur. II. tizimli dastur - tizimga xizmat ko'rsatuvchi yordamchi dastur. Mulohazalarni baholang.
\begin{enumerate}[label=\Alph*),leftmargin=9mm]
\item faqat I
\item faqat II
\item I va II
\item I ham, II ham emas
\end{enumerate}
\begin{tcolorbox}[enhanced,breakable,colback=teal!4,colframe=teal!60!black,boxrule=0.45pt,arc=1.2mm,left=2.7mm,right=2.7mm,top=1.8mm,bottom=1.8mm,before skip=2mm,after skip=2mm]
{\bfseries\sffamily\color{teal}To'g'ri javob:} \textbf{A) faqat I}
\par\smallskip
{\bfseries Izoh:} faqat I. Birinchi va ikkinchi mulohazaning ta'riflari alohida tekshirildi; natija: faqat I. Manba: ICT5, PDF 38-41; Qo'shimcha sahifa raqami (qaysi PDFga tegishli ekani aniqlanmagan; darslik sahifasi deb tasdiqlanmagan): 52
\end{tcolorbox}
\vspace{2mm}\hrule height 0.2pt\vspace{3mm}

% DB_ID 0c1d4092-5fae-32cb-a7db-3a23bf2f4a9a
\needspace{5\baselineskip}
{\bfseries\sffamily\color{ink} 574.} {scriptsizecolor{muted} Y1 | review | 0c1d4092}
\par
Minglab mijoz va buyurtmalar orasidan shart bo'yicha yozuv topish va hisobot chiqarish kerak. Eng mos dastur turi?
\begin{enumerate}[label=\Alph*),leftmargin=9mm]
\item rastr muharriri
\item media player
\item DBMS
\item oddiy kalkulyator
\end{enumerate}
\begin{tcolorbox}[enhanced,breakable,colback=teal!4,colframe=teal!60!black,boxrule=0.45pt,arc=1.2mm,left=2.7mm,right=2.7mm,top=1.8mm,bottom=1.8mm,before skip=2mm,after skip=2mm]
{\bfseries\sffamily\color{teal}To'g'ri javob:} \textbf{C) DBMS}
\par\smallskip
{\bfseries Izoh:} DBMS. DBMS tuzilmali yozuvlar, so'rovlar va hisobotlar uchun mo'ljallangan. Manba: CAM10/11, PDF 45-58; Qo'shimcha sahifa raqami (qaysi PDFga tegishli ekani aniqlanmagan; darslik sahifasi deb tasdiqlanmagan): 82
\end{tcolorbox}
\vspace{2mm}\hrule height 0.2pt\vspace{3mm}

% DB_ID 27f33b3a-6816-ab59-a9d4-5329d1977e7d
\needspace{5\baselineskip}
{\bfseries\sffamily\color{ink} 575.} {scriptsizecolor{muted} Y1 | review | 27f33b3a}
\par
Amaliy dasturlar va vazifaga mos tanlov mavzusida: I. CAD - aniq texnik chizma va model yaratish dasturi. II. matn protsessori - tuzilmali matnli hujjat yaratish va formatlash dasturi. Mulohazalarni baholang.
\begin{enumerate}[label=\Alph*),leftmargin=9mm]
\item faqat I
\item faqat II
\item I va II
\item I ham, II ham emas
\end{enumerate}
\begin{tcolorbox}[enhanced,breakable,colback=teal!4,colframe=teal!60!black,boxrule=0.45pt,arc=1.2mm,left=2.7mm,right=2.7mm,top=1.8mm,bottom=1.8mm,before skip=2mm,after skip=2mm]
{\bfseries\sffamily\color{teal}To'g'ri javob:} \textbf{C) I va II}
\par\smallskip
{\bfseries Izoh:} I va II. Birinchi va ikkinchi mulohazaning ta'riflari alohida tekshirildi; natija: I va II. Manba: ICT6, PDF 62-126; Qo'shimcha sahifa raqami (qaysi PDFga tegishli ekani aniqlanmagan; darslik sahifasi deb tasdiqlanmagan): 81
\end{tcolorbox}
\vspace{2mm}\hrule height 0.2pt\vspace{3mm}

% DB_ID 361ab604-929d-f6b2-1811-90327b85f6cd
\needspace{5\baselineskip}
{\bfseries\sffamily\color{ink} 576.} {scriptsizecolor{muted} Y1 | review | 361ab604}
\par
Dasturchi yozgan manba kodini kompyuter bajarishi mumkin bo'lgan ko'rinishga o'tkazish kerak. Bu ish uchun qaysi dastur turidan foydalaniladi?
\begin{enumerate}[label=\Alph*),leftmargin=9mm]
\item Translator: kodni boshqa bajariladigan yoki oraliq ko'rinishga tarjima qiladi
\item Fayl menejeri: manba koddagi har bir operatorni protsessor buyrug'iga o'giradi
\item Matn muharriri: manba kodni avtomatik tekshirib, faqat disk sektorlarini boshqaradi
\item Qurilma drayveri: dasturlash tilidagi barcha ko'rsatmalarni foydalanuvchiga tarjima qiladi
\end{enumerate}
\begin{tcolorbox}[enhanced,breakable,colback=teal!4,colframe=teal!60!black,boxrule=0.45pt,arc=1.2mm,left=2.7mm,right=2.7mm,top=1.8mm,bottom=1.8mm,before skip=2mm,after skip=2mm]
{\bfseries\sffamily\color{teal}To'g'ri javob:} \textbf{A) Translator: kodni boshqa bajariladigan yoki oraliq ko'rinishga tarjima qiladi}
\par\smallskip
{\bfseries Izoh:} Darslik kompilyator va interpretatorni tarjimon dasturlar sifatida tushuntiradi: ular manba kodni boshqa ko'rinishga o'tkazadi yoki bajarilishida talqin qiladi. [Manba: Cambridge+ 10-11-sinf, PDF 45-46.]
\end{tcolorbox}
\vspace{2mm}\hrule height 0.2pt\vspace{3mm}

% DB_ID 3c155ee9-cfd8-6a20-a802-1013b56287d8
\needspace{5\baselineskip}
{\bfseries\sffamily\color{ink} 577.} {scriptsizecolor{muted} Y1 | review | 3c155ee9}
\par
Drayver, utilita, firmware va translatorlar mavzusida: I. drayver - OT bilan aniq qurilma modeli o'rtasida buyruqlarni moslashtiruvchi dastur. II. utilita - tizimga xizmat ko'rsatuvchi yordamchi dastur. Mulohazalarni baholang.
\begin{enumerate}[label=\Alph*),leftmargin=9mm]
\item I va II
\item faqat I
\item faqat II
\item I ham, II ham emas
\end{enumerate}
\begin{tcolorbox}[enhanced,breakable,colback=teal!4,colframe=teal!60!black,boxrule=0.45pt,arc=1.2mm,left=2.7mm,right=2.7mm,top=1.8mm,bottom=1.8mm,before skip=2mm,after skip=2mm]
{\bfseries\sffamily\color{teal}To'g'ri javob:} \textbf{A) I va II}
\par\smallskip
{\bfseries Izoh:} I va II. Birinchi va ikkinchi mulohazaning ta'riflari alohida tekshirildi; natija: I va II. Manba: CAM9, PDF 14-15; Qo'shimcha sahifa raqami (qaysi PDFga tegishli ekani aniqlanmagan; darslik sahifasi deb tasdiqlanmagan): 70
\end{tcolorbox}
\vspace{2mm}\hrule height 0.2pt\vspace{3mm}

% DB_ID 4209d3d1-daa4-6f20-d2ec-55277a7a27b0
\needspace{5\baselineskip}
{\bfseries\sffamily\color{ink} 578.} {scriptsizecolor{muted} Y1 | review | 4209d3d1}
\par
Dasturiy ta'minotning tasnifi mavzusida: I. tizimli dastur - kompyuter resurslarini boshqarib, boshqa dasturlarga muhit yaratuvchi dastur. II. amaliy dastur - foydalanuvchining aniq vazifasini bajaruvchi dastur. Mulohazalarni baholang.
\begin{enumerate}[label=\Alph*),leftmargin=9mm]
\item faqat I
\item faqat II
\item I va II
\item I ham, II ham emas
\end{enumerate}
\begin{tcolorbox}[enhanced,breakable,colback=teal!4,colframe=teal!60!black,boxrule=0.45pt,arc=1.2mm,left=2.7mm,right=2.7mm,top=1.8mm,bottom=1.8mm,before skip=2mm,after skip=2mm]
{\bfseries\sffamily\color{teal}To'g'ri javob:} \textbf{C) I va II}
\par\smallskip
{\bfseries Izoh:} I va II. Birinchi va ikkinchi mulohazaning ta'riflari alohida tekshirildi; natija: I va II. Manba: CAM9, PDF 11-15; Qo'shimcha sahifa raqami (qaysi PDFga tegishli ekani aniqlanmagan; darslik sahifasi deb tasdiqlanmagan): 50
\end{tcolorbox}
\vspace{2mm}\hrule height 0.2pt\vspace{3mm}

% DB_ID 58378c4b-b18d-11f2-a7a6-5809f6b70d92
\needspace{5\baselineskip}
{\bfseries\sffamily\color{ink} 579.} {scriptsizecolor{muted} Y1 | review | 58378c4b}
\par
Dastur o'rnatish, yangilash, litsenziya va moslik mavzusida: I. open-source - manba kodi litsenziya asosida ochiq dastur. II. installation - dastur komponentlarini tizimda ishlashga tayyorlab joylashtirish. Mulohazalarni baholang.
\begin{enumerate}[label=\Alph*),leftmargin=9mm]
\item faqat I
\item faqat II
\item I va II
\item I ham, II ham emas
\end{enumerate}
\begin{tcolorbox}[enhanced,breakable,colback=teal!4,colframe=teal!60!black,boxrule=0.45pt,arc=1.2mm,left=2.7mm,right=2.7mm,top=1.8mm,bottom=1.8mm,before skip=2mm,after skip=2mm]
{\bfseries\sffamily\color{teal}To'g'ri javob:} \textbf{C) I va II}
\par\smallskip
{\bfseries Izoh:} I va II. Birinchi va ikkinchi mulohazaning ta'riflari alohida tekshirildi; natija: I va II. Manba: MAR, PDF 10-16; Qo'shimcha sahifa raqami (qaysi PDFga tegishli ekani aniqlanmagan; darslik sahifasi deb tasdiqlanmagan): 106
\end{tcolorbox}
\vspace{2mm}\hrule height 0.2pt\vspace{3mm}

% DB_ID 5ee966f2-6796-a038-58df-b0a3420557f5
\needspace{5\baselineskip}
{\bfseries\sffamily\color{ink} 580.} {scriptsizecolor{muted} Y1 | review | 5ee966f2}
\par
Dastur uchun minimal talab 8 GB RAM, tavsiya etilgani 16 GB. Kompyuterda 8 GB bor. Eng aniq xulosa qaysi?
\begin{enumerate}[label=\Alph*),leftmargin=9mm]
\item dastur mutlaqo o'rnatilmaydi
\item dastur albatta eng tez ishlaydi
\item ishga tushishi mumkin, lekin maqbul ishlashi kafolatlanmaydi
\item RAM talabi fayl formatidir
\end{enumerate}
\begin{tcolorbox}[enhanced,breakable,colback=teal!4,colframe=teal!60!black,boxrule=0.45pt,arc=1.2mm,left=2.7mm,right=2.7mm,top=1.8mm,bottom=1.8mm,before skip=2mm,after skip=2mm]
{\bfseries\sffamily\color{teal}To'g'ri javob:} \textbf{C) ishga tushishi mumkin, lekin maqbul ishlashi kafolatlanmaydi}
\par\smallskip
{\bfseries Izoh:} ishga tushishi mumkin, lekin maqbul ishlashi kafolatlanmaydi. Minimal talab ishga tushish chegarasi, tavsiya etilgani maqbul ishlashni ko'zlaydi. Manba: MAR, PDF 10-16; Qo'shimcha sahifa raqami (qaysi PDFga tegishli ekani aniqlanmagan; darslik sahifasi deb tasdiqlanmagan): 107
\end{tcolorbox}
\vspace{2mm}\hrule height 0.2pt\vspace{3mm}

% DB_ID 6456fa39-8a94-ceb2-6b47-33da1087bc27
\needspace{5\baselineskip}
{\bfseries\sffamily\color{ink} 581.} {scriptsizecolor{muted} Y1 | review | 6456fa39}
\par
Drayver, utilita, firmware va translatorlar mavzusida: I. firmware - qurilma ichidagi past darajali boshqaruv dasturi. II. drayver - qurilma ichidagi past darajali boshqaruv dasturi. Mulohazalarni baholang.
\begin{enumerate}[label=\Alph*),leftmargin=9mm]
\item faqat II
\item I va II
\item faqat I
\item I ham, II ham emas
\end{enumerate}
\begin{tcolorbox}[enhanced,breakable,colback=teal!4,colframe=teal!60!black,boxrule=0.45pt,arc=1.2mm,left=2.7mm,right=2.7mm,top=1.8mm,bottom=1.8mm,before skip=2mm,after skip=2mm]
{\bfseries\sffamily\color{teal}To'g'ri javob:} \textbf{C) faqat I}
\par\smallskip
{\bfseries Izoh:} faqat I. Birinchi va ikkinchi mulohazaning ta'riflari alohida tekshirildi; natija: faqat I. Manba: MAR, PDF 10-16; Qo'shimcha sahifa raqami (qaysi PDFga tegishli ekani aniqlanmagan; darslik sahifasi deb tasdiqlanmagan): 72
\end{tcolorbox}
\vspace{2mm}\hrule height 0.2pt\vspace{3mm}

% DB_ID 671c1e1d-01f7-79b2-c6c2-49798c2e1fd8
\needspace{5\baselineskip}
{\bfseries\sffamily\color{ink} 582.} {scriptsizecolor{muted} Y1 | review | 671c1e1d}
\par
Dastur o'rnatish, yangilash, litsenziya va moslik bo'yicha qatorlar ichidan "moslik" noto'g'ri izohlanganini aniqlang.
\begin{enumerate}[label=\Alph*),leftmargin=9mm]
\item moslik - dastur komponentlarini tizimda ishlashga tayyorlab joylashtirish
\item uninstall - dastur va ro'yxatdan o'tgan komponentlarini nazoratli olib tashlash
\item update - mavjud versiyaga tuzatish yoki kichik yaxshilanish o'rnatish
\item open-source - manba kodi litsenziya asosida ochiq dastur
\end{enumerate}
\begin{tcolorbox}[enhanced,breakable,colback=teal!4,colframe=teal!60!black,boxrule=0.45pt,arc=1.2mm,left=2.7mm,right=2.7mm,top=1.8mm,bottom=1.8mm,before skip=2mm,after skip=2mm]
{\bfseries\sffamily\color{teal}To'g'ri javob:} \textbf{A) moslik - dastur komponentlarini tizimda ishlashga tayyorlab joylashtirish}
\par\smallskip
{\bfseries Izoh:} moslik - dastur komponentlarini tizimda ishlashga tayyorlab joylashtirish. moslik uchun to'g'ri ta'rif: qurilma, OT, dastur va formatlarning birgalikda ishlay olish xususiyati. Manba: CAM10/11, PDF 46-61; Qo'shimcha sahifa raqami (qaysi PDFga tegishli ekani aniqlanmagan; darslik sahifasi deb tasdiqlanmagan): 105
\end{tcolorbox}
\vspace{2mm}\hrule height 0.2pt\vspace{3mm}

% DB_ID 6de72ba1-8401-30d6-23a7-09547a2bc0b3
\needspace{5\baselineskip}
{\bfseries\sffamily\color{ink} 583.} {scriptsizecolor{muted} Y1 | review | 6de72ba1}
\par
Amaliy dasturlar va vazifaga mos tanlov mavzusida: I. DBMS - tuzilmali ma'lumotlar bazasini yaratish, so'rash va boshqarish tizimi. II. matn protsessori - tuzilmali ma'lumotlar bazasini yaratish, so'rash va boshqarish tizimi. Mulohazalarni baholang.
\begin{enumerate}[label=\Alph*),leftmargin=9mm]
\item faqat I
\item faqat II
\item I va II
\item I ham, II ham emas
\end{enumerate}
\begin{tcolorbox}[enhanced,breakable,colback=teal!4,colframe=teal!60!black,boxrule=0.45pt,arc=1.2mm,left=2.7mm,right=2.7mm,top=1.8mm,bottom=1.8mm,before skip=2mm,after skip=2mm]
{\bfseries\sffamily\color{teal}To'g'ri javob:} \textbf{A) faqat I}
\par\smallskip
{\bfseries Izoh:} faqat I. Birinchi va ikkinchi mulohazaning ta'riflari alohida tekshirildi; natija: faqat I. Manba: CAM8, PDF 118-150; Qo'shimcha sahifa raqami (qaysi PDFga tegishli ekani aniqlanmagan; darslik sahifasi deb tasdiqlanmagan): 79
\end{tcolorbox}
\vspace{2mm}\hrule height 0.2pt\vspace{3mm}

% DB_ID 788f720b-e5e4-6c3d-3b44-db25a2ba859e
\needspace{5\baselineskip}
{\bfseries\sffamily\color{ink} 584.} {scriptsizecolor{muted} Y1 | review | 788f720b}
\par
Manba kodining barcha modullari alohida obyekt kodga aylantirildi. Ularni bitta bajariladigan faylga nima birlashtiradi?
\begin{enumerate}[label=\Alph*),leftmargin=9mm]
\item OCR
\item linker
\item defragmenter
\item device driver
\end{enumerate}
\begin{tcolorbox}[enhanced,breakable,colback=teal!4,colframe=teal!60!black,boxrule=0.45pt,arc=1.2mm,left=2.7mm,right=2.7mm,top=1.8mm,bottom=1.8mm,before skip=2mm,after skip=2mm]
{\bfseries\sffamily\color{teal}To'g'ri javob:} \textbf{B) linker}
\par\smallskip
{\bfseries Izoh:} linker. Linker obyekt modullari va kutubxonalarni bog'laydi. Manba: CAM9, PDF 14-15; Qo'shimcha sahifa raqami (qaysi PDFga tegishli ekani aniqlanmagan; darslik sahifasi deb tasdiqlanmagan): 75
\end{tcolorbox}
\vspace{2mm}\hrule height 0.2pt\vspace{3mm}

% DB_ID 9122555a-95b2-5123-b1d8-e8c1c3ca108f
\needspace{5\baselineskip}
{\bfseries\sffamily\color{ink} 585.} {scriptsizecolor{muted} Y1 | review | 9122555a}
\par
Amaliy dasturlar va vazifaga mos tanlov bo'yicha qatorlar ichidan "DTP" noto'g'ri izohlanganini aniqlang.
\begin{enumerate}[label=\Alph*),leftmargin=9mm]
\item DTP - tuzilmali matnli hujjat yaratish va formatlash dasturi
\item elektron jadval - formulalar asosida hisoblash, model va diagramma yaratish dasturi
\item DBMS - tuzilmali ma'lumotlar bazasini yaratish, so'rash va boshqarish tizimi
\item CAD - aniq texnik chizma va model yaratish dasturi
\end{enumerate}
\begin{tcolorbox}[enhanced,breakable,colback=teal!4,colframe=teal!60!black,boxrule=0.45pt,arc=1.2mm,left=2.7mm,right=2.7mm,top=1.8mm,bottom=1.8mm,before skip=2mm,after skip=2mm]
{\bfseries\sffamily\color{teal}To'g'ri javob:} \textbf{A) DTP - tuzilmali matnli hujjat yaratish va formatlash dasturi}
\par\smallskip
{\bfseries Izoh:} DTP - tuzilmali matnli hujjat yaratish va formatlash dasturi. DTP uchun to'g'ri ta'rif: professional sahifa maketi va nashr tayyorlash dasturi. Manba: CAM5, PDF 45-138; Qo'shimcha sahifa raqami (qaysi PDFga tegishli ekani aniqlanmagan; darslik sahifasi deb tasdiqlanmagan): 80
\end{tcolorbox}
\vspace{2mm}\hrule height 0.2pt\vspace{3mm}

% DB_ID a037a37e-97c8-39ca-ddbf-7e2266f58fcf
\needspace{5\baselineskip}
{\bfseries\sffamily\color{ink} 586.} {scriptsizecolor{muted} Y1 | review | a037a37e}
\par
Dastur o'rnatish, yangilash, litsenziya va moslik mavzusida: I. installation - dastur komponentlarini tizimda ishlashga tayyorlab joylashtirish. II. uninstall - dastur va ro'yxatdan o'tgan komponentlarini nazoratli olib tashlash. Mulohazalarni baholang.
\begin{enumerate}[label=\Alph*),leftmargin=9mm]
\item faqat I
\item faqat II
\item I va II
\item I ham, II ham emas
\end{enumerate}
\begin{tcolorbox}[enhanced,breakable,colback=teal!4,colframe=teal!60!black,boxrule=0.45pt,arc=1.2mm,left=2.7mm,right=2.7mm,top=1.8mm,bottom=1.8mm,before skip=2mm,after skip=2mm]
{\bfseries\sffamily\color{teal}To'g'ri javob:} \textbf{C) I va II}
\par\smallskip
{\bfseries Izoh:} I va II. Birinchi va ikkinchi mulohazaning ta'riflari alohida tekshirildi; natija: I va II. Manba: MAR, PDF 10-16; Qo'shimcha sahifa raqami (qaysi PDFga tegishli ekani aniqlanmagan; darslik sahifasi deb tasdiqlanmagan): 102
\end{tcolorbox}
\vspace{2mm}\hrule height 0.2pt\vspace{3mm}

% DB_ID b5eadaf7-afdf-cfe7-6a2f-a8091e332e40
\needspace{5\baselineskip}
{\bfseries\sffamily\color{ink} 587.} {scriptsizecolor{muted} Y1 | review | b5eadaf7}
\par
Amaliy dasturlar va vazifaga mos tanlov mavzusida: I. matn protsessori - tuzilmali matnli hujjat yaratish va formatlash dasturi. II. elektron jadval - formulalar asosida hisoblash, model va diagramma yaratish dasturi. Mulohazalarni baholang.
\begin{enumerate}[label=\Alph*),leftmargin=9mm]
\item faqat I
\item faqat II
\item I va II
\item I ham, II ham emas
\end{enumerate}
\begin{tcolorbox}[enhanced,breakable,colback=teal!4,colframe=teal!60!black,boxrule=0.45pt,arc=1.2mm,left=2.7mm,right=2.7mm,top=1.8mm,bottom=1.8mm,before skip=2mm,after skip=2mm]
{\bfseries\sffamily\color{teal}To'g'ri javob:} \textbf{C) I va II}
\par\smallskip
{\bfseries Izoh:} I va II. Birinchi va ikkinchi mulohazaning ta'riflari alohida tekshirildi; natija: I va II. Manba: CAM10/11, PDF 45-58; Qo'shimcha sahifa raqami (qaysi PDFga tegishli ekani aniqlanmagan; darslik sahifasi deb tasdiqlanmagan): 77
\end{tcolorbox}
\vspace{2mm}\hrule height 0.2pt\vspace{3mm}

% DB_ID b771f974-7b5e-c611-1ed5-c332fbeb63fe
\needspace{5\baselineskip}
{\bfseries\sffamily\color{ink} 588.} {scriptsizecolor{muted} Y1 | review | b771f974}
\par
Dastur o'rnatish, yangilash, litsenziya va moslik mavzusida: I. update - mavjud versiyaga tuzatish yoki kichik yaxshilanish o'rnatish. II. installation - mavjud versiyaga tuzatish yoki kichik yaxshilanish o'rnatish. Mulohazalarni baholang.
\begin{enumerate}[label=\Alph*),leftmargin=9mm]
\item faqat I
\item faqat II
\item I va II
\item I ham, II ham emas
\end{enumerate}
\begin{tcolorbox}[enhanced,breakable,colback=teal!4,colframe=teal!60!black,boxrule=0.45pt,arc=1.2mm,left=2.7mm,right=2.7mm,top=1.8mm,bottom=1.8mm,before skip=2mm,after skip=2mm]
{\bfseries\sffamily\color{teal}To'g'ri javob:} \textbf{A) faqat I}
\par\smallskip
{\bfseries Izoh:} faqat I. Birinchi va ikkinchi mulohazaning ta'riflari alohida tekshirildi; natija: faqat I. Manba: ICT8, PDF 106-110; Qo'shimcha sahifa raqami (qaysi PDFga tegishli ekani aniqlanmagan; darslik sahifasi deb tasdiqlanmagan): 104
\end{tcolorbox}
\vspace{2mm}\hrule height 0.2pt\vspace{3mm}

% DB_ID bd004add-7c32-7b3e-1cc2-403a277cc907
\needspace{5\baselineskip}
{\bfseries\sffamily\color{ink} 589.} {scriptsizecolor{muted} Y1 | review | bd004add}
\par
Yangi printer OTda ko'rinadi, ammo maxsus dupleks funksiyasi ishlamaydi. Eng ehtimoliy yechim qaysi?
\begin{enumerate}[label=\Alph*),leftmargin=9mm]
\item modelga mos drayverni o'rnatish
\item fayl kengaytmasini almashtirish
\item RAMni formatlash
\item papkani arxivlash
\end{enumerate}
\begin{tcolorbox}[enhanced,breakable,colback=teal!4,colframe=teal!60!black,boxrule=0.45pt,arc=1.2mm,left=2.7mm,right=2.7mm,top=1.8mm,bottom=1.8mm,before skip=2mm,after skip=2mm]
{\bfseries\sffamily\color{teal}To'g'ri javob:} \textbf{A) modelga mos drayverni o'rnatish}
\par\smallskip
{\bfseries Izoh:} modelga mos drayverni o'rnatish. Umumiy drayver asosiy ishlashni berishi, maxsus drayver qo'shimcha funksiyani ochishi mumkin. Manba: CAM9, PDF 14-15; Qo'shimcha sahifa raqami (qaysi PDFga tegishli ekani aniqlanmagan; darslik sahifasi deb tasdiqlanmagan): 75
\end{tcolorbox}
\vspace{2mm}\hrule height 0.2pt\vspace{3mm}

% DB_ID c8a7e590-401c-a343-e9d8-cfb82cf7bd35
\needspace{5\baselineskip}
{\bfseries\sffamily\color{ink} 590.} {scriptsizecolor{muted} Y1 | review | c8a7e590}
\par
Fayllarni siqish va mustaqil zaxira nusxalarini boshqarish uchun qaysi toifadagi dastur tanlanadi?
\begin{enumerate}[label=\Alph*),leftmargin=9mm]
\item Xizmat utilitalari ma'lumot siqish, zaxiralash va tizimga xizmat qilishga mos
\item Qurilma drayverlari faqat internetdagi sahifalarni tahrirlash va barcha fayllarning formatini bir xil usulda o'zgartirish uchun xizmat qiladi
\item Grafik protsessorlar tashqi xotiradagi barcha fayllarni alohida jadval bo'lib avtomatik siqish va turkumlash uchun mo'ljallangan
\item Matn protsessori OT resurslarini bevosita boshqarib, diskdagi barcha fayllarning zaxira nusxasini avtomatik yaratuvchi yagona dasturdir
\end{enumerate}
\begin{tcolorbox}[enhanced,breakable,colback=teal!4,colframe=teal!60!black,boxrule=0.45pt,arc=1.2mm,left=2.7mm,right=2.7mm,top=1.8mm,bottom=1.8mm,before skip=2mm,after skip=2mm]
{\bfseries\sffamily\color{teal}To'g'ri javob:} \textbf{A) Xizmat utilitalari ma'lumot siqish, zaxiralash va tizimga xizmat qilishga mos}
\par\smallskip
{\bfseries Izoh:} Ma'lumotni siqish va zaxiralash xizmat dasturlari yoki utilitalarning vazifalari jumlasiga kiradi. [Manba: Cambridge+ 10-11-sinf PDF 58-59; ICT 10-sinf PDF 28-29.]
\end{tcolorbox}
\vspace{2mm}\hrule height 0.2pt\vspace{3mm}

% DB_ID d502f14c-36f4-4263-6732-99707e082a7d
\needspace{5\baselineskip}
{\bfseries\sffamily\color{ink} 591.} {scriptsizecolor{muted} Y1 | review | d502f14c}
\par
O'quvchi matnli referatni tahrirlamoqda, laboratoriya mudiri esa kompyuterning qurilma resurslarini boshqarmoqda. Qaysi juftlik dasturlar rolini farqlaydi?
\begin{enumerate}[label=\Alph*),leftmargin=9mm]
\item Matn muharriri tizimli dastur, operatsion tizim esa rasm chizish vositasi
\item Matn muharriri qurilma drayveri, operatsion tizim esa mustaqil jadval muharriri
\item Matn muharriri amaliy dastur, operatsion tizim esa tizimli dastur
\item Matn muharriri fayl tizimi, operatsion tizim esa faqat printer utilitasi
\end{enumerate}
\begin{tcolorbox}[enhanced,breakable,colback=teal!4,colframe=teal!60!black,boxrule=0.45pt,arc=1.2mm,left=2.7mm,right=2.7mm,top=1.8mm,bottom=1.8mm,before skip=2mm,after skip=2mm]
{\bfseries\sffamily\color{teal}To'g'ri javob:} \textbf{C) Matn muharriri amaliy dastur, operatsion tizim esa tizimli dastur}
\par\smallskip
{\bfseries Izoh:} Amaliy dastur foydalanuvchining bevosita ishini bajaradi; operatsion tizim apparat va dastur resurslarini boshqaruvchi tizimli dasturiy ta'minotdir. [Manba: Cambridge+ 9-sinf, PDF 9-12.]
\end{tcolorbox}
\vspace{2mm}\hrule height 0.2pt\vspace{3mm}

% DB_ID d58dafb7-59bf-03b6-033c-8143c695204e
\needspace{5\baselineskip}
{\bfseries\sffamily\color{ink} 592.} {scriptsizecolor{muted} Y1 | review | d58dafb7}
\par
Mavjud dasturga faqat xavfsizlik xatosini tuzatuvchi kichik paket o'rnatildi. Ushbu amal nimani anglatadi?
\begin{enumerate}[label=\Alph*),leftmargin=9mm]
\item Dasturdagi barcha foydalanuvchi ma'lumotlarini mexanik diskning alohida sektorlariga doimiy yozish
\item Dastur manba kodini istalgan holda bir xil ko'rinishdagi yangi apparat qurilmasiga aylantirish
\item Mavjud dastur uchun xatoni tuzatuvchi yaxshilanish paketini o'rnatish
\item Dasturni butunlay olib tashlab, uning o'rniga operatsion tizimning barcha qurilma drayverlarini o'rnatish
\end{enumerate}
\begin{tcolorbox}[enhanced,breakable,colback=teal!4,colframe=teal!60!black,boxrule=0.45pt,arc=1.2mm,left=2.7mm,right=2.7mm,top=1.8mm,bottom=1.8mm,before skip=2mm,after skip=2mm]
{\bfseries\sffamily\color{teal}To'g'ri javob:} \textbf{C) Mavjud dastur uchun xatoni tuzatuvchi yaxshilanish paketini o'rnatish}
\par\smallskip
{\bfseries Izoh:} Darslikda mavjud dasturdagi nosozlikni tuzatish yoki kichik xususiyat qo'shish yaxshilash sifatida keltiriladi. [Manba: Cambridge+ 8-sinf PDF 106.]
\end{tcolorbox}
\vspace{2mm}\hrule height 0.2pt\vspace{3mm}

% DB_ID d89eca44-295a-e42c-a6a2-c8e1f5a50bfe
\needspace{5\baselineskip}
{\bfseries\sffamily\color{ink} 593.} {scriptsizecolor{muted} Y1 | review | d89eca44}
\par
Amaliy dasturlar va vazifaga mos tanlov bo'yicha qatorlar ichidan "elektron jadval" noto'g'ri izohlanganini aniqlang.
\begin{enumerate}[label=\Alph*),leftmargin=9mm]
\item DBMS - tuzilmali ma'lumotlar bazasini yaratish, so'rash va boshqarish tizimi
\item DTP - professional sahifa maketi va nashr tayyorlash dasturi
\item elektron jadval - tuzilmali matnli hujjat yaratish va formatlash dasturi
\item CAD - aniq texnik chizma va model yaratish dasturi
\end{enumerate}
\begin{tcolorbox}[enhanced,breakable,colback=teal!4,colframe=teal!60!black,boxrule=0.45pt,arc=1.2mm,left=2.7mm,right=2.7mm,top=1.8mm,bottom=1.8mm,before skip=2mm,after skip=2mm]
{\bfseries\sffamily\color{teal}To'g'ri javob:} \textbf{C) elektron jadval - tuzilmali matnli hujjat yaratish va formatlash dasturi}
\par\smallskip
{\bfseries Izoh:} elektron jadval - tuzilmali matnli hujjat yaratish va formatlash dasturi. elektron jadval uchun to'g'ri ta'rif: formulalar asosida hisoblash, model va diagramma yaratish dasturi. Manba: CAM6, PDF 34-205; Qo'shimcha sahifa raqami (qaysi PDFga tegishli ekani aniqlanmagan; darslik sahifasi deb tasdiqlanmagan): 78
\end{tcolorbox}
\vspace{2mm}\hrule height 0.2pt\vspace{3mm}

% DB_ID e8fe00dc-22ab-2707-782f-dc3e7a9c1aca
\needspace{5\baselineskip}
{\bfseries\sffamily\color{ink} 594.} {scriptsizecolor{muted} Y1 | review | e8fe00dc}
\par
Drayver, utilita, firmware va translatorlar bo'yicha qatorlar ichidan "kompilyator" noto'g'ri izohlanganini aniqlang.
\begin{enumerate}[label=\Alph*),leftmargin=9mm]
\item utilita - tizimga xizmat ko'rsatuvchi yordamchi dastur
\item firmware - qurilma ichidagi past darajali boshqaruv dasturi
\item kompilyator - OT bilan aniq qurilma modeli o'rtasida buyruqlarni moslashtiruvchi dastur
\item interpretator - kodni bajarish paytida bosqichma-bosqich tarjima qiluvchi translator
\end{enumerate}
\begin{tcolorbox}[enhanced,breakable,colback=teal!4,colframe=teal!60!black,boxrule=0.45pt,arc=1.2mm,left=2.7mm,right=2.7mm,top=1.8mm,bottom=1.8mm,before skip=2mm,after skip=2mm]
{\bfseries\sffamily\color{teal}To'g'ri javob:} \textbf{C) kompilyator - OT bilan aniq qurilma modeli o'rtasida buyruqlarni moslashtiruvchi dastur}
\par\smallskip
{\bfseries Izoh:} kompilyator - OT bilan aniq qurilma modeli o'rtasida buyruqlarni moslashtiruvchi dastur. kompilyator uchun to'g'ri ta'rif: butun manba kodini obyekt yoki bajariladigan kodga o'giruvchi translator. Manba: ICT8, PDF 106-110; Qo'shimcha sahifa raqami (qaysi PDFga tegishli ekani aniqlanmagan; darslik sahifasi deb tasdiqlanmagan): 73
\end{tcolorbox}
\vspace{2mm}\hrule height 0.2pt\vspace{3mm}

% DB_ID e9ce3df9-c9c8-8934-73d6-9681563cb30a
\needspace{5\baselineskip}
{\bfseries\sffamily\color{ink} 595.} {scriptsizecolor{muted} Y1 | review | e9ce3df9}
\par
Drayver, utilita, firmware va translatorlar mavzusida: I. interpretator - kodni bajarish paytida bosqichma-bosqich tarjima qiluvchi translator. II. drayver - OT bilan aniq qurilma modeli o'rtasida buyruqlarni moslashtiruvchi dastur. Mulohazalarni baholang.
\begin{enumerate}[label=\Alph*),leftmargin=9mm]
\item I va II
\item faqat I
\item faqat II
\item I ham, II ham emas
\end{enumerate}
\begin{tcolorbox}[enhanced,breakable,colback=teal!4,colframe=teal!60!black,boxrule=0.45pt,arc=1.2mm,left=2.7mm,right=2.7mm,top=1.8mm,bottom=1.8mm,before skip=2mm,after skip=2mm]
{\bfseries\sffamily\color{teal}To'g'ri javob:} \textbf{A) I va II}
\par\smallskip
{\bfseries Izoh:} I va II. Birinchi va ikkinchi mulohazaning ta'riflari alohida tekshirildi; natija: I va II. Manba: CAM9, PDF 14-15; Qo'shimcha sahifa raqami (qaysi PDFga tegishli ekani aniqlanmagan; darslik sahifasi deb tasdiqlanmagan): 74
\end{tcolorbox}
\vspace{2mm}\hrule height 0.2pt\vspace{3mm}

% DB_ID efb1fae0-ee6a-8aa0-6f77-3d416e598392
\needspace{5\baselineskip}
{\bfseries\sffamily\color{ink} 596.} {scriptsizecolor{muted} Y1 | review | efb1fae0}
\par
Dasturiy ta'minotning tasnifi mavzusida: I. umumiy maqsadli dastur - ko'p foydalanuvchi va vazifalarga mo'ljallangan tayyor dastur. II. tizimli dastur - kompyuter resurslarini boshqarib, boshqa dasturlarga muhit yaratuvchi dastur. Mulohazalarni baholang.
\begin{enumerate}[label=\Alph*),leftmargin=9mm]
\item faqat I
\item faqat II
\item I va II
\item I ham, II ham emas
\end{enumerate}
\begin{tcolorbox}[enhanced,breakable,colback=teal!4,colframe=teal!60!black,boxrule=0.45pt,arc=1.2mm,left=2.7mm,right=2.7mm,top=1.8mm,bottom=1.8mm,before skip=2mm,after skip=2mm]
{\bfseries\sffamily\color{teal}To'g'ri javob:} \textbf{C) I va II}
\par\smallskip
{\bfseries Izoh:} I va II. Birinchi va ikkinchi mulohazaning ta'riflari alohida tekshirildi; natija: I va II. Manba: CAM9, PDF 11-15; Qo'shimcha sahifa raqami (qaysi PDFga tegishli ekani aniqlanmagan; darslik sahifasi deb tasdiqlanmagan): 54
\end{tcolorbox}
\vspace{2mm}\hrule height 0.2pt\vspace{3mm}

% DB_ID f06b67da-1050-c625-27c4-da8e9d829b85
\needspace{5\baselineskip}
{\bfseries\sffamily\color{ink} 597.} {scriptsizecolor{muted} Y1 | review | f06b67da}
\par
Java manba dasturi .class baytkodiga kompilyatsiya qilindi. Dastur boshqa operatsion tizimda ishlashi uchun qaysi komponent kerak bo'lishi mumkin?
\begin{enumerate}[label=\Alph*),leftmargin=9mm]
\item Kompyuterga o'sha dastur kodining nusxasini printer orqali chiqarish
\item Kompyuterning asosiy plata drayverini fayl nomi bilan almashtirish
\item Kompyuterda tegishli Java Virtual Machine bajarish muhiti bo'lishi
\item Kompyuterdagi barcha .class fayllarni oddiy rastrli rasmga aylantirish
\end{enumerate}
\begin{tcolorbox}[enhanced,breakable,colback=teal!4,colframe=teal!60!black,boxrule=0.45pt,arc=1.2mm,left=2.7mm,right=2.7mm,top=1.8mm,bottom=1.8mm,before skip=2mm,after skip=2mm]
{\bfseries\sffamily\color{teal}To'g'ri javob:} \textbf{C) Kompyuterda tegishli Java Virtual Machine bajarish muhiti bo'lishi}
\par\smallskip
{\bfseries Izoh:} Darslik Java .java kodidan .class baytkodiga kompilyatsiya va JVM yordamida bajarilishni izohlaydi. JVM ning mos muhiti kerak bo'ladi. [Manba: Cambridge+ 10-11-sinf, PDF 46-bet.]
\end{tcolorbox}
\vspace{2mm}\hrule height 0.2pt\vspace{3mm}

% DB_ID f66eabae-a856-0cbd-43a2-849e15704a40
\needspace{5\baselineskip}
{\bfseries\sffamily\color{ink} 598.} {scriptsizecolor{muted} Y1 | review | f66eabae}
\par
Foydalanuvchi hisobot yozmoqda, operatsion tizim esa faylni diskka joylashtirmoqda. Dastur turlari qaysi tartibda?
\begin{enumerate}[label=\Alph*),leftmargin=9mm]
\item tizimli dastur; amaliy dastur
\item ikkalasi ham utilita
\item amaliy dastur; tizimli dastur
\item ikkalasi ham drayver
\end{enumerate}
\begin{tcolorbox}[enhanced,breakable,colback=teal!4,colframe=teal!60!black,boxrule=0.45pt,arc=1.2mm,left=2.7mm,right=2.7mm,top=1.8mm,bottom=1.8mm,before skip=2mm,after skip=2mm]
{\bfseries\sffamily\color{teal}To'g'ri javob:} \textbf{C) amaliy dastur; tizimli dastur}
\par\smallskip
{\bfseries Izoh:} amaliy dastur; tizimli dastur. Matn yozish foydalanuvchi vazifasi, fayl tizimini boshqarish OT vazifasidir. Manba: CAM9, PDF 11-15; Qo'shimcha sahifa raqami (qaysi PDFga tegishli ekani aniqlanmagan; darslik sahifasi deb tasdiqlanmagan): 54
\end{tcolorbox}
\vspace{2mm}\hrule height 0.2pt\vspace{3mm}

% DB_ID f84e20d1-032c-6ebd-7558-14617a4ee100
\needspace{5\baselineskip}
{\bfseries\sffamily\color{ink} 599.} {scriptsizecolor{muted} Y1 | review | f84e20d1}
\par
Kompaniya fayllar to'plamini siqilgan arxivga joyladi. Nima uchun bu har doim xavfsiz zaxira nusxa tayyorlandi degani emas?
\begin{enumerate}[label=\Alph*),leftmargin=9mm]
\item Arxivlash odatda dasturiy ta'minotni qo'shadi, zaxiralash esa faqat siqishni bildiradi
\item Arxivlash hajmni kamaytiradi; zaxiralash esa tiklash uchun mustaqil nusxa saqlaydi
\item Arxivlash faqat fayl nomlarini yangilaydi; zaxiralash esa faqat xatolarni to'g'rilaydi
\item Arxivlash faqat printerdan foydalanadi; zaxiralash esa faqat disk rasmini o'zgartiradi
\end{enumerate}
\begin{tcolorbox}[enhanced,breakable,colback=teal!4,colframe=teal!60!black,boxrule=0.45pt,arc=1.2mm,left=2.7mm,right=2.7mm,top=1.8mm,bottom=1.8mm,before skip=2mm,after skip=2mm]
{\bfseries\sffamily\color{teal}To'g'ri javob:} \textbf{B) Arxivlash hajmni kamaytiradi; zaxiralash esa tiklash uchun mustaqil nusxa saqlaydi}
\par\smallskip
{\bfseries Izoh:} Ma'lumot siqish ma'lumotning hajmini kamaytirishga qaratilgan; zaxiralash esa yo'qolishdan keyin tiklash uchun nusxani saqlashni nazarda tutadi. [Manba: Cambridge+ 10-11-sinf, PDF 58-bet.]
\end{tcolorbox}
\vspace{2mm}\hrule height 0.2pt\vspace{3mm}

% DB_ID 2b627f67-ab90-e289-ddeb-e4be53fa0694
\needspace{5\baselineskip}
{\bfseries\sffamily\color{ink} 600.} {scriptsizecolor{muted} Y1 | review | 2b627f67}
\par
Xodim bir dasturda hisobot matnini tahrir qildi, boshqa dasturda esa fayllarni siqib arxivga joyladi. Bu dasturlar qaysi guruhlarga kiradi?
\begin{enumerate}[label=\Alph*),leftmargin=9mm]
\item Matn muharriri tizimli dastur, faylni siqish vositasi esa alohida operatsion tizim hisoblanadi
\item Matn muharriri amaliy dastur, faylni siqish vositasi esa xizmat ko'rsatuvchi utilita hisoblanadi
\item Matn muharriri qurilma drayveri, faylni siqish vositasi esa kompyuterning bosh platasi hisoblanadi
\item Matn muharriri faqat kompilyator, faylni siqish vositasi esa faqat internet brauzeridir
\end{enumerate}
\begin{tcolorbox}[enhanced,breakable,colback=teal!4,colframe=teal!60!black,boxrule=0.45pt,arc=1.2mm,left=2.7mm,right=2.7mm,top=1.8mm,bottom=1.8mm,before skip=2mm,after skip=2mm]
{\bfseries\sffamily\color{teal}To'g'ri javob:} \textbf{B) Matn muharriri amaliy dastur, faylni siqish vositasi esa xizmat ko'rsatuvchi utilita hisoblanadi}
\par\smallskip
{\bfseries Izoh:} Matn muharriri foydalanuvchining ishini bevosita bajaradigan amaliy dastur. Fayllarni siqish esa utilita dasturining xizmat vazifasidir. [Manba: Cambridge+ 10-11-sinf, PDF 58; ICT 10-sinf, PDF 28-29.]
\end{tcolorbox}
\vspace{2mm}\hrule height 0.2pt\vspace{3mm}

% DB_ID 30f33d38-b242-91e9-10cf-a0d144b0970d
\needspace{5\baselineskip}
{\bfseries\sffamily\color{ink} 601.} {scriptsizecolor{muted} Y1 | review | 30f33d38}
\par
Dastur o'rnatish, yangilash, litsenziya va moslik mavzusida: I. moslik - dastur va ro'yxatdan o'tgan komponentlarini nazoratli olib tashlash. II. installation - qurilma, OT, dastur va formatlarning birgalikda ishlay olish xususiyati. Mulohazalarni baholang.
\begin{enumerate}[label=\Alph*),leftmargin=9mm]
\item faqat I
\item I ham, II ham emas
\item faqat II
\item I va II
\end{enumerate}
\begin{tcolorbox}[enhanced,breakable,colback=teal!4,colframe=teal!60!black,boxrule=0.45pt,arc=1.2mm,left=2.7mm,right=2.7mm,top=1.8mm,bottom=1.8mm,before skip=2mm,after skip=2mm]
{\bfseries\sffamily\color{teal}To'g'ri javob:} \textbf{B) I ham, II ham emas}
\par\smallskip
{\bfseries Izoh:} I ham, II ham emas. Birinchi va ikkinchi mulohazaning ta'riflari alohida tekshirildi; natija: I ham, II ham emas. Manba: CAM10/11, PDF 46-61; Qo'shimcha sahifa raqami (qaysi PDFga tegishli ekani aniqlanmagan; darslik sahifasi deb tasdiqlanmagan): 105
\end{tcolorbox}
\vspace{2mm}\hrule height 0.2pt\vspace{3mm}

% DB_ID 53fcbba8-d67d-5618-80b5-5633aab6f398
\needspace{5\baselineskip}
{\bfseries\sffamily\color{ink} 602.} {scriptsizecolor{muted} Y1 | review | 53fcbba8}
\par
Dasturiy ta'minotning tasnifi mavzusida: I. amaliy dastur - tizimga xizmat ko'rsatuvchi yordamchi dastur. II. tizimli dastur - kompyuter resurslarini boshqarib, boshqa dasturlarga muhit yaratuvchi dastur. Mulohazalarni baholang.
\begin{enumerate}[label=\Alph*),leftmargin=9mm]
\item faqat I
\item I va II
\item I ham, II ham emas
\item faqat II
\end{enumerate}
\begin{tcolorbox}[enhanced,breakable,colback=teal!4,colframe=teal!60!black,boxrule=0.45pt,arc=1.2mm,left=2.7mm,right=2.7mm,top=1.8mm,bottom=1.8mm,before skip=2mm,after skip=2mm]
{\bfseries\sffamily\color{teal}To'g'ri javob:} \textbf{D) faqat II}
\par\smallskip
{\bfseries Izoh:} faqat II. Birinchi va ikkinchi mulohazaning ta'riflari alohida tekshirildi; natija: faqat II. Manba: CAM10/11, PDF 44-58; Qo'shimcha sahifa raqami (qaysi PDFga tegishli ekani aniqlanmagan; darslik sahifasi deb tasdiqlanmagan): 51
\end{tcolorbox}
\vspace{2mm}\hrule height 0.2pt\vspace{3mm}

% DB_ID 598f9966-8162-0d9e-1d19-00023b6d4645
\needspace{5\baselineskip}
{\bfseries\sffamily\color{ink} 603.} {scriptsizecolor{muted} Y1 | review | 598f9966}
\par
Drayver, utilita, firmware va translatorlar mavzusida: I. utilita - qurilma ichidagi past darajali boshqaruv dasturi. II. drayver - OT bilan aniq qurilma modeli o'rtasida buyruqlarni moslashtiruvchi dastur. Mulohazalarni baholang.
\begin{enumerate}[label=\Alph*),leftmargin=9mm]
\item faqat I
\item faqat II
\item I va II
\item I ham, II ham emas
\end{enumerate}
\begin{tcolorbox}[enhanced,breakable,colback=teal!4,colframe=teal!60!black,boxrule=0.45pt,arc=1.2mm,left=2.7mm,right=2.7mm,top=1.8mm,bottom=1.8mm,before skip=2mm,after skip=2mm]
{\bfseries\sffamily\color{teal}To'g'ri javob:} \textbf{B) faqat II}
\par\smallskip
{\bfseries Izoh:} faqat II. Birinchi va ikkinchi mulohazaning ta'riflari alohida tekshirildi; natija: faqat II. Manba: CAM10/11, PDF 44-49; Qo'shimcha sahifa raqami (qaysi PDFga tegishli ekani aniqlanmagan; darslik sahifasi deb tasdiqlanmagan): 71
\end{tcolorbox}
\vspace{2mm}\hrule height 0.2pt\vspace{3mm}

% DB_ID 5f4e93ad-55d3-5cdd-7341-3a37281ebe24
\needspace{5\baselineskip}
{\bfseries\sffamily\color{ink} 604.} {scriptsizecolor{muted} Y1 | review | 5f4e93ad}
\par
Amaliy dasturlar va vazifaga mos tanlov mavzusida: I. DTP - formulalar asosida hisoblash, model va diagramma yaratish dasturi. II. matn protsessori - professional sahifa maketi va nashr tayyorlash dasturi. Mulohazalarni baholang.
\begin{enumerate}[label=\Alph*),leftmargin=9mm]
\item faqat I
\item I ham, II ham emas
\item faqat II
\item I va II
\end{enumerate}
\begin{tcolorbox}[enhanced,breakable,colback=teal!4,colframe=teal!60!black,boxrule=0.45pt,arc=1.2mm,left=2.7mm,right=2.7mm,top=1.8mm,bottom=1.8mm,before skip=2mm,after skip=2mm]
{\bfseries\sffamily\color{teal}To'g'ri javob:} \textbf{B) I ham, II ham emas}
\par\smallskip
{\bfseries Izoh:} I ham, II ham emas. Birinchi va ikkinchi mulohazaning ta'riflari alohida tekshirildi; natija: I ham, II ham emas. Manba: CAM5, PDF 45-138; Qo'shimcha sahifa raqami (qaysi PDFga tegishli ekani aniqlanmagan; darslik sahifasi deb tasdiqlanmagan): 80
\end{tcolorbox}
\vspace{2mm}\hrule height 0.2pt\vspace{3mm}

% DB_ID 6802df52-5e65-c879-6582-be89c27a9d3c
\needspace{5\baselineskip}
{\bfseries\sffamily\color{ink} 605.} {scriptsizecolor{muted} Y1 | review | 6802df52}
\par
Magnit qattiq diskdagi fayl bo'laklari bo'sh sektorlar orasiga tarqalgan. Ma'lumotlarni o'qishda ortiqcha joy qidirish kerak bo'lmoqda. Qaysi xizmat vositasi muammoni kamaytirishi mumkin?
\begin{enumerate}[label=\Alph*),leftmargin=9mm]
\item Vaqtinchalik fayllarni o'chirib, diskdagi bo'sh joy hajmini umumiy hisobda ko'paytirish
\item Fayllarni siqilgan arxivga yig'ib, barcha materiallarning saqlash hajmini kamaytirish
\item Diskdagi fayl tizimini qayta formatlab, uning barcha mavjud ma'lumotlarini yangidan yozish
\item Fayl bo'laklari joylashuvini tartibga soluvchi defragmentatsiya vositasidan foydalanish
\end{enumerate}
\begin{tcolorbox}[enhanced,breakable,colback=teal!4,colframe=teal!60!black,boxrule=0.45pt,arc=1.2mm,left=2.7mm,right=2.7mm,top=1.8mm,bottom=1.8mm,before skip=2mm,after skip=2mm]
{\bfseries\sffamily\color{teal}To'g'ri javob:} \textbf{D) Fayl bo'laklari joylashuvini tartibga soluvchi defragmentatsiya vositasidan foydalanish}
\par\smallskip
{\bfseries Izoh:} Cambridge+ darslikdagi HDD misolida fragmentatsiya va bo'laklarni tartiblash uchun defragmentatsiya yondashuvi tasvirlangan. [Manba: Cambridge+ 10-11-sinf, PDF 59-bet.]
\end{tcolorbox}
\vspace{2mm}\hrule height 0.2pt\vspace{3mm}

% DB_ID 80f4dcd4-0cdc-6018-5e6c-4c4c8cf5e9c8
\needspace{5\baselineskip}
{\bfseries\sffamily\color{ink} 606.} {scriptsizecolor{muted} Y1 | review | 80f4dcd4}
\par
Amaliy dasturlar va vazifaga mos tanlov mavzusida: I. elektron jadval - tuzilmali ma'lumotlar bazasini yaratish, so'rash va boshqarish tizimi. II. matn protsessori - tuzilmali matnli hujjat yaratish va formatlash dasturi. Mulohazalarni baholang.
\begin{enumerate}[label=\Alph*),leftmargin=9mm]
\item faqat I
\item I va II
\item I ham, II ham emas
\item faqat II
\end{enumerate}
\begin{tcolorbox}[enhanced,breakable,colback=teal!4,colframe=teal!60!black,boxrule=0.45pt,arc=1.2mm,left=2.7mm,right=2.7mm,top=1.8mm,bottom=1.8mm,before skip=2mm,after skip=2mm]
{\bfseries\sffamily\color{teal}To'g'ri javob:} \textbf{D) faqat II}
\par\smallskip
{\bfseries Izoh:} faqat II. Birinchi va ikkinchi mulohazaning ta'riflari alohida tekshirildi; natija: faqat II. Manba: CAM6, PDF 34-205; Qo'shimcha sahifa raqami (qaysi PDFga tegishli ekani aniqlanmagan; darslik sahifasi deb tasdiqlanmagan): 78
\end{tcolorbox}
\vspace{2mm}\hrule height 0.2pt\vspace{3mm}

% DB_ID 8e2a98a6-8f75-028c-9d68-961004e9b047
\needspace{5\baselineskip}
{\bfseries\sffamily\color{ink} 607.} {scriptsizecolor{muted} Y1 | review | 8e2a98a6}
\par
Dastur o'rnatish, yangilash, litsenziya va moslik mavzusida: I. uninstall - mavjud versiyaga tuzatish yoki kichik yaxshilanish o'rnatish. II. installation - dastur komponentlarini tizimda ishlashga tayyorlab joylashtirish. Mulohazalarni baholang.
\begin{enumerate}[label=\Alph*),leftmargin=9mm]
\item faqat I
\item I va II
\item I ham, II ham emas
\item faqat II
\end{enumerate}
\begin{tcolorbox}[enhanced,breakable,colback=teal!4,colframe=teal!60!black,boxrule=0.45pt,arc=1.2mm,left=2.7mm,right=2.7mm,top=1.8mm,bottom=1.8mm,before skip=2mm,after skip=2mm]
{\bfseries\sffamily\color{teal}To'g'ri javob:} \textbf{D) faqat II}
\par\smallskip
{\bfseries Izoh:} faqat II. Birinchi va ikkinchi mulohazaning ta'riflari alohida tekshirildi; natija: faqat II. Manba: CAM8, PDF 99-110; Qo'shimcha sahifa raqami (qaysi PDFga tegishli ekani aniqlanmagan; darslik sahifasi deb tasdiqlanmagan): 103
\end{tcolorbox}
\vspace{2mm}\hrule height 0.2pt\vspace{3mm}

% DB_ID a73e0776-c2e1-770a-03fb-3f85ff913a8d
\needspace{5\baselineskip}
{\bfseries\sffamily\color{ink} 608.} {scriptsizecolor{muted} Y1 | review | a73e0776}
\par
Klinikaga bemorlarni qabul qilishda noyob ichki qoidalari va maxsus hisobot shakllari kerak. Qaysi qaror dastur tanlashda eng asosli?
\begin{enumerate}[label=\Alph*),leftmargin=9mm]
\item Dastur narxi yuqori bo'lsa, talablarni tekshirmasdan eng qimmatini olish
\item Dastur interfeysi yorqin bo'lsa, uni har qanday tashkilot uchun tayyor deb olish
\item Talablarni yozib, tayyor yechim va buyurtma yechimning mosligini baholash
\item Talablar ko'p bo'lsa, foydalanuvchi ishini faqat qog'ozga qaytarib qo'yish
\end{enumerate}
\begin{tcolorbox}[enhanced,breakable,colback=teal!4,colframe=teal!60!black,boxrule=0.45pt,arc=1.2mm,left=2.7mm,right=2.7mm,top=1.8mm,bottom=1.8mm,before skip=2mm,after skip=2mm]
{\bfseries\sffamily\color{teal}To'g'ri javob:} \textbf{C) Talablarni yozib, tayyor yechim va buyurtma yechimning mosligini baholash}
\par\smallskip
{\bfseries Izoh:} Darslik tayyor dastur keng auditoriyaga mo'ljallanishini, buyurtma dastur esa aniq mijoz talabi uchun yaratilishini farqlaydi. Qaror talablar va xarajatlarga tayangan holda beriladi. [Manba: Cambridge+ 10-11-sinf, PDF 60-bet.]
\end{tcolorbox}
\vspace{2mm}\hrule height 0.2pt\vspace{3mm}

% DB_ID a9f84bf0-9abb-176a-8b31-9776082900e5
\needspace{5\baselineskip}
{\bfseries\sffamily\color{ink} 609.} {scriptsizecolor{muted} Y1 | review | a9f84bf0}
\par
Java dasturining .java manba kodi .class fayliga aylantirildi, so'ng dastur JVM orqali bajarildi. Qaysi izoh bu bosqichlarni to'g'ri tavsiflaydi?
\begin{enumerate}[label=\Alph*),leftmargin=9mm]
\item Java manba kodi dastlab operatsion tizim tomonidan bevosita protsessorning mashina kodiga aylantiriladi
\item Java uchun .class fayl kompilyatordan emas, faqat oddiy matn muharriridan olinadi
\item Java dasturi oraliq kod yaratmasdan, faqat manba matnini bevosita interpretatsiya qiladi
\item Kompilyator baytkod yaratadi, JVM esa dastur bajarilayotganda uni talqin qiladi
\end{enumerate}
\begin{tcolorbox}[enhanced,breakable,colback=teal!4,colframe=teal!60!black,boxrule=0.45pt,arc=1.2mm,left=2.7mm,right=2.7mm,top=1.8mm,bottom=1.8mm,before skip=2mm,after skip=2mm]
{\bfseries\sffamily\color{teal}To'g'ri javob:} \textbf{D) Kompilyator baytkod yaratadi, JVM esa dastur bajarilayotganda uni talqin qiladi}
\par\smallskip
{\bfseries Izoh:} Java manba kodi kompilyatsiya orqali baytkodga (.class) o'tkaziladi; JVM bajarish davomida ushbu oraliq kodni talqin qiladi. Manba: Cambridge+ 10-11-sinf, PDF 46
\end{tcolorbox}
\vspace{2mm}\hrule height 0.2pt\vspace{3mm}
\clearpage\section*{Qamrov va sifat nazorati}
\addcontentsline{toc}{section}{Qamrov va sifat nazorati}
Manba bazaning 609 ta amaldagi savoli kiritilgan (67 published va 542 review).
Javoblar va izohlar savollar bilan aynan identifikator bo'yicha bog'langan.
200 ta archived savol asosiy kitob tarkibidan chiqarilgan.
Tuzilishning avtomatik tekshiruvi ekspert tomonidan faktik tasdiqni anglatmaydi.
\par Keyingi bosqich: Microsoft Office (Word, Excel, PowerPoint).
\end{document}